% Master document for the rendered derivations.  The \input fragments are
% GENERATED by deriv/render.wl (see deriv/README.md, "Rendered output"); this file is not.
\documentclass[10pt,a4paper]{article}
\usepackage[T1]{fontenc}
\usepackage[utf8]{inputenc}
\usepackage{lmodern}
\usepackage[margin=2cm]{geometry}
\usepackage{amsmath,amssymb}
\usepackage{array,longtable,booktabs}
\usepackage[hidelinks]{hyperref}
% no build date or trailer id in the PDF, so that a rebuild of unchanged sources is identical
\ifdefined\pdfinfoomitdate \pdfinfoomitdate=1 \pdftrailerid{}\fi
\allowdisplaybreaks
\setcounter{MaxMatrixCols}{20}
\setcounter{tocdepth}{2}
\setlength{\parindent}{0pt}
\setlength{\parskip}{0.4em}

\title{Fluidum: the derivations of the generated matrices}
\author{}
\date{}

\begin{document}
\maketitle
\thispagestyle{empty}

Each section is written by the derivation script it describes, from the matrices, sources
and check results that script computed:
\begin{itemize}\setlength{\itemsep}{0pt}
\item \texttt{deriv/derive\_2p1d\_viscous.wls}
\item \texttt{deriv/derive\_hq\_2p1d.wls}
\item \texttt{deriv/derive\_two\_way.wls}
\item \texttt{deriv/derive\_1d\_matrices.wls}
\item \texttt{deriv/derive\_upstream.wls}
\item \texttt{deriv/derive\_drivers.wls}
\end{itemize}
Regenerate with \texttt{deriv/render\_all.sh}.

Conventions: boost-invariant Milne coordinates; every system is written source on the left,
$A_t\,\partial_\tau\phi + A_x\,\partial_x\phi + A_y\,\partial_y\phi + S = 0$; only nonzero
entries are listed; $u^\tau = \sqrt{1 + u_\perp^2}$ is written for the root wherever it occurs.

\tableofcontents
\clearpage

% GENERATED by deriv/render.wl from the derivation script.  Do not edit.

\section{Viscous hydrodynamics: Israel--Stewart with \texorpdfstring{$\delta_{\pi\pi}$}{\unichar{"03B4}\unichar{"03C0}\unichar{"03C0}}}

Script: \texttt{deriv/derive\_2p1d\_viscous.wls}

Writes: \texttt{src/Matrix/viscous\_generated.jl}

\subsection{Setup}

Milne coordinates, boost invariant. Two charts:

\begin{gather*}
x^{\mu} = ({\tau}, {r}, {\phi}, {\eta}),\qquad g_{\mu\nu} = \mathrm{diag}(-1, 1, {r}^2, {\tau}^2) \\
x^{\mu} = ({\tau}, {x}, {y}, {\eta}),\qquad g_{\mu\nu} = \mathrm{diag}(-1, 1, 1, {\tau}^2)
\end{gather*}
\begin{gather*}
T^{\mu\nu} = (e + p + \Pi)\,u^{\mu}u^{\nu} + (p + \Pi)\,g^{\mu\nu} + \pi^{\mu\nu},\qquad e = T\,p'(T) - p(T) \\
\nabla_{\mu}T^{\mu\nu} = 0 \\
\tau_{\pi}\,\Delta^{\mu\nu}_{\alpha\beta}\,D\pi^{\alpha\beta} + \pi^{\mu\nu} + 2\eta\,\sigma^{\mu\nu} + \delta_{\pi\pi}\,\pi^{\mu\nu}\theta = 0 \\
\tau_{\Pi}\,D\Pi + \Pi + \zeta\,\theta = 0
\end{gather*}
with

\begin{gather*}
D = u^{\mu}\nabla_{\mu},\qquad \theta = \nabla_{\mu}u^{\mu},\qquad \Delta^{\mu\nu} = g^{\mu\nu} + u^{\mu}u^{\nu} \\
\sigma^{\mu\nu} = \Delta^{\mu\nu}_{\alpha\beta}\nabla^{\alpha}u^{\beta},\qquad \Delta^{\mu\nu}_{\alpha\beta} = \tfrac{1}{2}\bigl(\Delta^{\mu}_{\alpha}\Delta^{\nu}_{\beta} + \Delta^{\mu}_{\beta}\Delta^{\nu}_{\alpha}\bigr) - \tfrac{1}{3}\Delta^{\mu\nu}\Delta_{\alpha\beta}
\end{gather*}
$\eta$, $\tau_{\pi}$, $\zeta$, $\tau_{\Pi}$ are not differentiated. Quasi-linear form, source on the left:

\begin{gather*}
A_t\,\partial_{\tau}\phi + A_r\,\partial_{r}\phi + S = 0\quad\text{(1+1D)} \\
A_t\,\partial_{\tau}\phi + A_x\,\partial_{x}\phi + A_y\,\partial_{y}\phi + S = 0\quad\text{(2+1D)}
\end{gather*}
The $\delta_{\pi\pi}$ term is linear, so each matrix and the source split as

\[
A = A^{\mathrm{IS}} + \delta_{\pi\pi}\,A^{(\delta)},\qquad S = S^{\mathrm{IS}} + \delta_{\pi\pi}\,S^{(\delta)}
\]
\subsection{1+1D cylindrical}

Fields, with $\pi^{rr}$ and the $\tau$ row of $\pi^{\mu\nu}$ fixed by $u_{\mu}\pi^{\mu\nu} = 0$, $\pi^{\mu}{}_{\mu} = 0$:

\begin{gather*}
\phi = ({T}, {u^{r}}, {\pi^{\phi}{}_{\phi}}, {\pi^{\eta}{}_{\eta}}, {\Pi}),\qquad \pi^{\phi}{}_{\phi} = r^{2}\pi^{\phi\phi},\quad \pi^{\eta}{}_{\eta} = \tau^{2}\pi^{\eta\eta} \\
\pi^{rr} = -\left(\left(u^{\tau}\right)^2 ({\pi^{\eta}{}_{\eta}}+{\pi^{\phi}{}_{\phi}})\right),\qquad u^{\tau} = \sqrt{\left(u^{r}\right)^2+1}
\end{gather*}
Rows: $\nabla_{\mu}T^{\mu\tau}$, $\nabla_{\mu}T^{\mu r}$, $E^{\phi\phi}$, $E^{\eta\eta}$, $E_{\Pi}$, where $E^{\mu\nu}$ and $E_{\Pi}$ are the left-hand sides of the shear and bulk equations.

\subsubsection{IS part}

\begin{align*}
(A^{\mathrm{IS}}_t)_{11} &= \left(u^{r}\right)^2 ({p''(T)} {T}+{p'(T)})+{p''(T)} {T} \\
(A^{\mathrm{IS}}_t)_{12} &= 2 {u^{r}} ({p'(T)} {T}+{\Pi}-{\pi^{\eta}{}_{\eta}}-{\pi^{\phi}{}_{\phi}}) \\
(A^{\mathrm{IS}}_t)_{13} &= -\left(u^{r}\right)^2 \\
(A^{\mathrm{IS}}_t)_{14} &= -\left(u^{r}\right)^2 \\
(A^{\mathrm{IS}}_t)_{15} &= \left(u^{r}\right)^2 \\
(A^{\mathrm{IS}}_t)_{21} &= {u^{r}} {u^{\tau}} ({p''(T)} {T}+{p'(T)}) \\
(A^{\mathrm{IS}}_t)_{22} &= \frac{\left(2 \left(u^{r}\right)^2+1\right) ({p'(T)} {T}+{\Pi}-{\pi^{\eta}{}_{\eta}}-{\pi^{\phi}{}_{\phi}})}{{u^{\tau}}} \\
(A^{\mathrm{IS}}_t)_{23} &= -{u^{r}} {u^{\tau}} \\
(A^{\mathrm{IS}}_t)_{24} &= -{u^{r}} {u^{\tau}} \\
(A^{\mathrm{IS}}_t)_{25} &= {u^{r}} {u^{\tau}} \\
(A^{\mathrm{IS}}_t)_{32} &= -\frac{2 {u^{r}} {\eta}}{3 {r}^2 {u^{\tau}}} \\
(A^{\mathrm{IS}}_t)_{33} &= \frac{{\tau_{\pi}} {u^{\tau}}}{{r}^2} \\
(A^{\mathrm{IS}}_t)_{42} &= -\frac{2 {u^{r}} {\eta}}{3 {\tau}^2 {u^{\tau}}} \\
(A^{\mathrm{IS}}_t)_{44} &= \frac{{\tau_{\pi}} {u^{\tau}}}{{\tau}^2} \\
(A^{\mathrm{IS}}_t)_{52} &= \frac{{u^{r}} {\zeta}}{{u^{\tau}}} \\
(A^{\mathrm{IS}}_t)_{55} &= {\tau_{\Pi}} {u^{\tau}}
\end{align*}
\begin{align*}
(A^{\mathrm{IS}}_r)_{11} &= {u^{r}} {u^{\tau}} ({p''(T)} {T}+{p'(T)}) \\
(A^{\mathrm{IS}}_r)_{12} &= \frac{\left(2 \left(u^{r}\right)^2+1\right) ({p'(T)} {T}+{\Pi}-{\pi^{\eta}{}_{\eta}}-{\pi^{\phi}{}_{\phi}})}{{u^{\tau}}} \\
(A^{\mathrm{IS}}_r)_{13} &= -{u^{r}} {u^{\tau}} \\
(A^{\mathrm{IS}}_r)_{14} &= -{u^{r}} {u^{\tau}} \\
(A^{\mathrm{IS}}_r)_{15} &= {u^{r}} {u^{\tau}} \\
(A^{\mathrm{IS}}_r)_{21} &= {p''(T)} \left(u^{r}\right)^2 {T}+{p'(T)} \left(u^{r}\right)^2+{p'(T)} \\
(A^{\mathrm{IS}}_r)_{22} &= 2 {u^{r}} ({p'(T)} {T}+{\Pi}-{\pi^{\eta}{}_{\eta}}-{\pi^{\phi}{}_{\phi}}) \\
(A^{\mathrm{IS}}_r)_{23} &= -\left(u^{r}\right)^2-1 \\
(A^{\mathrm{IS}}_r)_{24} &= -\left(u^{r}\right)^2-1 \\
(A^{\mathrm{IS}}_r)_{25} &= \left(u^{\tau}\right)^2 \\
(A^{\mathrm{IS}}_r)_{32} &= -\frac{2 {\eta}}{3 {r}^2} \\
(A^{\mathrm{IS}}_r)_{33} &= \frac{{\tau_{\pi}} {u^{r}}}{{r}^2} \\
(A^{\mathrm{IS}}_r)_{42} &= -\frac{2 {\eta}}{3 {\tau}^2} \\
(A^{\mathrm{IS}}_r)_{44} &= \frac{{\tau_{\pi}} {u^{r}}}{{\tau}^2} \\
(A^{\mathrm{IS}}_r)_{52} &= {\zeta} \\
(A^{\mathrm{IS}}_r)_{55} &= {\tau_{\Pi}} {u^{r}}
\end{align*}
\begin{align*}
S^{\mathrm{IS}}_{1} &= \frac{1}{{r} {\tau} {u^{\tau}}}\Bigl(\left(u^{\tau}\right)^2 ({p'(T)} {T}+{\Pi}) ({r} {u^{\tau}}+{\tau} {u^{r}}) - {\pi^{\phi}{}_{\phi}} {u^{r}} \left({r} {u^{r}} {u^{\tau}}+{\tau} \left(u^{r}\right)^2+{\tau}\right)\\
  &\quad - {\pi^{\eta}{}_{\eta}} \left({r} \left(\left(u^{r}\right)^2-1\right) {u^{\tau}}+{\tau} \left(\left(u^{r}\right)^3+{u^{r}}\right)\right)\Bigr) \\
S^{\mathrm{IS}}_{2} &= \frac{1}{{r} {\tau} {u^{\tau}}}\Bigl(- {\pi^{\eta}{}_{\eta}} \left(u^{\tau}\right)^2 ({r} {u^{r}}+{\tau} {u^{\tau}}) + {u^{r}} ({p'(T)} {T}+{\Pi}) \left({r} \left(u^{r}\right)^2+{r}+{\tau} {u^{r}} {u^{\tau}}\right)\\
  &\quad - {\pi^{\phi}{}_{\phi}} \left({r} \left(\left(u^{r}\right)^3+{u^{r}}\right)+{\tau} \left(\left(u^{r}\right)^2+2\right) {u^{\tau}}\right)\Bigr) \\
S^{\mathrm{IS}}_{3} &= \frac{{\pi^{\phi}{}_{\phi}}}{{r}^2}-\frac{2 {\eta} \left({r} \left(u^{r}\right)^2+{r}-2 {\tau} {u^{r}} {u^{\tau}}\right)}{3 {r}^3 {\tau} {u^{\tau}}} \\
S^{\mathrm{IS}}_{4} &= \frac{3 {\pi^{\eta}{}_{\eta}} {\tau}-\frac{2 {\tau} {u^{r}} {\eta}}{{r}}+4 {u^{\tau}} {\eta}}{3 {\tau}^3} \\
S^{\mathrm{IS}}_{5} &= {\Pi}+\frac{{u^{r}} {\zeta}}{{r}}+\frac{{u^{\tau}} {\zeta}}{{\tau}}
\end{align*}
\subsubsection{\texorpdfstring{$\delta_{\pi\pi}$}{\unichar{"03B4}\unichar{"03C0}\unichar{"03C0}} part}

\begin{align*}
(A^{(\delta)}_t)_{32} &= \frac{{\pi^{\phi}{}_{\phi}} {u^{r}}}{{r}^2 {u^{\tau}}} \\
(A^{(\delta)}_t)_{42} &= \frac{{\pi^{\eta}{}_{\eta}} {u^{r}}}{{\tau}^2 {u^{\tau}}}
\end{align*}
\begin{align*}
(A^{(\delta)}_r)_{32} &= \frac{{\pi^{\phi}{}_{\phi}}}{{r}^2} \\
(A^{(\delta)}_r)_{42} &= \frac{{\pi^{\eta}{}_{\eta}}}{{\tau}^2}
\end{align*}
\begin{align*}
S^{(\delta)}_{3} &= \frac{{\pi^{\phi}{}_{\phi}} ({r} {u^{\tau}}+{\tau} {u^{r}})}{{r}^3 {\tau}} \\
S^{(\delta)}_{4} &= \frac{{\pi^{\eta}{}_{\eta}} ({r} {u^{\tau}}+{\tau} {u^{r}})}{{r} {\tau}^3}
\end{align*}
\subsection{2+1D Cartesian}

Fields, with $\pi^{xx}$ and the $\tau$ row fixed by the same constraints:

\begin{gather*}
\phi = ({T}, {u^{x}}, {u^{y}}, {\pi^{yy}}, {\pi^{\eta}{}_{\eta}}, {\pi^{xy}}, {\Pi}) \\
\pi^{xx} = -\frac{{\pi^{\eta}{}_{\eta}} \left(u^{\tau}\right)^2-2 {\pi^{xy}} {u^{x}} {u^{y}}+{\pi^{yy}} \left(u^{x}\right)^2+{\pi^{yy}}}{\left(u^{y}\right)^2+1},\qquad u^{\tau} = \sqrt{\left(u^{x}\right)^2+\left(u^{y}\right)^2+1}
\end{gather*}
Rows: $\nabla_{\mu}T^{\mu\tau}$, $\nabla_{\mu}T^{\mu x}$, $\nabla_{\mu}T^{\mu y}$, $E^{yy}$, $E^{\eta\eta}$, $E^{xy}$, $E_{\Pi}$.

\subsubsection{IS part}

\begin{align*}
(A^{\mathrm{IS}}_t)_{11} &= \left(u^{\tau}\right)^2 ({p''(T)} {T}+{p'(T)})-{p'(T)} \\
(A^{\mathrm{IS}}_t)_{12} &= \frac{2 \left({p'(T)} {T} {u^{x}} \left(u^{y}\right)^2+{p'(T)} {T} {u^{x}}+{\Pi} {u^{x}} \left(\left(u^{y}\right)^2+1\right)-{\pi^{\eta}{}_{\eta}} {u^{x}}+{\pi^{xy}} {u^{y}}-{\pi^{yy}} {u^{x}}\right)}{\left(u^{y}\right)^2+1} \\
(A^{\mathrm{IS}}_t)_{13} &= \frac{2}{\left(\left(u^{y}\right)^2+1\right)^2}\Bigl({\pi^{xy}} \left({u^{x}}-{u^{x}} \left(u^{y}\right)^2\right)\\
  &\quad + {u^{y}}\Bigl({\pi^{yy}} + {p'(T)} {T} + {\pi^{\eta}{}_{\eta}} \left(u^{x}\right)^2 + {\pi^{yy}} \left(u^{x}\right)^2 + 2 {p'(T)} {T} \left(u^{y}\right)^2 + {p'(T)} {T} \left(u^{y}\right)^4\\
  &\quad + {\Pi} \left(\left(u^{y}\right)^2+1\right)^2\Bigr)\Bigr) \\
(A^{\mathrm{IS}}_t)_{14} &= \frac{\left(u^{y}\right)^2-\left(u^{x}\right)^2}{\left(u^{y}\right)^2+1} \\
(A^{\mathrm{IS}}_t)_{15} &= -\frac{\left(u^{x}\right)^2}{\left(u^{y}\right)^2+1} \\
(A^{\mathrm{IS}}_t)_{16} &= \frac{2 {u^{x}} {u^{y}}}{\left(u^{y}\right)^2+1} \\
(A^{\mathrm{IS}}_t)_{17} &= \left(u^{x}\right)^2+\left(u^{y}\right)^2 \\
(A^{\mathrm{IS}}_t)_{21} &= {u^{\tau}} {u^{x}} ({p''(T)} {T}+{p'(T)}) \\
(A^{\mathrm{IS}}_t)_{22} &= \frac{1}{\left(u^{\tau}\right)^3}\Bigl(\left(u^{\tau}\right)^4 ({p'(T)} {T}+{\Pi}) + \left(u^{\tau}\right)^2 \left(u^{x}\right)^2 ({p'(T)} {T}+{\Pi})\\
  &\quad + \frac{{u^{x}} \left({\pi^{\eta}{}_{\eta}} \left(u^{\tau}\right)^2 {u^{x}}-{\pi^{xy}} {u^{y}} \left(2 \left(u^{x}\right)^2+\left(u^{y}\right)^2+1\right)+{\pi^{yy}} \left(\left(u^{x}\right)^3+{u^{x}}\right)\right)}{\left(u^{y}\right)^2+1}\\
  &\quad - \frac{\left(u^{\tau}\right)^2 \left({\pi^{\eta}{}_{\eta}} \left(3 \left(u^{x}\right)^2+\left(u^{y}\right)^2+1\right)-4 {\pi^{xy}} {u^{x}} {u^{y}}+3 {\pi^{yy}} \left(u^{x}\right)^2+{\pi^{yy}}\right)}{\left(u^{y}\right)^2+1}\Bigr) \\
(A^{\mathrm{IS}}_t)_{23} &= \frac{1}{\left(u^{\tau}\right)^3}\Bigl(\left(u^{\tau}\right)^2 {u^{x}} {u^{y}} ({p'(T)} {T}+{\Pi}) + {u^{y}} \left(\frac{{u^{x}} \left({\pi^{\eta}{}_{\eta}} \left(u^{\tau}\right)^2-2 {\pi^{xy}} {u^{x}} {u^{y}}+{\pi^{yy}} \left(u^{x}\right)^2+{\pi^{yy}}\right)}{\left(u^{y}\right)^2+1}-{\pi^{xy}} {u^{y}}\right)\\
  &\quad + \frac{\left(u^{\tau}\right)^2 \left(2 {u^{x}} {u^{y}} \left({\pi^{\eta}{}_{\eta}} \left(u^{x}\right)^2+{\pi^{yy}} \left(u^{x}\right)^2+{\pi^{yy}}\right)+{\pi^{xy}} \left(\left(\left(u^{y}\right)^2+1\right)^2-2 \left(u^{x}\right)^2 \left(\left(u^{y}\right)^2-1\right)\right)\right)}{\left(\left(u^{y}\right)^2+1\right)^2}\Bigr) \\
(A^{\mathrm{IS}}_t)_{24} &= -\frac{{u^{x}} \left(\left(u^{x}\right)^2+1\right)}{{u^{\tau}} \left(\left(u^{y}\right)^2+1\right)} \\
(A^{\mathrm{IS}}_t)_{25} &= -\frac{{u^{\tau}} {u^{x}}}{\left(u^{y}\right)^2+1} \\
(A^{\mathrm{IS}}_t)_{26} &= \frac{{u^{y}} \left(\frac{2 \left(u^{x}\right)^2}{\left(u^{y}\right)^2+1}+1\right)}{{u^{\tau}}} \\
(A^{\mathrm{IS}}_t)_{27} &= {u^{\tau}} {u^{x}} \\
(A^{\mathrm{IS}}_t)_{31} &= {u^{\tau}} {u^{y}} ({p''(T)} {T}+{p'(T)}) \\
(A^{\mathrm{IS}}_t)_{32} &= \frac{{u^{x}} {u^{y}} \left({p'(T)} {T} \left(u^{\tau}\right)^2+{\Pi} \left(u^{\tau}\right)^2-{\pi^{yy}}\right)+{\pi^{xy}} \left(\left(u^{y}\right)^2+1\right)}{\left(u^{\tau}\right)^3} \\
(A^{\mathrm{IS}}_t)_{33} &= \frac{\left(u^{\tau}\right)^4 ({p'(T)} {T}+{\Pi})+\left(u^{\tau}\right)^2 \left(u^{y}\right)^2 ({p'(T)} {T}+{\Pi})-{u^{y}} ({\pi^{xy}} {u^{x}}+{\pi^{yy}} {u^{y}})+{\pi^{yy}} \left(u^{\tau}\right)^2}{\left(u^{\tau}\right)^3} \\
(A^{\mathrm{IS}}_t)_{34} &= \frac{{u^{y}}}{{u^{\tau}}} \\
(A^{\mathrm{IS}}_t)_{36} &= \frac{{u^{x}}}{{u^{\tau}}} \\
(A^{\mathrm{IS}}_t)_{37} &= {u^{\tau}} {u^{y}} \\
(A^{\mathrm{IS}}_t)_{42} &= -\frac{2 \left(3 {\pi^{xy}} {\tau_{\pi}} {u^{y}} \left(\left(u^{y}\right)^2+1\right)+{u^{x}} \left(\left(u^{y}\right)^2 ({\eta}-3 {\pi^{yy}} {\tau_{\pi}})+{\eta}\right)\right)}{3 {u^{\tau}}} \\
(A^{\mathrm{IS}}_t)_{43} &= \frac{2 {u^{y}} \left(3 {\pi^{xy}} {\tau_{\pi}} {u^{x}} {u^{y}}-3 {\pi^{yy}} {\tau_{\pi}} \left(\left(u^{x}\right)^2+1\right)+{\eta} \left(3 \left(u^{x}\right)^2+2 \left(u^{y}\right)^2+2\right)\right)}{3 {u^{\tau}}} \\
(A^{\mathrm{IS}}_t)_{44} &= {\tau_{\pi}} {u^{\tau}} \\
(A^{\mathrm{IS}}_t)_{52} &= -\frac{2 {u^{x}} {\eta}}{3 {\tau}^2 {u^{\tau}}} \\
(A^{\mathrm{IS}}_t)_{53} &= -\frac{2 {u^{y}} {\eta}}{3 {\tau}^2 {u^{\tau}}} \\
(A^{\mathrm{IS}}_t)_{55} &= \frac{{\tau_{\pi}} {u^{\tau}}}{{\tau}^2} \\
(A^{\mathrm{IS}}_t)_{62} &= \frac{1}{3 {u^{\tau}}}\Bigl(- 3 {\pi^{xy}} {\tau_{\pi}} {u^{x}} \left(2 \left(u^{y}\right)^2+1\right)\\
  &\quad + {u^{y}} \left(3 {\pi^{\eta}{}_{\eta}} {\tau_{\pi}} \left(u^{\tau}\right)^2+3 {\pi^{yy}} \left(2 {\tau_{\pi}} \left(u^{x}\right)^2+{\tau_{\pi}}\right)+{\eta} \left(\left(u^{x}\right)^2+3 \left(u^{y}\right)^2+3\right)\right)\Bigr) \\
(A^{\mathrm{IS}}_t)_{63} &= \frac{1}{3 {u^{\tau}} \left(\left(u^{y}\right)^2+1\right)}\Bigl(- 3 {\pi^{yy}} {\tau_{\pi}} {u^{x}} \left(\left(u^{x}\right)^2+1\right) \left(2 \left(u^{y}\right)^2+1\right) + 3 {\pi^{xy}} {\tau_{\pi}} {u^{y}} \left(\left(2 \left(u^{x}\right)^2-1\right) \left(u^{y}\right)^2-1\right)\\
  &\quad + {u^{x}} \left(\left(\left(u^{y}\right)^2+1\right) {\eta} \left(3 \left(u^{x}\right)^2+\left(u^{y}\right)^2+3\right)-3 {\pi^{\eta}{}_{\eta}} {\tau_{\pi}} \left(u^{\tau}\right)^2 \left(u^{y}\right)^2\right)\Bigr) \\
(A^{\mathrm{IS}}_t)_{66} &= {\tau_{\pi}} {u^{\tau}} \\
(A^{\mathrm{IS}}_t)_{72} &= \frac{{u^{x}} {\zeta}}{{u^{\tau}}} \\
(A^{\mathrm{IS}}_t)_{73} &= \frac{{u^{y}} {\zeta}}{{u^{\tau}}} \\
(A^{\mathrm{IS}}_t)_{77} &= {\tau_{\Pi}} {u^{\tau}}
\end{align*}
\begin{align*}
(A^{\mathrm{IS}}_x)_{11} &= {u^{\tau}} {u^{x}} ({p''(T)} {T}+{p'(T)}) \\
(A^{\mathrm{IS}}_x)_{12} &= \frac{1}{\left(u^{\tau}\right)^3}\Bigl(\left(u^{\tau}\right)^4 ({p'(T)} {T}+{\Pi}) + \left(u^{\tau}\right)^2 \left(u^{x}\right)^2 ({p'(T)} {T}+{\Pi})\\
  &\quad + \frac{{u^{x}} \left({\pi^{\eta}{}_{\eta}} \left(u^{\tau}\right)^2 {u^{x}}-{\pi^{xy}} {u^{y}} \left(2 \left(u^{x}\right)^2+\left(u^{y}\right)^2+1\right)+{\pi^{yy}} \left(\left(u^{x}\right)^3+{u^{x}}\right)\right)}{\left(u^{y}\right)^2+1}\\
  &\quad - \frac{\left(u^{\tau}\right)^2 \left({\pi^{\eta}{}_{\eta}} \left(3 \left(u^{x}\right)^2+\left(u^{y}\right)^2+1\right)-4 {\pi^{xy}} {u^{x}} {u^{y}}+3 {\pi^{yy}} \left(u^{x}\right)^2+{\pi^{yy}}\right)}{\left(u^{y}\right)^2+1}\Bigr) \\
(A^{\mathrm{IS}}_x)_{13} &= \frac{1}{\left(u^{\tau}\right)^3}\Bigl(\left(u^{\tau}\right)^2 {u^{x}} {u^{y}} ({p'(T)} {T}+{\Pi}) + {u^{y}} \left(\frac{{u^{x}} \left({\pi^{\eta}{}_{\eta}} \left(u^{\tau}\right)^2-2 {\pi^{xy}} {u^{x}} {u^{y}}+{\pi^{yy}} \left(u^{x}\right)^2+{\pi^{yy}}\right)}{\left(u^{y}\right)^2+1}-{\pi^{xy}} {u^{y}}\right)\\
  &\quad + \frac{\left(u^{\tau}\right)^2 \left(2 {u^{x}} {u^{y}} \left({\pi^{\eta}{}_{\eta}} \left(u^{x}\right)^2+{\pi^{yy}} \left(u^{x}\right)^2+{\pi^{yy}}\right)+{\pi^{xy}} \left(\left(\left(u^{y}\right)^2+1\right)^2-2 \left(u^{x}\right)^2 \left(\left(u^{y}\right)^2-1\right)\right)\right)}{\left(\left(u^{y}\right)^2+1\right)^2}\Bigr) \\
(A^{\mathrm{IS}}_x)_{14} &= -\frac{{u^{x}} \left(\left(u^{x}\right)^2+1\right)}{{u^{\tau}} \left(\left(u^{y}\right)^2+1\right)} \\
(A^{\mathrm{IS}}_x)_{15} &= -\frac{{u^{\tau}} {u^{x}}}{\left(u^{y}\right)^2+1} \\
(A^{\mathrm{IS}}_x)_{16} &= \frac{{u^{y}} \left(\frac{2 \left(u^{x}\right)^2}{\left(u^{y}\right)^2+1}+1\right)}{{u^{\tau}}} \\
(A^{\mathrm{IS}}_x)_{17} &= {u^{\tau}} {u^{x}} \\
(A^{\mathrm{IS}}_x)_{21} &= {p''(T)} {T} \left(u^{x}\right)^2+{p'(T)} \left(u^{x}\right)^2+{p'(T)} \\
(A^{\mathrm{IS}}_x)_{22} &= 2 {u^{x}} ({p'(T)} {T}+{\Pi})-\frac{2 ({\pi^{\eta}{}_{\eta}} {u^{x}}-{\pi^{xy}} {u^{y}}+{\pi^{yy}} {u^{x}})}{\left(u^{y}\right)^2+1} \\
(A^{\mathrm{IS}}_x)_{23} &= \frac{2 \left({u^{y}} \left({\pi^{\eta}{}_{\eta}} \left(u^{x}\right)^2+{\pi^{yy}} \left(u^{x}\right)^2+{\pi^{yy}}\right)+{\pi^{xy}} \left({u^{x}}-{u^{x}} \left(u^{y}\right)^2\right)\right)}{\left(\left(u^{y}\right)^2+1\right)^2} \\
(A^{\mathrm{IS}}_x)_{24} &= -\frac{\left(u^{x}\right)^2+1}{\left(u^{y}\right)^2+1} \\
(A^{\mathrm{IS}}_x)_{25} &= -\frac{\left(u^{\tau}\right)^2}{\left(u^{y}\right)^2+1} \\
(A^{\mathrm{IS}}_x)_{26} &= \frac{2 {u^{x}} {u^{y}}}{\left(u^{y}\right)^2+1} \\
(A^{\mathrm{IS}}_x)_{27} &= \left(u^{x}\right)^2+1 \\
(A^{\mathrm{IS}}_x)_{31} &= {u^{x}} {u^{y}} ({p''(T)} {T}+{p'(T)}) \\
(A^{\mathrm{IS}}_x)_{32} &= {u^{y}} ({p'(T)} {T}+{\Pi}) \\
(A^{\mathrm{IS}}_x)_{33} &= {u^{x}} ({p'(T)} {T}+{\Pi}) \\
(A^{\mathrm{IS}}_x)_{36} &= 1 \\
(A^{\mathrm{IS}}_x)_{37} &= {u^{x}} {u^{y}} \\
(A^{\mathrm{IS}}_x)_{42} &= -\frac{2 \left(3 {\pi^{xy}} {\tau_{\pi}} {u^{x}} {u^{y}} \left(\left(u^{y}\right)^2+1\right)-3 {\pi^{yy}} {\tau_{\pi}} \left(u^{x}\right)^2 \left(u^{y}\right)^2+\left(u^{\tau}\right)^2 \left(\left(u^{y}\right)^2+1\right) {\eta}\right)}{3 \left(u^{\tau}\right)^2} \\
(A^{\mathrm{IS}}_x)_{43} &= \frac{2 {u^{x}} {u^{y}} \left({\pi^{xy}} {\tau_{\pi}} {u^{x}} {u^{y}}-{\pi^{yy}} {\tau_{\pi}} \left(\left(u^{x}\right)^2+1\right)+\left(u^{\tau}\right)^2 {\eta}\right)}{\left(u^{\tau}\right)^2} \\
(A^{\mathrm{IS}}_x)_{44} &= {\tau_{\pi}} {u^{x}} \\
(A^{\mathrm{IS}}_x)_{52} &= -\frac{2 {\eta}}{3 {\tau}^2} \\
(A^{\mathrm{IS}}_x)_{55} &= \frac{{\tau_{\pi}} {u^{x}}}{{\tau}^2} \\
(A^{\mathrm{IS}}_x)_{62} &= \frac{{u^{x}} \left({u^{y}} \left(3 {\pi^{\eta}{}_{\eta}} {\tau_{\pi}} \left(u^{\tau}\right)^2+3 {\pi^{yy}} \left(2 {\tau_{\pi}} \left(u^{x}\right)^2+{\tau_{\pi}}\right)+\left(u^{\tau}\right)^2 {\eta}\right)-3 {\pi^{xy}} {\tau_{\pi}} {u^{x}} \left(2 \left(u^{y}\right)^2+1\right)\right)}{3 \left(u^{\tau}\right)^2} \\
(A^{\mathrm{IS}}_x)_{63} &= \frac{1}{\left(u^{\tau}\right)^2 \left(\left(u^{y}\right)^2+1\right)}\Bigl(- {\pi^{yy}} {\tau_{\pi}} \left(u^{x}\right)^2 \left(\left(u^{x}\right)^2+1\right) \left(2 \left(u^{y}\right)^2+1\right) + {\pi^{xy}} {\tau_{\pi}} {u^{x}} {u^{y}} \left(\left(2 \left(u^{x}\right)^2-1\right) \left(u^{y}\right)^2-1\right)\\
  &\quad - \left(u^{\tau}\right)^2 \left({\pi^{\eta}{}_{\eta}} {\tau_{\pi}} \left(u^{x}\right)^2 \left(u^{y}\right)^2-\left(\left(u^{x}\right)^2+1\right) \left(\left(u^{y}\right)^2+1\right) {\eta}\right)\Bigr) \\
(A^{\mathrm{IS}}_x)_{66} &= {\tau_{\pi}} {u^{x}} \\
(A^{\mathrm{IS}}_x)_{72} &= {\zeta} \\
(A^{\mathrm{IS}}_x)_{77} &= {\tau_{\Pi}} {u^{x}}
\end{align*}
\begin{align*}
(A^{\mathrm{IS}}_y)_{11} &= {u^{\tau}} {u^{y}} ({p''(T)} {T}+{p'(T)}) \\
(A^{\mathrm{IS}}_y)_{12} &= \frac{{u^{x}} {u^{y}} \left({p'(T)} {T} \left(u^{\tau}\right)^2+{\Pi} \left(u^{\tau}\right)^2-{\pi^{yy}}\right)+{\pi^{xy}} \left(\left(u^{y}\right)^2+1\right)}{\left(u^{\tau}\right)^3} \\
(A^{\mathrm{IS}}_y)_{13} &= \frac{\left(u^{\tau}\right)^4 ({p'(T)} {T}+{\Pi})+\left(u^{\tau}\right)^2 \left(u^{y}\right)^2 ({p'(T)} {T}+{\Pi})-{u^{y}} ({\pi^{xy}} {u^{x}}+{\pi^{yy}} {u^{y}})+{\pi^{yy}} \left(u^{\tau}\right)^2}{\left(u^{\tau}\right)^3} \\
(A^{\mathrm{IS}}_y)_{14} &= \frac{{u^{y}}}{{u^{\tau}}} \\
(A^{\mathrm{IS}}_y)_{16} &= \frac{{u^{x}}}{{u^{\tau}}} \\
(A^{\mathrm{IS}}_y)_{17} &= {u^{\tau}} {u^{y}} \\
(A^{\mathrm{IS}}_y)_{21} &= {u^{x}} {u^{y}} ({p''(T)} {T}+{p'(T)}) \\
(A^{\mathrm{IS}}_y)_{22} &= {u^{y}} ({p'(T)} {T}+{\Pi}) \\
(A^{\mathrm{IS}}_y)_{23} &= {u^{x}} ({p'(T)} {T}+{\Pi}) \\
(A^{\mathrm{IS}}_y)_{26} &= 1 \\
(A^{\mathrm{IS}}_y)_{27} &= {u^{x}} {u^{y}} \\
(A^{\mathrm{IS}}_y)_{31} &= {p''(T)} {T} \left(u^{y}\right)^2+{p'(T)} \left(u^{y}\right)^2+{p'(T)} \\
(A^{\mathrm{IS}}_y)_{33} &= 2 {u^{y}} ({p'(T)} {T}+{\Pi}) \\
(A^{\mathrm{IS}}_y)_{34} &= 1 \\
(A^{\mathrm{IS}}_y)_{37} &= \left(u^{y}\right)^2+1 \\
(A^{\mathrm{IS}}_y)_{42} &= -\frac{2 {\tau_{\pi}} \left(u^{y}\right)^2 \left({\pi^{xy}} \left(u^{y}\right)^2+{\pi^{xy}}-{\pi^{yy}} {u^{x}} {u^{y}}\right)}{\left(u^{\tau}\right)^2} \\
(A^{\mathrm{IS}}_y)_{43} &= \frac{6 {\pi^{xy}} {\tau_{\pi}} {u^{x}} \left(u^{y}\right)^3-6 {\pi^{yy}} {\tau_{\pi}} \left(\left(u^{x}\right)^2+1\right) \left(u^{y}\right)^2+4 \left(u^{\tau}\right)^2 \left(\left(u^{y}\right)^2+1\right) {\eta}}{3 \left(u^{\tau}\right)^2} \\
(A^{\mathrm{IS}}_y)_{44} &= {\tau_{\pi}} {u^{y}} \\
(A^{\mathrm{IS}}_y)_{53} &= -\frac{2 {\eta}}{3 {\tau}^2} \\
(A^{\mathrm{IS}}_y)_{55} &= \frac{{\tau_{\pi}} {u^{y}}}{{\tau}^2} \\
(A^{\mathrm{IS}}_y)_{62} &= \frac{1}{\left(u^{\tau}\right)^2}\Bigl(- {\pi^{xy}} {\tau_{\pi}} {u^{x}} {u^{y}} \left(2 \left(u^{y}\right)^2+1\right) + {\eta} + \left(u^{y}\right)^4 ({\pi^{\eta}{}_{\eta}} {\tau_{\pi}}+{\eta}) + \left(u^{y}\right)^2 ({\pi^{\eta}{}_{\eta}} {\tau_{\pi}}+{\pi^{yy}} {\tau_{\pi}}+2 {\eta})\\
  &\quad + \left(u^{x}\right)^2 \left(\left(u^{y}\right)^2 ({\pi^{\eta}{}_{\eta}} {\tau_{\pi}}+2 {\pi^{yy}} {\tau_{\pi}}+{\eta})+{\eta}\right)\Bigr) \\
(A^{\mathrm{IS}}_y)_{63} &= -\frac{1}{3 \left(u^{\tau}\right)^2 \left(\left(u^{y}\right)^2+1\right)}{u^{y}}\Bigl(3 {\pi^{yy}} {\tau_{\pi}} {u^{x}} \left(\left(u^{x}\right)^2+1\right) \left(2 \left(u^{y}\right)^2+1\right) + 3 {\pi^{xy}} {\tau_{\pi}} {u^{y}} \left(\left(1-2 \left(u^{x}\right)^2\right) \left(u^{y}\right)^2+1\right)\\
  &\quad + \left(u^{\tau}\right)^2 {u^{x}} \left(3 {\pi^{\eta}{}_{\eta}} {\tau_{\pi}} \left(u^{y}\right)^2-\left(\left(u^{y}\right)^2+1\right) {\eta}\right)\Bigr) \\
(A^{\mathrm{IS}}_y)_{66} &= {\tau_{\pi}} {u^{y}} \\
(A^{\mathrm{IS}}_y)_{73} &= {\zeta} \\
(A^{\mathrm{IS}}_y)_{77} &= {\tau_{\Pi}} {u^{y}}
\end{align*}
\begin{align*}
S^{\mathrm{IS}}_{1} &= \frac{1}{{\tau} \left(\left(u^{y}\right)^2+1\right)}\Bigl({p'(T)} {T} - {\pi^{yy}} \left(u^{x}\right)^2 + {p'(T)} {T} \left(u^{x}\right)^2 + 2 {\pi^{xy}} {u^{x}} {u^{y}} + {\pi^{yy}} \left(u^{y}\right)^2 + 2 {p'(T)} {T} \left(u^{y}\right)^2\\
  &\quad + {p'(T)} {T} \left(u^{x}\right)^2 \left(u^{y}\right)^2 + {p'(T)} {T} \left(u^{y}\right)^4 + {\Pi} \left(u^{\tau}\right)^2 \left(\left(u^{y}\right)^2+1\right) + {\pi^{\eta}{}_{\eta}} \left(-\left(u^{x}\right)^2+\left(u^{y}\right)^2+1\right)\Bigr) \\
S^{\mathrm{IS}}_{2} &= \frac{1}{{\tau} {u^{\tau}} \left(\left(u^{y}\right)^2+1\right)}\Bigl(- {\pi^{yy}} {u^{x}} + {p'(T)} {T} {u^{x}} - {\pi^{\eta}{}_{\eta}} \left(u^{\tau}\right)^2 {u^{x}} - {\pi^{yy}} \left(u^{x}\right)^3 + {p'(T)} {T} \left(u^{x}\right)^3 + {\pi^{xy}} {u^{y}} + 2 {\pi^{xy}} \left(u^{x}\right)^2 {u^{y}}\\
  &\quad + 2 {p'(T)} {T} {u^{x}} \left(u^{y}\right)^2 + {p'(T)} {T} \left(u^{x}\right)^3 \left(u^{y}\right)^2 + {\pi^{xy}} \left(u^{y}\right)^3 + {p'(T)} {T} {u^{x}} \left(u^{y}\right)^4\\
  &\quad + {\Pi} \left(u^{\tau}\right)^2 {u^{x}} \left(\left(u^{y}\right)^2+1\right)\Bigr) \\
S^{\mathrm{IS}}_{3} &= \frac{{u^{y}} \left({p'(T)} {T} \left(u^{\tau}\right)^2+{\Pi} \left(u^{\tau}\right)^2+{\pi^{yy}}\right)+{\pi^{xy}} {u^{x}}}{{\tau} {u^{\tau}}} \\
S^{\mathrm{IS}}_{4} &= {\pi^{yy}}-\frac{2 {u^{\tau}} \left(\left(u^{y}\right)^2+1\right) {\eta}}{3 {\tau}} \\
S^{\mathrm{IS}}_{5} &= \frac{3 {\pi^{\eta}{}_{\eta}} {\tau}+4 {u^{\tau}} {\eta}}{3 {\tau}^3} \\
S^{\mathrm{IS}}_{6} &= {\pi^{xy}}-\frac{2 {u^{\tau}} {u^{x}} {u^{y}} {\eta}}{3 {\tau}} \\
S^{\mathrm{IS}}_{7} &= {\Pi}+\frac{{u^{\tau}} {\zeta}}{{\tau}}
\end{align*}
\subsubsection{\texorpdfstring{$\delta_{\pi\pi}$}{\unichar{"03B4}\unichar{"03C0}\unichar{"03C0}} part}

\begin{align*}
(A^{(\delta)}_t)_{42} &= \frac{{\pi^{yy}} {u^{x}}}{{u^{\tau}}} \\
(A^{(\delta)}_t)_{43} &= \frac{{\pi^{yy}} {u^{y}}}{{u^{\tau}}} \\
(A^{(\delta)}_t)_{52} &= \frac{{\pi^{\eta}{}_{\eta}} {u^{x}}}{{\tau}^2 {u^{\tau}}} \\
(A^{(\delta)}_t)_{53} &= \frac{{\pi^{\eta}{}_{\eta}} {u^{y}}}{{\tau}^2 {u^{\tau}}} \\
(A^{(\delta)}_t)_{62} &= \frac{{\pi^{xy}} {u^{x}}}{{u^{\tau}}} \\
(A^{(\delta)}_t)_{63} &= \frac{{\pi^{xy}} {u^{y}}}{{u^{\tau}}}
\end{align*}
\begin{align*}
(A^{(\delta)}_x)_{42} &= {\pi^{yy}} \\
(A^{(\delta)}_x)_{52} &= \frac{{\pi^{\eta}{}_{\eta}}}{{\tau}^2} \\
(A^{(\delta)}_x)_{62} &= {\pi^{xy}}
\end{align*}
\begin{align*}
(A^{(\delta)}_y)_{43} &= {\pi^{yy}} \\
(A^{(\delta)}_y)_{53} &= \frac{{\pi^{\eta}{}_{\eta}}}{{\tau}^2} \\
(A^{(\delta)}_y)_{63} &= {\pi^{xy}}
\end{align*}
\begin{align*}
S^{(\delta)}_{4} &= \frac{{\pi^{yy}} {u^{\tau}}}{{\tau}} \\
S^{(\delta)}_{5} &= \frac{{\pi^{\eta}{}_{\eta}} {u^{\tau}}}{{\tau}^3} \\
S^{(\delta)}_{6} &= \frac{{\pi^{xy}} {u^{\tau}}}{{\tau}}
\end{align*}
\subsubsection{Co-rotating terms (diagnostics)}

Not part of the system above; exported alone so that a test can ask whether a 2+1D matrix carries them. The projected vorticity coupling, without a factor $\tau_{\pi}$ (\texttt{two\_d\_viscous\_vorticity}):

\[
2\,\pi_{\lambda}{}^{\langle\mu}\omega^{\nu\rangle\lambda},\qquad \omega^{\mu\nu} = \tfrac{1}{2}\Delta^{\mu\alpha}\Delta^{\nu\beta}(\nabla_{\alpha}u_{\beta} - \nabla_{\beta}u_{\alpha})
\]
and the unprojected commutator (\texttt{two\_d\_viscous\_notebook\_commutator}):

\[
\tau_{\pi}\bigl(\omega^{\mu}{}_{\lambda}\pi^{\lambda\nu} - \pi^{\mu\lambda}\omega_{\lambda}{}^{\nu}\bigr),\qquad \omega_{\mu\nu} = \tfrac{1}{2}(\nabla_{\mu}u_{\nu} - \nabla_{\nu}u_{\mu})
\]
Both carry no source (V-GV-src, V-GN-src). Nonzero entries of their matrices, $A^{(\omega)}$ for the first and $A^{(\mathrm{c})}$ for the second, rows and columns as above:

\[
\arraycolsep=2.5pt
A^{(\omega)}_t = \begin{pmatrix}\cdot & \cdot & \cdot & \cdot & \cdot & \cdot & \cdot \\ \cdot & \cdot & \cdot & \cdot & \cdot & \cdot & \cdot \\ \cdot & \cdot & \cdot & \cdot & \cdot & \cdot & \cdot \\ \cdot & \bullet & \bullet & \cdot & \cdot & \cdot & \cdot \\ \cdot & \cdot & \cdot & \cdot & \cdot & \cdot & \cdot \\ \cdot & \bullet & \bullet & \cdot & \cdot & \cdot & \cdot \\ \cdot & \cdot & \cdot & \cdot & \cdot & \cdot & \cdot\end{pmatrix},\qquad A^{(\omega)}_x = \begin{pmatrix}\cdot & \cdot & \cdot & \cdot & \cdot & \cdot & \cdot \\ \cdot & \cdot & \cdot & \cdot & \cdot & \cdot & \cdot \\ \cdot & \cdot & \cdot & \cdot & \cdot & \cdot & \cdot \\ \cdot & \bullet & \bullet & \cdot & \cdot & \cdot & \cdot \\ \cdot & \cdot & \cdot & \cdot & \cdot & \cdot & \cdot \\ \cdot & \bullet & \bullet & \cdot & \cdot & \cdot & \cdot \\ \cdot & \cdot & \cdot & \cdot & \cdot & \cdot & \cdot\end{pmatrix},\qquad A^{(\omega)}_y = \begin{pmatrix}\cdot & \cdot & \cdot & \cdot & \cdot & \cdot & \cdot \\ \cdot & \cdot & \cdot & \cdot & \cdot & \cdot & \cdot \\ \cdot & \cdot & \cdot & \cdot & \cdot & \cdot & \cdot \\ \cdot & \bullet & \bullet & \cdot & \cdot & \cdot & \cdot \\ \cdot & \cdot & \cdot & \cdot & \cdot & \cdot & \cdot \\ \cdot & \bullet & \bullet & \cdot & \cdot & \cdot & \cdot \\ \cdot & \cdot & \cdot & \cdot & \cdot & \cdot & \cdot\end{pmatrix}
\]
\[
\arraycolsep=2.5pt
A^{(\mathrm{c})}_t = \begin{pmatrix}\cdot & \cdot & \cdot & \cdot & \cdot & \cdot & \cdot \\ \cdot & \cdot & \cdot & \cdot & \cdot & \cdot & \cdot \\ \cdot & \cdot & \cdot & \cdot & \cdot & \cdot & \cdot \\ \cdot & \cdot & \bullet & \cdot & \cdot & \cdot & \cdot \\ \cdot & \cdot & \cdot & \cdot & \cdot & \cdot & \cdot \\ \cdot & \bullet & \bullet & \cdot & \cdot & \cdot & \cdot \\ \cdot & \cdot & \cdot & \cdot & \cdot & \cdot & \cdot\end{pmatrix},\qquad A^{(\mathrm{c})}_x = \begin{pmatrix}\cdot & \cdot & \cdot & \cdot & \cdot & \cdot & \cdot \\ \cdot & \cdot & \cdot & \cdot & \cdot & \cdot & \cdot \\ \cdot & \cdot & \cdot & \cdot & \cdot & \cdot & \cdot \\ \cdot & \cdot & \bullet & \cdot & \cdot & \cdot & \cdot \\ \cdot & \cdot & \cdot & \cdot & \cdot & \cdot & \cdot \\ \cdot & \bullet & \bullet & \cdot & \cdot & \cdot & \cdot \\ \cdot & \cdot & \cdot & \cdot & \cdot & \cdot & \cdot\end{pmatrix},\qquad A^{(\mathrm{c})}_y = \begin{pmatrix}\cdot & \cdot & \cdot & \cdot & \cdot & \cdot & \cdot \\ \cdot & \cdot & \cdot & \cdot & \cdot & \cdot & \cdot \\ \cdot & \cdot & \cdot & \cdot & \cdot & \cdot & \cdot \\ \cdot & \bullet & \bullet & \cdot & \cdot & \cdot & \cdot \\ \cdot & \cdot & \cdot & \cdot & \cdot & \cdot & \cdot \\ \cdot & \bullet & \bullet & \cdot & \cdot & \cdot & \cdot \\ \cdot & \cdot & \cdot & \cdot & \cdot & \cdot & \cdot\end{pmatrix}
\]
\subsection{Checks}

{\small
\begin{longtable}{@{}>{\raggedright\arraybackslash}p{0.27\textwidth}>{\raggedright\arraybackslash}p{0.60\textwidth}l@{}}
\toprule
id & what it tests & \\
\midrule
\endhead
\texttt{V-\allowbreak{}G0cyl} & cylindrical Milne: Nabla g =\allowbreak{}=\allowbreak{} 0 & PASS \\
\texttt{V-\allowbreak{}G0car} & Cartesian   Milne: Nabla g =\allowbreak{}=\allowbreak{} 0 & PASS \\
\texttt{V-\allowbreak{}G1cyl} & cylindrical: Nabla.V =\allowbreak{} X' + X/\allowbreak{}tau & PASS \\
\texttt{V-\allowbreak{}G1car} & Cartesian:   Nabla.V =\allowbreak{} X' + X/\allowbreak{}tau & PASS \\
\texttt{V-\allowbreak{}G2cyl-\allowbreak{}uu} & u.u =\allowbreak{}=\allowbreak{} -1 & PASS \\
\texttt{V-\allowbreak{}G2cyl-\allowbreak{}un} & u\_\allowbreak{}mu pi\textasciicircum{}\{mu nu\} =\allowbreak{}=\allowbreak{} 0 & PASS \\
\texttt{V-\allowbreak{}G2cyl-\allowbreak{}tr} & pi\textasciicircum{}mu\_\allowbreak{}mu =\allowbreak{}=\allowbreak{} 0 & PASS \\
\texttt{V-\allowbreak{}G3cyl} & cylindrical system is quasi-linear & PASS \\
\texttt{V-\allowbreak{}G6cyl} & cylindrical: (IS + dpp pi theta in eq. 2) =\allowbreak{}=\allowbreak{} IS + dpp * (the exported dpp part),\allowbreak{} equation by equation & PASS \\
\texttt{V-\allowbreak{}G4} & 1-D: derived =\allowbreak{}=\allowbreak{} shipped one\_\allowbreak{}d\_\allowbreak{}viscous\_\allowbreak{}matrix,\allowbreak{} A\_\allowbreak{}t,\allowbreak{} A\_\allowbreak{}r,\allowbreak{} src entry by entry,\allowbreak{} read from the .jl,\allowbreak{} SYMBOLIC dpp  [CONTROL] & PASS \\
\texttt{V-\allowbreak{}G4-\allowbreak{}0} & 1-D at dpp =\allowbreak{} 0 (plain IS): plain IS =\allowbreak{}=\allowbreak{} the shipped IS part & PASS \\
\texttt{V-\allowbreak{}G4-\allowbreak{}1} & 1-D at dpp =\allowbreak{} 4/\allowbreak{}3 tauS (default) & PASS \\
\texttt{V-\allowbreak{}G4m} & 1-D at dpp =\allowbreak{} 4/\allowbreak{}3 tauS =\allowbreak{}=\allowbreak{} the medium block of the 7-field :MIS one\_\allowbreak{}d\_\allowbreak{}viscous\_\allowbreak{}matrix\_\allowbreak{}fugacity,\allowbreak{} raw,\allowbreak{} read from the .jl & PASS \\
\texttt{V-\allowbreak{}G4m-\allowbreak{}ctl} & control: at dpp =\allowbreak{} 0 it is NOT the :MIS medium block & PASS \\
\texttt{V-\allowbreak{}G2car-\allowbreak{}uu} & u.u =\allowbreak{}=\allowbreak{} -1 & PASS \\
\texttt{V-\allowbreak{}G2car-\allowbreak{}un} & u\_\allowbreak{}mu pi\textasciicircum{}\{mu nu\} =\allowbreak{}=\allowbreak{} 0 & PASS \\
\texttt{V-\allowbreak{}G2car-\allowbreak{}tr} & pi\textasciicircum{}mu\_\allowbreak{}mu =\allowbreak{}=\allowbreak{} 0 & PASS \\
\texttt{V-\allowbreak{}G3car} & Cartesian system is quasi-linear & PASS \\
\texttt{V-\allowbreak{}G6car} & Cartesian: (IS + dpp pi theta in eq. 2) =\allowbreak{}=\allowbreak{} IS + dpp * (the exported dpp part),\allowbreak{} equation by equation & PASS \\
\texttt{V-\allowbreak{}G5-\allowbreak{}e} & energy row on Bjorken =\allowbreak{}=\allowbreak{} d\_\allowbreak{}tau e + (e+p+Pi+pi\textasciicircum{}eta\_\allowbreak{}eta)/\allowbreak{}tau & PASS \\
\texttt{V-\allowbreak{}G5-\allowbreak{}pi} & pi\textasciicircum{}eta\_\allowbreak{}eta row on Bjorken =\allowbreak{}=\allowbreak{} tauS pi' + pi + 4 eta/\allowbreak{}(3 tau) + dpp pi/\allowbreak{}tau  (theta =\allowbreak{} 1/\allowbreak{}tau) & PASS \\
\texttt{V-\allowbreak{}G5-\allowbreak{}Pi} & bulk row on Bjorken =\allowbreak{}=\allowbreak{} tauB Pi' + Pi + zeta/\allowbreak{}tau & PASS \\
\texttt{V-\allowbreak{}G5-\allowbreak{}mom} & the x,\allowbreak{} y and xy rows vanish identically on Bjorken flow & PASS \\
\texttt{V-\allowbreak{}GV-\allowbreak{}anti} & omega is antisymmetric & PASS \\
\texttt{V-\allowbreak{}GV-\allowbreak{}perp} & omega is transverse to u & PASS \\
\texttt{V-\allowbreak{}GV-\allowbreak{}src} & the vorticity term is purely a gradient term (no source) & PASS \\
\texttt{V-\allowbreak{}GN-\allowbreak{}anti} & the notebook's omega is antisymmetric & PASS \\
\texttt{V-\allowbreak{}GN-\allowbreak{}src} & the notebook's commutator is purely a gradient term (no source) & PASS \\
\bottomrule
\end{longtable}
}
28 passed, 0 failed (28 checks evaluated).

Not checked in this script: the slab limit (2+1D at $u^{y} = 0$, $\partial_{y} = 0$ against 1+1D at $r \to \infty$); it is a numerical comparison of the two generated routines.

Re-run with \texttt{wolframscript -file deriv/derive\_\allowbreak{}2p1d\_\allowbreak{}viscous.wls}.


\clearpage
% GENERATED by deriv/render.wl from the derivation script.  Do not edit.

\section{Charm first moment on a 2+1D background (\texorpdfstring{$\nabla\alpha$}{\unichar{"2207}\unichar{"03B1}} closure)}

Script: \texttt{deriv/derive\_hq\_2p1d.wls}

Writes: \texttt{two\_d\_HQ\_BG\_matrix} and \texttt{two\_d\_HQ\_live\_rows} in \texttt{src/Matrix/2d\_viscous\_HQ\_BG.jl}

\subsection{Setup}

Cartesian Milne coordinates, boost invariant, on a prescribed background $T$, $u^{x}$, $u^{y}$:

\[
x^{\mu} = ({\tau}, {x}, {y}, {\eta}),\qquad g_{\mu\nu} = \mathrm{diag}(-1, 1, 1, {\tau}^2),\qquad u^{\mu} = (u^{\tau}, u^{x}, u^{y}, 0),\quad u^{\tau} = \sqrt{\left(u^{x}\right)^2+\left(u^{y}\right)^2+1}
\]
Charm is a passive tracer:

\begin{gather*}
N^{\mu} = n\,u^{\mu} + \nu^{\mu},\qquad \nabla_{\mu}N^{\mu} = 0 \\
\tau_{n}\,\Delta^{\mu}{}_{\beta}\,D\nu^{\beta} + \nu^{\mu} + \kappa\,\Delta^{\mu\beta}\partial_{\beta}\alpha = 0 \\
n = n(T, \alpha),\qquad \kappa = D_{s}\,\partial_{\alpha}n,\qquad D = u^{\mu}\nabla_{\mu},\qquad \Delta^{\mu\nu} = g^{\mu\nu} + u^{\mu}u^{\nu}
\end{gather*}
Navier--Stokes limit: Fick's law $\nu^{\mu} = -\kappa\,\nabla^{\mu}_{\perp}\alpha$. Fields, with $\nu^{\tau}$ from $u_{\mu}\nu^{\mu} = 0$:

\[
\phi = ({\alpha}, {\nu^{x}}, {\nu^{y}}),\qquad \nu^{\tau} = \frac{{\nu^{x}} {u^{x}}+{\nu^{y}} {u^{y}}}{{u^{\tau}}}
\]
Rows: $\nabla_{\mu}N^{\mu}$ and the $x$, $y$ components of the current equation, in

\[
A_t\,\partial_{\tau}\phi + A_x\,\partial_{x}\phi + A_y\,\partial_{y}\phi + S = 0
\]
\subsection{Steps}

\[
\Gamma^{{\tau}}{}_{{\eta}{\eta}} = {\tau},\quad \Gamma^{{\eta}}{}_{{\tau}{\eta}} = \frac{1}{{\tau}}
\]
\begin{align*}
a^{{\tau}} &= \frac{{(\partial_{\tau} u^{x})} {u^{\tau}} {u^{x}}+{(\partial_{\tau} u^{y})} {u^{\tau}} {u^{y}}+{(\partial_{x} u^{x})} \left(u^{x}\right)^2+{(\partial_{x} u^{y})} {u^{x}} {u^{y}}+{(\partial_{y} u^{x})} {u^{x}} {u^{y}}+{(\partial_{y} u^{y})} \left(u^{y}\right)^2}{{u^{\tau}}} \\
a^{{x}} &= {(\partial_{\tau} u^{x})} {u^{\tau}}+{(\partial_{x} u^{x})} {u^{x}}+{(\partial_{y} u^{x})} {u^{y}} \\
a^{{y}} &= {(\partial_{\tau} u^{y})} {u^{\tau}}+{(\partial_{x} u^{y})} {u^{x}}+{(\partial_{y} u^{y})} {u^{y}} \\
\nu_{\lambda}a^{\lambda} &= {\nu^{x}} ({(\partial_{\tau} u^{x})} {u^{\tau}}+{(\partial_{x} u^{x})} {u^{x}}+{(\partial_{y} u^{x})} {u^{y}}) + {\nu^{y}} ({(\partial_{\tau} u^{y})} {u^{\tau}}+{(\partial_{x} u^{y})} {u^{x}}+{(\partial_{y} u^{y})} {u^{y}})\\
  &\quad - \frac{1}{\left(u^{\tau}\right)^2}{\nu^{x}} {u^{x}}+{\nu^{y}} {u^{y}}\Bigl({u^{\tau}} ({(\partial_{\tau} u^{x})} {u^{x}}+{(\partial_{\tau} u^{y})} {u^{y}}) + {u^{x}} ({(\partial_{x} u^{x})} {u^{x}}+{(\partial_{x} u^{y})} {u^{y}})\\
  &\quad + {u^{y}} ({(\partial_{y} u^{x})} {u^{x}}+{(\partial_{y} u^{y})} {u^{y}})\Bigr) \\
D\nu^{{\tau}} &= \frac{1}{\left(u^{\tau}\right)^3}\Bigl(\left(u^{\tau}\right)^3 ({(\partial_{{\tau}} \nu^{x})} {u^{x}}+{(\partial_{{\tau}} \nu^{y})} {u^{y}}+{(\partial_{\tau} u^{x})} {\nu^{x}}+{(\partial_{\tau} u^{y})} {\nu^{y}})\\
  &\quad + \left(u^{\tau}\right)^2 {u^{x}} ({(\partial_{{x}} \nu^{x})} {u^{x}}+{(\partial_{{x}} \nu^{y})} {u^{y}}+{(\partial_{x} u^{x})} {\nu^{x}}+{(\partial_{x} u^{y})} {\nu^{y}})\\
  &\quad + \left(u^{\tau}\right)^2 {u^{y}} ({(\partial_{{y}} \nu^{x})} {u^{x}}+{(\partial_{{y}} \nu^{y})} {u^{y}}+{(\partial_{y} u^{x})} {\nu^{x}}+{(\partial_{y} u^{y})} {\nu^{y}})\\
  &\quad - {u^{\tau}} ({(\partial_{\tau} u^{x})} {u^{x}}+{(\partial_{\tau} u^{y})} {u^{y}}) ({\nu^{x}} {u^{x}}+{\nu^{y}} {u^{y}})\\
  &\quad - {u^{x}} ({(\partial_{x} u^{x})} {u^{x}}+{(\partial_{x} u^{y})} {u^{y}}) ({\nu^{x}} {u^{x}}+{\nu^{y}} {u^{y}})\\
  &\quad - {u^{y}} ({(\partial_{y} u^{x})} {u^{x}}+{(\partial_{y} u^{y})} {u^{y}}) ({\nu^{x}} {u^{x}}+{\nu^{y}} {u^{y}})\Bigr) \\
D\nu^{{x}} &= {(\partial_{{\tau}} \nu^{x})} {u^{\tau}}+{(\partial_{{x}} \nu^{x})} {u^{x}}+{(\partial_{{y}} \nu^{x})} {u^{y}} \\
D\nu^{{y}} &= {(\partial_{{\tau}} \nu^{y})} {u^{\tau}}+{(\partial_{{x}} \nu^{y})} {u^{x}}+{(\partial_{{y}} \nu^{y})} {u^{y}}
\end{align*}
$\Delta^{\mu}{}_{\beta}D\nu^{\beta} = D\nu^{\mu} - u^{\mu}\nu_{\lambda}a^{\lambda}$, because $u_{\beta}\nu^{\beta} = 0$. Rows:

\begin{align*}
E_{1} = \nabla_{\mu}N^{\mu} &= {(\partial_{{x}} \nu^{x})} + {(\partial_{{y}} \nu^{y})} + {(\partial_{x} u^{x})} {n} + {(\partial_{y} u^{y})} {n} + {u^{\tau}} ({(\partial_{{\tau}} \alpha)} {(\partial_{\alpha} n)}+{(\partial_{\tau} T)} {(\partial_{T} n)})\\
  &\quad + {u^{x}} ({(\partial_{{x}} \alpha)} {(\partial_{\alpha} n)}+{(\partial_{x} T)} {(\partial_{T} n)}) + {u^{y}} ({(\partial_{{y}} \alpha)} {(\partial_{\alpha} n)}+{(\partial_{y} T)} {(\partial_{T} n)})\\
  &\quad + \frac{{(\partial_{{\tau}} \nu^{x})} {u^{x}}+{(\partial_{{\tau}} \nu^{y})} {u^{y}}+{(\partial_{\tau} u^{x})} {\nu^{x}}+{(\partial_{\tau} u^{y})} {\nu^{y}}}{{u^{\tau}}} + \frac{{n} ({(\partial_{\tau} u^{x})} {u^{x}}+{(\partial_{\tau} u^{y})} {u^{y}})}{{u^{\tau}}}\\
  &\quad - \frac{({(\partial_{\tau} u^{x})} {u^{x}}+{(\partial_{\tau} u^{y})} {u^{y}}) ({\nu^{x}} {u^{x}}+{\nu^{y}} {u^{y}})}{\left(u^{\tau}\right)^3} + \frac{{n} \left(u^{\tau}\right)^2+{\nu^{x}} {u^{x}}+{\nu^{y}} {u^{y}}}{{\tau} {u^{\tau}}} \\
E_{2} &= {\nu^{x}} + {\kappa} \left({(\partial_{{\tau}} \alpha)} {u^{\tau}} {u^{x}}+{(\partial_{{x}} \alpha)} \left(u^{x}\right)^2+{(\partial_{{x}} \alpha)}+{(\partial_{{y}} \alpha)} {u^{x}} {u^{y}}\right)\\
  &\quad + {\tau_{n}}\Bigl({(\partial_{{\tau}} \nu^{x})} {u^{\tau}} + {(\partial_{{x}} \nu^{x})} {u^{x}} + {(\partial_{{y}} \nu^{x})} {u^{y}}\\
  &\quad - {u^{x}}\Bigl({\nu^{x}} ({(\partial_{\tau} u^{x})} {u^{\tau}}+{(\partial_{x} u^{x})} {u^{x}}+{(\partial_{y} u^{x})} {u^{y}})\\
  &\quad + {\nu^{y}} ({(\partial_{\tau} u^{y})} {u^{\tau}}+{(\partial_{x} u^{y})} {u^{x}}+{(\partial_{y} u^{y})} {u^{y}})\\
  &\quad - \frac{1}{\left(u^{\tau}\right)^2}{\nu^{x}} {u^{x}}+{\nu^{y}} {u^{y}}\Bigl({u^{\tau}} ({(\partial_{\tau} u^{x})} {u^{x}}+{(\partial_{\tau} u^{y})} {u^{y}})\\
  &\quad + {u^{x}} ({(\partial_{x} u^{x})} {u^{x}}+{(\partial_{x} u^{y})} {u^{y}})\\
  &\quad + {u^{y}} ({(\partial_{y} u^{x})} {u^{x}}+{(\partial_{y} u^{y})} {u^{y}})\Bigr)\Bigr)\Bigr) \\
E_{3} &= {\nu^{y}} + {\kappa} \left({(\partial_{{\tau}} \alpha)} {u^{\tau}} {u^{y}}+{(\partial_{{x}} \alpha)} {u^{x}} {u^{y}}+{(\partial_{{y}} \alpha)} \left(u^{y}\right)^2+{(\partial_{{y}} \alpha)}\right)\\
  &\quad + {\tau_{n}}\Bigl({(\partial_{{\tau}} \nu^{y})} {u^{\tau}} + {(\partial_{{x}} \nu^{y})} {u^{x}} + {(\partial_{{y}} \nu^{y})} {u^{y}}\\
  &\quad - {u^{y}}\Bigl({\nu^{x}} ({(\partial_{\tau} u^{x})} {u^{\tau}}+{(\partial_{x} u^{x})} {u^{x}}+{(\partial_{y} u^{x})} {u^{y}})\\
  &\quad + {\nu^{y}} ({(\partial_{\tau} u^{y})} {u^{\tau}}+{(\partial_{x} u^{y})} {u^{x}}+{(\partial_{y} u^{y})} {u^{y}})\\
  &\quad - \frac{1}{\left(u^{\tau}\right)^2}{\nu^{x}} {u^{x}}+{\nu^{y}} {u^{y}}\Bigl({u^{\tau}} ({(\partial_{\tau} u^{x})} {u^{x}}+{(\partial_{\tau} u^{y})} {u^{y}})\\
  &\quad + {u^{x}} ({(\partial_{x} u^{x})} {u^{x}}+{(\partial_{x} u^{y})} {u^{y}})\\
  &\quad + {u^{y}} ({(\partial_{y} u^{x})} {u^{x}}+{(\partial_{y} u^{y})} {u^{y}})\Bigr)\Bigr)\Bigr)
\end{align*}
\[
A^{(k)}_{ij} = \frac{\partial E_{i}}{\partial(\partial_{k}\phi_{j})},\qquad S_{i} = E_{i}\bigl|_{\partial\phi = 0},\qquad k = \tau, x, y
\]
\subsection{Result}

\begin{align*}
(A_t)_{11} &= {(\partial_{\alpha} n)} {u^{\tau}} \\
(A_t)_{12} &= \frac{{u^{x}}}{{u^{\tau}}} \\
(A_t)_{13} &= \frac{{u^{y}}}{{u^{\tau}}} \\
(A_t)_{21} &= {\kappa} {u^{\tau}} {u^{x}} \\
(A_t)_{22} &= {\tau_{n}} {u^{\tau}} \\
(A_t)_{31} &= {\kappa} {u^{\tau}} {u^{y}} \\
(A_t)_{33} &= {\tau_{n}} {u^{\tau}}
\end{align*}
\begin{align*}
(A_x)_{11} &= {(\partial_{\alpha} n)} {u^{x}} \\
(A_x)_{12} &= 1 \\
(A_x)_{21} &= {\kappa} \left(\left(u^{x}\right)^2+1\right) \\
(A_x)_{22} &= {\tau_{n}} {u^{x}} \\
(A_x)_{31} &= {\kappa} {u^{x}} {u^{y}} \\
(A_x)_{33} &= {\tau_{n}} {u^{x}}
\end{align*}
\begin{align*}
(A_y)_{11} &= {(\partial_{\alpha} n)} {u^{y}} \\
(A_y)_{13} &= 1 \\
(A_y)_{21} &= {\kappa} {u^{x}} {u^{y}} \\
(A_y)_{22} &= {\tau_{n}} {u^{y}} \\
(A_y)_{31} &= {\kappa} \left(\left(u^{y}\right)^2+1\right) \\
(A_y)_{33} &= {\tau_{n}} {u^{y}}
\end{align*}
\begin{align*}
S_{1} &= \frac{1}{{\tau} \left(u^{\tau}\right)^3}\Bigl({(\partial_{\tau} u^{x})} {\nu^{x}} {\tau} \left(u^{\tau}\right)^2 + {(\partial_{\tau} u^{y})} {\nu^{y}} {\tau} \left(u^{\tau}\right)^2 + {(\partial_{x} u^{x})} {n} {\tau} \left(u^{\tau}\right)^3 + {(\partial_{y} u^{y})} {n} {\tau} \left(u^{\tau}\right)^3\\
  &\quad + {n} \left(u^{\tau}\right)^4 + {(\partial_{\tau} T)} {(\partial_{T} n)} {\tau} \left(u^{\tau}\right)^4 + {\nu^{x}} \left(u^{\tau}\right)^2 {u^{x}} + {(\partial_{\tau} u^{x})} {n} {\tau} \left(u^{\tau}\right)^2 {u^{x}}\\
  &\quad + {(\partial_{x} T)} {(\partial_{T} n)} {\tau} \left(u^{\tau}\right)^3 {u^{x}} - {(\partial_{\tau} u^{x})} {\nu^{x}} {\tau} \left(u^{x}\right)^2 + {\nu^{y}} \left(u^{\tau}\right)^2 {u^{y}} + {(\partial_{\tau} u^{y})} {n} {\tau} \left(u^{\tau}\right)^2 {u^{y}}\\
  &\quad + {(\partial_{y} T)} {(\partial_{T} n)} {\tau} \left(u^{\tau}\right)^3 {u^{y}} - {(\partial_{\tau} u^{y})} {\nu^{x}} {\tau} {u^{x}} {u^{y}} - {(\partial_{\tau} u^{x})} {\nu^{y}} {\tau} {u^{x}} {u^{y}}\\
  &\quad - {(\partial_{\tau} u^{y})} {\nu^{y}} {\tau} \left(u^{y}\right)^2\Bigr) \\
S_{2} &= \frac{1}{\left(u^{\tau}\right)^2}\Bigl(- {\nu^{y}} {\tau_{n}} {u^{x}}\Bigl({(\partial_{\tau} u^{y})} {u^{\tau}} + {(\partial_{x} u^{y})} \left(u^{x}\right)^3 + {(\partial_{y} u^{y})} {u^{y}}\\
  &\quad + \left(u^{x}\right)^2 ({(\partial_{\tau} u^{y})} {u^{\tau}}-{(\partial_{x} u^{x})} {u^{y}}+{(\partial_{y} u^{y})} {u^{y}})\\
  &\quad + {u^{x}} ({(\partial_{x} u^{y})}-{u^{y}} ({(\partial_{\tau} u^{x})} {u^{\tau}}+{(\partial_{y} u^{x})} {u^{y}}))\Bigr)\\
  &\quad + {\nu^{x}}\Bigl(1 + {(\partial_{x} u^{y})} {\tau_{n}} \left(u^{x}\right)^3 {u^{y}} + \left(u^{y}\right)^2\\
  &\quad - {\tau_{n}} {u^{x}} \left(\left(u^{y}\right)^2+1\right) ({(\partial_{\tau} u^{x})} {u^{\tau}}+{(\partial_{y} u^{x})} {u^{y}})\\
  &\quad + \left(u^{x}\right)^2\Bigl(1 + {(\partial_{\tau} u^{y})} {\tau_{n}} {u^{\tau}} {u^{y}} + {(\partial_{y} u^{y})} {\tau_{n}} \left(u^{y}\right)^2\\
  &\quad - {(\partial_{x} u^{x})} {\tau_{n}} \left(\left(u^{y}\right)^2+1\right)\Bigr)\Bigr)\Bigr) \\
S_{3} &= \frac{1}{\left(u^{\tau}\right)^2}\Bigl(- {\nu^{x}} {\tau_{n}} {u^{y}}\Bigl({(\partial_{\tau} u^{x})} {u^{\tau}} + {(\partial_{x} u^{x})} {u^{x}}\\
  &\quad + {u^{y}} ({(\partial_{y} u^{x})}-{u^{x}} ({(\partial_{\tau} u^{y})} {u^{\tau}}+{(\partial_{x} u^{y})} {u^{x}}))\\
  &\quad + \left(u^{y}\right)^2 ({(\partial_{\tau} u^{x})} {u^{\tau}}+{(\partial_{x} u^{x})} {u^{x}}-{(\partial_{y} u^{y})} {u^{x}}) + {(\partial_{y} u^{x})} \left(u^{y}\right)^3\Bigr)\\
  &\quad + {\nu^{y}}\Bigl(1 - {(\partial_{\tau} u^{y})} {\tau_{n}} {u^{\tau}} {u^{y}} - {(\partial_{x} u^{y})} {\tau_{n}} \left(u^{x}\right)^3 {u^{y}} + \left(u^{y}\right)^2 (1-{(\partial_{y} u^{y})} {\tau_{n}})\\
  &\quad + \left(u^{x}\right)^2 \left(-{(\partial_{\tau} u^{y})} {\tau_{n}} {u^{\tau}} {u^{y}}+{\tau_{n}} \left(u^{y}\right)^2 ({(\partial_{x} u^{x})}-{(\partial_{y} u^{y})})+1\right)\\
  &\quad + {\tau_{n}} {u^{x}} {u^{y}} ({u^{y}} ({(\partial_{\tau} u^{x})} {u^{\tau}}+{(\partial_{y} u^{x})} {u^{y}})-{(\partial_{x} u^{y})})\Bigr)\Bigr)
\end{align*}
\subsubsection{Live background}

The same three rows with the medium promoted to fields, $\phi_{L} = ({T}, {u^{x}}, {u^{y}}, {\alpha}, {\nu^{x}}, {\nu^{y}})$. Columns 4--6 are the matrices above (H-G6a); columns 1--3 and the remaining source:

\begin{align*}
(A^{L}_t)_{11} &= {(\partial_{T} n)} {u^{\tau}} \\
(A^{L}_t)_{12} &= \frac{{n} \left(u^{\tau}\right)^2 {u^{x}}+{\nu^{x}} \left(\left(u^{y}\right)^2+1\right)-{\nu^{y}} {u^{x}} {u^{y}}}{\left(u^{\tau}\right)^3} \\
(A^{L}_t)_{13} &= \frac{{n} \left(u^{\tau}\right)^2 {u^{y}}-{\nu^{x}} {u^{x}} {u^{y}}+{\nu^{y}} \left(\left(u^{x}\right)^2+1\right)}{\left(u^{\tau}\right)^3} \\
(A^{L}_t)_{22} &= \frac{{\tau_{n}} {u^{x}} \left({\nu^{y}} {u^{x}} {u^{y}}-{\nu^{x}} \left(\left(u^{y}\right)^2+1\right)\right)}{{u^{\tau}}} \\
(A^{L}_t)_{23} &= -\frac{{\tau_{n}} {u^{x}} \left(-{\nu^{x}} {u^{x}} {u^{y}}+{\nu^{y}} \left(u^{x}\right)^2+{\nu^{y}}\right)}{{u^{\tau}}} \\
(A^{L}_t)_{32} &= -\frac{{\tau_{n}} {u^{y}} \left({\nu^{x}} \left(u^{y}\right)^2+{\nu^{x}}-{\nu^{y}} {u^{x}} {u^{y}}\right)}{{u^{\tau}}} \\
(A^{L}_t)_{33} &= \frac{{\tau_{n}} {u^{y}} \left({\nu^{x}} {u^{x}} {u^{y}}-{\nu^{y}} \left(\left(u^{x}\right)^2+1\right)\right)}{{u^{\tau}}}
\end{align*}
\begin{align*}
(A^{L}_x)_{11} &= {(\partial_{T} n)} {u^{x}} \\
(A^{L}_x)_{12} &= {n} \\
(A^{L}_x)_{22} &= \frac{{\tau_{n}} \left(u^{x}\right)^2 \left({\nu^{y}} {u^{x}} {u^{y}}-{\nu^{x}} \left(\left(u^{y}\right)^2+1\right)\right)}{\left(u^{\tau}\right)^2} \\
(A^{L}_x)_{23} &= -\frac{{\tau_{n}} \left(u^{x}\right)^2 \left(-{\nu^{x}} {u^{x}} {u^{y}}+{\nu^{y}} \left(u^{x}\right)^2+{\nu^{y}}\right)}{\left(u^{\tau}\right)^2} \\
(A^{L}_x)_{32} &= \frac{{\tau_{n}} {u^{x}} {u^{y}} \left({\nu^{y}} {u^{x}} {u^{y}}-{\nu^{x}} \left(\left(u^{y}\right)^2+1\right)\right)}{\left(u^{\tau}\right)^2} \\
(A^{L}_x)_{33} &= \frac{{\tau_{n}} {u^{x}} {u^{y}} \left({\nu^{x}} {u^{x}} {u^{y}}-{\nu^{y}} \left(\left(u^{x}\right)^2+1\right)\right)}{\left(u^{\tau}\right)^2}
\end{align*}
\begin{align*}
(A^{L}_y)_{11} &= {(\partial_{T} n)} {u^{y}} \\
(A^{L}_y)_{13} &= {n} \\
(A^{L}_y)_{22} &= \frac{{\tau_{n}} {u^{x}} {u^{y}} \left({\nu^{y}} {u^{x}} {u^{y}}-{\nu^{x}} \left(\left(u^{y}\right)^2+1\right)\right)}{\left(u^{\tau}\right)^2} \\
(A^{L}_y)_{23} &= -\frac{{\tau_{n}} {u^{x}} {u^{y}} \left(-{\nu^{x}} {u^{x}} {u^{y}}+{\nu^{y}} \left(u^{x}\right)^2+{\nu^{y}}\right)}{\left(u^{\tau}\right)^2} \\
(A^{L}_y)_{32} &= -\frac{{\tau_{n}} \left(u^{y}\right)^2 \left({\nu^{x}} \left(u^{y}\right)^2+{\nu^{x}}-{\nu^{y}} {u^{x}} {u^{y}}\right)}{\left(u^{\tau}\right)^2} \\
(A^{L}_y)_{33} &= \frac{{\tau_{n}} \left(u^{y}\right)^2 \left({\nu^{x}} {u^{x}} {u^{y}}-{\nu^{y}} \left(\left(u^{x}\right)^2+1\right)\right)}{\left(u^{\tau}\right)^2}
\end{align*}
\begin{align*}
S^{L}_{1} &= \frac{{n} \left(u^{\tau}\right)^2+{\nu^{x}} {u^{x}}+{\nu^{y}} {u^{y}}}{{\tau} {u^{\tau}}} \\
S^{L}_{2} &= {\nu^{x}} \\
S^{L}_{3} &= {\nu^{y}}
\end{align*}
\subsection{Checks}

{\small
\begin{longtable}{@{}>{\raggedright\arraybackslash}p{0.27\textwidth}>{\raggedright\arraybackslash}p{0.60\textwidth}l@{}}
\toprule
id & what it tests & \\
\midrule
\endhead
\texttt{H-\allowbreak{}G0cyl} & cylindrical Milne: $\nabla_{\lambda}g_{\mu\nu}$ =\allowbreak{} 0 & PASS \\
\texttt{H-\allowbreak{}G0car} & Cartesian   Milne: $\nabla_{\lambda}g_{\mu\nu}$ =\allowbreak{} 0 & PASS \\
\texttt{H-\allowbreak{}G1cyl} & cylindrical: $\nabla_{\mu}V^{\mu} = \partial_{\tau}X + X/\tau$ for $V^{\mu} = (X(\tau), 0, 0, 0)$ & PASS \\
\texttt{H-\allowbreak{}G1car} & Cartesian:   $\nabla_{\mu}V^{\mu} = \partial_{\tau}X + X/\tau$ & PASS \\
\texttt{H-\allowbreak{}G2cyl-\allowbreak{}uu} & $u\cdot u$ =\allowbreak{} -1 & PASS \\
\texttt{H-\allowbreak{}G2cyl-\allowbreak{}un} & $u\cdot\nu$ =\allowbreak{} 0 & PASS \\
\texttt{H-\allowbreak{}G3cyl} & cylindrical system is quasi-linear & PASS \\
\texttt{H-\allowbreak{}G4-\allowbreak{}At} & cylindrical $A_{t}\times (u^{\tau})^{3} r\tau$ =\allowbreak{} Fluidum's shipped $A_{t}$ & PASS \\
\texttt{H-\allowbreak{}G4-\allowbreak{}Ax} & cylindrical $A_{x}\times (u^{\tau})^{3} r\tau$ =\allowbreak{} Fluidum's shipped $A_{x}$ & PASS \\
\texttt{H-\allowbreak{}G4-\allowbreak{}src} & cylindrical S $\times (u^{\tau})^{3} r\tau$ =\allowbreak{} Fluidum's shipped source & PASS \\
\texttt{H-\allowbreak{}G4src} & cylindrical $A_{t}, A_{x}, S\times (u^{\tau})^{3} r\tau$ =\allowbreak{} 1d\_\allowbreak{}viscous\_\allowbreak{}HQ\_\allowbreak{}cilindrical\_\allowbreak{}BG.jl read from the file & PASS \\
\texttt{H-\allowbreak{}G2car-\allowbreak{}uu} & $u\cdot u$ =\allowbreak{} -1 & PASS \\
\texttt{H-\allowbreak{}G2car-\allowbreak{}un} & $u\cdot\nu$ =\allowbreak{} 0 & PASS \\
\texttt{H-\allowbreak{}G3car} & Cartesian system is quasi-linear & PASS \\
\texttt{H-\allowbreak{}G5} & slab limit ($u^{y}\to 0$,\allowbreak{} $\partial_{y}\to 0$) decouples $\nu^{y}$ from $(\alpha, \nu^{x}$) & PASS \\
\texttt{H-\allowbreak{}G6a} & the $(\alpha, \nu^{x}, \nu^{y})$ columns of the 10-field rows =\allowbreak{} the background-mode matrices & PASS \\
\texttt{H-\allowbreak{}G6b} & $S^{L}$ + (medium columns).(background derivatives) =\allowbreak{} S: same system,\allowbreak{} re-sorted & PASS \\
\bottomrule
\end{longtable}
}
17 passed, 0 failed (17 checks evaluated).

Re-run with \texttt{wolframscript -file deriv/derive\_\allowbreak{}hq\_\allowbreak{}2p1d.wls}.


\clearpage
% GENERATED by deriv/render.wl from the derivation script.  Do not edit.

\section{Medium, charm and the two-way coupled system}

Script: \texttt{deriv/derive\_two\_way.wls}, with \texttt{derive\_two\_way\_gates.wls} and \texttt{derive\_two\_way\_export.wls}

Writes: the charm and two-way matrices of \texttt{src/Matrix} (one function per system, listed below)

\subsection{Setup}

Three charts, boost invariant: 0+1D Bjorken (fields of $\tau$ only), 1+1D cylindrical, 2+1D Cartesian Milne:

\begin{gather*}
x^{\mu} = ({\tau}, {x}, {y}, {\eta}),\qquad g_{\mu\nu} = \mathrm{diag}(-1, 1, 1, {\tau}^2),\quad u^{\mu} = (1, 0, 0, 0) \\
x^{\mu} = ({\tau}, {r}, {\phi}, {\eta}),\qquad g_{\mu\nu} = \mathrm{diag}(-1, 1, {r}^2, {\tau}^2) \\
x^{\mu} = ({\tau}, {x}, {y}, {\eta}),\qquad g_{\mu\nu} = \mathrm{diag}(-1, 1, 1, {\tau}^2)
\end{gather*}
Medium, with the EOS $P(T)$ kept symbolic:

\begin{gather*}
T_{m}^{\mu\nu} = (e + P + \Pi)\,u^{\mu}u^{\nu} + (P + \Pi)\,g^{\mu\nu} + \pi^{\mu\nu},\qquad e = T P'(T) - P(T) \\
\tau_{\pi}\,\Delta^{\mu\nu}_{\alpha\beta}D\pi^{\alpha\beta} + \pi^{\mu\nu} + 2\eta\,\sigma^{\mu\nu} + \delta_{\pi\pi}\,\pi^{\mu\nu}\theta = 0,\qquad \tau_{\Pi}D\Pi + \Pi + \zeta\theta + \delta_{\Pi\Pi}\Pi\theta = 0
\end{gather*}
($D$, $\theta$, $\Delta^{\mu\nu}$, $\Delta^{\mu\nu}_{\alpha\beta}$, $\sigma^{\mu\nu}$ as in the viscous section.) Charm, with $n = n(T, \alpha)$, $h = h(T)$, $P_{0} = nT$ and $\delta M^{\mu\nu} = \pi_{Q}^{\mu\nu} + \Pi_{Q}\Delta^{\mu\nu}$:

\begin{gather*}
\nabla_{\mu}(n\,u^{\mu} + \nu^{\mu}) = 0,\qquad \frac{D_{s}}{T}\,\Delta^{\mu}{}_{\lambda}\nabla_{\nu}T_{Q}^{\nu\lambda} + \nu^{\mu} = 0 \\
T_{Q}^{\mu\nu} = (n h - P_{0})\,u^{\mu}u^{\nu} + P_{0}\,\Delta^{\mu\nu} + h\,(u^{\mu}\nu^{\nu} + \nu^{\mu}u^{\nu}) \\
S^{\lambda\alpha\beta} = I_{31}\,(u^{\lambda}\Delta^{\alpha\beta} + 2\ \text{perm.}) + \mathcal{A}\,(u^{\lambda}u^{\alpha}\nu^{\beta} + 2\ \text{perm.}) + \mathcal{B}\,(\Delta^{\lambda\alpha}\nu^{\beta} + 2\ \text{perm.}) + \mathcal{C}\,(u^{\lambda}\delta M^{\alpha\beta} + 2\ \text{perm.}) \\
\mathcal{E}^{\mu\nu} = \Delta^{\mu}{}_{\alpha}\Delta^{\nu}{}_{\beta}\nabla_{\lambda}S^{\lambda\alpha\beta}:\qquad \frac{D_{s}}{2T}\Bigl(\mathcal{E}^{\mu\nu} - \tfrac{1}{3}\Delta^{\mu\nu}\mathcal{E}^{\lambda}{}_{\lambda}\Bigr) + \pi_{Q}^{\mu\nu} = 0,\qquad \frac{D_{s}}{2T}\,\tfrac{1}{3}\mathcal{E}^{\lambda}{}_{\lambda} + \Pi_{Q} = 0
\end{gather*}
``2 perm.'' are the two other placements of the lone index. $I_{31}(T, \alpha)$, $\mathcal{A}(T)$, $\mathcal{B}(T)$, $\mathcal{C}(T)$ stay symbolic (Boltzmann: $I_{31} = nTh$, $\mathcal{A} = m^{2} + 5Th$, $\mathcal{B} = Th$, $\mathcal{C} = m K_{4}/K_{3}$). In S1--S5 the second moment $\delta M$ is evolved but does not enter $T_{Q}$. In 2+1D the traceless rows also carry

\[
-(1 - w_{\omega})\,\tau_{M}W^{\mu\nu},\qquad \tau_{M} = \frac{D_{s}\mathcal{C}}{2T},\qquad W^{\mu\nu} = \pi_{Q}^{\mu}{}_{\lambda}\omega^{\lambda\nu} + \pi_{Q}^{\nu}{}_{\lambda}\omega^{\lambda\mu},\qquad \omega^{\mu\nu} = \tfrac{1}{2}(\Delta^{\mu\lambda}\nabla_{\lambda}u^{\nu} - \Delta^{\nu\lambda}\nabla_{\lambda}u^{\mu})
\]
so $w_{\omega} = 1$ is the covariant row, which is exported; $w_{\omega} = 0$ removes the corotation coupling. The five systems (S1, S2 and every 0+1D system are derived for the checks; S3, S4, S5 in 1+1D and 2+1D are exported, with $\delta_{\Pi\Pi} = 0$):

\begin{gather*}
\text{S1}:\ \nabla_{\mu}T_{m}^{\mu\nu} = 0,\ \pi = \Pi = 0\qquad \text{S2}:\ \text{S1} + \text{the }\pi, \Pi\text{ equations} \\
\text{S3}:\ \text{charm number} + \text{first moment, background }T, u\qquad \text{S4}:\ \text{S3} + \text{second moment, background }T, u \\
\text{S5}:\ \nabla_{\mu}(T_{m} + T_{Q})^{\mu\nu} = 0,\ \text{the }\pi, \Pi\text{ equations, the charm rows of S4; all fields evolved} \\
\text{S6}:\ \nabla_{\mu}(T_{m} + T_{Q} + \delta M)^{\mu\nu} = 0,\ \frac{D_{s}}{T}\Delta^{\mu}{}_{\lambda}\nabla_{\nu}(T_{Q} + \delta M)^{\nu\lambda} + \nu^{\mu} = 0,\ \text{the rest as S5}
\end{gather*}
Conformal: $P = a_{\mathrm{SB}}T^{4}$, no bulk. Source on the left, $A_t\,\partial_{\tau}\phi + A_r\,\partial_{r}\phi + S = 0$ in 1+1D and $A_t\,\partial_{\tau}\phi + A_x\,\partial_{x}\phi + A_y\,\partial_{y}\phi + S = 0$ in 2+1D.

\subsubsection{Fields}

Medium, charm first moment, charm second moment, chart by chart:

\begin{gather*}
\text{0+1D}:\ ({T}, {\pi^{\eta}{}_{\eta}}, {\Pi}),\ ({\alpha}),\ ({\pi_{Q}{}^{\eta}{}_{\eta}}, {\Pi_{Q}}) \\
\text{1+1D}:\ ({T}, {u^{r}}, {\pi^{\phi}{}_{\phi}}, {\pi^{\eta}{}_{\eta}}, {\Pi}),\ ({\alpha}, {\nu^{r}}),\ ({\pi_{Q}{}^{\phi}{}_{\phi}}, {\pi_{Q}{}^{\eta}{}_{\eta}}, {\Pi_{Q}}) \\
\text{2+1D}:\ ({T}, {u^{x}}, {u^{y}}, {\pi^{yy}}, {\pi^{\eta}{}_{\eta}}, {\pi^{xy}}, {\Pi}),\ ({\alpha}, {\nu^{x}}, {\nu^{y}}) \\
({\pi_{Q}^{yy}}, {\pi_{Q}{}^{\eta}{}_{\eta}}, {\pi_{Q}^{xy}}, {\Pi_{Q}})
\end{gather*}
Rows, in the same order: the fluid rows $\nabla_{\mu}T^{\mu\tau}$, $\nabla_{\mu}T^{\mu i}$ ($T_{m} + T_{Q}$ in S5); the shear rows $E^{\eta\eta}$ (0+1D), $E^{\phi\phi}$, $E^{\eta\eta}$ (1+1D), $E^{yy}$, $E^{\eta\eta}$, $E^{xy}$ (2+1D); the bulk row; $\nabla_{\mu}N^{\mu}$; the current components; the traceless second-moment rows, components $\eta\eta$ (0+1D), $\phi\phi$, $\eta\eta$ (1+1D), $yy$, $\eta\eta$, $xy$ (2+1D), as the medium's shear rows; the trace row.

In 0+1D the tensors are diagonal and $\nu = 0$:

\[
\pi^{\mu\nu} = \mathrm{diag}(0, -\frac{{\pi^{\eta}{}_{\eta}}}{2}, -\frac{{\pi^{\eta}{}_{\eta}}}{2}, \frac{{\pi^{\eta}{}_{\eta}}}{{\tau}^2}),\qquad \pi_{Q}^{\mu\nu} = \mathrm{diag}(0, -\frac{{\pi_{Q}{}^{\eta}{}_{\eta}}}{2}, -\frac{{\pi_{Q}{}^{\eta}{}_{\eta}}}{2}, \frac{{\pi_{Q}{}^{\eta}{}_{\eta}}}{{\tau}^2})
\]
In 1+1D the medium $\pi^{\mu\nu}$ is as in the viscous section, and so is the charm stress: $\pi_{Q}{}^{\phi}{}_{\phi}$ and $\pi_{Q}{}^{\eta}{}_{\eta}$ are independent, the radial component follows from the trace. With the orthonormal triad transverse to $u$,

\begin{gather*}
\pi_{Q}^{\mu\nu} = -(\pi_{Q}{}^{\phi}{}_{\phi} + \pi_{Q}{}^{\eta}{}_{\eta})\,\ell^{\mu}\ell^{\nu} + \pi_{Q}{}^{\phi}{}_{\phi}\,\hat\phi^{\mu}\hat\phi^{\nu} + \pi_{Q}{}^{\eta}{}_{\eta}\,\hat\eta^{\mu}\hat\eta^{\nu} \\
\ell^{\mu} = ({u^{r}}, {u^{\tau}}, 0, 0),\quad \hat\phi^{\mu} = (0, 0, \frac{1}{{r}}, 0),\quad \hat\eta^{\mu} = (0, 0, 0, \frac{1}{{\tau}})
\end{gather*}
In 2+1D the medium $\pi^{\mu\nu}$ is as in the viscous section, and so is the charm stress: $\pi_{Q}^{yy}$, $\pi_{Q}{}^{\eta}{}_{\eta}$, $\pi_{Q}^{xy}$ are independent, the $\tau$ row follows from $u_{\mu}\pi_{Q}^{\mu\nu} = 0$ and

\[
\pi_{Q}^{xx} = -\frac{-2 {\pi_{Q}^{xy}} {u^{x}} {u^{y}}+{\pi_{Q}^{yy}} \left(u^{x}\right)^2+{\pi_{Q}^{yy}}+{\pi_{Q}{}^{\eta}{}_{\eta}} \left(u^{x}\right)^2+{\pi_{Q}{}^{\eta}{}_{\eta}} \left(u^{y}\right)^2+{\pi_{Q}{}^{\eta}{}_{\eta}}}{\left(u^{y}\right)^2+1}
\]
\subsection{Steps}

\subsubsection{S3: charm first moment}

\[
\Gamma^{{\tau}}{}_{{\eta}{\eta}} = {\tau},\quad \Gamma^{{r}}{}_{{\phi}{\phi}} = -{r},\quad \Gamma^{{\phi}}{}_{{r}{\phi}} = \frac{1}{{r}},\quad \Gamma^{{\eta}}{}_{{\tau}{\eta}} = \frac{1}{{\tau}}
\]
\begin{align*}
N^{{\tau}} &= \frac{{n} \left(u^{r}\right)^2+{n}+{\nu^{r}} {u^{r}}}{{u^{\tau}}} \\
N^{{r}} &= {n} {u^{r}}+{\nu^{r}} \\
E_{\alpha} = \nabla_{\mu}N^{\mu} &= {(\partial_{r} \nu^{r})} + {(\partial_{r} u^{r})} {n} + {u^{r}} ({(\partial_{\alpha} n)} {(\partial_{r} \alpha)}+{(\partial_{T} n)} {(\partial_{r} T)}) + \frac{{n} {u^{r}}+{\nu^{r}}}{{r}} - \frac{{(\partial_{\tau} u^{r})} {\nu^{r}} \left(u^{r}\right)^2}{\left(u^{\tau}\right)^3}\\
  &\quad + \frac{{(\partial_{\tau} u^{r})} {\nu^{r}}}{{u^{\tau}}} + \frac{{(\partial_{\tau} \nu^{r})} {u^{r}}}{{u^{\tau}}} + \frac{{(\partial_{\tau} u^{r})} {n} {u^{r}}}{{u^{\tau}}} + \frac{{n} \left(u^{r}\right)^2+{n}+{\nu^{r}} {u^{r}}}{{\tau} {u^{\tau}}}\\
  &\quad + {u^{\tau}} ({(\partial_{\alpha} n)} {(\partial_{\tau} \alpha)}+{(\partial_{T} n)} {(\partial_{\tau} T)})
\end{align*}
\begin{align*}
T_{Q}^{{\tau}{\tau}} &= {h} \left({n} \left(u^{r}\right)^2+{n}+2 {\nu^{r}} {u^{r}}\right)-{n} {T} \\
T_{Q}^{{\tau}{r}} &= \frac{{h} \left({n} {u^{r}} \left(u^{\tau}\right)^2+2 {\nu^{r}} \left(u^{r}\right)^2+{\nu^{r}}\right)}{{u^{\tau}}} \\
T_{Q}^{{r}{r}} &= {n} \left({h} \left(u^{r}\right)^2+{T}\right)+2 {h} {\nu^{r}} {u^{r}} \\
T_{Q}^{{\phi}{\phi}} &= \frac{{n} {T}}{{r}^2} \\
T_{Q}^{{\eta}{\eta}} &= \frac{{n} {T}}{{\tau}^2}
\end{align*}
\begin{align*}
\nabla_{\mu}T_{Q}^{\mu {\tau}} &= {(\partial_{{\tau}} T_{Q}^{{\tau}{\tau}})}+{(\partial_{{r}} T_{Q}^{{\tau}{r}})}+\frac{{T_{Q}^{{\tau}{r}}}}{{r}}+\frac{{T_{Q}^{{\tau}{\tau}}}}{{\tau}}+{\tau} {T_{Q}^{{\eta}{\eta}}} \\
\nabla_{\mu}T_{Q}^{\mu {r}} &= {(\partial_{{\tau}} T_{Q}^{{\tau}{r}})}+{(\partial_{{r}} T_{Q}^{{r}{r}})}+\frac{{T_{Q}^{{r}{r}}}}{{r}}-{r} {T_{Q}^{{\phi}{\phi}}}+\frac{{T_{Q}^{{\tau}{r}}}}{{\tau}}
\end{align*}
\begin{align*}
\nabla_{\mu}T_{Q}^{\mu\tau} &= \frac{1}{{r} {\tau} \left(u^{\tau}\right)^3}\Bigl({r} {\tau} \left(u^{\tau}\right)^2\Bigl(- {u^{\tau}} ({(\partial_{\alpha} n)} {(\partial_{\tau} \alpha)} {T}+{(\partial_{\tau} T)} ({(\partial_{T} n)} {T}+{n}))\\
  &\quad + {h'(T)}\Bigl({(\partial_{r} T)} \left({n} \left(u^{r}\right)^3+{n} {u^{r}}+2 {\nu^{r}} \left(u^{r}\right)^2+{\nu^{r}}\right)\\
  &\quad + {(\partial_{\tau} T)} {u^{\tau}} \left({n} \left(u^{r}\right)^2+{n}+2 {\nu^{r}} {u^{r}}\right)\Bigr)\Bigr)\\
  &\quad + {h}\Bigl(\left(u^{\tau}\right)^2\Bigl({n} {\tau} {u^{r}} + {(\partial_{\alpha} n)} {(\partial_{r} \alpha)} {r} {\tau} {u^{r}} + {(\partial_{T} n)} {(\partial_{r} T)} {r} {\tau} {u^{r}}\\
  &\quad + {n} {\tau} \left(u^{r}\right)^3 + {(\partial_{\alpha} n)} {(\partial_{r} \alpha)} {r} {\tau} \left(u^{r}\right)^3 + {(\partial_{T} n)} {(\partial_{r} T)} {r} {\tau} \left(u^{r}\right)^3\\
  &\quad + {(\partial_{r} \nu^{r})} {r} {\tau} \left(2 \left(u^{r}\right)^2+1\right) + {(\partial_{r} u^{r})} {n} {r} {\tau} \left(2 \left(u^{r}\right)^2+1\right)\\
  &\quad + {n} {r} {u^{\tau}} + {(\partial_{\alpha} n)} {(\partial_{\tau} \alpha)} {r} {\tau} {u^{\tau}} + {(\partial_{T} n)} {(\partial_{\tau} T)} {r} {\tau} {u^{\tau}}\\
  &\quad + 2 {(\partial_{\tau} \nu^{r})} {r} {\tau} {u^{r}} {u^{\tau}} + 2 {(\partial_{\tau} u^{r})} {n} {r} {\tau} {u^{r}} {u^{\tau}} + {n} {r} \left(u^{r}\right)^2 {u^{\tau}}\\
  &\quad + {(\partial_{\alpha} n)} {(\partial_{\tau} \alpha)} {r} {\tau} \left(u^{r}\right)^2 {u^{\tau}} + {(\partial_{T} n)} {(\partial_{\tau} T)} {r} {\tau} \left(u^{r}\right)^2 {u^{\tau}}\Bigr)\\
  &\quad + {\nu^{r}}\Bigl({(\partial_{r} u^{r})} {r} {\tau} {u^{r}} \left(2 \left(u^{r}\right)^2+3\right) + 2 {r} {u^{r}} \left(u^{\tau}\right)^3\\
  &\quad + {\tau} \left(u^{\tau}\right)^2 \left(2 {(\partial_{\tau} u^{r})} {r} {u^{\tau}}+2 \left(u^{r}\right)^2+1\right)\Bigr)\Bigr)\Bigr) \\
\nabla_{\mu}T_{Q}^{\mu r} &= \frac{1}{{r} {\tau} \left(u^{\tau}\right)^3}\Bigl({r} {\tau} \left(u^{\tau}\right)^2\Bigl({u^{\tau}} ({(\partial_{\alpha} n)} {(\partial_{r} \alpha)} {T}+{(\partial_{r} T)} ({(\partial_{T} n)} {T}+{n}))\\
  &\quad + {h'(T)}\Bigl({(\partial_{\tau} T)} \left({n} \left(u^{r}\right)^3+{n} {u^{r}}+2 {\nu^{r}} \left(u^{r}\right)^2+{\nu^{r}}\right)\\
  &\quad + {(\partial_{r} T)} {u^{r}} {u^{\tau}} ({n} {u^{r}}+2 {\nu^{r}})\Bigr)\Bigr)\\
  &\quad + {h}\Bigl({\nu^{r}}\Bigl({(\partial_{\tau} u^{r})} {r} {\tau} {u^{r}} \left(2 \left(u^{r}\right)^2+3\right) + 2 {\tau} {u^{r}} \left(u^{\tau}\right)^3\\
  &\quad + {r} \left(u^{\tau}\right)^2 \left(2 {(\partial_{r} u^{r})} {\tau} {u^{\tau}}+2 \left(u^{r}\right)^2+1\right)\Bigr)\\
  &\quad + \left(u^{\tau}\right)^2\Bigl({(\partial_{\tau} \nu^{r})} {r} {\tau} \left(2 \left(u^{r}\right)^2+1\right) + {(\partial_{\tau} u^{r})} {n} {r} {\tau} \left(2 \left(u^{r}\right)^2+1\right)\\
  &\quad + {n} {u^{r}} \left({r} \left(2 {(\partial_{r} u^{r})} {\tau} {u^{\tau}}+\left(u^{r}\right)^2+1\right)+{\tau} {u^{r}} {u^{\tau}}\right)\\
  &\quad + {r} {\tau} {u^{r}}\Bigl(2 {(\partial_{r} \nu^{r})} {u^{\tau}}\\
  &\quad + {(\partial_{\alpha} n)} \left({(\partial_{r} \alpha)} {u^{r}} {u^{\tau}}+{(\partial_{\tau} \alpha)} \left(u^{r}\right)^2+{(\partial_{\tau} \alpha)}\right)\\
  &\quad + {(\partial_{T} n)} \left({(\partial_{r} T)} {u^{r}} {u^{\tau}}+{(\partial_{\tau} T)} \left(u^{r}\right)^2+{(\partial_{\tau} T)}\right)\Bigr)\Bigr)\Bigr)\Bigr) \\
\Delta^{r}{}_{\tau} &= -{u^{r}} {u^{\tau}} \\
\Delta^{r}{}_{r} &= \left(u^{\tau}\right)^2
\end{align*}
\[
E_{\nu} = \frac{D_{s}}{T}\bigl(\Delta^{r}{}_{\tau}\nabla_{\mu}T_{Q}^{\mu\tau} + \Delta^{r}{}_{r}\nabla_{\mu}T_{Q}^{\mu r}\bigr) + \nu^{r}
\]
\begin{align*}
E_{\nu} &= {\nu^{r}}\\
  &\quad + \frac{1}{{T} {u^{\tau}}}{D_{s}}\Bigl({h'(T)} {\nu^{r}} \left({(\partial_{r} T)} {u^{r}} {u^{\tau}}+{(\partial_{\tau} T)} \left(u^{r}\right)^2+{(\partial_{\tau} T)}\right)\\
  &\quad + \left(u^{\tau}\right)^2\Bigl({(\partial_{\tau} T)} {u^{r}} ({(\partial_{T} n)} {T}+{n}) + {(\partial_{r} T)} {u^{\tau}} ({(\partial_{T} n)} {T}+{n})\\
  &\quad + {(\partial_{\alpha} n)} {T} ({(\partial_{r} \alpha)} {u^{\tau}}+{(\partial_{\tau} \alpha)} {u^{r}})\Bigr)\Bigr)\\
  &\quad + \frac{1}{{r} {T} {\tau} \left(u^{\tau}\right)^3}{D_{s}} {h}\Bigl({\nu^{r}} {\tau} {u^{r}} \left(u^{\tau}\right)^3\\
  &\quad + {r} {\tau} \left(u^{\tau}\right)^2\Bigl({u^{r}} {u^{\tau}} ({(\partial_{r} \nu^{r})}+{(\partial_{r} u^{r})} {n}) + {(\partial_{\tau} \nu^{r})} \left(u^{\tau}\right)^2\\
  &\quad + {(\partial_{\tau} u^{r})} {n} \left(u^{\tau}\right)^2\Bigr)\\
  &\quad + {\nu^{r}} {r}\Bigl(1 + \left(u^{r}\right)^4 + 2 {(\partial_{r} u^{r})} {\tau} {u^{\tau}} + {(\partial_{\tau} u^{r})} {\tau} {u^{r}} \left(u^{\tau}\right)^2\\
  &\quad + \left(u^{r}\right)^2 ({(\partial_{r} u^{r})} {\tau} {u^{\tau}}+2)\Bigr)\Bigr)
\end{align*}
\subsubsection{S4: charm second moment}

$\mathcal{A}$ enters $S^{\lambda\alpha\beta}$ only through $u^{\lambda}u^{\alpha}\nu^{\beta}$: after the projection its derivative has a vanishing coefficient (G-aQ).

\[
\nabla_{\lambda}S^{\lambda\alpha\beta} = \partial_{\lambda}S^{\lambda\alpha\beta} + \Gamma^{\lambda}{}_{\lambda\sigma}S^{\sigma\alpha\beta} + \Gamma^{\alpha}{}_{\lambda\sigma}S^{\lambda\sigma\beta} + \Gamma^{\beta}{}_{\lambda\sigma}S^{\lambda\alpha\sigma}
\]
\begin{align*}
\mathcal{E}^{\phi\phi} &= \frac{1}{{r}^3 {\tau} \left(u^{\tau}\right)^3}\Bigl(- {\mathcal{B}} {(\partial_{\tau} u^{r})} {\nu^{r}} {r} {\tau} \left(u^{r}\right)^2 + {\mathcal{B}} {(\partial_{\tau} u^{r})} {\nu^{r}} {r} {\tau} \left(u^{\tau}\right)^2 + {\mathcal{B}} {\nu^{r}} {r} {u^{r}} \left(u^{\tau}\right)^2\\
  &\quad + {\mathcal{B}} {(\partial_{\tau} \nu^{r})} {r} {\tau} {u^{r}} \left(u^{\tau}\right)^2 + {(\partial_{\tau} u^{r})} {I_{31}} {r} {\tau} {u^{r}} \left(u^{\tau}\right)^2 + {\mathcal{B}'(T)} {(\partial_{\tau} T)} {\nu^{r}} {r} {\tau} {u^{r}} \left(u^{\tau}\right)^2\\
  &\quad + {\mathcal{C}} {(\partial_{\tau} u^{r})} {\Pi_{Q}} {r} {\tau} {u^{r}} \left(u^{\tau}\right)^2 + {\mathcal{C}} {(\partial_{\tau} u^{r})} {\pi_{Q}{}^{\phi}{}_{\phi}} {r} {\tau} {u^{r}} \left(u^{\tau}\right)^2 + 3 {\mathcal{B}} {\nu^{r}} {\tau} \left(u^{\tau}\right)^3\\
  &\quad + {\mathcal{B}} {(\partial_{r} \nu^{r})} {r} {\tau} \left(u^{\tau}\right)^3 + {(\partial_{r} u^{r})} {I_{31}} {r} {\tau} \left(u^{\tau}\right)^3 + {\mathcal{B}'(T)} {(\partial_{r} T)} {\nu^{r}} {r} {\tau} \left(u^{\tau}\right)^3\\
  &\quad + {\mathcal{C}} {(\partial_{r} u^{r})} {\Pi_{Q}} {r} {\tau} \left(u^{\tau}\right)^3 + {\mathcal{C}} {(\partial_{r} u^{r})} {\pi_{Q}{}^{\phi}{}_{\phi}} {r} {\tau} \left(u^{\tau}\right)^3 + 3 {I_{31}} {\tau} {u^{r}} \left(u^{\tau}\right)^3\\
  &\quad + 3 {\mathcal{C}} {\Pi_{Q}} {\tau} {u^{r}} \left(u^{\tau}\right)^3 + 3 {\mathcal{C}} {\pi_{Q}{}^{\phi}{}_{\phi}} {\tau} {u^{r}} \left(u^{\tau}\right)^3 + {(\partial_{\alpha} I_{31})} {(\partial_{r} \alpha)} {r} {\tau} {u^{r}} \left(u^{\tau}\right)^3\\
  &\quad + {\mathcal{C}} {(\partial_{r} \Pi_{Q})} {r} {\tau} {u^{r}} \left(u^{\tau}\right)^3 + {\mathcal{C}} {(\partial_{r} \pi_{Q}{}^{\phi}{}_{\phi})} {r} {\tau} {u^{r}} \left(u^{\tau}\right)^3 + {(\partial_{T} I_{31})} {(\partial_{r} T)} {r} {\tau} {u^{r}} \left(u^{\tau}\right)^3\\
  &\quad + {\mathcal{C}'(T)} {(\partial_{r} T)} {\Pi_{Q}} {r} {\tau} {u^{r}} \left(u^{\tau}\right)^3 + {\mathcal{C}'(T)} {(\partial_{r} T)} {\pi_{Q}{}^{\phi}{}_{\phi}} {r} {\tau} {u^{r}} \left(u^{\tau}\right)^3 + {I_{31}} {r} \left(u^{\tau}\right)^4\\
  &\quad + {\mathcal{C}} {\Pi_{Q}} {r} \left(u^{\tau}\right)^4 + {\mathcal{C}} {\pi_{Q}{}^{\phi}{}_{\phi}} {r} \left(u^{\tau}\right)^4 + {(\partial_{\alpha} I_{31})} {(\partial_{\tau} \alpha)} {r} {\tau} \left(u^{\tau}\right)^4 + {\mathcal{C}} {(\partial_{\tau} \Pi_{Q})} {r} {\tau} \left(u^{\tau}\right)^4\\
  &\quad + {\mathcal{C}} {(\partial_{\tau} \pi_{Q}{}^{\phi}{}_{\phi})} {r} {\tau} \left(u^{\tau}\right)^4 + {(\partial_{T} I_{31})} {(\partial_{\tau} T)} {r} {\tau} \left(u^{\tau}\right)^4 + {\mathcal{C}'(T)} {(\partial_{\tau} T)} {\Pi_{Q}} {r} {\tau} \left(u^{\tau}\right)^4\\
  &\quad + {\mathcal{C}'(T)} {(\partial_{\tau} T)} {\pi_{Q}{}^{\phi}{}_{\phi}} {r} {\tau} \left(u^{\tau}\right)^4\Bigr) \\
\mathcal{E}^{\eta\eta} &= \frac{1}{{r} {\tau}^3 \left(u^{\tau}\right)^3}\Bigl(- {\mathcal{B}} {(\partial_{\tau} u^{r})} {\nu^{r}} {r} {\tau} \left(u^{r}\right)^2 + {\mathcal{B}} {(\partial_{\tau} u^{r})} {\nu^{r}} {r} {\tau} \left(u^{\tau}\right)^2 - 2 {\mathcal{B}} {\nu^{r}} {r} {u^{r}} \left(u^{\tau}\right)^2\\
  &\quad + {\mathcal{B}} {(\partial_{\tau} \nu^{r})} {r} {\tau} {u^{r}} \left(u^{\tau}\right)^2 + {(\partial_{\tau} u^{r})} {I_{31}} {r} {\tau} {u^{r}} \left(u^{\tau}\right)^2 + {\mathcal{B}'(T)} {(\partial_{\tau} T)} {\nu^{r}} {r} {\tau} {u^{r}} \left(u^{\tau}\right)^2\\
  &\quad + {\mathcal{C}} {(\partial_{\tau} u^{r})} {r} {\tau} {u^{r}} \left(u^{\tau}\right)^2 ({\Pi_{Q}}+{\pi_{Q}{}^{\eta}{}_{\eta}}) + {\mathcal{B}} {(\partial_{r} \nu^{r})} {r} {\tau} \left(u^{\tau}\right)^3 + {(\partial_{r} u^{r})} {I_{31}} {r} {\tau} \left(u^{\tau}\right)^3\\
  &\quad + {\mathcal{B}'(T)} {(\partial_{r} T)} {\nu^{r}} {r} {\tau} \left(u^{\tau}\right)^3 + {\mathcal{C}} {(\partial_{r} u^{r})} {r} {\tau} \left(u^{\tau}\right)^3 ({\Pi_{Q}}+{\pi_{Q}{}^{\eta}{}_{\eta}})\\
  &\quad + {\mathcal{C}} {r} {\tau} {u^{r}} \left(u^{\tau}\right)^3 ({(\partial_{r} \Pi_{Q})}+{(\partial_{r} \pi_{Q}{}^{\eta}{}_{\eta})})\\
  &\quad + {r} {\tau} {u^{r}} \left(u^{\tau}\right)^3 ({(\partial_{\alpha} I_{31})} {(\partial_{r} \alpha)}+{(\partial_{T} I_{31})} {(\partial_{r} T)})\\
  &\quad + {\mathcal{C}'(T)} {(\partial_{r} T)} {r} {\tau} {u^{r}} \left(u^{\tau}\right)^3 ({\Pi_{Q}}+{\pi_{Q}{}^{\eta}{}_{\eta}})\\
  &\quad + {\tau} \left(u^{\tau}\right)^3 ({\mathcal{B}} {\nu^{r}}+{u^{r}} ({\mathcal{C}} ({\Pi_{Q}}+{\pi_{Q}{}^{\eta}{}_{\eta}})+{I_{31}})) - 2 {I_{31}} {r} \left(u^{\tau}\right)^4\\
  &\quad + {r} {\tau} \left(u^{\tau}\right)^4 ({(\partial_{\alpha} I_{31})} {(\partial_{\tau} \alpha)}+{(\partial_{T} I_{31})} {(\partial_{\tau} T)})\\
  &\quad + {\mathcal{C}'(T)} {(\partial_{\tau} T)} {r} {\tau} \left(u^{\tau}\right)^4 ({\Pi_{Q}}+{\pi_{Q}{}^{\eta}{}_{\eta}})\\
  &\quad - {\mathcal{C}} {r} \left(u^{\tau}\right)^4 (-{\tau} ({(\partial_{\tau} \Pi_{Q})}+{(\partial_{\tau} \pi_{Q}{}^{\eta}{}_{\eta})})+2 {\Pi_{Q}}+2 {\pi_{Q}{}^{\eta}{}_{\eta}})\\
  &\quad + 5 {r} \left(u^{\tau}\right)^2 \left({\mathcal{B}} {\nu^{r}} {u^{r}}+{\mathcal{C}} \left(u^{\tau}\right)^2 ({\Pi_{Q}}+{\pi_{Q}{}^{\eta}{}_{\eta}})+{I_{31}} \left(u^{\tau}\right)^2\right)\Bigr) \\
\mathcal{E}^{\lambda}{}_{\lambda} &= \frac{1}{{r} {\tau} \left(u^{\tau}\right)^3}\Bigl(3 {(\partial_{\alpha} I_{31})} {(\partial_{\tau} \alpha)} {r} {\tau} + 3 {(\partial_{T} I_{31})} {(\partial_{\tau} T)} {r} {\tau} + 2 {\mathcal{A}} {(\partial_{\tau} u^{r})} {\nu^{r}} {r} {\tau} + 5 {\mathcal{B}} {(\partial_{\tau} u^{r})} {\nu^{r}} {r} {\tau}\\
  &\quad + 3 {\mathcal{C}'(T)} {(\partial_{\tau} T)} {\Pi_{Q}} {r} {\tau} + 5 {\mathcal{B}} {\nu^{r}} {r} {u^{r}} + 5 {\mathcal{B}} {(\partial_{\tau} \nu^{r})} {r} {\tau} {u^{r}} + 5 {\mathcal{B}'(T)} {(\partial_{\tau} T)} {\nu^{r}} {r} {\tau} {u^{r}}\\
  &\quad + 6 {(\partial_{\alpha} I_{31})} {(\partial_{\tau} \alpha)} {r} {\tau} \left(u^{r}\right)^2 + 6 {(\partial_{T} I_{31})} {(\partial_{\tau} T)} {r} {\tau} \left(u^{r}\right)^2 + 2 {\mathcal{A}} {(\partial_{\tau} u^{r})} {\nu^{r}} {r} {\tau} \left(u^{r}\right)^2\\
  &\quad + 6 {\mathcal{C}'(T)} {(\partial_{\tau} T)} {\Pi_{Q}} {r} {\tau} \left(u^{r}\right)^2 + 5 {\mathcal{B}} {\nu^{r}} {r} \left(u^{r}\right)^3 + 5 {\mathcal{B}} {(\partial_{\tau} \nu^{r})} {r} {\tau} \left(u^{r}\right)^3\\
  &\quad + 5 {\mathcal{B}'(T)} {(\partial_{\tau} T)} {\nu^{r}} {r} {\tau} \left(u^{r}\right)^3 + 3 {(\partial_{\alpha} I_{31})} {(\partial_{\tau} \alpha)} {r} {\tau} \left(u^{r}\right)^4\\
  &\quad + 3 {(\partial_{T} I_{31})} {(\partial_{\tau} T)} {r} {\tau} \left(u^{r}\right)^4 + 3 {\mathcal{C}'(T)} {(\partial_{\tau} T)} {\Pi_{Q}} {r} {\tau} \left(u^{r}\right)^4 + 5 {\mathcal{B}} {\nu^{r}} {\tau} {u^{\tau}}\\
  &\quad + 5 {\mathcal{B}} {(\partial_{r} \nu^{r})} {r} {\tau} {u^{\tau}} + 5 {\mathcal{B}'(T)} {(\partial_{r} T)} {\nu^{r}} {r} {\tau} {u^{\tau}} + 3 {(\partial_{\alpha} I_{31})} {(\partial_{r} \alpha)} {r} {\tau} {u^{r}} {u^{\tau}}\\
  &\quad + 3 {(\partial_{T} I_{31})} {(\partial_{r} T)} {r} {\tau} {u^{r}} {u^{\tau}} + 2 {\mathcal{A}} {(\partial_{r} u^{r})} {\nu^{r}} {r} {\tau} {u^{r}} {u^{\tau}} + 3 {\mathcal{C}'(T)} {(\partial_{r} T)} {\Pi_{Q}} {r} {\tau} {u^{r}} {u^{\tau}}\\
  &\quad + 5 {\mathcal{B}} {\nu^{r}} {\tau} \left(u^{r}\right)^2 {u^{\tau}} + 5 {\mathcal{B}} {(\partial_{r} \nu^{r})} {r} {\tau} \left(u^{r}\right)^2 {u^{\tau}} + 5 {\mathcal{B}'(T)} {(\partial_{r} T)} {\nu^{r}} {r} {\tau} \left(u^{r}\right)^2 {u^{\tau}}\\
  &\quad + 3 {(\partial_{\alpha} I_{31})} {(\partial_{r} \alpha)} {r} {\tau} \left(u^{r}\right)^3 {u^{\tau}} + 3 {(\partial_{T} I_{31})} {(\partial_{r} T)} {r} {\tau} \left(u^{r}\right)^3 {u^{\tau}}\\
  &\quad + 3 {\mathcal{C}'(T)} {(\partial_{r} T)} {\Pi_{Q}} {r} {\tau} \left(u^{r}\right)^3 {u^{\tau}}\\
  &\quad + 5 {I_{31}} \left(u^{\tau}\right)^2 \left({r} \left({(\partial_{r} u^{r})} {\tau} {u^{\tau}}+{(\partial_{\tau} u^{r})} {\tau} {u^{r}}+\left(u^{r}\right)^2+1\right)+{\tau} {u^{r}} {u^{\tau}}\right)\\
  &\quad + {\mathcal{C}} \left(u^{\tau}\right)^2\Bigl(- 2 {\pi_{Q}{}^{\eta}{}_{\eta}} {r} \left({(\partial_{r} u^{r})} {\tau} {u^{\tau}}+{(\partial_{\tau} u^{r})} {\tau} {u^{r}}-\left(u^{r}\right)^2-1\right)\\
  &\quad + {\tau}\Bigl(- 2 {(\partial_{\tau} u^{r})} {\pi_{Q}{}^{\phi}{}_{\phi}} {r} {u^{r}}\\
  &\quad + {u^{\tau}} (3 {(\partial_{r} \Pi_{Q})} {r} {u^{r}}-2 {(\partial_{r} u^{r})} {\pi_{Q}{}^{\phi}{}_{\phi}} {r}+2 {\pi_{Q}{}^{\phi}{}_{\phi}} {u^{r}})\\
  &\quad + 3 {(\partial_{\tau} \Pi_{Q})} {r} \left(u^{\tau}\right)^2\Bigr)\\
  &\quad + 5 {\Pi_{Q}} \left({r} \left({(\partial_{r} u^{r})} {\tau} {u^{\tau}}+{(\partial_{\tau} u^{r})} {\tau} {u^{r}}+\left(u^{r}\right)^2+1\right)+{\tau} {u^{r}} {u^{\tau}}\right)\Bigr)\Bigr)
\end{align*}
\[
E_{\phi} = \frac{D_{s}}{2T}\Bigl(r^{2}\mathcal{E}^{\phi\phi} - \tfrac{1}{3}\mathcal{E}^{\lambda}{}_{\lambda}\Bigr) + \pi_{Q}{}^{\phi}{}_{\phi},\quad E_{\eta} = \frac{D_{s}}{2T}\Bigl(\tau^{2}\mathcal{E}^{\eta\eta} - \tfrac{1}{3}\mathcal{E}^{\lambda}{}_{\lambda}\Bigr) + \pi_{Q}{}^{\eta}{}_{\eta},\quad E_{\Pi_{Q}} = \frac{D_{s}}{6T}\mathcal{E}^{\lambda}{}_{\lambda} + \Pi_{Q}
\]
\subsubsection{S5: two-way}

\begin{align*}
T_{m}^{{\tau}{\tau}} &= {P'(T)} {T} \left(u^{\tau}\right)^2+\left(u^{r}\right)^2 ({\Pi}-{\pi^{\eta}{}_{\eta}}-{\pi^{\phi}{}_{\phi}})-{P} \\
T_{m}^{{\tau}{r}} &= {u^{r}} {u^{\tau}} ({P'(T)} {T}+{\Pi}-{\pi^{\eta}{}_{\eta}}-{\pi^{\phi}{}_{\phi}}) \\
T_{m}^{{r}{r}} &= \left(u^{r}\right)^2 ({P'(T)} {T}+{\Pi})+{\Pi}-\left(u^{\tau}\right)^2 ({\pi^{\eta}{}_{\eta}}+{\pi^{\phi}{}_{\phi}})+{P} \\
T_{m}^{{\phi}{\phi}} &= \frac{{\Pi}+{\pi^{\phi}{}_{\phi}}+{P}}{{r}^2} \\
T_{m}^{{\eta}{\eta}} &= \frac{{\Pi}+{\pi^{\eta}{}_{\eta}}+{P}}{{\tau}^2}
\end{align*}
\[
E_{\tau} = \nabla_{\mu}T_{m}^{\mu\tau} + \nabla_{\mu}T_{Q}^{\mu\tau},\qquad E_{r} = \nabla_{\mu}T_{m}^{\mu r} + \nabla_{\mu}T_{Q}^{\mu r}
\]
The shear and bulk rows as in the viscous section, the charm rows as S4.

\begin{align*}
E_{\tau} &= \frac{1}{{r} {\tau} \left(u^{\tau}\right)^3}\Bigl(\left(u^{\tau}\right)^2\Bigl({\Pi} {\tau} {u^{r}} - {\pi^{\eta}{}_{\eta}} {\tau} {u^{r}} - {\pi^{\phi}{}_{\phi}} {\tau} {u^{r}} + {(\partial_{r} \Pi)} {r} {\tau} {u^{r}} - {(\partial_{r} \pi^{\eta}{}_{\eta})} {r} {\tau} {u^{r}}\\
  &\quad - {(\partial_{r} \pi^{\phi}{}_{\phi})} {r} {\tau} {u^{r}} + {(\partial_{r} T)} {P'(T)} {r} {\tau} {u^{r}} + {P'(T)} {T} {\tau} {u^{r}}\\
  &\quad + {(\partial_{r} T)} {P''(T)} {r} {T} {\tau} {u^{r}} + {\Pi} {\tau} \left(u^{r}\right)^3 - {\pi^{\eta}{}_{\eta}} {\tau} \left(u^{r}\right)^3 - {\pi^{\phi}{}_{\phi}} {\tau} \left(u^{r}\right)^3\\
  &\quad + {(\partial_{r} \Pi)} {r} {\tau} \left(u^{r}\right)^3 - {(\partial_{r} \pi^{\eta}{}_{\eta})} {r} {\tau} \left(u^{r}\right)^3 - {(\partial_{r} \pi^{\phi}{}_{\phi})} {r} {\tau} \left(u^{r}\right)^3\\
  &\quad + {(\partial_{r} T)} {P'(T)} {r} {\tau} \left(u^{r}\right)^3 + {P'(T)} {T} {\tau} \left(u^{r}\right)^3\\
  &\quad + {(\partial_{r} T)} {P''(T)} {r} {T} {\tau} \left(u^{r}\right)^3\\
  &\quad + {(\partial_{r} u^{r})} {r} {\tau} \left(2 \left(u^{r}\right)^2+1\right) ({P'(T)} {T}+{\Pi}-{\pi^{\eta}{}_{\eta}}-{\pi^{\phi}{}_{\phi}}) + {\Pi} {r} {u^{\tau}}\\
  &\quad + {\pi^{\eta}{}_{\eta}} {r} {u^{\tau}} + {P'(T)} {r} {T} {u^{\tau}} - {(\partial_{\tau} T)} {n} {r} {\tau} {u^{\tau}} - {(\partial_{\alpha} n)} {(\partial_{\tau} \alpha)} {r} {T} {\tau} {u^{\tau}}\\
  &\quad - {(\partial_{T} n)} {(\partial_{\tau} T)} {r} {T} {\tau} {u^{\tau}} + {P''(T)} {(\partial_{\tau} T)} {r} {T} {\tau} {u^{\tau}} + 2 {(\partial_{\tau} u^{r})} {\Pi} {r} {\tau} {u^{r}} {u^{\tau}}\\
  &\quad - 2 {(\partial_{\tau} u^{r})} {\pi^{\eta}{}_{\eta}} {r} {\tau} {u^{r}} {u^{\tau}} - 2 {(\partial_{\tau} u^{r})} {\pi^{\phi}{}_{\phi}} {r} {\tau} {u^{r}} {u^{\tau}}\\
  &\quad + 2 {P'(T)} {(\partial_{\tau} u^{r})} {r} {T} {\tau} {u^{r}} {u^{\tau}} + {\Pi} {r} \left(u^{r}\right)^2 {u^{\tau}} - {\pi^{\eta}{}_{\eta}} {r} \left(u^{r}\right)^2 {u^{\tau}}\\
  &\quad - {\pi^{\phi}{}_{\phi}} {r} \left(u^{r}\right)^2 {u^{\tau}} + {P'(T)} {r} {T} \left(u^{r}\right)^2 {u^{\tau}} + {(\partial_{\tau} \Pi)} {r} {\tau} \left(u^{r}\right)^2 {u^{\tau}}\\
  &\quad - {(\partial_{\tau} \pi^{\eta}{}_{\eta})} {r} {\tau} \left(u^{r}\right)^2 {u^{\tau}} - {(\partial_{\tau} \pi^{\phi}{}_{\phi})} {r} {\tau} \left(u^{r}\right)^2 {u^{\tau}}\\
  &\quad + {P'(T)} {(\partial_{\tau} T)} {r} {\tau} \left(u^{r}\right)^2 {u^{\tau}} + {P''(T)} {(\partial_{\tau} T)} {r} {T} {\tau} \left(u^{r}\right)^2 {u^{\tau}}\\
  &\quad + {h'(T)} {r} {\tau}\Bigl({(\partial_{r} T)} \left({n} \left(u^{r}\right)^3+{n} {u^{r}}+2 {\nu^{r}} \left(u^{r}\right)^2+{\nu^{r}}\right)\\
  &\quad + {(\partial_{\tau} T)} {u^{\tau}} \left({n} \left(u^{r}\right)^2+{n}+2 {\nu^{r}} {u^{r}}\right)\Bigr)\Bigr)\\
  &\quad + {h}\Bigl(\left(u^{\tau}\right)^2\Bigl({n} {\tau} {u^{r}} + {(\partial_{\alpha} n)} {(\partial_{r} \alpha)} {r} {\tau} {u^{r}} + {(\partial_{T} n)} {(\partial_{r} T)} {r} {\tau} {u^{r}}\\
  &\quad + {n} {\tau} \left(u^{r}\right)^3 + {(\partial_{\alpha} n)} {(\partial_{r} \alpha)} {r} {\tau} \left(u^{r}\right)^3 + {(\partial_{T} n)} {(\partial_{r} T)} {r} {\tau} \left(u^{r}\right)^3\\
  &\quad + {(\partial_{r} \nu^{r})} {r} {\tau} \left(2 \left(u^{r}\right)^2+1\right) + {(\partial_{r} u^{r})} {n} {r} {\tau} \left(2 \left(u^{r}\right)^2+1\right)\\
  &\quad + {n} {r} {u^{\tau}} + {(\partial_{\alpha} n)} {(\partial_{\tau} \alpha)} {r} {\tau} {u^{\tau}} + {(\partial_{T} n)} {(\partial_{\tau} T)} {r} {\tau} {u^{\tau}}\\
  &\quad + 2 {(\partial_{\tau} \nu^{r})} {r} {\tau} {u^{r}} {u^{\tau}} + 2 {(\partial_{\tau} u^{r})} {n} {r} {\tau} {u^{r}} {u^{\tau}} + {n} {r} \left(u^{r}\right)^2 {u^{\tau}}\\
  &\quad + {(\partial_{\alpha} n)} {(\partial_{\tau} \alpha)} {r} {\tau} \left(u^{r}\right)^2 {u^{\tau}} + {(\partial_{T} n)} {(\partial_{\tau} T)} {r} {\tau} \left(u^{r}\right)^2 {u^{\tau}}\Bigr)\\
  &\quad + {\nu^{r}}\Bigl({(\partial_{r} u^{r})} {r} {\tau} {u^{r}} \left(2 \left(u^{r}\right)^2+3\right) + 2 {r} {u^{r}} \left(u^{\tau}\right)^3\\
  &\quad + {\tau} \left(u^{\tau}\right)^2 \left(2 {(\partial_{\tau} u^{r})} {r} {u^{\tau}}+2 \left(u^{r}\right)^2+1\right)\Bigr)\Bigr)\Bigr)
\end{align*}
\subsubsection{S6: fully coupled}

\begin{align*}
\delta M^{{\tau}{\tau}} &= \left(u^{r}\right)^2 ({\Pi_{Q}}-{\pi_{Q}{}^{\eta}{}_{\eta}}-{\pi_{Q}{}^{\phi}{}_{\phi}}) \\
\delta M^{{\tau}{r}} &= {u^{r}} {u^{\tau}} ({\Pi_{Q}}-{\pi_{Q}{}^{\eta}{}_{\eta}}-{\pi_{Q}{}^{\phi}{}_{\phi}}) \\
\delta M^{{r}{r}} &= \left(u^{\tau}\right)^2 ({\Pi_{Q}}-{\pi_{Q}{}^{\eta}{}_{\eta}}-{\pi_{Q}{}^{\phi}{}_{\phi}}) \\
\delta M^{{\phi}{\phi}} &= \frac{{\Pi_{Q}}+{\pi_{Q}{}^{\phi}{}_{\phi}}}{{r}^2} \\
\delta M^{{\eta}{\eta}} &= \frac{{\Pi_{Q}}+{\pi_{Q}{}^{\eta}{}_{\eta}}}{{\tau}^2}
\end{align*}
\begin{align*}
\nabla_{\mu}\delta M^{\mu {\tau}} &= {(\partial_{{\tau}} \delta M^{{\tau}{\tau}})}+{(\partial_{{r}} \delta M^{{\tau}{r}})}+\frac{{\delta M^{{\tau}{r}}}}{{r}}+\frac{{\delta M^{{\tau}{\tau}}}}{{\tau}}+{\tau} {\delta M^{{\eta}{\eta}}} \\
\nabla_{\mu}\delta M^{\mu {r}} &= {(\partial_{{\tau}} \delta M^{{\tau}{r}})}+{(\partial_{{r}} \delta M^{{r}{r}})}+\frac{{\delta M^{{r}{r}}}}{{r}}-{r} {\delta M^{{\phi}{\phi}}}+\frac{{\delta M^{{\tau}{r}}}}{{\tau}}
\end{align*}
\begin{align*}
\nabla_{\mu}\delta M^{\mu\tau} &= \frac{{\Pi_{Q}}+{\pi_{Q}{}^{\eta}{}_{\eta}}}{{\tau}} + 2 {(\partial_{\tau} u^{r})} {\Pi_{Q}} {u^{r}} - 2 {(\partial_{\tau} u^{r})} {u^{r}} ({\pi_{Q}{}^{\eta}{}_{\eta}}+{\pi_{Q}{}^{\phi}{}_{\phi}}) + {(\partial_{\tau} \Pi_{Q})} \left(u^{r}\right)^2\\
  &\quad - \left(u^{r}\right)^2 ({(\partial_{\tau} \pi_{Q}{}^{\eta}{}_{\eta})}+{(\partial_{\tau} \pi_{Q}{}^{\phi}{}_{\phi})}) + \frac{\left(u^{r}\right)^2 ({\Pi_{Q}}-{\pi_{Q}{}^{\eta}{}_{\eta}}-{\pi_{Q}{}^{\phi}{}_{\phi}})}{{\tau}} + \frac{{(\partial_{r} u^{r})} {\Pi_{Q}} \left(u^{r}\right)^2}{{u^{\tau}}}\\
  &\quad - \frac{{(\partial_{r} u^{r})} \left(u^{r}\right)^2 ({\pi_{Q}{}^{\eta}{}_{\eta}}+{\pi_{Q}{}^{\phi}{}_{\phi}})}{{u^{\tau}}} + {(\partial_{r} u^{r})} {\Pi_{Q}} {u^{\tau}} - {(\partial_{r} u^{r})} {u^{\tau}} ({\pi_{Q}{}^{\eta}{}_{\eta}}+{\pi_{Q}{}^{\phi}{}_{\phi}}) + {(\partial_{r} \Pi_{Q})} {u^{r}} {u^{\tau}}\\
  &\quad - {u^{r}} {u^{\tau}} ({(\partial_{r} \pi_{Q}{}^{\eta}{}_{\eta})}+{(\partial_{r} \pi_{Q}{}^{\phi}{}_{\phi})}) + \frac{{u^{r}} {u^{\tau}} ({\Pi_{Q}}-{\pi_{Q}{}^{\eta}{}_{\eta}}-{\pi_{Q}{}^{\phi}{}_{\phi}})}{{r}} \\
\nabla_{\mu}\delta M^{\mu r} &= - \frac{{\Pi_{Q}}+{\pi_{Q}{}^{\phi}{}_{\phi}}}{{r}} + 2 {(\partial_{r} u^{r})} {\Pi_{Q}} {u^{r}} - 2 {(\partial_{r} u^{r})} {u^{r}} ({\pi_{Q}{}^{\eta}{}_{\eta}}+{\pi_{Q}{}^{\phi}{}_{\phi}}) + \frac{{(\partial_{\tau} u^{r})} {\Pi_{Q}} \left(u^{r}\right)^2}{{u^{\tau}}}\\
  &\quad - \frac{{(\partial_{\tau} u^{r})} \left(u^{r}\right)^2 ({\pi_{Q}{}^{\eta}{}_{\eta}}+{\pi_{Q}{}^{\phi}{}_{\phi}})}{{u^{\tau}}} + {(\partial_{\tau} u^{r})} {\Pi_{Q}} {u^{\tau}} - {(\partial_{\tau} u^{r})} {u^{\tau}} ({\pi_{Q}{}^{\eta}{}_{\eta}}+{\pi_{Q}{}^{\phi}{}_{\phi}}) + {(\partial_{\tau} \Pi_{Q})} {u^{r}} {u^{\tau}}\\
  &\quad - {u^{r}} {u^{\tau}} ({(\partial_{\tau} \pi_{Q}{}^{\eta}{}_{\eta})}+{(\partial_{\tau} \pi_{Q}{}^{\phi}{}_{\phi})}) + \frac{{u^{r}} {u^{\tau}} ({\Pi_{Q}}-{\pi_{Q}{}^{\eta}{}_{\eta}}-{\pi_{Q}{}^{\phi}{}_{\phi}})}{{\tau}} + {(\partial_{r} \Pi_{Q})} \left(u^{\tau}\right)^2\\
  &\quad - \left(u^{\tau}\right)^2 ({(\partial_{r} \pi_{Q}{}^{\eta}{}_{\eta})}+{(\partial_{r} \pi_{Q}{}^{\phi}{}_{\phi})}) + \frac{\left(u^{\tau}\right)^2 ({\Pi_{Q}}-{\pi_{Q}{}^{\eta}{}_{\eta}}-{\pi_{Q}{}^{\phi}{}_{\phi}})}{{r}}
\end{align*}
\[
E_{\tau}^{\mathrm{S6}} = E_{\tau}^{\mathrm{S5}} + \nabla_{\mu}\delta M^{\mu\tau},\quad E_{r}^{\mathrm{S6}} = E_{r}^{\mathrm{S5}} + \nabla_{\mu}\delta M^{\mu r},\quad E_{\nu}^{\mathrm{S6}} = E_{\nu} + c_{M}\bigl(\Delta^{r}{}_{\tau}\nabla_{\mu}\delta M^{\mu\tau} + \Delta^{r}{}_{r}\nabla_{\mu}\delta M^{\mu r}\bigr),\quad c_{M} = \frac{D_{s}}{T}
\]
Every other row as S5.

\subsubsection{From the rows to the matrices}

\[
A^{(k)}_{ij} = \frac{\partial E_{i}}{\partial(\partial_{k}\phi_{j})},\qquad S_{i} = E_{i}\bigl|_{\partial\phi = 0},\qquad k = \tau, r
\]
\subsubsection{2+1D, S3}

\[
\Gamma^{{\tau}}{}_{{\eta}{\eta}} = {\tau},\quad \Gamma^{{\eta}}{}_{{\tau}{\eta}} = \frac{1}{{\tau}}
\]
\begin{align*}
N^{{\tau}} &= \frac{{n} \left(u^{\tau}\right)^2+{\nu^{x}} {u^{x}}+{\nu^{y}} {u^{y}}}{{u^{\tau}}} \\
N^{{x}} &= {n} {u^{x}}+{\nu^{x}} \\
N^{{y}} &= {n} {u^{y}}+{\nu^{y}} \\
E_{\alpha} = \nabla_{\mu}N^{\mu} &= {(\partial_{x} \nu^{x})} + {(\partial_{y} \nu^{y})} + {(\partial_{x} u^{x})} {n} + {(\partial_{y} u^{y})} {n} + {u^{\tau}} ({(\partial_{\alpha} n)} {(\partial_{\tau} \alpha)}+{(\partial_{T} n)} {(\partial_{\tau} T)})\\
  &\quad + {u^{x}} ({(\partial_{\alpha} n)} {(\partial_{x} \alpha)}+{(\partial_{T} n)} {(\partial_{x} T)}) + {u^{y}} ({(\partial_{\alpha} n)} {(\partial_{y} \alpha)}+{(\partial_{T} n)} {(\partial_{y} T)})\\
  &\quad + \frac{{(\partial_{\tau} \nu^{x})} {u^{x}}+{(\partial_{\tau} \nu^{y})} {u^{y}}+{(\partial_{\tau} u^{x})} {\nu^{x}}+{(\partial_{\tau} u^{y})} {\nu^{y}}}{{u^{\tau}}} + \frac{{n} ({(\partial_{\tau} u^{x})} {u^{x}}+{(\partial_{\tau} u^{y})} {u^{y}})}{{u^{\tau}}}\\
  &\quad - \frac{({(\partial_{\tau} u^{x})} {u^{x}}+{(\partial_{\tau} u^{y})} {u^{y}}) ({\nu^{x}} {u^{x}}+{\nu^{y}} {u^{y}})}{\left(u^{\tau}\right)^3} + \frac{{n} \left(u^{\tau}\right)^2+{\nu^{x}} {u^{x}}+{\nu^{y}} {u^{y}}}{{\tau} {u^{\tau}}}
\end{align*}
\begin{align*}
T_{Q}^{{\tau}{\tau}} &= {h} {n} \left(u^{\tau}\right)^2+2 {h} ({\nu^{x}} {u^{x}}+{\nu^{y}} {u^{y}})-{n} {T} \\
T_{Q}^{{\tau}{x}} &= \frac{{h} \left({n} \left(u^{\tau}\right)^2 {u^{x}}+{\nu^{x}} \left(2 \left(u^{x}\right)^2+\left(u^{y}\right)^2+1\right)+{\nu^{y}} {u^{x}} {u^{y}}\right)}{{u^{\tau}}} \\
T_{Q}^{{\tau}{y}} &= \frac{{h} \left({n} \left(u^{\tau}\right)^2 {u^{y}}+{\nu^{x}} {u^{x}} {u^{y}}+{\nu^{y}} \left(\left(u^{x}\right)^2+2 \left(u^{y}\right)^2+1\right)\right)}{{u^{\tau}}} \\
T_{Q}^{{x}{x}} &= {n} \left({h} \left(u^{x}\right)^2+{T}\right)+2 {h} {\nu^{x}} {u^{x}} \\
T_{Q}^{{x}{y}} &= {h} ({u^{y}} ({n} {u^{x}}+{\nu^{x}})+{\nu^{y}} {u^{x}}) \\
T_{Q}^{{y}{y}} &= {n} \left({h} \left(u^{y}\right)^2+{T}\right)+2 {h} {\nu^{y}} {u^{y}} \\
T_{Q}^{{\eta}{\eta}} &= \frac{{n} {T}}{{\tau}^2}
\end{align*}
\begin{align*}
\nabla_{\mu}T_{Q}^{\mu {\tau}} &= {(\partial_{{\tau}} T_{Q}^{{\tau}{\tau}})}+{(\partial_{{x}} T_{Q}^{{\tau}{x}})}+{(\partial_{{y}} T_{Q}^{{\tau}{y}})}+\frac{{T_{Q}^{{\tau}{\tau}}}}{{\tau}}+{\tau} {T_{Q}^{{\eta}{\eta}}} \\
\nabla_{\mu}T_{Q}^{\mu {x}} &= {(\partial_{{\tau}} T_{Q}^{{\tau}{x}})}+{(\partial_{{x}} T_{Q}^{{x}{x}})}+{(\partial_{{y}} T_{Q}^{{x}{y}})}+\frac{{T_{Q}^{{\tau}{x}}}}{{\tau}} \\
\nabla_{\mu}T_{Q}^{\mu {y}} &= {(\partial_{{\tau}} T_{Q}^{{\tau}{y}})}+{(\partial_{{x}} T_{Q}^{{x}{y}})}+{(\partial_{{y}} T_{Q}^{{y}{y}})}+\frac{{T_{Q}^{{\tau}{y}}}}{{\tau}}
\end{align*}
\begin{align*}
\nabla_{\mu}T_{Q}^{\mu {\tau}} &= {h} ({(\partial_{\alpha} n)} {(\partial_{\tau} \alpha)}+{(\partial_{T} n)} {(\partial_{\tau} T)}) - {(\partial_{\tau} T)} {n} + {h'(T)} {(\partial_{\tau} T)} {n}\\
  &\quad - {T} ({(\partial_{\alpha} n)} {(\partial_{\tau} \alpha)}+{(\partial_{T} n)} {(\partial_{\tau} T)}) + \frac{{n} {T}}{{\tau}} + {(\partial_{x} \nu^{x})} {h} {u^{\tau}} + {(\partial_{y} \nu^{y})} {h} {u^{\tau}} + {h'(T)} {(\partial_{x} T)} {\nu^{x}} {u^{\tau}}\\
  &\quad + {h'(T)} {(\partial_{y} T)} {\nu^{y}} {u^{\tau}} + {(\partial_{x} u^{x})} {n} {u^{\tau}} ({h}-{T}) + {(\partial_{y} u^{y})} {n} {u^{\tau}} ({h}-{T}) + {(\partial_{x} u^{x})} {n} {T} {u^{\tau}}\\
  &\quad + {(\partial_{y} u^{y})} {n} {T} {u^{\tau}} + \frac{{n} \left(u^{\tau}\right)^2 ({h}-{T})}{{\tau}} + {h} {u^{\tau}} {u^{x}} ({(\partial_{\alpha} n)} {(\partial_{x} \alpha)}+{(\partial_{T} n)} {(\partial_{x} T)})\\
  &\quad + {h'(T)} {(\partial_{x} T)} {n} {u^{\tau}} {u^{x}} + {h} \left(u^{x}\right)^2 ({(\partial_{\alpha} n)} {(\partial_{\tau} \alpha)}+{(\partial_{T} n)} {(\partial_{\tau} T)}) - {(\partial_{\tau} T)} {n} \left(u^{x}\right)^2\\
  &\quad + {h'(T)} {(\partial_{\tau} T)} {n} \left(u^{x}\right)^2 - {T} \left(u^{x}\right)^2 ({(\partial_{\alpha} n)} {(\partial_{\tau} \alpha)}+{(\partial_{T} n)} {(\partial_{\tau} T)})\\
  &\quad + {h} {u^{\tau}} {u^{y}} ({(\partial_{\alpha} n)} {(\partial_{y} \alpha)}+{(\partial_{T} n)} {(\partial_{y} T)}) + {h'(T)} {(\partial_{y} T)} {n} {u^{\tau}} {u^{y}}\\
  &\quad + {h} \left(u^{y}\right)^2 ({(\partial_{\alpha} n)} {(\partial_{\tau} \alpha)}+{(\partial_{T} n)} {(\partial_{\tau} T)}) - {(\partial_{\tau} T)} {n} \left(u^{y}\right)^2 + {h'(T)} {(\partial_{\tau} T)} {n} \left(u^{y}\right)^2\\
  &\quad - {T} \left(u^{y}\right)^2 ({(\partial_{\alpha} n)} {(\partial_{\tau} \alpha)}+{(\partial_{T} n)} {(\partial_{\tau} T)})\\
  &\quad + 2 {h} ({(\partial_{\tau} \nu^{x})} {u^{x}}+{(\partial_{\tau} \nu^{y})} {u^{y}}+{(\partial_{\tau} u^{x})} {\nu^{x}}+{(\partial_{\tau} u^{y})} {\nu^{y}})\\
  &\quad + 2 {n} ({h}-{T}) ({(\partial_{\tau} u^{x})} {u^{x}}+{(\partial_{\tau} u^{y})} {u^{y}}) + 2 {n} {T} ({(\partial_{\tau} u^{x})} {u^{x}}+{(\partial_{\tau} u^{y})} {u^{y}})\\
  &\quad + \frac{{h} {u^{x}} ({(\partial_{x} \nu^{x})} {u^{x}}+{(\partial_{x} \nu^{y})} {u^{y}}+{(\partial_{x} u^{x})} {\nu^{x}}+{(\partial_{x} u^{y})} {\nu^{y}})}{{u^{\tau}}} + \frac{{h} {\nu^{x}} ({(\partial_{x} u^{x})} {u^{x}}+{(\partial_{x} u^{y})} {u^{y}})}{{u^{\tau}}}\\
  &\quad + \frac{{n} {u^{x}} ({h}-{T}) ({(\partial_{x} u^{x})} {u^{x}}+{(\partial_{x} u^{y})} {u^{y}})}{{u^{\tau}}} + \frac{{n} {T} {u^{x}} ({(\partial_{x} u^{x})} {u^{x}}+{(\partial_{x} u^{y})} {u^{y}})}{{u^{\tau}}}\\
  &\quad + \frac{{h} {u^{y}} ({(\partial_{y} \nu^{x})} {u^{x}}+{(\partial_{y} \nu^{y})} {u^{y}}+{(\partial_{y} u^{x})} {\nu^{x}}+{(\partial_{y} u^{y})} {\nu^{y}})}{{u^{\tau}}} + \frac{{h} {\nu^{y}} ({(\partial_{y} u^{x})} {u^{x}}+{(\partial_{y} u^{y})} {u^{y}})}{{u^{\tau}}}\\
  &\quad + \frac{{n} {u^{y}} ({h}-{T}) ({(\partial_{y} u^{x})} {u^{x}}+{(\partial_{y} u^{y})} {u^{y}})}{{u^{\tau}}} + \frac{{n} {T} {u^{y}} ({(\partial_{y} u^{x})} {u^{x}}+{(\partial_{y} u^{y})} {u^{y}})}{{u^{\tau}}}\\
  &\quad + 2 {h'(T)} {(\partial_{\tau} T)} ({\nu^{x}} {u^{x}}+{\nu^{y}} {u^{y}}) + \frac{2 {h} ({\nu^{x}} {u^{x}}+{\nu^{y}} {u^{y}})}{{\tau}} + \frac{{(\partial_{x} u^{x})} {h} ({\nu^{x}} {u^{x}}+{\nu^{y}} {u^{y}})}{{u^{\tau}}}\\
  &\quad + \frac{{(\partial_{y} u^{y})} {h} ({\nu^{x}} {u^{x}}+{\nu^{y}} {u^{y}})}{{u^{\tau}}} + \frac{{h'(T)} {(\partial_{x} T)} {u^{x}} ({\nu^{x}} {u^{x}}+{\nu^{y}} {u^{y}})}{{u^{\tau}}}\\
  &\quad + \frac{{h'(T)} {(\partial_{y} T)} {u^{y}} ({\nu^{x}} {u^{x}}+{\nu^{y}} {u^{y}})}{{u^{\tau}}} - \frac{{h} {u^{x}} ({(\partial_{x} u^{x})} {u^{x}}+{(\partial_{x} u^{y})} {u^{y}}) ({\nu^{x}} {u^{x}}+{\nu^{y}} {u^{y}})}{\left(u^{\tau}\right)^3}\\
  &\quad - \frac{{h} {u^{y}} ({(\partial_{y} u^{x})} {u^{x}}+{(\partial_{y} u^{y})} {u^{y}}) ({\nu^{x}} {u^{x}}+{\nu^{y}} {u^{y}})}{\left(u^{\tau}\right)^3} + {(\partial_{\tau} T)} {n} \left(\left(u^{x}\right)^2+\left(u^{y}\right)^2\right)\\
  &\quad + {T} \left(\left(u^{x}\right)^2+\left(u^{y}\right)^2\right) ({(\partial_{\alpha} n)} {(\partial_{\tau} \alpha)}+{(\partial_{T} n)} {(\partial_{\tau} T)}) + \frac{{n} {T} \left(\left(u^{x}\right)^2+\left(u^{y}\right)^2\right)}{{\tau}} \\
\nabla_{\mu}T_{Q}^{\mu {x}} &= {(\partial_{y} u^{y})} {h} {\nu^{x}} + {(\partial_{y} u^{x})} {h} {\nu^{y}} + {(\partial_{\alpha} n)} {(\partial_{x} \alpha)} {T} + {(\partial_{\tau} \nu^{x})} {h} {u^{\tau}} + {(\partial_{\tau} u^{x})} {h} {n} {u^{\tau}}\\
  &\quad + {h'(T)} {(\partial_{\tau} T)} {\nu^{x}} {u^{\tau}} + \frac{{h} {\nu^{x}} {u^{\tau}}}{{\tau}} + 2 {(\partial_{x} \nu^{x})} {h} {u^{x}} + {(\partial_{y} \nu^{y})} {h} {u^{x}} + {(\partial_{y} u^{y})} {h} {n} {u^{x}}\\
  &\quad + {h'(T)} {(\partial_{y} T)} {\nu^{y}} {u^{x}} + \frac{3 {(\partial_{\tau} u^{x})} {h} {\nu^{x}} {u^{x}}}{{u^{\tau}}} + \frac{{(\partial_{\tau} u^{y})} {h} {\nu^{y}} {u^{x}}}{{u^{\tau}}} + {(\partial_{\alpha} n)} {(\partial_{\tau} \alpha)} {h} {u^{\tau}} {u^{x}}\\
  &\quad + {(\partial_{T} n)} {(\partial_{\tau} T)} {h} {u^{\tau}} {u^{x}} + {h'(T)} {(\partial_{\tau} T)} {n} {u^{\tau}} {u^{x}} + \frac{{h} {n} {u^{\tau}} {u^{x}}}{{\tau}} + {(\partial_{\alpha} n)} {(\partial_{x} \alpha)} {h} \left(u^{x}\right)^2\\
  &\quad + \frac{{(\partial_{\tau} \nu^{x})} {h} \left(u^{x}\right)^2}{{u^{\tau}}} + \frac{{(\partial_{\tau} u^{x})} {h} {n} \left(u^{x}\right)^2}{{u^{\tau}}} + \frac{{h'(T)} {(\partial_{\tau} T)} {\nu^{x}} \left(u^{x}\right)^2}{{u^{\tau}}} + \frac{{h} {\nu^{x}} \left(u^{x}\right)^2}{{\tau} {u^{\tau}}}\\
  &\quad - \frac{{(\partial_{\tau} u^{x})} {h} {\nu^{x}} \left(u^{x}\right)^3}{\left(u^{\tau}\right)^3} + 2 {(\partial_{x} u^{x})} {h} ({n} {u^{x}}+{\nu^{x}})\\
  &\quad + {(\partial_{x} T)} \left({h'(T)} {n} \left(u^{x}\right)^2+2 {h'(T)} {\nu^{x}} {u^{x}}+{(\partial_{T} n)} {h} \left(u^{x}\right)^2+{(\partial_{T} n)} {T}+{n}\right) + {(\partial_{y} \nu^{x})} {h} {u^{y}}\\
  &\quad + {(\partial_{y} u^{x})} {h} {n} {u^{y}} + {h'(T)} {(\partial_{y} T)} {\nu^{x}} {u^{y}} + \frac{{(\partial_{\tau} u^{y})} {h} {\nu^{x}} {u^{y}}}{{u^{\tau}}} + \frac{{(\partial_{\tau} u^{x})} {h} {\nu^{y}} {u^{y}}}{{u^{\tau}}} + {(\partial_{\alpha} n)} {(\partial_{y} \alpha)} {h} {u^{x}} {u^{y}}\\
  &\quad + {(\partial_{T} n)} {(\partial_{y} T)} {h} {u^{x}} {u^{y}} + {h'(T)} {(\partial_{y} T)} {n} {u^{x}} {u^{y}} + \frac{{(\partial_{\tau} \nu^{y})} {h} {u^{x}} {u^{y}}}{{u^{\tau}}} + \frac{{(\partial_{\tau} u^{y})} {h} {n} {u^{x}} {u^{y}}}{{u^{\tau}}}\\
  &\quad + \frac{{h'(T)} {(\partial_{\tau} T)} {\nu^{y}} {u^{x}} {u^{y}}}{{u^{\tau}}} + \frac{{h} {\nu^{y}} {u^{x}} {u^{y}}}{{\tau} {u^{\tau}}} - \frac{{(\partial_{\tau} u^{y})} {h} {\nu^{x}} \left(u^{x}\right)^2 {u^{y}}}{\left(u^{\tau}\right)^3} - \frac{{(\partial_{\tau} u^{x})} {h} {\nu^{y}} \left(u^{x}\right)^2 {u^{y}}}{\left(u^{\tau}\right)^3}\\
  &\quad - \frac{{(\partial_{\tau} u^{y})} {h} {\nu^{y}} {u^{x}} \left(u^{y}\right)^2}{\left(u^{\tau}\right)^3} \\
\nabla_{\mu}T_{Q}^{\mu {y}} &= {(\partial_{x} u^{x})} {h} {\nu^{y}} + 2 {(\partial_{y} u^{y})} {h} {\nu^{y}} + {(\partial_{\alpha} n)} {(\partial_{y} \alpha)} {T} + {(\partial_{\tau} \nu^{y})} {h} {u^{\tau}} + {(\partial_{\tau} u^{y})} {h} {n} {u^{\tau}}\\
  &\quad + {h'(T)} {(\partial_{\tau} T)} {\nu^{y}} {u^{\tau}} + \frac{{h} {\nu^{y}} {u^{\tau}}}{{\tau}} + {(\partial_{x} \nu^{y})} {h} {u^{x}} + {h'(T)} {(\partial_{x} T)} {\nu^{y}} {u^{x}} + \frac{{(\partial_{\tau} u^{y})} {h} {\nu^{x}} {u^{x}}}{{u^{\tau}}}\\
  &\quad + \frac{{(\partial_{\tau} u^{x})} {h} {\nu^{y}} {u^{x}}}{{u^{\tau}}} + {(\partial_{x} u^{y})} {h} ({n} {u^{x}}+{\nu^{x}}) + {(\partial_{x} \nu^{x})} {h} {u^{y}} + 2 {(\partial_{y} \nu^{y})} {h} {u^{y}} + {(\partial_{x} u^{x})} {h} {n} {u^{y}}\\
  &\quad + 2 {(\partial_{y} u^{y})} {h} {n} {u^{y}} + {h'(T)} {(\partial_{x} T)} {\nu^{x}} {u^{y}} + \frac{{(\partial_{\tau} u^{x})} {h} {\nu^{x}} {u^{y}}}{{u^{\tau}}} + \frac{3 {(\partial_{\tau} u^{y})} {h} {\nu^{y}} {u^{y}}}{{u^{\tau}}}\\
  &\quad + {(\partial_{\alpha} n)} {(\partial_{\tau} \alpha)} {h} {u^{\tau}} {u^{y}} + {(\partial_{T} n)} {(\partial_{\tau} T)} {h} {u^{\tau}} {u^{y}} + {h'(T)} {(\partial_{\tau} T)} {n} {u^{\tau}} {u^{y}} + \frac{{h} {n} {u^{\tau}} {u^{y}}}{{\tau}}\\
  &\quad + {(\partial_{\alpha} n)} {(\partial_{x} \alpha)} {h} {u^{x}} {u^{y}} + {(\partial_{T} n)} {(\partial_{x} T)} {h} {u^{x}} {u^{y}} + {h'(T)} {(\partial_{x} T)} {n} {u^{x}} {u^{y}} + \frac{{(\partial_{\tau} \nu^{x})} {h} {u^{x}} {u^{y}}}{{u^{\tau}}}\\
  &\quad + \frac{{(\partial_{\tau} u^{x})} {h} {n} {u^{x}} {u^{y}}}{{u^{\tau}}} + \frac{{h'(T)} {(\partial_{\tau} T)} {\nu^{x}} {u^{x}} {u^{y}}}{{u^{\tau}}} + \frac{{h} {\nu^{x}} {u^{x}} {u^{y}}}{{\tau} {u^{\tau}}} - \frac{{(\partial_{\tau} u^{x})} {h} {\nu^{x}} \left(u^{x}\right)^2 {u^{y}}}{\left(u^{\tau}\right)^3}\\
  &\quad + {(\partial_{\alpha} n)} {(\partial_{y} \alpha)} {h} \left(u^{y}\right)^2 + \frac{{(\partial_{\tau} \nu^{y})} {h} \left(u^{y}\right)^2}{{u^{\tau}}} + \frac{{(\partial_{\tau} u^{y})} {h} {n} \left(u^{y}\right)^2}{{u^{\tau}}} + \frac{{h'(T)} {(\partial_{\tau} T)} {\nu^{y}} \left(u^{y}\right)^2}{{u^{\tau}}}\\
  &\quad + \frac{{h} {\nu^{y}} \left(u^{y}\right)^2}{{\tau} {u^{\tau}}} - \frac{{(\partial_{\tau} u^{y})} {h} {\nu^{x}} {u^{x}} \left(u^{y}\right)^2}{\left(u^{\tau}\right)^3} - \frac{{(\partial_{\tau} u^{x})} {h} {\nu^{y}} {u^{x}} \left(u^{y}\right)^2}{\left(u^{\tau}\right)^3} - \frac{{(\partial_{\tau} u^{y})} {h} {\nu^{y}} \left(u^{y}\right)^3}{\left(u^{\tau}\right)^3}\\
  &\quad + {(\partial_{y} T)} \left({h'(T)} {n} \left(u^{y}\right)^2+2 {h'(T)} {\nu^{y}} {u^{y}}+{(\partial_{T} n)} {h} \left(u^{y}\right)^2+{(\partial_{T} n)} {T}+{n}\right)
\end{align*}
\begin{align*}
\Delta^{{x}}{}_{{\tau}} &= -{u^{\tau}} {u^{x}} \\
\Delta^{{x}}{}_{{x}} &= \left(u^{x}\right)^2+1 \\
\Delta^{{x}}{}_{{y}} &= {u^{x}} {u^{y}} \\
\Delta^{{y}}{}_{{\tau}} &= -{u^{\tau}} {u^{y}} \\
\Delta^{{y}}{}_{{x}} &= {u^{x}} {u^{y}} \\
\Delta^{{y}}{}_{{y}} &= \left(u^{y}\right)^2+1
\end{align*}
\[
E_{\nu^{i}} = \frac{D_{s}}{T}\Delta^{i}{}_{\lambda}\nabla_{\mu}T_{Q}^{\mu\lambda} + \nu^{i},\qquad i = x, y
\]
\begin{align*}
E_{\nu^{x}} &= {(\partial_{\alpha} n)} {D_{s}} {(\partial_{x} \alpha)} + {(\partial_{T} n)} {D_{s}} {(\partial_{x} T)} + {\nu^{x}} + \frac{{D_{s}} {(\partial_{x} T)} {n}}{{T}} + \frac{{(\partial_{\alpha} n)} {D_{s}} {(\partial_{\tau} \alpha)} {u^{x}}}{{u^{\tau}}}\\
  &\quad + \frac{{D_{s}} {(\partial_{\tau} T)} {u^{\tau}} {u^{x}} ({(\partial_{T} n)} {T}+{n})}{{T}} + {(\partial_{\alpha} n)} {D_{s}} {(\partial_{x} \alpha)} \left(u^{x}\right)^2 + {(\partial_{T} n)} {D_{s}} {(\partial_{x} T)} \left(u^{x}\right)^2\\
  &\quad + \frac{{D_{s}} {(\partial_{x} T)} {n} \left(u^{x}\right)^2}{{T}} + \frac{{(\partial_{\alpha} n)} {D_{s}} {(\partial_{\tau} \alpha)} \left(u^{x}\right)^3}{{u^{\tau}}} + {(\partial_{\alpha} n)} {D_{s}} {(\partial_{y} \alpha)} {u^{x}} {u^{y}} + {(\partial_{T} n)} {D_{s}} {(\partial_{y} T)} {u^{x}} {u^{y}}\\
  &\quad + \frac{{D_{s}} {(\partial_{y} T)} {n} {u^{x}} {u^{y}}}{{T}} + \frac{{(\partial_{\alpha} n)} {D_{s}} {(\partial_{\tau} \alpha)} {u^{x}} \left(u^{y}\right)^2}{{u^{\tau}}}\\
  &\quad + \frac{{h'(T)} {D_{s}} {\nu^{x}} ({(\partial_{\tau} T)} {u^{\tau}}+{(\partial_{x} T)} {u^{x}}+{(\partial_{y} T)} {u^{y}})}{{T}}\\
  &\quad + \frac{1}{{T} \left(u^{\tau}\right)^3}{D_{s}} {h}\Bigl({(\partial_{y} u^{x})} {\nu^{y}} {u^{\tau}} + {(\partial_{\tau} \nu^{x})} \left(u^{\tau}\right)^4 - {(\partial_{\tau} u^{y})} {\nu^{y}} {u^{x}} + {(\partial_{x} \nu^{x})} {u^{\tau}} {u^{x}}\\
  &\quad + {(\partial_{x} u^{x})} {n} {u^{\tau}} {u^{x}} - {(\partial_{x} u^{y})} {\nu^{y}} {u^{\tau}} \left(u^{x}\right)^2 + {(\partial_{y} u^{x})} {\nu^{y}} {u^{\tau}} \left(u^{x}\right)^2\\
  &\quad - 2 {(\partial_{\tau} u^{y})} {\nu^{y}} \left(u^{x}\right)^3 + {(\partial_{x} \nu^{x})} {u^{\tau}} \left(u^{x}\right)^3 + {(\partial_{x} u^{x})} {n} {u^{\tau}} \left(u^{x}\right)^3\\
  &\quad - {(\partial_{x} u^{y})} {\nu^{y}} {u^{\tau}} \left(u^{x}\right)^4 - {(\partial_{\tau} u^{y})} {\nu^{y}} \left(u^{x}\right)^5 + {(\partial_{y} \nu^{x})} {u^{\tau}} {u^{y}} + {(\partial_{y} u^{x})} {n} {u^{\tau}} {u^{y}}\\
  &\quad - {(\partial_{y} u^{y})} {\nu^{y}} {u^{\tau}} {u^{x}} {u^{y}} + {(\partial_{y} \nu^{x})} {u^{\tau}} \left(u^{x}\right)^2 {u^{y}} + {(\partial_{y} u^{x})} {n} {u^{\tau}} \left(u^{x}\right)^2 {u^{y}}\\
  &\quad + {(\partial_{x} u^{x})} {\nu^{y}} {u^{\tau}} \left(u^{x}\right)^3 {u^{y}} - {(\partial_{y} u^{y})} {\nu^{y}} {u^{\tau}} \left(u^{x}\right)^3 {u^{y}} + {(\partial_{y} u^{x})} {\nu^{y}} {u^{\tau}} \left(u^{y}\right)^2\\
  &\quad - {(\partial_{\tau} u^{y})} {\nu^{y}} {u^{x}} \left(u^{y}\right)^2 + {(\partial_{x} \nu^{x})} {u^{\tau}} {u^{x}} \left(u^{y}\right)^2 + {(\partial_{x} u^{x})} {n} {u^{\tau}} {u^{x}} \left(u^{y}\right)^2\\
  &\quad + {(\partial_{y} u^{x})} {\nu^{y}} {u^{\tau}} \left(u^{x}\right)^2 \left(u^{y}\right)^2 - {(\partial_{\tau} u^{y})} {\nu^{y}} \left(u^{x}\right)^3 \left(u^{y}\right)^2 + {(\partial_{y} \nu^{x})} {u^{\tau}} \left(u^{y}\right)^3\\
  &\quad + {(\partial_{y} u^{x})} {n} {u^{\tau}} \left(u^{y}\right)^3 + {(\partial_{\tau} u^{x})} \left(u^{\tau}\right)^2 \left({n} \left(u^{\tau}\right)^2+{\nu^{y}} \left(\left(u^{x}\right)^2+1\right) {u^{y}}\right)\Bigr)\\
  &\quad + \frac{1}{{T} {\tau} \left(u^{\tau}\right)^3}{D_{s}} {h} {\nu^{x}}\Bigl({(\partial_{x} u^{y})} {\tau} {u^{\tau}} \left(u^{x}\right)^3 {u^{y}} + \left(u^{x}\right)^4 ({(\partial_{\tau} u^{y})} {\tau} {u^{y}}+1)\\
  &\quad - {(\partial_{\tau} u^{x})} {\tau} \left(u^{\tau}\right)^2 {u^{x}} \left(\left(u^{y}\right)^2-1\right) - {(\partial_{y} u^{x})} {\tau} {u^{\tau}} {u^{x}} {u^{y}} \left(\left(u^{y}\right)^2+1\right)\\
  &\quad + \left(u^{y}\right)^2+1\Bigl(1 + 2 {(\partial_{x} u^{x})} {\tau} {u^{\tau}} + {(\partial_{y} u^{y})} {\tau} {u^{\tau}} + {(\partial_{\tau} u^{y})} {\tau} {u^{y}}\\
  &\quad + \left(u^{y}\right)^2\Bigr)\\
  &\quad + \left(u^{x}\right)^2\Bigl(2 + {(\partial_{x} u^{x})} {\tau} {u^{\tau}} + {(\partial_{y} u^{y})} {\tau} {u^{\tau}}\\
  &\quad + \left(u^{y}\right)^2 (-{(\partial_{x} u^{x})} {\tau} {u^{\tau}}+{(\partial_{y} u^{y})} {\tau} {u^{\tau}}+2)\\
  &\quad + {(\partial_{\tau} u^{y})} {\tau} {u^{y}} \left(\left(u^{y}\right)^2+2\right)\Bigr)\Bigr) \\
E_{\nu^{y}} &= {(\partial_{\alpha} n)} {D_{s}} {(\partial_{y} \alpha)} + {(\partial_{T} n)} {D_{s}} {(\partial_{y} T)} + {\nu^{y}} + \frac{{D_{s}} {(\partial_{y} T)} {n}}{{T}} + \frac{{(\partial_{\alpha} n)} {D_{s}} {(\partial_{\tau} \alpha)} {u^{y}}}{{u^{\tau}}}\\
  &\quad + \frac{{D_{s}} {(\partial_{\tau} T)} {u^{\tau}} {u^{y}} ({(\partial_{T} n)} {T}+{n})}{{T}} + {(\partial_{\alpha} n)} {D_{s}} {(\partial_{x} \alpha)} {u^{x}} {u^{y}} + {(\partial_{T} n)} {D_{s}} {(\partial_{x} T)} {u^{x}} {u^{y}}\\
  &\quad + \frac{{D_{s}} {(\partial_{x} T)} {n} {u^{x}} {u^{y}}}{{T}} + \frac{{(\partial_{\alpha} n)} {D_{s}} {(\partial_{\tau} \alpha)} \left(u^{x}\right)^2 {u^{y}}}{{u^{\tau}}} + {(\partial_{\alpha} n)} {D_{s}} {(\partial_{y} \alpha)} \left(u^{y}\right)^2\\
  &\quad + {(\partial_{T} n)} {D_{s}} {(\partial_{y} T)} \left(u^{y}\right)^2 + \frac{{D_{s}} {(\partial_{y} T)} {n} \left(u^{y}\right)^2}{{T}} + \frac{{(\partial_{\alpha} n)} {D_{s}} {(\partial_{\tau} \alpha)} \left(u^{y}\right)^3}{{u^{\tau}}}\\
  &\quad + \frac{{h'(T)} {D_{s}} {\nu^{y}} ({(\partial_{\tau} T)} {u^{\tau}}+{(\partial_{x} T)} {u^{x}}+{(\partial_{y} T)} {u^{y}})}{{T}}\\
  &\quad + \frac{1}{{T} {\tau} \left(u^{\tau}\right)^3}{D_{s}} {h} {\nu^{y}}\Bigl(1 + {(\partial_{x} u^{x})} {\tau} {u^{\tau}} + 2 {(\partial_{y} u^{y})} {\tau} {u^{\tau}} + {(\partial_{\tau} u^{y})} {\tau} {u^{y}} - {(\partial_{x} u^{y})} {\tau} {u^{\tau}} \left(u^{x}\right)^3 {u^{y}}\\
  &\quad + 2 \left(u^{y}\right)^2 + {(\partial_{x} u^{x})} {\tau} {u^{\tau}} \left(u^{y}\right)^2 + {(\partial_{y} u^{y})} {\tau} {u^{\tau}} \left(u^{y}\right)^2 + {(\partial_{\tau} u^{y})} {\tau} \left(u^{y}\right)^3\\
  &\quad + \left(u^{y}\right)^4 + \left(u^{x}\right)^4 (1-{(\partial_{\tau} u^{y})} {\tau} {u^{y}}) + {(\partial_{\tau} u^{x})} {\tau} \left(u^{\tau}\right)^2 {u^{x}} \left(\left(u^{y}\right)^2+1\right)\\
  &\quad + {\tau} {u^{\tau}} {u^{x}} {u^{y}} \left({(\partial_{y} u^{x})} \left(u^{y}\right)^2-{(\partial_{x} u^{y})}\right)\\
  &\quad + \left(u^{x}\right)^2\Bigl(2 + {(\partial_{x} u^{x})} {\tau} {u^{\tau}} + 2 {(\partial_{y} u^{y})} {\tau} {u^{\tau}}\\
  &\quad + \left(u^{y}\right)^2 ({(\partial_{x} u^{x})} {\tau} {u^{\tau}}-{(\partial_{y} u^{y})} {\tau} {u^{\tau}}+2)\\
  &\quad - {(\partial_{\tau} u^{y})} {\tau} \left(u^{y}\right)^3\Bigr)\Bigr)\\
  &\quad + \frac{1}{{T} \left(u^{\tau}\right)^3}{D_{s}} {h}\Bigl({(\partial_{x} u^{y})} {\nu^{x}} {u^{\tau}} + {(\partial_{\tau} \nu^{y})} \left(u^{\tau}\right)^4 + {(\partial_{x} \nu^{y})} {u^{\tau}} {u^{x}} + {(\partial_{x} u^{y})} {n} {u^{\tau}} {u^{x}}\\
  &\quad + {(\partial_{x} u^{y})} {\nu^{x}} {u^{\tau}} \left(u^{x}\right)^2 + {(\partial_{x} \nu^{y})} {u^{\tau}} \left(u^{x}\right)^3 + {(\partial_{x} u^{y})} {n} {u^{\tau}} \left(u^{x}\right)^3 - {(\partial_{\tau} u^{x})} {\nu^{x}} {u^{y}}\\
  &\quad + {(\partial_{y} \nu^{y})} {u^{\tau}} {u^{y}} + {(\partial_{y} u^{y})} {n} {u^{\tau}} {u^{y}} - {(\partial_{x} u^{x})} {\nu^{x}} {u^{\tau}} {u^{x}} {u^{y}} - {(\partial_{\tau} u^{x})} {\nu^{x}} \left(u^{x}\right)^2 {u^{y}}\\
  &\quad + {(\partial_{y} \nu^{y})} {u^{\tau}} \left(u^{x}\right)^2 {u^{y}} + {(\partial_{y} u^{y})} {n} {u^{\tau}} \left(u^{x}\right)^2 {u^{y}} + {(\partial_{x} u^{y})} {\nu^{x}} {u^{\tau}} \left(u^{y}\right)^2\\
  &\quad - {(\partial_{y} u^{x})} {\nu^{x}} {u^{\tau}} \left(u^{y}\right)^2 + {(\partial_{x} \nu^{y})} {u^{\tau}} {u^{x}} \left(u^{y}\right)^2 + {(\partial_{x} u^{y})} {n} {u^{\tau}} {u^{x}} \left(u^{y}\right)^2\\
  &\quad + {(\partial_{x} u^{y})} {\nu^{x}} {u^{\tau}} \left(u^{x}\right)^2 \left(u^{y}\right)^2 - 2 {(\partial_{\tau} u^{x})} {\nu^{x}} \left(u^{y}\right)^3 + {(\partial_{y} \nu^{y})} {u^{\tau}} \left(u^{y}\right)^3\\
  &\quad + {(\partial_{y} u^{y})} {n} {u^{\tau}} \left(u^{y}\right)^3 - {(\partial_{x} u^{x})} {\nu^{x}} {u^{\tau}} {u^{x}} \left(u^{y}\right)^3 + {(\partial_{y} u^{y})} {\nu^{x}} {u^{\tau}} {u^{x}} \left(u^{y}\right)^3\\
  &\quad - {(\partial_{\tau} u^{x})} {\nu^{x}} \left(u^{x}\right)^2 \left(u^{y}\right)^3 - {(\partial_{y} u^{x})} {\nu^{x}} {u^{\tau}} \left(u^{y}\right)^4 - {(\partial_{\tau} u^{x})} {\nu^{x}} \left(u^{y}\right)^5\\
  &\quad + {(\partial_{\tau} u^{y})} \left(u^{\tau}\right)^2 \left({n} \left(u^{\tau}\right)^2+{\nu^{x}} {u^{x}} \left(\left(u^{y}\right)^2+1\right)\right)\Bigr)
\end{align*}
\subsubsection{2+1D, S4}

\begin{align*}
S^{{\tau}{\tau}{\tau}} &= \frac{3}{{u^{\tau}} \left(\left(u^{y}\right)^2+1\right)}\Bigl({\mathcal{A}} {\nu^{x}} \left(u^{\tau}\right)^2 {u^{x}} - {\mathcal{C}} {\pi_{Q}^{yy}} \left(u^{x}\right)^2 + {I_{31}} \left(u^{\tau}\right)^2 \left(u^{x}\right)^2 + {\mathcal{C}} {\Pi_{Q}} \left(u^{\tau}\right)^2 \left(u^{x}\right)^2\\
  &\quad - {\mathcal{C}} {\pi_{Q}{}^{\eta}{}_{\eta}} \left(u^{\tau}\right)^2 \left(u^{x}\right)^2 + {\mathcal{B}} {\nu^{x}} \left(u^{x}\right)^3 - {\mathcal{C}} {\pi_{Q}^{yy}} \left(u^{x}\right)^4 + {\mathcal{A}} {\nu^{y}} \left(u^{\tau}\right)^2 {u^{y}}\\
  &\quad + 2 {\mathcal{C}} {\pi_{Q}^{xy}} {u^{x}} {u^{y}} + {\mathcal{B}} {\nu^{y}} \left(u^{x}\right)^2 {u^{y}} + 2 {\mathcal{C}} {\pi_{Q}^{xy}} \left(u^{x}\right)^3 {u^{y}} + {\mathcal{C}} {\pi_{Q}^{yy}} \left(u^{y}\right)^2\\
  &\quad + {I_{31}} \left(u^{\tau}\right)^2 \left(u^{y}\right)^2 + {\mathcal{C}} {\Pi_{Q}} \left(u^{\tau}\right)^2 \left(u^{y}\right)^2 + {\mathcal{B}} {\nu^{x}} {u^{x}} \left(u^{y}\right)^2 + {\mathcal{A}} {\nu^{x}} \left(u^{\tau}\right)^2 {u^{x}} \left(u^{y}\right)^2\\
  &\quad + {I_{31}} \left(u^{\tau}\right)^2 \left(u^{x}\right)^2 \left(u^{y}\right)^2 + {\mathcal{C}} {\Pi_{Q}} \left(u^{\tau}\right)^2 \left(u^{x}\right)^2 \left(u^{y}\right)^2 + {\mathcal{B}} {\nu^{x}} \left(u^{x}\right)^3 \left(u^{y}\right)^2\\
  &\quad + {\mathcal{B}} {\nu^{y}} \left(u^{y}\right)^3 + {\mathcal{A}} {\nu^{y}} \left(u^{\tau}\right)^2 \left(u^{y}\right)^3 + 2 {\mathcal{C}} {\pi_{Q}^{xy}} {u^{x}} \left(u^{y}\right)^3 + {\mathcal{B}} {\nu^{y}} \left(u^{x}\right)^2 \left(u^{y}\right)^3\\
  &\quad + {\mathcal{C}} {\pi_{Q}^{yy}} \left(u^{y}\right)^4 + {I_{31}} \left(u^{\tau}\right)^2 \left(u^{y}\right)^4 + {\mathcal{C}} {\Pi_{Q}} \left(u^{\tau}\right)^2 \left(u^{y}\right)^4 + {\mathcal{B}} {\nu^{x}} {u^{x}} \left(u^{y}\right)^4\\
  &\quad + {\mathcal{B}} {\nu^{y}} \left(u^{y}\right)^5\Bigr) \\
S^{{\tau}{\tau}{x}} &= \frac{1}{\left(u^{y}\right)^2+1}\Bigl(2 {\mathcal{C}} {\Pi_{Q}} {u^{x}} - 2 {\mathcal{C}} {\pi_{Q}^{yy}} {u^{x}} - 2 {\mathcal{C}} {\pi_{Q}{}^{\eta}{}_{\eta}} {u^{x}} + 3 {\mathcal{B}} {\nu^{x}} \left(u^{x}\right)^2 + 3 {\mathcal{C}} {\Pi_{Q}} \left(u^{x}\right)^3 - 3 {\mathcal{C}} {\pi_{Q}^{yy}} \left(u^{x}\right)^3\\
  &\quad - 3 {\mathcal{C}} {\pi_{Q}{}^{\eta}{}_{\eta}} \left(u^{x}\right)^3 + 2 {\mathcal{C}} {\pi_{Q}^{xy}} {u^{y}} + 2 {\mathcal{B}} {\nu^{y}} {u^{x}} {u^{y}} + 6 {\mathcal{C}} {\pi_{Q}^{xy}} \left(u^{x}\right)^2 {u^{y}} + {\mathcal{B}} {\nu^{x}} \left(u^{y}\right)^2\\
  &\quad + 5 {\mathcal{C}} {\Pi_{Q}} {u^{x}} \left(u^{y}\right)^2 + {\mathcal{C}} {\pi_{Q}^{yy}} {u^{x}} \left(u^{y}\right)^2 - 2 {\mathcal{C}} {\pi_{Q}{}^{\eta}{}_{\eta}} {u^{x}} \left(u^{y}\right)^2 + 3 {\mathcal{B}} {\nu^{x}} \left(u^{x}\right)^2 \left(u^{y}\right)^2\\
  &\quad + 3 {\mathcal{C}} {\Pi_{Q}} \left(u^{x}\right)^3 \left(u^{y}\right)^2 + 2 {\mathcal{C}} {\pi_{Q}^{xy}} \left(u^{y}\right)^3 + 2 {\mathcal{B}} {\nu^{y}} {u^{x}} \left(u^{y}\right)^3 + {\mathcal{B}} {\nu^{x}} \left(u^{y}\right)^4\\
  &\quad + 3 {\mathcal{C}} {\Pi_{Q}} {u^{x}} \left(u^{y}\right)^4 + {I_{31}} {u^{x}} \left(\left(u^{y}\right)^2+1\right) \left(3 \left(u^{x}\right)^2+3 \left(u^{y}\right)^2+2\right)\\
  &\quad + {\mathcal{A}} \left(\left(u^{y}\right)^2+1\right) \left({\nu^{x}} \left(3 \left(u^{x}\right)^2+\left(u^{y}\right)^2+1\right)+2 {\nu^{y}} {u^{x}} {u^{y}}\right)\Bigr) \\
S^{{\tau}{\tau}{y}} &= 2 {\mathcal{B}} {\nu^{x}} {u^{x}} {u^{y}} + {\mathcal{B}} {\nu^{y}} \left(\left(u^{x}\right)^2+3 \left(u^{y}\right)^2\right) + {I_{31}} {u^{y}} \left(3 \left(u^{x}\right)^2+3 \left(u^{y}\right)^2+2\right)\\
  &\quad + {\mathcal{A}} \left(2 {\nu^{x}} {u^{x}} {u^{y}}+{\nu^{y}} \left(\left(u^{x}\right)^2+3 \left(u^{y}\right)^2+1\right)\right)\\
  &\quad + \frac{1}{\left(u^{y}\right)^2+1}{\mathcal{C}}\Bigl(2 {\pi_{Q}^{xy}} \left(2 {u^{x}} \left(u^{y}\right)^2+{u^{x}}\right)\\
  &\quad + {u^{y}}\Bigl(- {\pi_{Q}{}^{\eta}{}_{\eta}} \left(u^{x}\right)^2 + {\pi_{Q}^{yy}} \left(-\left(u^{x}\right)^2+3 \left(u^{y}\right)^2+2\right)\\
  &\quad + {\Pi_{Q}} \left(\left(u^{y}\right)^2+1\right) \left(3 \left(u^{x}\right)^2+3 \left(u^{y}\right)^2+2\right)\Bigr)\Bigr) \\
S^{{\tau}{x}{x}} &= \frac{1}{{u^{\tau}} \left(\left(u^{y}\right)^2+1\right)}\Bigl({I_{31}} \left(u^{\tau}\right)^2 \left(3 \left(u^{x}\right)^2+1\right) \left(\left(u^{y}\right)^2+1\right)\\
  &\quad + {\mathcal{C}}\Bigl(- {\pi_{Q}{}^{\eta}{}_{\eta}} - 4 {\pi_{Q}{}^{\eta}{}_{\eta}} \left(u^{x}\right)^2 - 3 {\pi_{Q}{}^{\eta}{}_{\eta}} \left(u^{x}\right)^4 + 4 {\pi_{Q}^{xy}} {u^{x}} {u^{y}} + 6 {\pi_{Q}^{xy}} \left(u^{x}\right)^3 {u^{y}}\\
  &\quad - 2 {\pi_{Q}{}^{\eta}{}_{\eta}} \left(u^{y}\right)^2 - 4 {\pi_{Q}{}^{\eta}{}_{\eta}} \left(u^{x}\right)^2 \left(u^{y}\right)^2 + 4 {\pi_{Q}^{xy}} {u^{x}} \left(u^{y}\right)^3 - {\pi_{Q}{}^{\eta}{}_{\eta}} \left(u^{y}\right)^4\\
  &\quad + {\Pi_{Q}} \left(u^{\tau}\right)^2 \left(3 \left(u^{x}\right)^2+1\right) \left(\left(u^{y}\right)^2+1\right)\\
  &\quad - {\pi_{Q}^{yy}} \left(\left(u^{x}\right)^2+1\right) \left(3 \left(u^{x}\right)^2+\left(u^{y}\right)^2+1\right)\Bigr)\\
  &\quad + \left(u^{y}\right)^2+1\Bigl({\mathcal{A}} {u^{x}} \left({\nu^{x}} \left(3 \left(u^{x}\right)^2+2 \left(u^{y}\right)^2+2\right)+{\nu^{y}} {u^{x}} {u^{y}}\right)\\
  &\quad + {\mathcal{B}}\Bigl({\nu^{x}} {u^{x}} \left(3 \left(u^{x}\right)^2+2 \left(u^{y}\right)^2+3\right)+{\nu^{y}} \left(\left(u^{x}\right)^2+1\right) {u^{y}}\Bigr)\Bigr)\Bigr) \\
S^{{\tau}{x}{y}} &= \frac{1}{{u^{\tau}} \left(\left(u^{y}\right)^2+1\right)}\Bigl(\left(u^{y}\right)^2+1\Bigl(3 {I_{31}} \left(u^{\tau}\right)^2 {u^{x}} {u^{y}} + {\mathcal{A}} {\nu^{x}} {u^{y}} \left(2 \left(u^{x}\right)^2+\left(u^{y}\right)^2+1\right)\\
  &\quad + {\mathcal{B}} {\nu^{x}} {u^{y}} \left(2 \left(u^{x}\right)^2+\left(u^{y}\right)^2+1\right)\\
  &\quad + {\mathcal{A}} {\nu^{y}} \left(\left(u^{x}\right)^3+2 {u^{x}} \left(u^{y}\right)^2+{u^{x}}\right)\\
  &\quad + {\mathcal{B}} {\nu^{y}} \left(\left(u^{x}\right)^3+2 {u^{x}} \left(u^{y}\right)^2+{u^{x}}\right)\Bigr)\\
  &\quad + {\mathcal{C}}\Bigl({\pi_{Q}^{xy}} \left(2 \left(u^{y}\right)^2+1\right) \left(2 \left(u^{x}\right)^2+\left(u^{y}\right)^2+1\right)\\
  &\quad + {u^{x}} {u^{y}}\Bigl(- {\pi_{Q}{}^{\eta}{}_{\eta}} \left(u^{\tau}\right)^2 + 3 {\Pi_{Q}} \left(u^{\tau}\right)^2 \left(\left(u^{y}\right)^2+1\right)\\
  &\quad + {\pi_{Q}^{yy}} \left(\left(u^{y}\right)^2-\left(u^{x}\right)^2\right)\Bigr)\Bigr)\Bigr) \\
S^{{\tau}{y}{y}} &= \frac{1}{{u^{\tau}}}\Bigl({\mathcal{B}} {\nu^{x}} {u^{x}} + 2 {\mathcal{A}} {\nu^{y}} {u^{y}} + 3 {\mathcal{B}} {\nu^{y}} {u^{y}} + 2 {\mathcal{A}} {\nu^{y}} \left(u^{x}\right)^2 {u^{y}} + 2 {\mathcal{B}} {\nu^{y}} \left(u^{x}\right)^2 {u^{y}} + {\mathcal{A}} {\nu^{x}} {u^{x}} \left(u^{y}\right)^2\\
  &\quad + {\mathcal{B}} {\nu^{x}} {u^{x}} \left(u^{y}\right)^2 + 3 {\mathcal{A}} {\nu^{y}} \left(u^{y}\right)^3 + 3 {\mathcal{B}} {\nu^{y}} \left(u^{y}\right)^3 + {I_{31}} \left(u^{\tau}\right)^2 \left(3 \left(u^{y}\right)^2+1\right)\\
  &\quad + {\mathcal{C}} \left({\Pi_{Q}} \left(u^{\tau}\right)^2 \left(3 \left(u^{y}\right)^2+1\right)+2 {\pi_{Q}^{xy}} {u^{x}} {u^{y}}+{\pi_{Q}^{yy}} \left(u^{x}\right)^2+3 {\pi_{Q}^{yy}} \left(u^{y}\right)^2+{\pi_{Q}^{yy}}\right)\Bigr) \\
S^{{\tau}{\eta}{\eta}} &= \frac{{\mathcal{B}} ({\nu^{x}} {u^{x}}+{\nu^{y}} {u^{y}})+{\mathcal{C}} \left(u^{\tau}\right)^2 ({\Pi_{Q}}+{\pi_{Q}{}^{\eta}{}_{\eta}})+{I_{31}} \left(u^{\tau}\right)^2}{{\tau}^2 {u^{\tau}}} \\
S^{{x}{x}{x}} &= 3\Bigl({\mathcal{A}} {\nu^{x}} \left(u^{x}\right)^2 + {\mathcal{B}} {\nu^{x}} \left(\left(u^{x}\right)^2+1\right) + {I_{31}} {u^{x}} \left(\left(u^{x}\right)^2+1\right)\\
  &\quad + {\mathcal{C}} {u^{x}}\Bigl({\Pi_{Q}} + {\Pi_{Q}} \left(u^{x}\right)^2\\
  &\quad - \frac{-2 {\pi_{Q}^{xy}} {u^{x}} {u^{y}}+{\pi_{Q}^{yy}} \left(u^{x}\right)^2+{\pi_{Q}^{yy}}+{\pi_{Q}{}^{\eta}{}_{\eta}} \left(u^{x}\right)^2+{\pi_{Q}{}^{\eta}{}_{\eta}} \left(u^{y}\right)^2+{\pi_{Q}{}^{\eta}{}_{\eta}}}{\left(u^{y}\right)^2+1}\Bigr)\Bigr) \\
S^{{x}{x}{y}} &= {\mathcal{B}} {\nu^{y}} \left(\left(u^{x}\right)^2+1\right) + 2 {\mathcal{B}} {\nu^{x}} {u^{x}} {u^{y}} + 2 {I_{31}} \left(u^{x}\right)^2 {u^{y}} + {I_{31}} \left(\left(u^{x}\right)^2+1\right) {u^{y}}\\
  &\quad + {\mathcal{A}} {u^{x}} (2 {\nu^{x}} {u^{y}}+{\nu^{y}} {u^{x}}) + 2 {\mathcal{C}} {u^{x}} ({\Pi_{Q}} {u^{x}} {u^{y}}+{\pi_{Q}^{xy}})\\
  &\quad + {\mathcal{C}} {u^{y}}\Bigl({\Pi_{Q}} + {\Pi_{Q}} \left(u^{x}\right)^2\\
  &\quad - \frac{-2 {\pi_{Q}^{xy}} {u^{x}} {u^{y}}+{\pi_{Q}^{yy}} \left(u^{x}\right)^2+{\pi_{Q}^{yy}}+{\pi_{Q}{}^{\eta}{}_{\eta}} \left(u^{x}\right)^2+{\pi_{Q}{}^{\eta}{}_{\eta}} \left(u^{y}\right)^2+{\pi_{Q}{}^{\eta}{}_{\eta}}}{\left(u^{y}\right)^2+1}\Bigr) \\
S^{{x}{y}{y}} &= {\mathcal{C}} {\Pi_{Q}} {u^{x}} + {\mathcal{C}} {\pi_{Q}^{yy}} {u^{x}} + 2 {\mathcal{C}} {\pi_{Q}^{xy}} {u^{y}} + 2 {\mathcal{A}} {\nu^{y}} {u^{x}} {u^{y}} + {\mathcal{A}} {\nu^{x}} \left(u^{y}\right)^2 + 3 {\mathcal{C}} {\Pi_{Q}} {u^{x}} \left(u^{y}\right)^2\\
  &\quad + {\mathcal{B}} \left({\nu^{x}} \left(u^{y}\right)^2+{\nu^{x}}+2 {\nu^{y}} {u^{x}} {u^{y}}\right) + {I_{31}} \left(3 {u^{x}} \left(u^{y}\right)^2+{u^{x}}\right) \\
S^{{x}{\eta}{\eta}} &= \frac{{\mathcal{B}} {\nu^{x}}+{u^{x}} ({\mathcal{C}} ({\Pi_{Q}}+{\pi_{Q}{}^{\eta}{}_{\eta}})+{I_{31}})}{{\tau}^2} \\
S^{{y}{y}{y}} &= 3 \left({u^{y}} \left({\mathcal{A}} {\nu^{y}} {u^{y}}+{\mathcal{C}} \left({\Pi_{Q}} \left(u^{y}\right)^2+{\Pi_{Q}}+{\pi_{Q}^{yy}}\right)+{I_{31}} \left(\left(u^{y}\right)^2+1\right)\right)+{\mathcal{B}} {\nu^{y}} \left(\left(u^{y}\right)^2+1\right)\right) \\
S^{{y}{\eta}{\eta}} &= \frac{{\mathcal{B}} {\nu^{y}}+{u^{y}} ({\mathcal{C}} ({\Pi_{Q}}+{\pi_{Q}{}^{\eta}{}_{\eta}})+{I_{31}})}{{\tau}^2}
\end{align*}
\begin{align*}
\nabla_{\lambda}S^{\lambda {x}{x}} &= {(\partial_{{\tau}} S^{{\tau}{x}{x}})}+{(\partial_{{x}} S^{{x}{x}{x}})}+{(\partial_{{y}} S^{{x}{x}{y}})}+\frac{{S^{{\tau}{x}{x}}}}{{\tau}} \\
\nabla_{\lambda}S^{\lambda {x}{y}} &= {(\partial_{{\tau}} S^{{\tau}{x}{y}})}+{(\partial_{{x}} S^{{x}{x}{y}})}+{(\partial_{{y}} S^{{x}{y}{y}})}+\frac{{S^{{\tau}{x}{y}}}}{{\tau}} \\
\nabla_{\lambda}S^{\lambda {y}{y}} &= {(\partial_{{\tau}} S^{{\tau}{y}{y}})}+{(\partial_{{x}} S^{{x}{y}{y}})}+{(\partial_{{y}} S^{{y}{y}{y}})}+\frac{{S^{{\tau}{y}{y}}}}{{\tau}} \\
\nabla_{\lambda}S^{\lambda {\eta}{\eta}} &= {(\partial_{{\tau}} S^{{\tau}{\eta}{\eta}})}+{(\partial_{{x}} S^{{x}{\eta}{\eta}})}+{(\partial_{{y}} S^{{y}{\eta}{\eta}})}+\frac{5 {S^{{\tau}{\eta}{\eta}}}}{{\tau}}
\end{align*}
\[
\mathcal{E}^{ab} = \Delta^{a}{}_{\alpha}\Delta^{b}{}_{\beta}\nabla_{\lambda}S^{\lambda\alpha\beta},\qquad \mathcal{E}^{\lambda}{}_{\lambda} = \Delta_{ab}\mathcal{E}^{ab}
\]
\begin{gather*}
E_{ab} = \frac{D_{s}}{2T}\Bigl(\mathcal{E}^{ab} - \tfrac{1}{3}\Delta^{ab}\mathcal{E}^{\lambda}{}_{\lambda}\Bigr) + \pi_{Q}^{ab},\quad ab = yy,\ xy,\qquad E_{\eta} = \tau^{2}\frac{D_{s}}{2T}\Bigl(\mathcal{E}^{\eta\eta} - \tfrac{1}{3}\Delta^{\eta\eta}\mathcal{E}^{\lambda}{}_{\lambda}\Bigr) + \pi_{Q}{}^{\eta}{}_{\eta} \\
E_{\Pi_{Q}} = \frac{D_{s}}{6T}\mathcal{E}^{\lambda}{}_{\lambda} + \Pi_{Q}
\end{gather*}
\subsubsection{2+1D, S5}

\begin{align*}
T_{m}^{{\tau}{\tau}} &= \frac{1}{\left(u^{y}\right)^2+1}\Bigl({\Pi} \left(u^{x}\right)^2 - {\pi^{\eta}{}_{\eta}} \left(u^{x}\right)^2 - {\pi^{yy}} \left(u^{x}\right)^2 + 2 {\pi^{xy}} {u^{x}} {u^{y}} + {\Pi} \left(u^{y}\right)^2 + {\pi^{yy}} \left(u^{y}\right)^2\\
  &\quad + {\Pi} \left(u^{x}\right)^2 \left(u^{y}\right)^2 + {\Pi} \left(u^{y}\right)^4 - {P} \left(\left(u^{y}\right)^2+1\right) + {P'(T)} {T} \left(u^{\tau}\right)^2 \left(\left(u^{y}\right)^2+1\right)\Bigr) \\
T_{m}^{{\tau}{x}} &= \frac{1}{{u^{\tau}} \left(\left(u^{y}\right)^2+1\right)}\Bigl(- {\pi^{yy}} {u^{x}} + {\Pi} \left(u^{\tau}\right)^2 {u^{x}} - {\pi^{\eta}{}_{\eta}} \left(u^{\tau}\right)^2 {u^{x}} + {P'(T)} {T} \left(u^{\tau}\right)^2 {u^{x}} - {\pi^{yy}} \left(u^{x}\right)^3\\
  &\quad + {\pi^{xy}} {u^{y}} + 2 {\pi^{xy}} \left(u^{x}\right)^2 {u^{y}} + {\Pi} \left(u^{\tau}\right)^2 {u^{x}} \left(u^{y}\right)^2 + {P'(T)} {T} \left(u^{\tau}\right)^2 {u^{x}} \left(u^{y}\right)^2\\
  &\quad + {\pi^{xy}} \left(u^{y}\right)^3\Bigr) \\
T_{m}^{{\tau}{y}} &= \frac{\left(u^{\tau}\right)^2 {u^{y}} ({P'(T)} {T}+{\Pi})+{\pi^{xy}} {u^{x}}+{\pi^{yy}} {u^{y}}}{{u^{\tau}}} \\
T_{m}^{{x}{x}} &= \left(u^{x}\right)^2 ({P'(T)} {T}+{\Pi})+{\Pi}-\frac{{\pi^{\eta}{}_{\eta}} \left(u^{\tau}\right)^2-2 {\pi^{xy}} {u^{x}} {u^{y}}+{\pi^{yy}} \left(u^{x}\right)^2+{\pi^{yy}}}{\left(u^{y}\right)^2+1}+{P} \\
T_{m}^{{x}{y}} &= {u^{x}} {u^{y}} ({P'(T)} {T}+{\Pi})+{\pi^{xy}} \\
T_{m}^{{y}{y}} &= \left(u^{y}\right)^2 ({P'(T)} {T}+{\Pi})+{\Pi}+{\pi^{yy}}+{P} \\
T_{m}^{{\eta}{\eta}} &= \frac{{\Pi}+{\pi^{\eta}{}_{\eta}}+{P}}{{\tau}^2}
\end{align*}
\begin{gather*}
E_{\tau} = \nabla_{\mu}(T_{m} + T_{Q})^{\mu\tau},\quad E_{x} = \nabla_{\mu}(T_{m} + T_{Q})^{\mu x},\quad E_{y} = \nabla_{\mu}(T_{m} + T_{Q})^{\mu y} \\
A^{(k)}_{ij} = \frac{\partial E_{i}}{\partial(\partial_{k}\phi_{j})},\qquad S_{i} = E_{i}\bigl|_{\partial\phi = 0},\qquad k = \tau, x, y
\end{gather*}
\begin{align*}
E_{\tau} &= - {(\partial_{\tau} \Pi)} - {P'(T)} {(\partial_{\tau} T)} + \frac{{\Pi}+{\pi^{\eta}{}_{\eta}}+{P}}{{\tau}} + \frac{{n} {T}}{{\tau}} + {(\partial_{x} u^{x})} {n} {u^{\tau}} ({h}-{T})\\
  &\quad + {(\partial_{y} u^{y})} {n} {u^{\tau}} ({h}-{T}) + {(\partial_{x} u^{x})} {n} {T} {u^{\tau}} + {(\partial_{y} u^{y})} {n} {T} {u^{\tau}} + {(\partial_{x} u^{x})} {u^{\tau}} ({P'(T)} {T}+{\Pi})\\
  &\quad + {(\partial_{y} u^{y})} {u^{\tau}} ({P'(T)} {T}+{\Pi})\\
  &\quad + \left(u^{\tau}\right)^2 (({h'(T)}-1) {(\partial_{\tau} T)} {n}+{(\partial_{\alpha} n)} {(\partial_{\tau} \alpha)} ({h}-{T})+{(\partial_{T} n)} {(\partial_{\tau} T)} ({h}-{T}))\\
  &\quad + \left(u^{\tau}\right)^2 ({(\partial_{\tau} T)} ({P''(T)} {T}+{P'(T)})+{(\partial_{\tau} \Pi)}) + {(\partial_{x} T)} {n} {u^{\tau}} {u^{x}}\\
  &\quad + {u^{\tau}} {u^{x}} (({h'(T)}-1) {(\partial_{x} T)} {n}+{(\partial_{\alpha} n)} {(\partial_{x} \alpha)} ({h}-{T})+{(\partial_{T} n)} {(\partial_{x} T)} ({h}-{T}))\\
  &\quad + {T} {u^{\tau}} {u^{x}} ({(\partial_{\alpha} n)} {(\partial_{x} \alpha)}+{(\partial_{T} n)} {(\partial_{x} T)}) + {u^{\tau}} {u^{x}} ({(\partial_{x} T)} ({P''(T)} {T}+{P'(T)})+{(\partial_{x} \Pi)})\\
  &\quad + {(\partial_{y} T)} {n} {u^{\tau}} {u^{y}}\\
  &\quad + {u^{\tau}} {u^{y}} (({h'(T)}-1) {(\partial_{y} T)} {n}+{(\partial_{\alpha} n)} {(\partial_{y} \alpha)} ({h}-{T})+{(\partial_{T} n)} {(\partial_{y} T)} ({h}-{T}))\\
  &\quad + {T} {u^{\tau}} {u^{y}} ({(\partial_{\alpha} n)} {(\partial_{y} \alpha)}+{(\partial_{T} n)} {(\partial_{y} T)}) + {u^{\tau}} {u^{y}} ({(\partial_{y} T)} ({P''(T)} {T}+{P'(T)})+{(\partial_{y} \Pi)})\\
  &\quad + 2 {h} ({(\partial_{\tau} \nu^{x})} {u^{x}}+{(\partial_{\tau} \nu^{y})} {u^{y}}+{(\partial_{\tau} u^{x})} {\nu^{x}}+{(\partial_{\tau} u^{y})} {\nu^{y}})\\
  &\quad + 2 {n} ({h}-{T}) ({(\partial_{\tau} u^{x})} {u^{x}}+{(\partial_{\tau} u^{y})} {u^{y}}) + 2 {n} {T} ({(\partial_{\tau} u^{x})} {u^{x}}+{(\partial_{\tau} u^{y})} {u^{y}})\\
  &\quad + 2 ({P'(T)} {T}+{\Pi}) ({(\partial_{\tau} u^{x})} {u^{x}}+{(\partial_{\tau} u^{y})} {u^{y}}) + \frac{{n} {u^{x}} ({h}-{T}) ({(\partial_{x} u^{x})} {u^{x}}+{(\partial_{x} u^{y})} {u^{y}})}{{u^{\tau}}}\\
  &\quad + \frac{{n} {T} {u^{x}} ({(\partial_{x} u^{x})} {u^{x}}+{(\partial_{x} u^{y})} {u^{y}})}{{u^{\tau}}} + \frac{{u^{x}} ({P'(T)} {T}+{\Pi}) ({(\partial_{x} u^{x})} {u^{x}}+{(\partial_{x} u^{y})} {u^{y}})}{{u^{\tau}}}\\
  &\quad + \frac{{(\partial_{y} \pi^{xy})} {u^{x}}+{(\partial_{y} \pi^{yy})} {u^{y}}+{(\partial_{y} u^{x})} {\pi^{xy}}+{(\partial_{y} u^{y})} {\pi^{yy}}}{{u^{\tau}}}\\
  &\quad + \frac{{n} {u^{y}} ({h}-{T}) ({(\partial_{y} u^{x})} {u^{x}}+{(\partial_{y} u^{y})} {u^{y}})}{{u^{\tau}}} + \frac{{n} {T} {u^{y}} ({(\partial_{y} u^{x})} {u^{x}}+{(\partial_{y} u^{y})} {u^{y}})}{{u^{\tau}}}\\
  &\quad + \frac{{u^{y}} ({P'(T)} {T}+{\Pi}) ({(\partial_{y} u^{x})} {u^{x}}+{(\partial_{y} u^{y})} {u^{y}})}{{u^{\tau}}} + 2 {h'(T)} {(\partial_{\tau} T)} ({\nu^{x}} {u^{x}}+{\nu^{y}} {u^{y}})\\
  &\quad - \frac{({(\partial_{y} u^{x})} {u^{x}}+{(\partial_{y} u^{y})} {u^{y}}) ({\pi^{xy}} {u^{x}}+{\pi^{yy}} {u^{y}})}{\left(u^{\tau}\right)^3} + {(\partial_{\tau} T)} {n} \left(\left(u^{x}\right)^2+\left(u^{y}\right)^2\right)\\
  &\quad + {T} \left(\left(u^{x}\right)^2+\left(u^{y}\right)^2\right) ({(\partial_{\alpha} n)} {(\partial_{\tau} \alpha)}+{(\partial_{T} n)} {(\partial_{\tau} T)})\\
  &\quad + \frac{{h} {n} \left(u^{\tau}\right)^2+2 {h} ({\nu^{x}} {u^{x}}+{\nu^{y}} {u^{y}})-{n} {T}}{{\tau}}\\
  &\quad + \frac{({(\partial_{\tau} u^{x})} {u^{x}}+{(\partial_{\tau} u^{y})} {u^{y}}) \left({\pi^{\eta}{}_{\eta}} \left(u^{x}\right)^2-2 {\pi^{xy}} {u^{x}} {u^{y}}+{\pi^{yy}} \left(\left(u^{x}\right)^2-\left(u^{y}\right)^2\right)\right)}{\left(u^{\tau}\right)^2 \left(\left(u^{y}\right)^2+1\right)}\\
  &\quad + \frac{1}{{\tau} \left(\left(u^{y}\right)^2+1\right)}\Bigl({\Pi} \left(u^{x}\right)^2 - {\pi^{\eta}{}_{\eta}} \left(u^{x}\right)^2 - {\pi^{yy}} \left(u^{x}\right)^2 + 2 {\pi^{xy}} {u^{x}} {u^{y}} + {\Pi} \left(u^{y}\right)^2 + {\pi^{yy}} \left(u^{y}\right)^2\\
  &\quad + {\Pi} \left(u^{x}\right)^2 \left(u^{y}\right)^2 + {\Pi} \left(u^{y}\right)^4 - {P} \left(\left(u^{y}\right)^2+1\right)\\
  &\quad + {P'(T)} {T} \left(u^{\tau}\right)^2 \left(\left(u^{y}\right)^2+1\right)\Bigr) + \frac{{h'(T)} {(\partial_{x} T)} \left({\nu^{x}} \left(2 \left(u^{x}\right)^2+\left(u^{y}\right)^2+1\right)+{\nu^{y}} {u^{x}} {u^{y}}\right)}{{u^{\tau}}}\\
  &\quad + \frac{1}{\left(u^{\tau}\right)^3 \left(\left(u^{y}\right)^2+1\right)}{(\partial_{x} u^{x})} {u^{x}}+{(\partial_{x} u^{y})} {u^{y}}\Bigl({\pi^{\eta}{}_{\eta}} \left(u^{\tau}\right)^2 {u^{x}} + {\pi^{yy}} \left(\left(u^{x}\right)^3+{u^{x}}\right)\\
  &\quad - {\pi^{xy}} {u^{y}} \left(2 \left(u^{x}\right)^2+\left(u^{y}\right)^2+1\right)\Bigr) + \frac{{h'(T)} {(\partial_{y} T)} \left({\nu^{x}} {u^{x}} {u^{y}}+{\nu^{y}} \left(\left(u^{x}\right)^2+2 \left(u^{y}\right)^2+1\right)\right)}{{u^{\tau}}}\\
  &\quad + \frac{1}{\left(u^{\tau}\right)^3}{h}\Bigl({(\partial_{x} u^{y})} {\nu^{y}} {u^{x}} + {(\partial_{x} u^{y})} {\nu^{y}} \left(u^{x}\right)^3 + {(\partial_{x} u^{y})} {\nu^{x}} {u^{y}} + {(\partial_{x} \nu^{y})} {u^{x}} {u^{y}}\\
  &\quad + {(\partial_{x} \nu^{y})} \left(u^{x}\right)^3 {u^{y}} + {(\partial_{x} u^{y})} {\nu^{x}} \left(u^{y}\right)^3 + {(\partial_{x} \nu^{y})} {u^{x}} \left(u^{y}\right)^3\\
  &\quad + {(\partial_{x} u^{x})} {\nu^{y}} {u^{y}} \left(\left(u^{y}\right)^2+1\right) + {(\partial_{x} u^{x})} {\nu^{x}} {u^{x}} \left(2 \left(u^{x}\right)^2+3 \left(u^{y}\right)^2+3\right)\\
  &\quad + {(\partial_{x} \nu^{x})} \left(2 \left(u^{x}\right)^4+3 \left(u^{x}\right)^2 \left(\left(u^{y}\right)^2+1\right)+\left(\left(u^{y}\right)^2+1\right)^2\right)\Bigr)\\
  &\quad + \frac{1}{\left(u^{\tau}\right)^3}{h}\Bigl({(\partial_{y} u^{x})} {\nu^{y}} {u^{x}} + {(\partial_{y} u^{x})} {\nu^{y}} \left(u^{x}\right)^3 + {(\partial_{y} u^{y})} {\nu^{x}} \left(\left(u^{x}\right)^3+{u^{x}}\right) + {(\partial_{y} u^{x})} {\nu^{x}} {u^{y}}\\
  &\quad + {(\partial_{y} \nu^{x})} {u^{x}} {u^{y}} + {(\partial_{y} \nu^{x})} \left(u^{x}\right)^3 {u^{y}} + {(\partial_{y} u^{x})} {\nu^{x}} \left(u^{y}\right)^3 + {(\partial_{y} \nu^{x})} {u^{x}} \left(u^{y}\right)^3\\
  &\quad + {(\partial_{y} u^{y})} {\nu^{y}} {u^{y}} \left(3 \left(u^{x}\right)^2+2 \left(u^{y}\right)^2+3\right)\\
  &\quad + {(\partial_{y} \nu^{y})} \left(\left(u^{x}\right)^4+\left(u^{x}\right)^2 \left(3 \left(u^{y}\right)^2+2\right)+2 \left(u^{y}\right)^4+3 \left(u^{y}\right)^2+1\right)\Bigr)\\
  &\quad + \frac{1}{{u^{\tau}} \left(\left(u^{y}\right)^2+1\right)^2}\Bigl(- \left(u^{y}\right)^2+1\Bigl({(\partial_{x} \pi^{\eta}{}_{\eta})} \left(u^{\tau}\right)^2 {u^{x}} + {(\partial_{x} \pi^{yy})} \left(\left(u^{x}\right)^3+{u^{x}}\right)\\
  &\quad - {(\partial_{x} \pi^{xy})} {u^{y}} \left(2 \left(u^{x}\right)^2+\left(u^{y}\right)^2+1\right)\Bigr)\\
  &\quad - {(\partial_{x} u^{x})} \left(\left(u^{y}\right)^2+1\right)\Bigl({\pi^{yy}} + 3 {\pi^{yy}} \left(u^{x}\right)^2 - 4 {\pi^{xy}} {u^{x}} {u^{y}}\\
  &\quad + {\pi^{\eta}{}_{\eta}} \left(3 \left(u^{x}\right)^2+\left(u^{y}\right)^2+1\right)\Bigr)\\
  &\quad + {(\partial_{x} u^{y})}\Bigl(2 {u^{x}} {u^{y}} \left({\pi^{\eta}{}_{\eta}} \left(u^{x}\right)^2+{\pi^{yy}} \left(u^{x}\right)^2+{\pi^{yy}}\right)\\
  &\quad + {\pi^{xy}} \left(\left(\left(u^{y}\right)^2+1\right)^2-2 \left(u^{x}\right)^2 \left(\left(u^{y}\right)^2-1\right)\right)\Bigr)\Bigr)\\
  &\quad + \frac{1}{\left(u^{\tau}\right)^2 \left(\left(u^{y}\right)^2+1\right)^2}\Bigl(2 {(\partial_{\tau} u^{y})} {\pi^{xy}} {u^{x}} \left(-\left(\left(u^{x}\right)^2 \left(\left(u^{y}\right)^2-1\right)\right)+\left(u^{y}\right)^2+1\right)\\
  &\quad + \left(u^{\tau}\right)^2 \left(\left(u^{y}\right)^2+1\right)\Bigl(- {(\partial_{\tau} \pi^{\eta}{}_{\eta})} \left(u^{x}\right)^2 + 2 {(\partial_{\tau} \pi^{xy})} {u^{x}} {u^{y}}\\
  &\quad + {(\partial_{\tau} \pi^{yy})} \left(\left(u^{y}\right)^2-\left(u^{x}\right)^2\right)\Bigr)\\
  &\quad - {(\partial_{\tau} u^{x})} \left(\left(u^{y}\right)^2+1\right)\Bigl(- 2 {\pi^{xy}} {u^{y}} \left(2 \left(u^{x}\right)^2+\left(u^{y}\right)^2+1\right)\\
  &\quad + {\pi^{yy}} {u^{x}} \left(3 \left(u^{x}\right)^2+\left(u^{y}\right)^2+2\right)\\
  &\quad + {\pi^{\eta}{}_{\eta}} {u^{x}} \left(3 \left(u^{x}\right)^2+2 \left(u^{y}\right)^2+2\right)\Bigr)\\
  &\quad + {(\partial_{\tau} u^{y})} {u^{y}}\Bigl({\pi^{\eta}{}_{\eta}} \left(u^{x}\right)^2 \left(2 \left(u^{x}\right)^2+\left(u^{y}\right)^2+1\right)\\
  &\quad + {\pi^{yy}}\Bigl(2 + 2 \left(u^{x}\right)^4 + 3 \left(u^{y}\right)^2 + \left(u^{y}\right)^4\\
  &\quad + \left(u^{x}\right)^2 \left(\left(u^{y}\right)^2+3\right)\Bigr)\Bigr)\Bigr)
\end{align*}
\subsubsection{2+1D, S6}

\begin{align*}
\delta M^{{\tau}{\tau}} &= \frac{{\Pi_{Q}} \left(\left(u^{y}\right)^2+1\right) \left(\left(u^{x}\right)^2+\left(u^{y}\right)^2\right)+{u^{x}} (2 {\pi_{Q}^{xy}} {u^{y}}-{\pi_{Q}{}^{\eta}{}_{\eta}} {u^{x}})+{\pi_{Q}^{yy}} \left(\left(u^{y}\right)^2-\left(u^{x}\right)^2\right)}{\left(u^{y}\right)^2+1} \\
\delta M^{{\tau}{x}} &= \frac{1}{{u^{\tau}} \left(\left(u^{y}\right)^2+1\right)}\Bigl(- {\pi_{Q}^{yy}} {u^{x}} + {\Pi_{Q}} \left(u^{\tau}\right)^2 {u^{x}} - {\pi_{Q}{}^{\eta}{}_{\eta}} \left(u^{\tau}\right)^2 {u^{x}} - {\pi_{Q}^{yy}} \left(u^{x}\right)^3 + {\pi_{Q}^{xy}} {u^{y}}\\
  &\quad + 2 {\pi_{Q}^{xy}} \left(u^{x}\right)^2 {u^{y}} + {\Pi_{Q}} \left(u^{\tau}\right)^2 {u^{x}} \left(u^{y}\right)^2 + {\pi_{Q}^{xy}} \left(u^{y}\right)^3\Bigr) \\
\delta M^{{\tau}{y}} &= \frac{{u^{y}} \left({\Pi_{Q}} \left(u^{\tau}\right)^2+{\pi_{Q}^{yy}}\right)+{\pi_{Q}^{xy}} {u^{x}}}{{u^{\tau}}} \\
\delta M^{{x}{x}} &= {\Pi_{Q}} \left(u^{x}\right)^2+{\Pi_{Q}}-\frac{-2 {\pi_{Q}^{xy}} {u^{x}} {u^{y}}+{\pi_{Q}^{yy}} \left(u^{x}\right)^2+{\pi_{Q}^{yy}}+{\pi_{Q}{}^{\eta}{}_{\eta}} \left(u^{x}\right)^2+{\pi_{Q}{}^{\eta}{}_{\eta}} \left(u^{y}\right)^2+{\pi_{Q}{}^{\eta}{}_{\eta}}}{\left(u^{y}\right)^2+1} \\
\delta M^{{x}{y}} &= {\Pi_{Q}} {u^{x}} {u^{y}}+{\pi_{Q}^{xy}} \\
\delta M^{{y}{y}} &= {\Pi_{Q}} \left(u^{y}\right)^2+{\Pi_{Q}}+{\pi_{Q}^{yy}} \\
\delta M^{{\eta}{\eta}} &= \frac{{\Pi_{Q}}+{\pi_{Q}{}^{\eta}{}_{\eta}}}{{\tau}^2}
\end{align*}
\begin{align*}
\nabla_{\mu}\delta M^{\mu {\tau}} &= {(\partial_{{\tau}} \delta M^{{\tau}{\tau}})}+{(\partial_{{x}} \delta M^{{\tau}{x}})}+{(\partial_{{y}} \delta M^{{\tau}{y}})}+\frac{{\delta M^{{\tau}{\tau}}}}{{\tau}}+{\tau} {\delta M^{{\eta}{\eta}}} \\
\nabla_{\mu}\delta M^{\mu {x}} &= {(\partial_{{\tau}} \delta M^{{\tau}{x}})}+{(\partial_{{x}} \delta M^{{x}{x}})}+{(\partial_{{y}} \delta M^{{x}{y}})}+\frac{{\delta M^{{\tau}{x}}}}{{\tau}} \\
\nabla_{\mu}\delta M^{\mu {y}} &= {(\partial_{{\tau}} \delta M^{{\tau}{y}})}+{(\partial_{{x}} \delta M^{{x}{y}})}+{(\partial_{{y}} \delta M^{{y}{y}})}+\frac{{\delta M^{{\tau}{y}}}}{{\tau}}
\end{align*}
\begin{align*}
\nabla_{\mu}\delta M^{\mu {\tau}} &= \frac{{\Pi_{Q}}+{\pi_{Q}{}^{\eta}{}_{\eta}}}{{\tau}} + {(\partial_{x} u^{x})} {\Pi_{Q}} {u^{\tau}} + {(\partial_{y} u^{y})} {\Pi_{Q}} {u^{\tau}} + {(\partial_{x} \Pi_{Q})} {u^{\tau}} {u^{x}} + {(\partial_{y} \Pi_{Q})} {u^{\tau}} {u^{y}}\\
  &\quad + 2 {\Pi_{Q}} ({(\partial_{\tau} u^{x})} {u^{x}}+{(\partial_{\tau} u^{y})} {u^{y}}) + \frac{{\Pi_{Q}} {u^{x}} ({(\partial_{x} u^{x})} {u^{x}}+{(\partial_{x} u^{y})} {u^{y}})}{{u^{\tau}}}\\
  &\quad + \frac{{(\partial_{y} \pi_{Q}^{xy})} {u^{x}}+{(\partial_{y} \pi_{Q}^{yy})} {u^{y}}+{(\partial_{y} u^{x})} {\pi_{Q}^{xy}}+{(\partial_{y} u^{y})} {\pi_{Q}^{yy}}}{{u^{\tau}}} + \frac{{\Pi_{Q}} {u^{y}} ({(\partial_{y} u^{x})} {u^{x}}+{(\partial_{y} u^{y})} {u^{y}})}{{u^{\tau}}}\\
  &\quad - \frac{({(\partial_{y} u^{x})} {u^{x}}+{(\partial_{y} u^{y})} {u^{y}}) ({\pi_{Q}^{xy}} {u^{x}}+{\pi_{Q}^{yy}} {u^{y}})}{\left(u^{\tau}\right)^3} + {(\partial_{\tau} \Pi_{Q})} \left(\left(u^{x}\right)^2+\left(u^{y}\right)^2\right)\\
  &\quad + \frac{({(\partial_{\tau} u^{x})} {u^{x}}+{(\partial_{\tau} u^{y})} {u^{y}}) \left({u^{x}} ({\pi_{Q}{}^{\eta}{}_{\eta}} {u^{x}}-2 {\pi_{Q}^{xy}} {u^{y}})+{\pi_{Q}^{yy}} \left(\left(u^{x}\right)^2-\left(u^{y}\right)^2\right)\right)}{\left(u^{\tau}\right)^2 \left(\left(u^{y}\right)^2+1\right)}\\
  &\quad + \frac{{\Pi_{Q}} \left(\left(u^{y}\right)^2+1\right) \left(\left(u^{x}\right)^2+\left(u^{y}\right)^2\right)+{u^{x}} (2 {\pi_{Q}^{xy}} {u^{y}}-{\pi_{Q}{}^{\eta}{}_{\eta}} {u^{x}})+{\pi_{Q}^{yy}} \left(\left(u^{y}\right)^2-\left(u^{x}\right)^2\right)}{{\tau} \left(\left(u^{y}\right)^2+1\right)}\\
  &\quad + \frac{1}{\left(u^{\tau}\right)^3 \left(\left(u^{y}\right)^2+1\right)}{(\partial_{x} u^{x})} {u^{x}}+{(\partial_{x} u^{y})} {u^{y}}\Bigl({\pi_{Q}{}^{\eta}{}_{\eta}} \left(u^{\tau}\right)^2 {u^{x}} + {\pi_{Q}^{yy}} \left(\left(u^{x}\right)^3+{u^{x}}\right)\\
  &\quad - {\pi_{Q}^{xy}} {u^{y}} \left(2 \left(u^{x}\right)^2+\left(u^{y}\right)^2+1\right)\Bigr)\\
  &\quad + \frac{1}{{u^{\tau}} \left(\left(u^{y}\right)^2+1\right)^2}\Bigl(- \left(u^{y}\right)^2+1\Bigl({(\partial_{x} \pi_{Q}{}^{\eta}{}_{\eta})} {u^{x}} + {(\partial_{x} \pi_{Q}{}^{\eta}{}_{\eta})} \left(u^{x}\right)^3\\
  &\quad + {(\partial_{x} \pi_{Q}^{yy})} \left(\left(u^{x}\right)^3+{u^{x}}\right) - {(\partial_{x} \pi_{Q}^{xy})} {u^{y}}\\
  &\quad - 2 {(\partial_{x} \pi_{Q}^{xy})} \left(u^{x}\right)^2 {u^{y}} + {(\partial_{x} \pi_{Q}{}^{\eta}{}_{\eta})} {u^{x}} \left(u^{y}\right)^2\\
  &\quad - {(\partial_{x} \pi_{Q}^{xy})} \left(u^{y}\right)^3\\
  &\quad + {(\partial_{x} u^{x})}\Bigl({\pi_{Q}^{yy}} + {\pi_{Q}{}^{\eta}{}_{\eta}} + 3 {\pi_{Q}^{yy}} \left(u^{x}\right)^2 + 3 {\pi_{Q}{}^{\eta}{}_{\eta}} \left(u^{x}\right)^2\\
  &\quad - 4 {\pi_{Q}^{xy}} {u^{x}} {u^{y}} + {\pi_{Q}{}^{\eta}{}_{\eta}} \left(u^{y}\right)^2\Bigr)\Bigr)\\
  &\quad + {(\partial_{x} u^{y})}\Bigl(2 {u^{x}} {u^{y}} \left({\pi_{Q}^{yy}} \left(u^{x}\right)^2+{\pi_{Q}^{yy}}+{\pi_{Q}{}^{\eta}{}_{\eta}} \left(u^{x}\right)^2\right)\\
  &\quad + {\pi_{Q}^{xy}} \left(\left(\left(u^{y}\right)^2+1\right)^2-2 \left(u^{x}\right)^2 \left(\left(u^{y}\right)^2-1\right)\right)\Bigr)\Bigr)\\
  &\quad + \frac{1}{\left(u^{\tau}\right)^2 \left(\left(u^{y}\right)^2+1\right)^2}\Bigl(- \left(u^{y}\right)^2+1\Bigl(\left(u^{\tau}\right)^2\Bigl({u^{x}} ({(\partial_{\tau} \pi_{Q}{}^{\eta}{}_{\eta})} {u^{x}}-2 {(\partial_{\tau} \pi_{Q}^{xy})} {u^{y}})\\
  &\quad + {(\partial_{\tau} \pi_{Q}^{yy})} \left(\left(u^{x}\right)^2-\left(u^{y}\right)^2\right)\Bigr)\\
  &\quad + {(\partial_{\tau} u^{x})}\Bigl(- 2 {\pi_{Q}^{xy}} {u^{y}} \left(2 \left(u^{x}\right)^2+\left(u^{y}\right)^2+1\right)\\
  &\quad + {\pi_{Q}^{yy}} {u^{x}} \left(3 \left(u^{x}\right)^2+\left(u^{y}\right)^2+2\right)\\
  &\quad + {\pi_{Q}{}^{\eta}{}_{\eta}} {u^{x}} \left(3 \left(u^{x}\right)^2+2 \left(u^{y}\right)^2+2\right)\Bigr)\Bigr)\\
  &\quad + {(\partial_{\tau} u^{y})}\Bigl(2 {\pi_{Q}^{xy}} {u^{x}} \left(-\left(\left(u^{x}\right)^2 \left(\left(u^{y}\right)^2-1\right)\right)+\left(u^{y}\right)^2+1\right)\\
  &\quad + {u^{y}}\Bigl({\pi_{Q}{}^{\eta}{}_{\eta}} \left(u^{x}\right)^2 \left(2 \left(u^{x}\right)^2+\left(u^{y}\right)^2+1\right)\\
  &\quad + {\pi_{Q}^{yy}}\Bigl(2 + 2 \left(u^{x}\right)^4 + 3 \left(u^{y}\right)^2 + \left(u^{y}\right)^4\\
  &\quad + \left(u^{x}\right)^2 \left(\left(u^{y}\right)^2+3\right)\Bigr)\Bigr)\Bigr)\Bigr) \\
\nabla_{\mu}\delta M^{\mu {x}} &= \frac{1}{{\tau} \left(u^{\tau}\right)^3 \left(\left(u^{y}\right)^2+1\right)^2}\Bigl({(\partial_{\tau} u^{y})} {\pi_{Q}^{xy}} {\tau} \left(u^{\tau}\right)^2 - {(\partial_{\tau} u^{x})} {\pi_{Q}^{yy}} {\tau} \left(u^{\tau}\right)^2 - {(\partial_{\tau} u^{x})} {\pi_{Q}{}^{\eta}{}_{\eta}} {\tau} \left(u^{\tau}\right)^2\\
  &\quad + {(\partial_{x} \Pi_{Q})} {\tau} \left(u^{\tau}\right)^3 - {(\partial_{x} \pi_{Q}^{yy})} {\tau} \left(u^{\tau}\right)^3 - {(\partial_{x} \pi_{Q}{}^{\eta}{}_{\eta})} {\tau} \left(u^{\tau}\right)^3\\
  &\quad + {(\partial_{y} \pi_{Q}^{xy})} {\tau} \left(u^{\tau}\right)^3 + {(\partial_{\tau} u^{x})} {\Pi_{Q}} {\tau} \left(u^{\tau}\right)^4 - {\pi_{Q}^{yy}} \left(u^{\tau}\right)^2 {u^{x}}\\
  &\quad - {(\partial_{\tau} \pi_{Q}^{yy})} {\tau} \left(u^{\tau}\right)^2 {u^{x}} - {(\partial_{\tau} \pi_{Q}{}^{\eta}{}_{\eta})} {\tau} \left(u^{\tau}\right)^2 {u^{x}} + 2 {(\partial_{x} u^{x})} {\Pi_{Q}} {\tau} \left(u^{\tau}\right)^3 {u^{x}}\\
  &\quad + {(\partial_{y} u^{y})} {\Pi_{Q}} {\tau} \left(u^{\tau}\right)^3 {u^{x}} + 2 {(\partial_{x} u^{y})} {\pi_{Q}^{xy}} {\tau} \left(u^{\tau}\right)^3 {u^{x}}\\
  &\quad - 2 {(\partial_{x} u^{x})} {\pi_{Q}^{yy}} {\tau} \left(u^{\tau}\right)^3 {u^{x}} - 2 {(\partial_{x} u^{x})} {\pi_{Q}{}^{\eta}{}_{\eta}} {\tau} \left(u^{\tau}\right)^3 {u^{x}} + {\Pi_{Q}} \left(u^{\tau}\right)^4 {u^{x}}\\
  &\quad - {\pi_{Q}{}^{\eta}{}_{\eta}} \left(u^{\tau}\right)^4 {u^{x}} + {(\partial_{\tau} \Pi_{Q})} {\tau} \left(u^{\tau}\right)^4 {u^{x}} + {(\partial_{\tau} u^{x})} {\pi_{Q}^{yy}} {\tau} \left(u^{x}\right)^2\\
  &\quad + {(\partial_{\tau} u^{x})} {\Pi_{Q}} {\tau} \left(u^{\tau}\right)^2 \left(u^{x}\right)^2 + 2 {(\partial_{\tau} u^{y})} {\pi_{Q}^{xy}} {\tau} \left(u^{\tau}\right)^2 \left(u^{x}\right)^2\\
  &\quad - 3 {(\partial_{\tau} u^{x})} {\pi_{Q}^{yy}} {\tau} \left(u^{\tau}\right)^2 \left(u^{x}\right)^2 - 2 {(\partial_{\tau} u^{x})} {\pi_{Q}{}^{\eta}{}_{\eta}} {\tau} \left(u^{\tau}\right)^2 \left(u^{x}\right)^2\\
  &\quad + {(\partial_{x} \Pi_{Q})} {\tau} \left(u^{\tau}\right)^3 \left(u^{x}\right)^2 - {(\partial_{x} \pi_{Q}^{yy})} {\tau} \left(u^{\tau}\right)^3 \left(u^{x}\right)^2\\
  &\quad - {(\partial_{x} \pi_{Q}{}^{\eta}{}_{\eta})} {\tau} \left(u^{\tau}\right)^3 \left(u^{x}\right)^2 - {\pi_{Q}^{yy}} \left(u^{\tau}\right)^2 \left(u^{x}\right)^3 - {(\partial_{\tau} \pi_{Q}^{yy})} {\tau} \left(u^{\tau}\right)^2 \left(u^{x}\right)^3\\
  &\quad - {(\partial_{\tau} \pi_{Q}{}^{\eta}{}_{\eta})} {\tau} \left(u^{\tau}\right)^2 \left(u^{x}\right)^3 + {(\partial_{\tau} u^{x})} {\pi_{Q}^{yy}} {\tau} \left(u^{x}\right)^4 + {\pi_{Q}^{xy}} \left(u^{\tau}\right)^2 {u^{y}}\\
  &\quad + {(\partial_{\tau} \pi_{Q}^{xy})} {\tau} \left(u^{\tau}\right)^2 {u^{y}} + {(\partial_{y} u^{x})} {\Pi_{Q}} {\tau} \left(u^{\tau}\right)^3 {u^{y}} + 2 {(\partial_{x} u^{x})} {\pi_{Q}^{xy}} {\tau} \left(u^{\tau}\right)^3 {u^{y}}\\
  &\quad + 2 {(\partial_{x} u^{y})} {\pi_{Q}^{yy}} {\tau} \left(u^{\tau}\right)^3 {u^{y}} - 2 {(\partial_{x} u^{y})} {\pi_{Q}{}^{\eta}{}_{\eta}} {\tau} \left(u^{\tau}\right)^3 {u^{y}}\\
  &\quad + 2 {(\partial_{x} u^{y})} {\pi_{Q}{}^{\eta}{}_{\eta}} {\tau} \left(u^{\tau}\right)^5 {u^{y}} - {(\partial_{\tau} u^{x})} {\pi_{Q}^{xy}} {\tau} {u^{x}} {u^{y}} + {(\partial_{\tau} u^{y})} {\pi_{Q}^{yy}} {\tau} {u^{x}} {u^{y}}\\
  &\quad + {(\partial_{\tau} u^{y})} {\Pi_{Q}} {\tau} \left(u^{\tau}\right)^2 {u^{x}} {u^{y}} + 4 {(\partial_{\tau} u^{x})} {\pi_{Q}^{xy}} {\tau} \left(u^{\tau}\right)^2 {u^{x}} {u^{y}}\\
  &\quad + 2 {(\partial_{\tau} u^{y})} {\pi_{Q}^{yy}} {\tau} \left(u^{\tau}\right)^2 {u^{x}} {u^{y}} + {(\partial_{\tau} u^{y})} {\pi_{Q}{}^{\eta}{}_{\eta}} {\tau} \left(u^{\tau}\right)^2 {u^{x}} {u^{y}}\\
  &\quad + 2 {(\partial_{x} \pi_{Q}^{xy})} {\tau} \left(u^{\tau}\right)^3 {u^{x}} {u^{y}} + {(\partial_{y} \Pi_{Q})} {\tau} \left(u^{\tau}\right)^3 {u^{x}} {u^{y}}\\
  &\quad + 2 {\pi_{Q}^{xy}} \left(u^{\tau}\right)^2 \left(u^{x}\right)^2 {u^{y}} + 2 {(\partial_{\tau} \pi_{Q}^{xy})} {\tau} \left(u^{\tau}\right)^2 \left(u^{x}\right)^2 {u^{y}}\\
  &\quad + 2 {(\partial_{x} u^{y})} {\pi_{Q}^{yy}} {\tau} \left(u^{\tau}\right)^3 \left(u^{x}\right)^2 {u^{y}} - 2 {(\partial_{\tau} u^{x})} {\pi_{Q}^{xy}} {\tau} \left(u^{x}\right)^3 {u^{y}}\\
  &\quad + {(\partial_{\tau} u^{y})} {\pi_{Q}^{yy}} {\tau} \left(u^{x}\right)^3 {u^{y}} + 2 {(\partial_{\tau} u^{y})} {\pi_{Q}^{yy}} {\tau} \left(u^{\tau}\right)^2 \left(u^{x}\right)^3 {u^{y}}\\
  &\quad + 2 {(\partial_{\tau} u^{y})} {\pi_{Q}{}^{\eta}{}_{\eta}} {\tau} \left(u^{\tau}\right)^2 \left(u^{x}\right)^3 {u^{y}} - {(\partial_{\tau} u^{y})} {\pi_{Q}^{xy}} {\tau} \left(u^{y}\right)^2\\
  &\quad + 2 {(\partial_{\tau} u^{y})} {\pi_{Q}^{xy}} {\tau} \left(u^{\tau}\right)^2 \left(u^{y}\right)^2 - {(\partial_{\tau} u^{x})} {\pi_{Q}^{yy}} {\tau} \left(u^{\tau}\right)^2 \left(u^{y}\right)^2\\
  &\quad - 2 {(\partial_{\tau} u^{x})} {\pi_{Q}{}^{\eta}{}_{\eta}} {\tau} \left(u^{\tau}\right)^2 \left(u^{y}\right)^2 + 2 {(\partial_{x} \Pi_{Q})} {\tau} \left(u^{\tau}\right)^3 \left(u^{y}\right)^2\\
  &\quad - {(\partial_{x} \pi_{Q}^{yy})} {\tau} \left(u^{\tau}\right)^3 \left(u^{y}\right)^2 - 2 {(\partial_{x} \pi_{Q}{}^{\eta}{}_{\eta})} {\tau} \left(u^{\tau}\right)^3 \left(u^{y}\right)^2\\
  &\quad + 2 {(\partial_{y} \pi_{Q}^{xy})} {\tau} \left(u^{\tau}\right)^3 \left(u^{y}\right)^2 + 2 {(\partial_{\tau} u^{x})} {\Pi_{Q}} {\tau} \left(u^{\tau}\right)^4 \left(u^{y}\right)^2\\
  &\quad - {\pi_{Q}^{yy}} \left(u^{\tau}\right)^2 {u^{x}} \left(u^{y}\right)^2 - {(\partial_{\tau} \pi_{Q}^{yy})} {\tau} \left(u^{\tau}\right)^2 {u^{x}} \left(u^{y}\right)^2\\
  &\quad - 2 {(\partial_{\tau} \pi_{Q}{}^{\eta}{}_{\eta})} {\tau} \left(u^{\tau}\right)^2 {u^{x}} \left(u^{y}\right)^2 + 4 {(\partial_{x} u^{x})} {\Pi_{Q}} {\tau} \left(u^{\tau}\right)^3 {u^{x}} \left(u^{y}\right)^2\\
  &\quad + 2 {(\partial_{y} u^{y})} {\Pi_{Q}} {\tau} \left(u^{\tau}\right)^3 {u^{x}} \left(u^{y}\right)^2 - 2 {(\partial_{x} u^{y})} {\pi_{Q}^{xy}} {\tau} \left(u^{\tau}\right)^3 {u^{x}} \left(u^{y}\right)^2\\
  &\quad - 2 {(\partial_{x} u^{x})} {\pi_{Q}^{yy}} {\tau} \left(u^{\tau}\right)^3 {u^{x}} \left(u^{y}\right)^2 - 2 {(\partial_{x} u^{x})} {\pi_{Q}{}^{\eta}{}_{\eta}} {\tau} \left(u^{\tau}\right)^3 {u^{x}} \left(u^{y}\right)^2\\
  &\quad + 2 {\Pi_{Q}} \left(u^{\tau}\right)^4 {u^{x}} \left(u^{y}\right)^2 - {\pi_{Q}{}^{\eta}{}_{\eta}} \left(u^{\tau}\right)^4 {u^{x}} \left(u^{y}\right)^2\\
  &\quad + 2 {(\partial_{\tau} \Pi_{Q})} {\tau} \left(u^{\tau}\right)^4 {u^{x}} \left(u^{y}\right)^2 - 2 {(\partial_{\tau} u^{y})} {\pi_{Q}^{xy}} {\tau} \left(u^{x}\right)^2 \left(u^{y}\right)^2\\
  &\quad + {(\partial_{\tau} u^{x})} {\pi_{Q}^{yy}} {\tau} \left(u^{x}\right)^2 \left(u^{y}\right)^2 + 2 {(\partial_{\tau} u^{x})} {\Pi_{Q}} {\tau} \left(u^{\tau}\right)^2 \left(u^{x}\right)^2 \left(u^{y}\right)^2\\
  &\quad - 2 {(\partial_{\tau} u^{y})} {\pi_{Q}^{xy}} {\tau} \left(u^{\tau}\right)^2 \left(u^{x}\right)^2 \left(u^{y}\right)^2 - 3 {(\partial_{\tau} u^{x})} {\pi_{Q}^{yy}} {\tau} \left(u^{\tau}\right)^2 \left(u^{x}\right)^2 \left(u^{y}\right)^2\\
  &\quad - 2 {(\partial_{\tau} u^{x})} {\pi_{Q}{}^{\eta}{}_{\eta}} {\tau} \left(u^{\tau}\right)^2 \left(u^{x}\right)^2 \left(u^{y}\right)^2 + 2 {(\partial_{x} \Pi_{Q})} {\tau} \left(u^{\tau}\right)^3 \left(u^{x}\right)^2 \left(u^{y}\right)^2\\
  &\quad - {(\partial_{x} \pi_{Q}^{yy})} {\tau} \left(u^{\tau}\right)^3 \left(u^{x}\right)^2 \left(u^{y}\right)^2 - {(\partial_{x} \pi_{Q}{}^{\eta}{}_{\eta})} {\tau} \left(u^{\tau}\right)^3 \left(u^{x}\right)^2 \left(u^{y}\right)^2\\
  &\quad - {\pi_{Q}^{yy}} \left(u^{\tau}\right)^2 \left(u^{x}\right)^3 \left(u^{y}\right)^2 - {(\partial_{\tau} \pi_{Q}^{yy})} {\tau} \left(u^{\tau}\right)^2 \left(u^{x}\right)^3 \left(u^{y}\right)^2\\
  &\quad - {(\partial_{\tau} \pi_{Q}{}^{\eta}{}_{\eta})} {\tau} \left(u^{\tau}\right)^2 \left(u^{x}\right)^3 \left(u^{y}\right)^2 + {(\partial_{\tau} u^{x})} {\pi_{Q}^{yy}} {\tau} \left(u^{x}\right)^4 \left(u^{y}\right)^2\\
  &\quad + 2 {\pi_{Q}^{xy}} \left(u^{\tau}\right)^2 \left(u^{y}\right)^3 + 2 {(\partial_{\tau} \pi_{Q}^{xy})} {\tau} \left(u^{\tau}\right)^2 \left(u^{y}\right)^3\\
  &\quad + 2 {(\partial_{y} u^{x})} {\Pi_{Q}} {\tau} \left(u^{\tau}\right)^3 \left(u^{y}\right)^3 + 2 {(\partial_{x} u^{x})} {\pi_{Q}^{xy}} {\tau} \left(u^{\tau}\right)^3 \left(u^{y}\right)^3\\
  &\quad - 2 {(\partial_{x} u^{y})} {\pi_{Q}{}^{\eta}{}_{\eta}} {\tau} \left(u^{\tau}\right)^3 \left(u^{y}\right)^3 - 2 {(\partial_{\tau} u^{x})} {\pi_{Q}^{xy}} {\tau} {u^{x}} \left(u^{y}\right)^3\\
  &\quad + {(\partial_{\tau} u^{y})} {\pi_{Q}^{yy}} {\tau} {u^{x}} \left(u^{y}\right)^3 + 2 {(\partial_{\tau} u^{y})} {\Pi_{Q}} {\tau} \left(u^{\tau}\right)^2 {u^{x}} \left(u^{y}\right)^3\\
  &\quad + 4 {(\partial_{\tau} u^{x})} {\pi_{Q}^{xy}} {\tau} \left(u^{\tau}\right)^2 {u^{x}} \left(u^{y}\right)^3 + {(\partial_{\tau} u^{y})} {\pi_{Q}{}^{\eta}{}_{\eta}} {\tau} \left(u^{\tau}\right)^2 {u^{x}} \left(u^{y}\right)^3\\
  &\quad + 2 {(\partial_{x} \pi_{Q}^{xy})} {\tau} \left(u^{\tau}\right)^3 {u^{x}} \left(u^{y}\right)^3 + 2 {(\partial_{y} \Pi_{Q})} {\tau} \left(u^{\tau}\right)^3 {u^{x}} \left(u^{y}\right)^3\\
  &\quad + 2 {\pi_{Q}^{xy}} \left(u^{\tau}\right)^2 \left(u^{x}\right)^2 \left(u^{y}\right)^3 + 2 {(\partial_{\tau} \pi_{Q}^{xy})} {\tau} \left(u^{\tau}\right)^2 \left(u^{x}\right)^2 \left(u^{y}\right)^3\\
  &\quad - 2 {(\partial_{\tau} u^{x})} {\pi_{Q}^{xy}} {\tau} \left(u^{x}\right)^3 \left(u^{y}\right)^3 + {(\partial_{\tau} u^{y})} {\pi_{Q}^{yy}} {\tau} \left(u^{x}\right)^3 \left(u^{y}\right)^3\\
  &\quad - 2 {(\partial_{\tau} u^{y})} {\pi_{Q}^{xy}} {\tau} \left(u^{y}\right)^4 + {(\partial_{\tau} u^{y})} {\pi_{Q}^{xy}} {\tau} \left(u^{\tau}\right)^2 \left(u^{y}\right)^4\\
  &\quad - {(\partial_{\tau} u^{x})} {\pi_{Q}{}^{\eta}{}_{\eta}} {\tau} \left(u^{\tau}\right)^2 \left(u^{y}\right)^4 + {(\partial_{x} \Pi_{Q})} {\tau} \left(u^{\tau}\right)^3 \left(u^{y}\right)^4\\
  &\quad - {(\partial_{x} \pi_{Q}{}^{\eta}{}_{\eta})} {\tau} \left(u^{\tau}\right)^3 \left(u^{y}\right)^4 + {(\partial_{y} \pi_{Q}^{xy})} {\tau} \left(u^{\tau}\right)^3 \left(u^{y}\right)^4\\
  &\quad + {(\partial_{\tau} u^{x})} {\Pi_{Q}} {\tau} \left(u^{\tau}\right)^4 \left(u^{y}\right)^4 - {(\partial_{\tau} \pi_{Q}{}^{\eta}{}_{\eta})} {\tau} \left(u^{\tau}\right)^2 {u^{x}} \left(u^{y}\right)^4\\
  &\quad + 2 {(\partial_{x} u^{x})} {\Pi_{Q}} {\tau} \left(u^{\tau}\right)^3 {u^{x}} \left(u^{y}\right)^4 + {(\partial_{y} u^{y})} {\Pi_{Q}} {\tau} \left(u^{\tau}\right)^3 {u^{x}} \left(u^{y}\right)^4\\
  &\quad + {\Pi_{Q}} \left(u^{\tau}\right)^4 {u^{x}} \left(u^{y}\right)^4 + {(\partial_{\tau} \Pi_{Q})} {\tau} \left(u^{\tau}\right)^4 {u^{x}} \left(u^{y}\right)^4\\
  &\quad - 2 {(\partial_{\tau} u^{y})} {\pi_{Q}^{xy}} {\tau} \left(u^{x}\right)^2 \left(u^{y}\right)^4 + {(\partial_{\tau} u^{x})} {\Pi_{Q}} {\tau} \left(u^{\tau}\right)^2 \left(u^{x}\right)^2 \left(u^{y}\right)^4\\
  &\quad + {(\partial_{x} \Pi_{Q})} {\tau} \left(u^{\tau}\right)^3 \left(u^{x}\right)^2 \left(u^{y}\right)^4 + {\pi_{Q}^{xy}} \left(u^{\tau}\right)^2 \left(u^{y}\right)^5\\
  &\quad + {(\partial_{\tau} \pi_{Q}^{xy})} {\tau} \left(u^{\tau}\right)^2 \left(u^{y}\right)^5 + {(\partial_{y} u^{x})} {\Pi_{Q}} {\tau} \left(u^{\tau}\right)^3 \left(u^{y}\right)^5\\
  &\quad - {(\partial_{\tau} u^{x})} {\pi_{Q}^{xy}} {\tau} {u^{x}} \left(u^{y}\right)^5 + {(\partial_{\tau} u^{y})} {\Pi_{Q}} {\tau} \left(u^{\tau}\right)^2 {u^{x}} \left(u^{y}\right)^5\\
  &\quad + {(\partial_{y} \Pi_{Q})} {\tau} \left(u^{\tau}\right)^3 {u^{x}} \left(u^{y}\right)^5 - {(\partial_{\tau} u^{y})} {\pi_{Q}^{xy}} {\tau} \left(u^{y}\right)^6\Bigr) \\
\nabla_{\mu}\delta M^{\mu {y}} &= {(\partial_{x} \pi_{Q}^{xy})} + {(\partial_{y} \pi_{Q}^{yy})} + {(\partial_{\tau} u^{y})} {\Pi_{Q}} {u^{\tau}} + {(\partial_{x} u^{y})} {\Pi_{Q}} {u^{x}} + {(\partial_{x} u^{x})} {\Pi_{Q}} {u^{y}} + 2 {(\partial_{y} u^{y})} {\Pi_{Q}} {u^{y}}\\
  &\quad + {(\partial_{\tau} \Pi_{Q})} {u^{\tau}} {u^{y}} + {(\partial_{x} \Pi_{Q})} {u^{x}} {u^{y}} + \frac{{(\partial_{\tau} \pi_{Q}^{xy})} {u^{x}}+{(\partial_{\tau} \pi_{Q}^{yy})} {u^{y}}+{(\partial_{\tau} u^{x})} {\pi_{Q}^{xy}}+{(\partial_{\tau} u^{y})} {\pi_{Q}^{yy}}}{{u^{\tau}}}\\
  &\quad + \frac{{\Pi_{Q}} {u^{y}} ({(\partial_{\tau} u^{x})} {u^{x}}+{(\partial_{\tau} u^{y})} {u^{y}})}{{u^{\tau}}} - \frac{({(\partial_{\tau} u^{x})} {u^{x}}+{(\partial_{\tau} u^{y})} {u^{y}}) ({\pi_{Q}^{xy}} {u^{x}}+{\pi_{Q}^{yy}} {u^{y}})}{\left(u^{\tau}\right)^3}\\
  &\quad + \frac{{u^{y}} \left({\Pi_{Q}} \left(u^{\tau}\right)^2+{\pi_{Q}^{yy}}\right)+{\pi_{Q}^{xy}} {u^{x}}}{{\tau} {u^{\tau}}} + {(\partial_{y} \Pi_{Q})} \left(\left(u^{y}\right)^2+1\right)
\end{align*}
\[
E_{a}^{\mathrm{S6}} = E_{a}^{\mathrm{S5}} + \nabla_{\mu}\delta M^{\mu a},\quad a = \tau, x, y,\qquad E_{\nu^{i}}^{\mathrm{S6}} = E_{\nu^{i}} + c_{M}\Delta^{i}{}_{\lambda}\nabla_{\mu}\delta M^{\mu\lambda}
\]
Conformal:

\[
P = a_{\mathrm{SB}}T^{4},\quad P' = 4a_{\mathrm{SB}}T^{3},\quad P'' = 12a_{\mathrm{SB}}T^{2},\quad \Pi = \zeta = 0,\quad\text{drop the }\Pi\text{ row and column}
\]
\subsection{Result}

The 1+1D charm systems and the 1+1D conformal two-way system in full; the generic 1+1D two-way system and every 2+1D system as the pattern of nonzero entries ($\bullet$), rows and columns in the field order given, with their entries in the appendix.

\subsection{1+1D cylindrical}

\subsubsection{S3, charm first moment on a background}

\texttt{one\_\allowbreak{}d\_\allowbreak{}HQ\_\allowbreak{}BG\_\allowbreak{}consistent\_\allowbreak{}matrix} in \texttt{src/Matrix/1d\_\allowbreak{}viscous\_\allowbreak{}HQ\_\allowbreak{}cilindrical\_\allowbreak{}BG.jl}, fields:

\[
\phi = ({\alpha}, {\nu^{r}})
\]
\begin{align*}
(A_t)_{11} &= {(\partial_{\alpha} n)} {u^{\tau}} \\
(A_t)_{12} &= \frac{{u^{r}}}{{u^{\tau}}} \\
(A_t)_{21} &= {(\partial_{\alpha} n)} {D_{s}} {u^{r}} {u^{\tau}} \\
(A_t)_{22} &= \frac{{D_{s}} {h} {u^{\tau}}}{{T}} \\
(A_r)_{11} &= {(\partial_{\alpha} n)} {u^{r}} \\
(A_r)_{12} &= 1 \\
(A_r)_{21} &= {(\partial_{\alpha} n)} {D_{s}} \left(u^{\tau}\right)^2 \\
(A_r)_{22} &= \frac{{D_{s}} {h} {u^{r}}}{{T}} \\
S_{1} &= {(\partial_{r} u^{r})} {n} + {(\partial_{T} n)} {(\partial_{r} T)} {u^{r}} + \frac{{n} {u^{r}}+{\nu^{r}}}{{r}} - \frac{{(\partial_{\tau} u^{r})} {\nu^{r}} \left(u^{r}\right)^2}{\left(u^{\tau}\right)^3} + \frac{{(\partial_{\tau} u^{r})} {\nu^{r}}}{{u^{\tau}}} + \frac{{(\partial_{\tau} u^{r})} {n} {u^{r}}}{{u^{\tau}}}\\
  &\quad + \frac{{n} \left(u^{r}\right)^2+{n}+{\nu^{r}} {u^{r}}}{{\tau} {u^{\tau}}} + {(\partial_{T} n)} {(\partial_{\tau} T)} {u^{\tau}} \\
S_{2} &= {\nu^{r}}\\
  &\quad + \frac{1}{{T} {u^{\tau}}}{D_{s}}\Bigl(\left(u^{\tau}\right)^2 ({(\partial_{T} n)} {T}+{n}) ({(\partial_{r} T)} {u^{\tau}}+{(\partial_{\tau} T)} {u^{r}})\\
  &\quad + {h'(T)} {\nu^{r}} \left({(\partial_{r} T)} {u^{r}} {u^{\tau}}+{(\partial_{\tau} T)} \left(u^{r}\right)^2+{(\partial_{\tau} T)}\right)\Bigr)\\
  &\quad + \frac{1}{{r} {T} {\tau} \left(u^{\tau}\right)^3}{D_{s}} {h}\Bigl({\nu^{r}} {\tau} {u^{r}} \left(u^{\tau}\right)^3 + {n} {r} {\tau} \left(u^{\tau}\right)^2 \left({(\partial_{r} u^{r})} {u^{r}} {u^{\tau}}+{(\partial_{\tau} u^{r})} \left(u^{r}\right)^2+{(\partial_{\tau} u^{r})}\right)\\
  &\quad + {\nu^{r}} {r}\Bigl(1 + \left(u^{r}\right)^4 + 2 {(\partial_{r} u^{r})} {\tau} {u^{\tau}} + {(\partial_{\tau} u^{r})} {\tau} {u^{r}} \left(u^{\tau}\right)^2\\
  &\quad + \left(u^{r}\right)^2 ({(\partial_{r} u^{r})} {\tau} {u^{\tau}}+2)\Bigr)\Bigr)
\end{align*}
\subsubsection{S4, charm first and second moment on a background}

\texttt{one\_\allowbreak{}d\_\allowbreak{}HQ\_\allowbreak{}BG\_\allowbreak{}m2\_\allowbreak{}matrix} in \texttt{src/Matrix/1d\_\allowbreak{}viscous\_\allowbreak{}HQ\_\allowbreak{}cilindrical\_\allowbreak{}BG\_\allowbreak{}m2.jl}, fields:

\[
\phi = ({\alpha}, {\nu^{r}}, {\pi_{Q}{}^{\phi}{}_{\phi}}, {\pi_{Q}{}^{\eta}{}_{\eta}}, {\Pi_{Q}})
\]
\begin{align*}
(A_t)_{11} &= {(\partial_{\alpha} n)} {u^{\tau}} \\
(A_t)_{12} &= \frac{{u^{r}}}{{u^{\tau}}} \\
(A_t)_{21} &= {(\partial_{\alpha} n)} {D_{s}} {u^{r}} {u^{\tau}} \\
(A_t)_{22} &= \frac{{D_{s}} {h} {u^{\tau}}}{{T}} \\
(A_t)_{32} &= -\frac{{\mathcal{B}} {D_{s}} {u^{r}}}{3 {T} {u^{\tau}}} \\
(A_t)_{33} &= \frac{{\mathcal{C}} {D_{s}} {u^{\tau}}}{2 {T}} \\
(A_t)_{42} &= -\frac{{\mathcal{B}} {D_{s}} {u^{r}}}{3 {T} {u^{\tau}}} \\
(A_t)_{44} &= \frac{{\mathcal{C}} {D_{s}} {u^{\tau}}}{2 {T}} \\
(A_t)_{51} &= \frac{{(\partial_{\alpha} I_{31})} {D_{s}} {u^{\tau}}}{2 {T}} \\
(A_t)_{52} &= \frac{5 {\mathcal{B}} {D_{s}} {u^{r}}}{6 {T} {u^{\tau}}} \\
(A_t)_{55} &= \frac{{\mathcal{C}} {D_{s}} {u^{\tau}}}{2 {T}} \\
(A_r)_{11} &= {(\partial_{\alpha} n)} {u^{r}} \\
(A_r)_{12} &= 1 \\
(A_r)_{21} &= {(\partial_{\alpha} n)} {D_{s}} \left(u^{\tau}\right)^2 \\
(A_r)_{22} &= \frac{{D_{s}} {h} {u^{r}}}{{T}} \\
(A_r)_{32} &= -\frac{{\mathcal{B}} {D_{s}}}{3 {T}} \\
(A_r)_{33} &= \frac{{\mathcal{C}} {D_{s}} {u^{r}}}{2 {T}} \\
(A_r)_{42} &= -\frac{{\mathcal{B}} {D_{s}}}{3 {T}} \\
(A_r)_{44} &= \frac{{\mathcal{C}} {D_{s}} {u^{r}}}{2 {T}} \\
(A_r)_{51} &= \frac{{(\partial_{\alpha} I_{31})} {D_{s}} {u^{r}}}{2 {T}} \\
(A_r)_{52} &= \frac{5 {\mathcal{B}} {D_{s}}}{6 {T}} \\
(A_r)_{55} &= \frac{{\mathcal{C}} {D_{s}} {u^{r}}}{2 {T}} \\
S_{1} &= {(\partial_{r} u^{r})} {n} + {(\partial_{T} n)} {(\partial_{r} T)} {u^{r}} + \frac{{n} {u^{r}}+{\nu^{r}}}{{r}} - \frac{{(\partial_{\tau} u^{r})} {\nu^{r}} \left(u^{r}\right)^2}{\left(u^{\tau}\right)^3} + \frac{{(\partial_{\tau} u^{r})} {\nu^{r}}}{{u^{\tau}}} + \frac{{(\partial_{\tau} u^{r})} {n} {u^{r}}}{{u^{\tau}}}\\
  &\quad + \frac{{n} \left(u^{r}\right)^2+{n}+{\nu^{r}} {u^{r}}}{{\tau} {u^{\tau}}} + {(\partial_{T} n)} {(\partial_{\tau} T)} {u^{\tau}} \\
S_{2} &= {\nu^{r}}\\
  &\quad + \frac{1}{{T} {u^{\tau}}}{D_{s}}\Bigl(\left(u^{\tau}\right)^2 ({(\partial_{T} n)} {T}+{n}) ({(\partial_{r} T)} {u^{\tau}}+{(\partial_{\tau} T)} {u^{r}})\\
  &\quad + {h'(T)} {\nu^{r}} \left({(\partial_{r} T)} {u^{r}} {u^{\tau}}+{(\partial_{\tau} T)} \left(u^{r}\right)^2+{(\partial_{\tau} T)}\right)\Bigr)\\
  &\quad + \frac{1}{{r} {T} {\tau} \left(u^{\tau}\right)^3}{D_{s}} {h}\Bigl({\nu^{r}} {\tau} {u^{r}} \left(u^{\tau}\right)^3 + {n} {r} {\tau} \left(u^{\tau}\right)^2 \left({(\partial_{r} u^{r})} {u^{r}} {u^{\tau}}+{(\partial_{\tau} u^{r})} \left(u^{r}\right)^2+{(\partial_{\tau} u^{r})}\right)\\
  &\quad + {\nu^{r}} {r}\Bigl(1 + \left(u^{r}\right)^4 + 2 {(\partial_{r} u^{r})} {\tau} {u^{\tau}} + {(\partial_{\tau} u^{r})} {\tau} {u^{r}} \left(u^{\tau}\right)^2\\
  &\quad + \left(u^{r}\right)^2 ({(\partial_{r} u^{r})} {\tau} {u^{\tau}}+2)\Bigr)\Bigr) \\
S_{3} &= {\pi_{Q}{}^{\phi}{}_{\phi}}\\
  &\quad - \frac{1}{6 {r} {T} {\tau} \left(u^{\tau}\right)^3}{D_{s}}\Bigl(2 {\mathcal{A}} {(\partial_{\tau} u^{r})} {\nu^{r}} {r} {\tau} + 2 {\mathcal{B}} {(\partial_{\tau} u^{r})} {\nu^{r}} {r} {\tau} - 3 {\mathcal{C}'(T)} {(\partial_{\tau} T)} {\pi_{Q}{}^{\phi}{}_{\phi}} {r} {\tau}\\
  &\quad + 2 {\mathcal{B}} {\nu^{r}} {r} {u^{r}} + 2 {\mathcal{B}'(T)} {(\partial_{\tau} T)} {\nu^{r}} {r} {\tau} {u^{r}} + 2 {\mathcal{A}} {(\partial_{\tau} u^{r})} {\nu^{r}} {r} {\tau} \left(u^{r}\right)^2\\
  &\quad - 6 {\mathcal{C}'(T)} {(\partial_{\tau} T)} {\pi_{Q}{}^{\phi}{}_{\phi}} {r} {\tau} \left(u^{r}\right)^2 + 2 {\mathcal{B}} {\nu^{r}} {r} \left(u^{r}\right)^3\\
  &\quad + 2 {\mathcal{B}'(T)} {(\partial_{\tau} T)} {\nu^{r}} {r} {\tau} \left(u^{r}\right)^3 - 3 {\mathcal{C}'(T)} {(\partial_{\tau} T)} {\pi_{Q}{}^{\phi}{}_{\phi}} {r} {\tau} \left(u^{r}\right)^4 - 4 {\mathcal{B}} {\nu^{r}} {\tau} {u^{\tau}}\\
  &\quad + 2 {\mathcal{B}'(T)} {(\partial_{r} T)} {\nu^{r}} {r} {\tau} {u^{\tau}} + 2 {\mathcal{A}} {(\partial_{r} u^{r})} {\nu^{r}} {r} {\tau} {u^{r}} {u^{\tau}}\\
  &\quad - 3 {\mathcal{C}'(T)} {(\partial_{r} T)} {\pi_{Q}{}^{\phi}{}_{\phi}} {r} {\tau} {u^{r}} {u^{\tau}} - 4 {\mathcal{B}} {\nu^{r}} {\tau} \left(u^{r}\right)^2 {u^{\tau}}\\
  &\quad + 2 {\mathcal{B}'(T)} {(\partial_{r} T)} {\nu^{r}} {r} {\tau} \left(u^{r}\right)^2 {u^{\tau}} - 3 {\mathcal{C}'(T)} {(\partial_{r} T)} {\pi_{Q}{}^{\phi}{}_{\phi}} {r} {\tau} \left(u^{r}\right)^3 {u^{\tau}}\\
  &\quad + 2 {I_{31}} \left(u^{\tau}\right)^2 \left({r} \left({(\partial_{r} u^{r})} {\tau} {u^{\tau}}+{(\partial_{\tau} u^{r})} {\tau} {u^{r}}+\left(u^{r}\right)^2+1\right)-2 {\tau} {u^{r}} {u^{\tau}}\right)\\
  &\quad + {\mathcal{C}} \left(u^{\tau}\right)^2\Bigl(- 4 {\Pi_{Q}} {\tau} {u^{r}} {u^{\tau}} - 7 {\pi_{Q}{}^{\phi}{}_{\phi}} {\tau} {u^{r}} {u^{\tau}}\\
  &\quad - 2 {\pi_{Q}{}^{\eta}{}_{\eta}} {r} \left({(\partial_{r} u^{r})} {\tau} {u^{\tau}}+{(\partial_{\tau} u^{r})} {\tau} {u^{r}}-\left(u^{r}\right)^2-1\right)\\
  &\quad + 2 {\Pi_{Q}} {r} \left({(\partial_{r} u^{r})} {\tau} {u^{\tau}}+{(\partial_{\tau} u^{r})} {\tau} {u^{r}}+\left(u^{r}\right)^2+1\right)\\
  &\quad - {\pi_{Q}{}^{\phi}{}_{\phi}} {r} \left(5 {(\partial_{r} u^{r})} {\tau} {u^{\tau}}+5 {(\partial_{\tau} u^{r})} {\tau} {u^{r}}+3 \left(u^{r}\right)^2+3\right)\Bigr)\Bigr) \\
S_{4} &= {\pi_{Q}{}^{\eta}{}_{\eta}}\\
  &\quad + \frac{1}{6 {r} {T} {\tau} \left(u^{\tau}\right)^3}{D_{s}}\Bigl(- 2 {\mathcal{A}} {(\partial_{\tau} u^{r})} {\nu^{r}} {r} {\tau} - 2 {\mathcal{B}} {(\partial_{\tau} u^{r})} {\nu^{r}} {r} {\tau} + 3 {\mathcal{C}'(T)} {(\partial_{\tau} T)} {\pi_{Q}{}^{\eta}{}_{\eta}} {r} {\tau}\\
  &\quad + 4 {\mathcal{B}} {\nu^{r}} {r} {u^{r}} - 2 {\mathcal{B}'(T)} {(\partial_{\tau} T)} {\nu^{r}} {r} {\tau} {u^{r}} - 2 {\mathcal{A}} {(\partial_{\tau} u^{r})} {\nu^{r}} {r} {\tau} \left(u^{r}\right)^2\\
  &\quad + 6 {\mathcal{C}'(T)} {(\partial_{\tau} T)} {\pi_{Q}{}^{\eta}{}_{\eta}} {r} {\tau} \left(u^{r}\right)^2 + 4 {\mathcal{B}} {\nu^{r}} {r} \left(u^{r}\right)^3\\
  &\quad - 2 {\mathcal{B}'(T)} {(\partial_{\tau} T)} {\nu^{r}} {r} {\tau} \left(u^{r}\right)^3 + 3 {\mathcal{C}'(T)} {(\partial_{\tau} T)} {\pi_{Q}{}^{\eta}{}_{\eta}} {r} {\tau} \left(u^{r}\right)^4 - 2 {\mathcal{B}} {\nu^{r}} {\tau} {u^{\tau}}\\
  &\quad - 2 {\mathcal{B}'(T)} {(\partial_{r} T)} {\nu^{r}} {r} {\tau} {u^{\tau}} - 2 {\mathcal{A}} {(\partial_{r} u^{r})} {\nu^{r}} {r} {\tau} {u^{r}} {u^{\tau}}\\
  &\quad + 3 {\mathcal{C}'(T)} {(\partial_{r} T)} {\pi_{Q}{}^{\eta}{}_{\eta}} {r} {\tau} {u^{r}} {u^{\tau}} - 2 {\mathcal{B}} {\nu^{r}} {\tau} \left(u^{r}\right)^2 {u^{\tau}}\\
  &\quad - 2 {\mathcal{B}'(T)} {(\partial_{r} T)} {\nu^{r}} {r} {\tau} \left(u^{r}\right)^2 {u^{\tau}} + 3 {\mathcal{C}'(T)} {(\partial_{r} T)} {\pi_{Q}{}^{\eta}{}_{\eta}} {r} {\tau} \left(u^{r}\right)^3 {u^{\tau}}\\
  &\quad - 2 {I_{31}} \left(u^{\tau}\right)^2 \left({r} \left({(\partial_{r} u^{r})} {\tau} {u^{\tau}}+{(\partial_{\tau} u^{r})} {\tau} {u^{r}}-2 \left(u^{r}\right)^2-2\right)+{\tau} {u^{r}} {u^{\tau}}\right)\\
  &\quad + {\mathcal{C}} \left(u^{\tau}\right)^2\Bigl(2 {\pi_{Q}{}^{\phi}{}_{\phi}} {\tau} ({u^{\tau}} ({(\partial_{r} u^{r})} {r}-{u^{r}})+{(\partial_{\tau} u^{r})} {r} {u^{r}})\\
  &\quad - 2{\Pi_{Q}}\Bigl({\tau} {u^{r}} {u^{\tau}}\\
  &\quad + {r} \left({(\partial_{r} u^{r})} {\tau} {u^{\tau}}+{(\partial_{\tau} u^{r})} {\tau} {u^{r}}-2 \left(u^{r}\right)^2-2\right)\Bigr)\\
  &\quad + {\pi_{Q}{}^{\eta}{}_{\eta}}\Bigl(3 {\tau} {u^{r}} {u^{\tau}}\\
  &\quad + {r} \left(5 {(\partial_{r} u^{r})} {\tau} {u^{\tau}}+5 {(\partial_{\tau} u^{r})} {\tau} {u^{r}}+7 \left(u^{r}\right)^2+7\right)\Bigr)\Bigr)\Bigr) \\
S_{5} &= {\Pi_{Q}}\\
  &\quad + \frac{1}{6 {r} {T} {\tau} \left(u^{\tau}\right)^3}{D_{s}}\Bigl(3 {(\partial_{T} I_{31})} {(\partial_{\tau} T)} {r} {\tau} + 2 {\mathcal{A}} {(\partial_{\tau} u^{r})} {\nu^{r}} {r} {\tau} + 5 {\mathcal{B}} {(\partial_{\tau} u^{r})} {\nu^{r}} {r} {\tau}\\
  &\quad + 3 {\mathcal{C}'(T)} {(\partial_{\tau} T)} {\Pi_{Q}} {r} {\tau} + 5 {\mathcal{B}} {\nu^{r}} {r} {u^{r}} + 5 {\mathcal{B}'(T)} {(\partial_{\tau} T)} {\nu^{r}} {r} {\tau} {u^{r}}\\
  &\quad + 6 {(\partial_{T} I_{31})} {(\partial_{\tau} T)} {r} {\tau} \left(u^{r}\right)^2 + 2 {\mathcal{A}} {(\partial_{\tau} u^{r})} {\nu^{r}} {r} {\tau} \left(u^{r}\right)^2\\
  &\quad + 6 {\mathcal{C}'(T)} {(\partial_{\tau} T)} {\Pi_{Q}} {r} {\tau} \left(u^{r}\right)^2 + 5 {\mathcal{B}} {\nu^{r}} {r} \left(u^{r}\right)^3 + 5 {\mathcal{B}'(T)} {(\partial_{\tau} T)} {\nu^{r}} {r} {\tau} \left(u^{r}\right)^3\\
  &\quad + 3 {(\partial_{T} I_{31})} {(\partial_{\tau} T)} {r} {\tau} \left(u^{r}\right)^4 + 3 {\mathcal{C}'(T)} {(\partial_{\tau} T)} {\Pi_{Q}} {r} {\tau} \left(u^{r}\right)^4 + 5 {\mathcal{B}} {\nu^{r}} {\tau} {u^{\tau}}\\
  &\quad + 5 {\mathcal{B}'(T)} {(\partial_{r} T)} {\nu^{r}} {r} {\tau} {u^{\tau}} + 3 {(\partial_{T} I_{31})} {(\partial_{r} T)} {r} {\tau} {u^{r}} {u^{\tau}}\\
  &\quad + 2 {\mathcal{A}} {(\partial_{r} u^{r})} {\nu^{r}} {r} {\tau} {u^{r}} {u^{\tau}} + 3 {\mathcal{C}'(T)} {(\partial_{r} T)} {\Pi_{Q}} {r} {\tau} {u^{r}} {u^{\tau}} + 5 {\mathcal{B}} {\nu^{r}} {\tau} \left(u^{r}\right)^2 {u^{\tau}}\\
  &\quad + 5 {\mathcal{B}'(T)} {(\partial_{r} T)} {\nu^{r}} {r} {\tau} \left(u^{r}\right)^2 {u^{\tau}} + 3 {(\partial_{T} I_{31})} {(\partial_{r} T)} {r} {\tau} \left(u^{r}\right)^3 {u^{\tau}}\\
  &\quad + 3 {\mathcal{C}'(T)} {(\partial_{r} T)} {\Pi_{Q}} {r} {\tau} \left(u^{r}\right)^3 {u^{\tau}}\\
  &\quad + 5 {I_{31}} \left(u^{\tau}\right)^2 \left({r} \left({(\partial_{r} u^{r})} {\tau} {u^{\tau}}+{(\partial_{\tau} u^{r})} {\tau} {u^{r}}+\left(u^{r}\right)^2+1\right)+{\tau} {u^{r}} {u^{\tau}}\right)\\
  &\quad + {\mathcal{C}} \left(u^{\tau}\right)^2\Bigl(5{\Pi_{Q}}\Bigl({r} \left({(\partial_{r} u^{r})} {\tau} {u^{\tau}}+{(\partial_{\tau} u^{r})} {\tau} {u^{r}}+\left(u^{r}\right)^2+1\right)+{\tau} {u^{r}} {u^{\tau}}\Bigr)\\
  &\quad - 2\Bigl({\pi_{Q}{}^{\eta}{}_{\eta}} {r} \left({(\partial_{r} u^{r})} {\tau} {u^{\tau}}+{(\partial_{\tau} u^{r})} {\tau} {u^{r}}-\left(u^{r}\right)^2-1\right)\\
  &\quad + {\pi_{Q}{}^{\phi}{}_{\phi}} {\tau} ({u^{\tau}} ({(\partial_{r} u^{r})} {r}-{u^{r}})+{(\partial_{\tau} u^{r})} {r} {u^{r}})\Bigr)\Bigr)\Bigr)
\end{align*}
\subsubsection{S5, two-way}

\texttt{one\_\allowbreak{}d\_\allowbreak{}two\_\allowbreak{}way\_\allowbreak{}matrix} in \texttt{src/Matrix/1d\_\allowbreak{}two\_\allowbreak{}way\_\allowbreak{}cilindrical.jl}, fields:

\[
\phi = ({T}, {u^{r}}, {\pi^{\phi}{}_{\phi}}, {\pi^{\eta}{}_{\eta}}, {\Pi}, {\alpha}, {\nu^{r}}, {\pi_{Q}{}^{\phi}{}_{\phi}}, {\pi_{Q}{}^{\eta}{}_{\eta}}, {\Pi_{Q}})
\]
\[
\arraycolsep=2.5pt
A_t = \begin{pmatrix}\bullet & \bullet & \bullet & \bullet & \bullet & \bullet & \bullet & \cdot & \cdot & \cdot \\ \bullet & \bullet & \bullet & \bullet & \bullet & \bullet & \bullet & \cdot & \cdot & \cdot \\ \cdot & \bullet & \bullet & \cdot & \cdot & \cdot & \cdot & \cdot & \cdot & \cdot \\ \cdot & \bullet & \cdot & \bullet & \cdot & \cdot & \cdot & \cdot & \cdot & \cdot \\ \cdot & \bullet & \cdot & \cdot & \bullet & \cdot & \cdot & \cdot & \cdot & \cdot \\ \bullet & \bullet & \cdot & \cdot & \cdot & \bullet & \bullet & \cdot & \cdot & \cdot \\ \bullet & \bullet & \cdot & \cdot & \cdot & \bullet & \bullet & \cdot & \cdot & \cdot \\ \bullet & \bullet & \cdot & \cdot & \cdot & \cdot & \bullet & \bullet & \cdot & \cdot \\ \bullet & \bullet & \cdot & \cdot & \cdot & \cdot & \bullet & \cdot & \bullet & \cdot \\ \bullet & \bullet & \cdot & \cdot & \cdot & \bullet & \bullet & \cdot & \cdot & \bullet\end{pmatrix},\qquad A_r = \begin{pmatrix}\bullet & \bullet & \bullet & \bullet & \bullet & \bullet & \bullet & \cdot & \cdot & \cdot \\ \bullet & \bullet & \bullet & \bullet & \bullet & \bullet & \bullet & \cdot & \cdot & \cdot \\ \cdot & \bullet & \bullet & \cdot & \cdot & \cdot & \cdot & \cdot & \cdot & \cdot \\ \cdot & \bullet & \cdot & \bullet & \cdot & \cdot & \cdot & \cdot & \cdot & \cdot \\ \cdot & \bullet & \cdot & \cdot & \bullet & \cdot & \cdot & \cdot & \cdot & \cdot \\ \bullet & \bullet & \cdot & \cdot & \cdot & \bullet & \bullet & \cdot & \cdot & \cdot \\ \bullet & \bullet & \cdot & \cdot & \cdot & \bullet & \bullet & \cdot & \cdot & \cdot \\ \bullet & \bullet & \cdot & \cdot & \cdot & \cdot & \bullet & \bullet & \cdot & \cdot \\ \bullet & \bullet & \cdot & \cdot & \cdot & \cdot & \bullet & \cdot & \bullet & \cdot \\ \bullet & \bullet & \cdot & \cdot & \cdot & \bullet & \bullet & \cdot & \cdot & \bullet\end{pmatrix}
\]
\[
\arraycolsep=2.5pt
S = \begin{pmatrix}\bullet \\ \bullet \\ \bullet \\ \bullet \\ \bullet \\ \bullet \\ \bullet \\ \bullet \\ \bullet \\ \bullet\end{pmatrix}
\]
\subsubsection{S5, two-way, conformal}

\texttt{one\_\allowbreak{}d\_\allowbreak{}two\_\allowbreak{}way\_\allowbreak{}matrix\_\allowbreak{}conformal} in \texttt{src/Matrix/1d\_\allowbreak{}two\_\allowbreak{}way\_\allowbreak{}cilindrical.jl}, fields:

\[
\phi = ({T}, {u^{r}}, {\pi^{\phi}{}_{\phi}}, {\pi^{\eta}{}_{\eta}}, {\alpha}, {\nu^{r}}, {\pi_{Q}{}^{\phi}{}_{\phi}}, {\pi_{Q}{}^{\eta}{}_{\eta}}, {\Pi_{Q}})
\]
\begin{align*}
(A_t)_{11} &= 2 {h'(T)} {\nu^{r}} {u^{r}} + 4 {a_{\mathrm{SB}}} {T}^3 \left(4 \left(u^{r}\right)^2+3\right) + {n} \left({h'(T)} \left(u^{r}\right)^2+{h'(T)}-1\right)\\
  &\quad + {(\partial_{T} n)} \left({h} \left(u^{r}\right)^2+{h}-{T}\right) \\
(A_t)_{12} &= 2 {h} ({n} {u^{r}}+{\nu^{r}})-2 {u^{r}} \left(-4 {a_{\mathrm{SB}}} {T}^4+{\pi^{\eta}{}_{\eta}}+{\pi^{\phi}{}_{\phi}}\right) \\
(A_t)_{13} &= -\left(u^{r}\right)^2 \\
(A_t)_{14} &= -\left(u^{r}\right)^2 \\
(A_t)_{15} &= {(\partial_{\alpha} n)} \left({h} \left(u^{r}\right)^2+{h}-{T}\right) \\
(A_t)_{16} &= 2 {h} {u^{r}} \\
(A_t)_{21} &= \frac{{u^{r}} \left(u^{\tau}\right)^2 \left(16 {a_{\mathrm{SB}}} {T}^3+{(\partial_{T} n)} {h}\right)+{h'(T)} \left({n} {u^{r}} \left(u^{\tau}\right)^2+2 {\nu^{r}} \left(u^{r}\right)^2+{\nu^{r}}\right)}{{u^{\tau}}} \\
(A_t)_{22} &= \frac{1}{\left(u^{\tau}\right)^3}\Bigl(- \left(2 \left(u^{r}\right)^4+3 \left(u^{r}\right)^2+1\right) \left(-4 {a_{\mathrm{SB}}} {T}^4+{\pi^{\eta}{}_{\eta}}+{\pi^{\phi}{}_{\phi}}\right)\\
  &\quad + {h} \left(2 {n} \left(u^{r}\right)^4+3 {n} \left(u^{r}\right)^2+{n}+2 {\nu^{r}} \left(u^{r}\right)^3+3 {\nu^{r}} {u^{r}}\right)\Bigr) \\
(A_t)_{23} &= -{u^{r}} {u^{\tau}} \\
(A_t)_{24} &= -{u^{r}} {u^{\tau}} \\
(A_t)_{25} &= {(\partial_{\alpha} n)} {h} {u^{r}} {u^{\tau}} \\
(A_t)_{26} &= \frac{{h} \left(2 \left(u^{r}\right)^2+1\right)}{{u^{\tau}}} \\
(A_t)_{32} &= \frac{{u^{r}} (3 {\delta_{\pi\pi}} {\pi^{\phi}{}_{\phi}}-2 {\eta})}{3 {r}^2 {u^{\tau}}} \\
(A_t)_{33} &= \frac{{\tau_{\pi}} {u^{\tau}}}{{r}^2} \\
(A_t)_{42} &= \frac{{u^{r}} (3 {\delta_{\pi\pi}} {\pi^{\eta}{}_{\eta}}-2 {\eta})}{3 {\tau}^2 {u^{\tau}}} \\
(A_t)_{44} &= \frac{{\tau_{\pi}} {u^{\tau}}}{{\tau}^2} \\
(A_t)_{51} &= {(\partial_{T} n)} {u^{\tau}} \\
(A_t)_{52} &= \frac{{n} {u^{r}} \left(u^{\tau}\right)^2+{\nu^{r}}}{\left(u^{\tau}\right)^3} \\
(A_t)_{55} &= {(\partial_{\alpha} n)} {u^{\tau}} \\
(A_t)_{56} &= \frac{{u^{r}}}{{u^{\tau}}} \\
(A_t)_{61} &= \frac{{D_{s}} {u^{\tau}} ({h'(T)} {\nu^{r}}+{u^{r}} ({(\partial_{T} n)} {T}+{n}))}{{T}} \\
(A_t)_{62} &= \frac{{D_{s}} {h} \left({n} \left(u^{r}\right)^2+{n}+{\nu^{r}} {u^{r}}\right)}{{T} {u^{\tau}}} \\
(A_t)_{65} &= {(\partial_{\alpha} n)} {D_{s}} {u^{r}} {u^{\tau}} \\
(A_t)_{66} &= \frac{{D_{s}} {h} {u^{\tau}}}{{T}} \\
(A_t)_{71} &= \frac{{D_{s}} \left(3 {\mathcal{C}'(T)} {\pi_{Q}{}^{\phi}{}_{\phi}} \left(u^{\tau}\right)^2-2 {\mathcal{B}'(T)} {\nu^{r}} {u^{r}}\right)}{6 {T} {u^{\tau}}} \\
(A_t)_{72} &= -\frac{{D_{s}} \left(2 {\mathcal{A}} {\nu^{r}} \left(u^{\tau}\right)^2+2 {\mathcal{B}} {\nu^{r}}+{u^{r}} \left(u^{\tau}\right)^2 (2 {\mathcal{C}} {\Pi_{Q}}-2 {\mathcal{C}} {\pi_{Q}{}^{\eta}{}_{\eta}}-5 {\mathcal{C}} {\pi_{Q}{}^{\phi}{}_{\phi}}+2 {I_{31}})\right)}{6 {T} \left(u^{\tau}\right)^3} \\
(A_t)_{76} &= -\frac{{\mathcal{B}} {D_{s}} {u^{r}}}{3 {T} {u^{\tau}}} \\
(A_t)_{77} &= \frac{{\mathcal{C}} {D_{s}} {u^{\tau}}}{2 {T}} \\
(A_t)_{81} &= \frac{{D_{s}} \left(3 {\mathcal{C}'(T)} {\pi_{Q}{}^{\eta}{}_{\eta}} \left(u^{\tau}\right)^2-2 {\mathcal{B}'(T)} {\nu^{r}} {u^{r}}\right)}{6 {T} {u^{\tau}}} \\
(A_t)_{82} &= -\frac{{D_{s}} \left(2 {\mathcal{A}} {\nu^{r}} \left(u^{\tau}\right)^2+2 {\mathcal{B}} {\nu^{r}}+{u^{r}} \left(u^{\tau}\right)^2 (2 {\mathcal{C}} {\Pi_{Q}}-5 {\mathcal{C}} {\pi_{Q}{}^{\eta}{}_{\eta}}-2 {\mathcal{C}} {\pi_{Q}{}^{\phi}{}_{\phi}}+2 {I_{31}})\right)}{6 {T} \left(u^{\tau}\right)^3} \\
(A_t)_{86} &= -\frac{{\mathcal{B}} {D_{s}} {u^{r}}}{3 {T} {u^{\tau}}} \\
(A_t)_{88} &= \frac{{\mathcal{C}} {D_{s}} {u^{\tau}}}{2 {T}} \\
(A_t)_{91} &= \frac{{D_{s}} \left(5 {\mathcal{B}'(T)} {\nu^{r}} {u^{r}}+3 {\mathcal{C}'(T)} {\Pi_{Q}} \left(u^{\tau}\right)^2+3 {(\partial_{T} I_{31})} \left(u^{\tau}\right)^2\right)}{6 {T} {u^{\tau}}} \\
(A_t)_{92} &= \frac{{D_{s}} \left(2 {\mathcal{A}} {\nu^{r}} \left(u^{\tau}\right)^2+5 {\mathcal{B}} {\nu^{r}}+{u^{r}} \left(u^{\tau}\right)^2 (5 {\mathcal{C}} {\Pi_{Q}}-2 {\mathcal{C}} ({\pi_{Q}{}^{\eta}{}_{\eta}}+{\pi_{Q}{}^{\phi}{}_{\phi}})+5 {I_{31}})\right)}{6 {T} \left(u^{\tau}\right)^3} \\
(A_t)_{95} &= \frac{{(\partial_{\alpha} I_{31})} {D_{s}} {u^{\tau}}}{2 {T}} \\
(A_t)_{96} &= \frac{5 {\mathcal{B}} {D_{s}} {u^{r}}}{6 {T} {u^{\tau}}} \\
(A_t)_{99} &= \frac{{\mathcal{C}} {D_{s}} {u^{\tau}}}{2 {T}} \\
(A_r)_{11} &= \frac{{u^{r}} \left(u^{\tau}\right)^2 \left(16 {a_{\mathrm{SB}}} {T}^3+{(\partial_{T} n)} {h}\right)+{h'(T)} \left({n} {u^{r}} \left(u^{\tau}\right)^2+2 {\nu^{r}} \left(u^{r}\right)^2+{\nu^{r}}\right)}{{u^{\tau}}} \\
(A_r)_{12} &= \frac{1}{\left(u^{\tau}\right)^3}\Bigl(- \left(2 \left(u^{r}\right)^4+3 \left(u^{r}\right)^2+1\right) \left(-4 {a_{\mathrm{SB}}} {T}^4+{\pi^{\eta}{}_{\eta}}+{\pi^{\phi}{}_{\phi}}\right)\\
  &\quad + {h} \left(2 {n} \left(u^{r}\right)^4+3 {n} \left(u^{r}\right)^2+{n}+2 {\nu^{r}} \left(u^{r}\right)^3+3 {\nu^{r}} {u^{r}}\right)\Bigr) \\
(A_r)_{13} &= -{u^{r}} {u^{\tau}} \\
(A_r)_{14} &= -{u^{r}} {u^{\tau}} \\
(A_r)_{15} &= {(\partial_{\alpha} n)} {h} {u^{r}} {u^{\tau}} \\
(A_r)_{16} &= \frac{{h} \left(2 \left(u^{r}\right)^2+1\right)}{{u^{\tau}}} \\
(A_r)_{21} &= 16 {a_{\mathrm{SB}}} {T}^3 \left(u^{r}\right)^2+4 {a_{\mathrm{SB}}} {T}^3+{h'(T)} {n} \left(u^{r}\right)^2+2 {h'(T)} {\nu^{r}} {u^{r}}+{(\partial_{T} n)} {h} \left(u^{r}\right)^2+{(\partial_{T} n)} {T}+{n} \\
(A_r)_{22} &= 2 {h} ({n} {u^{r}}+{\nu^{r}})-2 {u^{r}} \left(-4 {a_{\mathrm{SB}}} {T}^4+{\pi^{\eta}{}_{\eta}}+{\pi^{\phi}{}_{\phi}}\right) \\
(A_r)_{23} &= -\left(u^{r}\right)^2-1 \\
(A_r)_{24} &= -\left(u^{r}\right)^2-1 \\
(A_r)_{25} &= {(\partial_{\alpha} n)} \left({h} \left(u^{r}\right)^2+{T}\right) \\
(A_r)_{26} &= 2 {h} {u^{r}} \\
(A_r)_{32} &= \frac{3 {\delta_{\pi\pi}} {\pi^{\phi}{}_{\phi}}-2 {\eta}}{3 {r}^2} \\
(A_r)_{33} &= \frac{{\tau_{\pi}} {u^{r}}}{{r}^2} \\
(A_r)_{42} &= \frac{3 {\delta_{\pi\pi}} {\pi^{\eta}{}_{\eta}}-2 {\eta}}{3 {\tau}^2} \\
(A_r)_{44} &= \frac{{\tau_{\pi}} {u^{r}}}{{\tau}^2} \\
(A_r)_{51} &= {(\partial_{T} n)} {u^{r}} \\
(A_r)_{52} &= {n} \\
(A_r)_{55} &= {(\partial_{\alpha} n)} {u^{r}} \\
(A_r)_{56} &= 1 \\
(A_r)_{61} &= \frac{{D_{s}} \left({h'(T)} {\nu^{r}} {u^{r}}+{(\partial_{T} n)} {T} \left(u^{\tau}\right)^2+{n} \left(u^{\tau}\right)^2\right)}{{T}} \\
(A_r)_{62} &= \frac{{D_{s}} {h} \left({n} {u^{r}} \left(u^{\tau}\right)^2+{\nu^{r}} \left(\left(u^{r}\right)^2+2\right)\right)}{{T} \left(u^{\tau}\right)^2} \\
(A_r)_{65} &= {(\partial_{\alpha} n)} {D_{s}} \left(u^{\tau}\right)^2 \\
(A_r)_{66} &= \frac{{D_{s}} {h} {u^{r}}}{{T}} \\
(A_r)_{71} &= -\frac{{D_{s}} (2 {\mathcal{B}'(T)} {\nu^{r}}-3 {\mathcal{C}'(T)} {\pi_{Q}{}^{\phi}{}_{\phi}} {u^{r}})}{6 {T}} \\
(A_r)_{72} &= -\frac{{D_{s}} \left(2 {\mathcal{A}} {\nu^{r}} {u^{r}}+{\mathcal{C}} \left(u^{\tau}\right)^2 (2 {\Pi_{Q}}-2 {\pi_{Q}{}^{\eta}{}_{\eta}}-5 {\pi_{Q}{}^{\phi}{}_{\phi}})+2 {I_{31}} \left(u^{\tau}\right)^2\right)}{6 {T} \left(u^{\tau}\right)^2} \\
(A_r)_{76} &= -\frac{{\mathcal{B}} {D_{s}}}{3 {T}} \\
(A_r)_{77} &= \frac{{\mathcal{C}} {D_{s}} {u^{r}}}{2 {T}} \\
(A_r)_{81} &= -\frac{{D_{s}} (2 {\mathcal{B}'(T)} {\nu^{r}}-3 {\mathcal{C}'(T)} {\pi_{Q}{}^{\eta}{}_{\eta}} {u^{r}})}{6 {T}} \\
(A_r)_{82} &= -\frac{{D_{s}} \left(2 {\mathcal{A}} {\nu^{r}} {u^{r}}+{\mathcal{C}} \left(u^{\tau}\right)^2 (2 {\Pi_{Q}}-5 {\pi_{Q}{}^{\eta}{}_{\eta}}-2 {\pi_{Q}{}^{\phi}{}_{\phi}})+2 {I_{31}} \left(u^{\tau}\right)^2\right)}{6 {T} \left(u^{\tau}\right)^2} \\
(A_r)_{86} &= -\frac{{\mathcal{B}} {D_{s}}}{3 {T}} \\
(A_r)_{88} &= \frac{{\mathcal{C}} {D_{s}} {u^{r}}}{2 {T}} \\
(A_r)_{91} &= \frac{{D_{s}} (5 {\mathcal{B}'(T)} {\nu^{r}}+3 {u^{r}} ({\mathcal{C}'(T)} {\Pi_{Q}}+{(\partial_{T} I_{31})}))}{6 {T}} \\
(A_r)_{92} &= \frac{{D_{s}} \left(2 {\mathcal{A}} {\nu^{r}} {u^{r}}+{\mathcal{C}} \left(u^{\tau}\right)^2 (5 {\Pi_{Q}}-2 ({\pi_{Q}{}^{\eta}{}_{\eta}}+{\pi_{Q}{}^{\phi}{}_{\phi}}))+5 {I_{31}} \left(u^{\tau}\right)^2\right)}{6 {T} \left(u^{\tau}\right)^2} \\
(A_r)_{95} &= \frac{{(\partial_{\alpha} I_{31})} {D_{s}} {u^{r}}}{2 {T}} \\
(A_r)_{96} &= \frac{5 {\mathcal{B}} {D_{s}}}{6 {T}} \\
(A_r)_{99} &= \frac{{\mathcal{C}} {D_{s}} {u^{r}}}{2 {T}} \\
S_{1} &= \frac{1}{{r} {\tau} {u^{\tau}}}\Bigl(- {\pi^{\phi}{}_{\phi}} {\tau} {u^{r}} + 4 {a_{\mathrm{SB}}} {T}^4 {\tau} {u^{r}} - {\pi^{\phi}{}_{\phi}} {\tau} \left(u^{r}\right)^3 + 4 {a_{\mathrm{SB}}} {T}^4 {\tau} \left(u^{r}\right)^3 + 4 {a_{\mathrm{SB}}} {r} {T}^4 {u^{\tau}}\\
  &\quad - {\pi^{\phi}{}_{\phi}} {r} \left(u^{r}\right)^2 {u^{\tau}} + 4 {a_{\mathrm{SB}}} {r} {T}^4 \left(u^{r}\right)^2 {u^{\tau}} - {\pi^{\eta}{}_{\eta}} \left({r} \left(\left(u^{r}\right)^2-1\right) {u^{\tau}}+{\tau} \left(\left(u^{r}\right)^3+{u^{r}}\right)\right)\\
  &\quad + {h} \left({n} \left(u^{\tau}\right)^2 ({r} {u^{\tau}}+{\tau} {u^{r}})+{\nu^{r}} \left(2 {r} {u^{r}} {u^{\tau}}+2 {\tau} \left(u^{r}\right)^2+{\tau}\right)\right)\Bigr) \\
S_{2} &= \frac{1}{{r} {\tau} {u^{\tau}}}\Bigl(- {\pi^{\phi}{}_{\phi}} {r} {u^{r}} + 4 {a_{\mathrm{SB}}} {r} {T}^4 {u^{r}} - {\pi^{\phi}{}_{\phi}} {r} \left(u^{r}\right)^3 + 4 {a_{\mathrm{SB}}} {r} {T}^4 \left(u^{r}\right)^3 - 2 {\pi^{\phi}{}_{\phi}} {\tau} {u^{\tau}} - {\pi^{\phi}{}_{\phi}} {\tau} \left(u^{r}\right)^2 {u^{\tau}}\\
  &\quad + 4 {a_{\mathrm{SB}}} {T}^4 {\tau} \left(u^{r}\right)^2 {u^{\tau}} - {\pi^{\eta}{}_{\eta}} \left(u^{\tau}\right)^2 ({r} {u^{r}}+{\tau} {u^{\tau}}) + {h} {n} {u^{r}} \left({r} \left(u^{r}\right)^2+{r}+{\tau} {u^{r}} {u^{\tau}}\right)\\
  &\quad + {h} {\nu^{r}} \left(2 {r} \left(u^{r}\right)^2+{r}+2 {\tau} {u^{r}} {u^{\tau}}\right)\Bigr) \\
S_{3} &= \frac{3 {\pi^{\phi}{}_{\phi}} ({r} ({\delta_{\pi\pi}} {u^{\tau}}+{\tau})+{\delta_{\pi\pi}} {\tau} {u^{r}})-2 {r} {u^{\tau}} {\eta}+4 {\tau} {u^{r}} {\eta}}{3 {r}^3 {\tau}} \\
S_{4} &= \frac{3 {\pi^{\eta}{}_{\eta}} ({r} ({\delta_{\pi\pi}} {u^{\tau}}+{\tau})+{\delta_{\pi\pi}} {\tau} {u^{r}})+4 {r} {u^{\tau}} {\eta}-2 {\tau} {u^{r}} {\eta}}{3 {r} {\tau}^3} \\
S_{5} &= \frac{{n} {u^{r}}+{\nu^{r}}}{{r}}+\frac{{n} \left(u^{r}\right)^2+{n}+{\nu^{r}} {u^{r}}}{{\tau} {u^{\tau}}} \\
S_{6} &= \frac{{D_{s}} {h} {\nu^{r}} \left({r} \left(u^{r}\right)^2+{r}+{\tau} {u^{r}} {u^{\tau}}\right)}{{r} {T} {\tau} {u^{\tau}}}+{\nu^{r}} \\
S_{7} &= {\pi_{Q}{}^{\phi}{}_{\phi}}\\
  &\quad + \frac{1}{6 {r} {T} {\tau} {u^{\tau}}}{D_{s}}\Bigl(- 2 {\mathcal{B}} {\nu^{r}} ({r} {u^{r}}-2 {\tau} {u^{\tau}}) - 2 {I_{31}} \left({r} \left(u^{r}\right)^2+{r}-2 {\tau} {u^{r}} {u^{\tau}}\right)\\
  &\quad + {\mathcal{C}}\Bigl(7 {\pi_{Q}{}^{\phi}{}_{\phi}} {\tau} {u^{r}} {u^{\tau}} - 2 {\pi_{Q}{}^{\eta}{}_{\eta}} {r} \left(u^{\tau}\right)^2 + 3 {\pi_{Q}{}^{\phi}{}_{\phi}} {r} \left(u^{\tau}\right)^2\\
  &\quad - 2 {\Pi_{Q}} \left({r} \left(u^{r}\right)^2+{r}-2 {\tau} {u^{r}} {u^{\tau}}\right)\Bigr)\Bigr) \\
S_{8} &= {\pi_{Q}{}^{\eta}{}_{\eta}}\\
  &\quad + \frac{1}{6 {r} {T} {\tau} {u^{\tau}}}{D_{s}}\Bigl(- 2 {I_{31}} {\tau} {u^{r}} {u^{\tau}} + 4 {I_{31}} {r} \left(u^{\tau}\right)^2 + 2 {\mathcal{B}} {\nu^{r}} (2 {r} {u^{r}}-{\tau} {u^{\tau}})\\
  &\quad + {\mathcal{C}}\Bigl(- 2 {\Pi_{Q}} {\tau} {u^{r}} {u^{\tau}} + 3 {\pi_{Q}{}^{\eta}{}_{\eta}} {\tau} {u^{r}} {u^{\tau}} - 2 {\pi_{Q}{}^{\phi}{}_{\phi}} {\tau} {u^{r}} {u^{\tau}} + 4 {\Pi_{Q}} {r} \left(u^{\tau}\right)^2\\
  &\quad + 7 {\pi_{Q}{}^{\eta}{}_{\eta}} {r} \left(u^{\tau}\right)^2\Bigr)\Bigr) \\
S_{9} &= {\Pi_{Q}}\\
  &\quad + \frac{1}{6 {r} {T} {\tau} {u^{\tau}}}{D_{s}}\Bigl(5 {\mathcal{B}} {\nu^{r}} ({r} {u^{r}}+{\tau} {u^{\tau}}) + 5 {I_{31}} \left({r} \left(u^{r}\right)^2+{r}+{\tau} {u^{r}} {u^{\tau}}\right)\\
  &\quad + {\mathcal{C}} \left(5 {\Pi_{Q}} \left({r} \left(u^{r}\right)^2+{r}+{\tau} {u^{r}} {u^{\tau}}\right)+2 {\pi_{Q}{}^{\eta}{}_{\eta}} {r} \left(u^{\tau}\right)^2+2 {\pi_{Q}{}^{\phi}{}_{\phi}} {\tau} {u^{r}} {u^{\tau}}\right)\Bigr)
\end{align*}
\subsubsection{S6, fully coupled}

\texttt{one\_\allowbreak{}d\_\allowbreak{}two\_\allowbreak{}way\_\allowbreak{}cM\_\allowbreak{}matrix} in \texttt{src/Matrix/1d\_\allowbreak{}two\_\allowbreak{}way\_\allowbreak{}cM\_\allowbreak{}cilindrical.jl}, fields:

\[
\phi = ({T}, {u^{r}}, {\pi^{\phi}{}_{\phi}}, {\pi^{\eta}{}_{\eta}}, {\Pi}, {\alpha}, {\nu^{r}}, {\pi_{Q}{}^{\phi}{}_{\phi}}, {\pi_{Q}{}^{\eta}{}_{\eta}}, {\Pi_{Q}})
\]
\[
\arraycolsep=2.5pt
A_t = \begin{pmatrix}\bullet & \bullet & \bullet & \bullet & \bullet & \bullet & \bullet & \bullet & \bullet & \bullet \\ \bullet & \bullet & \bullet & \bullet & \bullet & \bullet & \bullet & \bullet & \bullet & \bullet \\ \cdot & \bullet & \bullet & \cdot & \cdot & \cdot & \cdot & \cdot & \cdot & \cdot \\ \cdot & \bullet & \cdot & \bullet & \cdot & \cdot & \cdot & \cdot & \cdot & \cdot \\ \cdot & \bullet & \cdot & \cdot & \bullet & \cdot & \cdot & \cdot & \cdot & \cdot \\ \bullet & \bullet & \cdot & \cdot & \cdot & \bullet & \bullet & \cdot & \cdot & \cdot \\ \bullet & \bullet & \cdot & \cdot & \cdot & \bullet & \bullet & \bullet & \bullet & \bullet \\ \bullet & \bullet & \cdot & \cdot & \cdot & \cdot & \bullet & \bullet & \cdot & \cdot \\ \bullet & \bullet & \cdot & \cdot & \cdot & \cdot & \bullet & \cdot & \bullet & \cdot \\ \bullet & \bullet & \cdot & \cdot & \cdot & \bullet & \bullet & \cdot & \cdot & \bullet\end{pmatrix},\qquad A_r = \begin{pmatrix}\bullet & \bullet & \bullet & \bullet & \bullet & \bullet & \bullet & \bullet & \bullet & \bullet \\ \bullet & \bullet & \bullet & \bullet & \bullet & \bullet & \bullet & \bullet & \bullet & \bullet \\ \cdot & \bullet & \bullet & \cdot & \cdot & \cdot & \cdot & \cdot & \cdot & \cdot \\ \cdot & \bullet & \cdot & \bullet & \cdot & \cdot & \cdot & \cdot & \cdot & \cdot \\ \cdot & \bullet & \cdot & \cdot & \bullet & \cdot & \cdot & \cdot & \cdot & \cdot \\ \bullet & \bullet & \cdot & \cdot & \cdot & \bullet & \bullet & \cdot & \cdot & \cdot \\ \bullet & \bullet & \cdot & \cdot & \cdot & \bullet & \bullet & \bullet & \bullet & \bullet \\ \bullet & \bullet & \cdot & \cdot & \cdot & \cdot & \bullet & \bullet & \cdot & \cdot \\ \bullet & \bullet & \cdot & \cdot & \cdot & \cdot & \bullet & \cdot & \bullet & \cdot \\ \bullet & \bullet & \cdot & \cdot & \cdot & \bullet & \bullet & \cdot & \cdot & \bullet\end{pmatrix}
\]
\[
\arraycolsep=2.5pt
S = \begin{pmatrix}\bullet \\ \bullet \\ \bullet \\ \bullet \\ \bullet \\ \bullet \\ \bullet \\ \bullet \\ \bullet \\ \bullet\end{pmatrix}
\]
\subsubsection{S6, fully coupled, conformal}

\texttt{one\_\allowbreak{}d\_\allowbreak{}two\_\allowbreak{}way\_\allowbreak{}cM\_\allowbreak{}matrix\_\allowbreak{}conformal} in \texttt{src/Matrix/1d\_\allowbreak{}two\_\allowbreak{}way\_\allowbreak{}cM\_\allowbreak{}cilindrical.jl}, fields:

\[
\phi = ({T}, {u^{r}}, {\pi^{\phi}{}_{\phi}}, {\pi^{\eta}{}_{\eta}}, {\alpha}, {\nu^{r}}, {\pi_{Q}{}^{\phi}{}_{\phi}}, {\pi_{Q}{}^{\eta}{}_{\eta}}, {\Pi_{Q}})
\]
\begin{align*}
(A_t)_{11} &= 2 {h'(T)} {\nu^{r}} {u^{r}} + 4 {a_{\mathrm{SB}}} {T}^3 \left(4 \left(u^{r}\right)^2+3\right) + {n} \left({h'(T)} \left(u^{r}\right)^2+{h'(T)}-1\right)\\
  &\quad + {(\partial_{T} n)} \left({h} \left(u^{r}\right)^2+{h}-{T}\right) \\
(A_t)_{12} &= 2 \left({u^{r}} \left(4 {a_{\mathrm{SB}}} {T}^4+{\Pi_{Q}}-{\pi^{\eta}{}_{\eta}}-{\pi_{Q}{}^{\eta}{}_{\eta}}-{\pi^{\phi}{}_{\phi}}-{\pi_{Q}{}^{\phi}{}_{\phi}}\right)+{h} ({n} {u^{r}}+{\nu^{r}})\right) \\
(A_t)_{13} &= -\left(u^{r}\right)^2 \\
(A_t)_{14} &= -\left(u^{r}\right)^2 \\
(A_t)_{15} &= {(\partial_{\alpha} n)} \left({h} \left(u^{r}\right)^2+{h}-{T}\right) \\
(A_t)_{16} &= 2 {h} {u^{r}} \\
(A_t)_{17} &= -\left(u^{r}\right)^2 \\
(A_t)_{18} &= -\left(u^{r}\right)^2 \\
(A_t)_{19} &= \left(u^{r}\right)^2 \\
(A_t)_{21} &= \frac{{u^{r}} \left(u^{\tau}\right)^2 \left(16 {a_{\mathrm{SB}}} {T}^3+{(\partial_{T} n)} {h}\right)+{h'(T)} \left({n} {u^{r}} \left(u^{\tau}\right)^2+2 {\nu^{r}} \left(u^{r}\right)^2+{\nu^{r}}\right)}{{u^{\tau}}} \\
(A_t)_{22} &= \frac{1}{\left(u^{\tau}\right)^3}\Bigl(\left(2 \left(u^{r}\right)^4+3 \left(u^{r}\right)^2+1\right) \left(4 {a_{\mathrm{SB}}} {T}^4+{\Pi_{Q}}-{\pi^{\eta}{}_{\eta}}-{\pi_{Q}{}^{\eta}{}_{\eta}}-{\pi^{\phi}{}_{\phi}}-{\pi_{Q}{}^{\phi}{}_{\phi}}\right)\\
  &\quad + {h} \left(2 {n} \left(u^{r}\right)^4+3 {n} \left(u^{r}\right)^2+{n}+2 {\nu^{r}} \left(u^{r}\right)^3+3 {\nu^{r}} {u^{r}}\right)\Bigr) \\
(A_t)_{23} &= -{u^{r}} {u^{\tau}} \\
(A_t)_{24} &= -{u^{r}} {u^{\tau}} \\
(A_t)_{25} &= {(\partial_{\alpha} n)} {h} {u^{r}} {u^{\tau}} \\
(A_t)_{26} &= \frac{{h} \left(2 \left(u^{r}\right)^2+1\right)}{{u^{\tau}}} \\
(A_t)_{27} &= -{u^{r}} {u^{\tau}} \\
(A_t)_{28} &= -{u^{r}} {u^{\tau}} \\
(A_t)_{29} &= {u^{r}} {u^{\tau}} \\
(A_t)_{32} &= \frac{{u^{r}} (3 {\delta_{\pi\pi}} {\pi^{\phi}{}_{\phi}}-2 {\eta})}{3 {r}^2 {u^{\tau}}} \\
(A_t)_{33} &= \frac{{\tau_{\pi}} {u^{\tau}}}{{r}^2} \\
(A_t)_{42} &= \frac{{u^{r}} (3 {\delta_{\pi\pi}} {\pi^{\eta}{}_{\eta}}-2 {\eta})}{3 {\tau}^2 {u^{\tau}}} \\
(A_t)_{44} &= \frac{{\tau_{\pi}} {u^{\tau}}}{{\tau}^2} \\
(A_t)_{51} &= {(\partial_{T} n)} {u^{\tau}} \\
(A_t)_{52} &= \frac{{n} {u^{r}} \left(u^{\tau}\right)^2+{\nu^{r}}}{\left(u^{\tau}\right)^3} \\
(A_t)_{55} &= {(\partial_{\alpha} n)} {u^{\tau}} \\
(A_t)_{56} &= \frac{{u^{r}}}{{u^{\tau}}} \\
(A_t)_{61} &= \frac{{D_{s}} {u^{\tau}} ({h'(T)} {\nu^{r}}+{u^{r}} ({(\partial_{T} n)} {T}+{n}))}{{T}} \\
(A_t)_{62} &= \frac{{D_{s}} \left({h} \left({n} \left(u^{r}\right)^2+{n}+{\nu^{r}} {u^{r}}\right)+\left(u^{\tau}\right)^2 ({\Pi_{Q}}-{\pi_{Q}{}^{\eta}{}_{\eta}}-{\pi_{Q}{}^{\phi}{}_{\phi}})\right)}{{T} {u^{\tau}}} \\
(A_t)_{65} &= {(\partial_{\alpha} n)} {D_{s}} {u^{r}} {u^{\tau}} \\
(A_t)_{66} &= \frac{{D_{s}} {h} {u^{\tau}}}{{T}} \\
(A_t)_{67} &= -\frac{{D_{s}} {u^{r}} {u^{\tau}}}{{T}} \\
(A_t)_{68} &= -\frac{{D_{s}} {u^{r}} {u^{\tau}}}{{T}} \\
(A_t)_{69} &= \frac{{D_{s}} {u^{r}} {u^{\tau}}}{{T}} \\
(A_t)_{71} &= \frac{{D_{s}} \left(3 {\mathcal{C}'(T)} {\pi_{Q}{}^{\phi}{}_{\phi}} \left(u^{\tau}\right)^2-2 {\mathcal{B}'(T)} {\nu^{r}} {u^{r}}\right)}{6 {T} {u^{\tau}}} \\
(A_t)_{72} &= -\frac{{D_{s}} \left(2 {\mathcal{A}} {\nu^{r}} \left(u^{\tau}\right)^2+2 {\mathcal{B}} {\nu^{r}}+{u^{r}} \left(u^{\tau}\right)^2 (2 {\mathcal{C}} {\Pi_{Q}}-2 {\mathcal{C}} {\pi_{Q}{}^{\eta}{}_{\eta}}-5 {\mathcal{C}} {\pi_{Q}{}^{\phi}{}_{\phi}}+2 {I_{31}})\right)}{6 {T} \left(u^{\tau}\right)^3} \\
(A_t)_{76} &= -\frac{{\mathcal{B}} {D_{s}} {u^{r}}}{3 {T} {u^{\tau}}} \\
(A_t)_{77} &= \frac{{\mathcal{C}} {D_{s}} {u^{\tau}}}{2 {T}} \\
(A_t)_{81} &= \frac{{D_{s}} \left(3 {\mathcal{C}'(T)} {\pi_{Q}{}^{\eta}{}_{\eta}} \left(u^{\tau}\right)^2-2 {\mathcal{B}'(T)} {\nu^{r}} {u^{r}}\right)}{6 {T} {u^{\tau}}} \\
(A_t)_{82} &= -\frac{{D_{s}} \left(2 {\mathcal{A}} {\nu^{r}} \left(u^{\tau}\right)^2+2 {\mathcal{B}} {\nu^{r}}+{u^{r}} \left(u^{\tau}\right)^2 (2 {\mathcal{C}} {\Pi_{Q}}-5 {\mathcal{C}} {\pi_{Q}{}^{\eta}{}_{\eta}}-2 {\mathcal{C}} {\pi_{Q}{}^{\phi}{}_{\phi}}+2 {I_{31}})\right)}{6 {T} \left(u^{\tau}\right)^3} \\
(A_t)_{86} &= -\frac{{\mathcal{B}} {D_{s}} {u^{r}}}{3 {T} {u^{\tau}}} \\
(A_t)_{88} &= \frac{{\mathcal{C}} {D_{s}} {u^{\tau}}}{2 {T}} \\
(A_t)_{91} &= \frac{{D_{s}} \left(5 {\mathcal{B}'(T)} {\nu^{r}} {u^{r}}+3 {\mathcal{C}'(T)} {\Pi_{Q}} \left(u^{\tau}\right)^2+3 {(\partial_{T} I_{31})} \left(u^{\tau}\right)^2\right)}{6 {T} {u^{\tau}}} \\
(A_t)_{92} &= \frac{{D_{s}} \left(2 {\mathcal{A}} {\nu^{r}} \left(u^{\tau}\right)^2+5 {\mathcal{B}} {\nu^{r}}+{u^{r}} \left(u^{\tau}\right)^2 (5 {\mathcal{C}} {\Pi_{Q}}-2 {\mathcal{C}} ({\pi_{Q}{}^{\eta}{}_{\eta}}+{\pi_{Q}{}^{\phi}{}_{\phi}})+5 {I_{31}})\right)}{6 {T} \left(u^{\tau}\right)^3} \\
(A_t)_{95} &= \frac{{(\partial_{\alpha} I_{31})} {D_{s}} {u^{\tau}}}{2 {T}} \\
(A_t)_{96} &= \frac{5 {\mathcal{B}} {D_{s}} {u^{r}}}{6 {T} {u^{\tau}}} \\
(A_t)_{99} &= \frac{{\mathcal{C}} {D_{s}} {u^{\tau}}}{2 {T}} \\
(A_r)_{11} &= \frac{{u^{r}} \left(u^{\tau}\right)^2 \left(16 {a_{\mathrm{SB}}} {T}^3+{(\partial_{T} n)} {h}\right)+{h'(T)} \left({n} {u^{r}} \left(u^{\tau}\right)^2+2 {\nu^{r}} \left(u^{r}\right)^2+{\nu^{r}}\right)}{{u^{\tau}}} \\
(A_r)_{12} &= \frac{1}{\left(u^{\tau}\right)^3}\Bigl(\left(2 \left(u^{r}\right)^4+3 \left(u^{r}\right)^2+1\right) \left(4 {a_{\mathrm{SB}}} {T}^4+{\Pi_{Q}}-{\pi^{\eta}{}_{\eta}}-{\pi_{Q}{}^{\eta}{}_{\eta}}-{\pi^{\phi}{}_{\phi}}-{\pi_{Q}{}^{\phi}{}_{\phi}}\right)\\
  &\quad + {h} \left(2 {n} \left(u^{r}\right)^4+3 {n} \left(u^{r}\right)^2+{n}+2 {\nu^{r}} \left(u^{r}\right)^3+3 {\nu^{r}} {u^{r}}\right)\Bigr) \\
(A_r)_{13} &= -{u^{r}} {u^{\tau}} \\
(A_r)_{14} &= -{u^{r}} {u^{\tau}} \\
(A_r)_{15} &= {(\partial_{\alpha} n)} {h} {u^{r}} {u^{\tau}} \\
(A_r)_{16} &= \frac{{h} \left(2 \left(u^{r}\right)^2+1\right)}{{u^{\tau}}} \\
(A_r)_{17} &= -{u^{r}} {u^{\tau}} \\
(A_r)_{18} &= -{u^{r}} {u^{\tau}} \\
(A_r)_{19} &= {u^{r}} {u^{\tau}} \\
(A_r)_{21} &= 16 {a_{\mathrm{SB}}} {T}^3 \left(u^{r}\right)^2+4 {a_{\mathrm{SB}}} {T}^3+{h'(T)} {n} \left(u^{r}\right)^2+2 {h'(T)} {\nu^{r}} {u^{r}}+{(\partial_{T} n)} {h} \left(u^{r}\right)^2+{(\partial_{T} n)} {T}+{n} \\
(A_r)_{22} &= 2 \left({u^{r}} \left(4 {a_{\mathrm{SB}}} {T}^4+{\Pi_{Q}}-{\pi^{\eta}{}_{\eta}}-{\pi_{Q}{}^{\eta}{}_{\eta}}-{\pi^{\phi}{}_{\phi}}-{\pi_{Q}{}^{\phi}{}_{\phi}}\right)+{h} ({n} {u^{r}}+{\nu^{r}})\right) \\
(A_r)_{23} &= -\left(u^{r}\right)^2-1 \\
(A_r)_{24} &= -\left(u^{r}\right)^2-1 \\
(A_r)_{25} &= {(\partial_{\alpha} n)} \left({h} \left(u^{r}\right)^2+{T}\right) \\
(A_r)_{26} &= 2 {h} {u^{r}} \\
(A_r)_{27} &= -\left(u^{r}\right)^2-1 \\
(A_r)_{28} &= -\left(u^{r}\right)^2-1 \\
(A_r)_{29} &= \left(u^{\tau}\right)^2 \\
(A_r)_{32} &= \frac{3 {\delta_{\pi\pi}} {\pi^{\phi}{}_{\phi}}-2 {\eta}}{3 {r}^2} \\
(A_r)_{33} &= \frac{{\tau_{\pi}} {u^{r}}}{{r}^2} \\
(A_r)_{42} &= \frac{3 {\delta_{\pi\pi}} {\pi^{\eta}{}_{\eta}}-2 {\eta}}{3 {\tau}^2} \\
(A_r)_{44} &= \frac{{\tau_{\pi}} {u^{r}}}{{\tau}^2} \\
(A_r)_{51} &= {(\partial_{T} n)} {u^{r}} \\
(A_r)_{52} &= {n} \\
(A_r)_{55} &= {(\partial_{\alpha} n)} {u^{r}} \\
(A_r)_{56} &= 1 \\
(A_r)_{61} &= \frac{{D_{s}} \left({h'(T)} {\nu^{r}} {u^{r}}+{(\partial_{T} n)} {T} \left(u^{\tau}\right)^2+{n} \left(u^{\tau}\right)^2\right)}{{T}} \\
(A_r)_{62} &= \frac{{D_{s}} \left({h} {n} {u^{r}} \left(u^{\tau}\right)^2+{h} {\nu^{r}} \left(u^{r}\right)^2+2 {h} {\nu^{r}}+{\Pi_{Q}} {u^{r}} \left(u^{\tau}\right)^2-{\pi_{Q}{}^{\eta}{}_{\eta}} {u^{r}} \left(u^{\tau}\right)^2-{\pi_{Q}{}^{\phi}{}_{\phi}} {u^{r}} \left(u^{\tau}\right)^2\right)}{{T} \left(u^{\tau}\right)^2} \\
(A_r)_{65} &= {(\partial_{\alpha} n)} {D_{s}} \left(u^{\tau}\right)^2 \\
(A_r)_{66} &= \frac{{D_{s}} {h} {u^{r}}}{{T}} \\
(A_r)_{67} &= -\frac{{D_{s}} \left(u^{\tau}\right)^2}{{T}} \\
(A_r)_{68} &= -\frac{{D_{s}} \left(u^{\tau}\right)^2}{{T}} \\
(A_r)_{69} &= \frac{{D_{s}} \left(u^{\tau}\right)^2}{{T}} \\
(A_r)_{71} &= -\frac{{D_{s}} (2 {\mathcal{B}'(T)} {\nu^{r}}-3 {\mathcal{C}'(T)} {\pi_{Q}{}^{\phi}{}_{\phi}} {u^{r}})}{6 {T}} \\
(A_r)_{72} &= -\frac{{D_{s}} \left(2 {\mathcal{A}} {\nu^{r}} {u^{r}}+{\mathcal{C}} \left(u^{\tau}\right)^2 (2 {\Pi_{Q}}-2 {\pi_{Q}{}^{\eta}{}_{\eta}}-5 {\pi_{Q}{}^{\phi}{}_{\phi}})+2 {I_{31}} \left(u^{\tau}\right)^2\right)}{6 {T} \left(u^{\tau}\right)^2} \\
(A_r)_{76} &= -\frac{{\mathcal{B}} {D_{s}}}{3 {T}} \\
(A_r)_{77} &= \frac{{\mathcal{C}} {D_{s}} {u^{r}}}{2 {T}} \\
(A_r)_{81} &= -\frac{{D_{s}} (2 {\mathcal{B}'(T)} {\nu^{r}}-3 {\mathcal{C}'(T)} {\pi_{Q}{}^{\eta}{}_{\eta}} {u^{r}})}{6 {T}} \\
(A_r)_{82} &= -\frac{{D_{s}} \left(2 {\mathcal{A}} {\nu^{r}} {u^{r}}+{\mathcal{C}} \left(u^{\tau}\right)^2 (2 {\Pi_{Q}}-5 {\pi_{Q}{}^{\eta}{}_{\eta}}-2 {\pi_{Q}{}^{\phi}{}_{\phi}})+2 {I_{31}} \left(u^{\tau}\right)^2\right)}{6 {T} \left(u^{\tau}\right)^2} \\
(A_r)_{86} &= -\frac{{\mathcal{B}} {D_{s}}}{3 {T}} \\
(A_r)_{88} &= \frac{{\mathcal{C}} {D_{s}} {u^{r}}}{2 {T}} \\
(A_r)_{91} &= \frac{{D_{s}} (5 {\mathcal{B}'(T)} {\nu^{r}}+3 {u^{r}} ({\mathcal{C}'(T)} {\Pi_{Q}}+{(\partial_{T} I_{31})}))}{6 {T}} \\
(A_r)_{92} &= \frac{{D_{s}} \left(2 {\mathcal{A}} {\nu^{r}} {u^{r}}+{\mathcal{C}} \left(u^{\tau}\right)^2 (5 {\Pi_{Q}}-2 ({\pi_{Q}{}^{\eta}{}_{\eta}}+{\pi_{Q}{}^{\phi}{}_{\phi}}))+5 {I_{31}} \left(u^{\tau}\right)^2\right)}{6 {T} \left(u^{\tau}\right)^2} \\
(A_r)_{95} &= \frac{{(\partial_{\alpha} I_{31})} {D_{s}} {u^{r}}}{2 {T}} \\
(A_r)_{96} &= \frac{5 {\mathcal{B}} {D_{s}}}{6 {T}} \\
(A_r)_{99} &= \frac{{\mathcal{C}} {D_{s}} {u^{r}}}{2 {T}} \\
S_{1} &= \frac{1}{{r} {\tau} {u^{\tau}}}\Bigl(- {\pi^{\eta}{}_{\eta}} {\tau} {u^{r}} - {\pi_{Q}{}^{\eta}{}_{\eta}} {\tau} {u^{r}} - {\pi^{\phi}{}_{\phi}} {\tau} {u^{r}} - {\pi_{Q}{}^{\phi}{}_{\phi}} {\tau} {u^{r}} + 4 {a_{\mathrm{SB}}} {T}^4 {\tau} {u^{r}} - {\pi^{\eta}{}_{\eta}} {\tau} \left(u^{r}\right)^3\\
  &\quad - {\pi_{Q}{}^{\eta}{}_{\eta}} {\tau} \left(u^{r}\right)^3 - {\pi^{\phi}{}_{\phi}} {\tau} \left(u^{r}\right)^3 - {\pi_{Q}{}^{\phi}{}_{\phi}} {\tau} \left(u^{r}\right)^3 + 4 {a_{\mathrm{SB}}} {T}^4 {\tau} \left(u^{r}\right)^3 + {\pi^{\eta}{}_{\eta}} {r} {u^{\tau}} + {\pi_{Q}{}^{\eta}{}_{\eta}} {r} {u^{\tau}}\\
  &\quad + 4 {a_{\mathrm{SB}}} {r} {T}^4 {u^{\tau}} - {\pi^{\eta}{}_{\eta}} {r} \left(u^{r}\right)^2 {u^{\tau}} - {\pi_{Q}{}^{\eta}{}_{\eta}} {r} \left(u^{r}\right)^2 {u^{\tau}} - {\pi^{\phi}{}_{\phi}} {r} \left(u^{r}\right)^2 {u^{\tau}} - {\pi_{Q}{}^{\phi}{}_{\phi}} {r} \left(u^{r}\right)^2 {u^{\tau}}\\
  &\quad + 4 {a_{\mathrm{SB}}} {r} {T}^4 \left(u^{r}\right)^2 {u^{\tau}} + {h} {n} \left(u^{\tau}\right)^2 ({r} {u^{\tau}}+{\tau} {u^{r}}) + {\Pi_{Q}} \left(u^{\tau}\right)^2 ({r} {u^{\tau}}+{\tau} {u^{r}})\\
  &\quad + {h} {\nu^{r}} \left(2 {r} {u^{r}} {u^{\tau}}+2 {\tau} \left(u^{r}\right)^2+{\tau}\right)\Bigr) \\
S_{2} &= - \frac{{\Pi_{Q}}+{\pi_{Q}{}^{\phi}{}_{\phi}}}{{r}} - \frac{{n} {T}}{{r}} - \frac{{a_{\mathrm{SB}}} {T}^4+{\pi^{\phi}{}_{\phi}}}{{r}} + \frac{{h} {n} \left(u^{r}\right)^2+2 {h} {\nu^{r}} {u^{r}}+{n} {T}}{{r}}\\
  &\quad + \frac{{u^{r}} {u^{\tau}} ({\Pi_{Q}}-{\pi_{Q}{}^{\eta}{}_{\eta}}-{\pi_{Q}{}^{\phi}{}_{\phi}})}{{\tau}} - \frac{{u^{r}} {u^{\tau}} \left(-4 {a_{\mathrm{SB}}} {T}^4+{\pi^{\eta}{}_{\eta}}+{\pi^{\phi}{}_{\phi}}\right)}{{\tau}}\\
  &\quad + \frac{\left(u^{\tau}\right)^2 ({\Pi_{Q}}-{\pi_{Q}{}^{\eta}{}_{\eta}}-{\pi_{Q}{}^{\phi}{}_{\phi}})}{{r}} + \frac{4 {a_{\mathrm{SB}}} {T}^4 \left(u^{r}\right)^2+{a_{\mathrm{SB}}} {T}^4-\left(u^{\tau}\right)^2 ({\pi^{\eta}{}_{\eta}}+{\pi^{\phi}{}_{\phi}})}{{r}}\\
  &\quad + \frac{{h} \left({n} {u^{r}} \left(u^{\tau}\right)^2+2 {\nu^{r}} \left(u^{r}\right)^2+{\nu^{r}}\right)}{{\tau} {u^{\tau}}} \\
S_{3} &= \frac{3 {\pi^{\phi}{}_{\phi}} ({r} ({\delta_{\pi\pi}} {u^{\tau}}+{\tau})+{\delta_{\pi\pi}} {\tau} {u^{r}})-2 {r} {u^{\tau}} {\eta}+4 {\tau} {u^{r}} {\eta}}{3 {r}^3 {\tau}} \\
S_{4} &= \frac{3 {\pi^{\eta}{}_{\eta}} ({r} ({\delta_{\pi\pi}} {u^{\tau}}+{\tau})+{\delta_{\pi\pi}} {\tau} {u^{r}})+4 {r} {u^{\tau}} {\eta}-2 {\tau} {u^{r}} {\eta}}{3 {r} {\tau}^3} \\
S_{5} &= \frac{{n} {u^{r}}+{\nu^{r}}}{{r}}+\frac{{n} \left(u^{r}\right)^2+{n}+{\nu^{r}} {u^{r}}}{{\tau} {u^{\tau}}} \\
S_{6} &= \frac{1}{{r} {T} {\tau} {u^{\tau}}}\Bigl({\nu^{r}} {r} {T} {\tau} {u^{\tau}} - {D_{s}} \left(u^{\tau}\right)^2 (2 {\pi_{Q}{}^{\eta}{}_{\eta}} {r} {u^{r}}+{\pi_{Q}{}^{\eta}{}_{\eta}} {\tau} {u^{\tau}}+{\pi_{Q}{}^{\phi}{}_{\phi}} {r} {u^{r}}+2 {\pi_{Q}{}^{\phi}{}_{\phi}} {\tau} {u^{\tau}})\\
  &\quad + {D_{s}} {h} {\nu^{r}} \left({r} \left(u^{r}\right)^2+{r}+{\tau} {u^{r}} {u^{\tau}}\right)\Bigr) \\
S_{7} &= {\pi_{Q}{}^{\phi}{}_{\phi}}\\
  &\quad + \frac{1}{6 {r} {T} {\tau} {u^{\tau}}}{D_{s}}\Bigl(- 2 {\mathcal{B}} {\nu^{r}} ({r} {u^{r}}-2 {\tau} {u^{\tau}}) - 2 {I_{31}} \left({r} \left(u^{r}\right)^2+{r}-2 {\tau} {u^{r}} {u^{\tau}}\right)\\
  &\quad + {\mathcal{C}}\Bigl(7 {\pi_{Q}{}^{\phi}{}_{\phi}} {\tau} {u^{r}} {u^{\tau}} - 2 {\pi_{Q}{}^{\eta}{}_{\eta}} {r} \left(u^{\tau}\right)^2 + 3 {\pi_{Q}{}^{\phi}{}_{\phi}} {r} \left(u^{\tau}\right)^2\\
  &\quad - 2 {\Pi_{Q}} \left({r} \left(u^{r}\right)^2+{r}-2 {\tau} {u^{r}} {u^{\tau}}\right)\Bigr)\Bigr) \\
S_{8} &= {\pi_{Q}{}^{\eta}{}_{\eta}}\\
  &\quad + \frac{1}{6 {r} {T} {\tau} {u^{\tau}}}{D_{s}}\Bigl(- 2 {I_{31}} {\tau} {u^{r}} {u^{\tau}} + 4 {I_{31}} {r} \left(u^{\tau}\right)^2 + 2 {\mathcal{B}} {\nu^{r}} (2 {r} {u^{r}}-{\tau} {u^{\tau}})\\
  &\quad + {\mathcal{C}}\Bigl(- 2 {\Pi_{Q}} {\tau} {u^{r}} {u^{\tau}} + 3 {\pi_{Q}{}^{\eta}{}_{\eta}} {\tau} {u^{r}} {u^{\tau}} - 2 {\pi_{Q}{}^{\phi}{}_{\phi}} {\tau} {u^{r}} {u^{\tau}} + 4 {\Pi_{Q}} {r} \left(u^{\tau}\right)^2\\
  &\quad + 7 {\pi_{Q}{}^{\eta}{}_{\eta}} {r} \left(u^{\tau}\right)^2\Bigr)\Bigr) \\
S_{9} &= {\Pi_{Q}}\\
  &\quad + \frac{1}{6 {r} {T} {\tau} {u^{\tau}}}{D_{s}}\Bigl(5 {\mathcal{B}} {\nu^{r}} ({r} {u^{r}}+{\tau} {u^{\tau}}) + 5 {I_{31}} \left({r} \left(u^{r}\right)^2+{r}+{\tau} {u^{r}} {u^{\tau}}\right)\\
  &\quad + {\mathcal{C}} \left(5 {\Pi_{Q}} \left({r} \left(u^{r}\right)^2+{r}+{\tau} {u^{r}} {u^{\tau}}\right)+2 {\pi_{Q}{}^{\eta}{}_{\eta}} {r} \left(u^{\tau}\right)^2+2 {\pi_{Q}{}^{\phi}{}_{\phi}} {\tau} {u^{r}} {u^{\tau}}\right)\Bigr)
\end{align*}
\subsection{2+1D Cartesian}

\subsubsection{S3, charm first moment on a background}

\texttt{two\_\allowbreak{}d\_\allowbreak{}HQ\_\allowbreak{}BG\_\allowbreak{}consistent\_\allowbreak{}matrix} in \texttt{src/Matrix/2d\_\allowbreak{}viscous\_\allowbreak{}HQ\_\allowbreak{}BG.jl}, fields:

\[
\phi = ({\alpha}, {\nu^{x}}, {\nu^{y}})
\]
\[
\arraycolsep=2.5pt
A_t = \begin{pmatrix}\bullet & \bullet & \bullet \\ \bullet & \bullet & \cdot \\ \bullet & \cdot & \bullet\end{pmatrix},\qquad A_x = \begin{pmatrix}\bullet & \bullet & \cdot \\ \bullet & \bullet & \cdot \\ \bullet & \cdot & \bullet\end{pmatrix},\qquad A_y = \begin{pmatrix}\bullet & \cdot & \bullet \\ \bullet & \bullet & \cdot \\ \bullet & \cdot & \bullet\end{pmatrix},\qquad S = \begin{pmatrix}\bullet \\ \bullet \\ \bullet\end{pmatrix}
\]
\subsubsection{S4, charm first and second moment on a background}

\texttt{two\_\allowbreak{}d\_\allowbreak{}HQ\_\allowbreak{}BG\_\allowbreak{}m2\_\allowbreak{}matrix} in \texttt{src/Matrix/2d\_\allowbreak{}viscous\_\allowbreak{}HQ\_\allowbreak{}BG\_\allowbreak{}m2.jl}, fields:

\[
\phi = ({\alpha}, {\nu^{x}}, {\nu^{y}}, {\pi_{Q}^{yy}}, {\pi_{Q}{}^{\eta}{}_{\eta}}, {\pi_{Q}^{xy}}, {\Pi_{Q}})
\]
\[
\arraycolsep=2.5pt
A_t = \begin{pmatrix}\bullet & \bullet & \bullet & \cdot & \cdot & \cdot & \cdot \\ \bullet & \bullet & \cdot & \cdot & \cdot & \cdot & \cdot \\ \bullet & \cdot & \bullet & \cdot & \cdot & \cdot & \cdot \\ \cdot & \bullet & \bullet & \bullet & \cdot & \cdot & \cdot \\ \cdot & \bullet & \bullet & \cdot & \bullet & \cdot & \cdot \\ \cdot & \bullet & \bullet & \cdot & \cdot & \bullet & \cdot \\ \bullet & \bullet & \bullet & \cdot & \cdot & \cdot & \bullet\end{pmatrix},\qquad A_x = \begin{pmatrix}\bullet & \bullet & \cdot & \cdot & \cdot & \cdot & \cdot \\ \bullet & \bullet & \cdot & \cdot & \cdot & \cdot & \cdot \\ \bullet & \cdot & \bullet & \cdot & \cdot & \cdot & \cdot \\ \cdot & \bullet & \bullet & \bullet & \cdot & \cdot & \cdot \\ \cdot & \bullet & \cdot & \cdot & \bullet & \cdot & \cdot \\ \cdot & \bullet & \bullet & \cdot & \cdot & \bullet & \cdot \\ \bullet & \bullet & \cdot & \cdot & \cdot & \cdot & \bullet\end{pmatrix},\qquad A_y = \begin{pmatrix}\bullet & \cdot & \bullet & \cdot & \cdot & \cdot & \cdot \\ \bullet & \bullet & \cdot & \cdot & \cdot & \cdot & \cdot \\ \bullet & \cdot & \bullet & \cdot & \cdot & \cdot & \cdot \\ \cdot & \cdot & \bullet & \bullet & \cdot & \cdot & \cdot \\ \cdot & \cdot & \bullet & \cdot & \bullet & \cdot & \cdot \\ \cdot & \bullet & \bullet & \cdot & \cdot & \bullet & \cdot \\ \bullet & \cdot & \bullet & \cdot & \cdot & \cdot & \bullet\end{pmatrix}
\]
\[
\arraycolsep=2.5pt
S = \begin{pmatrix}\bullet \\ \bullet \\ \bullet \\ \bullet \\ \bullet \\ \bullet \\ \bullet\end{pmatrix}
\]
\subsubsection{S5, two-way}

\texttt{two\_\allowbreak{}d\_\allowbreak{}two\_\allowbreak{}way\_\allowbreak{}matrix} in \texttt{src/Matrix/2d\_\allowbreak{}two\_\allowbreak{}way.jl}, fields:

\[
\phi = ({T}, {u^{x}}, {u^{y}}, {\pi^{yy}}, {\pi^{\eta}{}_{\eta}}, {\pi^{xy}}, {\Pi}, {\alpha}, {\nu^{x}}, {\nu^{y}}, {\pi_{Q}^{yy}}, {\pi_{Q}{}^{\eta}{}_{\eta}}, {\pi_{Q}^{xy}}, {\Pi_{Q}})
\]
\[
\arraycolsep=2.5pt
A_t = \begin{pmatrix}\bullet & \bullet & \bullet & \bullet & \bullet & \bullet & \bullet & \bullet & \bullet & \bullet & \cdot & \cdot & \cdot & \cdot \\ \bullet & \bullet & \bullet & \bullet & \bullet & \bullet & \bullet & \bullet & \bullet & \bullet & \cdot & \cdot & \cdot & \cdot \\ \bullet & \bullet & \bullet & \bullet & \cdot & \bullet & \bullet & \bullet & \bullet & \bullet & \cdot & \cdot & \cdot & \cdot \\ \cdot & \bullet & \bullet & \bullet & \cdot & \cdot & \cdot & \cdot & \cdot & \cdot & \cdot & \cdot & \cdot & \cdot \\ \cdot & \bullet & \bullet & \cdot & \bullet & \cdot & \cdot & \cdot & \cdot & \cdot & \cdot & \cdot & \cdot & \cdot \\ \cdot & \bullet & \bullet & \cdot & \cdot & \bullet & \cdot & \cdot & \cdot & \cdot & \cdot & \cdot & \cdot & \cdot \\ \cdot & \bullet & \bullet & \cdot & \cdot & \cdot & \bullet & \cdot & \cdot & \cdot & \cdot & \cdot & \cdot & \cdot \\ \bullet & \bullet & \bullet & \cdot & \cdot & \cdot & \cdot & \bullet & \bullet & \bullet & \cdot & \cdot & \cdot & \cdot \\ \bullet & \bullet & \bullet & \cdot & \cdot & \cdot & \cdot & \bullet & \bullet & \cdot & \cdot & \cdot & \cdot & \cdot \\ \bullet & \bullet & \bullet & \cdot & \cdot & \cdot & \cdot & \bullet & \cdot & \bullet & \cdot & \cdot & \cdot & \cdot \\ \bullet & \bullet & \bullet & \cdot & \cdot & \cdot & \cdot & \cdot & \bullet & \bullet & \bullet & \cdot & \cdot & \cdot \\ \bullet & \bullet & \bullet & \cdot & \cdot & \cdot & \cdot & \cdot & \bullet & \bullet & \cdot & \bullet & \cdot & \cdot \\ \bullet & \bullet & \bullet & \cdot & \cdot & \cdot & \cdot & \cdot & \bullet & \bullet & \cdot & \cdot & \bullet & \cdot \\ \bullet & \bullet & \bullet & \cdot & \cdot & \cdot & \cdot & \bullet & \bullet & \bullet & \cdot & \cdot & \cdot & \bullet\end{pmatrix},\qquad A_x = \begin{pmatrix}\bullet & \bullet & \bullet & \bullet & \bullet & \bullet & \bullet & \bullet & \bullet & \bullet & \cdot & \cdot & \cdot & \cdot \\ \bullet & \bullet & \bullet & \bullet & \bullet & \bullet & \bullet & \bullet & \bullet & \cdot & \cdot & \cdot & \cdot & \cdot \\ \bullet & \bullet & \bullet & \cdot & \cdot & \bullet & \bullet & \bullet & \bullet & \bullet & \cdot & \cdot & \cdot & \cdot \\ \cdot & \bullet & \bullet & \bullet & \cdot & \cdot & \cdot & \cdot & \cdot & \cdot & \cdot & \cdot & \cdot & \cdot \\ \cdot & \bullet & \cdot & \cdot & \bullet & \cdot & \cdot & \cdot & \cdot & \cdot & \cdot & \cdot & \cdot & \cdot \\ \cdot & \bullet & \bullet & \cdot & \cdot & \bullet & \cdot & \cdot & \cdot & \cdot & \cdot & \cdot & \cdot & \cdot \\ \cdot & \bullet & \cdot & \cdot & \cdot & \cdot & \bullet & \cdot & \cdot & \cdot & \cdot & \cdot & \cdot & \cdot \\ \bullet & \bullet & \cdot & \cdot & \cdot & \cdot & \cdot & \bullet & \bullet & \cdot & \cdot & \cdot & \cdot & \cdot \\ \bullet & \bullet & \bullet & \cdot & \cdot & \cdot & \cdot & \bullet & \bullet & \cdot & \cdot & \cdot & \cdot & \cdot \\ \bullet & \bullet & \bullet & \cdot & \cdot & \cdot & \cdot & \bullet & \cdot & \bullet & \cdot & \cdot & \cdot & \cdot \\ \bullet & \bullet & \bullet & \cdot & \cdot & \cdot & \cdot & \cdot & \bullet & \bullet & \bullet & \cdot & \cdot & \cdot \\ \bullet & \bullet & \bullet & \cdot & \cdot & \cdot & \cdot & \cdot & \bullet & \cdot & \cdot & \bullet & \cdot & \cdot \\ \bullet & \bullet & \bullet & \cdot & \cdot & \cdot & \cdot & \cdot & \bullet & \bullet & \cdot & \cdot & \bullet & \cdot \\ \bullet & \bullet & \bullet & \cdot & \cdot & \cdot & \cdot & \bullet & \bullet & \cdot & \cdot & \cdot & \cdot & \bullet\end{pmatrix}
\]
\[
\arraycolsep=2.5pt
A_y = \begin{pmatrix}\bullet & \bullet & \bullet & \bullet & \cdot & \bullet & \bullet & \bullet & \bullet & \bullet & \cdot & \cdot & \cdot & \cdot \\ \bullet & \bullet & \bullet & \cdot & \cdot & \bullet & \bullet & \bullet & \bullet & \bullet & \cdot & \cdot & \cdot & \cdot \\ \bullet & \cdot & \bullet & \bullet & \cdot & \cdot & \bullet & \bullet & \cdot & \bullet & \cdot & \cdot & \cdot & \cdot \\ \cdot & \bullet & \bullet & \bullet & \cdot & \cdot & \cdot & \cdot & \cdot & \cdot & \cdot & \cdot & \cdot & \cdot \\ \cdot & \cdot & \bullet & \cdot & \bullet & \cdot & \cdot & \cdot & \cdot & \cdot & \cdot & \cdot & \cdot & \cdot \\ \cdot & \bullet & \bullet & \cdot & \cdot & \bullet & \cdot & \cdot & \cdot & \cdot & \cdot & \cdot & \cdot & \cdot \\ \cdot & \cdot & \bullet & \cdot & \cdot & \cdot & \bullet & \cdot & \cdot & \cdot & \cdot & \cdot & \cdot & \cdot \\ \bullet & \cdot & \bullet & \cdot & \cdot & \cdot & \cdot & \bullet & \cdot & \bullet & \cdot & \cdot & \cdot & \cdot \\ \bullet & \bullet & \bullet & \cdot & \cdot & \cdot & \cdot & \bullet & \bullet & \cdot & \cdot & \cdot & \cdot & \cdot \\ \bullet & \bullet & \bullet & \cdot & \cdot & \cdot & \cdot & \bullet & \cdot & \bullet & \cdot & \cdot & \cdot & \cdot \\ \bullet & \bullet & \bullet & \cdot & \cdot & \cdot & \cdot & \cdot & \cdot & \bullet & \bullet & \cdot & \cdot & \cdot \\ \bullet & \bullet & \bullet & \cdot & \cdot & \cdot & \cdot & \cdot & \cdot & \bullet & \cdot & \bullet & \cdot & \cdot \\ \bullet & \bullet & \bullet & \cdot & \cdot & \cdot & \cdot & \cdot & \bullet & \bullet & \cdot & \cdot & \bullet & \cdot \\ \bullet & \bullet & \bullet & \cdot & \cdot & \cdot & \cdot & \bullet & \cdot & \bullet & \cdot & \cdot & \cdot & \bullet\end{pmatrix},\qquad S = \begin{pmatrix}\bullet \\ \bullet \\ \bullet \\ \bullet \\ \bullet \\ \bullet \\ \bullet \\ \bullet \\ \bullet \\ \bullet \\ \bullet \\ \bullet \\ \bullet \\ \bullet\end{pmatrix}
\]
\subsubsection{S5, two-way, conformal}

\texttt{two\_\allowbreak{}d\_\allowbreak{}two\_\allowbreak{}way\_\allowbreak{}matrix\_\allowbreak{}conformal} in \texttt{src/Matrix/2d\_\allowbreak{}two\_\allowbreak{}way.jl}, fields:

\[
\phi = ({T}, {u^{x}}, {u^{y}}, {\pi^{yy}}, {\pi^{\eta}{}_{\eta}}, {\pi^{xy}}, {\alpha}, {\nu^{x}}, {\nu^{y}}, {\pi_{Q}^{yy}}, {\pi_{Q}{}^{\eta}{}_{\eta}}, {\pi_{Q}^{xy}}, {\Pi_{Q}})
\]
\[
\arraycolsep=2.5pt
A_t = \begin{pmatrix}\bullet & \bullet & \bullet & \bullet & \bullet & \bullet & \bullet & \bullet & \bullet & \cdot & \cdot & \cdot & \cdot \\ \bullet & \bullet & \bullet & \bullet & \bullet & \bullet & \bullet & \bullet & \bullet & \cdot & \cdot & \cdot & \cdot \\ \bullet & \bullet & \bullet & \bullet & \cdot & \bullet & \bullet & \bullet & \bullet & \cdot & \cdot & \cdot & \cdot \\ \cdot & \bullet & \bullet & \bullet & \cdot & \cdot & \cdot & \cdot & \cdot & \cdot & \cdot & \cdot & \cdot \\ \cdot & \bullet & \bullet & \cdot & \bullet & \cdot & \cdot & \cdot & \cdot & \cdot & \cdot & \cdot & \cdot \\ \cdot & \bullet & \bullet & \cdot & \cdot & \bullet & \cdot & \cdot & \cdot & \cdot & \cdot & \cdot & \cdot \\ \bullet & \bullet & \bullet & \cdot & \cdot & \cdot & \bullet & \bullet & \bullet & \cdot & \cdot & \cdot & \cdot \\ \bullet & \bullet & \bullet & \cdot & \cdot & \cdot & \bullet & \bullet & \cdot & \cdot & \cdot & \cdot & \cdot \\ \bullet & \bullet & \bullet & \cdot & \cdot & \cdot & \bullet & \cdot & \bullet & \cdot & \cdot & \cdot & \cdot \\ \bullet & \bullet & \bullet & \cdot & \cdot & \cdot & \cdot & \bullet & \bullet & \bullet & \cdot & \cdot & \cdot \\ \bullet & \bullet & \bullet & \cdot & \cdot & \cdot & \cdot & \bullet & \bullet & \cdot & \bullet & \cdot & \cdot \\ \bullet & \bullet & \bullet & \cdot & \cdot & \cdot & \cdot & \bullet & \bullet & \cdot & \cdot & \bullet & \cdot \\ \bullet & \bullet & \bullet & \cdot & \cdot & \cdot & \bullet & \bullet & \bullet & \cdot & \cdot & \cdot & \bullet\end{pmatrix},\qquad A_x = \begin{pmatrix}\bullet & \bullet & \bullet & \bullet & \bullet & \bullet & \bullet & \bullet & \bullet & \cdot & \cdot & \cdot & \cdot \\ \bullet & \bullet & \bullet & \bullet & \bullet & \bullet & \bullet & \bullet & \cdot & \cdot & \cdot & \cdot & \cdot \\ \bullet & \bullet & \bullet & \cdot & \cdot & \bullet & \bullet & \bullet & \bullet & \cdot & \cdot & \cdot & \cdot \\ \cdot & \bullet & \bullet & \bullet & \cdot & \cdot & \cdot & \cdot & \cdot & \cdot & \cdot & \cdot & \cdot \\ \cdot & \bullet & \cdot & \cdot & \bullet & \cdot & \cdot & \cdot & \cdot & \cdot & \cdot & \cdot & \cdot \\ \cdot & \bullet & \bullet & \cdot & \cdot & \bullet & \cdot & \cdot & \cdot & \cdot & \cdot & \cdot & \cdot \\ \bullet & \bullet & \cdot & \cdot & \cdot & \cdot & \bullet & \bullet & \cdot & \cdot & \cdot & \cdot & \cdot \\ \bullet & \bullet & \bullet & \cdot & \cdot & \cdot & \bullet & \bullet & \cdot & \cdot & \cdot & \cdot & \cdot \\ \bullet & \bullet & \bullet & \cdot & \cdot & \cdot & \bullet & \cdot & \bullet & \cdot & \cdot & \cdot & \cdot \\ \bullet & \bullet & \bullet & \cdot & \cdot & \cdot & \cdot & \bullet & \bullet & \bullet & \cdot & \cdot & \cdot \\ \bullet & \bullet & \bullet & \cdot & \cdot & \cdot & \cdot & \bullet & \cdot & \cdot & \bullet & \cdot & \cdot \\ \bullet & \bullet & \bullet & \cdot & \cdot & \cdot & \cdot & \bullet & \bullet & \cdot & \cdot & \bullet & \cdot \\ \bullet & \bullet & \bullet & \cdot & \cdot & \cdot & \bullet & \bullet & \cdot & \cdot & \cdot & \cdot & \bullet\end{pmatrix}
\]
\[
\arraycolsep=2.5pt
A_y = \begin{pmatrix}\bullet & \bullet & \bullet & \bullet & \cdot & \bullet & \bullet & \bullet & \bullet & \cdot & \cdot & \cdot & \cdot \\ \bullet & \bullet & \bullet & \cdot & \cdot & \bullet & \bullet & \bullet & \bullet & \cdot & \cdot & \cdot & \cdot \\ \bullet & \cdot & \bullet & \bullet & \cdot & \cdot & \bullet & \cdot & \bullet & \cdot & \cdot & \cdot & \cdot \\ \cdot & \bullet & \bullet & \bullet & \cdot & \cdot & \cdot & \cdot & \cdot & \cdot & \cdot & \cdot & \cdot \\ \cdot & \cdot & \bullet & \cdot & \bullet & \cdot & \cdot & \cdot & \cdot & \cdot & \cdot & \cdot & \cdot \\ \cdot & \bullet & \bullet & \cdot & \cdot & \bullet & \cdot & \cdot & \cdot & \cdot & \cdot & \cdot & \cdot \\ \bullet & \cdot & \bullet & \cdot & \cdot & \cdot & \bullet & \cdot & \bullet & \cdot & \cdot & \cdot & \cdot \\ \bullet & \bullet & \bullet & \cdot & \cdot & \cdot & \bullet & \bullet & \cdot & \cdot & \cdot & \cdot & \cdot \\ \bullet & \bullet & \bullet & \cdot & \cdot & \cdot & \bullet & \cdot & \bullet & \cdot & \cdot & \cdot & \cdot \\ \bullet & \bullet & \bullet & \cdot & \cdot & \cdot & \cdot & \cdot & \bullet & \bullet & \cdot & \cdot & \cdot \\ \bullet & \bullet & \bullet & \cdot & \cdot & \cdot & \cdot & \cdot & \bullet & \cdot & \bullet & \cdot & \cdot \\ \bullet & \bullet & \bullet & \cdot & \cdot & \cdot & \cdot & \bullet & \bullet & \cdot & \cdot & \bullet & \cdot \\ \bullet & \bullet & \bullet & \cdot & \cdot & \cdot & \bullet & \cdot & \bullet & \cdot & \cdot & \cdot & \bullet\end{pmatrix},\qquad S = \begin{pmatrix}\bullet \\ \bullet \\ \bullet \\ \bullet \\ \bullet \\ \bullet \\ \bullet \\ \bullet \\ \bullet \\ \bullet \\ \bullet \\ \bullet \\ \bullet\end{pmatrix}
\]
\subsubsection{S6, fully coupled}

\texttt{two\_\allowbreak{}d\_\allowbreak{}two\_\allowbreak{}way\_\allowbreak{}cM\_\allowbreak{}matrix} in \texttt{src/Matrix/2d\_\allowbreak{}two\_\allowbreak{}way\_\allowbreak{}cM.jl}, fields:

\[
\phi = ({T}, {u^{x}}, {u^{y}}, {\pi^{yy}}, {\pi^{\eta}{}_{\eta}}, {\pi^{xy}}, {\Pi}, {\alpha}, {\nu^{x}}, {\nu^{y}}, {\pi_{Q}^{yy}}, {\pi_{Q}{}^{\eta}{}_{\eta}}, {\pi_{Q}^{xy}}, {\Pi_{Q}})
\]
\[
\arraycolsep=2.5pt
A_t = \begin{pmatrix}\bullet & \bullet & \bullet & \bullet & \bullet & \bullet & \bullet & \bullet & \bullet & \bullet & \bullet & \bullet & \bullet & \bullet \\ \bullet & \bullet & \bullet & \bullet & \bullet & \bullet & \bullet & \bullet & \bullet & \bullet & \bullet & \bullet & \bullet & \bullet \\ \bullet & \bullet & \bullet & \bullet & \cdot & \bullet & \bullet & \bullet & \bullet & \bullet & \bullet & \cdot & \bullet & \bullet \\ \cdot & \bullet & \bullet & \bullet & \cdot & \cdot & \cdot & \cdot & \cdot & \cdot & \cdot & \cdot & \cdot & \cdot \\ \cdot & \bullet & \bullet & \cdot & \bullet & \cdot & \cdot & \cdot & \cdot & \cdot & \cdot & \cdot & \cdot & \cdot \\ \cdot & \bullet & \bullet & \cdot & \cdot & \bullet & \cdot & \cdot & \cdot & \cdot & \cdot & \cdot & \cdot & \cdot \\ \cdot & \bullet & \bullet & \cdot & \cdot & \cdot & \bullet & \cdot & \cdot & \cdot & \cdot & \cdot & \cdot & \cdot \\ \bullet & \bullet & \bullet & \cdot & \cdot & \cdot & \cdot & \bullet & \bullet & \bullet & \cdot & \cdot & \cdot & \cdot \\ \bullet & \bullet & \bullet & \cdot & \cdot & \cdot & \cdot & \bullet & \bullet & \cdot & \bullet & \bullet & \bullet & \bullet \\ \bullet & \bullet & \bullet & \cdot & \cdot & \cdot & \cdot & \bullet & \cdot & \bullet & \bullet & \cdot & \bullet & \bullet \\ \bullet & \bullet & \bullet & \cdot & \cdot & \cdot & \cdot & \cdot & \bullet & \bullet & \bullet & \cdot & \cdot & \cdot \\ \bullet & \bullet & \bullet & \cdot & \cdot & \cdot & \cdot & \cdot & \bullet & \bullet & \cdot & \bullet & \cdot & \cdot \\ \bullet & \bullet & \bullet & \cdot & \cdot & \cdot & \cdot & \cdot & \bullet & \bullet & \cdot & \cdot & \bullet & \cdot \\ \bullet & \bullet & \bullet & \cdot & \cdot & \cdot & \cdot & \bullet & \bullet & \bullet & \cdot & \cdot & \cdot & \bullet\end{pmatrix},\qquad A_x = \begin{pmatrix}\bullet & \bullet & \bullet & \bullet & \bullet & \bullet & \bullet & \bullet & \bullet & \bullet & \bullet & \bullet & \bullet & \bullet \\ \bullet & \bullet & \bullet & \bullet & \bullet & \bullet & \bullet & \bullet & \bullet & \cdot & \bullet & \bullet & \bullet & \bullet \\ \bullet & \bullet & \bullet & \cdot & \cdot & \bullet & \bullet & \bullet & \bullet & \bullet & \cdot & \cdot & \bullet & \bullet \\ \cdot & \bullet & \bullet & \bullet & \cdot & \cdot & \cdot & \cdot & \cdot & \cdot & \cdot & \cdot & \cdot & \cdot \\ \cdot & \bullet & \cdot & \cdot & \bullet & \cdot & \cdot & \cdot & \cdot & \cdot & \cdot & \cdot & \cdot & \cdot \\ \cdot & \bullet & \bullet & \cdot & \cdot & \bullet & \cdot & \cdot & \cdot & \cdot & \cdot & \cdot & \cdot & \cdot \\ \cdot & \bullet & \cdot & \cdot & \cdot & \cdot & \bullet & \cdot & \cdot & \cdot & \cdot & \cdot & \cdot & \cdot \\ \bullet & \bullet & \cdot & \cdot & \cdot & \cdot & \cdot & \bullet & \bullet & \cdot & \cdot & \cdot & \cdot & \cdot \\ \bullet & \bullet & \bullet & \cdot & \cdot & \cdot & \cdot & \bullet & \bullet & \cdot & \bullet & \bullet & \bullet & \bullet \\ \bullet & \bullet & \bullet & \cdot & \cdot & \cdot & \cdot & \bullet & \cdot & \bullet & \cdot & \cdot & \bullet & \bullet \\ \bullet & \bullet & \bullet & \cdot & \cdot & \cdot & \cdot & \cdot & \bullet & \bullet & \bullet & \cdot & \cdot & \cdot \\ \bullet & \bullet & \bullet & \cdot & \cdot & \cdot & \cdot & \cdot & \bullet & \cdot & \cdot & \bullet & \cdot & \cdot \\ \bullet & \bullet & \bullet & \cdot & \cdot & \cdot & \cdot & \cdot & \bullet & \bullet & \cdot & \cdot & \bullet & \cdot \\ \bullet & \bullet & \bullet & \cdot & \cdot & \cdot & \cdot & \bullet & \bullet & \cdot & \cdot & \cdot & \cdot & \bullet\end{pmatrix}
\]
\[
\arraycolsep=2.5pt
A_y = \begin{pmatrix}\bullet & \bullet & \bullet & \bullet & \cdot & \bullet & \bullet & \bullet & \bullet & \bullet & \bullet & \cdot & \bullet & \bullet \\ \bullet & \bullet & \bullet & \cdot & \cdot & \bullet & \bullet & \bullet & \bullet & \bullet & \cdot & \cdot & \bullet & \bullet \\ \bullet & \cdot & \bullet & \bullet & \cdot & \cdot & \bullet & \bullet & \cdot & \bullet & \bullet & \cdot & \cdot & \bullet \\ \cdot & \bullet & \bullet & \bullet & \cdot & \cdot & \cdot & \cdot & \cdot & \cdot & \cdot & \cdot & \cdot & \cdot \\ \cdot & \cdot & \bullet & \cdot & \bullet & \cdot & \cdot & \cdot & \cdot & \cdot & \cdot & \cdot & \cdot & \cdot \\ \cdot & \bullet & \bullet & \cdot & \cdot & \bullet & \cdot & \cdot & \cdot & \cdot & \cdot & \cdot & \cdot & \cdot \\ \cdot & \cdot & \bullet & \cdot & \cdot & \cdot & \bullet & \cdot & \cdot & \cdot & \cdot & \cdot & \cdot & \cdot \\ \bullet & \cdot & \bullet & \cdot & \cdot & \cdot & \cdot & \bullet & \cdot & \bullet & \cdot & \cdot & \cdot & \cdot \\ \bullet & \bullet & \bullet & \cdot & \cdot & \cdot & \cdot & \bullet & \bullet & \cdot & \cdot & \cdot & \bullet & \bullet \\ \bullet & \bullet & \bullet & \cdot & \cdot & \cdot & \cdot & \bullet & \cdot & \bullet & \bullet & \cdot & \cdot & \bullet \\ \bullet & \bullet & \bullet & \cdot & \cdot & \cdot & \cdot & \cdot & \cdot & \bullet & \bullet & \cdot & \cdot & \cdot \\ \bullet & \bullet & \bullet & \cdot & \cdot & \cdot & \cdot & \cdot & \cdot & \bullet & \cdot & \bullet & \cdot & \cdot \\ \bullet & \bullet & \bullet & \cdot & \cdot & \cdot & \cdot & \cdot & \bullet & \bullet & \cdot & \cdot & \bullet & \cdot \\ \bullet & \bullet & \bullet & \cdot & \cdot & \cdot & \cdot & \bullet & \cdot & \bullet & \cdot & \cdot & \cdot & \bullet\end{pmatrix},\qquad S = \begin{pmatrix}\bullet \\ \bullet \\ \bullet \\ \bullet \\ \bullet \\ \bullet \\ \bullet \\ \bullet \\ \bullet \\ \bullet \\ \bullet \\ \bullet \\ \bullet \\ \bullet\end{pmatrix}
\]
\subsubsection{S6, fully coupled, conformal}

\texttt{two\_\allowbreak{}d\_\allowbreak{}two\_\allowbreak{}way\_\allowbreak{}cM\_\allowbreak{}matrix\_\allowbreak{}conformal} in \texttt{src/Matrix/2d\_\allowbreak{}two\_\allowbreak{}way\_\allowbreak{}cM.jl}, fields:

\[
\phi = ({T}, {u^{x}}, {u^{y}}, {\pi^{yy}}, {\pi^{\eta}{}_{\eta}}, {\pi^{xy}}, {\alpha}, {\nu^{x}}, {\nu^{y}}, {\pi_{Q}^{yy}}, {\pi_{Q}{}^{\eta}{}_{\eta}}, {\pi_{Q}^{xy}}, {\Pi_{Q}})
\]
\[
\arraycolsep=2.5pt
A_t = \begin{pmatrix}\bullet & \bullet & \bullet & \bullet & \bullet & \bullet & \bullet & \bullet & \bullet & \bullet & \bullet & \bullet & \bullet \\ \bullet & \bullet & \bullet & \bullet & \bullet & \bullet & \bullet & \bullet & \bullet & \bullet & \bullet & \bullet & \bullet \\ \bullet & \bullet & \bullet & \bullet & \cdot & \bullet & \bullet & \bullet & \bullet & \bullet & \cdot & \bullet & \bullet \\ \cdot & \bullet & \bullet & \bullet & \cdot & \cdot & \cdot & \cdot & \cdot & \cdot & \cdot & \cdot & \cdot \\ \cdot & \bullet & \bullet & \cdot & \bullet & \cdot & \cdot & \cdot & \cdot & \cdot & \cdot & \cdot & \cdot \\ \cdot & \bullet & \bullet & \cdot & \cdot & \bullet & \cdot & \cdot & \cdot & \cdot & \cdot & \cdot & \cdot \\ \bullet & \bullet & \bullet & \cdot & \cdot & \cdot & \bullet & \bullet & \bullet & \cdot & \cdot & \cdot & \cdot \\ \bullet & \bullet & \bullet & \cdot & \cdot & \cdot & \bullet & \bullet & \cdot & \bullet & \bullet & \bullet & \bullet \\ \bullet & \bullet & \bullet & \cdot & \cdot & \cdot & \bullet & \cdot & \bullet & \bullet & \cdot & \bullet & \bullet \\ \bullet & \bullet & \bullet & \cdot & \cdot & \cdot & \cdot & \bullet & \bullet & \bullet & \cdot & \cdot & \cdot \\ \bullet & \bullet & \bullet & \cdot & \cdot & \cdot & \cdot & \bullet & \bullet & \cdot & \bullet & \cdot & \cdot \\ \bullet & \bullet & \bullet & \cdot & \cdot & \cdot & \cdot & \bullet & \bullet & \cdot & \cdot & \bullet & \cdot \\ \bullet & \bullet & \bullet & \cdot & \cdot & \cdot & \bullet & \bullet & \bullet & \cdot & \cdot & \cdot & \bullet\end{pmatrix},\qquad A_x = \begin{pmatrix}\bullet & \bullet & \bullet & \bullet & \bullet & \bullet & \bullet & \bullet & \bullet & \bullet & \bullet & \bullet & \bullet \\ \bullet & \bullet & \bullet & \bullet & \bullet & \bullet & \bullet & \bullet & \cdot & \bullet & \bullet & \bullet & \bullet \\ \bullet & \bullet & \bullet & \cdot & \cdot & \bullet & \bullet & \bullet & \bullet & \cdot & \cdot & \bullet & \bullet \\ \cdot & \bullet & \bullet & \bullet & \cdot & \cdot & \cdot & \cdot & \cdot & \cdot & \cdot & \cdot & \cdot \\ \cdot & \bullet & \cdot & \cdot & \bullet & \cdot & \cdot & \cdot & \cdot & \cdot & \cdot & \cdot & \cdot \\ \cdot & \bullet & \bullet & \cdot & \cdot & \bullet & \cdot & \cdot & \cdot & \cdot & \cdot & \cdot & \cdot \\ \bullet & \bullet & \cdot & \cdot & \cdot & \cdot & \bullet & \bullet & \cdot & \cdot & \cdot & \cdot & \cdot \\ \bullet & \bullet & \bullet & \cdot & \cdot & \cdot & \bullet & \bullet & \cdot & \bullet & \bullet & \bullet & \bullet \\ \bullet & \bullet & \bullet & \cdot & \cdot & \cdot & \bullet & \cdot & \bullet & \cdot & \cdot & \bullet & \bullet \\ \bullet & \bullet & \bullet & \cdot & \cdot & \cdot & \cdot & \bullet & \bullet & \bullet & \cdot & \cdot & \cdot \\ \bullet & \bullet & \bullet & \cdot & \cdot & \cdot & \cdot & \bullet & \cdot & \cdot & \bullet & \cdot & \cdot \\ \bullet & \bullet & \bullet & \cdot & \cdot & \cdot & \cdot & \bullet & \bullet & \cdot & \cdot & \bullet & \cdot \\ \bullet & \bullet & \bullet & \cdot & \cdot & \cdot & \bullet & \bullet & \cdot & \cdot & \cdot & \cdot & \bullet\end{pmatrix}
\]
\[
\arraycolsep=2.5pt
A_y = \begin{pmatrix}\bullet & \bullet & \bullet & \bullet & \cdot & \bullet & \bullet & \bullet & \bullet & \bullet & \cdot & \bullet & \bullet \\ \bullet & \bullet & \bullet & \cdot & \cdot & \bullet & \bullet & \bullet & \bullet & \cdot & \cdot & \bullet & \bullet \\ \bullet & \cdot & \bullet & \bullet & \cdot & \cdot & \bullet & \cdot & \bullet & \bullet & \cdot & \cdot & \bullet \\ \cdot & \bullet & \bullet & \bullet & \cdot & \cdot & \cdot & \cdot & \cdot & \cdot & \cdot & \cdot & \cdot \\ \cdot & \cdot & \bullet & \cdot & \bullet & \cdot & \cdot & \cdot & \cdot & \cdot & \cdot & \cdot & \cdot \\ \cdot & \bullet & \bullet & \cdot & \cdot & \bullet & \cdot & \cdot & \cdot & \cdot & \cdot & \cdot & \cdot \\ \bullet & \cdot & \bullet & \cdot & \cdot & \cdot & \bullet & \cdot & \bullet & \cdot & \cdot & \cdot & \cdot \\ \bullet & \bullet & \bullet & \cdot & \cdot & \cdot & \bullet & \bullet & \cdot & \cdot & \cdot & \bullet & \bullet \\ \bullet & \bullet & \bullet & \cdot & \cdot & \cdot & \bullet & \cdot & \bullet & \bullet & \cdot & \cdot & \bullet \\ \bullet & \bullet & \bullet & \cdot & \cdot & \cdot & \cdot & \cdot & \bullet & \bullet & \cdot & \cdot & \cdot \\ \bullet & \bullet & \bullet & \cdot & \cdot & \cdot & \cdot & \cdot & \bullet & \cdot & \bullet & \cdot & \cdot \\ \bullet & \bullet & \bullet & \cdot & \cdot & \cdot & \cdot & \bullet & \bullet & \cdot & \cdot & \bullet & \cdot \\ \bullet & \bullet & \bullet & \cdot & \cdot & \cdot & \bullet & \cdot & \bullet & \cdot & \cdot & \cdot & \bullet\end{pmatrix},\qquad S = \begin{pmatrix}\bullet \\ \bullet \\ \bullet \\ \bullet \\ \bullet \\ \bullet \\ \bullet \\ \bullet \\ \bullet \\ \bullet \\ \bullet \\ \bullet \\ \bullet\end{pmatrix}
\]
\subsection{Checks}

{\small
\begin{longtable}{@{}>{\raggedright\arraybackslash}p{0.27\textwidth}>{\raggedright\arraybackslash}p{0.60\textwidth}l@{}}
\toprule
id & what it tests & \\
\midrule
\endhead
\texttt{G0-\allowbreak{}0p1d} & $\nabla_{\lambda}g_{\mu\nu}$ =\allowbreak{} 0 & PASS \\
\texttt{G1-\allowbreak{}0p1d} & $\nabla_{\mu}(X, 0, 0, 0)^{\mu} = \partial_{\tau}X + X/\tau$ (Christoffel and $\sqrt{-g}$ routes) & PASS \\
\texttt{G2-\allowbreak{}0p1d} & $u\cdot u = -1$; $u_{\mu}\pi^{\mu\nu} = 0$,\allowbreak{} $\pi^{\mu}{}_{\mu} = 0$ (medium); $u\cdot\nu = 0$; $u_{\mu}\pi_{Q}^{\mu\nu} = 0$,\allowbreak{} $\pi_{Q}{}^{\mu}{}_{\mu} = 0$; both symmetric & PASS \\
\texttt{G-\allowbreak{}aQ-\allowbreak{}0p1d} & $\mathcal{A}'$ has identically vanishing coefficient in the second-moment rows (symbolic); dropped & PASS \\
\texttt{G3-\allowbreak{}0p1d} & every row (medium,\allowbreak{} charm-in-T,\allowbreak{} first moment,\allowbreak{} second moment,\allowbreak{} $\delta M$ in $T_{Q}$) is quasi-linear & PASS \\
\texttt{S3-\allowbreak{}free-\allowbreak{}0p1d} & S3 rows carry no second-moment field & PASS \\
\texttt{S34-\allowbreak{}nomedpi-\allowbreak{}0p1d} & charm rows never see the medium $\pi$ or $\Pi$ (so S3/\allowbreak{}S4 need only T,\allowbreak{} u as background) & PASS \\
\texttt{Gc-\allowbreak{}extract-\allowbreak{}0p1d} & $A_{t}\partial_{\tau}\phi + A_{x}\partial_{x}\phi + A_{y}\partial_{y}\phi + S$ =\allowbreak{} the raw rows (symbolic) & PASS \\
\texttt{Gc-\allowbreak{}ctrl-\allowbreak{}0p1d} & control: the exact random-point zero test rejects a 1e-9-perturbed identity and a nonzero row & PASS \\
\texttt{Gc-\allowbreak{}total-\allowbreak{}0p1d} & S5 fluid rows =\allowbreak{} $\nabla_{\mu}(T_{m} + T_{Q})^{\mu\nu}$,\allowbreak{} $\sqrt{-g}$ route (symbolic) & PASS \\
\texttt{Gc-\allowbreak{}number-\allowbreak{}0p1d} & charm row =\allowbreak{} $\nabla_{\mu}N^{\mu}$,\allowbreak{} $\sqrt{-g}$ route (symbolic) & PASS \\
\texttt{Gc-\allowbreak{}dM-\allowbreak{}0p1d} & the second moment is passive: $\delta M$ columns of the fluid + first-moment rows vanish identically,\allowbreak{} and their sources carry no $\delta M$ & PASS \\
\texttt{Gc-\allowbreak{}off-\allowbreak{}0p1d} & charm off $(n, \nu, \delta M\to 0)$: S5's medium block and medium sources =\allowbreak{} S2 exactly; the charm rows' medium columns and sources vanish,\allowbreak{} so charm =\allowbreak{} 0 is an invariant subspace on which S5 is S2 & PASS \\
\texttt{G6-\allowbreak{}extract-\allowbreak{}0p1d} & S6: $A_{t}\partial_{\tau}\phi + A_{x}\partial_{x}\phi + A_{y}\partial_{y}\phi + S$ =\allowbreak{} the raw rows (symbolic) & PASS \\
\texttt{G6-\allowbreak{}total-\allowbreak{}0p1d} & S6 fluid rows =\allowbreak{} $\nabla_{\mu}(T_{m} + T_{Q} + \delta M)^{\mu\nu}$,\allowbreak{} $\sqrt{-g}$ route (symbolic) & PASS \\
\texttt{G6-\allowbreak{}off-\allowbreak{}0p1d} & S6 at $\delta M$ =\allowbreak{} 0 =\allowbreak{} S5,\allowbreak{} row by row (symbolic) & PASS \\
\texttt{G6-\allowbreak{}m2-\allowbreak{}0p1d} & the second-moment rows of S6 are those of S5 (symbolic) & PASS \\
\texttt{G6-\allowbreak{}live-\allowbreak{}0p1d} & the $\delta M$ columns of the fluid and current rows are live (in 0+1D only through the source) & PASS \\
\texttt{G0-\allowbreak{}1p1d} & $\nabla_{\lambda}g_{\mu\nu}$ =\allowbreak{} 0 & PASS \\
\texttt{G1-\allowbreak{}1p1d} & $\nabla_{\mu}(X, 0, 0, 0)^{\mu} = \partial_{\tau}X + X/\tau$ (Christoffel and $\sqrt{-g}$ routes) & PASS \\
\texttt{G2-\allowbreak{}1p1d} & $u\cdot u = -1$; $u_{\mu}\pi^{\mu\nu} = 0$,\allowbreak{} $\pi^{\mu}{}_{\mu} = 0$ (medium); $u\cdot\nu = 0$; $u_{\mu}\pi_{Q}^{\mu\nu} = 0$,\allowbreak{} $\pi_{Q}{}^{\mu}{}_{\mu} = 0$; both symmetric & PASS \\
\texttt{G2-\allowbreak{}triad} & the $(\ell, \hat\phi, \hat\eta)$ triad is orthonormal and transverse & PASS \\
\texttt{G-\allowbreak{}aQ-\allowbreak{}1p1d} & $\mathcal{A}'$ has identically vanishing coefficient in the second-moment rows (symbolic); dropped & PASS \\
\texttt{G3-\allowbreak{}1p1d} & every row (medium,\allowbreak{} charm-in-T,\allowbreak{} first moment,\allowbreak{} second moment,\allowbreak{} $\delta M$ in $T_{Q}$) is quasi-linear & PASS \\
\texttt{S3-\allowbreak{}free-\allowbreak{}1p1d} & S3 rows carry no second-moment field & PASS \\
\texttt{S34-\allowbreak{}nomedpi-\allowbreak{}1p1d} & charm rows never see the medium $\pi$ or $\Pi$ (so S3/\allowbreak{}S4 need only T,\allowbreak{} u as background) & PASS \\
\texttt{Gc-\allowbreak{}extract-\allowbreak{}1p1d} & $A_{t}\partial_{\tau}\phi + A_{x}\partial_{x}\phi + A_{y}\partial_{y}\phi + S$ =\allowbreak{} the raw rows (symbolic) & PASS \\
\texttt{Gc-\allowbreak{}ctrl-\allowbreak{}1p1d} & control: the exact random-point zero test rejects a 1e-9-perturbed identity and a nonzero row & PASS \\
\texttt{Gc-\allowbreak{}total-\allowbreak{}1p1d} & S5 fluid rows =\allowbreak{} $\nabla_{\mu}(T_{m} + T_{Q})^{\mu\nu}$,\allowbreak{} $\sqrt{-g}$ route (symbolic) & PASS \\
\texttt{Gc-\allowbreak{}number-\allowbreak{}1p1d} & charm row =\allowbreak{} $\nabla_{\mu}N^{\mu}$,\allowbreak{} $\sqrt{-g}$ route (symbolic) & PASS \\
\texttt{Gc-\allowbreak{}dM-\allowbreak{}1p1d} & the second moment is passive: $\delta M$ columns of the fluid + first-moment rows vanish identically,\allowbreak{} and their sources carry no $\delta M$ & PASS \\
\texttt{Gc-\allowbreak{}off-\allowbreak{}1p1d} & charm off $(n, \nu, \delta M\to 0)$: S5's medium block and medium sources =\allowbreak{} S2 exactly; the charm rows' medium columns and sources vanish,\allowbreak{} so charm =\allowbreak{} 0 is an invariant subspace on which S5 is S2 & PASS \\
\texttt{G6-\allowbreak{}extract-\allowbreak{}1p1d} & S6: $A_{t}\partial_{\tau}\phi + A_{x}\partial_{x}\phi + A_{y}\partial_{y}\phi + S$ =\allowbreak{} the raw rows (symbolic) & PASS \\
\texttt{G6-\allowbreak{}total-\allowbreak{}1p1d} & S6 fluid rows =\allowbreak{} $\nabla_{\mu}(T_{m} + T_{Q} + \delta M)^{\mu\nu}$,\allowbreak{} $\sqrt{-g}$ route (symbolic) & PASS \\
\texttt{G6-\allowbreak{}off-\allowbreak{}1p1d} & S6 at $\delta M$ =\allowbreak{} 0 =\allowbreak{} S5,\allowbreak{} row by row (symbolic) & PASS \\
\texttt{G6-\allowbreak{}m2-\allowbreak{}1p1d} & the second-moment rows of S6 are those of S5 (symbolic) & PASS \\
\texttt{G6-\allowbreak{}live-\allowbreak{}1p1d} & the $\delta M$ columns of the fluid and current rows are live (in 0+1D only through the source) & PASS \\
\texttt{G0-\allowbreak{}2p1d} & $\nabla_{\lambda}g_{\mu\nu}$ =\allowbreak{} 0 & PASS \\
\texttt{G1-\allowbreak{}2p1d} & $\nabla_{\mu}(X, 0, 0, 0)^{\mu} = \partial_{\tau}X + X/\tau$ (Christoffel and $\sqrt{-g}$ routes) & PASS \\
\texttt{G2-\allowbreak{}2p1d} & $u\cdot u = -1$; $u_{\mu}\pi^{\mu\nu} = 0$,\allowbreak{} $\pi^{\mu}{}_{\mu} = 0$ (medium); $u\cdot\nu = 0$; $u_{\mu}\pi_{Q}^{\mu\nu} = 0$,\allowbreak{} $\pi_{Q}{}^{\mu}{}_{\mu} = 0$; both symmetric & PASS \\
\texttt{G-\allowbreak{}aQ-\allowbreak{}2p1d} & $\mathcal{A}'$ has identically vanishing coefficient in the second-moment rows (symbolic); dropped & PASS \\
\texttt{G3-\allowbreak{}2p1d} & every row (medium,\allowbreak{} charm-in-T,\allowbreak{} first moment,\allowbreak{} second moment,\allowbreak{} $\delta M$ in $T_{Q}$) is quasi-linear & PASS \\
\texttt{S3-\allowbreak{}free-\allowbreak{}2p1d} & S3 rows carry no second-moment field & PASS \\
\texttt{S34-\allowbreak{}nomedpi-\allowbreak{}2p1d} & charm rows never see the medium $\pi$ or $\Pi$ (so S3/\allowbreak{}S4 need only T,\allowbreak{} u as background) & PASS \\
\texttt{Gc-\allowbreak{}extract-\allowbreak{}2p1d} & $A_{t}\partial_{\tau}\phi + A_{x}\partial_{x}\phi + A_{y}\partial_{y}\phi + S$ =\allowbreak{} the raw rows (exact at 4 random rational points) & PASS \\
\texttt{Gc-\allowbreak{}ctrl-\allowbreak{}2p1d} & control: the exact random-point zero test rejects a 1e-9-perturbed identity and a nonzero row & PASS \\
\texttt{Gc-\allowbreak{}total-\allowbreak{}2p1d} & S5 fluid rows =\allowbreak{} $\nabla_{\mu}(T_{m} + T_{Q})^{\mu\nu}$,\allowbreak{} $\sqrt{-g}$ route (symbolic) & PASS \\
\texttt{Gc-\allowbreak{}number-\allowbreak{}2p1d} & charm row =\allowbreak{} $\nabla_{\mu}N^{\mu}$,\allowbreak{} $\sqrt{-g}$ route (symbolic) & PASS \\
\texttt{Gc-\allowbreak{}dM-\allowbreak{}2p1d} & the second moment is passive: $\delta M$ columns of the fluid + first-moment rows vanish identically,\allowbreak{} and their sources carry no $\delta M$ & PASS \\
\texttt{Gc-\allowbreak{}off-\allowbreak{}2p1d} & charm off $(n, \nu, \delta M\to 0)$: S5's medium block and medium sources =\allowbreak{} S2 exactly; the charm rows' medium columns and sources vanish,\allowbreak{} so charm =\allowbreak{} 0 is an invariant subspace on which S5 is S2 & PASS \\
\texttt{G6-\allowbreak{}extract-\allowbreak{}2p1d} & S6: $A_{t}\partial_{\tau}\phi + A_{x}\partial_{x}\phi + A_{y}\partial_{y}\phi + S$ =\allowbreak{} the raw rows (exact at 4 random rational points) & PASS \\
\texttt{G6-\allowbreak{}total-\allowbreak{}2p1d} & S6 fluid rows =\allowbreak{} $\nabla_{\mu}(T_{m} + T_{Q} + \delta M)^{\mu\nu}$,\allowbreak{} $\sqrt{-g}$ route (symbolic) & PASS \\
\texttt{G6-\allowbreak{}off-\allowbreak{}2p1d} & S6 at $\delta M$ =\allowbreak{} 0 =\allowbreak{} S5,\allowbreak{} row by row (symbolic) & PASS \\
\texttt{G6-\allowbreak{}m2-\allowbreak{}2p1d} & the second-moment rows of S6 are those of S5 (symbolic) & PASS \\
\texttt{G6-\allowbreak{}live-\allowbreak{}2p1d} & the $\delta M$ columns of the fluid and current rows are live (in 0+1D only through the source) & PASS \\
\texttt{R0-\allowbreak{}jl2wl} & Julia -\textgreater{} WL conversion of the shipped matrices (0) & PASS \\
\texttt{R-\allowbreak{}S2-\allowbreak{}1p1d-\allowbreak{}shipped} & S2 1+1D generic EOS,\allowbreak{} $\delta_{\pi\pi} = \delta_{\Pi\Pi} = 0$ =\allowbreak{} one\_\allowbreak{}d\_\allowbreak{}viscous\_\allowbreak{}matrix (1d\_\allowbreak{}viscous\_\allowbreak{}cilindrical.jl),\allowbreak{} entry by entry (55 entries,\allowbreak{} symbolic) & PASS \\
\texttt{R-\allowbreak{}S2-\allowbreak{}2p1d-\allowbreak{}generator} & S2 2+1D,\allowbreak{} $\delta_{\pi\pi} = \delta_{\Pi\Pi} = 0$ =\allowbreak{} two\_\allowbreak{}d\_\allowbreak{}viscous\_\allowbreak{}matrix\_\allowbreak{}derived (2d\_\allowbreak{}viscous.jl,\allowbreak{} from derive\_\allowbreak{}2p1d\_\allowbreak{}viscous.wls),\allowbreak{} entry by entry (154 entries,\allowbreak{} symbolic) & PASS \\
\texttt{R-\allowbreak{}S2-\allowbreak{}1p1d-\allowbreak{}MIS} & S2 1+1D with $\delta_{\pi\pi} = \tfrac{4}{3}\tau_{\pi}$,\allowbreak{} $\delta_{\Pi\Pi} = 0$ =\allowbreak{} medium block of one\_\allowbreak{}d\_\allowbreak{}viscous\_\allowbreak{}matrix\_\allowbreak{}fugacity (:MIS,\allowbreak{} read from the .jl),\allowbreak{} entry by entry (55 entries,\allowbreak{} symbolic) & PASS \\
\texttt{R-\allowbreak{}S2-\allowbreak{}1p1d-\allowbreak{}MIS-\allowbreak{}ctl} & control: $\delta_{\pi\pi} = 0$ does not reproduce the :MIS medium block (the two 1-D media differ by $\delta_{\pi\pi}$) & PASS \\
\texttt{R-\allowbreak{}S4S3-\allowbreak{}1p1d} & S4 rows 1-2 =\allowbreak{} S3 (so the first-moment block of S4/\allowbreak{}S5 is the consistent closure) & PASS \\
\texttt{X6-\allowbreak{}2to1-\allowbreak{}1} & S6 2+1D axisymmetric on the x axis: tau-momentum (symbolic) & PASS \\
\texttt{X6-\allowbreak{}2to1-\allowbreak{}2} & S6 2+1D axisymmetric on the x axis: x-momentum =\allowbreak{} r-momentum (symbolic) & PASS \\
\texttt{X6-\allowbreak{}2to1-\allowbreak{}3} & S6 2+1D axisymmetric on the x axis: y-momentum =\allowbreak{} 0 (symbolic) & PASS \\
\texttt{X6-\allowbreak{}2to1-\allowbreak{}4} & S6 2+1D axisymmetric on the x axis: $E^{yy}$ =\allowbreak{} r\textasciicircum{}2 $E^{\phi\phi}$ (symbolic) & PASS \\
\texttt{X6-\allowbreak{}2to1-\allowbreak{}5} & S6 2+1D axisymmetric on the x axis: $E^{\eta\eta}$ (symbolic) & PASS \\
\texttt{X6-\allowbreak{}2to1-\allowbreak{}6} & S6 2+1D axisymmetric on the x axis: $E^{xy}$ =\allowbreak{} 0 (symbolic) & PASS \\
\texttt{X6-\allowbreak{}2to1-\allowbreak{}7} & S6 2+1D axisymmetric on the x axis: bulk (symbolic) & PASS \\
\texttt{X6-\allowbreak{}2to1-\allowbreak{}8} & S6 2+1D axisymmetric on the x axis: charm number (symbolic) & PASS \\
\texttt{X6-\allowbreak{}2to1-\allowbreak{}9} & S6 2+1D axisymmetric on the x axis: $\nu^{x}$ row =\allowbreak{} $\nu^{r}$ row (symbolic) & PASS \\
\texttt{X6-\allowbreak{}2to1-\allowbreak{}10} & S6 2+1D axisymmetric on the x axis: $\nu^{y}$ row =\allowbreak{} 0 (symbolic) & PASS \\
\texttt{X6-\allowbreak{}2to1-\allowbreak{}11} & S6 2+1D axisymmetric on the x axis: $\pi_{Q}^{yy}$ row =\allowbreak{} $(\phi\phi)$ row (symbolic) & PASS \\
\texttt{X6-\allowbreak{}2to1-\allowbreak{}12} & S6 2+1D axisymmetric on the x axis: $\pi_{Q}{}^{\eta}{}_{\eta}$ row =\allowbreak{} $(\eta\eta)$ row (symbolic) & PASS \\
\texttt{X6-\allowbreak{}2to1-\allowbreak{}13} & S6 2+1D axisymmetric on the x axis: $\pi_{Q}^{xy}$ row =\allowbreak{} 0 (symbolic) & PASS \\
\texttt{X6-\allowbreak{}2to1-\allowbreak{}14} & S6 2+1D axisymmetric on the x axis: trace row (symbolic) & PASS \\
\texttt{X-\allowbreak{}2to1-\allowbreak{}1} & S5 2+1D axisymmetric on the x axis: tau-momentum (symbolic) & PASS \\
\texttt{X-\allowbreak{}2to1-\allowbreak{}2} & S5 2+1D axisymmetric on the x axis: x-momentum =\allowbreak{} r-momentum (symbolic) & PASS \\
\texttt{X-\allowbreak{}2to1-\allowbreak{}3} & S5 2+1D axisymmetric on the x axis: y-momentum =\allowbreak{} 0 (symbolic) & PASS \\
\texttt{X-\allowbreak{}2to1-\allowbreak{}4} & S5 2+1D axisymmetric on the x axis: $E^{yy}$ =\allowbreak{} r\textasciicircum{}2 $E^{\phi\phi}$ (symbolic) & PASS \\
\texttt{X-\allowbreak{}2to1-\allowbreak{}5} & S5 2+1D axisymmetric on the x axis: $E^{\eta\eta}$ (symbolic) & PASS \\
\texttt{X-\allowbreak{}2to1-\allowbreak{}6} & S5 2+1D axisymmetric on the x axis: $E^{xy}$ =\allowbreak{} 0 (symbolic) & PASS \\
\texttt{X-\allowbreak{}2to1-\allowbreak{}7} & S5 2+1D axisymmetric on the x axis: bulk (symbolic) & PASS \\
\texttt{X-\allowbreak{}2to1-\allowbreak{}8} & S5 2+1D axisymmetric on the x axis: charm number (symbolic) & PASS \\
\texttt{X-\allowbreak{}2to1-\allowbreak{}9} & S5 2+1D axisymmetric on the x axis: $\nu^{x}$ row =\allowbreak{} $\nu^{r}$ row (symbolic) & PASS \\
\texttt{X-\allowbreak{}2to1-\allowbreak{}10} & S5 2+1D axisymmetric on the x axis: $\nu^{y}$ row =\allowbreak{} 0 (symbolic) & PASS \\
\texttt{X-\allowbreak{}2to1-\allowbreak{}11} & S5 2+1D axisymmetric on the x axis: $\pi_{Q}^{yy}$ row =\allowbreak{} $(\phi\phi)$ row (symbolic) & PASS \\
\texttt{X-\allowbreak{}2to1-\allowbreak{}12} & S5 2+1D axisymmetric on the x axis: $\pi_{Q}{}^{\eta}{}_{\eta}$ row =\allowbreak{} $(\eta\eta)$ row (symbolic) & PASS \\
\texttt{X-\allowbreak{}2to1-\allowbreak{}13} & S5 2+1D axisymmetric on the x axis: $\pi_{Q}^{xy}$ row =\allowbreak{} 0 (symbolic) & PASS \\
\texttt{X-\allowbreak{}2to1-\allowbreak{}14} & S5 2+1D axisymmetric on the x axis: trace row (symbolic) & PASS \\
\texttt{X-\allowbreak{}ctrl} & control: the same test rejects a wrong pairing (2+1D tau row vs 1+1D r row) & PASS \\
\texttt{X-\allowbreak{}vort-\allowbreak{}axis} & the corotation coupling W vanishes identically on an axisymmetric state (so G-chart is blind to wv,\allowbreak{} as documented) & PASS \\
\texttt{X-\allowbreak{}1to0-\allowbreak{}1} & S5 1+1D at $u^{r} = 0$,\allowbreak{} no r-dependence: energy (symbolic) & PASS \\
\texttt{X-\allowbreak{}1to0-\allowbreak{}2} & S5 1+1D at $u^{r} = 0$,\allowbreak{} no r-dependence: r-momentum =\allowbreak{} 0 (symbolic) & PASS \\
\texttt{X-\allowbreak{}1to0-\allowbreak{}3} & S5 1+1D at $u^{r} = 0$,\allowbreak{} no r-dependence: $E^{\phi\phi}$ =\allowbreak{} -(tau\textasciicircum{}2/\allowbreak{}2r\textasciicircum{}2) $E^{\eta\eta}$ (symbolic) & PASS \\
\texttt{X-\allowbreak{}1to0-\allowbreak{}4} & S5 1+1D at $u^{r} = 0$,\allowbreak{} no r-dependence: $E^{\eta\eta}$ (symbolic) & PASS \\
\texttt{X-\allowbreak{}1to0-\allowbreak{}5} & S5 1+1D at $u^{r} = 0$,\allowbreak{} no r-dependence: bulk (symbolic) & PASS \\
\texttt{X-\allowbreak{}1to0-\allowbreak{}6} & S5 1+1D at $u^{r} = 0$,\allowbreak{} no r-dependence: charm number (symbolic) & PASS \\
\texttt{X-\allowbreak{}1to0-\allowbreak{}7} & S5 1+1D at $u^{r} = 0$,\allowbreak{} no r-dependence: $\nu^{r}$ row =\allowbreak{} 0 (symbolic) & PASS \\
\texttt{X-\allowbreak{}1to0-\allowbreak{}8} & S5 1+1D at $u^{r} = 0$,\allowbreak{} no r-dependence: $(\phi\phi)$ row =\allowbreak{} -1/\allowbreak{}2 $(\eta\eta)$ row (symbolic) & PASS \\
\texttt{X-\allowbreak{}1to0-\allowbreak{}9} & S5 1+1D at $u^{r} = 0$,\allowbreak{} no r-dependence: $(\eta\eta)$ row (symbolic) & PASS \\
\texttt{X-\allowbreak{}1to0-\allowbreak{}10} & S5 1+1D at $u^{r} = 0$,\allowbreak{} no r-dependence: trace row (symbolic) & PASS \\
\texttt{X6-\allowbreak{}1to0-\allowbreak{}1} & S6 1+1D at $u^{r} = 0$,\allowbreak{} no r-dependence: energy (symbolic) & PASS \\
\texttt{X6-\allowbreak{}1to0-\allowbreak{}2} & S6 1+1D at $u^{r} = 0$,\allowbreak{} no r-dependence: r-momentum =\allowbreak{} 0 (symbolic) & PASS \\
\texttt{X6-\allowbreak{}1to0-\allowbreak{}3} & S6 1+1D at $u^{r} = 0$,\allowbreak{} no r-dependence: $E^{\phi\phi}$ =\allowbreak{} -(tau\textasciicircum{}2/\allowbreak{}2r\textasciicircum{}2) $E^{\eta\eta}$ (symbolic) & PASS \\
\texttt{X6-\allowbreak{}1to0-\allowbreak{}4} & S6 1+1D at $u^{r} = 0$,\allowbreak{} no r-dependence: $E^{\eta\eta}$ (symbolic) & PASS \\
\texttt{X6-\allowbreak{}1to0-\allowbreak{}5} & S6 1+1D at $u^{r} = 0$,\allowbreak{} no r-dependence: bulk (symbolic) & PASS \\
\texttt{X6-\allowbreak{}1to0-\allowbreak{}6} & S6 1+1D at $u^{r} = 0$,\allowbreak{} no r-dependence: charm number (symbolic) & PASS \\
\texttt{X6-\allowbreak{}1to0-\allowbreak{}7} & S6 1+1D at $u^{r} = 0$,\allowbreak{} no r-dependence: $\nu^{r}$ row =\allowbreak{} 0 (symbolic) & PASS \\
\texttt{X6-\allowbreak{}1to0-\allowbreak{}8} & S6 1+1D at $u^{r} = 0$,\allowbreak{} no r-dependence: $(\phi\phi)$ row =\allowbreak{} -1/\allowbreak{}2 $(\eta\eta)$ row (symbolic) & PASS \\
\texttt{X6-\allowbreak{}1to0-\allowbreak{}9} & S6 1+1D at $u^{r} = 0$,\allowbreak{} no r-dependence: $(\eta\eta)$ row (symbolic) & PASS \\
\texttt{X6-\allowbreak{}1to0-\allowbreak{}10} & S6 1+1D at $u^{r} = 0$,\allowbreak{} no r-dependence: trace row (symbolic) & PASS \\
\texttt{B-\allowbreak{}ideal-\allowbreak{}conf} & conformal ideal Bjorken: $\partial_{\tau}T = -T/(3\tau)$ & PASS \\
\texttt{B-\allowbreak{}ideal-\allowbreak{}gen} & generic-EOS ideal Bjorken: $\partial_{\tau}T = -P'/(P''\tau) = -c_{s}^{2}T/\tau,\ c_{s}^{2} = P'/(TP'')$ & PASS \\
\texttt{B-\allowbreak{}IS-\allowbreak{}gen} & generic is Bjorken: $\partial_{\tau}T = -(TP' + \Pi + \pi)/(\tau TP''),\ \tau_{\pi}\partial_{\tau}\pi = -\pi - 4\eta/(3\tau) - \delta_{\pi\pi}\pi/\tau,\ \tau_{\Pi}\partial_{\tau}\Pi = -\Pi - \zeta/\tau - \delta_{\Pi\Pi}\Pi/\tau$ & PASS \\
\texttt{B-\allowbreak{}IS-\allowbreak{}conf} & conformal is Bjorken: $\partial_{\tau}T = -T/(3\tau) - \pi/(12aT^{3}\tau),\ \tau_{\pi}\partial_{\tau}\pi = -\pi - 4\eta/(3\tau) - \delta_{\pi\pi}\pi/\tau$  [$\delta_{\pi\pi} = \tfrac{4}{3}\tau_{\pi}$: the DNMR $\delta_{\pi\pi}$ form; no $\tau_{\pi\pi}, \lambda_{1}$] & PASS \\
\texttt{B-\allowbreak{}NS-\allowbreak{}conf} & first-order (NS) limit $\pi = -4\eta/(3\tau)$: $\partial_{\tau}T = -T/(3\tau) + (4/9)(\eta/s)/\tau^{2},\ s = 4aT^{3}$ & PASS \\
\texttt{B-\allowbreak{}twoway-\allowbreak{}energy} & S5 Bjorken energy row: $\partial_{\tau}(e + e_{Q}) + (e + P + \Pi + \pi + nh)/\tau$ & PASS \\
\texttt{B-\allowbreak{}S6-\allowbreak{}energy} & S6 Bjorken energy row: $\partial_{\tau}(e + e_{Q}) + (e + P + \Pi + \pi + nh + \pi_{Q}{}^{\eta}{}_{\eta} + \Pi_{Q})/\tau$ & PASS \\
\texttt{B-\allowbreak{}charm-\allowbreak{}number} & Bjorken charm number: $\tau^{-1}\partial_{\tau}(n\tau) = 0$ & PASS \\
\texttt{T-\allowbreak{}rest-\allowbreak{}1p1d} & $\nu = 0, \delta M = 0, \nabla\alpha = 0$,\allowbreak{} medium at rest: transverse momentum exchange $\Delta^{\mu}{}_{\nu}\nabla_{\lambda}T_{Q}^{\lambda\nu}$ =\allowbreak{} 0 and the nu row is satisfied & PASS \\
\texttt{T-\allowbreak{}drag-\allowbreak{}1p1d} & drag: on shell,\allowbreak{} $\nabla_{\mu}T_{Q}^{\mu i}$ =\allowbreak{} -$(T/D_{s})\nu^{i}$ (drag on charm) and the fluid row then gives $\nabla_{\mu}T_{m}^{\mu i}$ =\allowbreak{} +$(T/D_{s})\nu^{i}$ (Newton III),\allowbreak{} exact & PASS \\
\texttt{T-\allowbreak{}split-\allowbreak{}1p1d} & fluid momentum rows =\allowbreak{} $\nabla_{\mu}T_{m}^{\mu\nu} + \nabla_{\mu}T_{Q}^{\mu\nu}$ identically & PASS \\
\texttt{T6-\allowbreak{}drag-\allowbreak{}1p1d} & S6 drag: on shell,\allowbreak{} $\nabla_{\mu}(T_{Q} + \delta M)^{\mu i} = -(T/D_{s})\nu^{i}$,\allowbreak{} so the medium gains $(T/D_{s})\nu^{i}$ (Newton III),\allowbreak{} exact & PASS \\
\texttt{T6-\allowbreak{}static} & static $\delta M$ at rest: $\nabla_{\mu}\delta M^{\mu r} = \partial_{r}(\pi_{rr} + \Pi_{Q}) + (\pi_{rr} - \pi^{\phi}{}_{\phi})/r$ & PASS \\
\texttt{T-\allowbreak{}rest-\allowbreak{}2p1d} & $\nu = 0, \delta M = 0, \nabla\alpha = 0$,\allowbreak{} medium at rest: transverse momentum exchange $\Delta^{\mu}{}_{\nu}\nabla_{\lambda}T_{Q}^{\lambda\nu}$ =\allowbreak{} 0 and the nu row is satisfied & PASS \\
\texttt{T-\allowbreak{}drag-\allowbreak{}2p1d} & drag: on shell,\allowbreak{} $\nabla_{\mu}T_{Q}^{\mu i}$ =\allowbreak{} -$(T/D_{s})\nu^{i}$ (drag on charm) and the fluid row then gives $\nabla_{\mu}T_{m}^{\mu i}$ =\allowbreak{} +$(T/D_{s})\nu^{i}$ (Newton III),\allowbreak{} exact & PASS \\
\texttt{T-\allowbreak{}split-\allowbreak{}2p1d} & fluid momentum rows =\allowbreak{} $\nabla_{\mu}T_{m}^{\mu\nu} + \nabla_{\mu}T_{Q}^{\mu\nu}$ identically & PASS \\
\texttt{T6-\allowbreak{}drag-\allowbreak{}2p1d} & S6 drag: on shell,\allowbreak{} $\nabla_{\mu}(T_{Q} + \delta M)^{\mu i} = -(T/D_{s})\nu^{i}$,\allowbreak{} so the medium gains $(T/D_{s})\nu^{i}$ (Newton III),\allowbreak{} exact & PASS \\
\texttt{S6-\allowbreak{}speeds-\allowbreak{}1p1d} & S6 is hyperbolic and subluminal at all 72 (state,\allowbreak{} direction) pairs,\allowbreak{} exactly,\allowbreak{} with max speed 0.9901 & PASS \\
\texttt{S6-\allowbreak{}speeds-\allowbreak{}2p1d} & S6 is hyperbolic and subluminal at all 144 (state,\allowbreak{} direction) pairs,\allowbreak{} exactly,\allowbreak{} with max speed 0.9903 & PASS \\
\texttt{E-\allowbreak{}same-\allowbreak{}one\_\allowbreak{}d\_\allowbreak{}HQ\_\allowbreak{}BG\_\allowbreak{}consistent\_\allowbreak{}matrix} & simplified exported entries =\allowbreak{} derived entries (symbolic) & PASS \\
\texttt{E-\allowbreak{}free-\allowbreak{}one\_\allowbreak{}d\_\allowbreak{}HQ\_\allowbreak{}BG\_\allowbreak{}consistent\_\allowbreak{}matrix} & every free symbol is an argument or a field & PASS \\
\texttt{E-\allowbreak{}same-\allowbreak{}one\_\allowbreak{}d\_\allowbreak{}HQ\_\allowbreak{}BG\_\allowbreak{}m2\_\allowbreak{}matrix} & simplified exported entries =\allowbreak{} derived entries (symbolic) & PASS \\
\texttt{E-\allowbreak{}free-\allowbreak{}one\_\allowbreak{}d\_\allowbreak{}HQ\_\allowbreak{}BG\_\allowbreak{}m2\_\allowbreak{}matrix} & every free symbol is an argument or a field & PASS \\
\texttt{E-\allowbreak{}same-\allowbreak{}one\_\allowbreak{}d\_\allowbreak{}two\_\allowbreak{}way\_\allowbreak{}matrix} & simplified exported entries =\allowbreak{} derived entries (symbolic) & PASS \\
\texttt{E-\allowbreak{}free-\allowbreak{}one\_\allowbreak{}d\_\allowbreak{}two\_\allowbreak{}way\_\allowbreak{}matrix} & every free symbol is an argument or a field & PASS \\
\texttt{E-\allowbreak{}same-\allowbreak{}one\_\allowbreak{}d\_\allowbreak{}two\_\allowbreak{}way\_\allowbreak{}matrix\_\allowbreak{}conformal} & simplified exported entries =\allowbreak{} derived entries (symbolic) & PASS \\
\texttt{E-\allowbreak{}free-\allowbreak{}one\_\allowbreak{}d\_\allowbreak{}two\_\allowbreak{}way\_\allowbreak{}matrix\_\allowbreak{}conformal} & every free symbol is an argument or a field & PASS \\
\texttt{E-\allowbreak{}same-\allowbreak{}two\_\allowbreak{}d\_\allowbreak{}HQ\_\allowbreak{}BG\_\allowbreak{}consistent\_\allowbreak{}matrix} & simplified exported entries =\allowbreak{} derived entries (symbolic) & PASS \\
\texttt{E-\allowbreak{}free-\allowbreak{}two\_\allowbreak{}d\_\allowbreak{}HQ\_\allowbreak{}BG\_\allowbreak{}consistent\_\allowbreak{}matrix} & every free symbol is an argument or a field & PASS \\
\texttt{E-\allowbreak{}same-\allowbreak{}two\_\allowbreak{}d\_\allowbreak{}HQ\_\allowbreak{}BG\_\allowbreak{}m2\_\allowbreak{}matrix} & simplified exported entries =\allowbreak{} derived entries (exact at 3 random rational points) & PASS \\
\texttt{E-\allowbreak{}free-\allowbreak{}two\_\allowbreak{}d\_\allowbreak{}HQ\_\allowbreak{}BG\_\allowbreak{}m2\_\allowbreak{}matrix} & every free symbol is an argument or a field & PASS \\
\texttt{E-\allowbreak{}same-\allowbreak{}two\_\allowbreak{}d\_\allowbreak{}two\_\allowbreak{}way\_\allowbreak{}matrix} & simplified exported entries =\allowbreak{} derived entries (exact at 3 random rational points) & PASS \\
\texttt{E-\allowbreak{}free-\allowbreak{}two\_\allowbreak{}d\_\allowbreak{}two\_\allowbreak{}way\_\allowbreak{}matrix} & every free symbol is an argument or a field & PASS \\
\texttt{E-\allowbreak{}same-\allowbreak{}two\_\allowbreak{}d\_\allowbreak{}two\_\allowbreak{}way\_\allowbreak{}matrix\_\allowbreak{}conformal} & simplified exported entries =\allowbreak{} derived entries (exact at 3 random rational points) & PASS \\
\texttt{E-\allowbreak{}free-\allowbreak{}two\_\allowbreak{}d\_\allowbreak{}two\_\allowbreak{}way\_\allowbreak{}matrix\_\allowbreak{}conformal} & every free symbol is an argument or a field & PASS \\
\texttt{E-\allowbreak{}same-\allowbreak{}one\_\allowbreak{}d\_\allowbreak{}two\_\allowbreak{}way\_\allowbreak{}cM\_\allowbreak{}matrix} & simplified exported entries =\allowbreak{} derived entries (symbolic) & PASS \\
\texttt{E-\allowbreak{}free-\allowbreak{}one\_\allowbreak{}d\_\allowbreak{}two\_\allowbreak{}way\_\allowbreak{}cM\_\allowbreak{}matrix} & every free symbol is an argument or a field & PASS \\
\texttt{E-\allowbreak{}same-\allowbreak{}one\_\allowbreak{}d\_\allowbreak{}two\_\allowbreak{}way\_\allowbreak{}cM\_\allowbreak{}matrix\_\allowbreak{}conformal} & simplified exported entries =\allowbreak{} derived entries (symbolic) & PASS \\
\texttt{E-\allowbreak{}free-\allowbreak{}one\_\allowbreak{}d\_\allowbreak{}two\_\allowbreak{}way\_\allowbreak{}cM\_\allowbreak{}matrix\_\allowbreak{}conformal} & every free symbol is an argument or a field & PASS \\
\texttt{E-\allowbreak{}same-\allowbreak{}two\_\allowbreak{}d\_\allowbreak{}two\_\allowbreak{}way\_\allowbreak{}cM\_\allowbreak{}matrix} & simplified exported entries =\allowbreak{} derived entries (exact at 3 random rational points) & PASS \\
\texttt{E-\allowbreak{}free-\allowbreak{}two\_\allowbreak{}d\_\allowbreak{}two\_\allowbreak{}way\_\allowbreak{}cM\_\allowbreak{}matrix} & every free symbol is an argument or a field & PASS \\
\texttt{E-\allowbreak{}same-\allowbreak{}two\_\allowbreak{}d\_\allowbreak{}two\_\allowbreak{}way\_\allowbreak{}cM\_\allowbreak{}matrix\_\allowbreak{}conformal} & simplified exported entries =\allowbreak{} derived entries (exact at 3 random rational points) & PASS \\
\texttt{E-\allowbreak{}free-\allowbreak{}two\_\allowbreak{}d\_\allowbreak{}two\_\allowbreak{}way\_\allowbreak{}cM\_\allowbreak{}matrix\_\allowbreak{}conformal} & every free symbol is an argument or a field & PASS \\
\texttt{E-\allowbreak{}jl} & the generated Julia (12 functions) parses,\allowbreak{} runs,\allowbreak{} and reproduces its 24-state CSVs: max rel err 1.4298582449629782e-15 & PASS \\
\bottomrule
\end{longtable}
}
155 passed, 0 failed (155 checks evaluated).

Not a check: the characteristic speeds of S5 are computed and reported (\texttt{two\_way\_char\_speeds.tsv} in the scratch directory).

Re-run with \texttt{wolframscript -file deriv/derive\_\allowbreak{}two\_\allowbreak{}way.wls}.


\clearpage
% GENERATED by deriv/render.wl from the derivation script.  Do not edit.

\section{The hand-written 1+1D matrices}

Script: \texttt{deriv/derive\_1d\_matrices.wls}

Writes nothing into \texttt{src/Matrix}: it compares the derivation with the hand-written \texttt{.jl} files, read from the file.

\subsection{Setup}

\[
x^{\mu} = ({\tau}, {r}, {\phi}, {\eta}),\qquad g_{\mu\nu} = \mathrm{diag}(-1, 1, {r}^2, {\tau}^2),\qquad u^{\mu} = (u^{\tau}, u^{r}, 0, 0)
\]
Medium as in the viscous section. Charm:

\[
\nabla_{\mu}(n\,u^{\mu} + \nu^{\mu}) = 0,\qquad \tau_{n}\,\Delta^{\mu}{}_{\beta}D\nu^{\beta} + \nu^{\mu} + \kappa\,\Delta^{\mu\beta}\partial_{\beta}\alpha = 0,\qquad \nu^{\tau} = \frac{u^{r}\nu^{r}}{u^{\tau}}
\]
Rows $\nabla_{\mu}N^{\mu}$ and the $r$ component of the current equation. Variants:

\begin{gather*}
\text{MIS}:\ n = n(T, \alpha),\ \delta_{\pi\pi} = \tfrac{4}{3}\tau_{\pi},\quad \phi = (T, u^{r}, \pi^{\phi}{}_{\phi}, \pi^{\eta}{}_{\eta}, \Pi, \alpha, \nu^{r}) \\
\text{gamma}:\ n = \gamma\,n(T),\ \kappa\partial\alpha\to\kappa\partial\gamma,\ \delta_{\pi\pi} = \tfrac{4}{3}\tau_{\pi},\quad \phi = (T, u^{r}, \pi^{\phi}{}_{\phi}, \pi^{\eta}{}_{\eta}, \Pi, \gamma, \nu^{r}) \\
\text{ideal}:\ \pi = \Pi = 0,\ \nu = 0,\quad \phi = (T, u^{r}, n)\ \text{or}\ (T, u^{r}, \alpha)
\end{gather*}
\subsection{Steps}

\[
\Gamma^{{\tau}}{}_{{\eta}{\eta}} = {\tau},\quad \Gamma^{{r}}{}_{{\phi}{\phi}} = -{r},\quad \Gamma^{{\phi}}{}_{{r}{\phi}} = \frac{1}{{r}},\quad \Gamma^{{\eta}}{}_{{\tau}{\eta}} = \frac{1}{{\tau}}
\]
\begin{align*}
a^{{\tau}} &= {u^{r}} \left({(\partial_{{\tau}} u^{r})}+\frac{{(\partial_{{r}} u^{r})} {u^{r}}}{{u^{\tau}}}\right) \\
a^{{r}} &= {(\partial_{{\tau}} u^{r})} {u^{\tau}}+{(\partial_{{r}} u^{r})} {u^{r}} \\
\nu_{\lambda}a^{\lambda} &= \frac{{\nu^{r}} ({(\partial_{{\tau}} u^{r})} {u^{\tau}}+{(\partial_{{r}} u^{r})} {u^{r}})}{\left(u^{\tau}\right)^2} \\
D\nu^{{\tau}} &= \frac{{u^{\tau}} \left({(\partial_{{\tau}} \nu^{r})} \left(u^{r}\right)^3+{(\partial_{{\tau}} \nu^{r})} {u^{r}}+{(\partial_{{\tau}} u^{r})} {\nu^{r}}\right)+{(\partial_{{r}} \nu^{r})} \left(\left(u^{r}\right)^4+\left(u^{r}\right)^2\right)+{(\partial_{{r}} u^{r})} {\nu^{r}} {u^{r}}}{\left(u^{\tau}\right)^3} \\
D\nu^{{r}} &= {(\partial_{{\tau}} \nu^{r})} {u^{\tau}}+{(\partial_{{r}} \nu^{r})} {u^{r}}
\end{align*}
\[
E_{\nu} = \tau_{n}\bigl(D\nu^{r} - u^{r}\nu_{\lambda}a^{\lambda}\bigr) + \nu^{r} + \kappa\bigl(\Delta^{r\tau}\partial_{\tau}\alpha + \Delta^{rr}\partial_{r}\alpha\bigr),\qquad \Delta^{r\tau} = u^{r}u^{\tau},\ \Delta^{rr} = 1 + (u^{r})^{2}
\]
\begin{align*}
E_{\alpha} &= {(\partial_{{r}} \nu^{r})} + {(\partial_{{r}} u^{r})} {n} + {u^{r}} ({(\partial_{{r}} \alpha)} {(\partial_{\alpha} n)}+{(\partial_{{r}} T)} {(\partial_{T} n)}) + \frac{{n} {u^{r}}+{\nu^{r}}}{{r}} - \frac{{(\partial_{{\tau}} u^{r})} {\nu^{r}} \left(u^{r}\right)^2}{\left(u^{\tau}\right)^3}\\
  &\quad + \frac{{(\partial_{{\tau}} u^{r})} {\nu^{r}}}{{u^{\tau}}} + \frac{{(\partial_{{\tau}} \nu^{r})} {u^{r}}}{{u^{\tau}}} + \frac{{(\partial_{{\tau}} u^{r})} {n} {u^{r}}}{{u^{\tau}}} + \frac{{n} \left(u^{r}\right)^2+{n}+{\nu^{r}} {u^{r}}}{{\tau} {u^{\tau}}}\\
  &\quad + {u^{\tau}} ({(\partial_{{\tau}} \alpha)} {(\partial_{\alpha} n)}+{(\partial_{{\tau}} T)} {(\partial_{T} n)}) \\
E_{\nu} &= {(\partial_{{r}} \nu^{r})} {\tau_{n}} {u^{r}} + {u^{\tau}} ({(\partial_{{\tau}} \alpha)} {\kappa} {u^{r}}+{(\partial_{{\tau}} \nu^{r})} {\tau_{n}}) + {(\partial_{{r}} \alpha)} {\kappa} \left(u^{\tau}\right)^2\\
  &\quad + \frac{{\nu^{r}} \left(-{(\partial_{{\tau}} u^{r})} {\tau_{n}} {u^{r}} {u^{\tau}}+\left(u^{r}\right)^2 (1-{(\partial_{{r}} u^{r})} {\tau_{n}})+1\right)}{\left(u^{\tau}\right)^2}
\end{align*}
\[
A^{(k)}_{ij} = \frac{\partial E_{i}}{\partial(\partial_{k}\phi_{j})},\qquad S_{i} = E_{i}\bigl|_{\partial\phi = 0},\qquad k = \tau, r
\]
\subsection{Result}

\subsubsection{\texttt{one\_d\_viscous\_matrix\_fugacity\_BG\_only} in \texttt{src/Matrix/1d\_viscous\_HQ\_cilindrical\_BG.jl}}

Background $T$, $u^{r}$; $\phi = (\alpha, \nu^{r})$; the file multiplies every row by $(u^{\tau})^{3}r\tau$ (M-BG). As derived:

\begin{align*}
(A_t)_{11} &= {(\partial_{\alpha} n)} {u^{\tau}} \\
(A_t)_{12} &= \frac{{u^{r}}}{{u^{\tau}}} \\
(A_t)_{21} &= {\kappa} {u^{r}} {u^{\tau}} \\
(A_t)_{22} &= {\tau_{n}} {u^{\tau}} \\
(A_r)_{11} &= {(\partial_{\alpha} n)} {u^{r}} \\
(A_r)_{12} &= 1 \\
(A_r)_{21} &= {\kappa} \left(u^{\tau}\right)^2 \\
(A_r)_{22} &= {\tau_{n}} {u^{r}} \\
S_{1} &= \frac{1}{{r} {\tau} \left(u^{\tau}\right)^3}\Bigl({(\partial_{\tau} u^{r})} {\nu^{r}} {r} {\tau} + {\nu^{r}} \left(u^{\tau}\right)^2 ({r} {u^{r}}+{\tau} {u^{\tau}})\\
  &\quad + {(\partial_{T} n)} {r} {\tau} \left(u^{\tau}\right)^2 \left({(\partial_{r} T)} {u^{r}} {u^{\tau}}+{(\partial_{\tau} T)} \left(u^{r}\right)^2+{(\partial_{\tau} T)}\right)\\
  &\quad + {n} \left(u^{\tau}\right)^2 \left({r} \left({(\partial_{r} u^{r})} {\tau} {u^{\tau}}+{(\partial_{\tau} u^{r})} {\tau} {u^{r}}+\left(u^{r}\right)^2+1\right)+{\tau} {u^{r}} {u^{\tau}}\right)\Bigr) \\
S_{2} &= \frac{{\nu^{r}} \left(\left(u^{r}\right)^2 (1-{(\partial_{r} u^{r})} {\tau_{n}})-{(\partial_{\tau} u^{r})} {\tau_{n}} {u^{r}} {u^{\tau}}+1\right)}{\left(u^{\tau}\right)^2}
\end{align*}
\subsubsection{\texttt{one\_\allowbreak{}d\_\allowbreak{}viscous\_\allowbreak{}matrix\_\allowbreak{}fugacity} in \texttt{src/Matrix/1d\_\allowbreak{}viscous\_\allowbreak{}HQ\_\allowbreak{}cilindrical\_\allowbreak{}fugacity.jl}}

Derived, charm rows:

\begin{align*}
(A_t)_{61} &= {(\partial_{T} n)} {u^{\tau}} \\
(A_t)_{62} &= \frac{{n} {u^{r}} \left(u^{\tau}\right)^2+{\nu^{r}}}{\left(u^{\tau}\right)^3} \\
(A_t)_{66} &= {(\partial_{\alpha} n)} {u^{\tau}} \\
(A_t)_{67} &= \frac{{u^{r}}}{{u^{\tau}}} \\
(A_t)_{72} &= -\frac{{\nu^{r}} {\tau_{n}} {u^{r}}}{{u^{\tau}}} \\
(A_t)_{76} &= {\kappa} {u^{r}} {u^{\tau}} \\
(A_t)_{77} &= {\tau_{n}} {u^{\tau}} \\
(A_r)_{61} &= {(\partial_{T} n)} {u^{r}} \\
(A_r)_{62} &= {n} \\
(A_r)_{66} &= {(\partial_{\alpha} n)} {u^{r}} \\
(A_r)_{67} &= 1 \\
(A_r)_{72} &= -\frac{{\nu^{r}} {\tau_{n}} \left(u^{r}\right)^2}{\left(u^{\tau}\right)^2} \\
(A_r)_{76} &= {\kappa} \left(u^{\tau}\right)^2 \\
(A_r)_{77} &= {\tau_{n}} {u^{r}} \\
S_{6} &= \frac{{n} \left({r} \left(u^{r}\right)^2+{r}+{\tau} {u^{r}} {u^{\tau}}\right)+{\nu^{r}} ({r} {u^{r}}+{\tau} {u^{\tau}})}{{r} {\tau} {u^{\tau}}} \\
S_{7} &= {\nu^{r}}
\end{align*}
With $c$ the charm rows and $m$ the medium rows, $\mathrm{row}_{c}\to\mathrm{row}_{c} - L\,\mathrm{row}_{m}$, $L = A_t^{cm}(A_t^{mm})^{-1}$, removes $\partial_{\tau}$ of the medium from the charm rows:

\begin{align*}
\mathcal{D}_{1} &= 3 \left(p'(T)\right)^2 {\tau_{\Pi}} {\tau_{\pi}} \left(u^{r}\right)^2 {T}\\
  &\quad + 3 {p'(T)} {\tau_{\Pi}} {\tau_{\pi}} \left(-{p''(T)} \left(u^{r}\right)^2 {T}^2-{p''(T)} {T}^2+{\Pi} \left(u^{r}\right)^2-{\pi^{\eta}{}_{\eta}} \left(u^{r}\right)^2-{\pi^{\phi}{}_{\phi}} \left(u^{r}\right)^2\right)\\
  &\quad - {p''(T)} {T}\Bigl(- 3 {\pi^{\phi}{}_{\phi}} {\tau_{\Pi}} {\tau_{\pi}} + {\pi^{\phi}{}_{\phi}} {\tau_{\Pi}} {\tau_{\pi}} \left(u^{r}\right)^2 + {\pi^{\eta}{}_{\eta}} {\tau_{\Pi}} {\tau_{\pi}} \left(\left(u^{r}\right)^2-3\right) + 3 {\Pi} {\tau_{\Pi}} {\tau_{\pi}} \left(u^{\tau}\right)^2\\
  &\quad - 4 {\tau_{\Pi}} \left(u^{r}\right)^2 {\eta} - 3 {\tau_{\pi}} \left(u^{r}\right)^2 {\zeta}\Bigr) \\
\mathcal{D}_{2} &= - 3 \left(p'(T)\right)^2 {\tau_{\Pi}} {\tau_{\pi}} \left(u^{r}\right)^2 {T}\\
  &\quad + 3 {p'(T)} {\tau_{\Pi}} {\tau_{\pi}} \left({p''(T)} \left(u^{r}\right)^2 {T}^2+{p''(T)} {T}^2-{\Pi} \left(u^{r}\right)^2+{\pi^{\eta}{}_{\eta}} \left(u^{r}\right)^2+{\pi^{\phi}{}_{\phi}} \left(u^{r}\right)^2\right)\\
  &\quad + {p''(T)} {T}\Bigl(- 3 {\pi^{\phi}{}_{\phi}} {\tau_{\Pi}} {\tau_{\pi}} + {\pi^{\phi}{}_{\phi}} {\tau_{\Pi}} {\tau_{\pi}} \left(u^{r}\right)^2 + {\pi^{\eta}{}_{\eta}} {\tau_{\Pi}} {\tau_{\pi}} \left(\left(u^{r}\right)^2-3\right) + 3 {\Pi} {\tau_{\Pi}} {\tau_{\pi}} \left(u^{\tau}\right)^2\\
  &\quad - 4 {\tau_{\Pi}} \left(u^{r}\right)^2 {\eta} - 3 {\tau_{\pi}} \left(u^{r}\right)^2 {\zeta}\Bigr) \\
\mathcal{D}_{3} &= 3 {p'(T)} {\tau_{\Pi}} {\tau_{\pi}} \left(u^{r}\right)^2 (-{p'(T)} {T}-{\Pi}+{\pi^{\eta}{}_{\eta}}+{\pi^{\phi}{}_{\phi}})\\
  &\quad + {p''(T)} {T}\Bigl(3 {p'(T)} {\tau_{\Pi}} {\tau_{\pi}} {T} \left(u^{\tau}\right)^2\\
  &\quad + {\tau_{\Pi}} \left(\left(u^{r}\right)^2 ({\tau_{\pi}} (3 {\Pi}+{\pi^{\eta}{}_{\eta}}+{\pi^{\phi}{}_{\phi}})-4 {\eta})+3 {\Pi} {\tau_{\pi}}-3 {\tau_{\pi}} ({\pi^{\eta}{}_{\eta}}+{\pi^{\phi}{}_{\phi}})\right)\\
  &\quad - 3 {\tau_{\pi}} \left(u^{r}\right)^2 {\zeta}\Bigr) \\
\mathcal{D}_{4} &= 3 {p'(T)} {\tau_{\Pi}} {\tau_{\pi}} \left(u^{r}\right)^2 ({p'(T)} {T}+{\Pi}-{\pi^{\eta}{}_{\eta}}-{\pi^{\phi}{}_{\phi}})\\
  &\quad - {p''(T)} {T}\Bigl(3 {p'(T)} {\tau_{\Pi}} {\tau_{\pi}} {T} \left(u^{\tau}\right)^2\\
  &\quad + {\tau_{\Pi}} \left(\left(u^{r}\right)^2 ({\tau_{\pi}} (3 {\Pi}+{\pi^{\eta}{}_{\eta}}+{\pi^{\phi}{}_{\phi}})-4 {\eta})+3 {\Pi} {\tau_{\pi}}-3 {\tau_{\pi}} ({\pi^{\eta}{}_{\eta}}+{\pi^{\phi}{}_{\phi}})\right)\\
  &\quad - 3 {\tau_{\pi}} \left(u^{r}\right)^2 {\zeta}\Bigr)
\end{align*}
\begin{align*}
L_{61} &= \frac{1}{{u^{\tau}} {\mathcal{D}_{1}}}\Bigl(3 {\tau_{\Pi}} {\tau_{\pi}} {u^{r}} ({p''(T)} {T}+{p'(T)}) \left({n} {u^{r}} \left(u^{\tau}\right)^2+{\nu^{r}}\right)\\
  &\quad + {(\partial_{T} n)} \left(u^{\tau}\right)^2\Bigl(3 {\pi^{\phi}{}_{\phi}} {\tau_{\Pi}} {\tau_{\pi}} + 2 {\pi^{\phi}{}_{\phi}} {\tau_{\Pi}} {\tau_{\pi}} \left(u^{r}\right)^2 - 3 {\Pi} {\tau_{\Pi}} {\tau_{\pi}} \left(2 \left(u^{r}\right)^2+1\right)\\
  &\quad + {\pi^{\eta}{}_{\eta}} {\tau_{\Pi}} {\tau_{\pi}} \left(2 \left(u^{r}\right)^2+3\right) - 3 {p'(T)} {\tau_{\Pi}} {\tau_{\pi}} {T} - 6 {p'(T)} {\tau_{\Pi}} {\tau_{\pi}} \left(u^{r}\right)^2 {T}\\
  &\quad + 4 {\tau_{\Pi}} \left(u^{r}\right)^2 {\eta} + 3 {\tau_{\pi}} \left(u^{r}\right)^2 {\zeta}\Bigr)\Bigr) \\
L_{62} &= \frac{1}{\left(u^{\tau}\right)^2 {\mathcal{D}_{1}}}\Bigl(- 3 {\tau_{\Pi}} {\tau_{\pi}} \left({n} {u^{r}} \left(u^{\tau}\right)^2+{\nu^{r}}\right) \left({p''(T)} {T} \left(u^{\tau}\right)^2+{p'(T)} \left(u^{r}\right)^2\right)\\
  &\quad + {(\partial_{T} n)} {u^{r}} \left(u^{\tau}\right)^2\Bigl(- 6 {\pi^{\phi}{}_{\phi}} {\tau_{\Pi}} {\tau_{\pi}} - 2 {\pi^{\phi}{}_{\phi}} {\tau_{\Pi}} {\tau_{\pi}} \left(u^{r}\right)^2 - 2 {\pi^{\eta}{}_{\eta}} {\tau_{\Pi}} {\tau_{\pi}} \left(\left(u^{r}\right)^2+3\right)\\
  &\quad + 6 {p'(T)} {\tau_{\Pi}} {\tau_{\pi}} {T} + 6 {p'(T)} {\tau_{\Pi}} {\tau_{\pi}} \left(u^{r}\right)^2 {T} + 6 {\Pi} {\tau_{\Pi}} {\tau_{\pi}} \left(u^{\tau}\right)^2\\
  &\quad - 4 {\tau_{\Pi}} \left(u^{r}\right)^2 {\eta} - 3 {\tau_{\pi}} \left(u^{r}\right)^2 {\zeta}\Bigr)\Bigr) \\
L_{63} &= \frac{3 {r}^2 {\tau_{\Pi}} {u^{r}} \left({(\partial_{T} n)} {u^{r}} \left(u^{\tau}\right)^2 ({p'(T)} {T}+{\Pi}-{\pi^{\eta}{}_{\eta}}-{\pi^{\phi}{}_{\phi}})-{p''(T)} {T} \left({n} {u^{r}} \left(u^{\tau}\right)^2+{\nu^{r}}\right)\right)}{\left(u^{\tau}\right)^2 {\mathcal{D}_{1}}} \\
L_{64} &= \frac{3 {\tau}^2 {\tau_{\Pi}} {u^{r}} \left({(\partial_{T} n)} {u^{r}} \left(u^{\tau}\right)^2 ({p'(T)} {T}+{\Pi}-{\pi^{\eta}{}_{\eta}}-{\pi^{\phi}{}_{\phi}})-{p''(T)} {T} \left({n} {u^{r}} \left(u^{\tau}\right)^2+{\nu^{r}}\right)\right)}{\left(u^{\tau}\right)^2 {\mathcal{D}_{1}}} \\
L_{65} &= \frac{3 {\tau_{\pi}} {u^{r}} \left({p''(T)} {T} \left({n} {u^{r}} \left(u^{\tau}\right)^2+{\nu^{r}}\right)-{(\partial_{T} n)} {u^{r}} \left(u^{\tau}\right)^2 ({p'(T)} {T}+{\Pi}-{\pi^{\eta}{}_{\eta}}-{\pi^{\phi}{}_{\phi}})\right)}{\left(u^{\tau}\right)^2 {\mathcal{D}_{1}}} \\
L_{71} &= -\frac{3 {\nu^{r}} {\tau_{\Pi}} {\tau_{n}} {\tau_{\pi}} \left(u^{r}\right)^2 {u^{\tau}} ({p''(T)} {T}+{p'(T)})}{{\mathcal{D}_{1}}} \\
L_{72} &= \frac{3 {\nu^{r}} {\tau_{\Pi}} {\tau_{n}} {\tau_{\pi}} {u^{r}} \left({p''(T)} {T} \left(u^{\tau}\right)^2+{p'(T)} \left(u^{r}\right)^2\right)}{{\mathcal{D}_{1}}} \\
L_{73} &= -\frac{3 {p''(T)} {\nu^{r}} {r}^2 {\tau_{\Pi}} {\tau_{n}} \left(u^{r}\right)^2 {T}}{{\mathcal{D}_{2}}} \\
L_{74} &= -\frac{3 {p''(T)} {\nu^{r}} {\tau}^2 {\tau_{\Pi}} {\tau_{n}} \left(u^{r}\right)^2 {T}}{{\mathcal{D}_{2}}} \\
L_{75} &= \frac{3 {p''(T)} {\nu^{r}} {\tau_{n}} {\tau_{\pi}} \left(u^{r}\right)^2 {T}}{{\mathcal{D}_{2}}}
\end{align*}
The file, charm rows (M-V7e, M-G7e, M-I3ne, M-I3ae):

\begin{align*}
(A_t)_{66} &= {(\partial_{\alpha} n)} {u^{\tau}} \\
(A_t)_{67} &= \frac{{u^{r}}}{{u^{\tau}}} \\
(A_t)_{76} &= {\kappa} {u^{r}} {u^{\tau}} \\
(A_t)_{77} &= {\tau_{n}} {u^{\tau}} \\
(A_r)_{61} &= -\frac{3}{\left(u^{\tau}\right)^2 {\mathcal{D}_{3}}}{p'(T)} {\tau_{\Pi}} {\tau_{\pi}}\Bigl({(\partial_{T} n)} {u^{r}} \left(u^{\tau}\right)^2 (-{p'(T)} {T}-{\Pi}+{\pi^{\eta}{}_{\eta}}+{\pi^{\phi}{}_{\phi}})\\
  &\quad + {p''(T)} {T} \left({n} {u^{r}} \left(u^{\tau}\right)^2+{\nu^{r}}\right)\Bigr) \\
(A_r)_{62} &= \frac{1}{\left(u^{\tau}\right)^2 {\mathcal{D}_{3}}}\Bigl(3 {p'(T)} {\nu^{r}} {\tau_{\Pi}} {\tau_{\pi}} {u^{r}} ({p'(T)} {T}+{\Pi}-{\pi^{\eta}{}_{\eta}}-{\pi^{\phi}{}_{\phi}})\\
  &\quad - 3 {(\partial_{T} n)} {\tau_{\Pi}} {\tau_{\pi}} \left(u^{\tau}\right)^2 (-{p'(T)} {T}-{\Pi}+{\pi^{\eta}{}_{\eta}}+{\pi^{\phi}{}_{\phi}})^2\\
  &\quad + {p''(T)} {T}\Bigl(3 {n} {\tau_{\Pi}} {\tau_{\pi}} \left(u^{\tau}\right)^2 ({p'(T)} {T}+{\Pi}-{\pi^{\eta}{}_{\eta}}-{\pi^{\phi}{}_{\phi}})\\
  &\quad - {\nu^{r}} {u^{r}}\Bigl(3 {\Pi} {\tau_{\Pi}} {\tau_{\pi}} + {\pi^{\eta}{}_{\eta}} {\tau_{\Pi}} {\tau_{\pi}} + {\pi^{\phi}{}_{\phi}} {\tau_{\Pi}} {\tau_{\pi}} + 3 {p'(T)} {\tau_{\Pi}} {\tau_{\pi}} {T}\\
  &\quad - 4 {\tau_{\Pi}} {\eta} - 3 {\tau_{\pi}} {\zeta}\Bigr)\Bigr)\Bigr) \\
(A_r)_{63} &= \frac{3 {\tau_{\Pi}} {\tau_{\pi}} \left({p''(T)} {T} \left({n} {u^{r}} \left(u^{\tau}\right)^2+{\nu^{r}}\right)+{(\partial_{T} n)} {u^{r}} \left(u^{\tau}\right)^2 (-{p'(T)} {T}-{\Pi}+{\pi^{\eta}{}_{\eta}}+{\pi^{\phi}{}_{\phi}})\right)}{\left(u^{\tau}\right)^2 {\mathcal{D}_{3}}} \\
(A_r)_{64} &= \frac{3 {\tau_{\Pi}} {\tau_{\pi}} \left({p''(T)} {T} \left({n} {u^{r}} \left(u^{\tau}\right)^2+{\nu^{r}}\right)+{(\partial_{T} n)} {u^{r}} \left(u^{\tau}\right)^2 (-{p'(T)} {T}-{\Pi}+{\pi^{\eta}{}_{\eta}}+{\pi^{\phi}{}_{\phi}})\right)}{\left(u^{\tau}\right)^2 {\mathcal{D}_{3}}} \\
(A_r)_{65} &= -\frac{3 {\tau_{\Pi}} {\tau_{\pi}} \left({p''(T)} {T} \left({n} {u^{r}} \left(u^{\tau}\right)^2+{\nu^{r}}\right)+{(\partial_{T} n)} {u^{r}} \left(u^{\tau}\right)^2 (-{p'(T)} {T}-{\Pi}+{\pi^{\eta}{}_{\eta}}+{\pi^{\phi}{}_{\phi}})\right)}{\left(u^{\tau}\right)^2 {\mathcal{D}_{3}}} \\
(A_r)_{66} &= {(\partial_{\alpha} n)} {u^{r}} \\
(A_r)_{67} &= 1 \\
(A_r)_{71} &= \frac{3 {p''(T)} {p'(T)} {\nu^{r}} {\tau_{\Pi}} {\tau_{n}} {\tau_{\pi}} {u^{r}} {T}}{{\mathcal{D}_{3}}} \\
(A_r)_{72} &= \frac{1}{\left(u^{\tau}\right)^2 {\mathcal{D}_{3}}}{\nu^{r}} {\tau_{n}} \left(u^{r}\right)^2\Bigl(3 {p'(T)} {\tau_{\Pi}} {\tau_{\pi}} (-{\Pi}+{\pi^{\eta}{}_{\eta}}+{\pi^{\phi}{}_{\phi}}) - 3 \left(p'(T)\right)^2 {\tau_{\Pi}} {\tau_{\pi}} {T}\\
  &\quad + {p''(T)} {T} (4 {\tau_{\Pi}} {\tau_{\pi}} ({\pi^{\eta}{}_{\eta}}+{\pi^{\phi}{}_{\phi}})-4 {\tau_{\Pi}} {\eta}-3 {\tau_{\pi}} {\zeta})\Bigr) \\
(A_r)_{73} &= -\frac{3 {p''(T)} {\nu^{r}} {\tau_{\Pi}} {\tau_{n}} {\tau_{\pi}} {u^{r}} {T}}{{\mathcal{D}_{3}}} \\
(A_r)_{74} &= -\frac{3 {p''(T)} {\nu^{r}} {\tau_{\Pi}} {\tau_{n}} {\tau_{\pi}} {u^{r}} {T}}{{\mathcal{D}_{3}}} \\
(A_r)_{75} &= \frac{3 {p''(T)} {\nu^{r}} {\tau_{\Pi}} {\tau_{n}} {\tau_{\pi}} {u^{r}} {T}}{{\mathcal{D}_{3}}} \\
(A_r)_{76} &= {\kappa} \left(u^{\tau}\right)^2 \\
(A_r)_{77} &= {\tau_{n}} {u^{r}} \\
S_{6} &= \frac{1}{{r} {\tau} \left(u^{\tau}\right)^3 {\mathcal{D}_{4}}}\Bigl(3{p'(T)} {\tau_{\Pi}} {\tau_{\pi}} {u^{r}}\Bigl({p'(T)} {\nu^{r}} {T} \left({r} \left(\left(u^{r}\right)^4-1\right)+{\tau} \left(u^{r}\right)^3 {u^{\tau}}\right)\\
  &\quad - {n} {u^{r}} \left(u^{\tau}\right)^2\Bigl({r} \left(u^{\tau}\right)^2 (2 {\pi^{\eta}{}_{\eta}}+{\pi^{\phi}{}_{\phi}})+{\tau} {u^{r}} {u^{\tau}} ({\pi^{\eta}{}_{\eta}}+2 {\pi^{\phi}{}_{\phi}})\Bigr)\\
  &\quad + {\nu^{r}}\Bigl({\tau} {u^{r}} {u^{\tau}} \left(\left(u^{r}\right)^2 ({\Pi}-{\pi^{\eta}{}_{\eta}}-{\pi^{\phi}{}_{\phi}})-{\pi^{\eta}{}_{\eta}}-2 {\pi^{\phi}{}_{\phi}}\right)\\
  &\quad - {r} \left(u^{\tau}\right)^2\Bigl(-{\Pi} \left(u^{r}\right)^2+{\Pi}+{\pi^{\eta}{}_{\eta}} \left(u^{r}\right)^2+{\pi^{\eta}{}_{\eta}}+{\pi^{\phi}{}_{\phi}} \left(u^{r}\right)^2\Bigr)\Bigr)\Bigr)\\
  &\quad + {(\partial_{T} n)} \left(u^{\tau}\right)^2\Bigl(3 \left(p'(T)\right)^2 {\tau_{\Pi}} {\tau_{\pi}} {T}^2 \left(u^{\tau}\right)^2 \left({r} \left(u^{r}\right)^2+{r}+{\tau} {u^{r}} {u^{\tau}}\right)\\
  &\quad - 3{p'(T)} {T}\Bigl(2{\tau} {\tau_{\Pi}} {u^{r}} {u^{\tau}}\Bigl(\left(u^{r}\right)^2 ({\eta}-{\tau_{\pi}} ({\Pi}+{\pi^{\phi}{}_{\phi}}))-{\tau_{\pi}} ({\Pi}+{\pi^{\phi}{}_{\phi}})\Bigr)\\
  &\quad + {r}\Bigl(- {\Pi} {\tau} {\tau_{\pi}} \left(u^{r}\right)^2 {u^{\tau}}\\
  &\quad + {\tau_{\Pi}}\Bigl(- 2 {\Pi} {\tau_{\pi}} + {\pi^{\phi}{}_{\phi}} {\tau_{\pi}} + 2 \left(u^{r}\right)^4 ({\eta}-{\tau_{\pi}} ({\Pi}+{\pi^{\eta}{}_{\eta}}))\\
  &\quad + \left(u^{r}\right)^2\Bigl({\tau_{\pi}} (-4 {\Pi}-2 {\pi^{\eta}{}_{\eta}}+{\pi^{\phi}{}_{\phi}})+{\tau} {u^{\tau}} ({\pi^{\eta}{}_{\eta}}+{\pi^{\phi}{}_{\phi}})+2 {\eta}\Bigr)\Bigr)\Bigr)\Bigr)\\
  &\quad - {\tau} {u^{r}} {u^{\tau}}\Bigl(- 3 {\tau_{\Pi}} {\tau_{\pi}} ({\Pi}-{\pi^{\eta}{}_{\eta}}-{\pi^{\phi}{}_{\phi}}) ({\Pi}+{\pi^{\eta}{}_{\eta}}+3 {\pi^{\phi}{}_{\phi}})\\
  &\quad + {\tau_{\Pi}} \left(u^{r}\right)^2\Bigl(- {\tau_{\pi}} \left(3 {\Pi}^2+6 {\Pi} {\pi^{\phi}{}_{\phi}}+\left(\pi^{\eta}{}_{\eta}\right)^2-\left(\pi^{\phi}{}_{\phi}\right)^2\right)\\
  &\quad + 2 {\eta} (3 {\Pi}-{\pi^{\eta}{}_{\eta}}+{\pi^{\phi}{}_{\phi}})\Bigr)\\
  &\quad + 3 {\tau_{\pi}} \left(u^{r}\right)^2 {\zeta} ({\pi^{\eta}{}_{\eta}}+2 {\pi^{\phi}{}_{\phi}})\Bigr)\\
  &\quad - {r}\Bigl({\tau_{\Pi}}\Bigl(- 3 {\tau_{\pi}} ({\Pi}+{\pi^{\eta}{}_{\eta}}) ({\Pi}-{\pi^{\eta}{}_{\eta}}-{\pi^{\phi}{}_{\phi}})\\
  &\quad + \left(u^{r}\right)^2\Bigl(- {\tau_{\pi}}\Bigl(6 {\Pi}^2 - 3 {\Pi} ({\pi^{\phi}{}_{\phi}}-2 {\pi^{\eta}{}_{\eta}})\\
  &\quad + ({\pi^{\phi}{}_{\phi}}-4 {\pi^{\eta}{}_{\eta}}) ({\pi^{\eta}{}_{\eta}}+{\pi^{\phi}{}_{\phi}})\Bigr)\\
  &\quad + 3 {\tau} {u^{\tau}} ({\pi^{\eta}{}_{\eta}}+{\pi^{\phi}{}_{\phi}}) ({\Pi}-{\pi^{\eta}{}_{\eta}}-{\pi^{\phi}{}_{\phi}})\\
  &\quad + 2 {\eta} (3 {\Pi}+{\pi^{\eta}{}_{\eta}}-{\pi^{\phi}{}_{\phi}})\Bigr)\\
  &\quad + \left(u^{r}\right)^4\Bigl(- 3 {\Pi}^2 {\tau_{\pi}} + ({\pi^{\eta}{}_{\eta}}-{\pi^{\phi}{}_{\phi}}) ({\pi^{\eta}{}_{\eta}} {\tau_{\pi}}+{\pi^{\phi}{}_{\phi}} {\tau_{\pi}}+2 {\eta})\\
  &\quad + {\Pi} (6 {\eta}-6 {\pi^{\eta}{}_{\eta}} {\tau_{\pi}})\Bigr)\Bigr)\\
  &\quad + 3{\tau_{\pi}} \left(u^{r}\right)^2\Bigl({\Pi} {\tau} {u^{\tau}} (-{\Pi}+{\pi^{\eta}{}_{\eta}}+{\pi^{\phi}{}_{\phi}})+\left(u^{\tau}\right)^2 {\zeta} (2 {\pi^{\eta}{}_{\eta}}+{\pi^{\phi}{}_{\phi}})\Bigr)\Bigr)\Bigr)\\
  &\quad - {p''(T)} {T}\Bigl(3{p'(T)} {\tau_{\Pi}} {\tau_{\pi}} {T} \left(u^{\tau}\right)^4\Bigl({n} {r} \left(u^{\tau}\right)^2+{n} {\tau} {u^{r}} {u^{\tau}}+{\nu^{r}} {r} {u^{r}}+{\nu^{r}} {\tau} {u^{\tau}}\Bigr)\\
  &\quad - 3{n} \left(u^{\tau}\right)^2\Bigl({\tau} {\tau_{\Pi}} {u^{r}} {u^{\tau}}\Bigl(\left(u^{r}\right)^2 (2 {\eta}-{\tau_{\pi}} ({\Pi}+{\pi^{\phi}{}_{\phi}}))-{\tau_{\pi}} ({\Pi}+{\pi^{\phi}{}_{\phi}})\Bigr)\\
  &\quad + {r}\Bigl(- {\Pi} {\tau} {\tau_{\pi}} \left(u^{r}\right)^2 {u^{\tau}}\\
  &\quad + {\tau_{\Pi}}\Bigl({\tau_{\pi}} (-{\Pi}+{\pi^{\eta}{}_{\eta}}+{\pi^{\phi}{}_{\phi}})\\
  &\quad + \left(u^{r}\right)^4 (2 {\eta}-{\tau_{\pi}} ({\Pi}+{\pi^{\eta}{}_{\eta}}))\\
  &\quad + \left(u^{r}\right)^2 ({\tau_{\pi}} ({\pi^{\phi}{}_{\phi}}-2 {\Pi})+{\tau} {u^{\tau}} ({\pi^{\eta}{}_{\eta}}+{\pi^{\phi}{}_{\phi}})+2 {\eta})\Bigr)\Bigr)\Bigr)\\
  &\quad + {\nu^{r}}\Bigl({\tau} {u^{\tau}}\Bigl({\tau_{\Pi}}\Bigl(3 {\tau_{\pi}} ({\Pi}+{\pi^{\phi}{}_{\phi}})\\
  &\quad + \left(u^{r}\right)^4 ({\tau_{\pi}} (3 {\Pi}+{\pi^{\eta}{}_{\eta}}+{\pi^{\phi}{}_{\phi}})-4 {\eta})\\
  &\quad - 3 \left(u^{r}\right)^2 (-2 {\Pi} {\tau_{\pi}}+{\pi^{\eta}{}_{\eta}} {\tau_{\pi}}+2 {\eta})\Bigr) - 3 {\tau_{\pi}} \left(u^{r}\right)^4 {\zeta}\Bigr)\\
  &\quad + {r} {u^{r}}\Bigl(- 3 {\tau} {\tau_{\Pi}} {u^{\tau}} ({\pi^{\eta}{}_{\eta}}+{\pi^{\phi}{}_{\phi}})\\
  &\quad + {\tau_{\Pi}} {\tau_{\pi}} \left(u^{\tau}\right)^2\Bigl(3 {\Pi} \left(u^{\tau}\right)^2+\left(u^{r}\right)^2 ({\pi^{\eta}{}_{\eta}}+{\pi^{\phi}{}_{\phi}})-{\pi^{\eta}{}_{\eta}}-4 {\pi^{\phi}{}_{\phi}}\Bigr)\\
  &\quad - 2 {\tau_{\Pi}} \left(2 \left(u^{r}\right)^4+3 \left(u^{r}\right)^2+1\right) {\eta}\\
  &\quad + 3 {\tau_{\pi}} \left({\Pi} {\tau} {u^{\tau}}+\left(u^{r}\right)^4 (-{\zeta})+{\zeta}\right)\Bigr)\Bigr)\Bigr)\Bigr) \\
S_{7} &= \frac{1}{{r} {\tau} {\mathcal{D}_{3}}}{\nu^{r}}\Bigl(- 3{p'(T)} {\tau_{\Pi}} {\tau_{\pi}} \left(u^{r}\right)^2\Bigl({r} {\tau} ({\Pi}-{\pi^{\eta}{}_{\eta}}-{\pi^{\phi}{}_{\phi}}) + {\tau} {\tau_{n}} {u^{r}} ({\Pi}+{\pi^{\phi}{}_{\phi}})\\
  &\quad + {r} {\tau_{n}} {u^{\tau}} ({\Pi}+{\pi^{\eta}{}_{\eta}}) + {p'(T)} {T} ({r} ({\tau}+{\tau_{n}} {u^{\tau}})+{\tau} {\tau_{n}} {u^{r}})\Bigr)\\
  &\quad + {p''(T)} {T}\Bigl(3 {p'(T)} {r} {\tau} {\tau_{\Pi}} {\tau_{\pi}} {T} \left(u^{\tau}\right)^2\\
  &\quad + {\tau} {\tau_{n}} {u^{r}}\Bigl({\pi^{\eta}{}_{\eta}} {\tau_{\Pi}} {\tau_{\pi}} \left(\left(u^{r}\right)^2-3\right) - 2 {\pi^{\phi}{}_{\phi}} {\tau_{\Pi}} {\tau_{\pi}} \left(\left(u^{r}\right)^2+3\right)\\
  &\quad + \left(u^{r}\right)^2 (2 {\tau_{\Pi}} {\eta}-3 {\tau_{\pi}} {\zeta})\Bigr)\\
  &\quad + {r}\Bigl({\tau} {\tau_{\Pi}}\Bigl(3 {\Pi} {\tau_{\pi}} - 3 {\tau_{\pi}} ({\pi^{\eta}{}_{\eta}}+{\pi^{\phi}{}_{\phi}})\\
  &\quad + \left(u^{r}\right)^2\Bigl({\tau_{\pi}} (3 {\Pi}+{\pi^{\eta}{}_{\eta}}+{\pi^{\phi}{}_{\phi}})+3 {\tau_{n}} ({\pi^{\eta}{}_{\eta}}+{\pi^{\phi}{}_{\phi}})-4 {\eta}\Bigr)\Bigr)\\
  &\quad - 3 {\tau} {\tau_{\pi}} \left(u^{r}\right)^2 ({\Pi} {\tau_{n}}+{\zeta})\\
  &\quad + {\tau_{n}} \left(u^{r}\right)^2 {u^{\tau}} ({\tau_{\Pi}} {\tau_{\pi}} ({\pi^{\phi}{}_{\phi}}-2 {\pi^{\eta}{}_{\eta}})+2 {\tau_{\Pi}} {\eta}-3 {\tau_{\pi}} {\zeta})\Bigr)\Bigr)\Bigr)
\end{align*}
Rows 1--5 are the viscous section's system at $\delta_{\pi\pi} = \tfrac{4}{3}\tau_{\pi}$ (M-V7m).

\subsubsection{\texttt{one\_\allowbreak{}d\_\allowbreak{}viscous\_\allowbreak{}matrix\_\allowbreak{}gamma} in \texttt{src/Matrix/1d\_\allowbreak{}viscous\_\allowbreak{}HQ\_\allowbreak{}cilindrical\_\allowbreak{}gamma.jl}}

Derived, charm rows:

\begin{align*}
(A_t)_{61} &= {(\partial_{T} n)} {\gamma} {u^{\tau}} \\
(A_t)_{62} &= \frac{{\gamma} {n} {u^{r}} \left(u^{\tau}\right)^2+{\nu^{r}}}{\left(u^{\tau}\right)^3} \\
(A_t)_{66} &= {n} {u^{\tau}} \\
(A_t)_{67} &= \frac{{u^{r}}}{{u^{\tau}}} \\
(A_t)_{72} &= -\frac{{\nu^{r}} {\tau_{n}} {u^{r}}}{{u^{\tau}}} \\
(A_t)_{76} &= {\kappa} {u^{r}} {u^{\tau}} \\
(A_t)_{77} &= {\tau_{n}} {u^{\tau}} \\
(A_r)_{61} &= {(\partial_{T} n)} {\gamma} {u^{r}} \\
(A_r)_{62} &= {\gamma} {n} \\
(A_r)_{66} &= {n} {u^{r}} \\
(A_r)_{67} &= 1 \\
(A_r)_{72} &= -\frac{{\nu^{r}} {\tau_{n}} \left(u^{r}\right)^2}{\left(u^{\tau}\right)^2} \\
(A_r)_{76} &= {\kappa} \left(u^{\tau}\right)^2 \\
(A_r)_{77} &= {\tau_{n}} {u^{r}} \\
S_{6} &= \frac{{\gamma} {n} \left({r} \left(u^{r}\right)^2+{r}+{\tau} {u^{r}} {u^{\tau}}\right)+{\nu^{r}} ({r} {u^{r}}+{\tau} {u^{\tau}})}{{r} {\tau} {u^{\tau}}} \\
S_{7} &= {\nu^{r}}
\end{align*}
With $c$ the charm rows and $m$ the medium rows, $\mathrm{row}_{c}\to\mathrm{row}_{c} - L\,\mathrm{row}_{m}$, $L = A_t^{cm}(A_t^{mm})^{-1}$, removes $\partial_{\tau}$ of the medium from the charm rows:

\begin{align*}
\mathcal{D}_{1} &= 3 \left(p'(T)\right)^2 {\tau_{\Pi}} {\tau_{\pi}} \left(u^{r}\right)^2 {T}\\
  &\quad + 3 {p'(T)} {\tau_{\Pi}} {\tau_{\pi}} \left(-{p''(T)} \left(u^{r}\right)^2 {T}^2-{p''(T)} {T}^2+{\Pi} \left(u^{r}\right)^2-{\pi^{\eta}{}_{\eta}} \left(u^{r}\right)^2-{\pi^{\phi}{}_{\phi}} \left(u^{r}\right)^2\right)\\
  &\quad - {p''(T)} {T}\Bigl(- 3 {\pi^{\phi}{}_{\phi}} {\tau_{\Pi}} {\tau_{\pi}} + {\pi^{\phi}{}_{\phi}} {\tau_{\Pi}} {\tau_{\pi}} \left(u^{r}\right)^2 + {\pi^{\eta}{}_{\eta}} {\tau_{\Pi}} {\tau_{\pi}} \left(\left(u^{r}\right)^2-3\right) + 3 {\Pi} {\tau_{\Pi}} {\tau_{\pi}} \left(u^{\tau}\right)^2\\
  &\quad - 4 {\tau_{\Pi}} \left(u^{r}\right)^2 {\eta} - 3 {\tau_{\pi}} \left(u^{r}\right)^2 {\zeta}\Bigr) \\
\mathcal{D}_{2} &= - 3 \left(p'(T)\right)^2 {\tau_{\Pi}} {\tau_{\pi}} \left(u^{r}\right)^2 {T}\\
  &\quad + 3 {p'(T)} {\tau_{\Pi}} {\tau_{\pi}} \left({p''(T)} \left(u^{r}\right)^2 {T}^2+{p''(T)} {T}^2-{\Pi} \left(u^{r}\right)^2+{\pi^{\eta}{}_{\eta}} \left(u^{r}\right)^2+{\pi^{\phi}{}_{\phi}} \left(u^{r}\right)^2\right)\\
  &\quad + {p''(T)} {T}\Bigl(- 3 {\pi^{\phi}{}_{\phi}} {\tau_{\Pi}} {\tau_{\pi}} + {\pi^{\phi}{}_{\phi}} {\tau_{\Pi}} {\tau_{\pi}} \left(u^{r}\right)^2 + {\pi^{\eta}{}_{\eta}} {\tau_{\Pi}} {\tau_{\pi}} \left(\left(u^{r}\right)^2-3\right) + 3 {\Pi} {\tau_{\Pi}} {\tau_{\pi}} \left(u^{\tau}\right)^2\\
  &\quad - 4 {\tau_{\Pi}} \left(u^{r}\right)^2 {\eta} - 3 {\tau_{\pi}} \left(u^{r}\right)^2 {\zeta}\Bigr) \\
\mathcal{D}_{3} &= 3 {p'(T)} {\tau_{\Pi}} {\tau_{\pi}} \left(u^{r}\right)^2 (-{p'(T)} {T}-{\Pi}+{\pi^{\eta}{}_{\eta}}+{\pi^{\phi}{}_{\phi}})\\
  &\quad + {p''(T)} {T}\Bigl(3 {p'(T)} {\tau_{\Pi}} {\tau_{\pi}} {T} \left(u^{\tau}\right)^2\\
  &\quad + {\tau_{\Pi}} \left(\left(u^{r}\right)^2 ({\tau_{\pi}} (3 {\Pi}+{\pi^{\eta}{}_{\eta}}+{\pi^{\phi}{}_{\phi}})-4 {\eta})+3 {\Pi} {\tau_{\pi}}-3 {\tau_{\pi}} ({\pi^{\eta}{}_{\eta}}+{\pi^{\phi}{}_{\phi}})\right)\\
  &\quad - 3 {\tau_{\pi}} \left(u^{r}\right)^2 {\zeta}\Bigr) \\
\mathcal{D}_{4} &= 3 {p'(T)} {\tau_{\Pi}} {\tau_{\pi}} \left(u^{r}\right)^2 ({p'(T)} {T}+{\Pi}-{\pi^{\eta}{}_{\eta}}-{\pi^{\phi}{}_{\phi}})\\
  &\quad - {p''(T)} {T}\Bigl(3 {p'(T)} {\tau_{\Pi}} {\tau_{\pi}} {T} \left(u^{\tau}\right)^2\\
  &\quad + {\tau_{\Pi}} \left(\left(u^{r}\right)^2 ({\tau_{\pi}} (3 {\Pi}+{\pi^{\eta}{}_{\eta}}+{\pi^{\phi}{}_{\phi}})-4 {\eta})+3 {\Pi} {\tau_{\pi}}-3 {\tau_{\pi}} ({\pi^{\eta}{}_{\eta}}+{\pi^{\phi}{}_{\phi}})\right)\\
  &\quad - 3 {\tau_{\pi}} \left(u^{r}\right)^2 {\zeta}\Bigr)
\end{align*}
\begin{align*}
L_{61} &= \frac{1}{{u^{\tau}} {\mathcal{D}_{1}}}\Bigl(3 {\tau_{\Pi}} {\tau_{\pi}} {u^{r}} ({p''(T)} {T}+{p'(T)}) \left({\gamma} {n} {u^{r}} \left(u^{\tau}\right)^2+{\nu^{r}}\right)\\
  &\quad + {(\partial_{T} n)} {\gamma} \left(u^{\tau}\right)^2\Bigl(3 {\pi^{\phi}{}_{\phi}} {\tau_{\Pi}} {\tau_{\pi}} + 2 {\pi^{\phi}{}_{\phi}} {\tau_{\Pi}} {\tau_{\pi}} \left(u^{r}\right)^2 - 3 {\Pi} {\tau_{\Pi}} {\tau_{\pi}} \left(2 \left(u^{r}\right)^2+1\right)\\
  &\quad + {\pi^{\eta}{}_{\eta}} {\tau_{\Pi}} {\tau_{\pi}} \left(2 \left(u^{r}\right)^2+3\right) - 3 {p'(T)} {\tau_{\Pi}} {\tau_{\pi}} {T}\\
  &\quad - 6 {p'(T)} {\tau_{\Pi}} {\tau_{\pi}} \left(u^{r}\right)^2 {T} + 4 {\tau_{\Pi}} \left(u^{r}\right)^2 {\eta} + 3 {\tau_{\pi}} \left(u^{r}\right)^2 {\zeta}\Bigr)\Bigr) \\
L_{62} &= \frac{1}{\left(u^{\tau}\right)^2 {\mathcal{D}_{1}}}\Bigl(- 3 {\tau_{\Pi}} {\tau_{\pi}} \left({p''(T)} {T} \left(u^{\tau}\right)^2+{p'(T)} \left(u^{r}\right)^2\right) \left({\gamma} {n} {u^{r}} \left(u^{\tau}\right)^2+{\nu^{r}}\right)\\
  &\quad + {(\partial_{T} n)} {\gamma} {u^{r}} \left(u^{\tau}\right)^2\Bigl(- 6 {\pi^{\phi}{}_{\phi}} {\tau_{\Pi}} {\tau_{\pi}} - 2 {\pi^{\phi}{}_{\phi}} {\tau_{\Pi}} {\tau_{\pi}} \left(u^{r}\right)^2 - 2 {\pi^{\eta}{}_{\eta}} {\tau_{\Pi}} {\tau_{\pi}} \left(\left(u^{r}\right)^2+3\right)\\
  &\quad + 6 {p'(T)} {\tau_{\Pi}} {\tau_{\pi}} {T} + 6 {p'(T)} {\tau_{\Pi}} {\tau_{\pi}} \left(u^{r}\right)^2 {T} + 6 {\Pi} {\tau_{\Pi}} {\tau_{\pi}} \left(u^{\tau}\right)^2\\
  &\quad - 4 {\tau_{\Pi}} \left(u^{r}\right)^2 {\eta} - 3 {\tau_{\pi}} \left(u^{r}\right)^2 {\zeta}\Bigr)\Bigr) \\
L_{63} &= \frac{3 {r}^2 {\tau_{\Pi}} {u^{r}} \left({(\partial_{T} n)} {\gamma} {u^{r}} \left(u^{\tau}\right)^2 ({p'(T)} {T}+{\Pi}-{\pi^{\eta}{}_{\eta}}-{\pi^{\phi}{}_{\phi}})-{p''(T)} {T} \left({\gamma} {n} {u^{r}} \left(u^{\tau}\right)^2+{\nu^{r}}\right)\right)}{\left(u^{\tau}\right)^2 {\mathcal{D}_{1}}} \\
L_{64} &= \frac{3 {\tau}^2 {\tau_{\Pi}} {u^{r}} \left({(\partial_{T} n)} {\gamma} {u^{r}} \left(u^{\tau}\right)^2 ({p'(T)} {T}+{\Pi}-{\pi^{\eta}{}_{\eta}}-{\pi^{\phi}{}_{\phi}})-{p''(T)} {T} \left({\gamma} {n} {u^{r}} \left(u^{\tau}\right)^2+{\nu^{r}}\right)\right)}{\left(u^{\tau}\right)^2 {\mathcal{D}_{1}}} \\
L_{65} &= \frac{3 {\tau_{\pi}} {u^{r}} \left({p''(T)} {T} \left({\gamma} {n} {u^{r}} \left(u^{\tau}\right)^2+{\nu^{r}}\right)-{(\partial_{T} n)} {\gamma} {u^{r}} \left(u^{\tau}\right)^2 ({p'(T)} {T}+{\Pi}-{\pi^{\eta}{}_{\eta}}-{\pi^{\phi}{}_{\phi}})\right)}{\left(u^{\tau}\right)^2 {\mathcal{D}_{1}}} \\
L_{71} &= -\frac{3 {\nu^{r}} {\tau_{\Pi}} {\tau_{n}} {\tau_{\pi}} \left(u^{r}\right)^2 {u^{\tau}} ({p''(T)} {T}+{p'(T)})}{{\mathcal{D}_{1}}} \\
L_{72} &= \frac{3 {\nu^{r}} {\tau_{\Pi}} {\tau_{n}} {\tau_{\pi}} {u^{r}} \left({p''(T)} {T} \left(u^{\tau}\right)^2+{p'(T)} \left(u^{r}\right)^2\right)}{{\mathcal{D}_{1}}} \\
L_{73} &= -\frac{3 {p''(T)} {\nu^{r}} {r}^2 {\tau_{\Pi}} {\tau_{n}} \left(u^{r}\right)^2 {T}}{{\mathcal{D}_{2}}} \\
L_{74} &= -\frac{3 {p''(T)} {\nu^{r}} {\tau}^2 {\tau_{\Pi}} {\tau_{n}} \left(u^{r}\right)^2 {T}}{{\mathcal{D}_{2}}} \\
L_{75} &= \frac{3 {p''(T)} {\nu^{r}} {\tau_{n}} {\tau_{\pi}} \left(u^{r}\right)^2 {T}}{{\mathcal{D}_{2}}}
\end{align*}
The file, charm rows (M-V7e, M-G7e, M-I3ne, M-I3ae):

\begin{align*}
(A_t)_{66} &= {n} {u^{\tau}} \\
(A_t)_{67} &= \frac{{u^{r}}}{{u^{\tau}}} \\
(A_t)_{76} &= {\kappa} {u^{r}} {u^{\tau}} \\
(A_t)_{77} &= {\tau_{n}} {u^{\tau}} \\
(A_r)_{61} &= -\frac{3}{\left(u^{\tau}\right)^2 {\mathcal{D}_{3}}}{p'(T)} {\tau_{\Pi}} {\tau_{\pi}}\Bigl({(\partial_{T} n)} {\gamma} {u^{r}} \left(u^{\tau}\right)^2 (-{p'(T)} {T}-{\Pi}+{\pi^{\eta}{}_{\eta}}+{\pi^{\phi}{}_{\phi}})\\
  &\quad + {p''(T)} {T} \left({\gamma} {n} {u^{r}} \left(u^{\tau}\right)^2+{\nu^{r}}\right)\Bigr) \\
(A_r)_{62} &= \frac{1}{\left(u^{\tau}\right)^2 {\mathcal{D}_{3}}}\Bigl(3 {p'(T)} {\nu^{r}} {\tau_{\Pi}} {\tau_{\pi}} {u^{r}} ({p'(T)} {T}+{\Pi}-{\pi^{\eta}{}_{\eta}}-{\pi^{\phi}{}_{\phi}})\\
  &\quad - 3 {(\partial_{T} n)} {\gamma} {\tau_{\Pi}} {\tau_{\pi}} \left(u^{\tau}\right)^2 (-{p'(T)} {T}-{\Pi}+{\pi^{\eta}{}_{\eta}}+{\pi^{\phi}{}_{\phi}})^2\\
  &\quad + {p''(T)} {T}\Bigl(3 {\gamma} {n} {\tau_{\Pi}} {\tau_{\pi}} \left(u^{\tau}\right)^2 ({p'(T)} {T}+{\Pi}-{\pi^{\eta}{}_{\eta}}-{\pi^{\phi}{}_{\phi}})\\
  &\quad - {\nu^{r}} {u^{r}}\Bigl(3 {\Pi} {\tau_{\Pi}} {\tau_{\pi}} + {\pi^{\eta}{}_{\eta}} {\tau_{\Pi}} {\tau_{\pi}} + {\pi^{\phi}{}_{\phi}} {\tau_{\Pi}} {\tau_{\pi}} + 3 {p'(T)} {\tau_{\Pi}} {\tau_{\pi}} {T}\\
  &\quad - 4 {\tau_{\Pi}} {\eta} - 3 {\tau_{\pi}} {\zeta}\Bigr)\Bigr)\Bigr) \\
(A_r)_{63} &= \frac{3 {\tau_{\Pi}} {\tau_{\pi}} \left({p''(T)} {T} \left({\gamma} {n} {u^{r}} \left(u^{\tau}\right)^2+{\nu^{r}}\right)+{(\partial_{T} n)} {\gamma} {u^{r}} \left(u^{\tau}\right)^2 (-{p'(T)} {T}-{\Pi}+{\pi^{\eta}{}_{\eta}}+{\pi^{\phi}{}_{\phi}})\right)}{\left(u^{\tau}\right)^2 {\mathcal{D}_{3}}} \\
(A_r)_{64} &= \frac{3 {\tau_{\Pi}} {\tau_{\pi}} \left({p''(T)} {T} \left({\gamma} {n} {u^{r}} \left(u^{\tau}\right)^2+{\nu^{r}}\right)+{(\partial_{T} n)} {\gamma} {u^{r}} \left(u^{\tau}\right)^2 (-{p'(T)} {T}-{\Pi}+{\pi^{\eta}{}_{\eta}}+{\pi^{\phi}{}_{\phi}})\right)}{\left(u^{\tau}\right)^2 {\mathcal{D}_{3}}} \\
(A_r)_{65} &= -\frac{3 {\tau_{\Pi}} {\tau_{\pi}} \left({p''(T)} {T} \left({\gamma} {n} {u^{r}} \left(u^{\tau}\right)^2+{\nu^{r}}\right)+{(\partial_{T} n)} {\gamma} {u^{r}} \left(u^{\tau}\right)^2 (-{p'(T)} {T}-{\Pi}+{\pi^{\eta}{}_{\eta}}+{\pi^{\phi}{}_{\phi}})\right)}{\left(u^{\tau}\right)^2 {\mathcal{D}_{3}}} \\
(A_r)_{66} &= {n} {u^{r}} \\
(A_r)_{67} &= 1 \\
(A_r)_{71} &= \frac{3 {p''(T)} {p'(T)} {\nu^{r}} {\tau_{\Pi}} {\tau_{n}} {\tau_{\pi}} {u^{r}} {T}}{{\mathcal{D}_{3}}} \\
(A_r)_{72} &= \frac{1}{\left(u^{\tau}\right)^2 {\mathcal{D}_{3}}}{\nu^{r}} {\tau_{n}} \left(u^{r}\right)^2\Bigl(3 {p'(T)} {\tau_{\Pi}} {\tau_{\pi}} (-{\Pi}+{\pi^{\eta}{}_{\eta}}+{\pi^{\phi}{}_{\phi}}) - 3 \left(p'(T)\right)^2 {\tau_{\Pi}} {\tau_{\pi}} {T}\\
  &\quad + {p''(T)} {T} (4 {\tau_{\Pi}} {\tau_{\pi}} ({\pi^{\eta}{}_{\eta}}+{\pi^{\phi}{}_{\phi}})-4 {\tau_{\Pi}} {\eta}-3 {\tau_{\pi}} {\zeta})\Bigr) \\
(A_r)_{73} &= -\frac{3 {p''(T)} {\nu^{r}} {\tau_{\Pi}} {\tau_{n}} {\tau_{\pi}} {u^{r}} {T}}{{\mathcal{D}_{3}}} \\
(A_r)_{74} &= -\frac{3 {p''(T)} {\nu^{r}} {\tau_{\Pi}} {\tau_{n}} {\tau_{\pi}} {u^{r}} {T}}{{\mathcal{D}_{3}}} \\
(A_r)_{75} &= \frac{3 {p''(T)} {\nu^{r}} {\tau_{\Pi}} {\tau_{n}} {\tau_{\pi}} {u^{r}} {T}}{{\mathcal{D}_{3}}} \\
(A_r)_{76} &= {\kappa} \left(u^{\tau}\right)^2 \\
(A_r)_{77} &= {\tau_{n}} {u^{r}} \\
S_{6} &= \frac{1}{{r} {\tau} \left(u^{\tau}\right)^3 {\mathcal{D}_{4}}}\Bigl(3{p'(T)} {\tau_{\Pi}} {\tau_{\pi}} {u^{r}}\Bigl({p'(T)} {\nu^{r}} {T} \left({r} \left(\left(u^{r}\right)^4-1\right)+{\tau} \left(u^{r}\right)^3 {u^{\tau}}\right)\\
  &\quad - {\gamma} {n} {u^{r}} \left(u^{\tau}\right)^2\Bigl({r} \left(u^{\tau}\right)^2 (2 {\pi^{\eta}{}_{\eta}}+{\pi^{\phi}{}_{\phi}})+{\tau} {u^{r}} {u^{\tau}} ({\pi^{\eta}{}_{\eta}}+2 {\pi^{\phi}{}_{\phi}})\Bigr)\\
  &\quad + {\nu^{r}}\Bigl({\tau} {u^{r}} {u^{\tau}} \left(\left(u^{r}\right)^2 ({\Pi}-{\pi^{\eta}{}_{\eta}}-{\pi^{\phi}{}_{\phi}})-{\pi^{\eta}{}_{\eta}}-2 {\pi^{\phi}{}_{\phi}}\right)\\
  &\quad - {r} \left(u^{\tau}\right)^2\Bigl(-{\Pi} \left(u^{r}\right)^2+{\Pi}+{\pi^{\eta}{}_{\eta}} \left(u^{r}\right)^2+{\pi^{\eta}{}_{\eta}}+{\pi^{\phi}{}_{\phi}} \left(u^{r}\right)^2\Bigr)\Bigr)\Bigr)\\
  &\quad + {(\partial_{T} n)} {\gamma} \left(u^{\tau}\right)^2\Bigl(3 \left(p'(T)\right)^2 {\tau_{\Pi}} {\tau_{\pi}} {T}^2 \left(u^{\tau}\right)^2 \left({r} \left(u^{r}\right)^2+{r}+{\tau} {u^{r}} {u^{\tau}}\right)\\
  &\quad - 3{p'(T)} {T}\Bigl(2{\tau} {\tau_{\Pi}} {u^{r}} {u^{\tau}}\Bigl(\left(u^{r}\right)^2 ({\eta}-{\tau_{\pi}} ({\Pi}+{\pi^{\phi}{}_{\phi}}))-{\tau_{\pi}} ({\Pi}+{\pi^{\phi}{}_{\phi}})\Bigr)\\
  &\quad + {r}\Bigl(- {\Pi} {\tau} {\tau_{\pi}} \left(u^{r}\right)^2 {u^{\tau}}\\
  &\quad + {\tau_{\Pi}}\Bigl(- 2 {\Pi} {\tau_{\pi}} + {\pi^{\phi}{}_{\phi}} {\tau_{\pi}} + 2 \left(u^{r}\right)^4 ({\eta}-{\tau_{\pi}} ({\Pi}+{\pi^{\eta}{}_{\eta}}))\\
  &\quad + \left(u^{r}\right)^2\Bigl({\tau_{\pi}} (-4 {\Pi}-2 {\pi^{\eta}{}_{\eta}}+{\pi^{\phi}{}_{\phi}})+{\tau} {u^{\tau}} ({\pi^{\eta}{}_{\eta}}+{\pi^{\phi}{}_{\phi}})+2 {\eta}\Bigr)\Bigr)\Bigr)\Bigr)\\
  &\quad - {\tau} {u^{r}} {u^{\tau}}\Bigl(- 3 {\tau_{\Pi}} {\tau_{\pi}} ({\Pi}-{\pi^{\eta}{}_{\eta}}-{\pi^{\phi}{}_{\phi}}) ({\Pi}+{\pi^{\eta}{}_{\eta}}+3 {\pi^{\phi}{}_{\phi}})\\
  &\quad + {\tau_{\Pi}} \left(u^{r}\right)^2\Bigl(- {\tau_{\pi}} \left(3 {\Pi}^2+6 {\Pi} {\pi^{\phi}{}_{\phi}}+\left(\pi^{\eta}{}_{\eta}\right)^2-\left(\pi^{\phi}{}_{\phi}\right)^2\right)\\
  &\quad + 2 {\eta} (3 {\Pi}-{\pi^{\eta}{}_{\eta}}+{\pi^{\phi}{}_{\phi}})\Bigr)\\
  &\quad + 3 {\tau_{\pi}} \left(u^{r}\right)^2 {\zeta} ({\pi^{\eta}{}_{\eta}}+2 {\pi^{\phi}{}_{\phi}})\Bigr)\\
  &\quad - {r}\Bigl({\tau_{\Pi}}\Bigl(- 3 {\tau_{\pi}} ({\Pi}+{\pi^{\eta}{}_{\eta}}) ({\Pi}-{\pi^{\eta}{}_{\eta}}-{\pi^{\phi}{}_{\phi}})\\
  &\quad + \left(u^{r}\right)^2\Bigl(- {\tau_{\pi}}\Bigl(6 {\Pi}^2 - 3 {\Pi} ({\pi^{\phi}{}_{\phi}}-2 {\pi^{\eta}{}_{\eta}})\\
  &\quad + ({\pi^{\phi}{}_{\phi}}-4 {\pi^{\eta}{}_{\eta}}) ({\pi^{\eta}{}_{\eta}}+{\pi^{\phi}{}_{\phi}})\Bigr)\\
  &\quad + 3 {\tau} {u^{\tau}} ({\pi^{\eta}{}_{\eta}}+{\pi^{\phi}{}_{\phi}}) ({\Pi}-{\pi^{\eta}{}_{\eta}}-{\pi^{\phi}{}_{\phi}})\\
  &\quad + 2 {\eta} (3 {\Pi}+{\pi^{\eta}{}_{\eta}}-{\pi^{\phi}{}_{\phi}})\Bigr)\\
  &\quad + \left(u^{r}\right)^4\Bigl(- 3 {\Pi}^2 {\tau_{\pi}} + ({\pi^{\eta}{}_{\eta}}-{\pi^{\phi}{}_{\phi}}) ({\pi^{\eta}{}_{\eta}} {\tau_{\pi}}+{\pi^{\phi}{}_{\phi}} {\tau_{\pi}}+2 {\eta})\\
  &\quad + {\Pi} (6 {\eta}-6 {\pi^{\eta}{}_{\eta}} {\tau_{\pi}})\Bigr)\Bigr)\\
  &\quad + 3{\tau_{\pi}} \left(u^{r}\right)^2\Bigl({\Pi} {\tau} {u^{\tau}} (-{\Pi}+{\pi^{\eta}{}_{\eta}}+{\pi^{\phi}{}_{\phi}})+\left(u^{\tau}\right)^2 {\zeta} (2 {\pi^{\eta}{}_{\eta}}+{\pi^{\phi}{}_{\phi}})\Bigr)\Bigr)\Bigr)\\
  &\quad - {p''(T)} {T}\Bigl(3{p'(T)} {\tau_{\Pi}} {\tau_{\pi}} {T} \left(u^{\tau}\right)^4\Bigl({\gamma} {n} \left({r} \left(u^{r}\right)^2+{r}+{\tau} {u^{r}} {u^{\tau}}\right)+{\nu^{r}} ({r} {u^{r}}+{\tau} {u^{\tau}})\Bigr)\\
  &\quad - 3{\gamma} {n} \left(u^{\tau}\right)^2\Bigl({\tau} {\tau_{\Pi}} {u^{r}} {u^{\tau}}\Bigl(\left(u^{r}\right)^2 (2 {\eta}-{\tau_{\pi}} ({\Pi}+{\pi^{\phi}{}_{\phi}}))-{\tau_{\pi}} ({\Pi}+{\pi^{\phi}{}_{\phi}})\Bigr)\\
  &\quad + {r}\Bigl(- {\Pi} {\tau} {\tau_{\pi}} \left(u^{r}\right)^2 {u^{\tau}}\\
  &\quad + {\tau_{\Pi}}\Bigl({\tau_{\pi}} (-{\Pi}+{\pi^{\eta}{}_{\eta}}+{\pi^{\phi}{}_{\phi}})\\
  &\quad + \left(u^{r}\right)^4 (2 {\eta}-{\tau_{\pi}} ({\Pi}+{\pi^{\eta}{}_{\eta}}))\\
  &\quad + \left(u^{r}\right)^2 ({\tau_{\pi}} ({\pi^{\phi}{}_{\phi}}-2 {\Pi})+{\tau} {u^{\tau}} ({\pi^{\eta}{}_{\eta}}+{\pi^{\phi}{}_{\phi}})+2 {\eta})\Bigr)\Bigr)\Bigr)\\
  &\quad + {\nu^{r}}\Bigl({\tau} {u^{\tau}}\Bigl({\tau_{\Pi}}\Bigl(3 {\tau_{\pi}} ({\Pi}+{\pi^{\phi}{}_{\phi}})\\
  &\quad + \left(u^{r}\right)^4 ({\tau_{\pi}} (3 {\Pi}+{\pi^{\eta}{}_{\eta}}+{\pi^{\phi}{}_{\phi}})-4 {\eta})\\
  &\quad - 3 \left(u^{r}\right)^2 (-2 {\Pi} {\tau_{\pi}}+{\pi^{\eta}{}_{\eta}} {\tau_{\pi}}+2 {\eta})\Bigr) - 3 {\tau_{\pi}} \left(u^{r}\right)^4 {\zeta}\Bigr)\\
  &\quad + {r} {u^{r}}\Bigl(- 3 {\tau} {\tau_{\Pi}} {u^{\tau}} ({\pi^{\eta}{}_{\eta}}+{\pi^{\phi}{}_{\phi}})\\
  &\quad + {\tau_{\Pi}} {\tau_{\pi}} \left(u^{\tau}\right)^2\Bigl(3 {\Pi} \left(u^{\tau}\right)^2+\left(u^{r}\right)^2 ({\pi^{\eta}{}_{\eta}}+{\pi^{\phi}{}_{\phi}})-{\pi^{\eta}{}_{\eta}}-4 {\pi^{\phi}{}_{\phi}}\Bigr)\\
  &\quad - 2 {\tau_{\Pi}} \left(2 \left(u^{r}\right)^4+3 \left(u^{r}\right)^2+1\right) {\eta}\\
  &\quad + 3 {\tau_{\pi}} \left({\Pi} {\tau} {u^{\tau}}+\left(u^{r}\right)^4 (-{\zeta})+{\zeta}\right)\Bigr)\Bigr)\Bigr)\Bigr) \\
S_{7} &= \frac{1}{{r} {\tau} {\mathcal{D}_{3}}}{\nu^{r}}\Bigl(- 3{p'(T)} {\tau_{\Pi}} {\tau_{\pi}} \left(u^{r}\right)^2\Bigl({r} {\tau} ({\Pi}-{\pi^{\eta}{}_{\eta}}-{\pi^{\phi}{}_{\phi}}) + {\tau} {\tau_{n}} {u^{r}} ({\Pi}+{\pi^{\phi}{}_{\phi}})\\
  &\quad + {r} {\tau_{n}} {u^{\tau}} ({\Pi}+{\pi^{\eta}{}_{\eta}}) + {p'(T)} {T} ({r} ({\tau}+{\tau_{n}} {u^{\tau}})+{\tau} {\tau_{n}} {u^{r}})\Bigr)\\
  &\quad + {p''(T)} {T}\Bigl(3 {p'(T)} {r} {\tau} {\tau_{\Pi}} {\tau_{\pi}} {T} \left(u^{\tau}\right)^2\\
  &\quad + {\tau} {\tau_{n}} {u^{r}}\Bigl({\pi^{\eta}{}_{\eta}} {\tau_{\Pi}} {\tau_{\pi}} \left(\left(u^{r}\right)^2-3\right) - 2 {\pi^{\phi}{}_{\phi}} {\tau_{\Pi}} {\tau_{\pi}} \left(\left(u^{r}\right)^2+3\right)\\
  &\quad + \left(u^{r}\right)^2 (2 {\tau_{\Pi}} {\eta}-3 {\tau_{\pi}} {\zeta})\Bigr)\\
  &\quad + {r}\Bigl({\tau} {\tau_{\Pi}}\Bigl(3 {\Pi} {\tau_{\pi}} - 3 {\tau_{\pi}} ({\pi^{\eta}{}_{\eta}}+{\pi^{\phi}{}_{\phi}})\\
  &\quad + \left(u^{r}\right)^2\Bigl({\tau_{\pi}} (3 {\Pi}+{\pi^{\eta}{}_{\eta}}+{\pi^{\phi}{}_{\phi}})+3 {\tau_{n}} ({\pi^{\eta}{}_{\eta}}+{\pi^{\phi}{}_{\phi}})-4 {\eta}\Bigr)\Bigr)\\
  &\quad - 3 {\tau} {\tau_{\pi}} \left(u^{r}\right)^2 ({\Pi} {\tau_{n}}+{\zeta})\\
  &\quad + {\tau_{n}} \left(u^{r}\right)^2 {u^{\tau}} ({\tau_{\Pi}} {\tau_{\pi}} ({\pi^{\phi}{}_{\phi}}-2 {\pi^{\eta}{}_{\eta}})+2 {\tau_{\Pi}} {\eta}-3 {\tau_{\pi}} {\zeta})\Bigr)\Bigr)\Bigr)
\end{align*}
$\partial_{T}n$ above is $\gamma$-free: $n$ and $\partial_{T}n$ are those of $n(T)$ at $\gamma = 1$.

\subsubsection{Ideal medium, rows 1--2}

\begin{align*}
(A_t)_{11} &= \left(u^{r}\right)^2 ({p''(T)} {T}+{p'(T)})+{p''(T)} {T} \\
(A_t)_{12} &= 2 {p'(T)} {u^{r}} {T} \\
(A_t)_{21} &= {u^{r}} {u^{\tau}} ({p''(T)} {T}+{p'(T)}) \\
(A_t)_{22} &= \frac{{p'(T)} \left(2 \left(u^{r}\right)^2+1\right) {T}}{{u^{\tau}}} \\
(A_r)_{11} &= {u^{r}} {u^{\tau}} ({p''(T)} {T}+{p'(T)}) \\
(A_r)_{12} &= \frac{{p'(T)} \left(2 \left(u^{r}\right)^2+1\right) {T}}{{u^{\tau}}} \\
(A_r)_{21} &= {p''(T)} \left(u^{r}\right)^2 {T}+{p'(T)} \left(u^{r}\right)^2+{p'(T)} \\
(A_r)_{22} &= 2 {p'(T)} {u^{r}} {T} \\
S_{1} &= \frac{{p'(T)} {T} \left({r} \left(u^{r}\right)^2+{r}+{\tau} {u^{r}} {u^{\tau}}\right)}{{r} {\tau}} \\
S_{2} &= \frac{{p'(T)} {u^{r}} {T} ({r} {u^{\tau}}+{\tau} {u^{r}})}{{r} {\tau}}
\end{align*}
\subsubsection{\texttt{one\_\allowbreak{}d\_\allowbreak{}ideal\_\allowbreak{}matrix\_\allowbreak{}ruwen4} in \texttt{src/Matrix/1d\_\allowbreak{}ideal\_\allowbreak{}HQ\_\allowbreak{}cilindrical\_\allowbreak{}density.jl}}

Derived, charm rows:

\begin{align*}
(A_t)_{32} &= \frac{{n} {u^{r}}}{{u^{\tau}}} \\
(A_t)_{33} &= {u^{\tau}} \\
(A_r)_{32} &= {n} \\
(A_r)_{33} &= {u^{r}} \\
S_{3} &= {n} \left(\frac{{u^{r}}}{{r}}+\frac{{u^{\tau}}}{{\tau}}\right)
\end{align*}
With $c$ the charm rows and $m$ the medium rows, $\mathrm{row}_{c}\to\mathrm{row}_{c} - L\,\mathrm{row}_{m}$, $L = A_t^{cm}(A_t^{mm})^{-1}$, removes $\partial_{\tau}$ of the medium from the charm rows:

\begin{align*}
L_{31} &= \frac{{n} \left(u^{r}\right)^2 {u^{\tau}} ({p''(T)} {T}+{p'(T)})}{{p'(T)} {T} \left({p'(T)} \left(u^{r}\right)^2-{p''(T)} {T} \left(u^{\tau}\right)^2\right)} \\
L_{32} &= \frac{{n} {u^{r}} \left({p''(T)} {T} \left(u^{\tau}\right)^2+{p'(T)} \left(u^{r}\right)^2\right)}{{p'(T)} {T} \left({p''(T)} {T} \left(u^{\tau}\right)^2-{p'(T)} \left(u^{r}\right)^2\right)}
\end{align*}
The file, charm rows (M-V7e, M-G7e, M-I3ne, M-I3ae):

\begin{align*}
(A_t)_{33} &= {u^{\tau}} \\
(A_r)_{31} &= -\frac{{p''(T)} {n} {u^{r}}}{{p''(T)} {T} \left(u^{\tau}\right)^2-{p'(T)} \left(u^{r}\right)^2} \\
(A_r)_{32} &= \frac{{p''(T)} {n} {T}}{{p''(T)} {T} \left(u^{\tau}\right)^2-{p'(T)} \left(u^{r}\right)^2} \\
(A_r)_{33} &= {u^{r}} \\
S_{3} &= \frac{{p''(T)} {n} {T} \left(u^{\tau}\right)^2 ({r} {u^{\tau}}+{\tau} {u^{r}})}{{r} {\tau} \left({p''(T)} {T} \left(u^{\tau}\right)^2-{p'(T)} \left(u^{r}\right)^2\right)}
\end{align*}
\subsubsection{\texttt{one\_\allowbreak{}d\_\allowbreak{}ideal\_\allowbreak{}matrix\_\allowbreak{}ruwen5} in \texttt{src/Matrix/1d\_\allowbreak{}ideal\_\allowbreak{}HQ\_\allowbreak{}cilindrical\_\allowbreak{}fugacity.jl}}

Derived, charm rows:

\begin{align*}
(A_t)_{31} &= {(\partial_{T} n)} {u^{\tau}} \\
(A_t)_{32} &= \frac{{n} {u^{r}}}{{u^{\tau}}} \\
(A_t)_{33} &= {(\partial_{\alpha} n)} {u^{\tau}} \\
(A_r)_{31} &= {(\partial_{T} n)} {u^{r}} \\
(A_r)_{32} &= {n} \\
(A_r)_{33} &= {(\partial_{\alpha} n)} {u^{r}} \\
S_{3} &= {n} \left(\frac{{u^{r}}}{{r}}+\frac{{u^{\tau}}}{{\tau}}\right)
\end{align*}
With $c$ the charm rows and $m$ the medium rows, $\mathrm{row}_{c}\to\mathrm{row}_{c} - L\,\mathrm{row}_{m}$, $L = A_t^{cm}(A_t^{mm})^{-1}$, removes $\partial_{\tau}$ of the medium from the charm rows:

\begin{align*}
L_{31} &= \frac{{u^{\tau}} \left({p''(T)} {n} \left(u^{r}\right)^2 {T}-{(\partial_{T} n)} {p'(T)} \left(2 \left(u^{r}\right)^2+1\right) {T}+{p'(T)} {n} \left(u^{r}\right)^2\right)}{{p'(T)} {T} \left({p'(T)} \left(u^{r}\right)^2-{p''(T)} {T} \left(u^{\tau}\right)^2\right)} \\
L_{32} &= -\frac{{u^{r}} \left({p''(T)} {n} {T} \left(u^{\tau}\right)^2-2 {(\partial_{T} n)} {p'(T)} {T} \left(u^{\tau}\right)^2+{p'(T)} {n} \left(u^{r}\right)^2\right)}{{p'(T)} {T} \left({p'(T)} \left(u^{r}\right)^2-{p''(T)} {T} \left(u^{\tau}\right)^2\right)}
\end{align*}
The file, charm rows (M-V7e, M-G7e, M-I3ne, M-I3ae):

\begin{align*}
(A_t)_{33} &= {(\partial_{\alpha} n)} {u^{\tau}} \\
(A_r)_{31} &= \frac{{u^{r}} ({p''(T)} {n}-{(\partial_{T} n)} {p'(T)})}{{p'(T)} \left(u^{r}\right)^2-{p''(T)} {T} \left(u^{\tau}\right)^2} \\
(A_r)_{32} &= \frac{{T} ({(\partial_{T} n)} {p'(T)}-{p''(T)} {n})}{{p'(T)} \left(u^{r}\right)^2-{p''(T)} {T} \left(u^{\tau}\right)^2} \\
(A_r)_{33} &= {(\partial_{\alpha} n)} {u^{r}} \\
S_{3} &= \frac{{T} \left(u^{\tau}\right)^2 ({(\partial_{T} n)} {p'(T)}-{p''(T)} {n}) ({r} {u^{\tau}}+{\tau} {u^{r}})}{{r} {\tau} \left({p'(T)} \left(u^{r}\right)^2-{p''(T)} {T} \left(u^{\tau}\right)^2\right)}
\end{align*}
\subsection{Checks}

{\small
\begin{longtable}{@{}>{\raggedright\arraybackslash}p{0.27\textwidth}>{\raggedright\arraybackslash}p{0.60\textwidth}l@{}}
\toprule
id & what it tests & \\
\midrule
\endhead
\texttt{M-\allowbreak{}G0} & cylindrical Milne: $\nabla_{\lambda}g_{\mu\nu}$ =\allowbreak{} 0 & PASS \\
\texttt{M-\allowbreak{}G1} & $\nabla_{\mu}V^{\mu} = \partial_{\tau}X + X/\tau$ for $V^{\mu} = (X(\tau), 0, 0, 0)$: the Christoffels are in & PASS \\
\texttt{M-\allowbreak{}G2} & $u\cdot u = -1$,\allowbreak{} pi traceless and transverse & PASS \\
\texttt{M-\allowbreak{}V5q} & quasi-linear & PASS \\
\texttt{M-\allowbreak{}V5} & derived is + $\delta_{\pi\pi}\pi\theta$ =\allowbreak{} the shipped 5-field $A_{t}, A_{r}, S$,\allowbreak{} raw entry by entry,\allowbreak{} read from the .jl,\allowbreak{} symbolic $\delta_{\pi\pi}$ & PASS \\
\texttt{M-\allowbreak{}V5-\allowbreak{}0} & FLUIDUM\_\allowbreak{}DELTA\_\allowbreak{}PIPI=\allowbreak{}0: $\delta_{\pi\pi} = 0$,\allowbreak{} plain is =\allowbreak{}=\allowbreak{} the shipped is part & PASS \\
\texttt{M-\allowbreak{}V5-\allowbreak{}1} & FLUIDUM\_\allowbreak{}DELTA\_\allowbreak{}PIPI=\allowbreak{}1 (default): $\delta_{\pi\pi} = \tfrac{4}{3}\tau_{\pi}$ & PASS \\
\texttt{M-\allowbreak{}BGq} & quasi-linear & PASS \\
\texttt{M-\allowbreak{}BG} & derived first moment$\times (u^{\tau})^{3} r\tau$ =\allowbreak{} 1d\_\allowbreak{}viscous\_\allowbreak{}HQ\_\allowbreak{}cilindrical\_\allowbreak{}BG.jl,\allowbreak{} raw,\allowbreak{} read from the .jl & PASS \\
\texttt{M-\allowbreak{}V7m-\allowbreak{}dpp} & the medium carries $\delta_{\pi\pi} = \tfrac{4}{3}\tau_{\pi}$ (DNMR $\delta_{\pi\pi}$) & PASS \\
\texttt{M-\allowbreak{}V7m} & medium rows/\allowbreak{}cols 1-5 =\allowbreak{} IS + $\tfrac{4}{3}\tau_{\pi}\pi\theta$,\allowbreak{} raw entry by entry & PASS \\
\texttt{M-\allowbreak{}V7m-\allowbreak{}ctl} & control: plain is ($\delta_{\pi\pi} = 0$) does not reproduce the 7-field medium block & PASS \\
\texttt{M-\allowbreak{}V5=V7m} & at the default $\delta_{\pi\pi} = \tfrac{4}{3}\tau_{\pi}$ the 5-field kernel is the 7-field :MIS medium block,\allowbreak{} raw entry by entry (both read from the .jl) & PASS \\
\texttt{M-\allowbreak{}V7q} & quasi-linear & PASS \\
\texttt{M-\allowbreak{}V7c} & $A_{t}^{-1}A_{r}$ and $A_{t}^{-1}S$ =\allowbreak{} derived,\allowbreak{} exactly (rational arithmetic),\allowbreak{} 8 random states & PASS \\
\texttt{M-\allowbreak{}V7e} & charm rows 6-7 of the .jl =\allowbreak{} derived rows with $\partial_{\tau}T, \partial_{\tau}u^{r}$ eliminated with rows 1-5,\allowbreak{} raw,\allowbreak{} exact,\allowbreak{} 8 random states & PASS \\
\texttt{M-\allowbreak{}V7e-\allowbreak{}ctl} & control: without the elimination the charm rows of $A_{r}$ differ & PASS \\
\texttt{M-\allowbreak{}G7} & $n = \gamma n_{0}(T)$,\allowbreak{} $\kappa$0 d gamma: solved form =\allowbreak{} derived exactly,\allowbreak{} 8 random states & PASS \\
\texttt{M-\allowbreak{}G7e} & charm rows 6-7 of the .jl =\allowbreak{} derived rows with $\partial_{\tau}T, \partial_{\tau}u^{r}$ eliminated,\allowbreak{} raw,\allowbreak{} exact,\allowbreak{} 8 random states & PASS \\
\texttt{M-\allowbreak{}I3n} & ruwen4 $(T, u^{r}, n)$: solved form =\allowbreak{} derived exactly & PASS \\
\texttt{M-\allowbreak{}I3ne} & charm row 3 of the .jl =\allowbreak{} derived row with $\partial_{\tau}T, \partial_{\tau}u^{r}$ eliminated,\allowbreak{} raw,\allowbreak{} exact,\allowbreak{} 8 random states & PASS \\
\texttt{M-\allowbreak{}I3a} & ruwen5 $(T, u^{r}, \alpha)$: solved form =\allowbreak{} derived exactly & PASS \\
\texttt{M-\allowbreak{}I3ae} & charm row 3 of the .jl =\allowbreak{} derived row with $\partial_{\tau}T, \partial_{\tau}u^{r}$ eliminated,\allowbreak{} raw,\allowbreak{} exact,\allowbreak{} 8 random states & PASS \\
\bottomrule
\end{longtable}
}
23 passed, 0 failed (23 checks evaluated).

Re-run with \texttt{wolframscript -file deriv/derive\_\allowbreak{}1d\_\allowbreak{}matrices.wls}.


\clearpage
% GENERATED by deriv/render.wl from the derivation script.  Do not edit.

\section{The matrices no production path relies on}

Script: \texttt{deriv/derive\_upstream.wls}

Writes nothing into \texttt{src/Matrix}.

\subsection{\texttt{two\_d\_viscous\_matrix} (\texttt{legacy = true})}

Against the first section's matrix at $\delta_{\pi\pi} = 0$, $\phi = (T, u^{x}, u^{y}, \pi^{yy}, \pi^{\eta}{}_{\eta}, \pi^{xy}, \Pi)$. The source and 137 entries agree. The ten that differ, with $\delta$ = legacy $-$ derived:

\begin{align*}
(A_t)_{43}^{\mathrm{legacy}} &= \frac{3 {\pi^{xy}} \left(2 {\tau_{\pi}} {u^{x}} \left(u^{y}\right)^2+{u^{x}}\right)+{\pi^{yy}} {u^{y}} \left(3-6 {\tau_{\pi}} \left(\left(u^{x}\right)^2+1\right)\right)+2 {u^{y}} {\eta} \left(3 \left(u^{x}\right)^2+2 \left(u^{y}\right)^2+2\right)}{3 {u^{\tau}}} \\
(A_t)_{43}^{\mathrm{derived}} &= \frac{2 {u^{y}} \left(3 {\pi^{xy}} {\tau_{\pi}} {u^{x}} {u^{y}}-3 {\pi^{yy}} {\tau_{\pi}} \left(\left(u^{x}\right)^2+1\right)+{\eta} \left(3 \left(u^{x}\right)^2+2 \left(u^{y}\right)^2+2\right)\right)}{3 {u^{\tau}}} \\
\delta(A_t)_{43} &= \frac{{\pi^{xy}} {u^{x}}+{\pi^{yy}} {u^{y}}}{{u^{\tau}}} \\
(A_t)_{62}^{\mathrm{legacy}} &= \frac{1}{6 {u^{\tau}}}\Bigl(- 3 {\pi^{xy}} {u^{x}} \left({\tau_{\pi}} \left(4 \left(u^{y}\right)^2+2\right)-1\right)\\
  &\quad + {u^{y}} \left(6 {\pi^{\eta}{}_{\eta}} {\tau_{\pi}} \left(u^{\tau}\right)^2+3 {\pi^{yy}} \left({\tau_{\pi}} \left(4 \left(u^{x}\right)^2+2\right)+1\right)-4 \left(u^{x}\right)^2 {\eta}\right)\Bigr) \\
(A_t)_{62}^{\mathrm{derived}} &= \frac{1}{3 {u^{\tau}}}\Bigl(- 3 {\pi^{xy}} {\tau_{\pi}} {u^{x}} \left(2 \left(u^{y}\right)^2+1\right)\\
  &\quad + {u^{y}} \left(3 {\pi^{\eta}{}_{\eta}} {\tau_{\pi}} \left(u^{\tau}\right)^2+3 {\pi^{yy}} \left(2 {\tau_{\pi}} \left(u^{x}\right)^2+{\tau_{\pi}}\right)+{\eta} \left(\left(u^{x}\right)^2+3 \left(u^{y}\right)^2+3\right)\right)\Bigr) \\
\delta(A_t)_{62} &= \frac{{\pi^{xy}} {u^{x}}+{u^{y}} \left({\pi^{yy}}-2 \left(u^{\tau}\right)^2 {\eta}\right)}{2 {u^{\tau}}} \\
(A_t)_{63}^{\mathrm{legacy}} &= \frac{1}{6 {u^{\tau}} \left(\left(u^{y}\right)^2+1\right)}\Bigl(- 3 {\pi^{xy}} (2 {\tau_{\pi}}-1) {u^{y}} \left(\left(u^{y}\right)^2+1\right) + 6 {\pi^{xy}} \left(u^{x}\right)^2 {u^{y}} \left(2 {\tau_{\pi}} \left(u^{y}\right)^2+1\right)\\
  &\quad - 3\left(u^{x}\right)^3\Bigl({\pi^{\eta}{}_{\eta}} + {\pi^{yy}} + 2 {\pi^{yy}} {\tau_{\pi}} + 2 {\tau_{\pi}} \left(u^{y}\right)^2 ({\pi^{\eta}{}_{\eta}}+2 {\pi^{yy}})\\
  &\quad - 4 \left(\left(u^{y}\right)^2+1\right) {\eta}\Bigr)\\
  &\quad + {u^{x}}\Bigl(- 3 ({\pi^{\eta}{}_{\eta}}+2 {\pi^{yy}} {\tau_{\pi}}+{\pi^{yy}})\\
  &\quad - 3 \left(u^{y}\right)^2 \left(2 {\pi^{\eta}{}_{\eta}} {\tau_{\pi}} \left(\left(u^{y}\right)^2+1\right)+{\pi^{\eta}{}_{\eta}}+4 {\pi^{yy}} {\tau_{\pi}}\right)\\
  &\quad + 4 \left(2 \left(u^{y}\right)^4+5 \left(u^{y}\right)^2+3\right) {\eta}\Bigr)\Bigr) \\
(A_t)_{63}^{\mathrm{derived}} &= \frac{1}{3 {u^{\tau}} \left(\left(u^{y}\right)^2+1\right)}\Bigl(- 3 {\pi^{yy}} {\tau_{\pi}} {u^{x}} \left(\left(u^{x}\right)^2+1\right) \left(2 \left(u^{y}\right)^2+1\right)\\
  &\quad + 3 {\pi^{xy}} {\tau_{\pi}} {u^{y}} \left(\left(2 \left(u^{x}\right)^2-1\right) \left(u^{y}\right)^2-1\right)\\
  &\quad + {u^{x}} \left(\left(\left(u^{y}\right)^2+1\right) {\eta} \left(3 \left(u^{x}\right)^2+\left(u^{y}\right)^2+3\right)-3 {\pi^{\eta}{}_{\eta}} {\tau_{\pi}} \left(u^{\tau}\right)^2 \left(u^{y}\right)^2\right)\Bigr) \\
\delta(A_t)_{63} &= \frac{1}{2 {u^{\tau}} \left(\left(u^{y}\right)^2+1\right)}\Bigl(- {\pi^{\eta}{}_{\eta}} \left(u^{\tau}\right)^2 {u^{x}} - {\pi^{yy}} \left(\left(u^{x}\right)^3+{u^{x}}\right) + {\pi^{xy}} {u^{y}} + 2 {\pi^{xy}} \left(u^{x}\right)^2 {u^{y}}\\
  &\quad + {\pi^{xy}} \left(u^{y}\right)^3 + 2 {u^{x}} {\eta} + 2 \left(u^{x}\right)^3 {\eta} + 4 {u^{x}} \left(u^{y}\right)^2 {\eta} + 2 \left(u^{x}\right)^3 \left(u^{y}\right)^2 {\eta}\\
  &\quad + 2 {u^{x}} \left(u^{y}\right)^4 {\eta}\Bigr) \\
(A_x)_{43}^{\mathrm{legacy}} &= \frac{2 {u^{x}} {u^{y}} \left({\pi^{xy}} {\tau_{\pi}} {u^{x}} {u^{y}}-{\pi^{yy}} {\tau_{\pi}} \left(\left(u^{x}\right)^2+1\right)+\left(u^{\tau}\right)^2 {\eta}\right)}{\left(u^{\tau}\right)^2}+{\pi^{xy}} \\
(A_x)_{43}^{\mathrm{derived}} &= \frac{2 {u^{x}} {u^{y}} \left({\pi^{xy}} {\tau_{\pi}} {u^{x}} {u^{y}}-{\pi^{yy}} {\tau_{\pi}} \left(\left(u^{x}\right)^2+1\right)+\left(u^{\tau}\right)^2 {\eta}\right)}{\left(u^{\tau}\right)^2} \\
\delta(A_x)_{43} &= {\pi^{xy}} \\
(A_x)_{62}^{\mathrm{legacy}} &= \frac{1}{6 \left(u^{\tau}\right)^2}{u^{x}}\Bigl(3 {\pi^{xy}} {u^{x}} + 3 {\pi^{yy}} {u^{y}} \left({\tau_{\pi}} \left(4 \left(u^{x}\right)^2+2\right)+1\right) - 6 {\pi^{xy}} {\tau_{\pi}} {u^{x}} \left(2 \left(u^{y}\right)^2+1\right)\\
  &\quad + 2 \left(u^{\tau}\right)^2 {u^{y}} (3 {\pi^{\eta}{}_{\eta}} {\tau_{\pi}}-2 {\eta})\Bigr) \\
(A_x)_{62}^{\mathrm{derived}} &= \frac{{u^{x}} \left({u^{y}} \left(3 {\pi^{\eta}{}_{\eta}} {\tau_{\pi}} \left(u^{\tau}\right)^2+3 {\pi^{yy}} \left(2 {\tau_{\pi}} \left(u^{x}\right)^2+{\tau_{\pi}}\right)+\left(u^{\tau}\right)^2 {\eta}\right)-3 {\pi^{xy}} {\tau_{\pi}} {u^{x}} \left(2 \left(u^{y}\right)^2+1\right)\right)}{3 \left(u^{\tau}\right)^2} \\
\delta(A_x)_{62} &= \frac{{u^{x}} \left({\pi^{xy}} {u^{x}}+{u^{y}} \left({\pi^{yy}}-2 \left(u^{\tau}\right)^2 {\eta}\right)\right)}{2 \left(u^{\tau}\right)^2} \\
(A_x)_{63}^{\mathrm{legacy}} &= \frac{1}{2 \left(u^{\tau}\right)^2 \left(\left(u^{y}\right)^2+1\right)}\Bigl(- {\pi^{xy}} (2 {\tau_{\pi}}-3) {u^{x}} {u^{y}} \left(\left(u^{y}\right)^2+1\right)\\
  &\quad - \left(\left(u^{y}\right)^2+1\right) \left({\pi^{\eta}{}_{\eta}} \left(u^{y}\right)^2+{\pi^{\eta}{}_{\eta}}+2 {\pi^{yy}}\right)\\
  &\quad + 2 {\pi^{xy}} \left(u^{x}\right)^3 {u^{y}} \left(2 {\tau_{\pi}} \left(u^{y}\right)^2+1\right)\\
  &\quad - \left(u^{x}\right)^4 \left(2 {\pi^{\eta}{}_{\eta}} {\tau_{\pi}} \left(u^{y}\right)^2+{\pi^{\eta}{}_{\eta}}+{\pi^{yy}} {\tau_{\pi}} \left(4 \left(u^{y}\right)^2+2\right)+{\pi^{yy}}\right)\\
  &\quad - \left(u^{x}\right)^2\Bigl(2 {\pi^{\eta}{}_{\eta}} \left(\left(u^{y}\right)^2+1\right) \left({\tau_{\pi}} \left(u^{y}\right)^2+1\right)\\
  &\quad + {\pi^{yy}} \left(2 (2 {\tau_{\pi}}+1) \left(u^{y}\right)^2+2 {\tau_{\pi}}+3\right)\Bigr)\\
  &\quad + 4 \left(u^{\tau}\right)^2 \left(\left(u^{x}\right)^2+1\right) \left(\left(u^{y}\right)^2+1\right) {\eta}\Bigr) \\
(A_x)_{63}^{\mathrm{derived}} &= \frac{1}{\left(u^{\tau}\right)^2 \left(\left(u^{y}\right)^2+1\right)}\Bigl(- {\pi^{yy}} {\tau_{\pi}} \left(u^{x}\right)^2 \left(\left(u^{x}\right)^2+1\right) \left(2 \left(u^{y}\right)^2+1\right)\\
  &\quad + {\pi^{xy}} {\tau_{\pi}} {u^{x}} {u^{y}} \left(\left(2 \left(u^{x}\right)^2-1\right) \left(u^{y}\right)^2-1\right)\\
  &\quad - \left(u^{\tau}\right)^2 \left({\pi^{\eta}{}_{\eta}} {\tau_{\pi}} \left(u^{x}\right)^2 \left(u^{y}\right)^2-\left(\left(u^{x}\right)^2+1\right) \left(\left(u^{y}\right)^2+1\right) {\eta}\right)\Bigr) \\
\delta(A_x)_{63} &= \frac{1}{2 \left(u^{\tau}\right)^2 \left(\left(u^{y}\right)^2+1\right)}\Bigl(- {\pi^{\eta}{}_{\eta}} \left(u^{\tau}\right)^4 + 3 {\pi^{xy}} {u^{x}} {u^{y}} + 2 {\pi^{xy}} \left(u^{x}\right)^3 {u^{y}} + 3 {\pi^{xy}} {u^{x}} \left(u^{y}\right)^3\\
  &\quad - {\pi^{yy}} \left(\left(u^{x}\right)^2+1\right) \left(\left(u^{x}\right)^2+2 \left(u^{y}\right)^2+2\right) + 2 {\eta} + 4 \left(u^{x}\right)^2 {\eta}\\
  &\quad + 2 \left(u^{x}\right)^4 {\eta} + 4 \left(u^{y}\right)^2 {\eta} + 6 \left(u^{x}\right)^2 \left(u^{y}\right)^2 {\eta} + 2 \left(u^{x}\right)^4 \left(u^{y}\right)^2 {\eta}\\
  &\quad + 2 \left(u^{y}\right)^4 {\eta} + 2 \left(u^{x}\right)^2 \left(u^{y}\right)^4 {\eta}\Bigr) \\
(A_y)_{42}^{\mathrm{legacy}} &= -\frac{\left(2 {\tau_{\pi}} \left(u^{y}\right)^2+1\right) \left({\pi^{xy}} \left(u^{y}\right)^2+{\pi^{xy}}-{\pi^{yy}} {u^{x}} {u^{y}}\right)}{\left(u^{\tau}\right)^2} \\
(A_y)_{42}^{\mathrm{derived}} &= -\frac{2 {\tau_{\pi}} \left(u^{y}\right)^2 \left({\pi^{xy}} \left(u^{y}\right)^2+{\pi^{xy}}-{\pi^{yy}} {u^{x}} {u^{y}}\right)}{\left(u^{\tau}\right)^2} \\
\delta(A_y)_{42} &= \frac{{\pi^{yy}} {u^{x}} {u^{y}}-{\pi^{xy}} \left(\left(u^{y}\right)^2+1\right)}{\left(u^{\tau}\right)^2} \\
(A_y)_{43}^{\mathrm{legacy}} &= \frac{{u^{y}} \left({\pi^{xy}} \left(2 {\tau_{\pi}} {u^{x}} \left(u^{y}\right)^2+{u^{x}}\right)+{\pi^{yy}} {u^{y}} \left(1-2 {\tau_{\pi}} \left(\left(u^{x}\right)^2+1\right)\right)\right)}{\left(u^{\tau}\right)^2}+\frac{4}{3} \left(\left(u^{y}\right)^2+1\right) {\eta} \\
(A_y)_{43}^{\mathrm{derived}} &= \frac{6 {\pi^{xy}} {\tau_{\pi}} {u^{x}} \left(u^{y}\right)^3-6 {\pi^{yy}} {\tau_{\pi}} \left(\left(u^{x}\right)^2+1\right) \left(u^{y}\right)^2+4 \left(u^{\tau}\right)^2 \left(\left(u^{y}\right)^2+1\right) {\eta}}{3 \left(u^{\tau}\right)^2} \\
\delta(A_y)_{43} &= \frac{{u^{y}} ({\pi^{xy}} {u^{x}}+{\pi^{yy}} {u^{y}})}{\left(u^{\tau}\right)^2} \\
(A_y)_{62}^{\mathrm{legacy}} &= \frac{1}{2 \left(u^{\tau}\right)^2}\Bigl({\pi^{\eta}{}_{\eta}} + 2 {\pi^{yy}} + \left(u^{x}\right)^2 ({\pi^{\eta}{}_{\eta}}+2 {\pi^{yy}}) \left(2 {\tau_{\pi}} \left(u^{y}\right)^2+1\right)\\
  &\quad + \left(u^{y}\right)^2 \left(2 {\pi^{\eta}{}_{\eta}} {\tau_{\pi}} \left(\left(u^{y}\right)^2+1\right)+{\pi^{\eta}{}_{\eta}}+2 {\pi^{yy}} {\tau_{\pi}}+{\pi^{yy}}\right)\\
  &\quad - {\pi^{xy}} {u^{x}} {u^{y}} \left({\tau_{\pi}} \left(4 \left(u^{y}\right)^2+2\right)+1\right)\Bigr) \\
(A_y)_{62}^{\mathrm{derived}} &= \frac{1}{\left(u^{\tau}\right)^2}\Bigl(- {\pi^{xy}} {\tau_{\pi}} {u^{x}} {u^{y}} \left(2 \left(u^{y}\right)^2+1\right) + {\eta} + \left(u^{y}\right)^4 ({\pi^{\eta}{}_{\eta}} {\tau_{\pi}}+{\eta})\\
  &\quad + \left(u^{y}\right)^2 ({\pi^{\eta}{}_{\eta}} {\tau_{\pi}}+{\pi^{yy}} {\tau_{\pi}}+2 {\eta}) + \left(u^{x}\right)^2 \left(\left(u^{y}\right)^2 ({\pi^{\eta}{}_{\eta}} {\tau_{\pi}}+2 {\pi^{yy}} {\tau_{\pi}}+{\eta})+{\eta}\right)\Bigr) \\
\delta(A_y)_{62} &= \frac{1}{2 \left(u^{\tau}\right)^2}\Bigl({\pi^{\eta}{}_{\eta}} \left(u^{\tau}\right)^2 - {\pi^{xy}} {u^{x}} {u^{y}} + {\pi^{yy}} \left(2 \left(u^{x}\right)^2+\left(u^{y}\right)^2+2\right) - 2 {\eta} - 2 \left(u^{x}\right)^2 {\eta}\\
  &\quad - 4 \left(u^{y}\right)^2 {\eta} - 2 \left(u^{x}\right)^2 \left(u^{y}\right)^2 {\eta} - 2 \left(u^{y}\right)^4 {\eta}\Bigr) \\
(A_y)_{63}^{\mathrm{legacy}} &= \frac{1}{6 \left(u^{\tau}\right)^2 \left(\left(u^{y}\right)^2+1\right)}{u^{y}}\Bigl(- 3 {\pi^{xy}} (2 {\tau_{\pi}}-1) {u^{y}} \left(\left(u^{y}\right)^2+1\right)\\
  &\quad + 6 {\pi^{xy}} \left(u^{x}\right)^2 {u^{y}} \left(2 {\tau_{\pi}} \left(u^{y}\right)^2+1\right)\\
  &\quad + \left(u^{x}\right)^3\Bigl(- 3 ({\pi^{\eta}{}_{\eta}}+2 {\pi^{yy}} {\tau_{\pi}}+{\pi^{yy}}) - 6 {\tau_{\pi}} \left(u^{y}\right)^2 ({\pi^{\eta}{}_{\eta}}+2 {\pi^{yy}})\\
  &\quad + 8 \left(\left(u^{y}\right)^2+1\right) {\eta}\Bigr)\\
  &\quad + {u^{x}}\Bigl(- 3\Bigl({\pi^{yy}} + {\pi^{yy}} {\tau_{\pi}} \left(4 \left(u^{y}\right)^2+2\right)\\
  &\quad + {\pi^{\eta}{}_{\eta}} \left(\left(u^{y}\right)^2+1\right) \left(2 {\tau_{\pi}} \left(u^{y}\right)^2+1\right)\Bigr)\\
  &\quad + 8 \left(\left(u^{y}\right)^2+1\right)^2 {\eta}\Bigr)\Bigr) \\
(A_y)_{63}^{\mathrm{derived}} &= -\frac{1}{3 \left(u^{\tau}\right)^2 \left(\left(u^{y}\right)^2+1\right)}{u^{y}}\Bigl(3 {\pi^{yy}} {\tau_{\pi}} {u^{x}} \left(\left(u^{x}\right)^2+1\right) \left(2 \left(u^{y}\right)^2+1\right)\\
  &\quad + 3 {\pi^{xy}} {\tau_{\pi}} {u^{y}} \left(\left(1-2 \left(u^{x}\right)^2\right) \left(u^{y}\right)^2+1\right)\\
  &\quad + \left(u^{\tau}\right)^2 {u^{x}} \left(3 {\pi^{\eta}{}_{\eta}} {\tau_{\pi}} \left(u^{y}\right)^2-\left(\left(u^{y}\right)^2+1\right) {\eta}\right)\Bigr) \\
\delta(A_y)_{63} &= \frac{1}{2 \left(u^{\tau}\right)^2 \left(\left(u^{y}\right)^2+1\right)}{u^{y}}\Bigl(- {\pi^{\eta}{}_{\eta}} \left(u^{\tau}\right)^2 {u^{x}} - {\pi^{yy}} \left(\left(u^{x}\right)^3+{u^{x}}\right) + {\pi^{xy}} {u^{y}} + 2 {\pi^{xy}} \left(u^{x}\right)^2 {u^{y}}\\
  &\quad + {\pi^{xy}} \left(u^{y}\right)^3 + 2 {u^{x}} {\eta} + 2 \left(u^{x}\right)^3 {\eta} + 4 {u^{x}} \left(u^{y}\right)^2 {\eta} + 2 \left(u^{x}\right)^3 \left(u^{y}\right)^2 {\eta}\\
  &\quad + 2 {u^{x}} \left(u^{y}\right)^4 {\eta}\Bigr)
\end{align*}
\subsection{\texttt{two\_d\_viscous\_HQ\_matrix}}

$\phi = (T, u^{x}, u^{y}, \pi^{yy}, \pi^{\eta}{}_{\eta}, \pi^{xy}, \Pi, \alpha, \nu^{x}, \nu^{y})$. Rows 1--7: the derived matrix at $\delta_{\pi\pi} = 0$. Rows 8--10:

\[
\nabla_{\mu}(n u^{\mu} + \nu^{\mu}) = 0,\qquad \tau_{n}\Delta^{\mu}{}_{\beta}D\nu^{\beta} + \nu^{\mu} + \kappa\Delta^{\mu\beta}\partial_{\beta}\alpha = 0,\qquad \kappa = D_{s}n
\]
The steps are those of the second section. With the medium as fields, in columns 1, 2, 3, 8, 9, 10:

\begin{align*}
(A_{t})_{8,1} &= {(\partial_{T}n)} {u^{\tau}} \\
(A_{t})_{8,2} &= \frac{{n} \left(u^{\tau}\right)^2 {u^{x}}+{\nu^{x}} \left(\left(u^{y}\right)^2+1\right)-{\nu^{y}} {u^{x}} {u^{y}}}{\left(u^{\tau}\right)^3} \\
(A_{t})_{8,3} &= \frac{{n} \left(u^{\tau}\right)^2 {u^{y}}-{\nu^{x}} {u^{x}} {u^{y}}+{\nu^{y}} \left(\left(u^{x}\right)^2+1\right)}{\left(u^{\tau}\right)^3} \\
(A_{t})_{8,8} &= {(\partial_{\alpha}n)} {u^{\tau}} \\
(A_{t})_{8,9} &= \frac{{u^{x}}}{{u^{\tau}}} \\
(A_{t})_{8,10} &= \frac{{u^{y}}}{{u^{\tau}}} \\
(A_{t})_{9,2} &= \frac{{\tau_{n}} {u^{x}} \left({\nu^{y}} {u^{x}} {u^{y}}-{\nu^{x}} \left(\left(u^{y}\right)^2+1\right)\right)}{{u^{\tau}}} \\
(A_{t})_{9,3} &= -\frac{{\tau_{n}} {u^{x}} \left(-{\nu^{x}} {u^{x}} {u^{y}}+{\nu^{y}} \left(u^{x}\right)^2+{\nu^{y}}\right)}{{u^{\tau}}} \\
(A_{t})_{9,8} &= {D_{s}} {n} {u^{\tau}} {u^{x}} \\
(A_{t})_{9,9} &= {\tau_{n}} {u^{\tau}} \\
(A_{t})_{10,2} &= -\frac{{\tau_{n}} {u^{y}} \left({\nu^{x}} \left(u^{y}\right)^2+{\nu^{x}}-{\nu^{y}} {u^{x}} {u^{y}}\right)}{{u^{\tau}}} \\
(A_{t})_{10,3} &= \frac{{\tau_{n}} {u^{y}} \left({\nu^{x}} {u^{x}} {u^{y}}-{\nu^{y}} \left(\left(u^{x}\right)^2+1\right)\right)}{{u^{\tau}}} \\
(A_{t})_{10,8} &= {D_{s}} {n} {u^{\tau}} {u^{y}} \\
(A_{t})_{10,10} &= {\tau_{n}} {u^{\tau}} \\
(A_{x})_{8,1} &= {(\partial_{T}n)} {u^{x}} \\
(A_{x})_{8,2} &= {n} \\
(A_{x})_{8,8} &= {(\partial_{\alpha}n)} {u^{x}} \\
(A_{x})_{8,9} &= 1 \\
(A_{x})_{9,2} &= \frac{{\tau_{n}} \left(u^{x}\right)^2 \left({\nu^{y}} {u^{x}} {u^{y}}-{\nu^{x}} \left(\left(u^{y}\right)^2+1\right)\right)}{\left(u^{\tau}\right)^2} \\
(A_{x})_{9,3} &= -\frac{{\tau_{n}} \left(u^{x}\right)^2 \left(-{\nu^{x}} {u^{x}} {u^{y}}+{\nu^{y}} \left(u^{x}\right)^2+{\nu^{y}}\right)}{\left(u^{\tau}\right)^2} \\
(A_{x})_{9,8} &= {D_{s}} {n} \left(\left(u^{x}\right)^2+1\right) \\
(A_{x})_{9,9} &= {\tau_{n}} {u^{x}} \\
(A_{x})_{10,2} &= \frac{{\tau_{n}} {u^{x}} {u^{y}} \left({\nu^{y}} {u^{x}} {u^{y}}-{\nu^{x}} \left(\left(u^{y}\right)^2+1\right)\right)}{\left(u^{\tau}\right)^2} \\
(A_{x})_{10,3} &= \frac{{\tau_{n}} {u^{x}} {u^{y}} \left({\nu^{x}} {u^{x}} {u^{y}}-{\nu^{y}} \left(\left(u^{x}\right)^2+1\right)\right)}{\left(u^{\tau}\right)^2} \\
(A_{x})_{10,8} &= {D_{s}} {n} {u^{x}} {u^{y}} \\
(A_{x})_{10,10} &= {\tau_{n}} {u^{x}} \\
(A_{y})_{8,1} &= {(\partial_{T}n)} {u^{y}} \\
(A_{y})_{8,3} &= {n} \\
(A_{y})_{8,8} &= {(\partial_{\alpha}n)} {u^{y}} \\
(A_{y})_{8,10} &= 1 \\
(A_{y})_{9,2} &= \frac{{\tau_{n}} {u^{x}} {u^{y}} \left({\nu^{y}} {u^{x}} {u^{y}}-{\nu^{x}} \left(\left(u^{y}\right)^2+1\right)\right)}{\left(u^{\tau}\right)^2} \\
(A_{y})_{9,3} &= -\frac{{\tau_{n}} {u^{x}} {u^{y}} \left(-{\nu^{x}} {u^{x}} {u^{y}}+{\nu^{y}} \left(u^{x}\right)^2+{\nu^{y}}\right)}{\left(u^{\tau}\right)^2} \\
(A_{y})_{9,8} &= {D_{s}} {n} {u^{x}} {u^{y}} \\
(A_{y})_{9,9} &= {\tau_{n}} {u^{y}} \\
(A_{y})_{10,2} &= -\frac{{\tau_{n}} \left(u^{y}\right)^2 \left({\nu^{x}} \left(u^{y}\right)^2+{\nu^{x}}-{\nu^{y}} {u^{x}} {u^{y}}\right)}{\left(u^{\tau}\right)^2} \\
(A_{y})_{10,3} &= \frac{{\tau_{n}} \left(u^{y}\right)^2 \left({\nu^{x}} {u^{x}} {u^{y}}-{\nu^{y}} \left(\left(u^{x}\right)^2+1\right)\right)}{\left(u^{\tau}\right)^2} \\
(A_{y})_{10,8} &= {D_{s}} {n} \left(\left(u^{y}\right)^2+1\right) \\
(A_{y})_{10,10} &= {\tau_{n}} {u^{y}} \\
S_{8} &= \frac{{n} \left(u^{\tau}\right)^2+{\nu^{x}} {u^{x}}+{\nu^{y}} {u^{y}}}{{\tau} {u^{\tau}}} \\
S_{9} &= {\nu^{x}} \\
S_{10} &= {\nu^{y}}
\end{align*}
Its wrapper is \texttt{matrix2d\_visc\_HQ!} (D-HQup).

\subsection{\texttt{one\_d\_ideal\_matrix\_ruwen7}}

\begin{gather*}
x^{\mu} = ({\tau}, {r}, {\phi}, {\eta}),\qquad g_{\mu\nu} = \mathrm{diag}(-1, 1, {r}^2, {\tau}^2),\qquad u^{\mu} = (u^{\tau}, u^{r}, 0, 0) \\
P_{\mathrm{tot}} = P(T) + P_{h}(T, \alpha),\qquad T^{\mu\nu} = (e + P_{\mathrm{tot}})u^{\mu}u^{\nu} + P_{\mathrm{tot}}g^{\mu\nu},\qquad \nabla_{\mu}T^{\mu\nu} = 0,\qquad \nabla_{\mu}(n u^{\mu}) = 0 \\
e + P_{\mathrm{tot}} = Ts + \mu n = T\partial_{T}P_{\mathrm{tot}}\bigl|_{\mu} + \mu\partial_{\mu}P_{\mathrm{tot}}\bigl|_{T} = T\partial_{T}P_{\mathrm{tot}}\bigl|_{\alpha},\qquad \mu = \alpha T
\end{gather*}
$\phi = (T, u^{r}, \alpha)$, the $T^{\mu\nu}$ as written, the same extraction. The file, raw (U-P7):

\begin{align*}
(A_t)_{11} &= \left(u^{r}\right)^2 ({p'(T)}+{T} ({p''(T)}+{(\partial^{2}_{T}P_{h})})+{(\partial_{T}P_{h})})+{T} ({p''(T)}+{(\partial^{2}_{T}P_{h})}) \\
(A_t)_{12} &= 2 {T} {u^{r}} ({p'(T)}+{(\partial_{T}P_{h})}) \\
(A_t)_{13} &= {(\partial_{T}\partial_{\alpha}P_{h})} {T} \left(u^{\tau}\right)^2-{(\partial_{\alpha}P_{h})} \\
(A_t)_{21} &= {u^{r}} {u^{\tau}} ({p'(T)}+{T} ({p''(T)}+{(\partial^{2}_{T}P_{h})})+{(\partial_{T}P_{h})}) \\
(A_t)_{22} &= \frac{{T} \left(2 \left(u^{r}\right)^2+1\right) ({p'(T)}+{(\partial_{T}P_{h})})}{{u^{\tau}}} \\
(A_t)_{23} &= {(\partial_{T}\partial_{\alpha}P_{h})} {T} {u^{r}} {u^{\tau}} \\
(A_t)_{31} &= {(\partial_{T}n)} {u^{\tau}} \\
(A_t)_{32} &= \frac{{n} {u^{r}}}{{u^{\tau}}} \\
(A_t)_{33} &= {(\partial_{\alpha}n)} {u^{\tau}} \\
(A_r)_{11} &= {u^{r}} {u^{\tau}} ({p'(T)}+{T} ({p''(T)}+{(\partial^{2}_{T}P_{h})})+{(\partial_{T}P_{h})}) \\
(A_r)_{12} &= \frac{{T} \left(2 \left(u^{r}\right)^2+1\right) ({p'(T)}+{(\partial_{T}P_{h})})}{{u^{\tau}}} \\
(A_r)_{13} &= {(\partial_{T}\partial_{\alpha}P_{h})} {T} {u^{r}} {u^{\tau}} \\
(A_r)_{21} &= {p'(T)} \left(u^{\tau}\right)^2+{T} \left(u^{r}\right)^2 ({p''(T)}+{(\partial^{2}_{T}P_{h})})+{(\partial_{T}P_{h})} \left(u^{\tau}\right)^2 \\
(A_r)_{22} &= 2 {T} {u^{r}} ({p'(T)}+{(\partial_{T}P_{h})}) \\
(A_r)_{23} &= {(\partial_{\alpha}P_{h})}+{(\partial_{T}\partial_{\alpha}P_{h})} {T} \left(u^{r}\right)^2 \\
(A_r)_{31} &= {(\partial_{T}n)} {u^{r}} \\
(A_r)_{32} &= {n} \\
(A_r)_{33} &= {(\partial_{\alpha}n)} {u^{r}} \\
S_{1} &= \frac{{T} {u^{\tau}} ({p'(T)}+{(\partial_{T}P_{h})}) ({r} {u^{\tau}}+{\tau} {u^{r}})}{{r} {\tau}} \\
S_{2} &= \frac{{T} {u^{r}} ({p'(T)}+{(\partial_{T}P_{h})}) \left({r} \left(u^{r}\right)^2+{r}+{\tau} {u^{r}} {u^{\tau}}\right)}{{r} {\tau} {u^{\tau}}} \\
S_{3} &= {n} \left(\frac{{u^{r}}}{{r}}+\frac{{u^{\tau}}}{{\tau}}\right)
\end{align*}
Its wrapper is \texttt{matrix1d\_ideal\_HQ\_br\_p\_fugacity!} (D-r7). For a Boltzmann gas of unit-charge species $P_{h} = nT$:

\[
\partial_{T}P_{h} = n + T\partial_{T}n,\quad \partial_{\alpha}P_{h} = T\partial_{\alpha}n,\quad \partial^{2}_{T}P_{h} = 2\partial_{T}n + T\partial^{2}_{T}n,\quad \partial_{T}\partial_{\alpha}P_{h} = \partial_{\alpha}n + T\partial_{T}\partial_{\alpha}n,\quad \partial^{2}_{\alpha}P_{h} = T\partial^{2}_{\alpha}n
\]
\subsection{\texttt{one\_d\_ideal\_matrix\_ruwen6}}

The same system in $\phi = (T, u^{r}, n)$, with $n = \partial_{\alpha}P_{h}/T$ and $\alpha = \alpha(T, n)$:

\[
\partial_{n}\alpha\bigl|_{T} = \frac{T}{\partial^{2}_{\alpha}P_{h}},\qquad \partial_{T}\alpha\bigl|_{n} = \frac{\partial_{\alpha}P_{h}/T - \partial_{T}\partial_{\alpha}P_{h}}{\partial^{2}_{\alpha}P_{h}},\qquad \partial_{k}\alpha = \partial_{T}\alpha\,\partial_{k}T + \partial_{n}\alpha\,\partial_{k}n
\]
Derived:

\begin{align*}
(A_t)_{11} &= - {p'(T)} - {(\partial_{T}P_{h})} - \frac{{(\partial_{\alpha}P_{h})} ({(\partial_{\alpha}P_{h})}-{(\partial_{T}\partial_{\alpha}P_{h})} {T})}{{(\partial^{2}_{\alpha}P_{h})} {T}} + \left(u^{\tau}\right)^2 ({p'(T)}+{(\partial_{T}P_{h})})\\
  &\quad + {T} \left(u^{\tau}\right)^2 \left({p''(T)}+\frac{{(\partial_{T}\partial_{\alpha}P_{h})} ({(\partial_{\alpha}P_{h})}-{(\partial_{T}\partial_{\alpha}P_{h})} {T})}{{(\partial^{2}_{\alpha}P_{h})} {T}}+{(\partial^{2}_{T}P_{h})}\right) \\
(A_t)_{12} &= 2 {T} {u^{r}} ({p'(T)}+{(\partial_{T}P_{h})}) \\
(A_t)_{13} &= \frac{{T} \left({(\partial_{T}\partial_{\alpha}P_{h})} {T} \left(u^{\tau}\right)^2-{(\partial_{\alpha}P_{h})}\right)}{{(\partial^{2}_{\alpha}P_{h})}} \\
(A_t)_{21} &= {u^{r}} {u^{\tau}} \left({p'(T)}+{T} ({p''(T)}+{(\partial^{2}_{T}P_{h})})+\frac{{(\partial_{T}\partial_{\alpha}P_{h})} ({(\partial_{\alpha}P_{h})}-{(\partial_{T}\partial_{\alpha}P_{h})} {T})}{{(\partial^{2}_{\alpha}P_{h})}}+{(\partial_{T}P_{h})}\right) \\
(A_t)_{22} &= \frac{{T} \left(2 \left(u^{r}\right)^2+1\right) ({p'(T)}+{(\partial_{T}P_{h})})}{{u^{\tau}}} \\
(A_t)_{23} &= \frac{{(\partial_{T}\partial_{\alpha}P_{h})} {T}^2 {u^{r}} {u^{\tau}}}{{(\partial^{2}_{\alpha}P_{h})}} \\
(A_t)_{32} &= \frac{{u^{r}} {n}}{{u^{\tau}}} \\
(A_t)_{33} &= {u^{\tau}} \\
(A_r)_{11} &= {u^{r}} {u^{\tau}} \left({p'(T)}+{T} ({p''(T)}+{(\partial^{2}_{T}P_{h})})+\frac{{(\partial_{T}\partial_{\alpha}P_{h})} ({(\partial_{\alpha}P_{h})}-{(\partial_{T}\partial_{\alpha}P_{h})} {T})}{{(\partial^{2}_{\alpha}P_{h})}}+{(\partial_{T}P_{h})}\right) \\
(A_r)_{12} &= \frac{{T} \left(2 \left(u^{r}\right)^2+1\right) ({p'(T)}+{(\partial_{T}P_{h})})}{{u^{\tau}}} \\
(A_r)_{13} &= \frac{{(\partial_{T}\partial_{\alpha}P_{h})} {T}^2 {u^{r}} {u^{\tau}}}{{(\partial^{2}_{\alpha}P_{h})}} \\
(A_r)_{21} &= {p'(T)} + {(\partial_{T}P_{h})} + \frac{{(\partial_{\alpha}P_{h})} ({(\partial_{\alpha}P_{h})}-{(\partial_{T}\partial_{\alpha}P_{h})} {T})}{{(\partial^{2}_{\alpha}P_{h})} {T}} + \left(u^{r}\right)^2 ({p'(T)}+{(\partial_{T}P_{h})})\\
  &\quad + {T} \left(u^{r}\right)^2 \left({p''(T)}+\frac{{(\partial_{T}\partial_{\alpha}P_{h})} ({(\partial_{\alpha}P_{h})}-{(\partial_{T}\partial_{\alpha}P_{h})} {T})}{{(\partial^{2}_{\alpha}P_{h})} {T}}+{(\partial^{2}_{T}P_{h})}\right) \\
(A_r)_{22} &= 2 {T} {u^{r}} ({p'(T)}+{(\partial_{T}P_{h})}) \\
(A_r)_{23} &= \frac{{T} \left({(\partial_{\alpha}P_{h})}+{(\partial_{T}\partial_{\alpha}P_{h})} {T} \left(u^{r}\right)^2\right)}{{(\partial^{2}_{\alpha}P_{h})}} \\
(A_r)_{32} &= {n} \\
(A_r)_{33} &= {u^{r}} \\
S_{1} &= \frac{{T} {u^{\tau}} ({p'(T)}+{(\partial_{T}P_{h})}) ({r} {u^{\tau}}+{\tau} {u^{r}})}{{r} {\tau}} \\
S_{2} &= \frac{{T} {u^{r}} ({p'(T)}+{(\partial_{T}P_{h})}) \left({r} \left(u^{r}\right)^2+{r}+{\tau} {u^{r}} {u^{\tau}}\right)}{{r} {\tau} {u^{\tau}}} \\
S_{3} &= {n} \left(\frac{{u^{r}}}{{r}}+\frac{{u^{\tau}}}{{\tau}}\right)
\end{align*}
The file agrees, raw (U-P6-cols, U-P6-row3, U-P6-c3). Its wrapper takes $P_{h} = nT$ from \texttt{charm\_pressure} and $\alpha$ from \texttt{charm\_alpha} (D-r6).

\subsection{Checks}

{\small
\begin{longtable}{@{}>{\raggedright\arraybackslash}p{0.27\textwidth}>{\raggedright\arraybackslash}p{0.60\textwidth}l@{}}
\toprule
id & what it tests & \\
\midrule
\endhead
\texttt{U-\allowbreak{}L-\allowbreak{}count} & exactly ten of the 147 entries of $A_{t}$,\allowbreak{} $A_{x}$,\allowbreak{} A\_\allowbreak{}y differ from the derived matrix ($\delta_{\pi\pi} = 0$) & PASS \\
\texttt{U-\allowbreak{}L-\allowbreak{}place} & they are in the velocity columns of the $\pi^{yy}$ and $\pi^{xy}$ rows & PASS \\
\texttt{U-\allowbreak{}L-\allowbreak{}src} & the source is right & PASS \\
\texttt{U-\allowbreak{}H-\allowbreak{}med} & its medium block =\allowbreak{} the derived matrix at $\delta_{\pi\pi} = 0$,\allowbreak{} raw (fixed 2026-10-30: it carried the ten legacy entries) & PASS \\
\texttt{U-\allowbreak{}H-\allowbreak{}pass} & the medium does not see charm & PASS \\
\texttt{U-\allowbreak{}H-\allowbreak{}charm} & rows 8-10 =\allowbreak{} N and Fick,\allowbreak{} $\tau_{n}\Delta D\nu + \nu + D_{s}n\,\Delta\nabla\alpha = 0$,\allowbreak{} raw & PASS \\
\texttt{U-\allowbreak{}H-\allowbreak{}charmsrc} & and their source,\allowbreak{} raw & PASS \\
\texttt{U-\allowbreak{}H-\allowbreak{}kappa} & the argument $D_{s}$ enters as $\kappa = D_{s}n$,\allowbreak{} which is $D_{s}$ $\partial_{\alpha}n$ for one unit-charge Boltzmann species & PASS \\
\texttt{U-\allowbreak{}P7} & ruwen7 =\allowbreak{} the derivation with $e + P_{\mathrm{tot}} = T\partial_{T}P_{\mathrm{tot}}$,\allowbreak{} raw,\allowbreak{} entry by entry & PASS \\
\texttt{U-\allowbreak{}P6-\allowbreak{}chain} & implicit function: $\partial_{T}\alpha$ at fixed n and $\partial_{n}\alpha$ at fixed T from $n = \partial_{\alpha}P_{h}/T$ & PASS \\
\texttt{U-\allowbreak{}P6-\allowbreak{}cols} & ruwen6,\allowbreak{} columns 1-2: =\allowbreak{} the derivation,\allowbreak{} raw & PASS \\
\texttt{U-\allowbreak{}P6-\allowbreak{}row3} & ruwen6,\allowbreak{} row 3 and the source: =\allowbreak{} the derivation,\allowbreak{} raw & PASS \\
\texttt{U-\allowbreak{}P6-\allowbreak{}c3} & ruwen6,\allowbreak{} column 3 of rows 1-2: =\allowbreak{} the derivation,\allowbreak{} raw (fixed 2026-10-30: $A_{t}$ and $A_{r}$ were exchanged there) & PASS \\
\texttt{U-\allowbreak{}P6-\allowbreak{}ctl} & control: $A_{t}$ and $A_{r}$ differ in that column,\allowbreak{} so the exchange would be seen & PASS \\
\bottomrule
\end{longtable}
}
14 passed, 0 failed (14 checks evaluated).

Re-run with \texttt{wolframscript -file deriv/derive\_\allowbreak{}upstream.wls}.


\clearpage
% GENERATED by deriv/render.wl from the derivation script.  Do not edit.

\section{From the matrices to the solver}

Script: \texttt{deriv/derive\_drivers.wls}, which runs \texttt{deriv/audit\_matrices.jl}

\subsection{The solved form}

\[
\partial_{\tau}\phi = -A_t^{-1}A_i\,\partial_{i}\phi - A_t^{-1}S
\]
Each driver evaluates the EOS and the transport at the state, calls a kernel, and returns $A_t^{-1}A_i$, $A_t^{-1}S$; the solver upwinds $A_t^{-1}A_i\partial_{i}\phi$. $\eta$, $\tau_{\pi}$, $\zeta$, $\tau_{\Pi}$ are inputs, not derived. $\delta_{\pi\pi} = \tfrac{4}{3}\tau_{\pi}$ unless switched off.

\subsection{\texorpdfstring{$\kappa$}{kappa} and \texorpdfstring{$\tau_{n}$}{taun}}

From the consistent first moment of the section on the two-way system, at rest ($u^{r} = 0$, $\nu^{r} = 0$ in the coefficients):

\begin{gather*}
\frac{D_{s}}{T}\Delta^{r}{}_{\lambda}\nabla_{\mu}T_{Q}^{\mu\lambda} + \nu^{r}\ \to\ \frac{D_{s}h}{T}\partial_{\tau}\nu^{r} + \nu^{r} + D_{s}\partial_{\alpha}n\,\partial_{r}\alpha + \frac{D_{s}}{T}(n + T\partial_{T}n)\partial_{r}T \\
\tau_{n} = \frac{D_{s}h}{T},\qquad \kappa = D_{s}\partial_{\alpha}n
\end{gather*}
The Fick closure keeps these two coefficients and drops the $\partial_{r}T$ term. In the code $\tau_{n}$ is written with a Bessel combination:

\begin{gather*}
K_{5} = K_{3} + \frac{8}{z}K_{4},\quad K_{4} = K_{2} + \frac{6}{z}K_{3},\quad K_{1} = K_{3} - \frac{4}{z}K_{2}\ \Rightarrow\ 2K_{1} - 3K_{3} + K_{5} = \frac{48}{z^{2}}K_{3} \\
\tau_{n} = \frac{\sum_{i}g_{i}q_{i}^{2}\,2\pi D_{s}T\,m_{i}^{5}\,e^{q_{i}\alpha}(2K_{1} - 3K_{3} + K_{5})/(192\pi^{3}T^{3})}{\sum_{i}q_{i}^{2}n_{i}} = \frac{D_{s}}{T}\,\frac{\sum_{i}q_{i}^{2}n_{i}h_{i}}{\sum_{i}q_{i}^{2}n_{i}} = \frac{D_{s}h}{T},\qquad h_{i} = m_{i}\frac{K_{3}(z_{i})}{K_{2}(z_{i})}
\end{gather*}
\subsection{The charm closures}

\[
f = g\,e^{q\alpha - E/T},\qquad dP = \frac{d^{3}p}{(2\pi)^{3}E},\qquad I_{nq} = \frac{1}{(2q + 1)!!}\int dP\,E^{n - 2q}p^{2q}f
\]
Angular averages, then the matching in the rest frame:

\begin{gather*}
p_{x}^{2}\to\frac{p^{2}}{3},\qquad p_{x}^{4}\to\frac{p^{4}}{5},\qquad p_{x}^{2}p_{y}^{2}\to\frac{p^{4}}{15} \\
\delta f_{\nu} = f\,\frac{\nu_{a}p^{a}}{P_{0}}:\quad \int dP\,p_{x}\delta f_{\nu} = \frac{I_{21}}{P_{0}}\nu^{x} = \nu^{x},\quad S^{\tau\tau x} = \frac{I_{41}}{P_{0}}\nu^{x} = \mathcal{A}\nu^{x},\quad S^{xxx} = 3\frac{I_{42}}{P_{0}}\nu^{x} = 3\mathcal{B}\nu^{x} \\
\delta f_{\pi} = f\,\frac{\pi_{ab}p^{a}p^{b}}{2I_{42}}:\quad \int dP\,p_{x}p_{y}\delta f_{\pi} = \pi^{xy},\quad S^{\tau xy} = \frac{I_{52}}{I_{42}}\pi^{xy} = \mathcal{C}\pi^{xy} \\
S^{\tau xx}\bigl|_{f} = I_{31},\qquad \delta M^{ab} = \pi_{Q}^{ab} + \Pi_{Q}\Delta^{ab}\ \text{with the same }\mathcal{C}\ \text{(closure)}
\end{gather*}
One Boltzmann species:

\[
n = \frac{g}{2\pi^{2}}m^{2}T\,K_{2}(z)e^{q\alpha},\ \ I_{21} = P_{0} = nT,\ \ I_{31} = P_{0}h,\ \ I_{41} = P_{0}(m^{2} + 5Th),\ \ I_{42} = P_{0}Th,\ \ I_{52} = I_{42}\,m\frac{K_{4}}{K_{3}}
\]
A list of species at one $(T, \alpha)$, with $w_{i} = q_{i}^{2}n_{i}$, as in \texttt{charm\_eos}:

\begin{gather*}
h = \frac{\sum w_{i}h_{i}}{\sum w_{i}},\quad \mathcal{A} = \frac{\sum w_{i}m_{i}^{2}}{\sum w_{i}} + 5Th,\quad \mathcal{B} = Th,\quad \mathcal{C} = \frac{\sum w_{i}h_{i}\mathcal{C}_{i}}{\sum w_{i}h_{i}},\quad I_{31} = nTh \\
\partial_{T}h_{i} = -\frac{m^{2}}{T^{2}}\frac{-\tfrac{1}{2}(K_{2} + K_{4})K_{2} + \tfrac{1}{2}K_{3}(K_{1} + K_{3})}{K_{2}^{2}},\quad \partial_{T}\mathcal{C}_{i} = -\frac{m^{2}}{T^{2}}\frac{-\tfrac{1}{2}(K_{3} + K_{5})K_{3} + \tfrac{1}{2}K_{4}(K_{2} + K_{4})}{K_{3}^{2}} \\
\partial_{T}n_{i} = \frac{g}{2\pi^{2}}e^{q\alpha}\frac{m^{2}(mK_{1} + 2TK_{2} + mK_{3})}{2T},\quad \partial_{\alpha}I_{31} = \partial_{\alpha}n\,Th,\quad \partial_{T}I_{31} = \partial_{T}n\,Th + nh + nT\partial_{T}h
\end{gather*}
A state of charge $|q| = 2$ enters with these weights; a charge-correct second moment would also carry $\partial_{\alpha}\mathcal{C}$, $\partial_{\alpha}h$.

\subsection{The vacuum regulators}

Two forms. Two-way drivers, $\epsilon$ = \texttt{dmn\_eps}:

\[
E_{\alpha}\to\nabla_{\mu}N^{\mu} + \epsilon\,u^{\mu}\partial_{\mu}\alpha:\qquad (A_t)_{\alpha\alpha}\to(A_t)_{\alpha\alpha} + \epsilon u^{\tau},\quad (A_i)_{\alpha\alpha}\to(A_i)_{\alpha\alpha} + \epsilon u^{i}
\]
Background and 7-field drivers, $\partial_{\alpha}n\to\partial_{\alpha}n + \epsilon$: the same term in the number row. In the consistent rows it also reaches $\kappa$, since $\kappa$ is formed inside the kernel:

\[
\kappa\to D_{s}(\partial_{\alpha}n + \epsilon)\quad\text{(consistent, second moment)},\qquad \kappa\ \text{unchanged}\quad\text{(Fick)}
\]
It is not conservative:

\[
\partial_{\tau}(\tau rN^{\tau}) + \partial_{r}(\tau rN^{r}) = -\tau r\,\epsilon\,u^{\mu}\partial_{\mu}\alpha
\]
\texttt{matrix1d\_visc\_gamma\_HQ!} accepts \texttt{dmn\_eps}, but its kernel has no $\partial_{\alpha}n$: the number row carries $n$, so the regulator does nothing there. $\alpha$ is clamped to $|\alpha|\le\alpha_{\max}$ where the EOS is read.

\subsection{The cold edge}

\begin{gather*}
x = \frac{T - (T_{\mathrm{edge}} - w_{T})}{w_{T}},\qquad s(x) = x^{3}(10 - 15x + 6x^{2}) \\
s(0) = s'(0) = s''(0) = s'(1) = s''(1) = 0,\quad s(1) = 1 \\
D_{s}\to f(T)D_{s},\qquad f = f_{0} + (1 - f_{0})s\ \text{in the window},\quad f = 1\ (T\ge T_{\mathrm{edge}}),\quad f = f_{0}\ (T\le T_{\mathrm{edge}} - w_{T}) \\
w = 1 - s = \frac{1 - f}{1 - f_{0}}
\end{gather*}
Only $D_{s}$ changes, and it is not differentiated, as in the matrices. Its gradient is dropped; this is a decoupling prescription, not a derived term.

\subsection{The corona update}

Below $T_{\mathrm{edge}}$ the quasi-linear $\dot\alpha$ is blended with one that gives a flux-form number update. $V = r$ in 1+1D, $V = 1$ in 2+1D; rates of $T$, $u$, $\nu$ are those the matrices gave:

\begin{gather*}
N^{\tau} = nu^{\tau} + \nu^{\tau},\qquad N^{i} = nu^{i} + \nu^{i},\qquad \nu^{\tau} = \frac{u_{i}\nu^{i}}{u^{\tau}} \\
\frac{d}{d\tau}(\tau VN^{\tau}) = VN^{\tau} + \tau V\bigl(\partial_{T}n\,\dot T\,u^{\tau} + \partial_{\alpha}n\,\dot\alpha\,u^{\tau} + n\dot u^{\tau} + \dot\nu^{\tau}\bigr) \equiv \mathcal{T} \\
\dot\alpha_{\mathrm{flux}} = \frac{(\mathcal{T} - VN^{\tau})/(\tau V) - \partial_{T}n\,\dot T\,u^{\tau} - n\dot u^{\tau} - \dot\nu^{\tau}}{u^{\tau}(\partial_{\alpha}n + \epsilon)} \\
\dot u^{\tau} = \frac{u_{i}\dot u^{i}}{u^{\tau}},\qquad \dot\nu^{\tau} = \frac{\dot u_{i}\nu^{i} + u_{i}\dot\nu^{i}}{u^{\tau}} - \frac{\nu^{\tau}\dot u^{\tau}}{u^{\tau}} \\
\mathcal{T} = -\frac{G_{i+1/2} - G_{i-1/2}}{\Delta r},\qquad G_{i+1/2} = \tau r_{i+1/2}F_{i+1/2},\qquad G_{1/2} = 0\ \text{on the axis}\quad\text{(1+1D)} \\
\mathcal{T} = -\tau\Bigl(\frac{F^{x}_{i+1/2,j} - F^{x}_{i-1/2,j}}{\Delta x} + \frac{F^{y}_{i,j+1/2} - F^{y}_{i,j-1/2}}{\Delta y}\Bigr)\quad\text{(2+1D)} \\
F_{i+1/2} = \frac{N^{i}_{L} + N^{i}_{R}}{2} - a\frac{N^{\tau}_{R} - N^{\tau}_{L}}{2},\qquad a = \min\Bigl(1,\ \max\Bigl(\frac{|N^{i}_{L}|}{N^{\tau}_{L}},\ \frac{|N^{i}_{R}|}{N^{\tau}_{R}}\Bigr)\Bigr) \\
\dot\alpha = (1 - w)\dot\alpha_{\mathrm{ql}} + w\,\dot\alpha_{\mathrm{flux}}
\end{gather*}
For $w = 1$ and $\epsilon = 0$ the sum over cells telescopes to the boundary fluxes:

\[
\sum_{i}\Delta r\,\frac{d}{d\tau}(\tau r_{i}N^{\tau}_{i}) = -(G_{\mathrm{end}} - G_{1/2})
\]
The upwind cap and the blend are numerical choices, not limits of the equations. The update is a discretisation of $E_{\alpha}$, so the matrices are unchanged.

\subsection{Checks}

R-: symbolic, in the script. A-: \texttt{audit\_matrices.jl} rebuilds each system numerically with ForwardDiff and compares the solved form at random states, with a control that multiplies one coefficient by 1.2. D-: it calls the shipped drivers at physical states and evaluates the corona update on a grid.

{\small
\begin{longtable}{@{}>{\raggedright\arraybackslash}p{0.27\textwidth}>{\raggedright\arraybackslash}p{0.60\textwidth}l@{}}
\toprule
id & what it tests & \\
\midrule
\endhead
\texttt{R-\allowbreak{}fick} & at rest the row is $\tau_{n}\partial_{\tau}\nu^{r} + \kappa\partial_{r}\alpha + (D_{s}/T)(n + T\partial_{T}n)\partial_{r}T$ & PASS \\
\texttt{R-\allowbreak{}taun} & the coefficient of $D\nu^{r}$ is $\tau_{n} = D_{s}h/T$ & PASS \\
\texttt{R-\allowbreak{}kappa} & the coefficient of $\partial_{r}\alpha$ at rest is $\kappa = D_{s}\partial_{\alpha}n$ & PASS \\
\texttt{R-\allowbreak{}bess} & $2K_{1} - 3K_{3} + K_{5} = 48K_{3}/z^{2}$,\allowbreak{} so the code's $\tau_{n}$ is $D_{s}h/T$ & PASS \\
\texttt{R-\allowbreak{}tausp} & tau\_\allowbreak{}diffusion\_\allowbreak{}hadron,\allowbreak{} one species: the prefactor of the code gives $(D_{s}/T)mK_{3}/K_{2}$ & PASS \\
\texttt{R-\allowbreak{}clos} & momentum integrals at $z = 3,\ 15/2,\ 19/2$: $n,\ I_{21} = P_{0},\ I_{31} = P_{0}h,\ I_{41} = P_{0}\mathcal{A},\ I_{42} = P_{0}\mathcal{B},\ I_{52} = I_{42}\mathcal{C}$ & PASS \\
\texttt{R-\allowbreak{}ang} & angular averages: $p_{x}^{2}\to p^{2}/3,\ p_{x}^{4}\to p^{4}/5,\ p_{x}^{2}p_{y}^{2}\to p^{4}/15$ & PASS \\
\texttt{R-\allowbreak{}match} & $\int dP\,p_{x}\delta f_{\nu} = \nu^{x},\ S^{\tau\tau x} = (I_{41}/P_{0})\nu^{x},\ S^{xxx} = 3(I_{42}/P_{0})\nu^{x},\ \int dP\,p_{x}p_{y}\delta f_{\pi} = \pi^{xy},\ S^{\tau xy} = (I_{52}/I_{42})\pi^{xy}$ & PASS \\
\texttt{R-\allowbreak{}dh} & $\partial_{T}h$ of charm\_\allowbreak{}eos & PASS \\
\texttt{R-\allowbreak{}dC} & $\partial_{T}\mathcal{C}$ of charm\_\allowbreak{}eos & PASS \\
\texttt{R-\allowbreak{}dn} & $\partial_{T}n$ of charm\_\allowbreak{}eos,\allowbreak{} per species & PASS \\
\texttt{R-\allowbreak{}I31} & $\partial_{T}I_{31}$ and $\partial_{\alpha}I_{31}$ of $I_{31} = nTh$ & PASS \\
\texttt{R-\allowbreak{}reg} & two-way: $\epsilon\,u\cdot\partial\alpha$ adds $\epsilon u^{\tau}, \epsilon u^{r}$ to the $\alpha$ column of the number row & PASS \\
\texttt{R-\allowbreak{}regbg} & background drivers: $\partial_{\alpha}n\to\partial_{\alpha}n + \epsilon$ in the number row is the same term & PASS \\
\texttt{R-\allowbreak{}regk} & and through $\kappa = D_{s}\partial_{\alpha}n$ it adds $D_{s}\epsilon\,\partial\alpha$ to the consistent current row; the two-way drivers avoid this & PASS \\
\texttt{R-\allowbreak{}leak} & $\tau rN$ is conserved by the number row: the regulated row gives $\partial_{\tau}(\tau rN^{\tau}) + \partial_{r}(\tau rN^{r}) = -\tau r\epsilon\,u\cdot\partial\alpha$ & PASS \\
\texttt{R-\allowbreak{}step} & $s(x) = x^{3}(10 - 15x + 6x^{2})$ is C2: s,\allowbreak{} s',\allowbreak{} s'' fixed at both ends & PASS \\
\texttt{R-\allowbreak{}stepw} & $f = f_{0} + (1 - f_{0})s,\ w = 1 - s = (1 - f)/(1 - f_{0})$ & PASS \\
\texttt{R-\allowbreak{}cor} & corona\_\allowbreak{}$\alpha$\_\allowbreak{}rate =\allowbreak{} the solution of $\partial_{\tau}(\tau VN^{\tau})$ =\allowbreak{} target,\allowbreak{} with its $u^{\tau}$ and n$u^{\tau}$ rates,\allowbreak{} symbolically & PASS \\
\texttt{R-\allowbreak{}tele} & the flux differences telescope & PASS \\
\texttt{R-\allowbreak{}llf} & the LLF flux is consistent,\allowbreak{} and antisymmetric under exchanging the two sides with a & PASS \\
\texttt{R-\allowbreak{}pos} & with a \textgreater{}=\allowbreak{} |N\textasciicircum{}r|/\allowbreak{}N\textasciicircum{}tau the flux is upwind where the drift equals a: the outflow of a vacuum cell vanishes & PASS \\
\texttt{R-\allowbreak{}audit} & audit\_\allowbreak{}matrices.jl ran and every check passed: 71 checks & PASS \\
\texttt{A-\allowbreak{}visc2d} & two\_\allowbreak{}d\_\allowbreak{}viscous\_\allowbreak{}matrix\_\allowbreak{}derived & PASS \\
\texttt{A-\allowbreak{}BGf2d} & two\_\allowbreak{}d\_\allowbreak{}HQ\_\allowbreak{}BG\_\allowbreak{}matrix & PASS \\
\texttt{A-\allowbreak{}live} & two\_\allowbreak{}d\_\allowbreak{}HQ\_\allowbreak{}live\_\allowbreak{}rows & PASS \\
\texttt{A-\allowbreak{}BGc1d} & one\_\allowbreak{}d\_\allowbreak{}HQ\_\allowbreak{}BG\_\allowbreak{}consistent\_\allowbreak{}matrix & PASS \\
\texttt{A-\allowbreak{}BGm1d} & one\_\allowbreak{}d\_\allowbreak{}HQ\_\allowbreak{}BG\_\allowbreak{}m2\_\allowbreak{}matrix & PASS \\
\texttt{A-\allowbreak{}BGc2d} & two\_\allowbreak{}d\_\allowbreak{}HQ\_\allowbreak{}BG\_\allowbreak{}consistent\_\allowbreak{}matrix & PASS \\
\texttt{A-\allowbreak{}BGm2d} & two\_\allowbreak{}d\_\allowbreak{}HQ\_\allowbreak{}BG\_\allowbreak{}m2\_\allowbreak{}matrix & PASS \\
\texttt{A-\allowbreak{}tw1d} & one\_\allowbreak{}d\_\allowbreak{}two\_\allowbreak{}way\_\allowbreak{}matrix & PASS \\
\texttt{A-\allowbreak{}tw1dc} & one\_\allowbreak{}d\_\allowbreak{}two\_\allowbreak{}way\_\allowbreak{}matrix\_\allowbreak{}conformal & PASS \\
\texttt{A-\allowbreak{}tw2d} & two\_\allowbreak{}d\_\allowbreak{}two\_\allowbreak{}way\_\allowbreak{}matrix & PASS \\
\texttt{A-\allowbreak{}tw2dc} & two\_\allowbreak{}d\_\allowbreak{}two\_\allowbreak{}way\_\allowbreak{}matrix\_\allowbreak{}conformal & PASS \\
\texttt{A-\allowbreak{}cm1d} & one\_\allowbreak{}d\_\allowbreak{}two\_\allowbreak{}way\_\allowbreak{}cM\_\allowbreak{}matrix & PASS \\
\texttt{A-\allowbreak{}cm1dc} & one\_\allowbreak{}d\_\allowbreak{}two\_\allowbreak{}way\_\allowbreak{}cM\_\allowbreak{}matrix\_\allowbreak{}conformal & PASS \\
\texttt{A-\allowbreak{}cm2d} & two\_\allowbreak{}d\_\allowbreak{}two\_\allowbreak{}way\_\allowbreak{}cM\_\allowbreak{}matrix & PASS \\
\texttt{A-\allowbreak{}cm2dc} & two\_\allowbreak{}d\_\allowbreak{}two\_\allowbreak{}way\_\allowbreak{}cM\_\allowbreak{}matrix\_\allowbreak{}conformal & PASS \\
\texttt{A-\allowbreak{}cm0} & control: the S5 kernel differs from the system with $\delta M$ in $T_{Q}$ & PASS \\
\texttt{A-\allowbreak{}visc1d} & one\_\allowbreak{}d\_\allowbreak{}viscous\_\allowbreak{}matrix & PASS \\
\texttt{A-\allowbreak{}MIS} & one\_\allowbreak{}d\_\allowbreak{}viscous\_\allowbreak{}matrix\_\allowbreak{}fugacity & PASS \\
\texttt{A-\allowbreak{}gamma} & one\_\allowbreak{}d\_\allowbreak{}viscous\_\allowbreak{}matrix\_\allowbreak{}gamma & PASS \\
\texttt{A-\allowbreak{}BGf1d} & one\_\allowbreak{}d\_\allowbreak{}viscous\_\allowbreak{}matrix\_\allowbreak{}fugacity\_\allowbreak{}BG\_\allowbreak{}only & PASS \\
\texttt{A-\allowbreak{}r4} & one\_\allowbreak{}d\_\allowbreak{}ideal\_\allowbreak{}matrix\_\allowbreak{}ruwen4 & PASS \\
\texttt{A-\allowbreak{}r5} & one\_\allowbreak{}d\_\allowbreak{}ideal\_\allowbreak{}matrix\_\allowbreak{}ruwen5 & PASS \\
\texttt{A-\allowbreak{}legacy} & two\_\allowbreak{}d\_\allowbreak{}viscous\_\allowbreak{}matrix differs from the equations (its ten wrong entries) & PASS \\
\texttt{A-\allowbreak{}HQ} & two\_\allowbreak{}d\_\allowbreak{}viscous\_\allowbreak{}HQ\_\allowbreak{}matrix: derived medium at $\delta_{\pi\pi} = 0$,\allowbreak{} Fick charm rows with $\kappa = D_{s}n$ & PASS \\
\texttt{A-\allowbreak{}r7} & one\_\allowbreak{}d\_\allowbreak{}ideal\_\allowbreak{}matrix\_\allowbreak{}ruwen7,\allowbreak{} $e + P = T\partial_{T}P_{\mathrm{tot}}$ at fixed $\alpha$ & PASS \\
\texttt{A-\allowbreak{}r7mu} & ruwen7 differs from $e + P = T\partial_{T}P_{\mathrm{tot}} + \alpha\partial_{\alpha}P_{\mathrm{tot}}$ (the wrong enthalpy) & PASS \\
\texttt{A-\allowbreak{}r6} & one\_\allowbreak{}d\_\allowbreak{}ideal\_\allowbreak{}matrix\_\allowbreak{}ruwen6 =\allowbreak{} $n = \partial_{\alpha}P_{h}/T$ & PASS \\
\texttt{A-\allowbreak{}r6s} & control: ruwen6 with column 3,\allowbreak{} rows 1-2 of $A_{t}$ and $A_{x}$ exchanged differs & PASS \\
\texttt{D-\allowbreak{}visc1} & matrix1d\_\allowbreak{}visc!: EOS,\allowbreak{} $\eta, \tau_{\pi}, \zeta, \tau_{\Pi}$,\allowbreak{} $\delta_{\pi\pi}$ =\allowbreak{} 4/\allowbreak{}3 tau\_\allowbreak{}pi into the kernel & PASS \\
\texttt{D-\allowbreak{}visc2} & matrix2d\_\allowbreak{}visc!: the same,\allowbreak{} derived matrix & PASS \\
\texttt{D-\allowbreak{}MIS} & matrix1d\_\allowbreak{}visc\_\allowbreak{}HQ!,\allowbreak{} $\epsilon = 0$: $\kappa, \tau_{n}$ from the list; the p slot holds the entropy and is unused & PASS \\
\texttt{D-\allowbreak{}MISe} & matrix1d\_\allowbreak{}visc\_\allowbreak{}HQ!,\allowbreak{} $\epsilon = 10^{-3}$: only $\partial_{\alpha}n$ of the number row,\allowbreak{} not $\kappa$ & PASS \\
\texttt{D-\allowbreak{}gamma} & matrix1d\_\allowbreak{}visc\_\allowbreak{}gamma\_\allowbreak{}HQ!: $n, \kappa$,\allowbreak{} $\tau_{n}$ at $\alpha = 0$ & PASS \\
\texttt{D-\allowbreak{}BGf1} & matrix1d\_\allowbreak{}visc\_\allowbreak{}HQ\_\allowbreak{}BG!,\allowbreak{} $\epsilon = 10^{-3}$: $\partial_{\alpha}n$ only,\allowbreak{} not $\kappa$ & PASS \\
\texttt{D-\allowbreak{}BGc1} & matrix1d\_\allowbreak{}visc\_\allowbreak{}HQ\_\allowbreak{}BG\_\allowbreak{}consistent!,\allowbreak{} $\epsilon = 10^{-3}$: enters $\partial_{\alpha}n$ in both rows & PASS \\
\texttt{D-\allowbreak{}BGm1} & matrix1d\_\allowbreak{}visc\_\allowbreak{}HQ\_\allowbreak{}BG\_\allowbreak{}consistent\_\allowbreak{}m2!,\allowbreak{} $\epsilon = 10^{-3}$ & PASS \\
\texttt{D-\allowbreak{}BGf2} & matrix2d\_\allowbreak{}HQ\_\allowbreak{}BG!,\allowbreak{} $\epsilon = 10^{-3}$ & PASS \\
\texttt{D-\allowbreak{}BGc2} & matrix2d\_\allowbreak{}HQ\_\allowbreak{}BG\_\allowbreak{}consistent!,\allowbreak{} $\epsilon = 10^{-3}$ & PASS \\
\texttt{D-\allowbreak{}BGm2} & matrix2d\_\allowbreak{}HQ\_\allowbreak{}BG\_\allowbreak{}m2!,\allowbreak{} $\epsilon = 10^{-3}$ & PASS \\
\texttt{D-\allowbreak{}live} & matrix2d\_\allowbreak{}visc\_\allowbreak{}HQ\_\allowbreak{}BG!,\allowbreak{} $\epsilon = 0$: derived medium with $\delta_{\pi\pi}$,\allowbreak{} passive medium,\allowbreak{} live charm rows & PASS \\
\texttt{D-\allowbreak{}tw1} & matrix1d\_\allowbreak{}two\_\allowbreak{}way!,\allowbreak{} $\epsilon = 0$,\allowbreak{} no edge & PASS \\
\texttt{D-\allowbreak{}tw1c} & matrix1d\_\allowbreak{}two\_\allowbreak{}way!,\allowbreak{} conformal & PASS \\
\texttt{D-\allowbreak{}tw2} & matrix2d\_\allowbreak{}two\_\allowbreak{}way!,\allowbreak{} $\epsilon = 0$,\allowbreak{} no edge & PASS \\
\texttt{D-\allowbreak{}tw2c} & matrix2d\_\allowbreak{}two\_\allowbreak{}way!,\allowbreak{} conformal & PASS \\
\texttt{D-\allowbreak{}r4} & matrix1d\_\allowbreak{}ideal\_\allowbreak{}HQ\_\allowbreak{}density!: the medium EOS into ruwen4 & PASS \\
\texttt{D-\allowbreak{}r5} & matrix1d\_\allowbreak{}ideal\_\allowbreak{}HQ\_\allowbreak{}fugacity!: the medium EOS and n from the list into ruwen5 & PASS \\
\texttt{D-\allowbreak{}r7} & matrix1d\_\allowbreak{}ideal\_\allowbreak{}HQ\_\allowbreak{}br\_\allowbreak{}p\_\allowbreak{}fugacity!: P\_\allowbreak{}h =\allowbreak{} n T,\allowbreak{} n from the list & PASS \\
\texttt{D-\allowbreak{}r6} & matrix1d\_\allowbreak{}ideal\_\allowbreak{}HQ\_\allowbreak{}br\_\allowbreak{}p\_\allowbreak{}density!: $\alpha$ from n,\allowbreak{} P\_\allowbreak{}h =\allowbreak{} n T & PASS \\
\texttt{D-\allowbreak{}HQup} & matrix2d\_\allowbreak{}visc\_\allowbreak{}HQ!: phi\_\allowbreak{}8 is the fugacity,\allowbreak{} P in the pressure slot,\allowbreak{} $D_{s}$ n =\allowbreak{} $\kappa$ & PASS \\
\texttt{D-\allowbreak{}hess} & the Hessian of a charm list =\allowbreak{} differences of its gradient & PASS \\
\texttt{D-\allowbreak{}Ph} & d\_\allowbreak{}$\alpha$ P\_\allowbreak{}h /\allowbreak{} T =\allowbreak{} n,\allowbreak{} and charm\_\allowbreak{}$\alpha$ inverts n & PASS \\
\texttt{D-\allowbreak{}hqp} & hq\_\allowbreak{}pressure runs: the density of one charm quark & PASS \\
\texttt{D-\allowbreak{}mv} & matrix1d\_\allowbreak{}visc runs & PASS \\
\texttt{D-\allowbreak{}cm1} & matrix1d\_\allowbreak{}two\_\allowbreak{}way\_\allowbreak{}cM!,\allowbreak{} $\epsilon = 0$,\allowbreak{} no edge & PASS \\
\texttt{D-\allowbreak{}cm1c} & matrix1d\_\allowbreak{}two\_\allowbreak{}way\_\allowbreak{}cM!,\allowbreak{} conformal & PASS \\
\texttt{D-\allowbreak{}cm2} & matrix2d\_\allowbreak{}two\_\allowbreak{}way\_\allowbreak{}cM!,\allowbreak{} $\epsilon = 0$,\allowbreak{} no edge & PASS \\
\texttt{D-\allowbreak{}cm2c} & matrix2d\_\allowbreak{}two\_\allowbreak{}way\_\allowbreak{}cM!,\allowbreak{} conformal & PASS \\
\texttt{D-\allowbreak{}edge} & matrix1d\_\allowbreak{}two\_\allowbreak{}way! at T below $T_{\mathrm{edge}}$: $D_{s}\to f(T)D_{s}$ and nothing else & PASS \\
\texttt{D-\allowbreak{}edge2} & matrix2d\_\allowbreak{}two\_\allowbreak{}way! at T below $T_{\mathrm{edge}}$: $D_{s}\to f(T)D_{s}$ and nothing else & PASS \\
\texttt{D-\allowbreak{}edgew} & $w = (1 - f)/(1 - f_{0})$ inside the window & PASS \\
\texttt{D-\allowbreak{}edgeb} & f =\allowbreak{} 1 above $T_{\mathrm{edge}}$,\allowbreak{} $f = f_{0}$ below $T_{\mathrm{edge}} - w_{T}$; the weight the other way round & PASS \\
\texttt{D-\allowbreak{}twe} & matrix1d\_\allowbreak{}two\_\allowbreak{}way!,\allowbreak{} $\epsilon = 10^{-3}$: + $\epsilon\,u\cdot\partial\alpha$ on the number row's $\alpha$ column,\allowbreak{} nothing else & PASS \\
\texttt{D-\allowbreak{}kappa} & $\kappa\text{ of the code} = D_{s}\partial_{\alpha}n$,\allowbreak{} one species and a list & PASS \\
\texttt{D-\allowbreak{}taun} & $\tau_{n}\text{ of the code} = D_{s}h/T$,\allowbreak{} one species and a list & PASS \\
\texttt{D-\allowbreak{}I31} & $I_{31} = nTh$,\allowbreak{} $\mathcal{B} = Th$ & PASS \\
\texttt{D-\allowbreak{}mom} & $I_{31}/n,\ I_{41}/P_{0},\ I_{42}/P_{0},\ I_{52}/I_{42}$ of a Boltzmann species,\allowbreak{} momentum quadrature =\allowbreak{} $Th,\ \mathcal{A},\ \mathcal{B},\ \mathcal{C}$ & PASS \\
\texttt{D-\allowbreak{}cor} & corona\_\allowbreak{}update\_\allowbreak{}1d!,\allowbreak{} T below the window: $\partial_{\tau}(\tau rN^{\tau}) = -(G_{i+1/2} - G_{i-1/2})/\Delta r$,\allowbreak{} LLF,\allowbreak{} with the updated $\alpha$ rate & PASS \\
\texttt{D-\allowbreak{}tele} & the flux form telescopes: the sum of r dr dN/\allowbreak{}dtau over the cells is the boundary flux difference & PASS \\
\texttt{D-\allowbreak{}cor0} & above $T_{\mathrm{edge}}$ the update changes nothing & PASS \\
\texttt{D-\allowbreak{}corb} & inside the window: $\dot\alpha = (1 - w)\dot\alpha_{\mathrm{ql}} + w\dot\alpha_{\mathrm{flux}}$ & PASS \\
\texttt{D-\allowbreak{}cor2} & corona\_\allowbreak{}update\_\allowbreak{}2d!,\allowbreak{} T below the window: $\partial_{\tau}(\tau N^{\tau})$ =\allowbreak{} the LLF flux differences & PASS \\
\bottomrule
\end{longtable}
}
94 passed, 0 failed (94 checks evaluated).

Re-run with \texttt{wolframscript -file deriv/derive\_\allowbreak{}drivers.wls}.


\clearpage
\appendix
% GENERATED by deriv/render.wl from the derivation script.  Do not edit.

\section{The entries of the remaining two-way and charm matrices}

As computed by \texttt{deriv/derive\_two\_way.wls}: the entries of its generated blocks, each compared with the derived expression (E-same). Field names and $u^{\tau}$ as in the section on the two-way system.

\subsection{1+1D cylindrical}

\subsubsection{S5, two-way}

\texttt{one\_\allowbreak{}d\_\allowbreak{}two\_\allowbreak{}way\_\allowbreak{}matrix} in \texttt{src/Matrix/1d\_\allowbreak{}two\_\allowbreak{}way\_\allowbreak{}cilindrical.jl}, fields:

\[
\phi = ({T}, {u^{r}}, {\pi^{\phi}{}_{\phi}}, {\pi^{\eta}{}_{\eta}}, {\Pi}, {\alpha}, {\nu^{r}}, {\pi_{Q}{}^{\phi}{}_{\phi}}, {\pi_{Q}{}^{\eta}{}_{\eta}}, {\Pi_{Q}})
\]
\begin{align*}
(A_t)_{1,1} &= {P''(T)} {T} + 2 {h'(T)} {\nu^{r}} {u^{r}} + {P'(T)} \left(u^{r}\right)^2 + {P''(T)} {T} \left(u^{r}\right)^2\\
  &\quad + {n} \left({h'(T)} \left(u^{r}\right)^2+{h'(T)}-1\right) + {(\partial_{T} n)} \left({h} \left(u^{r}\right)^2+{h}-{T}\right) \\
(A_t)_{1,2} &= 2 ({u^{r}} ({P'(T)} {T}+{\Pi}-{\pi^{\eta}{}_{\eta}}-{\pi^{\phi}{}_{\phi}})+{h} ({n} {u^{r}}+{\nu^{r}})) \\
(A_t)_{1,3} &= -\left(u^{r}\right)^2 \\
(A_t)_{1,4} &= -\left(u^{r}\right)^2 \\
(A_t)_{1,5} &= \left(u^{r}\right)^2 \\
(A_t)_{1,6} &= {(\partial_{\alpha} n)} \left({h} \left(u^{r}\right)^2+{h}-{T}\right) \\
(A_t)_{1,7} &= 2 {h} {u^{r}} \\
(A_t)_{2,1} &= \frac{{h'(T)} \left({n} {u^{r}} \left(u^{\tau}\right)^2+2 {\nu^{r}} \left(u^{r}\right)^2+{\nu^{r}}\right)+{u^{r}} \left(u^{\tau}\right)^2 ({(\partial_{T} n)} {h}+{P''(T)} {T}+{P'(T)})}{{u^{\tau}}} \\
(A_t)_{2,2} &= \frac{1}{\left(u^{\tau}\right)^3}\Bigl(\left(2 \left(u^{r}\right)^4+3 \left(u^{r}\right)^2+1\right) ({P'(T)} {T}+{\Pi}-{\pi^{\eta}{}_{\eta}}-{\pi^{\phi}{}_{\phi}})\\
  &\quad + {h} \left(2 {n} \left(u^{r}\right)^4+3 {n} \left(u^{r}\right)^2+{n}+2 {\nu^{r}} \left(u^{r}\right)^3+3 {\nu^{r}} {u^{r}}\right)\Bigr) \\
(A_t)_{2,3} &= -{u^{r}} {u^{\tau}} \\
(A_t)_{2,4} &= -{u^{r}} {u^{\tau}} \\
(A_t)_{2,5} &= {u^{r}} {u^{\tau}} \\
(A_t)_{2,6} &= {(\partial_{\alpha} n)} {h} {u^{r}} {u^{\tau}} \\
(A_t)_{2,7} &= \frac{{h} \left(2 \left(u^{r}\right)^2+1\right)}{{u^{\tau}}} \\
(A_t)_{3,2} &= \frac{{u^{r}} (3 {\delta_{\pi\pi}} {\pi^{\phi}{}_{\phi}}-2 {\eta})}{3 {r}^2 {u^{\tau}}} \\
(A_t)_{3,3} &= \frac{{\tau_{\pi}} {u^{\tau}}}{{r}^2} \\
(A_t)_{4,2} &= \frac{{u^{r}} (3 {\delta_{\pi\pi}} {\pi^{\eta}{}_{\eta}}-2 {\eta})}{3 {\tau}^2 {u^{\tau}}} \\
(A_t)_{4,4} &= \frac{{\tau_{\pi}} {u^{\tau}}}{{\tau}^2} \\
(A_t)_{5,2} &= \frac{{u^{r}} {\zeta}}{{u^{\tau}}} \\
(A_t)_{5,5} &= {\tau_{\Pi}} {u^{\tau}} \\
(A_t)_{6,1} &= {(\partial_{T} n)} {u^{\tau}} \\
(A_t)_{6,2} &= \frac{{n} {u^{r}} \left(u^{\tau}\right)^2+{\nu^{r}}}{\left(u^{\tau}\right)^3} \\
(A_t)_{6,6} &= {(\partial_{\alpha} n)} {u^{\tau}} \\
(A_t)_{6,7} &= \frac{{u^{r}}}{{u^{\tau}}} \\
(A_t)_{7,1} &= \frac{{D_{s}} {u^{\tau}} ({h'(T)} {\nu^{r}}+{u^{r}} ({(\partial_{T} n)} {T}+{n}))}{{T}} \\
(A_t)_{7,2} &= \frac{{D_{s}} {h} \left({n} \left(u^{r}\right)^2+{n}+{\nu^{r}} {u^{r}}\right)}{{T} {u^{\tau}}} \\
(A_t)_{7,6} &= {(\partial_{\alpha} n)} {D_{s}} {u^{r}} {u^{\tau}} \\
(A_t)_{7,7} &= \frac{{D_{s}} {h} {u^{\tau}}}{{T}} \\
(A_t)_{8,1} &= \frac{{D_{s}} \left(3 {\mathcal{C}'(T)} {\pi_{Q}{}^{\phi}{}_{\phi}} \left(u^{\tau}\right)^2-2 {\mathcal{B}'(T)} {\nu^{r}} {u^{r}}\right)}{6 {T} {u^{\tau}}} \\
(A_t)_{8,2} &= -\frac{{D_{s}} \left(2 {\mathcal{A}} {\nu^{r}} \left(u^{\tau}\right)^2+2 {\mathcal{B}} {\nu^{r}}+{u^{r}} \left(u^{\tau}\right)^2 (2 {\mathcal{C}} {\Pi_{Q}}-2 {\mathcal{C}} {\pi_{Q}{}^{\eta}{}_{\eta}}-5 {\mathcal{C}} {\pi_{Q}{}^{\phi}{}_{\phi}}+2 {I_{31}})\right)}{6 {T} \left(u^{\tau}\right)^3} \\
(A_t)_{8,7} &= -\frac{{\mathcal{B}} {D_{s}} {u^{r}}}{3 {T} {u^{\tau}}} \\
(A_t)_{8,8} &= \frac{{\mathcal{C}} {D_{s}} {u^{\tau}}}{2 {T}} \\
(A_t)_{9,1} &= \frac{{D_{s}} \left(3 {\mathcal{C}'(T)} {\pi_{Q}{}^{\eta}{}_{\eta}} \left(u^{\tau}\right)^2-2 {\mathcal{B}'(T)} {\nu^{r}} {u^{r}}\right)}{6 {T} {u^{\tau}}} \\
(A_t)_{9,2} &= -\frac{{D_{s}} \left(2 {\mathcal{A}} {\nu^{r}} \left(u^{\tau}\right)^2+2 {\mathcal{B}} {\nu^{r}}+{u^{r}} \left(u^{\tau}\right)^2 (2 {\mathcal{C}} {\Pi_{Q}}-5 {\mathcal{C}} {\pi_{Q}{}^{\eta}{}_{\eta}}-2 {\mathcal{C}} {\pi_{Q}{}^{\phi}{}_{\phi}}+2 {I_{31}})\right)}{6 {T} \left(u^{\tau}\right)^3} \\
(A_t)_{9,7} &= -\frac{{\mathcal{B}} {D_{s}} {u^{r}}}{3 {T} {u^{\tau}}} \\
(A_t)_{9,9} &= \frac{{\mathcal{C}} {D_{s}} {u^{\tau}}}{2 {T}} \\
(A_t)_{10,1} &= \frac{{D_{s}} \left(5 {\mathcal{B}'(T)} {\nu^{r}} {u^{r}}+3 {\mathcal{C}'(T)} {\Pi_{Q}} \left(u^{\tau}\right)^2+3 {(\partial_{T} I_{31})} \left(u^{\tau}\right)^2\right)}{6 {T} {u^{\tau}}} \\
(A_t)_{10,2} &= \frac{{D_{s}} \left(2 {\mathcal{A}} {\nu^{r}} \left(u^{\tau}\right)^2+5 {\mathcal{B}} {\nu^{r}}+{u^{r}} \left(u^{\tau}\right)^2 (5 {\mathcal{C}} {\Pi_{Q}}-2 {\mathcal{C}} ({\pi_{Q}{}^{\eta}{}_{\eta}}+{\pi_{Q}{}^{\phi}{}_{\phi}})+5 {I_{31}})\right)}{6 {T} \left(u^{\tau}\right)^3} \\
(A_t)_{10,6} &= \frac{{(\partial_{\alpha} I_{31})} {D_{s}} {u^{\tau}}}{2 {T}} \\
(A_t)_{10,7} &= \frac{5 {\mathcal{B}} {D_{s}} {u^{r}}}{6 {T} {u^{\tau}}} \\
(A_t)_{10,10} &= \frac{{\mathcal{C}} {D_{s}} {u^{\tau}}}{2 {T}} \\
(A_r)_{1,1} &= \frac{{h'(T)} \left({n} {u^{r}} \left(u^{\tau}\right)^2+2 {\nu^{r}} \left(u^{r}\right)^2+{\nu^{r}}\right)+{u^{r}} \left(u^{\tau}\right)^2 ({(\partial_{T} n)} {h}+{P''(T)} {T}+{P'(T)})}{{u^{\tau}}} \\
(A_r)_{1,2} &= \frac{1}{\left(u^{\tau}\right)^3}\Bigl(\left(2 \left(u^{r}\right)^4+3 \left(u^{r}\right)^2+1\right) ({P'(T)} {T}+{\Pi}-{\pi^{\eta}{}_{\eta}}-{\pi^{\phi}{}_{\phi}})\\
  &\quad + {h} \left(2 {n} \left(u^{r}\right)^4+3 {n} \left(u^{r}\right)^2+{n}+2 {\nu^{r}} \left(u^{r}\right)^3+3 {\nu^{r}} {u^{r}}\right)\Bigr) \\
(A_r)_{1,3} &= -{u^{r}} {u^{\tau}} \\
(A_r)_{1,4} &= -{u^{r}} {u^{\tau}} \\
(A_r)_{1,5} &= {u^{r}} {u^{\tau}} \\
(A_r)_{1,6} &= {(\partial_{\alpha} n)} {h} {u^{r}} {u^{\tau}} \\
(A_r)_{1,7} &= \frac{{h} \left(2 \left(u^{r}\right)^2+1\right)}{{u^{\tau}}} \\
(A_r)_{2,1} &= {P'(T)} + {n} + {(\partial_{T} n)} {T} + 2 {h'(T)} {\nu^{r}} {u^{r}} + {P'(T)} \left(u^{r}\right)^2 + {(\partial_{T} n)} {h} \left(u^{r}\right)^2 + {h'(T)} {n} \left(u^{r}\right)^2\\
  &\quad + {P''(T)} {T} \left(u^{r}\right)^2 \\
(A_r)_{2,2} &= 2 ({u^{r}} ({P'(T)} {T}+{\Pi}-{\pi^{\eta}{}_{\eta}}-{\pi^{\phi}{}_{\phi}})+{h} ({n} {u^{r}}+{\nu^{r}})) \\
(A_r)_{2,3} &= -\left(u^{r}\right)^2-1 \\
(A_r)_{2,4} &= -\left(u^{r}\right)^2-1 \\
(A_r)_{2,5} &= \left(u^{\tau}\right)^2 \\
(A_r)_{2,6} &= {(\partial_{\alpha} n)} \left({h} \left(u^{r}\right)^2+{T}\right) \\
(A_r)_{2,7} &= 2 {h} {u^{r}} \\
(A_r)_{3,2} &= \frac{3 {\delta_{\pi\pi}} {\pi^{\phi}{}_{\phi}}-2 {\eta}}{3 {r}^2} \\
(A_r)_{3,3} &= \frac{{\tau_{\pi}} {u^{r}}}{{r}^2} \\
(A_r)_{4,2} &= \frac{3 {\delta_{\pi\pi}} {\pi^{\eta}{}_{\eta}}-2 {\eta}}{3 {\tau}^2} \\
(A_r)_{4,4} &= \frac{{\tau_{\pi}} {u^{r}}}{{\tau}^2} \\
(A_r)_{5,2} &= {\zeta} \\
(A_r)_{5,5} &= {\tau_{\Pi}} {u^{r}} \\
(A_r)_{6,1} &= {(\partial_{T} n)} {u^{r}} \\
(A_r)_{6,2} &= {n} \\
(A_r)_{6,6} &= {(\partial_{\alpha} n)} {u^{r}} \\
(A_r)_{6,7} &= 1 \\
(A_r)_{7,1} &= \frac{{D_{s}} \left({h'(T)} {\nu^{r}} {u^{r}}+{(\partial_{T} n)} {T} \left(u^{\tau}\right)^2+{n} \left(u^{\tau}\right)^2\right)}{{T}} \\
(A_r)_{7,2} &= \frac{{D_{s}} {h} \left({n} {u^{r}} \left(u^{\tau}\right)^2+{\nu^{r}} \left(\left(u^{r}\right)^2+2\right)\right)}{{T} \left(u^{\tau}\right)^2} \\
(A_r)_{7,6} &= {(\partial_{\alpha} n)} {D_{s}} \left(u^{\tau}\right)^2 \\
(A_r)_{7,7} &= \frac{{D_{s}} {h} {u^{r}}}{{T}} \\
(A_r)_{8,1} &= -\frac{{D_{s}} (2 {\mathcal{B}'(T)} {\nu^{r}}-3 {\mathcal{C}'(T)} {\pi_{Q}{}^{\phi}{}_{\phi}} {u^{r}})}{6 {T}} \\
(A_r)_{8,2} &= -\frac{{D_{s}} \left(2 {\mathcal{A}} {\nu^{r}} {u^{r}}+{\mathcal{C}} \left(u^{\tau}\right)^2 (2 {\Pi_{Q}}-2 {\pi_{Q}{}^{\eta}{}_{\eta}}-5 {\pi_{Q}{}^{\phi}{}_{\phi}})+2 {I_{31}} \left(u^{\tau}\right)^2\right)}{6 {T} \left(u^{\tau}\right)^2} \\
(A_r)_{8,7} &= -\frac{{\mathcal{B}} {D_{s}}}{3 {T}} \\
(A_r)_{8,8} &= \frac{{\mathcal{C}} {D_{s}} {u^{r}}}{2 {T}} \\
(A_r)_{9,1} &= -\frac{{D_{s}} (2 {\mathcal{B}'(T)} {\nu^{r}}-3 {\mathcal{C}'(T)} {\pi_{Q}{}^{\eta}{}_{\eta}} {u^{r}})}{6 {T}} \\
(A_r)_{9,2} &= -\frac{{D_{s}} \left(2 {\mathcal{A}} {\nu^{r}} {u^{r}}+{\mathcal{C}} \left(u^{\tau}\right)^2 (2 {\Pi_{Q}}-5 {\pi_{Q}{}^{\eta}{}_{\eta}}-2 {\pi_{Q}{}^{\phi}{}_{\phi}})+2 {I_{31}} \left(u^{\tau}\right)^2\right)}{6 {T} \left(u^{\tau}\right)^2} \\
(A_r)_{9,7} &= -\frac{{\mathcal{B}} {D_{s}}}{3 {T}} \\
(A_r)_{9,9} &= \frac{{\mathcal{C}} {D_{s}} {u^{r}}}{2 {T}} \\
(A_r)_{10,1} &= \frac{{D_{s}} (5 {\mathcal{B}'(T)} {\nu^{r}}+3 {u^{r}} ({\mathcal{C}'(T)} {\Pi_{Q}}+{(\partial_{T} I_{31})}))}{6 {T}} \\
(A_r)_{10,2} &= \frac{{D_{s}} \left(2 {\mathcal{A}} {\nu^{r}} {u^{r}}+{\mathcal{C}} \left(u^{\tau}\right)^2 (5 {\Pi_{Q}}-2 ({\pi_{Q}{}^{\eta}{}_{\eta}}+{\pi_{Q}{}^{\phi}{}_{\phi}}))+5 {I_{31}} \left(u^{\tau}\right)^2\right)}{6 {T} \left(u^{\tau}\right)^2} \\
(A_r)_{10,6} &= \frac{{(\partial_{\alpha} I_{31})} {D_{s}} {u^{r}}}{2 {T}} \\
(A_r)_{10,7} &= \frac{5 {\mathcal{B}} {D_{s}}}{6 {T}} \\
(A_r)_{10,10} &= \frac{{\mathcal{C}} {D_{s}} {u^{r}}}{2 {T}} \\
S_{1} &= \frac{1}{{r} {\tau} {u^{\tau}}}\Bigl(- {\pi^{\eta}{}_{\eta}} {\tau} {u^{r}} - {\pi^{\phi}{}_{\phi}} {\tau} {u^{r}} + {P'(T)} {T} {\tau} {u^{r}} - {\pi^{\eta}{}_{\eta}} {\tau} \left(u^{r}\right)^3 - {\pi^{\phi}{}_{\phi}} {\tau} \left(u^{r}\right)^3 + {P'(T)} {T} {\tau} \left(u^{r}\right)^3\\
  &\quad + {\pi^{\eta}{}_{\eta}} {r} {u^{\tau}} + {P'(T)} {r} {T} {u^{\tau}} - {\pi^{\eta}{}_{\eta}} {r} \left(u^{r}\right)^2 {u^{\tau}} - {\pi^{\phi}{}_{\phi}} {r} \left(u^{r}\right)^2 {u^{\tau}} + {P'(T)} {r} {T} \left(u^{r}\right)^2 {u^{\tau}}\\
  &\quad + {h} {n} \left(u^{\tau}\right)^2 ({r} {u^{\tau}}+{\tau} {u^{r}}) + {\Pi} \left(u^{\tau}\right)^2 ({r} {u^{\tau}}+{\tau} {u^{r}}) + {h} {\nu^{r}} \left(2 {r} {u^{r}} {u^{\tau}}+2 {\tau} \left(u^{r}\right)^2+{\tau}\right)\Bigr) \\
S_{2} &= - \frac{{\Pi}+{\pi^{\phi}{}_{\phi}}+{P}}{{r}} - \frac{{n} {T}}{{r}} + \frac{{h} {n} \left(u^{r}\right)^2+2 {h} {\nu^{r}} {u^{r}}+{n} {T}}{{r}} + \frac{{u^{r}} {u^{\tau}} ({P'(T)} {T}+{\Pi}-{\pi^{\eta}{}_{\eta}}-{\pi^{\phi}{}_{\phi}})}{{\tau}}\\
  &\quad + \frac{\left(u^{r}\right)^2 ({P'(T)} {T}+{\Pi})+{\Pi}-\left(u^{\tau}\right)^2 ({\pi^{\eta}{}_{\eta}}+{\pi^{\phi}{}_{\phi}})+{P}}{{r}} + \frac{{h} \left({n} {u^{r}} \left(u^{\tau}\right)^2+2 {\nu^{r}} \left(u^{r}\right)^2+{\nu^{r}}\right)}{{\tau} {u^{\tau}}} \\
S_{3} &= \frac{3 {\pi^{\phi}{}_{\phi}} ({r} ({\delta_{\pi\pi}} {u^{\tau}}+{\tau})+{\delta_{\pi\pi}} {\tau} {u^{r}})-2 {r} {u^{\tau}} {\eta}+4 {\tau} {u^{r}} {\eta}}{3 {r}^3 {\tau}} \\
S_{4} &= \frac{3 {\pi^{\eta}{}_{\eta}} ({r} ({\delta_{\pi\pi}} {u^{\tau}}+{\tau})+{\delta_{\pi\pi}} {\tau} {u^{r}})+4 {r} {u^{\tau}} {\eta}-2 {\tau} {u^{r}} {\eta}}{3 {r} {\tau}^3} \\
S_{5} &= {\Pi}+\frac{{u^{r}} {\zeta}}{{r}}+\frac{{u^{\tau}} {\zeta}}{{\tau}} \\
S_{6} &= \frac{{n} {u^{r}}+{\nu^{r}}}{{r}}+\frac{{n} \left(u^{r}\right)^2+{n}+{\nu^{r}} {u^{r}}}{{\tau} {u^{\tau}}} \\
S_{7} &= \frac{{D_{s}} {h} {\nu^{r}} \left({r} \left(u^{r}\right)^2+{r}+{\tau} {u^{r}} {u^{\tau}}\right)}{{r} {T} {\tau} {u^{\tau}}}+{\nu^{r}} \\
S_{8} &= {\pi_{Q}{}^{\phi}{}_{\phi}}\\
  &\quad + \frac{1}{6 {r} {T} {\tau} {u^{\tau}}}{D_{s}}\Bigl(- 2 {\mathcal{B}} {\nu^{r}} ({r} {u^{r}}-2 {\tau} {u^{\tau}}) - 2 {I_{31}} \left({r} \left(u^{r}\right)^2+{r}-2 {\tau} {u^{r}} {u^{\tau}}\right)\\
  &\quad + {\mathcal{C}}\Bigl(7 {\pi_{Q}{}^{\phi}{}_{\phi}} {\tau} {u^{r}} {u^{\tau}} - 2 {\pi_{Q}{}^{\eta}{}_{\eta}} {r} \left(u^{\tau}\right)^2 + 3 {\pi_{Q}{}^{\phi}{}_{\phi}} {r} \left(u^{\tau}\right)^2\\
  &\quad - 2 {\Pi_{Q}} \left({r} \left(u^{r}\right)^2+{r}-2 {\tau} {u^{r}} {u^{\tau}}\right)\Bigr)\Bigr) \\
S_{9} &= {\pi_{Q}{}^{\eta}{}_{\eta}}\\
  &\quad + \frac{1}{6 {r} {T} {\tau} {u^{\tau}}}{D_{s}}\Bigl(- 2 {I_{31}} {\tau} {u^{r}} {u^{\tau}} + 4 {I_{31}} {r} \left(u^{\tau}\right)^2 + 2 {\mathcal{B}} {\nu^{r}} (2 {r} {u^{r}}-{\tau} {u^{\tau}})\\
  &\quad + {\mathcal{C}}\Bigl(- 2 {\Pi_{Q}} {\tau} {u^{r}} {u^{\tau}} + 3 {\pi_{Q}{}^{\eta}{}_{\eta}} {\tau} {u^{r}} {u^{\tau}} - 2 {\pi_{Q}{}^{\phi}{}_{\phi}} {\tau} {u^{r}} {u^{\tau}} + 4 {\Pi_{Q}} {r} \left(u^{\tau}\right)^2\\
  &\quad + 7 {\pi_{Q}{}^{\eta}{}_{\eta}} {r} \left(u^{\tau}\right)^2\Bigr)\Bigr) \\
S_{10} &= {\Pi_{Q}}\\
  &\quad + \frac{1}{6 {r} {T} {\tau} {u^{\tau}}}{D_{s}}\Bigl(5 {\mathcal{B}} {\nu^{r}} ({r} {u^{r}}+{\tau} {u^{\tau}}) + 5 {I_{31}} \left({r} \left(u^{r}\right)^2+{r}+{\tau} {u^{r}} {u^{\tau}}\right)\\
  &\quad + {\mathcal{C}} \left(5 {\Pi_{Q}} \left({r} \left(u^{r}\right)^2+{r}+{\tau} {u^{r}} {u^{\tau}}\right)+2 {\pi_{Q}{}^{\eta}{}_{\eta}} {r} \left(u^{\tau}\right)^2+2 {\pi_{Q}{}^{\phi}{}_{\phi}} {\tau} {u^{r}} {u^{\tau}}\right)\Bigr)
\end{align*}
\subsubsection{S6, fully coupled}

\texttt{one\_\allowbreak{}d\_\allowbreak{}two\_\allowbreak{}way\_\allowbreak{}cM\_\allowbreak{}matrix} in \texttt{src/Matrix/1d\_\allowbreak{}two\_\allowbreak{}way\_\allowbreak{}cM\_\allowbreak{}cilindrical.jl}, fields:

\[
\phi = ({T}, {u^{r}}, {\pi^{\phi}{}_{\phi}}, {\pi^{\eta}{}_{\eta}}, {\Pi}, {\alpha}, {\nu^{r}}, {\pi_{Q}{}^{\phi}{}_{\phi}}, {\pi_{Q}{}^{\eta}{}_{\eta}}, {\Pi_{Q}})
\]
\begin{align*}
(A_t)_{1,1} &= {P''(T)} {T} + 2 {h'(T)} {\nu^{r}} {u^{r}} + {P'(T)} \left(u^{r}\right)^2 + {P''(T)} {T} \left(u^{r}\right)^2\\
  &\quad + {n} \left({h'(T)} \left(u^{r}\right)^2+{h'(T)}-1\right) + {(\partial_{T} n)} \left({h} \left(u^{r}\right)^2+{h}-{T}\right) \\
(A_t)_{1,2} &= 2 ({u^{r}} ({P'(T)} {T}+{\Pi}+{\Pi_{Q}}-{\pi^{\eta}{}_{\eta}}-{\pi_{Q}{}^{\eta}{}_{\eta}}-{\pi^{\phi}{}_{\phi}}-{\pi_{Q}{}^{\phi}{}_{\phi}})+{h} ({n} {u^{r}}+{\nu^{r}})) \\
(A_t)_{1,3} &= -\left(u^{r}\right)^2 \\
(A_t)_{1,4} &= -\left(u^{r}\right)^2 \\
(A_t)_{1,5} &= \left(u^{r}\right)^2 \\
(A_t)_{1,6} &= {(\partial_{\alpha} n)} \left({h} \left(u^{r}\right)^2+{h}-{T}\right) \\
(A_t)_{1,7} &= 2 {h} {u^{r}} \\
(A_t)_{1,8} &= -\left(u^{r}\right)^2 \\
(A_t)_{1,9} &= -\left(u^{r}\right)^2 \\
(A_t)_{1,10} &= \left(u^{r}\right)^2 \\
(A_t)_{2,1} &= \frac{{h'(T)} \left({n} {u^{r}} \left(u^{\tau}\right)^2+2 {\nu^{r}} \left(u^{r}\right)^2+{\nu^{r}}\right)+{u^{r}} \left(u^{\tau}\right)^2 ({(\partial_{T} n)} {h}+{P''(T)} {T}+{P'(T)})}{{u^{\tau}}} \\
(A_t)_{2,2} &= \frac{1}{\left(u^{\tau}\right)^3}\Bigl(\left(2 \left(u^{r}\right)^4+3 \left(u^{r}\right)^2+1\right) ({P'(T)} {T}+{\Pi}+{\Pi_{Q}}-{\pi^{\eta}{}_{\eta}}-{\pi_{Q}{}^{\eta}{}_{\eta}}-{\pi^{\phi}{}_{\phi}}-{\pi_{Q}{}^{\phi}{}_{\phi}})\\
  &\quad + {h} \left(2 {n} \left(u^{r}\right)^4+3 {n} \left(u^{r}\right)^2+{n}+2 {\nu^{r}} \left(u^{r}\right)^3+3 {\nu^{r}} {u^{r}}\right)\Bigr) \\
(A_t)_{2,3} &= -{u^{r}} {u^{\tau}} \\
(A_t)_{2,4} &= -{u^{r}} {u^{\tau}} \\
(A_t)_{2,5} &= {u^{r}} {u^{\tau}} \\
(A_t)_{2,6} &= {(\partial_{\alpha} n)} {h} {u^{r}} {u^{\tau}} \\
(A_t)_{2,7} &= \frac{{h} \left(2 \left(u^{r}\right)^2+1\right)}{{u^{\tau}}} \\
(A_t)_{2,8} &= -{u^{r}} {u^{\tau}} \\
(A_t)_{2,9} &= -{u^{r}} {u^{\tau}} \\
(A_t)_{2,10} &= {u^{r}} {u^{\tau}} \\
(A_t)_{3,2} &= \frac{{u^{r}} (3 {\delta_{\pi\pi}} {\pi^{\phi}{}_{\phi}}-2 {\eta})}{3 {r}^2 {u^{\tau}}} \\
(A_t)_{3,3} &= \frac{{\tau_{\pi}} {u^{\tau}}}{{r}^2} \\
(A_t)_{4,2} &= \frac{{u^{r}} (3 {\delta_{\pi\pi}} {\pi^{\eta}{}_{\eta}}-2 {\eta})}{3 {\tau}^2 {u^{\tau}}} \\
(A_t)_{4,4} &= \frac{{\tau_{\pi}} {u^{\tau}}}{{\tau}^2} \\
(A_t)_{5,2} &= \frac{{u^{r}} {\zeta}}{{u^{\tau}}} \\
(A_t)_{5,5} &= {\tau_{\Pi}} {u^{\tau}} \\
(A_t)_{6,1} &= {(\partial_{T} n)} {u^{\tau}} \\
(A_t)_{6,2} &= \frac{{n} {u^{r}} \left(u^{\tau}\right)^2+{\nu^{r}}}{\left(u^{\tau}\right)^3} \\
(A_t)_{6,6} &= {(\partial_{\alpha} n)} {u^{\tau}} \\
(A_t)_{6,7} &= \frac{{u^{r}}}{{u^{\tau}}} \\
(A_t)_{7,1} &= \frac{{D_{s}} {u^{\tau}} ({h'(T)} {\nu^{r}}+{u^{r}} ({(\partial_{T} n)} {T}+{n}))}{{T}} \\
(A_t)_{7,2} &= \frac{{D_{s}} \left({h} \left({n} \left(u^{r}\right)^2+{n}+{\nu^{r}} {u^{r}}\right)+\left(u^{\tau}\right)^2 ({\Pi_{Q}}-{\pi_{Q}{}^{\eta}{}_{\eta}}-{\pi_{Q}{}^{\phi}{}_{\phi}})\right)}{{T} {u^{\tau}}} \\
(A_t)_{7,6} &= {(\partial_{\alpha} n)} {D_{s}} {u^{r}} {u^{\tau}} \\
(A_t)_{7,7} &= \frac{{D_{s}} {h} {u^{\tau}}}{{T}} \\
(A_t)_{7,8} &= -\frac{{D_{s}} {u^{r}} {u^{\tau}}}{{T}} \\
(A_t)_{7,9} &= -\frac{{D_{s}} {u^{r}} {u^{\tau}}}{{T}} \\
(A_t)_{7,10} &= \frac{{D_{s}} {u^{r}} {u^{\tau}}}{{T}} \\
(A_t)_{8,1} &= \frac{{D_{s}} \left(3 {\mathcal{C}'(T)} {\pi_{Q}{}^{\phi}{}_{\phi}} \left(u^{\tau}\right)^2-2 {\mathcal{B}'(T)} {\nu^{r}} {u^{r}}\right)}{6 {T} {u^{\tau}}} \\
(A_t)_{8,2} &= -\frac{{D_{s}} \left(2 {\mathcal{A}} {\nu^{r}} \left(u^{\tau}\right)^2+2 {\mathcal{B}} {\nu^{r}}+{u^{r}} \left(u^{\tau}\right)^2 (2 {\mathcal{C}} {\Pi_{Q}}-2 {\mathcal{C}} {\pi_{Q}{}^{\eta}{}_{\eta}}-5 {\mathcal{C}} {\pi_{Q}{}^{\phi}{}_{\phi}}+2 {I_{31}})\right)}{6 {T} \left(u^{\tau}\right)^3} \\
(A_t)_{8,7} &= -\frac{{\mathcal{B}} {D_{s}} {u^{r}}}{3 {T} {u^{\tau}}} \\
(A_t)_{8,8} &= \frac{{\mathcal{C}} {D_{s}} {u^{\tau}}}{2 {T}} \\
(A_t)_{9,1} &= \frac{{D_{s}} \left(3 {\mathcal{C}'(T)} {\pi_{Q}{}^{\eta}{}_{\eta}} \left(u^{\tau}\right)^2-2 {\mathcal{B}'(T)} {\nu^{r}} {u^{r}}\right)}{6 {T} {u^{\tau}}} \\
(A_t)_{9,2} &= -\frac{{D_{s}} \left(2 {\mathcal{A}} {\nu^{r}} \left(u^{\tau}\right)^2+2 {\mathcal{B}} {\nu^{r}}+{u^{r}} \left(u^{\tau}\right)^2 (2 {\mathcal{C}} {\Pi_{Q}}-5 {\mathcal{C}} {\pi_{Q}{}^{\eta}{}_{\eta}}-2 {\mathcal{C}} {\pi_{Q}{}^{\phi}{}_{\phi}}+2 {I_{31}})\right)}{6 {T} \left(u^{\tau}\right)^3} \\
(A_t)_{9,7} &= -\frac{{\mathcal{B}} {D_{s}} {u^{r}}}{3 {T} {u^{\tau}}} \\
(A_t)_{9,9} &= \frac{{\mathcal{C}} {D_{s}} {u^{\tau}}}{2 {T}} \\
(A_t)_{10,1} &= \frac{{D_{s}} \left(5 {\mathcal{B}'(T)} {\nu^{r}} {u^{r}}+3 {\mathcal{C}'(T)} {\Pi_{Q}} \left(u^{\tau}\right)^2+3 {(\partial_{T} I_{31})} \left(u^{\tau}\right)^2\right)}{6 {T} {u^{\tau}}} \\
(A_t)_{10,2} &= \frac{{D_{s}} \left(2 {\mathcal{A}} {\nu^{r}} \left(u^{\tau}\right)^2+5 {\mathcal{B}} {\nu^{r}}+{u^{r}} \left(u^{\tau}\right)^2 (5 {\mathcal{C}} {\Pi_{Q}}-2 {\mathcal{C}} ({\pi_{Q}{}^{\eta}{}_{\eta}}+{\pi_{Q}{}^{\phi}{}_{\phi}})+5 {I_{31}})\right)}{6 {T} \left(u^{\tau}\right)^3} \\
(A_t)_{10,6} &= \frac{{(\partial_{\alpha} I_{31})} {D_{s}} {u^{\tau}}}{2 {T}} \\
(A_t)_{10,7} &= \frac{5 {\mathcal{B}} {D_{s}} {u^{r}}}{6 {T} {u^{\tau}}} \\
(A_t)_{10,10} &= \frac{{\mathcal{C}} {D_{s}} {u^{\tau}}}{2 {T}} \\
(A_r)_{1,1} &= \frac{{h'(T)} \left({n} {u^{r}} \left(u^{\tau}\right)^2+2 {\nu^{r}} \left(u^{r}\right)^2+{\nu^{r}}\right)+{u^{r}} \left(u^{\tau}\right)^2 ({(\partial_{T} n)} {h}+{P''(T)} {T}+{P'(T)})}{{u^{\tau}}} \\
(A_r)_{1,2} &= \frac{1}{\left(u^{\tau}\right)^3}\Bigl(\left(2 \left(u^{r}\right)^4+3 \left(u^{r}\right)^2+1\right) ({P'(T)} {T}+{\Pi}+{\Pi_{Q}}-{\pi^{\eta}{}_{\eta}}-{\pi_{Q}{}^{\eta}{}_{\eta}}-{\pi^{\phi}{}_{\phi}}-{\pi_{Q}{}^{\phi}{}_{\phi}})\\
  &\quad + {h} \left(2 {n} \left(u^{r}\right)^4+3 {n} \left(u^{r}\right)^2+{n}+2 {\nu^{r}} \left(u^{r}\right)^3+3 {\nu^{r}} {u^{r}}\right)\Bigr) \\
(A_r)_{1,3} &= -{u^{r}} {u^{\tau}} \\
(A_r)_{1,4} &= -{u^{r}} {u^{\tau}} \\
(A_r)_{1,5} &= {u^{r}} {u^{\tau}} \\
(A_r)_{1,6} &= {(\partial_{\alpha} n)} {h} {u^{r}} {u^{\tau}} \\
(A_r)_{1,7} &= \frac{{h} \left(2 \left(u^{r}\right)^2+1\right)}{{u^{\tau}}} \\
(A_r)_{1,8} &= -{u^{r}} {u^{\tau}} \\
(A_r)_{1,9} &= -{u^{r}} {u^{\tau}} \\
(A_r)_{1,10} &= {u^{r}} {u^{\tau}} \\
(A_r)_{2,1} &= {P'(T)} + {n} + {(\partial_{T} n)} {T} + 2 {h'(T)} {\nu^{r}} {u^{r}} + {P'(T)} \left(u^{r}\right)^2 + {(\partial_{T} n)} {h} \left(u^{r}\right)^2 + {h'(T)} {n} \left(u^{r}\right)^2\\
  &\quad + {P''(T)} {T} \left(u^{r}\right)^2 \\
(A_r)_{2,2} &= 2 ({u^{r}} ({P'(T)} {T}+{\Pi}+{\Pi_{Q}}-{\pi^{\eta}{}_{\eta}}-{\pi_{Q}{}^{\eta}{}_{\eta}}-{\pi^{\phi}{}_{\phi}}-{\pi_{Q}{}^{\phi}{}_{\phi}})+{h} ({n} {u^{r}}+{\nu^{r}})) \\
(A_r)_{2,3} &= -\left(u^{r}\right)^2-1 \\
(A_r)_{2,4} &= -\left(u^{r}\right)^2-1 \\
(A_r)_{2,5} &= \left(u^{\tau}\right)^2 \\
(A_r)_{2,6} &= {(\partial_{\alpha} n)} \left({h} \left(u^{r}\right)^2+{T}\right) \\
(A_r)_{2,7} &= 2 {h} {u^{r}} \\
(A_r)_{2,8} &= -\left(u^{r}\right)^2-1 \\
(A_r)_{2,9} &= -\left(u^{r}\right)^2-1 \\
(A_r)_{2,10} &= \left(u^{\tau}\right)^2 \\
(A_r)_{3,2} &= \frac{3 {\delta_{\pi\pi}} {\pi^{\phi}{}_{\phi}}-2 {\eta}}{3 {r}^2} \\
(A_r)_{3,3} &= \frac{{\tau_{\pi}} {u^{r}}}{{r}^2} \\
(A_r)_{4,2} &= \frac{3 {\delta_{\pi\pi}} {\pi^{\eta}{}_{\eta}}-2 {\eta}}{3 {\tau}^2} \\
(A_r)_{4,4} &= \frac{{\tau_{\pi}} {u^{r}}}{{\tau}^2} \\
(A_r)_{5,2} &= {\zeta} \\
(A_r)_{5,5} &= {\tau_{\Pi}} {u^{r}} \\
(A_r)_{6,1} &= {(\partial_{T} n)} {u^{r}} \\
(A_r)_{6,2} &= {n} \\
(A_r)_{6,6} &= {(\partial_{\alpha} n)} {u^{r}} \\
(A_r)_{6,7} &= 1 \\
(A_r)_{7,1} &= \frac{{D_{s}} \left({h'(T)} {\nu^{r}} {u^{r}}+{(\partial_{T} n)} {T} \left(u^{\tau}\right)^2+{n} \left(u^{\tau}\right)^2\right)}{{T}} \\
(A_r)_{7,2} &= \frac{{D_{s}} \left({h} {n} {u^{r}} \left(u^{\tau}\right)^2+{h} {\nu^{r}} \left(u^{r}\right)^2+2 {h} {\nu^{r}}+{\Pi_{Q}} {u^{r}} \left(u^{\tau}\right)^2-{\pi_{Q}{}^{\eta}{}_{\eta}} {u^{r}} \left(u^{\tau}\right)^2-{\pi_{Q}{}^{\phi}{}_{\phi}} {u^{r}} \left(u^{\tau}\right)^2\right)}{{T} \left(u^{\tau}\right)^2} \\
(A_r)_{7,6} &= {(\partial_{\alpha} n)} {D_{s}} \left(u^{\tau}\right)^2 \\
(A_r)_{7,7} &= \frac{{D_{s}} {h} {u^{r}}}{{T}} \\
(A_r)_{7,8} &= -\frac{{D_{s}} \left(u^{\tau}\right)^2}{{T}} \\
(A_r)_{7,9} &= -\frac{{D_{s}} \left(u^{\tau}\right)^2}{{T}} \\
(A_r)_{7,10} &= \frac{{D_{s}} \left(u^{\tau}\right)^2}{{T}} \\
(A_r)_{8,1} &= -\frac{{D_{s}} (2 {\mathcal{B}'(T)} {\nu^{r}}-3 {\mathcal{C}'(T)} {\pi_{Q}{}^{\phi}{}_{\phi}} {u^{r}})}{6 {T}} \\
(A_r)_{8,2} &= -\frac{{D_{s}} \left(2 {\mathcal{A}} {\nu^{r}} {u^{r}}+{\mathcal{C}} \left(u^{\tau}\right)^2 (2 {\Pi_{Q}}-2 {\pi_{Q}{}^{\eta}{}_{\eta}}-5 {\pi_{Q}{}^{\phi}{}_{\phi}})+2 {I_{31}} \left(u^{\tau}\right)^2\right)}{6 {T} \left(u^{\tau}\right)^2} \\
(A_r)_{8,7} &= -\frac{{\mathcal{B}} {D_{s}}}{3 {T}} \\
(A_r)_{8,8} &= \frac{{\mathcal{C}} {D_{s}} {u^{r}}}{2 {T}} \\
(A_r)_{9,1} &= -\frac{{D_{s}} (2 {\mathcal{B}'(T)} {\nu^{r}}-3 {\mathcal{C}'(T)} {\pi_{Q}{}^{\eta}{}_{\eta}} {u^{r}})}{6 {T}} \\
(A_r)_{9,2} &= -\frac{{D_{s}} \left(2 {\mathcal{A}} {\nu^{r}} {u^{r}}+{\mathcal{C}} \left(u^{\tau}\right)^2 (2 {\Pi_{Q}}-5 {\pi_{Q}{}^{\eta}{}_{\eta}}-2 {\pi_{Q}{}^{\phi}{}_{\phi}})+2 {I_{31}} \left(u^{\tau}\right)^2\right)}{6 {T} \left(u^{\tau}\right)^2} \\
(A_r)_{9,7} &= -\frac{{\mathcal{B}} {D_{s}}}{3 {T}} \\
(A_r)_{9,9} &= \frac{{\mathcal{C}} {D_{s}} {u^{r}}}{2 {T}} \\
(A_r)_{10,1} &= \frac{{D_{s}} (5 {\mathcal{B}'(T)} {\nu^{r}}+3 {u^{r}} ({\mathcal{C}'(T)} {\Pi_{Q}}+{(\partial_{T} I_{31})}))}{6 {T}} \\
(A_r)_{10,2} &= \frac{{D_{s}} \left(2 {\mathcal{A}} {\nu^{r}} {u^{r}}+{\mathcal{C}} \left(u^{\tau}\right)^2 (5 {\Pi_{Q}}-2 ({\pi_{Q}{}^{\eta}{}_{\eta}}+{\pi_{Q}{}^{\phi}{}_{\phi}}))+5 {I_{31}} \left(u^{\tau}\right)^2\right)}{6 {T} \left(u^{\tau}\right)^2} \\
(A_r)_{10,6} &= \frac{{(\partial_{\alpha} I_{31})} {D_{s}} {u^{r}}}{2 {T}} \\
(A_r)_{10,7} &= \frac{5 {\mathcal{B}} {D_{s}}}{6 {T}} \\
(A_r)_{10,10} &= \frac{{\mathcal{C}} {D_{s}} {u^{r}}}{2 {T}} \\
S_{1} &= \frac{{\Pi_{Q}}+{\pi_{Q}{}^{\eta}{}_{\eta}}}{{\tau}} + \frac{{\Pi}+{\pi^{\eta}{}_{\eta}}+{P}}{{\tau}} + \frac{{n} {T}}{{\tau}} + \frac{\left(u^{r}\right)^2 ({\Pi_{Q}}-{\pi_{Q}{}^{\eta}{}_{\eta}}-{\pi_{Q}{}^{\phi}{}_{\phi}})}{{\tau}}\\
  &\quad + \frac{{h} \left({n} \left(u^{r}\right)^2+{n}+2 {\nu^{r}} {u^{r}}\right)-{n} {T}}{{\tau}} + \frac{{u^{r}} {u^{\tau}} ({\Pi_{Q}}-{\pi_{Q}{}^{\eta}{}_{\eta}}-{\pi_{Q}{}^{\phi}{}_{\phi}})}{{r}}\\
  &\quad + \frac{{u^{r}} {u^{\tau}} ({P'(T)} {T}+{\Pi}-{\pi^{\eta}{}_{\eta}}-{\pi^{\phi}{}_{\phi}})}{{r}} + \frac{{P'(T)} {T} \left(u^{\tau}\right)^2+\left(u^{r}\right)^2 ({\Pi}-{\pi^{\eta}{}_{\eta}}-{\pi^{\phi}{}_{\phi}})-{P}}{{\tau}}\\
  &\quad + \frac{{h} \left({n} {u^{r}} \left(u^{\tau}\right)^2+2 {\nu^{r}} \left(u^{r}\right)^2+{\nu^{r}}\right)}{{r} {u^{\tau}}} \\
S_{2} &= - \frac{{\Pi_{Q}}+{\pi_{Q}{}^{\phi}{}_{\phi}}}{{r}} - \frac{{\Pi}+{\pi^{\phi}{}_{\phi}}+{P}}{{r}} - \frac{{n} {T}}{{r}} + \frac{{h} {n} \left(u^{r}\right)^2+2 {h} {\nu^{r}} {u^{r}}+{n} {T}}{{r}}\\
  &\quad + \frac{{u^{r}} {u^{\tau}} ({\Pi_{Q}}-{\pi_{Q}{}^{\eta}{}_{\eta}}-{\pi_{Q}{}^{\phi}{}_{\phi}})}{{\tau}} + \frac{{u^{r}} {u^{\tau}} ({P'(T)} {T}+{\Pi}-{\pi^{\eta}{}_{\eta}}-{\pi^{\phi}{}_{\phi}})}{{\tau}}\\
  &\quad + \frac{\left(u^{\tau}\right)^2 ({\Pi_{Q}}-{\pi_{Q}{}^{\eta}{}_{\eta}}-{\pi_{Q}{}^{\phi}{}_{\phi}})}{{r}} + \frac{\left(u^{r}\right)^2 ({P'(T)} {T}+{\Pi})+{\Pi}-\left(u^{\tau}\right)^2 ({\pi^{\eta}{}_{\eta}}+{\pi^{\phi}{}_{\phi}})+{P}}{{r}}\\
  &\quad + \frac{{h} \left({n} {u^{r}} \left(u^{\tau}\right)^2+2 {\nu^{r}} \left(u^{r}\right)^2+{\nu^{r}}\right)}{{\tau} {u^{\tau}}} \\
S_{3} &= \frac{3 {\pi^{\phi}{}_{\phi}} ({r} ({\delta_{\pi\pi}} {u^{\tau}}+{\tau})+{\delta_{\pi\pi}} {\tau} {u^{r}})-2 {r} {u^{\tau}} {\eta}+4 {\tau} {u^{r}} {\eta}}{3 {r}^3 {\tau}} \\
S_{4} &= \frac{3 {\pi^{\eta}{}_{\eta}} ({r} ({\delta_{\pi\pi}} {u^{\tau}}+{\tau})+{\delta_{\pi\pi}} {\tau} {u^{r}})+4 {r} {u^{\tau}} {\eta}-2 {\tau} {u^{r}} {\eta}}{3 {r} {\tau}^3} \\
S_{5} &= {\Pi}+\frac{{u^{r}} {\zeta}}{{r}}+\frac{{u^{\tau}} {\zeta}}{{\tau}} \\
S_{6} &= \frac{{n} {u^{r}}+{\nu^{r}}}{{r}}+\frac{{n} \left(u^{r}\right)^2+{n}+{\nu^{r}} {u^{r}}}{{\tau} {u^{\tau}}} \\
S_{7} &= \frac{1}{{r} {T} {\tau} {u^{\tau}}}\Bigl({\nu^{r}} {r} {T} {\tau} {u^{\tau}} - {D_{s}} \left(u^{\tau}\right)^2 (2 {\pi_{Q}{}^{\eta}{}_{\eta}} {r} {u^{r}}+{\pi_{Q}{}^{\eta}{}_{\eta}} {\tau} {u^{\tau}}+{\pi_{Q}{}^{\phi}{}_{\phi}} {r} {u^{r}}+2 {\pi_{Q}{}^{\phi}{}_{\phi}} {\tau} {u^{\tau}})\\
  &\quad + {D_{s}} {h} {\nu^{r}} \left({r} \left(u^{r}\right)^2+{r}+{\tau} {u^{r}} {u^{\tau}}\right)\Bigr) \\
S_{8} &= {\pi_{Q}{}^{\phi}{}_{\phi}}\\
  &\quad + \frac{1}{6 {r} {T} {\tau} {u^{\tau}}}{D_{s}}\Bigl(- 2 {\mathcal{B}} {\nu^{r}} ({r} {u^{r}}-2 {\tau} {u^{\tau}}) - 2 {I_{31}} \left({r} \left(u^{r}\right)^2+{r}-2 {\tau} {u^{r}} {u^{\tau}}\right)\\
  &\quad + {\mathcal{C}}\Bigl(7 {\pi_{Q}{}^{\phi}{}_{\phi}} {\tau} {u^{r}} {u^{\tau}} - 2 {\pi_{Q}{}^{\eta}{}_{\eta}} {r} \left(u^{\tau}\right)^2 + 3 {\pi_{Q}{}^{\phi}{}_{\phi}} {r} \left(u^{\tau}\right)^2\\
  &\quad - 2 {\Pi_{Q}} \left({r} \left(u^{r}\right)^2+{r}-2 {\tau} {u^{r}} {u^{\tau}}\right)\Bigr)\Bigr) \\
S_{9} &= {\pi_{Q}{}^{\eta}{}_{\eta}}\\
  &\quad + \frac{1}{6 {r} {T} {\tau} {u^{\tau}}}{D_{s}}\Bigl(- 2 {I_{31}} {\tau} {u^{r}} {u^{\tau}} + 4 {I_{31}} {r} \left(u^{\tau}\right)^2 + 2 {\mathcal{B}} {\nu^{r}} (2 {r} {u^{r}}-{\tau} {u^{\tau}})\\
  &\quad + {\mathcal{C}}\Bigl(- 2 {\Pi_{Q}} {\tau} {u^{r}} {u^{\tau}} + 3 {\pi_{Q}{}^{\eta}{}_{\eta}} {\tau} {u^{r}} {u^{\tau}} - 2 {\pi_{Q}{}^{\phi}{}_{\phi}} {\tau} {u^{r}} {u^{\tau}} + 4 {\Pi_{Q}} {r} \left(u^{\tau}\right)^2\\
  &\quad + 7 {\pi_{Q}{}^{\eta}{}_{\eta}} {r} \left(u^{\tau}\right)^2\Bigr)\Bigr) \\
S_{10} &= {\Pi_{Q}}\\
  &\quad + \frac{1}{6 {r} {T} {\tau} {u^{\tau}}}{D_{s}}\Bigl(5 {\mathcal{B}} {\nu^{r}} ({r} {u^{r}}+{\tau} {u^{\tau}}) + 5 {I_{31}} \left({r} \left(u^{r}\right)^2+{r}+{\tau} {u^{r}} {u^{\tau}}\right)\\
  &\quad + {\mathcal{C}} \left(5 {\Pi_{Q}} \left({r} \left(u^{r}\right)^2+{r}+{\tau} {u^{r}} {u^{\tau}}\right)+2 {\pi_{Q}{}^{\eta}{}_{\eta}} {r} \left(u^{\tau}\right)^2+2 {\pi_{Q}{}^{\phi}{}_{\phi}} {\tau} {u^{r}} {u^{\tau}}\right)\Bigr)
\end{align*}
\subsection{2+1D Cartesian}

\subsubsection{S3, charm first moment on a background}

\texttt{two\_\allowbreak{}d\_\allowbreak{}HQ\_\allowbreak{}BG\_\allowbreak{}consistent\_\allowbreak{}matrix} in \texttt{src/Matrix/2d\_\allowbreak{}viscous\_\allowbreak{}HQ\_\allowbreak{}BG.jl}, fields:

\[
\phi = ({\alpha}, {\nu^{x}}, {\nu^{y}})
\]
\begin{align*}
(A_t)_{11} &= {(\partial_{\alpha} n)} {u^{\tau}} \\
(A_t)_{12} &= \frac{{u^{x}}}{{u^{\tau}}} \\
(A_t)_{13} &= \frac{{u^{y}}}{{u^{\tau}}} \\
(A_t)_{21} &= {(\partial_{\alpha} n)} {D_{s}} {u^{\tau}} {u^{x}} \\
(A_t)_{22} &= \frac{{D_{s}} {h} {u^{\tau}}}{{T}} \\
(A_t)_{31} &= {(\partial_{\alpha} n)} {D_{s}} {u^{\tau}} {u^{y}} \\
(A_t)_{33} &= \frac{{D_{s}} {h} {u^{\tau}}}{{T}} \\
(A_x)_{11} &= {(\partial_{\alpha} n)} {u^{x}} \\
(A_x)_{12} &= 1 \\
(A_x)_{21} &= {(\partial_{\alpha} n)} {D_{s}} \left(\left(u^{x}\right)^2+1\right) \\
(A_x)_{22} &= \frac{{D_{s}} {h} {u^{x}}}{{T}} \\
(A_x)_{31} &= {(\partial_{\alpha} n)} {D_{s}} {u^{x}} {u^{y}} \\
(A_x)_{33} &= \frac{{D_{s}} {h} {u^{x}}}{{T}} \\
(A_y)_{11} &= {(\partial_{\alpha} n)} {u^{y}} \\
(A_y)_{13} &= 1 \\
(A_y)_{21} &= {(\partial_{\alpha} n)} {D_{s}} {u^{x}} {u^{y}} \\
(A_y)_{22} &= \frac{{D_{s}} {h} {u^{y}}}{{T}} \\
(A_y)_{31} &= {(\partial_{\alpha} n)} {D_{s}} \left(\left(u^{y}\right)^2+1\right) \\
(A_y)_{33} &= \frac{{D_{s}} {h} {u^{y}}}{{T}} \\
S_{1} &= {(\partial_{x} u^{x})} {n} + {(\partial_{y} u^{y})} {n} + \frac{{(\partial_{\tau} u^{x})} {\nu^{x}}+{(\partial_{\tau} u^{y})} {\nu^{y}}}{{u^{\tau}}} + {(\partial_{T} n)} {(\partial_{\tau} T)} {u^{\tau}} + {(\partial_{T} n)} {(\partial_{x} T)} {u^{x}}\\
  &\quad + {(\partial_{T} n)} {(\partial_{y} T)} {u^{y}} + \frac{{n} ({(\partial_{\tau} u^{x})} {u^{x}}+{(\partial_{\tau} u^{y})} {u^{y}})}{{u^{\tau}}} - \frac{({(\partial_{\tau} u^{x})} {u^{x}}+{(\partial_{\tau} u^{y})} {u^{y}}) ({\nu^{x}} {u^{x}}+{\nu^{y}} {u^{y}})}{\left(u^{\tau}\right)^3}\\
  &\quad + \frac{{n} \left(u^{\tau}\right)^2+{\nu^{x}} {u^{x}}+{\nu^{y}} {u^{y}}}{{\tau} {u^{\tau}}} \\
S_{2} &= {\nu^{x}}\\
  &\quad + \frac{1}{{T} {u^{\tau}}}{D_{s}}\Bigl({h'(T)} {\nu^{x}} \left({(\partial_{\tau} T)} \left(u^{\tau}\right)^2+{u^{\tau}} ({(\partial_{x} T)} {u^{x}}+{(\partial_{y} T)} {u^{y}})\right)\\
  &\quad + {(\partial_{T} n)} {T}+{n}\Bigl({(\partial_{\tau} T)} \left(u^{\tau}\right)^2 {u^{x}}+{u^{\tau}} \left({(\partial_{x} T)} \left(u^{x}\right)^2+{(\partial_{x} T)}+{(\partial_{y} T)} {u^{x}} {u^{y}}\right)\Bigr)\Bigr)\\
  &\quad + \frac{1}{{T} {\tau} \left(u^{\tau}\right)^3}{D_{s}} {h}\Bigl({\tau}\Bigl(- {(\partial_{\tau} u^{y})} {\nu^{y}} \left(u^{\tau}\right)^2 {u^{x}} \left(\left(u^{x}\right)^2+1\right)\\
  &\quad + {(\partial_{\tau} u^{x})} \left(u^{\tau}\right)^2 \left({n} \left(u^{\tau}\right)^2+{\nu^{y}} \left(\left(u^{x}\right)^2+1\right) {u^{y}}\right)\\
  &\quad + {u^{\tau}}\Bigl({(\partial_{x} u^{x})} {n} \left(u^{\tau}\right)^2 {u^{x}} + {(\partial_{y} u^{x})} {n} \left(u^{\tau}\right)^2 {u^{y}}\\
  &\quad + {(\partial_{x} u^{x})} {\nu^{y}} \left(u^{x}\right)^3 {u^{y}}\\
  &\quad - {\nu^{y}} {u^{x}} \left(\left(u^{x}\right)^2+1\right) ({(\partial_{x} u^{y})} {u^{x}}+{(\partial_{y} u^{y})} {u^{y}})\\
  &\quad + {(\partial_{y} u^{x})} {\nu^{y}} \left(\left(u^{x}\right)^2+1\right) \left(\left(u^{y}\right)^2+1\right)\Bigr)\Bigr)\\
  &\quad + {\nu^{x}}\Bigl({(\partial_{x} u^{y})} {\tau} {u^{\tau}} \left(u^{x}\right)^3 {u^{y}} + \left(u^{x}\right)^4 ({(\partial_{\tau} u^{y})} {\tau} {u^{y}}+1)\\
  &\quad - {(\partial_{\tau} u^{x})} {\tau} \left(u^{\tau}\right)^2 {u^{x}} \left(\left(u^{y}\right)^2-1\right)\\
  &\quad - {(\partial_{y} u^{x})} {\tau} {u^{\tau}} {u^{x}} {u^{y}} \left(\left(u^{y}\right)^2+1\right)\\
  &\quad + \left(u^{y}\right)^2+1\Bigl(1 + 2 {(\partial_{x} u^{x})} {\tau} {u^{\tau}} + {(\partial_{y} u^{y})} {\tau} {u^{\tau}} + {(\partial_{\tau} u^{y})} {\tau} {u^{y}}\\
  &\quad + \left(u^{y}\right)^2\Bigr)\\
  &\quad + \left(u^{x}\right)^2\Bigl(2 + {(\partial_{x} u^{x})} {\tau} {u^{\tau}} + {(\partial_{y} u^{y})} {\tau} {u^{\tau}}\\
  &\quad + \left(u^{y}\right)^2 (-{(\partial_{x} u^{x})} {\tau} {u^{\tau}}+{(\partial_{y} u^{y})} {\tau} {u^{\tau}}+2)\\
  &\quad + {(\partial_{\tau} u^{y})} {\tau} {u^{y}} \left(\left(u^{y}\right)^2+2\right)\Bigr)\Bigr)\Bigr) \\
S_{3} &= {\nu^{y}}\\
  &\quad + \frac{1}{{T} {u^{\tau}}}{D_{s}}\Bigl({h'(T)} {\nu^{y}} \left({(\partial_{\tau} T)} \left(u^{\tau}\right)^2+{u^{\tau}} ({(\partial_{x} T)} {u^{x}}+{(\partial_{y} T)} {u^{y}})\right)\\
  &\quad + {(\partial_{T} n)} {T}+{n}\Bigl({(\partial_{\tau} T)} \left(u^{\tau}\right)^2 {u^{y}}+{u^{\tau}} \left({(\partial_{x} T)} {u^{x}} {u^{y}}+{(\partial_{y} T)} \left(u^{y}\right)^2+{(\partial_{y} T)}\right)\Bigr)\Bigr)\\
  &\quad + \frac{1}{{T} {\tau} \left(u^{\tau}\right)^3}{D_{s}} {h}\Bigl({\nu^{y}}\Bigl(1 + {(\partial_{x} u^{x})} {\tau} {u^{\tau}} + 2 {(\partial_{y} u^{y})} {\tau} {u^{\tau}} + {(\partial_{\tau} u^{y})} {\tau} {u^{y}}\\
  &\quad - {(\partial_{x} u^{y})} {\tau} {u^{\tau}} \left(u^{x}\right)^3 {u^{y}} + 2 \left(u^{y}\right)^2 + {(\partial_{x} u^{x})} {\tau} {u^{\tau}} \left(u^{y}\right)^2\\
  &\quad + {(\partial_{y} u^{y})} {\tau} {u^{\tau}} \left(u^{y}\right)^2 + {(\partial_{\tau} u^{y})} {\tau} \left(u^{y}\right)^3 + \left(u^{y}\right)^4\\
  &\quad + \left(u^{x}\right)^4 (1-{(\partial_{\tau} u^{y})} {\tau} {u^{y}}) + {(\partial_{\tau} u^{x})} {\tau} \left(u^{\tau}\right)^2 {u^{x}} \left(\left(u^{y}\right)^2+1\right)\\
  &\quad + {\tau} {u^{\tau}} {u^{x}} {u^{y}} \left({(\partial_{y} u^{x})} \left(u^{y}\right)^2-{(\partial_{x} u^{y})}\right)\\
  &\quad + \left(u^{x}\right)^2\Bigl(2 + {(\partial_{x} u^{x})} {\tau} {u^{\tau}} + 2 {(\partial_{y} u^{y})} {\tau} {u^{\tau}}\\
  &\quad + \left(u^{y}\right)^2 ({(\partial_{x} u^{x})} {\tau} {u^{\tau}}-{(\partial_{y} u^{y})} {\tau} {u^{\tau}}+2)\\
  &\quad - {(\partial_{\tau} u^{y})} {\tau} \left(u^{y}\right)^3\Bigr)\Bigr)\\
  &\quad + {\tau}\Bigl(- {(\partial_{\tau} u^{x})} {\nu^{x}} \left(u^{\tau}\right)^2 {u^{y}} \left(\left(u^{y}\right)^2+1\right)\\
  &\quad + {(\partial_{\tau} u^{y})} \left(u^{\tau}\right)^2 \left({n} \left(u^{\tau}\right)^2+{\nu^{x}} {u^{x}} \left(\left(u^{y}\right)^2+1\right)\right)\\
  &\quad + {u^{\tau}}\Bigl({(\partial_{x} u^{y})} {n} \left(u^{\tau}\right)^2 {u^{x}} + {(\partial_{y} u^{y})} {n} \left(u^{\tau}\right)^2 {u^{y}}\\
  &\quad + {(\partial_{y} u^{y})} {\nu^{x}} {u^{x}} \left(u^{y}\right)^3\\
  &\quad + {(\partial_{x} u^{y})} {\nu^{x}} \left(\left(u^{x}\right)^2+1\right) \left(\left(u^{y}\right)^2+1\right)\\
  &\quad - {\nu^{x}} {u^{y}} \left(\left(u^{y}\right)^2+1\right) ({(\partial_{x} u^{x})} {u^{x}}+{(\partial_{y} u^{x})} {u^{y}})\Bigr)\Bigr)\Bigr)
\end{align*}
\subsubsection{S4, charm first and second moment on a background}

\texttt{two\_\allowbreak{}d\_\allowbreak{}HQ\_\allowbreak{}BG\_\allowbreak{}m2\_\allowbreak{}matrix} in \texttt{src/Matrix/2d\_\allowbreak{}viscous\_\allowbreak{}HQ\_\allowbreak{}BG\_\allowbreak{}m2.jl}, fields:

\[
\phi = ({\alpha}, {\nu^{x}}, {\nu^{y}}, {\pi_{Q}^{yy}}, {\pi_{Q}{}^{\eta}{}_{\eta}}, {\pi_{Q}^{xy}}, {\Pi_{Q}})
\]
\begin{align*}
(A_t)_{11} &= {(\partial_{\alpha} n)} {u^{\tau}} \\
(A_t)_{12} &= \frac{{u^{x}}}{{u^{\tau}}} \\
(A_t)_{13} &= \frac{{u^{y}}}{{u^{\tau}}} \\
(A_t)_{21} &= {(\partial_{\alpha} n)} {D_{s}} {u^{\tau}} {u^{x}} \\
(A_t)_{22} &= \frac{{D_{s}} {h} {u^{\tau}}}{{T}} \\
(A_t)_{31} &= {(\partial_{\alpha} n)} {D_{s}} {u^{\tau}} {u^{y}} \\
(A_t)_{33} &= \frac{{D_{s}} {h} {u^{\tau}}}{{T}} \\
(A_t)_{42} &= -\frac{{\mathcal{B}} {D_{s}} {u^{x}} \left(\left(u^{y}\right)^2+1\right)}{3 {T} {u^{\tau}}} \\
(A_t)_{43} &= \frac{{\mathcal{B}} {D_{s}} {u^{y}} \left(3 \left(u^{x}\right)^2+2 \left(u^{y}\right)^2+2\right)}{3 {T} {u^{\tau}}} \\
(A_t)_{44} &= \frac{{\mathcal{C}} {D_{s}} {u^{\tau}}}{2 {T}} \\
(A_t)_{52} &= -\frac{{\mathcal{B}} {D_{s}} {u^{x}}}{3 {T} {u^{\tau}}} \\
(A_t)_{53} &= -\frac{{\mathcal{B}} {D_{s}} {u^{y}}}{3 {T} {u^{\tau}}} \\
(A_t)_{55} &= \frac{{\mathcal{C}} {D_{s}} {u^{\tau}}}{2 {T}} \\
(A_t)_{62} &= \frac{{\mathcal{B}} {D_{s}} {u^{y}} \left(\left(u^{x}\right)^2+3 \left(u^{y}\right)^2+3\right)}{6 {T} {u^{\tau}}} \\
(A_t)_{63} &= \frac{{\mathcal{B}} {D_{s}} {u^{x}} \left(3 \left(u^{x}\right)^2+\left(u^{y}\right)^2+3\right)}{6 {T} {u^{\tau}}} \\
(A_t)_{66} &= \frac{{\mathcal{C}} {D_{s}} {u^{\tau}}}{2 {T}} \\
(A_t)_{71} &= \frac{{(\partial_{\alpha} I_{31})} {D_{s}} {u^{\tau}}}{2 {T}} \\
(A_t)_{72} &= \frac{5 {\mathcal{B}} {D_{s}} {u^{x}}}{6 {T} {u^{\tau}}} \\
(A_t)_{73} &= \frac{5 {\mathcal{B}} {D_{s}} {u^{y}}}{6 {T} {u^{\tau}}} \\
(A_t)_{77} &= \frac{{\mathcal{C}} {D_{s}} {u^{\tau}}}{2 {T}} \\
(A_x)_{11} &= {(\partial_{\alpha} n)} {u^{x}} \\
(A_x)_{12} &= 1 \\
(A_x)_{21} &= {(\partial_{\alpha} n)} {D_{s}} \left(\left(u^{x}\right)^2+1\right) \\
(A_x)_{22} &= \frac{{D_{s}} {h} {u^{x}}}{{T}} \\
(A_x)_{31} &= {(\partial_{\alpha} n)} {D_{s}} {u^{x}} {u^{y}} \\
(A_x)_{33} &= \frac{{D_{s}} {h} {u^{x}}}{{T}} \\
(A_x)_{42} &= -\frac{{\mathcal{B}} {D_{s}} \left(\left(u^{y}\right)^2+1\right)}{3 {T}} \\
(A_x)_{43} &= \frac{{\mathcal{B}} {D_{s}} {u^{x}} {u^{y}}}{{T}} \\
(A_x)_{44} &= \frac{{\mathcal{C}} {D_{s}} {u^{x}}}{2 {T}} \\
(A_x)_{52} &= -\frac{{\mathcal{B}} {D_{s}}}{3 {T}} \\
(A_x)_{55} &= \frac{{\mathcal{C}} {D_{s}} {u^{x}}}{2 {T}} \\
(A_x)_{62} &= \frac{{\mathcal{B}} {D_{s}} {u^{x}} {u^{y}}}{6 {T}} \\
(A_x)_{63} &= \frac{{\mathcal{B}} {D_{s}} \left(\left(u^{x}\right)^2+1\right)}{2 {T}} \\
(A_x)_{66} &= \frac{{\mathcal{C}} {D_{s}} {u^{x}}}{2 {T}} \\
(A_x)_{71} &= \frac{{(\partial_{\alpha} I_{31})} {D_{s}} {u^{x}}}{2 {T}} \\
(A_x)_{72} &= \frac{5 {\mathcal{B}} {D_{s}}}{6 {T}} \\
(A_x)_{77} &= \frac{{\mathcal{C}} {D_{s}} {u^{x}}}{2 {T}} \\
(A_y)_{11} &= {(\partial_{\alpha} n)} {u^{y}} \\
(A_y)_{13} &= 1 \\
(A_y)_{21} &= {(\partial_{\alpha} n)} {D_{s}} {u^{x}} {u^{y}} \\
(A_y)_{22} &= \frac{{D_{s}} {h} {u^{y}}}{{T}} \\
(A_y)_{31} &= {(\partial_{\alpha} n)} {D_{s}} \left(\left(u^{y}\right)^2+1\right) \\
(A_y)_{33} &= \frac{{D_{s}} {h} {u^{y}}}{{T}} \\
(A_y)_{43} &= \frac{2 {\mathcal{B}} {D_{s}} \left(\left(u^{y}\right)^2+1\right)}{3 {T}} \\
(A_y)_{44} &= \frac{{\mathcal{C}} {D_{s}} {u^{y}}}{2 {T}} \\
(A_y)_{53} &= -\frac{{\mathcal{B}} {D_{s}}}{3 {T}} \\
(A_y)_{55} &= \frac{{\mathcal{C}} {D_{s}} {u^{y}}}{2 {T}} \\
(A_y)_{62} &= \frac{{\mathcal{B}} {D_{s}} \left(\left(u^{y}\right)^2+1\right)}{2 {T}} \\
(A_y)_{63} &= \frac{{\mathcal{B}} {D_{s}} {u^{x}} {u^{y}}}{6 {T}} \\
(A_y)_{66} &= \frac{{\mathcal{C}} {D_{s}} {u^{y}}}{2 {T}} \\
(A_y)_{71} &= \frac{{(\partial_{\alpha} I_{31})} {D_{s}} {u^{y}}}{2 {T}} \\
(A_y)_{73} &= \frac{5 {\mathcal{B}} {D_{s}}}{6 {T}} \\
(A_y)_{77} &= \frac{{\mathcal{C}} {D_{s}} {u^{y}}}{2 {T}} \\
S_{1} &= {(\partial_{x} u^{x})} {n} + {(\partial_{y} u^{y})} {n} + \frac{{(\partial_{\tau} u^{x})} {\nu^{x}}+{(\partial_{\tau} u^{y})} {\nu^{y}}}{{u^{\tau}}} + {(\partial_{T} n)} {(\partial_{\tau} T)} {u^{\tau}} + {(\partial_{T} n)} {(\partial_{x} T)} {u^{x}}\\
  &\quad + {(\partial_{T} n)} {(\partial_{y} T)} {u^{y}} + \frac{{n} ({(\partial_{\tau} u^{x})} {u^{x}}+{(\partial_{\tau} u^{y})} {u^{y}})}{{u^{\tau}}} - \frac{({(\partial_{\tau} u^{x})} {u^{x}}+{(\partial_{\tau} u^{y})} {u^{y}}) ({\nu^{x}} {u^{x}}+{\nu^{y}} {u^{y}})}{\left(u^{\tau}\right)^3}\\
  &\quad + \frac{{n} \left(u^{\tau}\right)^2+{\nu^{x}} {u^{x}}+{\nu^{y}} {u^{y}}}{{\tau} {u^{\tau}}} \\
S_{2} &= {\nu^{x}}\\
  &\quad + \frac{1}{{T} {u^{\tau}}}{D_{s}}\Bigl({h'(T)} {\nu^{x}} \left({(\partial_{\tau} T)} \left(u^{\tau}\right)^2+{u^{\tau}} ({(\partial_{x} T)} {u^{x}}+{(\partial_{y} T)} {u^{y}})\right)\\
  &\quad + {(\partial_{T} n)} {T}+{n}\Bigl({(\partial_{\tau} T)} \left(u^{\tau}\right)^2 {u^{x}}+{u^{\tau}} \left({(\partial_{x} T)} \left(u^{x}\right)^2+{(\partial_{x} T)}+{(\partial_{y} T)} {u^{x}} {u^{y}}\right)\Bigr)\Bigr)\\
  &\quad + \frac{1}{{T} {\tau} \left(u^{\tau}\right)^3}{D_{s}} {h}\Bigl({\tau}\Bigl(- {(\partial_{\tau} u^{y})} {\nu^{y}} \left(u^{\tau}\right)^2 {u^{x}} \left(\left(u^{x}\right)^2+1\right)\\
  &\quad + {(\partial_{\tau} u^{x})} \left(u^{\tau}\right)^2 \left({n} \left(u^{\tau}\right)^2+{\nu^{y}} \left(\left(u^{x}\right)^2+1\right) {u^{y}}\right)\\
  &\quad + {u^{\tau}}\Bigl({(\partial_{x} u^{x})} {n} \left(u^{\tau}\right)^2 {u^{x}} + {(\partial_{y} u^{x})} {n} \left(u^{\tau}\right)^2 {u^{y}}\\
  &\quad + {(\partial_{x} u^{x})} {\nu^{y}} \left(u^{x}\right)^3 {u^{y}}\\
  &\quad - {\nu^{y}} {u^{x}} \left(\left(u^{x}\right)^2+1\right) ({(\partial_{x} u^{y})} {u^{x}}+{(\partial_{y} u^{y})} {u^{y}})\\
  &\quad + {(\partial_{y} u^{x})} {\nu^{y}} \left(\left(u^{x}\right)^2+1\right) \left(\left(u^{y}\right)^2+1\right)\Bigr)\Bigr)\\
  &\quad + {\nu^{x}}\Bigl({(\partial_{x} u^{y})} {\tau} {u^{\tau}} \left(u^{x}\right)^3 {u^{y}} + \left(u^{x}\right)^4 ({(\partial_{\tau} u^{y})} {\tau} {u^{y}}+1)\\
  &\quad - {(\partial_{\tau} u^{x})} {\tau} \left(u^{\tau}\right)^2 {u^{x}} \left(\left(u^{y}\right)^2-1\right)\\
  &\quad - {(\partial_{y} u^{x})} {\tau} {u^{\tau}} {u^{x}} {u^{y}} \left(\left(u^{y}\right)^2+1\right)\\
  &\quad + \left(u^{y}\right)^2+1\Bigl(1 + 2 {(\partial_{x} u^{x})} {\tau} {u^{\tau}} + {(\partial_{y} u^{y})} {\tau} {u^{\tau}} + {(\partial_{\tau} u^{y})} {\tau} {u^{y}}\\
  &\quad + \left(u^{y}\right)^2\Bigr)\\
  &\quad + \left(u^{x}\right)^2\Bigl(2 + {(\partial_{x} u^{x})} {\tau} {u^{\tau}} + {(\partial_{y} u^{y})} {\tau} {u^{\tau}}\\
  &\quad + \left(u^{y}\right)^2 (-{(\partial_{x} u^{x})} {\tau} {u^{\tau}}+{(\partial_{y} u^{y})} {\tau} {u^{\tau}}+2)\\
  &\quad + {(\partial_{\tau} u^{y})} {\tau} {u^{y}} \left(\left(u^{y}\right)^2+2\right)\Bigr)\Bigr)\Bigr) \\
S_{3} &= {\nu^{y}}\\
  &\quad + \frac{1}{{T} {u^{\tau}}}{D_{s}}\Bigl({h'(T)} {\nu^{y}} \left({(\partial_{\tau} T)} \left(u^{\tau}\right)^2+{u^{\tau}} ({(\partial_{x} T)} {u^{x}}+{(\partial_{y} T)} {u^{y}})\right)\\
  &\quad + {(\partial_{T} n)} {T}+{n}\Bigl({(\partial_{\tau} T)} \left(u^{\tau}\right)^2 {u^{y}}+{u^{\tau}} \left({(\partial_{x} T)} {u^{x}} {u^{y}}+{(\partial_{y} T)} \left(u^{y}\right)^2+{(\partial_{y} T)}\right)\Bigr)\Bigr)\\
  &\quad + \frac{1}{{T} {\tau} \left(u^{\tau}\right)^3}{D_{s}} {h}\Bigl({\nu^{y}}\Bigl(1 + {(\partial_{x} u^{x})} {\tau} {u^{\tau}} + 2 {(\partial_{y} u^{y})} {\tau} {u^{\tau}} + {(\partial_{\tau} u^{y})} {\tau} {u^{y}}\\
  &\quad - {(\partial_{x} u^{y})} {\tau} {u^{\tau}} \left(u^{x}\right)^3 {u^{y}} + 2 \left(u^{y}\right)^2 + {(\partial_{x} u^{x})} {\tau} {u^{\tau}} \left(u^{y}\right)^2\\
  &\quad + {(\partial_{y} u^{y})} {\tau} {u^{\tau}} \left(u^{y}\right)^2 + {(\partial_{\tau} u^{y})} {\tau} \left(u^{y}\right)^3 + \left(u^{y}\right)^4\\
  &\quad + \left(u^{x}\right)^4 (1-{(\partial_{\tau} u^{y})} {\tau} {u^{y}}) + {(\partial_{\tau} u^{x})} {\tau} \left(u^{\tau}\right)^2 {u^{x}} \left(\left(u^{y}\right)^2+1\right)\\
  &\quad + {\tau} {u^{\tau}} {u^{x}} {u^{y}} \left({(\partial_{y} u^{x})} \left(u^{y}\right)^2-{(\partial_{x} u^{y})}\right)\\
  &\quad + \left(u^{x}\right)^2\Bigl(2 + {(\partial_{x} u^{x})} {\tau} {u^{\tau}} + 2 {(\partial_{y} u^{y})} {\tau} {u^{\tau}}\\
  &\quad + \left(u^{y}\right)^2 ({(\partial_{x} u^{x})} {\tau} {u^{\tau}}-{(\partial_{y} u^{y})} {\tau} {u^{\tau}}+2)\\
  &\quad - {(\partial_{\tau} u^{y})} {\tau} \left(u^{y}\right)^3\Bigr)\Bigr)\\
  &\quad + {\tau}\Bigl(- {(\partial_{\tau} u^{x})} {\nu^{x}} \left(u^{\tau}\right)^2 {u^{y}} \left(\left(u^{y}\right)^2+1\right)\\
  &\quad + {(\partial_{\tau} u^{y})} \left(u^{\tau}\right)^2 \left({n} \left(u^{\tau}\right)^2+{\nu^{x}} {u^{x}} \left(\left(u^{y}\right)^2+1\right)\right)\\
  &\quad + {u^{\tau}}\Bigl({(\partial_{x} u^{y})} {n} \left(u^{\tau}\right)^2 {u^{x}} + {(\partial_{y} u^{y})} {n} \left(u^{\tau}\right)^2 {u^{y}}\\
  &\quad + {(\partial_{y} u^{y})} {\nu^{x}} {u^{x}} \left(u^{y}\right)^3\\
  &\quad + {(\partial_{x} u^{y})} {\nu^{x}} \left(\left(u^{x}\right)^2+1\right) \left(\left(u^{y}\right)^2+1\right)\\
  &\quad - {\nu^{x}} {u^{y}} \left(\left(u^{y}\right)^2+1\right) ({(\partial_{x} u^{x})} {u^{x}}+{(\partial_{y} u^{x})} {u^{y}})\Bigr)\Bigr)\Bigr) \\
S_{4} &= \frac{1}{6 {T} {\tau} \left(u^{\tau}\right)^4}\Bigl(- 2 {D_{s}} {(\partial_{x} u^{x})} {I_{31}} {\tau} + 4 {D_{s}} {(\partial_{y} u^{y})} {I_{31}} {\tau} - 2 {\mathcal{B}'(T)} {D_{s}} {(\partial_{x} T)} {\nu^{x}} {\tau}\\
  &\quad + 4 {\mathcal{B}'(T)} {D_{s}} {(\partial_{y} T)} {\nu^{y}} {\tau} - 2 {\mathcal{C}} {D_{s}} {(\partial_{x} u^{x})} {\Pi_{Q}} {\tau} + 4 {\mathcal{C}} {D_{s}} {(\partial_{y} u^{y})} {\Pi_{Q}} {\tau}\\
  &\quad + 4 {\mathcal{C}} {D_{s}} {(\partial_{x} u^{y})} {\pi_{Q}^{xy}} {\tau} - 2 {\mathcal{C}} {D_{s}} {(\partial_{y} u^{x})} {\pi_{Q}^{xy}} {\tau} + 5 {\mathcal{C}} {D_{s}} {(\partial_{x} u^{x})} {\pi_{Q}^{yy}} {\tau}\\
  &\quad + 7 {\mathcal{C}} {D_{s}} {(\partial_{y} u^{y})} {\pi_{Q}^{yy}} {\tau} + 2 {\mathcal{C}} {D_{s}} {(\partial_{x} u^{x})} {\pi_{Q}{}^{\eta}{}_{\eta}} {\tau} + 6 {\pi_{Q}^{yy}} {T} {\tau} - 2 {D_{s}} {I_{31}} {u^{\tau}} - 2 {\mathcal{C}} {D_{s}} {\Pi_{Q}} {u^{\tau}}\\
  &\quad + 3 {\mathcal{C}} {D_{s}} {\pi_{Q}^{yy}} {u^{\tau}} - 2 {\mathcal{C}} {D_{s}} {\pi_{Q}{}^{\eta}{}_{\eta}} {u^{\tau}} - 2 {\mathcal{A}} {D_{s}} {(\partial_{\tau} u^{x})} {\nu^{x}} {\tau} {u^{\tau}} - 2 {\mathcal{B}} {D_{s}} {(\partial_{\tau} u^{x})} {\nu^{x}} {\tau} {u^{\tau}}\\
  &\quad + 4 {\mathcal{A}} {D_{s}} {(\partial_{\tau} u^{y})} {\nu^{y}} {\tau} {u^{\tau}} + 4 {\mathcal{B}} {D_{s}} {(\partial_{\tau} u^{y})} {\nu^{y}} {\tau} {u^{\tau}} + 3 {\mathcal{C}'(T)} {D_{s}} {(\partial_{\tau} T)} {\pi_{Q}^{yy}} {\tau} {u^{\tau}}\\
  &\quad - 2 {\mathcal{A}} {D_{s}} {(\partial_{x} u^{x})} {\nu^{x}} {\tau} {u^{x}} + 4 {\mathcal{A}} {D_{s}} {(\partial_{x} u^{y})} {\nu^{y}} {\tau} {u^{x}} + 6 {\mathcal{B}} {D_{s}} {(\partial_{x} u^{y})} {\nu^{y}} {\tau} {u^{x}}\\
  &\quad + 3 {\mathcal{C}'(T)} {D_{s}} {(\partial_{x} T)} {\pi_{Q}^{yy}} {\tau} {u^{x}} - 2 {\mathcal{B}} {D_{s}} {\nu^{x}} {u^{\tau}} {u^{x}} - 2 {D_{s}} {(\partial_{\tau} u^{x})} {I_{31}} {\tau} {u^{\tau}} {u^{x}}\\
  &\quad - 2 {\mathcal{B}'(T)} {D_{s}} {(\partial_{\tau} T)} {\nu^{x}} {\tau} {u^{\tau}} {u^{x}} - 2 {\mathcal{C}} {D_{s}} {(\partial_{\tau} u^{x})} {\Pi_{Q}} {\tau} {u^{\tau}} {u^{x}} + 4 {\mathcal{C}} {D_{s}} {(\partial_{\tau} u^{y})} {\pi_{Q}^{xy}} {\tau} {u^{\tau}} {u^{x}}\\
  &\quad + 5 {\mathcal{C}} {D_{s}} {(\partial_{\tau} u^{x})} {\pi_{Q}^{yy}} {\tau} {u^{\tau}} {u^{x}} + 2 {\mathcal{C}} {D_{s}} {(\partial_{\tau} u^{x})} {\pi_{Q}{}^{\eta}{}_{\eta}} {\tau} {u^{\tau}} {u^{x}} - 4 {D_{s}} {(\partial_{x} u^{x})} {I_{31}} {\tau} \left(u^{x}\right)^2\\
  &\quad + 8 {D_{s}} {(\partial_{y} u^{y})} {I_{31}} {\tau} \left(u^{x}\right)^2 - 4 {\mathcal{B}'(T)} {D_{s}} {(\partial_{x} T)} {\nu^{x}} {\tau} \left(u^{x}\right)^2\\
  &\quad + 8 {\mathcal{B}'(T)} {D_{s}} {(\partial_{y} T)} {\nu^{y}} {\tau} \left(u^{x}\right)^2 - 4 {\mathcal{C}} {D_{s}} {(\partial_{x} u^{x})} {\Pi_{Q}} {\tau} \left(u^{x}\right)^2 + 8 {\mathcal{C}} {D_{s}} {(\partial_{y} u^{y})} {\Pi_{Q}} {\tau} \left(u^{x}\right)^2\\
  &\quad + 8 {\mathcal{C}} {D_{s}} {(\partial_{x} u^{y})} {\pi_{Q}^{xy}} {\tau} \left(u^{x}\right)^2 - 2 {\mathcal{C}} {D_{s}} {(\partial_{y} u^{x})} {\pi_{Q}^{xy}} {\tau} \left(u^{x}\right)^2 + 10 {\mathcal{C}} {D_{s}} {(\partial_{x} u^{x})} {\pi_{Q}^{yy}} {\tau} \left(u^{x}\right)^2\\
  &\quad + 14 {\mathcal{C}} {D_{s}} {(\partial_{y} u^{y})} {\pi_{Q}^{yy}} {\tau} \left(u^{x}\right)^2 + 4 {\mathcal{C}} {D_{s}} {(\partial_{x} u^{x})} {\pi_{Q}{}^{\eta}{}_{\eta}} {\tau} \left(u^{x}\right)^2 + 12 {\pi_{Q}^{yy}} {T} {\tau} \left(u^{x}\right)^2\\
  &\quad - 4 {D_{s}} {I_{31}} {u^{\tau}} \left(u^{x}\right)^2 - 4 {\mathcal{C}} {D_{s}} {\Pi_{Q}} {u^{\tau}} \left(u^{x}\right)^2 + 6 {\mathcal{C}} {D_{s}} {\pi_{Q}^{yy}} {u^{\tau}} \left(u^{x}\right)^2 - 4 {\mathcal{C}} {D_{s}} {\pi_{Q}{}^{\eta}{}_{\eta}} {u^{\tau}} \left(u^{x}\right)^2\\
  &\quad - 2 {\mathcal{A}} {D_{s}} {(\partial_{\tau} u^{x})} {\nu^{x}} {\tau} {u^{\tau}} \left(u^{x}\right)^2 + 8 {\mathcal{A}} {D_{s}} {(\partial_{\tau} u^{y})} {\nu^{y}} {\tau} {u^{\tau}} \left(u^{x}\right)^2\\
  &\quad + 10 {\mathcal{B}} {D_{s}} {(\partial_{\tau} u^{y})} {\nu^{y}} {\tau} {u^{\tau}} \left(u^{x}\right)^2 + 6 {\mathcal{C}'(T)} {D_{s}} {(\partial_{\tau} T)} {\pi_{Q}^{yy}} {\tau} {u^{\tau}} \left(u^{x}\right)^2\\
  &\quad - 2 {\mathcal{A}} {D_{s}} {(\partial_{x} u^{x})} {\nu^{x}} {\tau} \left(u^{x}\right)^3 + 8 {\mathcal{A}} {D_{s}} {(\partial_{x} u^{y})} {\nu^{y}} {\tau} \left(u^{x}\right)^3 + 12 {\mathcal{B}} {D_{s}} {(\partial_{x} u^{y})} {\nu^{y}} {\tau} \left(u^{x}\right)^3\\
  &\quad + 6 {\mathcal{C}'(T)} {D_{s}} {(\partial_{x} T)} {\pi_{Q}^{yy}} {\tau} \left(u^{x}\right)^3 - 2 {\mathcal{B}} {D_{s}} {\nu^{x}} {u^{\tau}} \left(u^{x}\right)^3 - 2 {D_{s}} {(\partial_{\tau} u^{x})} {I_{31}} {\tau} {u^{\tau}} \left(u^{x}\right)^3\\
  &\quad - 2 {\mathcal{B}'(T)} {D_{s}} {(\partial_{\tau} T)} {\nu^{x}} {\tau} {u^{\tau}} \left(u^{x}\right)^3 - 2 {\mathcal{C}} {D_{s}} {(\partial_{\tau} u^{x})} {\Pi_{Q}} {\tau} {u^{\tau}} \left(u^{x}\right)^3\\
  &\quad + 4 {\mathcal{C}} {D_{s}} {(\partial_{\tau} u^{y})} {\pi_{Q}^{xy}} {\tau} {u^{\tau}} \left(u^{x}\right)^3 + 5 {\mathcal{C}} {D_{s}} {(\partial_{\tau} u^{x})} {\pi_{Q}^{yy}} {\tau} {u^{\tau}} \left(u^{x}\right)^3\\
  &\quad + 2 {\mathcal{C}} {D_{s}} {(\partial_{\tau} u^{x})} {\pi_{Q}{}^{\eta}{}_{\eta}} {\tau} {u^{\tau}} \left(u^{x}\right)^3 - 2 {D_{s}} {(\partial_{x} u^{x})} {I_{31}} {\tau} \left(u^{x}\right)^4 + 4 {D_{s}} {(\partial_{y} u^{y})} {I_{31}} {\tau} \left(u^{x}\right)^4\\
  &\quad - 2 {\mathcal{B}'(T)} {D_{s}} {(\partial_{x} T)} {\nu^{x}} {\tau} \left(u^{x}\right)^4 + 4 {\mathcal{B}'(T)} {D_{s}} {(\partial_{y} T)} {\nu^{y}} {\tau} \left(u^{x}\right)^4\\
  &\quad - 2 {\mathcal{C}} {D_{s}} {(\partial_{x} u^{x})} {\Pi_{Q}} {\tau} \left(u^{x}\right)^4 + 4 {\mathcal{C}} {D_{s}} {(\partial_{y} u^{y})} {\Pi_{Q}} {\tau} \left(u^{x}\right)^4 + 4 {\mathcal{C}} {D_{s}} {(\partial_{x} u^{y})} {\pi_{Q}^{xy}} {\tau} \left(u^{x}\right)^4\\
  &\quad + 5 {\mathcal{C}} {D_{s}} {(\partial_{x} u^{x})} {\pi_{Q}^{yy}} {\tau} \left(u^{x}\right)^4 + 7 {\mathcal{C}} {D_{s}} {(\partial_{y} u^{y})} {\pi_{Q}^{yy}} {\tau} \left(u^{x}\right)^4 + 2 {\mathcal{C}} {D_{s}} {(\partial_{x} u^{x})} {\pi_{Q}{}^{\eta}{}_{\eta}} {\tau} \left(u^{x}\right)^4\\
  &\quad + 6 {\pi_{Q}^{yy}} {T} {\tau} \left(u^{x}\right)^4 - 2 {D_{s}} {I_{31}} {u^{\tau}} \left(u^{x}\right)^4 - 2 {\mathcal{C}} {D_{s}} {\Pi_{Q}} {u^{\tau}} \left(u^{x}\right)^4 + 3 {\mathcal{C}} {D_{s}} {\pi_{Q}^{yy}} {u^{\tau}} \left(u^{x}\right)^4\\
  &\quad - 2 {\mathcal{C}} {D_{s}} {\pi_{Q}{}^{\eta}{}_{\eta}} {u^{\tau}} \left(u^{x}\right)^4 + 4 {\mathcal{A}} {D_{s}} {(\partial_{\tau} u^{y})} {\nu^{y}} {\tau} {u^{\tau}} \left(u^{x}\right)^4 + 6 {\mathcal{B}} {D_{s}} {(\partial_{\tau} u^{y})} {\nu^{y}} {\tau} {u^{\tau}} \left(u^{x}\right)^4\\
  &\quad + 3 {\mathcal{C}'(T)} {D_{s}} {(\partial_{\tau} T)} {\pi_{Q}^{yy}} {\tau} {u^{\tau}} \left(u^{x}\right)^4 + 4 {\mathcal{A}} {D_{s}} {(\partial_{x} u^{y})} {\nu^{y}} {\tau} \left(u^{x}\right)^5\\
  &\quad + 6 {\mathcal{B}} {D_{s}} {(\partial_{x} u^{y})} {\nu^{y}} {\tau} \left(u^{x}\right)^5 + 3 {\mathcal{C}'(T)} {D_{s}} {(\partial_{x} T)} {\pi_{Q}^{yy}} {\tau} \left(u^{x}\right)^5 - 2 {\mathcal{A}} {D_{s}} {(\partial_{y} u^{x})} {\nu^{x}} {\tau} {u^{y}}\\
  &\quad - 6 {\mathcal{B}} {D_{s}} {(\partial_{y} u^{x})} {\nu^{x}} {\tau} {u^{y}} + 4 {\mathcal{A}} {D_{s}} {(\partial_{y} u^{y})} {\nu^{y}} {\tau} {u^{y}} + 3 {\mathcal{C}'(T)} {D_{s}} {(\partial_{y} T)} {\pi_{Q}^{yy}} {\tau} {u^{y}}\\
  &\quad - 2 {\mathcal{B}} {D_{s}} {\nu^{y}} {u^{\tau}} {u^{y}} + 4 {D_{s}} {(\partial_{\tau} u^{y})} {I_{31}} {\tau} {u^{\tau}} {u^{y}} + 4 {\mathcal{B}'(T)} {D_{s}} {(\partial_{\tau} T)} {\nu^{y}} {\tau} {u^{\tau}} {u^{y}}\\
  &\quad + 4 {\mathcal{C}} {D_{s}} {(\partial_{\tau} u^{y})} {\Pi_{Q}} {\tau} {u^{\tau}} {u^{y}} - 8 {\mathcal{C}} {D_{s}} {(\partial_{\tau} u^{x})} {\pi_{Q}^{xy}} {\tau} {u^{\tau}} {u^{y}} + {\mathcal{C}} {D_{s}} {(\partial_{\tau} u^{y})} {\pi_{Q}^{yy}} {\tau} {u^{\tau}} {u^{y}}\\
  &\quad + 6 {D_{s}} {(\partial_{x} u^{y})} {I_{31}} {\tau} {u^{x}} {u^{y}} + 6 {\mathcal{B}'(T)} {D_{s}} {(\partial_{x} T)} {\nu^{y}} {\tau} {u^{x}} {u^{y}} + 6 {\mathcal{C}} {D_{s}} {(\partial_{x} u^{y})} {\Pi_{Q}} {\tau} {u^{x}} {u^{y}}\\
  &\quad - 8 {\mathcal{C}} {D_{s}} {(\partial_{x} u^{x})} {\pi_{Q}^{xy}} {\tau} {u^{x}} {u^{y}} + 2 {\mathcal{C}} {D_{s}} {(\partial_{y} u^{y})} {\pi_{Q}^{xy}} {\tau} {u^{x}} {u^{y}} - 8 {\mathcal{C}} {D_{s}} {(\partial_{x} u^{y})} {\pi_{Q}^{yy}} {\tau} {u^{x}} {u^{y}}\\
  &\quad + 2 {\mathcal{C}} {D_{s}} {(\partial_{y} u^{x})} {\pi_{Q}^{yy}} {\tau} {u^{x}} {u^{y}} - 2 {\mathcal{C}} {D_{s}} {(\partial_{x} u^{y})} {\pi_{Q}{}^{\eta}{}_{\eta}} {\tau} {u^{x}} {u^{y}} + 2 {\mathcal{A}} {D_{s}} {(\partial_{\tau} u^{y})} {\nu^{x}} {\tau} {u^{\tau}} {u^{x}} {u^{y}}\\
  &\quad + 2 {\mathcal{B}} {D_{s}} {(\partial_{\tau} u^{y})} {\nu^{x}} {\tau} {u^{\tau}} {u^{x}} {u^{y}} + 2 {\mathcal{A}} {D_{s}} {(\partial_{\tau} u^{x})} {\nu^{y}} {\tau} {u^{\tau}} {u^{x}} {u^{y}} + 2 {\mathcal{B}} {D_{s}} {(\partial_{\tau} u^{x})} {\nu^{y}} {\tau} {u^{\tau}} {u^{x}} {u^{y}}\\
  &\quad + 2 {\mathcal{A}} {D_{s}} {(\partial_{x} u^{y})} {\nu^{x}} {\tau} \left(u^{x}\right)^2 {u^{y}} - 2 {\mathcal{A}} {D_{s}} {(\partial_{y} u^{x})} {\nu^{x}} {\tau} \left(u^{x}\right)^2 {u^{y}}\\
  &\quad - 6 {\mathcal{B}} {D_{s}} {(\partial_{y} u^{x})} {\nu^{x}} {\tau} \left(u^{x}\right)^2 {u^{y}} + 2 {\mathcal{A}} {D_{s}} {(\partial_{x} u^{x})} {\nu^{y}} {\tau} \left(u^{x}\right)^2 {u^{y}}\\
  &\quad + 8 {\mathcal{A}} {D_{s}} {(\partial_{y} u^{y})} {\nu^{y}} {\tau} \left(u^{x}\right)^2 {u^{y}} + 6 {\mathcal{C}'(T)} {D_{s}} {(\partial_{y} T)} {\pi_{Q}^{yy}} {\tau} \left(u^{x}\right)^2 {u^{y}}\\
  &\quad - 2 {\mathcal{B}} {D_{s}} {\nu^{y}} {u^{\tau}} \left(u^{x}\right)^2 {u^{y}} + 10 {D_{s}} {(\partial_{\tau} u^{y})} {I_{31}} {\tau} {u^{\tau}} \left(u^{x}\right)^2 {u^{y}}\\
  &\quad + 10 {\mathcal{B}'(T)} {D_{s}} {(\partial_{\tau} T)} {\nu^{y}} {\tau} {u^{\tau}} \left(u^{x}\right)^2 {u^{y}} + 10 {\mathcal{C}} {D_{s}} {(\partial_{\tau} u^{y})} {\Pi_{Q}} {\tau} {u^{\tau}} \left(u^{x}\right)^2 {u^{y}}\\
  &\quad - 8 {\mathcal{C}} {D_{s}} {(\partial_{\tau} u^{x})} {\pi_{Q}^{xy}} {\tau} {u^{\tau}} \left(u^{x}\right)^2 {u^{y}} - 7 {\mathcal{C}} {D_{s}} {(\partial_{\tau} u^{y})} {\pi_{Q}^{yy}} {\tau} {u^{\tau}} \left(u^{x}\right)^2 {u^{y}}\\
  &\quad - 2 {\mathcal{C}} {D_{s}} {(\partial_{\tau} u^{y})} {\pi_{Q}{}^{\eta}{}_{\eta}} {\tau} {u^{\tau}} \left(u^{x}\right)^2 {u^{y}} + 12 {D_{s}} {(\partial_{x} u^{y})} {I_{31}} {\tau} \left(u^{x}\right)^3 {u^{y}}\\
  &\quad + 12 {\mathcal{B}'(T)} {D_{s}} {(\partial_{x} T)} {\nu^{y}} {\tau} \left(u^{x}\right)^3 {u^{y}} + 12 {\mathcal{C}} {D_{s}} {(\partial_{x} u^{y})} {\Pi_{Q}} {\tau} \left(u^{x}\right)^3 {u^{y}}\\
  &\quad - 8 {\mathcal{C}} {D_{s}} {(\partial_{x} u^{x})} {\pi_{Q}^{xy}} {\tau} \left(u^{x}\right)^3 {u^{y}} + 2 {\mathcal{C}} {D_{s}} {(\partial_{y} u^{y})} {\pi_{Q}^{xy}} {\tau} \left(u^{x}\right)^3 {u^{y}}\\
  &\quad - 16 {\mathcal{C}} {D_{s}} {(\partial_{x} u^{y})} {\pi_{Q}^{yy}} {\tau} \left(u^{x}\right)^3 {u^{y}} + 2 {\mathcal{C}} {D_{s}} {(\partial_{y} u^{x})} {\pi_{Q}^{yy}} {\tau} \left(u^{x}\right)^3 {u^{y}}\\
  &\quad - 4 {\mathcal{C}} {D_{s}} {(\partial_{x} u^{y})} {\pi_{Q}{}^{\eta}{}_{\eta}} {\tau} \left(u^{x}\right)^3 {u^{y}} + 2 {\mathcal{A}} {D_{s}} {(\partial_{\tau} u^{y})} {\nu^{x}} {\tau} {u^{\tau}} \left(u^{x}\right)^3 {u^{y}}\\
  &\quad + 2 {\mathcal{A}} {D_{s}} {(\partial_{\tau} u^{x})} {\nu^{y}} {\tau} {u^{\tau}} \left(u^{x}\right)^3 {u^{y}} + 2 {\mathcal{A}} {D_{s}} {(\partial_{x} u^{y})} {\nu^{x}} {\tau} \left(u^{x}\right)^4 {u^{y}}\\
  &\quad + 2 {\mathcal{A}} {D_{s}} {(\partial_{x} u^{x})} {\nu^{y}} {\tau} \left(u^{x}\right)^4 {u^{y}} + 4 {\mathcal{A}} {D_{s}} {(\partial_{y} u^{y})} {\nu^{y}} {\tau} \left(u^{x}\right)^4 {u^{y}}\\
  &\quad + 3 {\mathcal{C}'(T)} {D_{s}} {(\partial_{y} T)} {\pi_{Q}^{yy}} {\tau} \left(u^{x}\right)^4 {u^{y}} + 6 {D_{s}} {(\partial_{\tau} u^{y})} {I_{31}} {\tau} {u^{\tau}} \left(u^{x}\right)^4 {u^{y}}\\
  &\quad + 6 {\mathcal{B}'(T)} {D_{s}} {(\partial_{\tau} T)} {\nu^{y}} {\tau} {u^{\tau}} \left(u^{x}\right)^4 {u^{y}} + 6 {\mathcal{C}} {D_{s}} {(\partial_{\tau} u^{y})} {\Pi_{Q}} {\tau} {u^{\tau}} \left(u^{x}\right)^4 {u^{y}}\\
  &\quad - 8 {\mathcal{C}} {D_{s}} {(\partial_{\tau} u^{y})} {\pi_{Q}^{yy}} {\tau} {u^{\tau}} \left(u^{x}\right)^4 {u^{y}} - 2 {\mathcal{C}} {D_{s}} {(\partial_{\tau} u^{y})} {\pi_{Q}{}^{\eta}{}_{\eta}} {\tau} {u^{\tau}} \left(u^{x}\right)^4 {u^{y}}\\
  &\quad + 6 {D_{s}} {(\partial_{x} u^{y})} {I_{31}} {\tau} \left(u^{x}\right)^5 {u^{y}} + 6 {\mathcal{B}'(T)} {D_{s}} {(\partial_{x} T)} {\nu^{y}} {\tau} \left(u^{x}\right)^5 {u^{y}}\\
  &\quad + 6 {\mathcal{C}} {D_{s}} {(\partial_{x} u^{y})} {\Pi_{Q}} {\tau} \left(u^{x}\right)^5 {u^{y}} - 8 {\mathcal{C}} {D_{s}} {(\partial_{x} u^{y})} {\pi_{Q}^{yy}} {\tau} \left(u^{x}\right)^5 {u^{y}}\\
  &\quad - 2 {\mathcal{C}} {D_{s}} {(\partial_{x} u^{y})} {\pi_{Q}{}^{\eta}{}_{\eta}} {\tau} \left(u^{x}\right)^5 {u^{y}} - 6 {D_{s}} {(\partial_{x} u^{x})} {I_{31}} {\tau} \left(u^{y}\right)^2 + 12 {D_{s}} {(\partial_{y} u^{y})} {I_{31}} {\tau} \left(u^{y}\right)^2\\
  &\quad - 6 {\mathcal{B}'(T)} {D_{s}} {(\partial_{x} T)} {\nu^{x}} {\tau} \left(u^{y}\right)^2 + 12 {\mathcal{B}'(T)} {D_{s}} {(\partial_{y} T)} {\nu^{y}} {\tau} \left(u^{y}\right)^2\\
  &\quad - 6 {\mathcal{C}} {D_{s}} {(\partial_{x} u^{x})} {\Pi_{Q}} {\tau} \left(u^{y}\right)^2 + 12 {\mathcal{C}} {D_{s}} {(\partial_{y} u^{y})} {\Pi_{Q}} {\tau} \left(u^{y}\right)^2 + 8 {\mathcal{C}} {D_{s}} {(\partial_{x} u^{y})} {\pi_{Q}^{xy}} {\tau} \left(u^{y}\right)^2\\
  &\quad - 12 {\mathcal{C}} {D_{s}} {(\partial_{y} u^{x})} {\pi_{Q}^{xy}} {\tau} \left(u^{y}\right)^2 + 10 {\mathcal{C}} {D_{s}} {(\partial_{x} u^{x})} {\pi_{Q}^{yy}} {\tau} \left(u^{y}\right)^2\\
  &\quad + 8 {\mathcal{C}} {D_{s}} {(\partial_{y} u^{y})} {\pi_{Q}^{yy}} {\tau} \left(u^{y}\right)^2 + 6 {\mathcal{C}} {D_{s}} {(\partial_{x} u^{x})} {\pi_{Q}{}^{\eta}{}_{\eta}} {\tau} \left(u^{y}\right)^2 + 12 {\pi_{Q}^{yy}} {T} {\tau} \left(u^{y}\right)^2\\
  &\quad - 6 {D_{s}} {I_{31}} {u^{\tau}} \left(u^{y}\right)^2 - 6 {\mathcal{C}} {D_{s}} {\Pi_{Q}} {u^{\tau}} \left(u^{y}\right)^2 + 6 {\mathcal{C}} {D_{s}} {\pi_{Q}^{yy}} {u^{\tau}} \left(u^{y}\right)^2 - 6 {\mathcal{C}} {D_{s}} {\pi_{Q}{}^{\eta}{}_{\eta}} {u^{\tau}} \left(u^{y}\right)^2\\
  &\quad - 6 {\mathcal{A}} {D_{s}} {(\partial_{\tau} u^{x})} {\nu^{x}} {\tau} {u^{\tau}} \left(u^{y}\right)^2 - 10 {\mathcal{B}} {D_{s}} {(\partial_{\tau} u^{x})} {\nu^{x}} {\tau} {u^{\tau}} \left(u^{y}\right)^2\\
  &\quad + 8 {\mathcal{A}} {D_{s}} {(\partial_{\tau} u^{y})} {\nu^{y}} {\tau} {u^{\tau}} \left(u^{y}\right)^2 + 4 {\mathcal{B}} {D_{s}} {(\partial_{\tau} u^{y})} {\nu^{y}} {\tau} {u^{\tau}} \left(u^{y}\right)^2\\
  &\quad + 6 {\mathcal{C}'(T)} {D_{s}} {(\partial_{\tau} T)} {\pi_{Q}^{yy}} {\tau} {u^{\tau}} \left(u^{y}\right)^2 - 6 {\mathcal{A}} {D_{s}} {(\partial_{x} u^{x})} {\nu^{x}} {\tau} {u^{x}} \left(u^{y}\right)^2\\
  &\quad - 6 {\mathcal{B}} {D_{s}} {(\partial_{x} u^{x})} {\nu^{x}} {\tau} {u^{x}} \left(u^{y}\right)^2 + 2 {\mathcal{A}} {D_{s}} {(\partial_{y} u^{y})} {\nu^{x}} {\tau} {u^{x}} \left(u^{y}\right)^2\\
  &\quad + 6 {\mathcal{B}} {D_{s}} {(\partial_{y} u^{y})} {\nu^{x}} {\tau} {u^{x}} \left(u^{y}\right)^2 + 8 {\mathcal{A}} {D_{s}} {(\partial_{x} u^{y})} {\nu^{y}} {\tau} {u^{x}} \left(u^{y}\right)^2\\
  &\quad + 6 {\mathcal{B}} {D_{s}} {(\partial_{x} u^{y})} {\nu^{y}} {\tau} {u^{x}} \left(u^{y}\right)^2 + 2 {\mathcal{A}} {D_{s}} {(\partial_{y} u^{x})} {\nu^{y}} {\tau} {u^{x}} \left(u^{y}\right)^2\\
  &\quad + 6 {\mathcal{B}} {D_{s}} {(\partial_{y} u^{x})} {\nu^{y}} {\tau} {u^{x}} \left(u^{y}\right)^2 + 6 {\mathcal{C}'(T)} {D_{s}} {(\partial_{x} T)} {\pi_{Q}^{yy}} {\tau} {u^{x}} \left(u^{y}\right)^2\\
  &\quad - 4 {\mathcal{B}} {D_{s}} {\nu^{x}} {u^{\tau}} {u^{x}} \left(u^{y}\right)^2 - 4 {D_{s}} {(\partial_{\tau} u^{x})} {I_{31}} {\tau} {u^{\tau}} {u^{x}} \left(u^{y}\right)^2\\
  &\quad - 4 {\mathcal{B}'(T)} {D_{s}} {(\partial_{\tau} T)} {\nu^{x}} {\tau} {u^{\tau}} {u^{x}} \left(u^{y}\right)^2 - 4 {\mathcal{C}} {D_{s}} {(\partial_{\tau} u^{x})} {\Pi_{Q}} {\tau} {u^{\tau}} {u^{x}} \left(u^{y}\right)^2\\
  &\quad + 12 {\mathcal{C}} {D_{s}} {(\partial_{\tau} u^{y})} {\pi_{Q}^{xy}} {\tau} {u^{\tau}} {u^{x}} \left(u^{y}\right)^2 + 13 {\mathcal{C}} {D_{s}} {(\partial_{\tau} u^{x})} {\pi_{Q}^{yy}} {\tau} {u^{\tau}} {u^{x}} \left(u^{y}\right)^2\\
  &\quad + 4 {\mathcal{C}} {D_{s}} {(\partial_{\tau} u^{x})} {\pi_{Q}{}^{\eta}{}_{\eta}} {\tau} {u^{\tau}} {u^{x}} \left(u^{y}\right)^2 - 8 {D_{s}} {(\partial_{x} u^{x})} {I_{31}} {\tau} \left(u^{x}\right)^2 \left(u^{y}\right)^2\\
  &\quad + 16 {D_{s}} {(\partial_{y} u^{y})} {I_{31}} {\tau} \left(u^{x}\right)^2 \left(u^{y}\right)^2 - 8 {\mathcal{B}'(T)} {D_{s}} {(\partial_{x} T)} {\nu^{x}} {\tau} \left(u^{x}\right)^2 \left(u^{y}\right)^2\\
  &\quad + 16 {\mathcal{B}'(T)} {D_{s}} {(\partial_{y} T)} {\nu^{y}} {\tau} \left(u^{x}\right)^2 \left(u^{y}\right)^2 - 8 {\mathcal{C}} {D_{s}} {(\partial_{x} u^{x})} {\Pi_{Q}} {\tau} \left(u^{x}\right)^2 \left(u^{y}\right)^2\\
  &\quad + 16 {\mathcal{C}} {D_{s}} {(\partial_{y} u^{y})} {\Pi_{Q}} {\tau} \left(u^{x}\right)^2 \left(u^{y}\right)^2 + 16 {\mathcal{C}} {D_{s}} {(\partial_{x} u^{y})} {\pi_{Q}^{xy}} {\tau} \left(u^{x}\right)^2 \left(u^{y}\right)^2\\
  &\quad - 10 {\mathcal{C}} {D_{s}} {(\partial_{y} u^{x})} {\pi_{Q}^{xy}} {\tau} \left(u^{x}\right)^2 \left(u^{y}\right)^2 + 18 {\mathcal{C}} {D_{s}} {(\partial_{x} u^{x})} {\pi_{Q}^{yy}} {\tau} \left(u^{x}\right)^2 \left(u^{y}\right)^2\\
  &\quad + 8 {\mathcal{C}} {D_{s}} {(\partial_{x} u^{x})} {\pi_{Q}{}^{\eta}{}_{\eta}} {\tau} \left(u^{x}\right)^2 \left(u^{y}\right)^2 + 12 {\pi_{Q}^{yy}} {T} {\tau} \left(u^{x}\right)^2 \left(u^{y}\right)^2\\
  &\quad - 8 {D_{s}} {I_{31}} {u^{\tau}} \left(u^{x}\right)^2 \left(u^{y}\right)^2 - 8 {\mathcal{C}} {D_{s}} {\Pi_{Q}} {u^{\tau}} \left(u^{x}\right)^2 \left(u^{y}\right)^2 + 6 {\mathcal{C}} {D_{s}} {\pi_{Q}^{yy}} {u^{\tau}} \left(u^{x}\right)^2 \left(u^{y}\right)^2\\
  &\quad - 8 {\mathcal{C}} {D_{s}} {\pi_{Q}{}^{\eta}{}_{\eta}} {u^{\tau}} \left(u^{x}\right)^2 \left(u^{y}\right)^2 - 4 {\mathcal{A}} {D_{s}} {(\partial_{\tau} u^{x})} {\nu^{x}} {\tau} {u^{\tau}} \left(u^{x}\right)^2 \left(u^{y}\right)^2\\
  &\quad - 6 {\mathcal{B}} {D_{s}} {(\partial_{\tau} u^{x})} {\nu^{x}} {\tau} {u^{\tau}} \left(u^{x}\right)^2 \left(u^{y}\right)^2 + 6 {\mathcal{A}} {D_{s}} {(\partial_{\tau} u^{y})} {\nu^{y}} {\tau} {u^{\tau}} \left(u^{x}\right)^2 \left(u^{y}\right)^2\\
  &\quad - 2 {\mathcal{B}} {D_{s}} {(\partial_{\tau} u^{y})} {\nu^{y}} {\tau} {u^{\tau}} \left(u^{x}\right)^2 \left(u^{y}\right)^2 + 6 {\mathcal{C}'(T)} {D_{s}} {(\partial_{\tau} T)} {\pi_{Q}^{yy}} {\tau} {u^{\tau}} \left(u^{x}\right)^2 \left(u^{y}\right)^2\\
  &\quad - 4 {\mathcal{A}} {D_{s}} {(\partial_{x} u^{x})} {\nu^{x}} {\tau} \left(u^{x}\right)^3 \left(u^{y}\right)^2 - 6 {\mathcal{B}} {D_{s}} {(\partial_{x} u^{x})} {\nu^{x}} {\tau} \left(u^{x}\right)^3 \left(u^{y}\right)^2\\
  &\quad + 2 {\mathcal{A}} {D_{s}} {(\partial_{y} u^{y})} {\nu^{x}} {\tau} \left(u^{x}\right)^3 \left(u^{y}\right)^2 + 6 {\mathcal{B}} {D_{s}} {(\partial_{y} u^{y})} {\nu^{x}} {\tau} \left(u^{x}\right)^3 \left(u^{y}\right)^2\\
  &\quad + 6 {\mathcal{A}} {D_{s}} {(\partial_{x} u^{y})} {\nu^{y}} {\tau} \left(u^{x}\right)^3 \left(u^{y}\right)^2 + 2 {\mathcal{A}} {D_{s}} {(\partial_{y} u^{x})} {\nu^{y}} {\tau} \left(u^{x}\right)^3 \left(u^{y}\right)^2\\
  &\quad + 6 {\mathcal{B}} {D_{s}} {(\partial_{y} u^{x})} {\nu^{y}} {\tau} \left(u^{x}\right)^3 \left(u^{y}\right)^2 + 6 {\mathcal{C}'(T)} {D_{s}} {(\partial_{x} T)} {\pi_{Q}^{yy}} {\tau} \left(u^{x}\right)^3 \left(u^{y}\right)^2\\
  &\quad - 2 {\mathcal{B}} {D_{s}} {\nu^{x}} {u^{\tau}} \left(u^{x}\right)^3 \left(u^{y}\right)^2 - 2 {D_{s}} {(\partial_{\tau} u^{x})} {I_{31}} {\tau} {u^{\tau}} \left(u^{x}\right)^3 \left(u^{y}\right)^2\\
  &\quad - 2 {\mathcal{B}'(T)} {D_{s}} {(\partial_{\tau} T)} {\nu^{x}} {\tau} {u^{\tau}} \left(u^{x}\right)^3 \left(u^{y}\right)^2 - 2 {\mathcal{C}} {D_{s}} {(\partial_{\tau} u^{x})} {\Pi_{Q}} {\tau} {u^{\tau}} \left(u^{x}\right)^3 \left(u^{y}\right)^2\\
  &\quad + 8 {\mathcal{C}} {D_{s}} {(\partial_{\tau} u^{y})} {\pi_{Q}^{xy}} {\tau} {u^{\tau}} \left(u^{x}\right)^3 \left(u^{y}\right)^2 + 8 {\mathcal{C}} {D_{s}} {(\partial_{\tau} u^{x})} {\pi_{Q}^{yy}} {\tau} {u^{\tau}} \left(u^{x}\right)^3 \left(u^{y}\right)^2\\
  &\quad + 2 {\mathcal{C}} {D_{s}} {(\partial_{\tau} u^{x})} {\pi_{Q}{}^{\eta}{}_{\eta}} {\tau} {u^{\tau}} \left(u^{x}\right)^3 \left(u^{y}\right)^2 - 2 {D_{s}} {(\partial_{x} u^{x})} {I_{31}} {\tau} \left(u^{x}\right)^4 \left(u^{y}\right)^2\\
  &\quad + 4 {D_{s}} {(\partial_{y} u^{y})} {I_{31}} {\tau} \left(u^{x}\right)^4 \left(u^{y}\right)^2 - 2 {\mathcal{B}'(T)} {D_{s}} {(\partial_{x} T)} {\nu^{x}} {\tau} \left(u^{x}\right)^4 \left(u^{y}\right)^2\\
  &\quad + 4 {\mathcal{B}'(T)} {D_{s}} {(\partial_{y} T)} {\nu^{y}} {\tau} \left(u^{x}\right)^4 \left(u^{y}\right)^2 - 2 {\mathcal{C}} {D_{s}} {(\partial_{x} u^{x})} {\Pi_{Q}} {\tau} \left(u^{x}\right)^4 \left(u^{y}\right)^2\\
  &\quad + 4 {\mathcal{C}} {D_{s}} {(\partial_{y} u^{y})} {\Pi_{Q}} {\tau} \left(u^{x}\right)^4 \left(u^{y}\right)^2 + 8 {\mathcal{C}} {D_{s}} {(\partial_{x} u^{y})} {\pi_{Q}^{xy}} {\tau} \left(u^{x}\right)^4 \left(u^{y}\right)^2\\
  &\quad + 8 {\mathcal{C}} {D_{s}} {(\partial_{x} u^{x})} {\pi_{Q}^{yy}} {\tau} \left(u^{x}\right)^4 \left(u^{y}\right)^2 - 8 {\mathcal{C}} {D_{s}} {(\partial_{y} u^{y})} {\pi_{Q}^{yy}} {\tau} \left(u^{x}\right)^4 \left(u^{y}\right)^2\\
  &\quad + 2 {\mathcal{C}} {D_{s}} {(\partial_{x} u^{x})} {\pi_{Q}{}^{\eta}{}_{\eta}} {\tau} \left(u^{x}\right)^4 \left(u^{y}\right)^2 - 2 {D_{s}} {I_{31}} {u^{\tau}} \left(u^{x}\right)^4 \left(u^{y}\right)^2\\
  &\quad - 2 {\mathcal{C}} {D_{s}} {\Pi_{Q}} {u^{\tau}} \left(u^{x}\right)^4 \left(u^{y}\right)^2 - 2 {\mathcal{C}} {D_{s}} {\pi_{Q}{}^{\eta}{}_{\eta}} {u^{\tau}} \left(u^{x}\right)^4 \left(u^{y}\right)^2\\
  &\quad - 2 {\mathcal{A}} {D_{s}} {(\partial_{\tau} u^{y})} {\nu^{y}} {\tau} {u^{\tau}} \left(u^{x}\right)^4 \left(u^{y}\right)^2 - 6 {\mathcal{B}} {D_{s}} {(\partial_{\tau} u^{y})} {\nu^{y}} {\tau} {u^{\tau}} \left(u^{x}\right)^4 \left(u^{y}\right)^2\\
  &\quad - 2 {\mathcal{A}} {D_{s}} {(\partial_{x} u^{y})} {\nu^{y}} {\tau} \left(u^{x}\right)^5 \left(u^{y}\right)^2 - 6 {\mathcal{B}} {D_{s}} {(\partial_{x} u^{y})} {\nu^{y}} {\tau} \left(u^{x}\right)^5 \left(u^{y}\right)^2\\
  &\quad - 6 {\mathcal{A}} {D_{s}} {(\partial_{y} u^{x})} {\nu^{x}} {\tau} \left(u^{y}\right)^3 - 18 {\mathcal{B}} {D_{s}} {(\partial_{y} u^{x})} {\nu^{x}} {\tau} \left(u^{y}\right)^3 + 8 {\mathcal{A}} {D_{s}} {(\partial_{y} u^{y})} {\nu^{y}} {\tau} \left(u^{y}\right)^3\\
  &\quad + 6 {\mathcal{C}'(T)} {D_{s}} {(\partial_{y} T)} {\pi_{Q}^{yy}} {\tau} \left(u^{y}\right)^3 - 4 {\mathcal{B}} {D_{s}} {\nu^{y}} {u^{\tau}} \left(u^{y}\right)^3 + 8 {D_{s}} {(\partial_{\tau} u^{y})} {I_{31}} {\tau} {u^{\tau}} \left(u^{y}\right)^3\\
  &\quad + 8 {\mathcal{B}'(T)} {D_{s}} {(\partial_{\tau} T)} {\nu^{y}} {\tau} {u^{\tau}} \left(u^{y}\right)^3 + 8 {\mathcal{C}} {D_{s}} {(\partial_{\tau} u^{y})} {\Pi_{Q}} {\tau} {u^{\tau}} \left(u^{y}\right)^3\\
  &\quad - 16 {\mathcal{C}} {D_{s}} {(\partial_{\tau} u^{x})} {\pi_{Q}^{xy}} {\tau} {u^{\tau}} \left(u^{y}\right)^3 + {\mathcal{C}} {D_{s}} {(\partial_{\tau} u^{y})} {\pi_{Q}^{yy}} {\tau} {u^{\tau}} \left(u^{y}\right)^3\\
  &\quad + 12 {D_{s}} {(\partial_{x} u^{y})} {I_{31}} {\tau} {u^{x}} \left(u^{y}\right)^3 + 12 {\mathcal{B}'(T)} {D_{s}} {(\partial_{x} T)} {\nu^{y}} {\tau} {u^{x}} \left(u^{y}\right)^3\\
  &\quad + 12 {\mathcal{C}} {D_{s}} {(\partial_{x} u^{y})} {\Pi_{Q}} {\tau} {u^{x}} \left(u^{y}\right)^3 - 16 {\mathcal{C}} {D_{s}} {(\partial_{x} u^{x})} {\pi_{Q}^{xy}} {\tau} {u^{x}} \left(u^{y}\right)^3\\
  &\quad + 10 {\mathcal{C}} {D_{s}} {(\partial_{y} u^{y})} {\pi_{Q}^{xy}} {\tau} {u^{x}} \left(u^{y}\right)^3 - 8 {\mathcal{C}} {D_{s}} {(\partial_{x} u^{y})} {\pi_{Q}^{yy}} {\tau} {u^{x}} \left(u^{y}\right)^3\\
  &\quad + 10 {\mathcal{C}} {D_{s}} {(\partial_{y} u^{x})} {\pi_{Q}^{yy}} {\tau} {u^{x}} \left(u^{y}\right)^3 - 4 {\mathcal{C}} {D_{s}} {(\partial_{x} u^{y})} {\pi_{Q}{}^{\eta}{}_{\eta}} {\tau} {u^{x}} \left(u^{y}\right)^3\\
  &\quad + 4 {\mathcal{A}} {D_{s}} {(\partial_{\tau} u^{y})} {\nu^{x}} {\tau} {u^{\tau}} {u^{x}} \left(u^{y}\right)^3 + 8 {\mathcal{B}} {D_{s}} {(\partial_{\tau} u^{y})} {\nu^{x}} {\tau} {u^{\tau}} {u^{x}} \left(u^{y}\right)^3\\
  &\quad + 4 {\mathcal{A}} {D_{s}} {(\partial_{\tau} u^{x})} {\nu^{y}} {\tau} {u^{\tau}} {u^{x}} \left(u^{y}\right)^3 + 8 {\mathcal{B}} {D_{s}} {(\partial_{\tau} u^{x})} {\nu^{y}} {\tau} {u^{\tau}} {u^{x}} \left(u^{y}\right)^3\\
  &\quad + 4 {\mathcal{A}} {D_{s}} {(\partial_{x} u^{y})} {\nu^{x}} {\tau} \left(u^{x}\right)^2 \left(u^{y}\right)^3 + 6 {\mathcal{B}} {D_{s}} {(\partial_{x} u^{y})} {\nu^{x}} {\tau} \left(u^{x}\right)^2 \left(u^{y}\right)^3\\
  &\quad - 4 {\mathcal{A}} {D_{s}} {(\partial_{y} u^{x})} {\nu^{x}} {\tau} \left(u^{x}\right)^2 \left(u^{y}\right)^3 - 12 {\mathcal{B}} {D_{s}} {(\partial_{y} u^{x})} {\nu^{x}} {\tau} \left(u^{x}\right)^2 \left(u^{y}\right)^3\\
  &\quad + 4 {\mathcal{A}} {D_{s}} {(\partial_{x} u^{x})} {\nu^{y}} {\tau} \left(u^{x}\right)^2 \left(u^{y}\right)^3 + 6 {\mathcal{B}} {D_{s}} {(\partial_{x} u^{x})} {\nu^{y}} {\tau} \left(u^{x}\right)^2 \left(u^{y}\right)^3\\
  &\quad + 6 {\mathcal{A}} {D_{s}} {(\partial_{y} u^{y})} {\nu^{y}} {\tau} \left(u^{x}\right)^2 \left(u^{y}\right)^3 - 6 {\mathcal{B}} {D_{s}} {(\partial_{y} u^{y})} {\nu^{y}} {\tau} \left(u^{x}\right)^2 \left(u^{y}\right)^3\\
  &\quad + 6 {\mathcal{C}'(T)} {D_{s}} {(\partial_{y} T)} {\pi_{Q}^{yy}} {\tau} \left(u^{x}\right)^2 \left(u^{y}\right)^3 - 2 {\mathcal{B}} {D_{s}} {\nu^{y}} {u^{\tau}} \left(u^{x}\right)^2 \left(u^{y}\right)^3\\
  &\quad + 10 {D_{s}} {(\partial_{\tau} u^{y})} {I_{31}} {\tau} {u^{\tau}} \left(u^{x}\right)^2 \left(u^{y}\right)^3 + 10 {\mathcal{B}'(T)} {D_{s}} {(\partial_{\tau} T)} {\nu^{y}} {\tau} {u^{\tau}} \left(u^{x}\right)^2 \left(u^{y}\right)^3\\
  &\quad + 10 {\mathcal{C}} {D_{s}} {(\partial_{\tau} u^{y})} {\Pi_{Q}} {\tau} {u^{\tau}} \left(u^{x}\right)^2 \left(u^{y}\right)^3 - 8 {\mathcal{C}} {D_{s}} {(\partial_{\tau} u^{x})} {\pi_{Q}^{xy}} {\tau} {u^{\tau}} \left(u^{x}\right)^2 \left(u^{y}\right)^3\\
  &\quad - 8 {\mathcal{C}} {D_{s}} {(\partial_{\tau} u^{y})} {\pi_{Q}^{yy}} {\tau} {u^{\tau}} \left(u^{x}\right)^2 \left(u^{y}\right)^3 - 2 {\mathcal{C}} {D_{s}} {(\partial_{\tau} u^{y})} {\pi_{Q}{}^{\eta}{}_{\eta}} {\tau} {u^{\tau}} \left(u^{x}\right)^2 \left(u^{y}\right)^3\\
  &\quad + 12 {D_{s}} {(\partial_{x} u^{y})} {I_{31}} {\tau} \left(u^{x}\right)^3 \left(u^{y}\right)^3 + 12 {\mathcal{B}'(T)} {D_{s}} {(\partial_{x} T)} {\nu^{y}} {\tau} \left(u^{x}\right)^3 \left(u^{y}\right)^3\\
  &\quad + 12 {\mathcal{C}} {D_{s}} {(\partial_{x} u^{y})} {\Pi_{Q}} {\tau} \left(u^{x}\right)^3 \left(u^{y}\right)^3 - 8 {\mathcal{C}} {D_{s}} {(\partial_{x} u^{x})} {\pi_{Q}^{xy}} {\tau} \left(u^{x}\right)^3 \left(u^{y}\right)^3\\
  &\quad + 8 {\mathcal{C}} {D_{s}} {(\partial_{y} u^{y})} {\pi_{Q}^{xy}} {\tau} \left(u^{x}\right)^3 \left(u^{y}\right)^3 - 8 {\mathcal{C}} {D_{s}} {(\partial_{x} u^{y})} {\pi_{Q}^{yy}} {\tau} \left(u^{x}\right)^3 \left(u^{y}\right)^3\\
  &\quad + 8 {\mathcal{C}} {D_{s}} {(\partial_{y} u^{x})} {\pi_{Q}^{yy}} {\tau} \left(u^{x}\right)^3 \left(u^{y}\right)^3 - 4 {\mathcal{C}} {D_{s}} {(\partial_{x} u^{y})} {\pi_{Q}{}^{\eta}{}_{\eta}} {\tau} \left(u^{x}\right)^3 \left(u^{y}\right)^3\\
  &\quad + 2 {\mathcal{A}} {D_{s}} {(\partial_{\tau} u^{y})} {\nu^{x}} {\tau} {u^{\tau}} \left(u^{x}\right)^3 \left(u^{y}\right)^3 + 6 {\mathcal{B}} {D_{s}} {(\partial_{\tau} u^{y})} {\nu^{x}} {\tau} {u^{\tau}} \left(u^{x}\right)^3 \left(u^{y}\right)^3\\
  &\quad + 2 {\mathcal{A}} {D_{s}} {(\partial_{\tau} u^{x})} {\nu^{y}} {\tau} {u^{\tau}} \left(u^{x}\right)^3 \left(u^{y}\right)^3 + 6 {\mathcal{B}} {D_{s}} {(\partial_{\tau} u^{x})} {\nu^{y}} {\tau} {u^{\tau}} \left(u^{x}\right)^3 \left(u^{y}\right)^3\\
  &\quad + 2 {\mathcal{A}} {D_{s}} {(\partial_{x} u^{y})} {\nu^{x}} {\tau} \left(u^{x}\right)^4 \left(u^{y}\right)^3 + 6 {\mathcal{B}} {D_{s}} {(\partial_{x} u^{y})} {\nu^{x}} {\tau} \left(u^{x}\right)^4 \left(u^{y}\right)^3\\
  &\quad + 2 {\mathcal{A}} {D_{s}} {(\partial_{x} u^{x})} {\nu^{y}} {\tau} \left(u^{x}\right)^4 \left(u^{y}\right)^3 + 6 {\mathcal{B}} {D_{s}} {(\partial_{x} u^{x})} {\nu^{y}} {\tau} \left(u^{x}\right)^4 \left(u^{y}\right)^3\\
  &\quad - 2 {\mathcal{A}} {D_{s}} {(\partial_{y} u^{y})} {\nu^{y}} {\tau} \left(u^{x}\right)^4 \left(u^{y}\right)^3 - 6 {\mathcal{B}} {D_{s}} {(\partial_{y} u^{y})} {\nu^{y}} {\tau} \left(u^{x}\right)^4 \left(u^{y}\right)^3\\
  &\quad - 6 {D_{s}} {(\partial_{x} u^{x})} {I_{31}} {\tau} \left(u^{y}\right)^4 + 12 {D_{s}} {(\partial_{y} u^{y})} {I_{31}} {\tau} \left(u^{y}\right)^4 - 6 {\mathcal{B}'(T)} {D_{s}} {(\partial_{x} T)} {\nu^{x}} {\tau} \left(u^{y}\right)^4\\
  &\quad + 12 {\mathcal{B}'(T)} {D_{s}} {(\partial_{y} T)} {\nu^{y}} {\tau} \left(u^{y}\right)^4 - 6 {\mathcal{C}} {D_{s}} {(\partial_{x} u^{x})} {\Pi_{Q}} {\tau} \left(u^{y}\right)^4\\
  &\quad + 12 {\mathcal{C}} {D_{s}} {(\partial_{y} u^{y})} {\Pi_{Q}} {\tau} \left(u^{y}\right)^4 + 4 {\mathcal{C}} {D_{s}} {(\partial_{x} u^{y})} {\pi_{Q}^{xy}} {\tau} \left(u^{y}\right)^4 - 18 {\mathcal{C}} {D_{s}} {(\partial_{y} u^{x})} {\pi_{Q}^{xy}} {\tau} \left(u^{y}\right)^4\\
  &\quad + 5 {\mathcal{C}} {D_{s}} {(\partial_{x} u^{x})} {\pi_{Q}^{yy}} {\tau} \left(u^{y}\right)^4 + {\mathcal{C}} {D_{s}} {(\partial_{y} u^{y})} {\pi_{Q}^{yy}} {\tau} \left(u^{y}\right)^4 + 6 {\mathcal{C}} {D_{s}} {(\partial_{x} u^{x})} {\pi_{Q}{}^{\eta}{}_{\eta}} {\tau} \left(u^{y}\right)^4\\
  &\quad + 6 {\pi_{Q}^{yy}} {T} {\tau} \left(u^{y}\right)^4 - 6 {D_{s}} {I_{31}} {u^{\tau}} \left(u^{y}\right)^4 - 6 {\mathcal{C}} {D_{s}} {\Pi_{Q}} {u^{\tau}} \left(u^{y}\right)^4 + 3 {\mathcal{C}} {D_{s}} {\pi_{Q}^{yy}} {u^{\tau}} \left(u^{y}\right)^4\\
  &\quad - 6 {\mathcal{C}} {D_{s}} {\pi_{Q}{}^{\eta}{}_{\eta}} {u^{\tau}} \left(u^{y}\right)^4 - 6 {\mathcal{A}} {D_{s}} {(\partial_{\tau} u^{x})} {\nu^{x}} {\tau} {u^{\tau}} \left(u^{y}\right)^4 - 14 {\mathcal{B}} {D_{s}} {(\partial_{\tau} u^{x})} {\nu^{x}} {\tau} {u^{\tau}} \left(u^{y}\right)^4\\
  &\quad + 4 {\mathcal{A}} {D_{s}} {(\partial_{\tau} u^{y})} {\nu^{y}} {\tau} {u^{\tau}} \left(u^{y}\right)^4 + 3 {\mathcal{C}'(T)} {D_{s}} {(\partial_{\tau} T)} {\pi_{Q}^{yy}} {\tau} {u^{\tau}} \left(u^{y}\right)^4\\
  &\quad - 6 {\mathcal{A}} {D_{s}} {(\partial_{x} u^{x})} {\nu^{x}} {\tau} {u^{x}} \left(u^{y}\right)^4 - 12 {\mathcal{B}} {D_{s}} {(\partial_{x} u^{x})} {\nu^{x}} {\tau} {u^{x}} \left(u^{y}\right)^4\\
  &\quad + 4 {\mathcal{A}} {D_{s}} {(\partial_{y} u^{y})} {\nu^{x}} {\tau} {u^{x}} \left(u^{y}\right)^4 + 12 {\mathcal{B}} {D_{s}} {(\partial_{y} u^{y})} {\nu^{x}} {\tau} {u^{x}} \left(u^{y}\right)^4\\
  &\quad + 4 {\mathcal{A}} {D_{s}} {(\partial_{x} u^{y})} {\nu^{y}} {\tau} {u^{x}} \left(u^{y}\right)^4 + 4 {\mathcal{A}} {D_{s}} {(\partial_{y} u^{x})} {\nu^{y}} {\tau} {u^{x}} \left(u^{y}\right)^4\\
  &\quad + 12 {\mathcal{B}} {D_{s}} {(\partial_{y} u^{x})} {\nu^{y}} {\tau} {u^{x}} \left(u^{y}\right)^4 + 3 {\mathcal{C}'(T)} {D_{s}} {(\partial_{x} T)} {\pi_{Q}^{yy}} {\tau} {u^{x}} \left(u^{y}\right)^4\\
  &\quad - 2 {\mathcal{B}} {D_{s}} {\nu^{x}} {u^{\tau}} {u^{x}} \left(u^{y}\right)^4 - 2 {D_{s}} {(\partial_{\tau} u^{x})} {I_{31}} {\tau} {u^{\tau}} {u^{x}} \left(u^{y}\right)^4\\
  &\quad - 2 {\mathcal{B}'(T)} {D_{s}} {(\partial_{\tau} T)} {\nu^{x}} {\tau} {u^{\tau}} {u^{x}} \left(u^{y}\right)^4 - 2 {\mathcal{C}} {D_{s}} {(\partial_{\tau} u^{x})} {\Pi_{Q}} {\tau} {u^{\tau}} {u^{x}} \left(u^{y}\right)^4\\
  &\quad + 8 {\mathcal{C}} {D_{s}} {(\partial_{\tau} u^{y})} {\pi_{Q}^{xy}} {\tau} {u^{\tau}} {u^{x}} \left(u^{y}\right)^4 + 8 {\mathcal{C}} {D_{s}} {(\partial_{\tau} u^{x})} {\pi_{Q}^{yy}} {\tau} {u^{\tau}} {u^{x}} \left(u^{y}\right)^4\\
  &\quad + 2 {\mathcal{C}} {D_{s}} {(\partial_{\tau} u^{x})} {\pi_{Q}{}^{\eta}{}_{\eta}} {\tau} {u^{\tau}} {u^{x}} \left(u^{y}\right)^4 - 4 {D_{s}} {(\partial_{x} u^{x})} {I_{31}} {\tau} \left(u^{x}\right)^2 \left(u^{y}\right)^4\\
  &\quad + 8 {D_{s}} {(\partial_{y} u^{y})} {I_{31}} {\tau} \left(u^{x}\right)^2 \left(u^{y}\right)^4 - 4 {\mathcal{B}'(T)} {D_{s}} {(\partial_{x} T)} {\nu^{x}} {\tau} \left(u^{x}\right)^2 \left(u^{y}\right)^4\\
  &\quad + 8 {\mathcal{B}'(T)} {D_{s}} {(\partial_{y} T)} {\nu^{y}} {\tau} \left(u^{x}\right)^2 \left(u^{y}\right)^4 - 4 {\mathcal{C}} {D_{s}} {(\partial_{x} u^{x})} {\Pi_{Q}} {\tau} \left(u^{x}\right)^2 \left(u^{y}\right)^4\\
  &\quad + 8 {\mathcal{C}} {D_{s}} {(\partial_{y} u^{y})} {\Pi_{Q}} {\tau} \left(u^{x}\right)^2 \left(u^{y}\right)^4 + 8 {\mathcal{C}} {D_{s}} {(\partial_{x} u^{y})} {\pi_{Q}^{xy}} {\tau} \left(u^{x}\right)^2 \left(u^{y}\right)^4\\
  &\quad - 8 {\mathcal{C}} {D_{s}} {(\partial_{y} u^{x})} {\pi_{Q}^{xy}} {\tau} \left(u^{x}\right)^2 \left(u^{y}\right)^4 + 8 {\mathcal{C}} {D_{s}} {(\partial_{x} u^{x})} {\pi_{Q}^{yy}} {\tau} \left(u^{x}\right)^2 \left(u^{y}\right)^4\\
  &\quad - 8 {\mathcal{C}} {D_{s}} {(\partial_{y} u^{y})} {\pi_{Q}^{yy}} {\tau} \left(u^{x}\right)^2 \left(u^{y}\right)^4 + 4 {\mathcal{C}} {D_{s}} {(\partial_{x} u^{x})} {\pi_{Q}{}^{\eta}{}_{\eta}} {\tau} \left(u^{x}\right)^2 \left(u^{y}\right)^4\\
  &\quad - 4 {D_{s}} {I_{31}} {u^{\tau}} \left(u^{x}\right)^2 \left(u^{y}\right)^4 - 4 {\mathcal{C}} {D_{s}} {\Pi_{Q}} {u^{\tau}} \left(u^{x}\right)^2 \left(u^{y}\right)^4 - 4 {\mathcal{C}} {D_{s}} {\pi_{Q}{}^{\eta}{}_{\eta}} {u^{\tau}} \left(u^{x}\right)^2 \left(u^{y}\right)^4\\
  &\quad - 2 {\mathcal{A}} {D_{s}} {(\partial_{\tau} u^{x})} {\nu^{x}} {\tau} {u^{\tau}} \left(u^{x}\right)^2 \left(u^{y}\right)^4 - 6 {\mathcal{B}} {D_{s}} {(\partial_{\tau} u^{x})} {\nu^{x}} {\tau} {u^{\tau}} \left(u^{x}\right)^2 \left(u^{y}\right)^4\\
  &\quad - 2 {\mathcal{A}} {D_{s}} {(\partial_{\tau} u^{y})} {\nu^{y}} {\tau} {u^{\tau}} \left(u^{x}\right)^2 \left(u^{y}\right)^4 - 6 {\mathcal{B}} {D_{s}} {(\partial_{\tau} u^{y})} {\nu^{y}} {\tau} {u^{\tau}} \left(u^{x}\right)^2 \left(u^{y}\right)^4\\
  &\quad - 2 {\mathcal{A}} {D_{s}} {(\partial_{x} u^{x})} {\nu^{x}} {\tau} \left(u^{x}\right)^3 \left(u^{y}\right)^4 - 6 {\mathcal{B}} {D_{s}} {(\partial_{x} u^{x})} {\nu^{x}} {\tau} \left(u^{x}\right)^3 \left(u^{y}\right)^4\\
  &\quad + 2 {\mathcal{A}} {D_{s}} {(\partial_{y} u^{y})} {\nu^{x}} {\tau} \left(u^{x}\right)^3 \left(u^{y}\right)^4 + 6 {\mathcal{B}} {D_{s}} {(\partial_{y} u^{y})} {\nu^{x}} {\tau} \left(u^{x}\right)^3 \left(u^{y}\right)^4\\
  &\quad - 2 {\mathcal{A}} {D_{s}} {(\partial_{x} u^{y})} {\nu^{y}} {\tau} \left(u^{x}\right)^3 \left(u^{y}\right)^4 - 6 {\mathcal{B}} {D_{s}} {(\partial_{x} u^{y})} {\nu^{y}} {\tau} \left(u^{x}\right)^3 \left(u^{y}\right)^4\\
  &\quad + 2 {\mathcal{A}} {D_{s}} {(\partial_{y} u^{x})} {\nu^{y}} {\tau} \left(u^{x}\right)^3 \left(u^{y}\right)^4 + 6 {\mathcal{B}} {D_{s}} {(\partial_{y} u^{x})} {\nu^{y}} {\tau} \left(u^{x}\right)^3 \left(u^{y}\right)^4\\
  &\quad - 6 {\mathcal{A}} {D_{s}} {(\partial_{y} u^{x})} {\nu^{x}} {\tau} \left(u^{y}\right)^5 - 18 {\mathcal{B}} {D_{s}} {(\partial_{y} u^{x})} {\nu^{x}} {\tau} \left(u^{y}\right)^5 + 4 {\mathcal{A}} {D_{s}} {(\partial_{y} u^{y})} {\nu^{y}} {\tau} \left(u^{y}\right)^5\\
  &\quad + 3 {\mathcal{C}'(T)} {D_{s}} {(\partial_{y} T)} {\pi_{Q}^{yy}} {\tau} \left(u^{y}\right)^5 - 2 {\mathcal{B}} {D_{s}} {\nu^{y}} {u^{\tau}} \left(u^{y}\right)^5 + 4 {D_{s}} {(\partial_{\tau} u^{y})} {I_{31}} {\tau} {u^{\tau}} \left(u^{y}\right)^5\\
  &\quad + 4 {\mathcal{B}'(T)} {D_{s}} {(\partial_{\tau} T)} {\nu^{y}} {\tau} {u^{\tau}} \left(u^{y}\right)^5 + 4 {\mathcal{C}} {D_{s}} {(\partial_{\tau} u^{y})} {\Pi_{Q}} {\tau} {u^{\tau}} \left(u^{y}\right)^5\\
  &\quad - 8 {\mathcal{C}} {D_{s}} {(\partial_{\tau} u^{x})} {\pi_{Q}^{xy}} {\tau} {u^{\tau}} \left(u^{y}\right)^5 + 6 {D_{s}} {(\partial_{x} u^{y})} {I_{31}} {\tau} {u^{x}} \left(u^{y}\right)^5\\
  &\quad + 6 {\mathcal{B}'(T)} {D_{s}} {(\partial_{x} T)} {\nu^{y}} {\tau} {u^{x}} \left(u^{y}\right)^5 + 6 {\mathcal{C}} {D_{s}} {(\partial_{x} u^{y})} {\Pi_{Q}} {\tau} {u^{x}} \left(u^{y}\right)^5\\
  &\quad - 8 {\mathcal{C}} {D_{s}} {(\partial_{x} u^{x})} {\pi_{Q}^{xy}} {\tau} {u^{x}} \left(u^{y}\right)^5 + 8 {\mathcal{C}} {D_{s}} {(\partial_{y} u^{y})} {\pi_{Q}^{xy}} {\tau} {u^{x}} \left(u^{y}\right)^5\\
  &\quad + 8 {\mathcal{C}} {D_{s}} {(\partial_{y} u^{x})} {\pi_{Q}^{yy}} {\tau} {u^{x}} \left(u^{y}\right)^5 - 2 {\mathcal{C}} {D_{s}} {(\partial_{x} u^{y})} {\pi_{Q}{}^{\eta}{}_{\eta}} {\tau} {u^{x}} \left(u^{y}\right)^5\\
  &\quad + 2 {\mathcal{A}} {D_{s}} {(\partial_{\tau} u^{y})} {\nu^{x}} {\tau} {u^{\tau}} {u^{x}} \left(u^{y}\right)^5 + 6 {\mathcal{B}} {D_{s}} {(\partial_{\tau} u^{y})} {\nu^{x}} {\tau} {u^{\tau}} {u^{x}} \left(u^{y}\right)^5\\
  &\quad + 2 {\mathcal{A}} {D_{s}} {(\partial_{\tau} u^{x})} {\nu^{y}} {\tau} {u^{\tau}} {u^{x}} \left(u^{y}\right)^5 + 6 {\mathcal{B}} {D_{s}} {(\partial_{\tau} u^{x})} {\nu^{y}} {\tau} {u^{\tau}} {u^{x}} \left(u^{y}\right)^5\\
  &\quad + 2 {\mathcal{A}} {D_{s}} {(\partial_{x} u^{y})} {\nu^{x}} {\tau} \left(u^{x}\right)^2 \left(u^{y}\right)^5 + 6 {\mathcal{B}} {D_{s}} {(\partial_{x} u^{y})} {\nu^{x}} {\tau} \left(u^{x}\right)^2 \left(u^{y}\right)^5\\
  &\quad - 2 {\mathcal{A}} {D_{s}} {(\partial_{y} u^{x})} {\nu^{x}} {\tau} \left(u^{x}\right)^2 \left(u^{y}\right)^5 - 6 {\mathcal{B}} {D_{s}} {(\partial_{y} u^{x})} {\nu^{x}} {\tau} \left(u^{x}\right)^2 \left(u^{y}\right)^5\\
  &\quad + 2 {\mathcal{A}} {D_{s}} {(\partial_{x} u^{x})} {\nu^{y}} {\tau} \left(u^{x}\right)^2 \left(u^{y}\right)^5 + 6 {\mathcal{B}} {D_{s}} {(\partial_{x} u^{x})} {\nu^{y}} {\tau} \left(u^{x}\right)^2 \left(u^{y}\right)^5\\
  &\quad - 2 {\mathcal{A}} {D_{s}} {(\partial_{y} u^{y})} {\nu^{y}} {\tau} \left(u^{x}\right)^2 \left(u^{y}\right)^5 - 6 {\mathcal{B}} {D_{s}} {(\partial_{y} u^{y})} {\nu^{y}} {\tau} \left(u^{x}\right)^2 \left(u^{y}\right)^5\\
  &\quad - 2 {D_{s}} {(\partial_{x} u^{x})} {I_{31}} {\tau} \left(u^{y}\right)^6 + 4 {D_{s}} {(\partial_{y} u^{y})} {I_{31}} {\tau} \left(u^{y}\right)^6 - 2 {\mathcal{B}'(T)} {D_{s}} {(\partial_{x} T)} {\nu^{x}} {\tau} \left(u^{y}\right)^6\\
  &\quad + 4 {\mathcal{B}'(T)} {D_{s}} {(\partial_{y} T)} {\nu^{y}} {\tau} \left(u^{y}\right)^6 - 2 {\mathcal{C}} {D_{s}} {(\partial_{x} u^{x})} {\Pi_{Q}} {\tau} \left(u^{y}\right)^6 + 4 {\mathcal{C}} {D_{s}} {(\partial_{y} u^{y})} {\Pi_{Q}} {\tau} \left(u^{y}\right)^6\\
  &\quad - 8 {\mathcal{C}} {D_{s}} {(\partial_{y} u^{x})} {\pi_{Q}^{xy}} {\tau} \left(u^{y}\right)^6 + 2 {\mathcal{C}} {D_{s}} {(\partial_{x} u^{x})} {\pi_{Q}{}^{\eta}{}_{\eta}} {\tau} \left(u^{y}\right)^6 - 2 {D_{s}} {I_{31}} {u^{\tau}} \left(u^{y}\right)^6\\
  &\quad - 2 {\mathcal{C}} {D_{s}} {\Pi_{Q}} {u^{\tau}} \left(u^{y}\right)^6 - 2 {\mathcal{C}} {D_{s}} {\pi_{Q}{}^{\eta}{}_{\eta}} {u^{\tau}} \left(u^{y}\right)^6 - 2 {\mathcal{A}} {D_{s}} {(\partial_{\tau} u^{x})} {\nu^{x}} {\tau} {u^{\tau}} \left(u^{y}\right)^6\\
  &\quad - 6 {\mathcal{B}} {D_{s}} {(\partial_{\tau} u^{x})} {\nu^{x}} {\tau} {u^{\tau}} \left(u^{y}\right)^6 - 2 {\mathcal{A}} {D_{s}} {(\partial_{x} u^{x})} {\nu^{x}} {\tau} {u^{x}} \left(u^{y}\right)^6\\
  &\quad - 6 {\mathcal{B}} {D_{s}} {(\partial_{x} u^{x})} {\nu^{x}} {\tau} {u^{x}} \left(u^{y}\right)^6 + 2 {\mathcal{A}} {D_{s}} {(\partial_{y} u^{y})} {\nu^{x}} {\tau} {u^{x}} \left(u^{y}\right)^6\\
  &\quad + 6 {\mathcal{B}} {D_{s}} {(\partial_{y} u^{y})} {\nu^{x}} {\tau} {u^{x}} \left(u^{y}\right)^6 + 2 {\mathcal{A}} {D_{s}} {(\partial_{y} u^{x})} {\nu^{y}} {\tau} {u^{x}} \left(u^{y}\right)^6\\
  &\quad + 6 {\mathcal{B}} {D_{s}} {(\partial_{y} u^{x})} {\nu^{y}} {\tau} {u^{x}} \left(u^{y}\right)^6 - 2 {\mathcal{A}} {D_{s}} {(\partial_{y} u^{x})} {\nu^{x}} {\tau} \left(u^{y}\right)^7 - 6 {\mathcal{B}} {D_{s}} {(\partial_{y} u^{x})} {\nu^{x}} {\tau} \left(u^{y}\right)^7\Bigr) \\
S_{5} &= \frac{1}{6 {T} {\tau} \left(u^{\tau}\right)^4 \left(\left(u^{y}\right)^2+1\right)}\Bigl(- 2 {D_{s}} {(\partial_{x} u^{x})} {I_{31}} {\tau} - 2 {D_{s}} {(\partial_{y} u^{y})} {I_{31}} {\tau} - 2 {\mathcal{B}'(T)} {D_{s}} {(\partial_{x} T)} {\nu^{x}} {\tau}\\
  &\quad - 2 {\mathcal{B}'(T)} {D_{s}} {(\partial_{y} T)} {\nu^{y}} {\tau} - 2 {\mathcal{C}} {D_{s}} {(\partial_{x} u^{x})} {\Pi_{Q}} {\tau} - 2 {\mathcal{C}} {D_{s}} {(\partial_{y} u^{y})} {\Pi_{Q}} {\tau}\\
  &\quad - 2 {\mathcal{C}} {D_{s}} {(\partial_{x} u^{y})} {\pi_{Q}^{xy}} {\tau} - 2 {\mathcal{C}} {D_{s}} {(\partial_{y} u^{x})} {\pi_{Q}^{xy}} {\tau} + 2 {\mathcal{C}} {D_{s}} {(\partial_{x} u^{x})} {\pi_{Q}^{yy}} {\tau}\\
  &\quad - 2 {\mathcal{C}} {D_{s}} {(\partial_{y} u^{y})} {\pi_{Q}^{yy}} {\tau} + 5 {\mathcal{C}} {D_{s}} {(\partial_{x} u^{x})} {\pi_{Q}{}^{\eta}{}_{\eta}} {\tau} + 3 {\mathcal{C}} {D_{s}} {(\partial_{y} u^{y})} {\pi_{Q}{}^{\eta}{}_{\eta}} {\tau}\\
  &\quad + 6 {\pi_{Q}{}^{\eta}{}_{\eta}} {T} {\tau} + 4 {D_{s}} {I_{31}} {u^{\tau}} + 4 {\mathcal{C}} {D_{s}} {\Pi_{Q}} {u^{\tau}} + 7 {\mathcal{C}} {D_{s}} {\pi_{Q}{}^{\eta}{}_{\eta}} {u^{\tau}}\\
  &\quad - 2 {\mathcal{A}} {D_{s}} {(\partial_{\tau} u^{x})} {\nu^{x}} {\tau} {u^{\tau}} - 2 {\mathcal{B}} {D_{s}} {(\partial_{\tau} u^{x})} {\nu^{x}} {\tau} {u^{\tau}} - 2 {\mathcal{A}} {D_{s}} {(\partial_{\tau} u^{y})} {\nu^{y}} {\tau} {u^{\tau}}\\
  &\quad - 2 {\mathcal{B}} {D_{s}} {(\partial_{\tau} u^{y})} {\nu^{y}} {\tau} {u^{\tau}} + 3 {\mathcal{C}'(T)} {D_{s}} {(\partial_{\tau} T)} {\pi_{Q}{}^{\eta}{}_{\eta}} {\tau} {u^{\tau}}\\
  &\quad - 2 {\mathcal{A}} {D_{s}} {(\partial_{x} u^{x})} {\nu^{x}} {\tau} {u^{x}} - 2 {\mathcal{A}} {D_{s}} {(\partial_{x} u^{y})} {\nu^{y}} {\tau} {u^{x}}\\
  &\quad + 3 {\mathcal{C}'(T)} {D_{s}} {(\partial_{x} T)} {\pi_{Q}{}^{\eta}{}_{\eta}} {\tau} {u^{x}} + 4 {\mathcal{B}} {D_{s}} {\nu^{x}} {u^{\tau}} {u^{x}} - 2 {D_{s}} {(\partial_{\tau} u^{x})} {I_{31}} {\tau} {u^{\tau}} {u^{x}}\\
  &\quad - 2 {\mathcal{B}'(T)} {D_{s}} {(\partial_{\tau} T)} {\nu^{x}} {\tau} {u^{\tau}} {u^{x}} - 2 {\mathcal{C}} {D_{s}} {(\partial_{\tau} u^{x})} {\Pi_{Q}} {\tau} {u^{\tau}} {u^{x}}\\
  &\quad - 2 {\mathcal{C}} {D_{s}} {(\partial_{\tau} u^{y})} {\pi_{Q}^{xy}} {\tau} {u^{\tau}} {u^{x}} + 2 {\mathcal{C}} {D_{s}} {(\partial_{\tau} u^{x})} {\pi_{Q}^{yy}} {\tau} {u^{\tau}} {u^{x}}\\
  &\quad + 5 {\mathcal{C}} {D_{s}} {(\partial_{\tau} u^{x})} {\pi_{Q}{}^{\eta}{}_{\eta}} {\tau} {u^{\tau}} {u^{x}} - 4 {D_{s}} {(\partial_{x} u^{x})} {I_{31}} {\tau} \left(u^{x}\right)^2\\
  &\quad - 4 {D_{s}} {(\partial_{y} u^{y})} {I_{31}} {\tau} \left(u^{x}\right)^2 - 4 {\mathcal{B}'(T)} {D_{s}} {(\partial_{x} T)} {\nu^{x}} {\tau} \left(u^{x}\right)^2\\
  &\quad - 4 {\mathcal{B}'(T)} {D_{s}} {(\partial_{y} T)} {\nu^{y}} {\tau} \left(u^{x}\right)^2 - 4 {\mathcal{C}} {D_{s}} {(\partial_{x} u^{x})} {\Pi_{Q}} {\tau} \left(u^{x}\right)^2\\
  &\quad - 4 {\mathcal{C}} {D_{s}} {(\partial_{y} u^{y})} {\Pi_{Q}} {\tau} \left(u^{x}\right)^2 - 4 {\mathcal{C}} {D_{s}} {(\partial_{x} u^{y})} {\pi_{Q}^{xy}} {\tau} \left(u^{x}\right)^2\\
  &\quad - 2 {\mathcal{C}} {D_{s}} {(\partial_{y} u^{x})} {\pi_{Q}^{xy}} {\tau} \left(u^{x}\right)^2 + 4 {\mathcal{C}} {D_{s}} {(\partial_{x} u^{x})} {\pi_{Q}^{yy}} {\tau} \left(u^{x}\right)^2\\
  &\quad - 4 {\mathcal{C}} {D_{s}} {(\partial_{y} u^{y})} {\pi_{Q}^{yy}} {\tau} \left(u^{x}\right)^2 + 10 {\mathcal{C}} {D_{s}} {(\partial_{x} u^{x})} {\pi_{Q}{}^{\eta}{}_{\eta}} {\tau} \left(u^{x}\right)^2\\
  &\quad + 6 {\mathcal{C}} {D_{s}} {(\partial_{y} u^{y})} {\pi_{Q}{}^{\eta}{}_{\eta}} {\tau} \left(u^{x}\right)^2 + 12 {\pi_{Q}{}^{\eta}{}_{\eta}} {T} {\tau} \left(u^{x}\right)^2 + 8 {D_{s}} {I_{31}} {u^{\tau}} \left(u^{x}\right)^2\\
  &\quad + 8 {\mathcal{C}} {D_{s}} {\Pi_{Q}} {u^{\tau}} \left(u^{x}\right)^2 + 14 {\mathcal{C}} {D_{s}} {\pi_{Q}{}^{\eta}{}_{\eta}} {u^{\tau}} \left(u^{x}\right)^2\\
  &\quad - 2 {\mathcal{A}} {D_{s}} {(\partial_{\tau} u^{x})} {\nu^{x}} {\tau} {u^{\tau}} \left(u^{x}\right)^2 - 4 {\mathcal{A}} {D_{s}} {(\partial_{\tau} u^{y})} {\nu^{y}} {\tau} {u^{\tau}} \left(u^{x}\right)^2\\
  &\quad - 2 {\mathcal{B}} {D_{s}} {(\partial_{\tau} u^{y})} {\nu^{y}} {\tau} {u^{\tau}} \left(u^{x}\right)^2 + 6 {\mathcal{C}'(T)} {D_{s}} {(\partial_{\tau} T)} {\pi_{Q}{}^{\eta}{}_{\eta}} {\tau} {u^{\tau}} \left(u^{x}\right)^2\\
  &\quad - 2 {\mathcal{A}} {D_{s}} {(\partial_{x} u^{x})} {\nu^{x}} {\tau} \left(u^{x}\right)^3 - 4 {\mathcal{A}} {D_{s}} {(\partial_{x} u^{y})} {\nu^{y}} {\tau} \left(u^{x}\right)^3\\
  &\quad + 6 {\mathcal{C}'(T)} {D_{s}} {(\partial_{x} T)} {\pi_{Q}{}^{\eta}{}_{\eta}} {\tau} \left(u^{x}\right)^3 + 4 {\mathcal{B}} {D_{s}} {\nu^{x}} {u^{\tau}} \left(u^{x}\right)^3\\
  &\quad - 2 {D_{s}} {(\partial_{\tau} u^{x})} {I_{31}} {\tau} {u^{\tau}} \left(u^{x}\right)^3 - 2 {\mathcal{B}'(T)} {D_{s}} {(\partial_{\tau} T)} {\nu^{x}} {\tau} {u^{\tau}} \left(u^{x}\right)^3\\
  &\quad - 2 {\mathcal{C}} {D_{s}} {(\partial_{\tau} u^{x})} {\Pi_{Q}} {\tau} {u^{\tau}} \left(u^{x}\right)^3 - 2 {\mathcal{C}} {D_{s}} {(\partial_{\tau} u^{y})} {\pi_{Q}^{xy}} {\tau} {u^{\tau}} \left(u^{x}\right)^3\\
  &\quad + 2 {\mathcal{C}} {D_{s}} {(\partial_{\tau} u^{x})} {\pi_{Q}^{yy}} {\tau} {u^{\tau}} \left(u^{x}\right)^3 + 5 {\mathcal{C}} {D_{s}} {(\partial_{\tau} u^{x})} {\pi_{Q}{}^{\eta}{}_{\eta}} {\tau} {u^{\tau}} \left(u^{x}\right)^3\\
  &\quad - 2 {D_{s}} {(\partial_{x} u^{x})} {I_{31}} {\tau} \left(u^{x}\right)^4 - 2 {D_{s}} {(\partial_{y} u^{y})} {I_{31}} {\tau} \left(u^{x}\right)^4\\
  &\quad - 2 {\mathcal{B}'(T)} {D_{s}} {(\partial_{x} T)} {\nu^{x}} {\tau} \left(u^{x}\right)^4 - 2 {\mathcal{B}'(T)} {D_{s}} {(\partial_{y} T)} {\nu^{y}} {\tau} \left(u^{x}\right)^4\\
  &\quad - 2 {\mathcal{C}} {D_{s}} {(\partial_{x} u^{x})} {\Pi_{Q}} {\tau} \left(u^{x}\right)^4 - 2 {\mathcal{C}} {D_{s}} {(\partial_{y} u^{y})} {\Pi_{Q}} {\tau} \left(u^{x}\right)^4\\
  &\quad - 2 {\mathcal{C}} {D_{s}} {(\partial_{x} u^{y})} {\pi_{Q}^{xy}} {\tau} \left(u^{x}\right)^4 + 2 {\mathcal{C}} {D_{s}} {(\partial_{x} u^{x})} {\pi_{Q}^{yy}} {\tau} \left(u^{x}\right)^4\\
  &\quad - 2 {\mathcal{C}} {D_{s}} {(\partial_{y} u^{y})} {\pi_{Q}^{yy}} {\tau} \left(u^{x}\right)^4 + 5 {\mathcal{C}} {D_{s}} {(\partial_{x} u^{x})} {\pi_{Q}{}^{\eta}{}_{\eta}} {\tau} \left(u^{x}\right)^4\\
  &\quad + 3 {\mathcal{C}} {D_{s}} {(\partial_{y} u^{y})} {\pi_{Q}{}^{\eta}{}_{\eta}} {\tau} \left(u^{x}\right)^4 + 6 {\pi_{Q}{}^{\eta}{}_{\eta}} {T} {\tau} \left(u^{x}\right)^4 + 4 {D_{s}} {I_{31}} {u^{\tau}} \left(u^{x}\right)^4\\
  &\quad + 4 {\mathcal{C}} {D_{s}} {\Pi_{Q}} {u^{\tau}} \left(u^{x}\right)^4 + 7 {\mathcal{C}} {D_{s}} {\pi_{Q}{}^{\eta}{}_{\eta}} {u^{\tau}} \left(u^{x}\right)^4 - 2 {\mathcal{A}} {D_{s}} {(\partial_{\tau} u^{y})} {\nu^{y}} {\tau} {u^{\tau}} \left(u^{x}\right)^4\\
  &\quad + 3 {\mathcal{C}'(T)} {D_{s}} {(\partial_{\tau} T)} {\pi_{Q}{}^{\eta}{}_{\eta}} {\tau} {u^{\tau}} \left(u^{x}\right)^4 - 2 {\mathcal{A}} {D_{s}} {(\partial_{x} u^{y})} {\nu^{y}} {\tau} \left(u^{x}\right)^5\\
  &\quad + 3 {\mathcal{C}'(T)} {D_{s}} {(\partial_{x} T)} {\pi_{Q}{}^{\eta}{}_{\eta}} {\tau} \left(u^{x}\right)^5 - 2 {\mathcal{A}} {D_{s}} {(\partial_{y} u^{x})} {\nu^{x}} {\tau} {u^{y}}\\
  &\quad - 2 {\mathcal{A}} {D_{s}} {(\partial_{y} u^{y})} {\nu^{y}} {\tau} {u^{y}} + 3 {\mathcal{C}'(T)} {D_{s}} {(\partial_{y} T)} {\pi_{Q}{}^{\eta}{}_{\eta}} {\tau} {u^{y}} + 4 {\mathcal{B}} {D_{s}} {\nu^{y}} {u^{\tau}} {u^{y}}\\
  &\quad - 2 {D_{s}} {(\partial_{\tau} u^{y})} {I_{31}} {\tau} {u^{\tau}} {u^{y}} - 2 {\mathcal{B}'(T)} {D_{s}} {(\partial_{\tau} T)} {\nu^{y}} {\tau} {u^{\tau}} {u^{y}}\\
  &\quad - 2 {\mathcal{C}} {D_{s}} {(\partial_{\tau} u^{y})} {\Pi_{Q}} {\tau} {u^{\tau}} {u^{y}} - 2 {\mathcal{C}} {D_{s}} {(\partial_{\tau} u^{x})} {\pi_{Q}^{xy}} {\tau} {u^{\tau}} {u^{y}}\\
  &\quad - 2 {\mathcal{C}} {D_{s}} {(\partial_{\tau} u^{y})} {\pi_{Q}^{yy}} {\tau} {u^{\tau}} {u^{y}} + 3 {\mathcal{C}} {D_{s}} {(\partial_{\tau} u^{y})} {\pi_{Q}{}^{\eta}{}_{\eta}} {\tau} {u^{\tau}} {u^{y}}\\
  &\quad - 2 {\mathcal{C}} {D_{s}} {(\partial_{x} u^{x})} {\pi_{Q}^{xy}} {\tau} {u^{x}} {u^{y}} + 2 {\mathcal{C}} {D_{s}} {(\partial_{y} u^{y})} {\pi_{Q}^{xy}} {\tau} {u^{x}} {u^{y}}\\
  &\quad - 2 {\mathcal{C}} {D_{s}} {(\partial_{x} u^{y})} {\pi_{Q}^{yy}} {\tau} {u^{x}} {u^{y}} + 2 {\mathcal{C}} {D_{s}} {(\partial_{y} u^{x})} {\pi_{Q}^{yy}} {\tau} {u^{x}} {u^{y}}\\
  &\quad - 2 {\mathcal{C}} {D_{s}} {(\partial_{x} u^{y})} {\pi_{Q}{}^{\eta}{}_{\eta}} {\tau} {u^{x}} {u^{y}} + 2 {\mathcal{A}} {D_{s}} {(\partial_{\tau} u^{y})} {\nu^{x}} {\tau} {u^{\tau}} {u^{x}} {u^{y}}\\
  &\quad + 2 {\mathcal{B}} {D_{s}} {(\partial_{\tau} u^{y})} {\nu^{x}} {\tau} {u^{\tau}} {u^{x}} {u^{y}} + 2 {\mathcal{A}} {D_{s}} {(\partial_{\tau} u^{x})} {\nu^{y}} {\tau} {u^{\tau}} {u^{x}} {u^{y}}\\
  &\quad + 2 {\mathcal{B}} {D_{s}} {(\partial_{\tau} u^{x})} {\nu^{y}} {\tau} {u^{\tau}} {u^{x}} {u^{y}} + 2 {\mathcal{A}} {D_{s}} {(\partial_{x} u^{y})} {\nu^{x}} {\tau} \left(u^{x}\right)^2 {u^{y}}\\
  &\quad - 2 {\mathcal{A}} {D_{s}} {(\partial_{y} u^{x})} {\nu^{x}} {\tau} \left(u^{x}\right)^2 {u^{y}} + 2 {\mathcal{A}} {D_{s}} {(\partial_{x} u^{x})} {\nu^{y}} {\tau} \left(u^{x}\right)^2 {u^{y}}\\
  &\quad - 4 {\mathcal{A}} {D_{s}} {(\partial_{y} u^{y})} {\nu^{y}} {\tau} \left(u^{x}\right)^2 {u^{y}} + 6 {\mathcal{C}'(T)} {D_{s}} {(\partial_{y} T)} {\pi_{Q}{}^{\eta}{}_{\eta}} {\tau} \left(u^{x}\right)^2 {u^{y}}\\
  &\quad + 4 {\mathcal{B}} {D_{s}} {\nu^{y}} {u^{\tau}} \left(u^{x}\right)^2 {u^{y}} - 2 {D_{s}} {(\partial_{\tau} u^{y})} {I_{31}} {\tau} {u^{\tau}} \left(u^{x}\right)^2 {u^{y}}\\
  &\quad - 2 {\mathcal{B}'(T)} {D_{s}} {(\partial_{\tau} T)} {\nu^{y}} {\tau} {u^{\tau}} \left(u^{x}\right)^2 {u^{y}} - 2 {\mathcal{C}} {D_{s}} {(\partial_{\tau} u^{y})} {\Pi_{Q}} {\tau} {u^{\tau}} \left(u^{x}\right)^2 {u^{y}}\\
  &\quad - 2 {\mathcal{C}} {D_{s}} {(\partial_{\tau} u^{x})} {\pi_{Q}^{xy}} {\tau} {u^{\tau}} \left(u^{x}\right)^2 {u^{y}} - 4 {\mathcal{C}} {D_{s}} {(\partial_{\tau} u^{y})} {\pi_{Q}^{yy}} {\tau} {u^{\tau}} \left(u^{x}\right)^2 {u^{y}}\\
  &\quad + {\mathcal{C}} {D_{s}} {(\partial_{\tau} u^{y})} {\pi_{Q}{}^{\eta}{}_{\eta}} {\tau} {u^{\tau}} \left(u^{x}\right)^2 {u^{y}} - 2 {\mathcal{C}} {D_{s}} {(\partial_{x} u^{x})} {\pi_{Q}^{xy}} {\tau} \left(u^{x}\right)^3 {u^{y}}\\
  &\quad + 2 {\mathcal{C}} {D_{s}} {(\partial_{y} u^{y})} {\pi_{Q}^{xy}} {\tau} \left(u^{x}\right)^3 {u^{y}} - 4 {\mathcal{C}} {D_{s}} {(\partial_{x} u^{y})} {\pi_{Q}^{yy}} {\tau} \left(u^{x}\right)^3 {u^{y}}\\
  &\quad + 2 {\mathcal{C}} {D_{s}} {(\partial_{y} u^{x})} {\pi_{Q}^{yy}} {\tau} \left(u^{x}\right)^3 {u^{y}} - 4 {\mathcal{C}} {D_{s}} {(\partial_{x} u^{y})} {\pi_{Q}{}^{\eta}{}_{\eta}} {\tau} \left(u^{x}\right)^3 {u^{y}}\\
  &\quad + 2 {\mathcal{A}} {D_{s}} {(\partial_{\tau} u^{y})} {\nu^{x}} {\tau} {u^{\tau}} \left(u^{x}\right)^3 {u^{y}} + 2 {\mathcal{A}} {D_{s}} {(\partial_{\tau} u^{x})} {\nu^{y}} {\tau} {u^{\tau}} \left(u^{x}\right)^3 {u^{y}}\\
  &\quad + 2 {\mathcal{A}} {D_{s}} {(\partial_{x} u^{y})} {\nu^{x}} {\tau} \left(u^{x}\right)^4 {u^{y}} + 2 {\mathcal{A}} {D_{s}} {(\partial_{x} u^{x})} {\nu^{y}} {\tau} \left(u^{x}\right)^4 {u^{y}}\\
  &\quad - 2 {\mathcal{A}} {D_{s}} {(\partial_{y} u^{y})} {\nu^{y}} {\tau} \left(u^{x}\right)^4 {u^{y}} + 3 {\mathcal{C}'(T)} {D_{s}} {(\partial_{y} T)} {\pi_{Q}{}^{\eta}{}_{\eta}} {\tau} \left(u^{x}\right)^4 {u^{y}}\\
  &\quad - 2 {\mathcal{C}} {D_{s}} {(\partial_{\tau} u^{y})} {\pi_{Q}^{yy}} {\tau} {u^{\tau}} \left(u^{x}\right)^4 {u^{y}} - 2 {\mathcal{C}} {D_{s}} {(\partial_{\tau} u^{y})} {\pi_{Q}{}^{\eta}{}_{\eta}} {\tau} {u^{\tau}} \left(u^{x}\right)^4 {u^{y}}\\
  &\quad - 2 {\mathcal{C}} {D_{s}} {(\partial_{x} u^{y})} {\pi_{Q}^{yy}} {\tau} \left(u^{x}\right)^5 {u^{y}} - 2 {\mathcal{C}} {D_{s}} {(\partial_{x} u^{y})} {\pi_{Q}{}^{\eta}{}_{\eta}} {\tau} \left(u^{x}\right)^5 {u^{y}}\\
  &\quad - 6 {D_{s}} {(\partial_{x} u^{x})} {I_{31}} {\tau} \left(u^{y}\right)^2 - 6 {D_{s}} {(\partial_{y} u^{y})} {I_{31}} {\tau} \left(u^{y}\right)^2\\
  &\quad - 6 {\mathcal{B}'(T)} {D_{s}} {(\partial_{x} T)} {\nu^{x}} {\tau} \left(u^{y}\right)^2 - 6 {\mathcal{B}'(T)} {D_{s}} {(\partial_{y} T)} {\nu^{y}} {\tau} \left(u^{y}\right)^2\\
  &\quad - 6 {\mathcal{C}} {D_{s}} {(\partial_{x} u^{x})} {\Pi_{Q}} {\tau} \left(u^{y}\right)^2 - 6 {\mathcal{C}} {D_{s}} {(\partial_{y} u^{y})} {\Pi_{Q}} {\tau} \left(u^{y}\right)^2\\
  &\quad - 4 {\mathcal{C}} {D_{s}} {(\partial_{x} u^{y})} {\pi_{Q}^{xy}} {\tau} \left(u^{y}\right)^2 - 6 {\mathcal{C}} {D_{s}} {(\partial_{y} u^{x})} {\pi_{Q}^{xy}} {\tau} \left(u^{y}\right)^2\\
  &\quad + 4 {\mathcal{C}} {D_{s}} {(\partial_{x} u^{x})} {\pi_{Q}^{yy}} {\tau} \left(u^{y}\right)^2 - 4 {\mathcal{C}} {D_{s}} {(\partial_{y} u^{y})} {\pi_{Q}^{yy}} {\tau} \left(u^{y}\right)^2\\
  &\quad + 15 {\mathcal{C}} {D_{s}} {(\partial_{x} u^{x})} {\pi_{Q}{}^{\eta}{}_{\eta}} {\tau} \left(u^{y}\right)^2 + 9 {\mathcal{C}} {D_{s}} {(\partial_{y} u^{y})} {\pi_{Q}{}^{\eta}{}_{\eta}} {\tau} \left(u^{y}\right)^2\\
  &\quad + 18 {\pi_{Q}{}^{\eta}{}_{\eta}} {T} {\tau} \left(u^{y}\right)^2 + 12 {D_{s}} {I_{31}} {u^{\tau}} \left(u^{y}\right)^2 + 12 {\mathcal{C}} {D_{s}} {\Pi_{Q}} {u^{\tau}} \left(u^{y}\right)^2\\
  &\quad + 21 {\mathcal{C}} {D_{s}} {\pi_{Q}{}^{\eta}{}_{\eta}} {u^{\tau}} \left(u^{y}\right)^2 - 6 {\mathcal{A}} {D_{s}} {(\partial_{\tau} u^{x})} {\nu^{x}} {\tau} {u^{\tau}} \left(u^{y}\right)^2\\
  &\quad - 4 {\mathcal{B}} {D_{s}} {(\partial_{\tau} u^{x})} {\nu^{x}} {\tau} {u^{\tau}} \left(u^{y}\right)^2 - 4 {\mathcal{A}} {D_{s}} {(\partial_{\tau} u^{y})} {\nu^{y}} {\tau} {u^{\tau}} \left(u^{y}\right)^2\\
  &\quad - 2 {\mathcal{B}} {D_{s}} {(\partial_{\tau} u^{y})} {\nu^{y}} {\tau} {u^{\tau}} \left(u^{y}\right)^2 + 9 {\mathcal{C}'(T)} {D_{s}} {(\partial_{\tau} T)} {\pi_{Q}{}^{\eta}{}_{\eta}} {\tau} {u^{\tau}} \left(u^{y}\right)^2\\
  &\quad - 6 {\mathcal{A}} {D_{s}} {(\partial_{x} u^{x})} {\nu^{x}} {\tau} {u^{x}} \left(u^{y}\right)^2 + 2 {\mathcal{A}} {D_{s}} {(\partial_{y} u^{y})} {\nu^{x}} {\tau} {u^{x}} \left(u^{y}\right)^2\\
  &\quad - 4 {\mathcal{A}} {D_{s}} {(\partial_{x} u^{y})} {\nu^{y}} {\tau} {u^{x}} \left(u^{y}\right)^2 + 2 {\mathcal{A}} {D_{s}} {(\partial_{y} u^{x})} {\nu^{y}} {\tau} {u^{x}} \left(u^{y}\right)^2\\
  &\quad + 9 {\mathcal{C}'(T)} {D_{s}} {(\partial_{x} T)} {\pi_{Q}{}^{\eta}{}_{\eta}} {\tau} {u^{x}} \left(u^{y}\right)^2 + 8 {\mathcal{B}} {D_{s}} {\nu^{x}} {u^{\tau}} {u^{x}} \left(u^{y}\right)^2\\
  &\quad - 4 {D_{s}} {(\partial_{\tau} u^{x})} {I_{31}} {\tau} {u^{\tau}} {u^{x}} \left(u^{y}\right)^2 - 4 {\mathcal{B}'(T)} {D_{s}} {(\partial_{\tau} T)} {\nu^{x}} {\tau} {u^{\tau}} {u^{x}} \left(u^{y}\right)^2\\
  &\quad - 4 {\mathcal{C}} {D_{s}} {(\partial_{\tau} u^{x})} {\Pi_{Q}} {\tau} {u^{\tau}} {u^{x}} \left(u^{y}\right)^2 + 4 {\mathcal{C}} {D_{s}} {(\partial_{\tau} u^{x})} {\pi_{Q}^{yy}} {\tau} {u^{\tau}} {u^{x}} \left(u^{y}\right)^2\\
  &\quad + 10 {\mathcal{C}} {D_{s}} {(\partial_{\tau} u^{x})} {\pi_{Q}{}^{\eta}{}_{\eta}} {\tau} {u^{\tau}} {u^{x}} \left(u^{y}\right)^2 - 8 {D_{s}} {(\partial_{x} u^{x})} {I_{31}} {\tau} \left(u^{x}\right)^2 \left(u^{y}\right)^2\\
  &\quad - 8 {D_{s}} {(\partial_{y} u^{y})} {I_{31}} {\tau} \left(u^{x}\right)^2 \left(u^{y}\right)^2 - 8 {\mathcal{B}'(T)} {D_{s}} {(\partial_{x} T)} {\nu^{x}} {\tau} \left(u^{x}\right)^2 \left(u^{y}\right)^2\\
  &\quad - 8 {\mathcal{B}'(T)} {D_{s}} {(\partial_{y} T)} {\nu^{y}} {\tau} \left(u^{x}\right)^2 \left(u^{y}\right)^2 - 8 {\mathcal{C}} {D_{s}} {(\partial_{x} u^{x})} {\Pi_{Q}} {\tau} \left(u^{x}\right)^2 \left(u^{y}\right)^2\\
  &\quad - 8 {\mathcal{C}} {D_{s}} {(\partial_{y} u^{y})} {\Pi_{Q}} {\tau} \left(u^{x}\right)^2 \left(u^{y}\right)^2 - 2 {\mathcal{C}} {D_{s}} {(\partial_{x} u^{y})} {\pi_{Q}^{xy}} {\tau} \left(u^{x}\right)^2 \left(u^{y}\right)^2\\
  &\quad - 4 {\mathcal{C}} {D_{s}} {(\partial_{y} u^{x})} {\pi_{Q}^{xy}} {\tau} \left(u^{x}\right)^2 \left(u^{y}\right)^2 + 6 {\mathcal{C}} {D_{s}} {(\partial_{x} u^{x})} {\pi_{Q}^{yy}} {\tau} \left(u^{x}\right)^2 \left(u^{y}\right)^2\\
  &\quad - 6 {\mathcal{C}} {D_{s}} {(\partial_{y} u^{y})} {\pi_{Q}^{yy}} {\tau} \left(u^{x}\right)^2 \left(u^{y}\right)^2 + 20 {\mathcal{C}} {D_{s}} {(\partial_{x} u^{x})} {\pi_{Q}{}^{\eta}{}_{\eta}} {\tau} \left(u^{x}\right)^2 \left(u^{y}\right)^2\\
  &\quad + 12 {\mathcal{C}} {D_{s}} {(\partial_{y} u^{y})} {\pi_{Q}{}^{\eta}{}_{\eta}} {\tau} \left(u^{x}\right)^2 \left(u^{y}\right)^2 + 24 {\pi_{Q}{}^{\eta}{}_{\eta}} {T} {\tau} \left(u^{x}\right)^2 \left(u^{y}\right)^2\\
  &\quad + 16 {D_{s}} {I_{31}} {u^{\tau}} \left(u^{x}\right)^2 \left(u^{y}\right)^2 + 16 {\mathcal{C}} {D_{s}} {\Pi_{Q}} {u^{\tau}} \left(u^{x}\right)^2 \left(u^{y}\right)^2\\
  &\quad + 28 {\mathcal{C}} {D_{s}} {\pi_{Q}{}^{\eta}{}_{\eta}} {u^{\tau}} \left(u^{x}\right)^2 \left(u^{y}\right)^2 - 4 {\mathcal{A}} {D_{s}} {(\partial_{\tau} u^{x})} {\nu^{x}} {\tau} {u^{\tau}} \left(u^{x}\right)^2 \left(u^{y}\right)^2\\
  &\quad - 6 {\mathcal{A}} {D_{s}} {(\partial_{\tau} u^{y})} {\nu^{y}} {\tau} {u^{\tau}} \left(u^{x}\right)^2 \left(u^{y}\right)^2 - 2 {\mathcal{B}} {D_{s}} {(\partial_{\tau} u^{y})} {\nu^{y}} {\tau} {u^{\tau}} \left(u^{x}\right)^2 \left(u^{y}\right)^2\\
  &\quad + 12 {\mathcal{C}'(T)} {D_{s}} {(\partial_{\tau} T)} {\pi_{Q}{}^{\eta}{}_{\eta}} {\tau} {u^{\tau}} \left(u^{x}\right)^2 \left(u^{y}\right)^2\\
  &\quad - 4 {\mathcal{A}} {D_{s}} {(\partial_{x} u^{x})} {\nu^{x}} {\tau} \left(u^{x}\right)^3 \left(u^{y}\right)^2 + 2 {\mathcal{A}} {D_{s}} {(\partial_{y} u^{y})} {\nu^{x}} {\tau} \left(u^{x}\right)^3 \left(u^{y}\right)^2\\
  &\quad - 6 {\mathcal{A}} {D_{s}} {(\partial_{x} u^{y})} {\nu^{y}} {\tau} \left(u^{x}\right)^3 \left(u^{y}\right)^2 + 2 {\mathcal{A}} {D_{s}} {(\partial_{y} u^{x})} {\nu^{y}} {\tau} \left(u^{x}\right)^3 \left(u^{y}\right)^2\\
  &\quad + 12 {\mathcal{C}'(T)} {D_{s}} {(\partial_{x} T)} {\pi_{Q}{}^{\eta}{}_{\eta}} {\tau} \left(u^{x}\right)^3 \left(u^{y}\right)^2 + 4 {\mathcal{B}} {D_{s}} {\nu^{x}} {u^{\tau}} \left(u^{x}\right)^3 \left(u^{y}\right)^2\\
  &\quad - 2 {D_{s}} {(\partial_{\tau} u^{x})} {I_{31}} {\tau} {u^{\tau}} \left(u^{x}\right)^3 \left(u^{y}\right)^2\\
  &\quad - 2 {\mathcal{B}'(T)} {D_{s}} {(\partial_{\tau} T)} {\nu^{x}} {\tau} {u^{\tau}} \left(u^{x}\right)^3 \left(u^{y}\right)^2\\
  &\quad - 2 {\mathcal{C}} {D_{s}} {(\partial_{\tau} u^{x})} {\Pi_{Q}} {\tau} {u^{\tau}} \left(u^{x}\right)^3 \left(u^{y}\right)^2\\
  &\quad + 2 {\mathcal{C}} {D_{s}} {(\partial_{\tau} u^{y})} {\pi_{Q}^{xy}} {\tau} {u^{\tau}} \left(u^{x}\right)^3 \left(u^{y}\right)^2\\
  &\quad + 2 {\mathcal{C}} {D_{s}} {(\partial_{\tau} u^{x})} {\pi_{Q}^{yy}} {\tau} {u^{\tau}} \left(u^{x}\right)^3 \left(u^{y}\right)^2\\
  &\quad + 5 {\mathcal{C}} {D_{s}} {(\partial_{\tau} u^{x})} {\pi_{Q}{}^{\eta}{}_{\eta}} {\tau} {u^{\tau}} \left(u^{x}\right)^3 \left(u^{y}\right)^2 - 2 {D_{s}} {(\partial_{x} u^{x})} {I_{31}} {\tau} \left(u^{x}\right)^4 \left(u^{y}\right)^2\\
  &\quad - 2 {D_{s}} {(\partial_{y} u^{y})} {I_{31}} {\tau} \left(u^{x}\right)^4 \left(u^{y}\right)^2 - 2 {\mathcal{B}'(T)} {D_{s}} {(\partial_{x} T)} {\nu^{x}} {\tau} \left(u^{x}\right)^4 \left(u^{y}\right)^2\\
  &\quad - 2 {\mathcal{B}'(T)} {D_{s}} {(\partial_{y} T)} {\nu^{y}} {\tau} \left(u^{x}\right)^4 \left(u^{y}\right)^2 - 2 {\mathcal{C}} {D_{s}} {(\partial_{x} u^{x})} {\Pi_{Q}} {\tau} \left(u^{x}\right)^4 \left(u^{y}\right)^2\\
  &\quad - 2 {\mathcal{C}} {D_{s}} {(\partial_{y} u^{y})} {\Pi_{Q}} {\tau} \left(u^{x}\right)^4 \left(u^{y}\right)^2 + 2 {\mathcal{C}} {D_{s}} {(\partial_{x} u^{y})} {\pi_{Q}^{xy}} {\tau} \left(u^{x}\right)^4 \left(u^{y}\right)^2\\
  &\quad + 2 {\mathcal{C}} {D_{s}} {(\partial_{x} u^{x})} {\pi_{Q}^{yy}} {\tau} \left(u^{x}\right)^4 \left(u^{y}\right)^2 - 2 {\mathcal{C}} {D_{s}} {(\partial_{y} u^{y})} {\pi_{Q}^{yy}} {\tau} \left(u^{x}\right)^4 \left(u^{y}\right)^2\\
  &\quad + 5 {\mathcal{C}} {D_{s}} {(\partial_{x} u^{x})} {\pi_{Q}{}^{\eta}{}_{\eta}} {\tau} \left(u^{x}\right)^4 \left(u^{y}\right)^2 + 3 {\mathcal{C}} {D_{s}} {(\partial_{y} u^{y})} {\pi_{Q}{}^{\eta}{}_{\eta}} {\tau} \left(u^{x}\right)^4 \left(u^{y}\right)^2\\
  &\quad + 6 {\pi_{Q}{}^{\eta}{}_{\eta}} {T} {\tau} \left(u^{x}\right)^4 \left(u^{y}\right)^2 + 4 {D_{s}} {I_{31}} {u^{\tau}} \left(u^{x}\right)^4 \left(u^{y}\right)^2\\
  &\quad + 4 {\mathcal{C}} {D_{s}} {\Pi_{Q}} {u^{\tau}} \left(u^{x}\right)^4 \left(u^{y}\right)^2 + 7 {\mathcal{C}} {D_{s}} {\pi_{Q}{}^{\eta}{}_{\eta}} {u^{\tau}} \left(u^{x}\right)^4 \left(u^{y}\right)^2\\
  &\quad - 2 {\mathcal{A}} {D_{s}} {(\partial_{\tau} u^{y})} {\nu^{y}} {\tau} {u^{\tau}} \left(u^{x}\right)^4 \left(u^{y}\right)^2\\
  &\quad + 3 {\mathcal{C}'(T)} {D_{s}} {(\partial_{\tau} T)} {\pi_{Q}{}^{\eta}{}_{\eta}} {\tau} {u^{\tau}} \left(u^{x}\right)^4 \left(u^{y}\right)^2\\
  &\quad - 2 {\mathcal{A}} {D_{s}} {(\partial_{x} u^{y})} {\nu^{y}} {\tau} \left(u^{x}\right)^5 \left(u^{y}\right)^2\\
  &\quad + 3 {\mathcal{C}'(T)} {D_{s}} {(\partial_{x} T)} {\pi_{Q}{}^{\eta}{}_{\eta}} {\tau} \left(u^{x}\right)^5 \left(u^{y}\right)^2 - 6 {\mathcal{A}} {D_{s}} {(\partial_{y} u^{x})} {\nu^{x}} {\tau} \left(u^{y}\right)^3\\
  &\quad - 4 {\mathcal{A}} {D_{s}} {(\partial_{y} u^{y})} {\nu^{y}} {\tau} \left(u^{y}\right)^3 + 9 {\mathcal{C}'(T)} {D_{s}} {(\partial_{y} T)} {\pi_{Q}{}^{\eta}{}_{\eta}} {\tau} \left(u^{y}\right)^3\\
  &\quad + 8 {\mathcal{B}} {D_{s}} {\nu^{y}} {u^{\tau}} \left(u^{y}\right)^3 - 4 {D_{s}} {(\partial_{\tau} u^{y})} {I_{31}} {\tau} {u^{\tau}} \left(u^{y}\right)^3\\
  &\quad - 4 {\mathcal{B}'(T)} {D_{s}} {(\partial_{\tau} T)} {\nu^{y}} {\tau} {u^{\tau}} \left(u^{y}\right)^3 - 4 {\mathcal{C}} {D_{s}} {(\partial_{\tau} u^{y})} {\Pi_{Q}} {\tau} {u^{\tau}} \left(u^{y}\right)^3\\
  &\quad - 4 {\mathcal{C}} {D_{s}} {(\partial_{\tau} u^{x})} {\pi_{Q}^{xy}} {\tau} {u^{\tau}} \left(u^{y}\right)^3 - 2 {\mathcal{C}} {D_{s}} {(\partial_{\tau} u^{y})} {\pi_{Q}^{yy}} {\tau} {u^{\tau}} \left(u^{y}\right)^3\\
  &\quad + 6 {\mathcal{C}} {D_{s}} {(\partial_{\tau} u^{y})} {\pi_{Q}{}^{\eta}{}_{\eta}} {\tau} {u^{\tau}} \left(u^{y}\right)^3 - 4 {\mathcal{C}} {D_{s}} {(\partial_{x} u^{x})} {\pi_{Q}^{xy}} {\tau} {u^{x}} \left(u^{y}\right)^3\\
  &\quad + 4 {\mathcal{C}} {D_{s}} {(\partial_{y} u^{y})} {\pi_{Q}^{xy}} {\tau} {u^{x}} \left(u^{y}\right)^3 - 2 {\mathcal{C}} {D_{s}} {(\partial_{x} u^{y})} {\pi_{Q}^{yy}} {\tau} {u^{x}} \left(u^{y}\right)^3\\
  &\quad + 4 {\mathcal{C}} {D_{s}} {(\partial_{y} u^{x})} {\pi_{Q}^{yy}} {\tau} {u^{x}} \left(u^{y}\right)^3 - 4 {\mathcal{C}} {D_{s}} {(\partial_{x} u^{y})} {\pi_{Q}{}^{\eta}{}_{\eta}} {\tau} {u^{x}} \left(u^{y}\right)^3\\
  &\quad + 4 {\mathcal{A}} {D_{s}} {(\partial_{\tau} u^{y})} {\nu^{x}} {\tau} {u^{\tau}} {u^{x}} \left(u^{y}\right)^3 + 2 {\mathcal{B}} {D_{s}} {(\partial_{\tau} u^{y})} {\nu^{x}} {\tau} {u^{\tau}} {u^{x}} \left(u^{y}\right)^3\\
  &\quad + 4 {\mathcal{A}} {D_{s}} {(\partial_{\tau} u^{x})} {\nu^{y}} {\tau} {u^{\tau}} {u^{x}} \left(u^{y}\right)^3 + 2 {\mathcal{B}} {D_{s}} {(\partial_{\tau} u^{x})} {\nu^{y}} {\tau} {u^{\tau}} {u^{x}} \left(u^{y}\right)^3\\
  &\quad + 4 {\mathcal{A}} {D_{s}} {(\partial_{x} u^{y})} {\nu^{x}} {\tau} \left(u^{x}\right)^2 \left(u^{y}\right)^3 - 4 {\mathcal{A}} {D_{s}} {(\partial_{y} u^{x})} {\nu^{x}} {\tau} \left(u^{x}\right)^2 \left(u^{y}\right)^3\\
  &\quad + 4 {\mathcal{A}} {D_{s}} {(\partial_{x} u^{x})} {\nu^{y}} {\tau} \left(u^{x}\right)^2 \left(u^{y}\right)^3 - 6 {\mathcal{A}} {D_{s}} {(\partial_{y} u^{y})} {\nu^{y}} {\tau} \left(u^{x}\right)^2 \left(u^{y}\right)^3\\
  &\quad + 12 {\mathcal{C}'(T)} {D_{s}} {(\partial_{y} T)} {\pi_{Q}{}^{\eta}{}_{\eta}} {\tau} \left(u^{x}\right)^2 \left(u^{y}\right)^3 + 4 {\mathcal{B}} {D_{s}} {\nu^{y}} {u^{\tau}} \left(u^{x}\right)^2 \left(u^{y}\right)^3\\
  &\quad - 2 {D_{s}} {(\partial_{\tau} u^{y})} {I_{31}} {\tau} {u^{\tau}} \left(u^{x}\right)^2 \left(u^{y}\right)^3\\
  &\quad - 2 {\mathcal{B}'(T)} {D_{s}} {(\partial_{\tau} T)} {\nu^{y}} {\tau} {u^{\tau}} \left(u^{x}\right)^2 \left(u^{y}\right)^3\\
  &\quad - 2 {\mathcal{C}} {D_{s}} {(\partial_{\tau} u^{y})} {\Pi_{Q}} {\tau} {u^{\tau}} \left(u^{x}\right)^2 \left(u^{y}\right)^3\\
  &\quad - 2 {\mathcal{C}} {D_{s}} {(\partial_{\tau} u^{x})} {\pi_{Q}^{xy}} {\tau} {u^{\tau}} \left(u^{x}\right)^2 \left(u^{y}\right)^3\\
  &\quad - 2 {\mathcal{C}} {D_{s}} {(\partial_{\tau} u^{y})} {\pi_{Q}^{yy}} {\tau} {u^{\tau}} \left(u^{x}\right)^2 \left(u^{y}\right)^3\\
  &\quad + {\mathcal{C}} {D_{s}} {(\partial_{\tau} u^{y})} {\pi_{Q}{}^{\eta}{}_{\eta}} {\tau} {u^{\tau}} \left(u^{x}\right)^2 \left(u^{y}\right)^3 - 2 {\mathcal{C}} {D_{s}} {(\partial_{x} u^{x})} {\pi_{Q}^{xy}} {\tau} \left(u^{x}\right)^3 \left(u^{y}\right)^3\\
  &\quad + 2 {\mathcal{C}} {D_{s}} {(\partial_{y} u^{y})} {\pi_{Q}^{xy}} {\tau} \left(u^{x}\right)^3 \left(u^{y}\right)^3 - 2 {\mathcal{C}} {D_{s}} {(\partial_{x} u^{y})} {\pi_{Q}^{yy}} {\tau} \left(u^{x}\right)^3 \left(u^{y}\right)^3\\
  &\quad + 2 {\mathcal{C}} {D_{s}} {(\partial_{y} u^{x})} {\pi_{Q}^{yy}} {\tau} \left(u^{x}\right)^3 \left(u^{y}\right)^3 - 4 {\mathcal{C}} {D_{s}} {(\partial_{x} u^{y})} {\pi_{Q}{}^{\eta}{}_{\eta}} {\tau} \left(u^{x}\right)^3 \left(u^{y}\right)^3\\
  &\quad + 2 {\mathcal{A}} {D_{s}} {(\partial_{\tau} u^{y})} {\nu^{x}} {\tau} {u^{\tau}} \left(u^{x}\right)^3 \left(u^{y}\right)^3 + 2 {\mathcal{A}} {D_{s}} {(\partial_{\tau} u^{x})} {\nu^{y}} {\tau} {u^{\tau}} \left(u^{x}\right)^3 \left(u^{y}\right)^3\\
  &\quad + 2 {\mathcal{A}} {D_{s}} {(\partial_{x} u^{y})} {\nu^{x}} {\tau} \left(u^{x}\right)^4 \left(u^{y}\right)^3 + 2 {\mathcal{A}} {D_{s}} {(\partial_{x} u^{x})} {\nu^{y}} {\tau} \left(u^{x}\right)^4 \left(u^{y}\right)^3\\
  &\quad - 2 {\mathcal{A}} {D_{s}} {(\partial_{y} u^{y})} {\nu^{y}} {\tau} \left(u^{x}\right)^4 \left(u^{y}\right)^3\\
  &\quad + 3 {\mathcal{C}'(T)} {D_{s}} {(\partial_{y} T)} {\pi_{Q}{}^{\eta}{}_{\eta}} {\tau} \left(u^{x}\right)^4 \left(u^{y}\right)^3 - 6 {D_{s}} {(\partial_{x} u^{x})} {I_{31}} {\tau} \left(u^{y}\right)^4\\
  &\quad - 6 {D_{s}} {(\partial_{y} u^{y})} {I_{31}} {\tau} \left(u^{y}\right)^4 - 6 {\mathcal{B}'(T)} {D_{s}} {(\partial_{x} T)} {\nu^{x}} {\tau} \left(u^{y}\right)^4\\
  &\quad - 6 {\mathcal{B}'(T)} {D_{s}} {(\partial_{y} T)} {\nu^{y}} {\tau} \left(u^{y}\right)^4 - 6 {\mathcal{C}} {D_{s}} {(\partial_{x} u^{x})} {\Pi_{Q}} {\tau} \left(u^{y}\right)^4\\
  &\quad - 6 {\mathcal{C}} {D_{s}} {(\partial_{y} u^{y})} {\Pi_{Q}} {\tau} \left(u^{y}\right)^4 - 2 {\mathcal{C}} {D_{s}} {(\partial_{x} u^{y})} {\pi_{Q}^{xy}} {\tau} \left(u^{y}\right)^4\\
  &\quad - 6 {\mathcal{C}} {D_{s}} {(\partial_{y} u^{x})} {\pi_{Q}^{xy}} {\tau} \left(u^{y}\right)^4 + 2 {\mathcal{C}} {D_{s}} {(\partial_{x} u^{x})} {\pi_{Q}^{yy}} {\tau} \left(u^{y}\right)^4\\
  &\quad - 2 {\mathcal{C}} {D_{s}} {(\partial_{y} u^{y})} {\pi_{Q}^{yy}} {\tau} \left(u^{y}\right)^4 + 15 {\mathcal{C}} {D_{s}} {(\partial_{x} u^{x})} {\pi_{Q}{}^{\eta}{}_{\eta}} {\tau} \left(u^{y}\right)^4\\
  &\quad + 9 {\mathcal{C}} {D_{s}} {(\partial_{y} u^{y})} {\pi_{Q}{}^{\eta}{}_{\eta}} {\tau} \left(u^{y}\right)^4 + 18 {\pi_{Q}{}^{\eta}{}_{\eta}} {T} {\tau} \left(u^{y}\right)^4 + 12 {D_{s}} {I_{31}} {u^{\tau}} \left(u^{y}\right)^4\\
  &\quad + 12 {\mathcal{C}} {D_{s}} {\Pi_{Q}} {u^{\tau}} \left(u^{y}\right)^4 + 21 {\mathcal{C}} {D_{s}} {\pi_{Q}{}^{\eta}{}_{\eta}} {u^{\tau}} \left(u^{y}\right)^4\\
  &\quad - 6 {\mathcal{A}} {D_{s}} {(\partial_{\tau} u^{x})} {\nu^{x}} {\tau} {u^{\tau}} \left(u^{y}\right)^4 - 2 {\mathcal{B}} {D_{s}} {(\partial_{\tau} u^{x})} {\nu^{x}} {\tau} {u^{\tau}} \left(u^{y}\right)^4\\
  &\quad - 2 {\mathcal{A}} {D_{s}} {(\partial_{\tau} u^{y})} {\nu^{y}} {\tau} {u^{\tau}} \left(u^{y}\right)^4 + 9 {\mathcal{C}'(T)} {D_{s}} {(\partial_{\tau} T)} {\pi_{Q}{}^{\eta}{}_{\eta}} {\tau} {u^{\tau}} \left(u^{y}\right)^4\\
  &\quad - 6 {\mathcal{A}} {D_{s}} {(\partial_{x} u^{x})} {\nu^{x}} {\tau} {u^{x}} \left(u^{y}\right)^4 + 4 {\mathcal{A}} {D_{s}} {(\partial_{y} u^{y})} {\nu^{x}} {\tau} {u^{x}} \left(u^{y}\right)^4\\
  &\quad - 2 {\mathcal{A}} {D_{s}} {(\partial_{x} u^{y})} {\nu^{y}} {\tau} {u^{x}} \left(u^{y}\right)^4 + 4 {\mathcal{A}} {D_{s}} {(\partial_{y} u^{x})} {\nu^{y}} {\tau} {u^{x}} \left(u^{y}\right)^4\\
  &\quad + 9 {\mathcal{C}'(T)} {D_{s}} {(\partial_{x} T)} {\pi_{Q}{}^{\eta}{}_{\eta}} {\tau} {u^{x}} \left(u^{y}\right)^4 + 4 {\mathcal{B}} {D_{s}} {\nu^{x}} {u^{\tau}} {u^{x}} \left(u^{y}\right)^4\\
  &\quad - 2 {D_{s}} {(\partial_{\tau} u^{x})} {I_{31}} {\tau} {u^{\tau}} {u^{x}} \left(u^{y}\right)^4 - 2 {\mathcal{B}'(T)} {D_{s}} {(\partial_{\tau} T)} {\nu^{x}} {\tau} {u^{\tau}} {u^{x}} \left(u^{y}\right)^4\\
  &\quad - 2 {\mathcal{C}} {D_{s}} {(\partial_{\tau} u^{x})} {\Pi_{Q}} {\tau} {u^{\tau}} {u^{x}} \left(u^{y}\right)^4 + 2 {\mathcal{C}} {D_{s}} {(\partial_{\tau} u^{y})} {\pi_{Q}^{xy}} {\tau} {u^{\tau}} {u^{x}} \left(u^{y}\right)^4\\
  &\quad + 2 {\mathcal{C}} {D_{s}} {(\partial_{\tau} u^{x})} {\pi_{Q}^{yy}} {\tau} {u^{\tau}} {u^{x}} \left(u^{y}\right)^4 + 5 {\mathcal{C}} {D_{s}} {(\partial_{\tau} u^{x})} {\pi_{Q}{}^{\eta}{}_{\eta}} {\tau} {u^{\tau}} {u^{x}} \left(u^{y}\right)^4\\
  &\quad - 4 {D_{s}} {(\partial_{x} u^{x})} {I_{31}} {\tau} \left(u^{x}\right)^2 \left(u^{y}\right)^4 - 4 {D_{s}} {(\partial_{y} u^{y})} {I_{31}} {\tau} \left(u^{x}\right)^2 \left(u^{y}\right)^4\\
  &\quad - 4 {\mathcal{B}'(T)} {D_{s}} {(\partial_{x} T)} {\nu^{x}} {\tau} \left(u^{x}\right)^2 \left(u^{y}\right)^4\\
  &\quad - 4 {\mathcal{B}'(T)} {D_{s}} {(\partial_{y} T)} {\nu^{y}} {\tau} \left(u^{x}\right)^2 \left(u^{y}\right)^4 - 4 {\mathcal{C}} {D_{s}} {(\partial_{x} u^{x})} {\Pi_{Q}} {\tau} \left(u^{x}\right)^2 \left(u^{y}\right)^4\\
  &\quad - 4 {\mathcal{C}} {D_{s}} {(\partial_{y} u^{y})} {\Pi_{Q}} {\tau} \left(u^{x}\right)^2 \left(u^{y}\right)^4 + 2 {\mathcal{C}} {D_{s}} {(\partial_{x} u^{y})} {\pi_{Q}^{xy}} {\tau} \left(u^{x}\right)^2 \left(u^{y}\right)^4\\
  &\quad - 2 {\mathcal{C}} {D_{s}} {(\partial_{y} u^{x})} {\pi_{Q}^{xy}} {\tau} \left(u^{x}\right)^2 \left(u^{y}\right)^4 + 2 {\mathcal{C}} {D_{s}} {(\partial_{x} u^{x})} {\pi_{Q}^{yy}} {\tau} \left(u^{x}\right)^2 \left(u^{y}\right)^4\\
  &\quad - 2 {\mathcal{C}} {D_{s}} {(\partial_{y} u^{y})} {\pi_{Q}^{yy}} {\tau} \left(u^{x}\right)^2 \left(u^{y}\right)^4 + 10 {\mathcal{C}} {D_{s}} {(\partial_{x} u^{x})} {\pi_{Q}{}^{\eta}{}_{\eta}} {\tau} \left(u^{x}\right)^2 \left(u^{y}\right)^4\\
  &\quad + 6 {\mathcal{C}} {D_{s}} {(\partial_{y} u^{y})} {\pi_{Q}{}^{\eta}{}_{\eta}} {\tau} \left(u^{x}\right)^2 \left(u^{y}\right)^4 + 12 {\pi_{Q}{}^{\eta}{}_{\eta}} {T} {\tau} \left(u^{x}\right)^2 \left(u^{y}\right)^4\\
  &\quad + 8 {D_{s}} {I_{31}} {u^{\tau}} \left(u^{x}\right)^2 \left(u^{y}\right)^4 + 8 {\mathcal{C}} {D_{s}} {\Pi_{Q}} {u^{\tau}} \left(u^{x}\right)^2 \left(u^{y}\right)^4\\
  &\quad + 14 {\mathcal{C}} {D_{s}} {\pi_{Q}{}^{\eta}{}_{\eta}} {u^{\tau}} \left(u^{x}\right)^2 \left(u^{y}\right)^4 - 2 {\mathcal{A}} {D_{s}} {(\partial_{\tau} u^{x})} {\nu^{x}} {\tau} {u^{\tau}} \left(u^{x}\right)^2 \left(u^{y}\right)^4\\
  &\quad - 2 {\mathcal{A}} {D_{s}} {(\partial_{\tau} u^{y})} {\nu^{y}} {\tau} {u^{\tau}} \left(u^{x}\right)^2 \left(u^{y}\right)^4\\
  &\quad + 6 {\mathcal{C}'(T)} {D_{s}} {(\partial_{\tau} T)} {\pi_{Q}{}^{\eta}{}_{\eta}} {\tau} {u^{\tau}} \left(u^{x}\right)^2 \left(u^{y}\right)^4\\
  &\quad - 2 {\mathcal{A}} {D_{s}} {(\partial_{x} u^{x})} {\nu^{x}} {\tau} \left(u^{x}\right)^3 \left(u^{y}\right)^4 + 2 {\mathcal{A}} {D_{s}} {(\partial_{y} u^{y})} {\nu^{x}} {\tau} \left(u^{x}\right)^3 \left(u^{y}\right)^4\\
  &\quad - 2 {\mathcal{A}} {D_{s}} {(\partial_{x} u^{y})} {\nu^{y}} {\tau} \left(u^{x}\right)^3 \left(u^{y}\right)^4 + 2 {\mathcal{A}} {D_{s}} {(\partial_{y} u^{x})} {\nu^{y}} {\tau} \left(u^{x}\right)^3 \left(u^{y}\right)^4\\
  &\quad + 6 {\mathcal{C}'(T)} {D_{s}} {(\partial_{x} T)} {\pi_{Q}{}^{\eta}{}_{\eta}} {\tau} \left(u^{x}\right)^3 \left(u^{y}\right)^4 - 6 {\mathcal{A}} {D_{s}} {(\partial_{y} u^{x})} {\nu^{x}} {\tau} \left(u^{y}\right)^5\\
  &\quad - 2 {\mathcal{A}} {D_{s}} {(\partial_{y} u^{y})} {\nu^{y}} {\tau} \left(u^{y}\right)^5 + 9 {\mathcal{C}'(T)} {D_{s}} {(\partial_{y} T)} {\pi_{Q}{}^{\eta}{}_{\eta}} {\tau} \left(u^{y}\right)^5\\
  &\quad + 4 {\mathcal{B}} {D_{s}} {\nu^{y}} {u^{\tau}} \left(u^{y}\right)^5 - 2 {D_{s}} {(\partial_{\tau} u^{y})} {I_{31}} {\tau} {u^{\tau}} \left(u^{y}\right)^5\\
  &\quad - 2 {\mathcal{B}'(T)} {D_{s}} {(\partial_{\tau} T)} {\nu^{y}} {\tau} {u^{\tau}} \left(u^{y}\right)^5 - 2 {\mathcal{C}} {D_{s}} {(\partial_{\tau} u^{y})} {\Pi_{Q}} {\tau} {u^{\tau}} \left(u^{y}\right)^5\\
  &\quad - 2 {\mathcal{C}} {D_{s}} {(\partial_{\tau} u^{x})} {\pi_{Q}^{xy}} {\tau} {u^{\tau}} \left(u^{y}\right)^5 + 3 {\mathcal{C}} {D_{s}} {(\partial_{\tau} u^{y})} {\pi_{Q}{}^{\eta}{}_{\eta}} {\tau} {u^{\tau}} \left(u^{y}\right)^5\\
  &\quad - 2 {\mathcal{C}} {D_{s}} {(\partial_{x} u^{x})} {\pi_{Q}^{xy}} {\tau} {u^{x}} \left(u^{y}\right)^5 + 2 {\mathcal{C}} {D_{s}} {(\partial_{y} u^{y})} {\pi_{Q}^{xy}} {\tau} {u^{x}} \left(u^{y}\right)^5\\
  &\quad + 2 {\mathcal{C}} {D_{s}} {(\partial_{y} u^{x})} {\pi_{Q}^{yy}} {\tau} {u^{x}} \left(u^{y}\right)^5 - 2 {\mathcal{C}} {D_{s}} {(\partial_{x} u^{y})} {\pi_{Q}{}^{\eta}{}_{\eta}} {\tau} {u^{x}} \left(u^{y}\right)^5\\
  &\quad + 2 {\mathcal{A}} {D_{s}} {(\partial_{\tau} u^{y})} {\nu^{x}} {\tau} {u^{\tau}} {u^{x}} \left(u^{y}\right)^5 + 2 {\mathcal{A}} {D_{s}} {(\partial_{\tau} u^{x})} {\nu^{y}} {\tau} {u^{\tau}} {u^{x}} \left(u^{y}\right)^5\\
  &\quad + 2 {\mathcal{A}} {D_{s}} {(\partial_{x} u^{y})} {\nu^{x}} {\tau} \left(u^{x}\right)^2 \left(u^{y}\right)^5 - 2 {\mathcal{A}} {D_{s}} {(\partial_{y} u^{x})} {\nu^{x}} {\tau} \left(u^{x}\right)^2 \left(u^{y}\right)^5\\
  &\quad + 2 {\mathcal{A}} {D_{s}} {(\partial_{x} u^{x})} {\nu^{y}} {\tau} \left(u^{x}\right)^2 \left(u^{y}\right)^5 - 2 {\mathcal{A}} {D_{s}} {(\partial_{y} u^{y})} {\nu^{y}} {\tau} \left(u^{x}\right)^2 \left(u^{y}\right)^5\\
  &\quad + 6 {\mathcal{C}'(T)} {D_{s}} {(\partial_{y} T)} {\pi_{Q}{}^{\eta}{}_{\eta}} {\tau} \left(u^{x}\right)^2 \left(u^{y}\right)^5 - 2 {D_{s}} {(\partial_{x} u^{x})} {I_{31}} {\tau} \left(u^{y}\right)^6\\
  &\quad - 2 {D_{s}} {(\partial_{y} u^{y})} {I_{31}} {\tau} \left(u^{y}\right)^6 - 2 {\mathcal{B}'(T)} {D_{s}} {(\partial_{x} T)} {\nu^{x}} {\tau} \left(u^{y}\right)^6\\
  &\quad - 2 {\mathcal{B}'(T)} {D_{s}} {(\partial_{y} T)} {\nu^{y}} {\tau} \left(u^{y}\right)^6 - 2 {\mathcal{C}} {D_{s}} {(\partial_{x} u^{x})} {\Pi_{Q}} {\tau} \left(u^{y}\right)^6\\
  &\quad - 2 {\mathcal{C}} {D_{s}} {(\partial_{y} u^{y})} {\Pi_{Q}} {\tau} \left(u^{y}\right)^6 - 2 {\mathcal{C}} {D_{s}} {(\partial_{y} u^{x})} {\pi_{Q}^{xy}} {\tau} \left(u^{y}\right)^6\\
  &\quad + 5 {\mathcal{C}} {D_{s}} {(\partial_{x} u^{x})} {\pi_{Q}{}^{\eta}{}_{\eta}} {\tau} \left(u^{y}\right)^6 + 3 {\mathcal{C}} {D_{s}} {(\partial_{y} u^{y})} {\pi_{Q}{}^{\eta}{}_{\eta}} {\tau} \left(u^{y}\right)^6\\
  &\quad + 6 {\pi_{Q}{}^{\eta}{}_{\eta}} {T} {\tau} \left(u^{y}\right)^6 + 4 {D_{s}} {I_{31}} {u^{\tau}} \left(u^{y}\right)^6 + 4 {\mathcal{C}} {D_{s}} {\Pi_{Q}} {u^{\tau}} \left(u^{y}\right)^6\\
  &\quad + 7 {\mathcal{C}} {D_{s}} {\pi_{Q}{}^{\eta}{}_{\eta}} {u^{\tau}} \left(u^{y}\right)^6 - 2 {\mathcal{A}} {D_{s}} {(\partial_{\tau} u^{x})} {\nu^{x}} {\tau} {u^{\tau}} \left(u^{y}\right)^6\\
  &\quad + 3 {\mathcal{C}'(T)} {D_{s}} {(\partial_{\tau} T)} {\pi_{Q}{}^{\eta}{}_{\eta}} {\tau} {u^{\tau}} \left(u^{y}\right)^6 - 2 {\mathcal{A}} {D_{s}} {(\partial_{x} u^{x})} {\nu^{x}} {\tau} {u^{x}} \left(u^{y}\right)^6\\
  &\quad + 2 {\mathcal{A}} {D_{s}} {(\partial_{y} u^{y})} {\nu^{x}} {\tau} {u^{x}} \left(u^{y}\right)^6 + 2 {\mathcal{A}} {D_{s}} {(\partial_{y} u^{x})} {\nu^{y}} {\tau} {u^{x}} \left(u^{y}\right)^6\\
  &\quad + 3 {\mathcal{C}'(T)} {D_{s}} {(\partial_{x} T)} {\pi_{Q}{}^{\eta}{}_{\eta}} {\tau} {u^{x}} \left(u^{y}\right)^6 - 2 {\mathcal{A}} {D_{s}} {(\partial_{y} u^{x})} {\nu^{x}} {\tau} \left(u^{y}\right)^7\\
  &\quad + 3 {\mathcal{C}'(T)} {D_{s}} {(\partial_{y} T)} {\pi_{Q}{}^{\eta}{}_{\eta}} {\tau} \left(u^{y}\right)^7\Bigr) \\
S_{6} &= \frac{1}{6 {T} {\tau} \left(u^{\tau}\right)^4 \left(\left(u^{y}\right)^2+1\right)}\Bigl(3 {D_{s}} {(\partial_{x} u^{y})} {I_{31}} {\tau} + 3 {D_{s}} {(\partial_{y} u^{x})} {I_{31}} {\tau} + 3 {\mathcal{B}'(T)} {D_{s}} {(\partial_{y} T)} {\nu^{x}} {\tau}\\
  &\quad + 3 {\mathcal{B}'(T)} {D_{s}} {(\partial_{x} T)} {\nu^{y}} {\tau} + 3 {\mathcal{C}} {D_{s}} {(\partial_{x} u^{y})} {\Pi_{Q}} {\tau} + 3 {\mathcal{C}} {D_{s}} {(\partial_{y} u^{x})} {\Pi_{Q}} {\tau}\\
  &\quad + 6 {\mathcal{C}} {D_{s}} {(\partial_{x} u^{x})} {\pi_{Q}^{xy}} {\tau} + 6 {\mathcal{C}} {D_{s}} {(\partial_{y} u^{y})} {\pi_{Q}^{xy}} {\tau} - 3 {\mathcal{C}} {D_{s}} {(\partial_{x} u^{y})} {\pi_{Q}^{yy}} {\tau}\\
  &\quad + 3 {\mathcal{C}} {D_{s}} {(\partial_{y} u^{x})} {\pi_{Q}^{yy}} {\tau} - 3 {\mathcal{C}} {D_{s}} {(\partial_{x} u^{y})} {\pi_{Q}{}^{\eta}{}_{\eta}} {\tau} + 6 {\pi_{Q}^{xy}} {T} {\tau} + 3 {\mathcal{C}} {D_{s}} {\pi_{Q}^{xy}} {u^{\tau}}\\
  &\quad + 3 {\mathcal{A}} {D_{s}} {(\partial_{\tau} u^{y})} {\nu^{x}} {\tau} {u^{\tau}} + 3 {\mathcal{B}} {D_{s}} {(\partial_{\tau} u^{y})} {\nu^{x}} {\tau} {u^{\tau}} + 3 {\mathcal{A}} {D_{s}} {(\partial_{\tau} u^{x})} {\nu^{y}} {\tau} {u^{\tau}}\\
  &\quad + 3 {\mathcal{B}} {D_{s}} {(\partial_{\tau} u^{x})} {\nu^{y}} {\tau} {u^{\tau}} + 3 {\mathcal{C}'(T)} {D_{s}} {(\partial_{\tau} T)} {\pi_{Q}^{xy}} {\tau} {u^{\tau}}\\
  &\quad + 3 {\mathcal{A}} {D_{s}} {(\partial_{x} u^{y})} {\nu^{x}} {\tau} {u^{x}} + 3 {\mathcal{B}} {D_{s}} {(\partial_{x} u^{y})} {\nu^{x}} {\tau} {u^{x}} - 3 {\mathcal{B}} {D_{s}} {(\partial_{y} u^{x})} {\nu^{x}} {\tau} {u^{x}}\\
  &\quad + 3 {\mathcal{A}} {D_{s}} {(\partial_{x} u^{x})} {\nu^{y}} {\tau} {u^{x}} + 3 {\mathcal{B}} {D_{s}} {(\partial_{x} u^{x})} {\nu^{y}} {\tau} {u^{x}} - 3 {\mathcal{B}} {D_{s}} {(\partial_{y} u^{y})} {\nu^{y}} {\tau} {u^{x}}\\
  &\quad + 3 {\mathcal{C}'(T)} {D_{s}} {(\partial_{x} T)} {\pi_{Q}^{xy}} {\tau} {u^{x}} + 3 {D_{s}} {(\partial_{\tau} u^{y})} {I_{31}} {\tau} {u^{\tau}} {u^{x}}\\
  &\quad + 3 {\mathcal{B}'(T)} {D_{s}} {(\partial_{\tau} T)} {\nu^{y}} {\tau} {u^{\tau}} {u^{x}} + 3 {\mathcal{C}} {D_{s}} {(\partial_{\tau} u^{y})} {\Pi_{Q}} {\tau} {u^{\tau}} {u^{x}}\\
  &\quad + 3 {\mathcal{C}} {D_{s}} {(\partial_{\tau} u^{x})} {\pi_{Q}^{xy}} {\tau} {u^{\tau}} {u^{x}} - 6 {\mathcal{C}} {D_{s}} {(\partial_{\tau} u^{y})} {\pi_{Q}^{yy}} {\tau} {u^{\tau}} {u^{x}}\\
  &\quad - 3 {\mathcal{C}} {D_{s}} {(\partial_{\tau} u^{y})} {\pi_{Q}{}^{\eta}{}_{\eta}} {\tau} {u^{\tau}} {u^{x}} + 9 {D_{s}} {(\partial_{x} u^{y})} {I_{31}} {\tau} \left(u^{x}\right)^2\\
  &\quad + 6 {D_{s}} {(\partial_{y} u^{x})} {I_{31}} {\tau} \left(u^{x}\right)^2 + 6 {\mathcal{B}'(T)} {D_{s}} {(\partial_{y} T)} {\nu^{x}} {\tau} \left(u^{x}\right)^2\\
  &\quad + 9 {\mathcal{B}'(T)} {D_{s}} {(\partial_{x} T)} {\nu^{y}} {\tau} \left(u^{x}\right)^2 + 9 {\mathcal{C}} {D_{s}} {(\partial_{x} u^{y})} {\Pi_{Q}} {\tau} \left(u^{x}\right)^2\\
  &\quad + 6 {\mathcal{C}} {D_{s}} {(\partial_{y} u^{x})} {\Pi_{Q}} {\tau} \left(u^{x}\right)^2 + 9 {\mathcal{C}} {D_{s}} {(\partial_{x} u^{x})} {\pi_{Q}^{xy}} {\tau} \left(u^{x}\right)^2\\
  &\quad + 12 {\mathcal{C}} {D_{s}} {(\partial_{y} u^{y})} {\pi_{Q}^{xy}} {\tau} \left(u^{x}\right)^2 - 12 {\mathcal{C}} {D_{s}} {(\partial_{x} u^{y})} {\pi_{Q}^{yy}} {\tau} \left(u^{x}\right)^2\\
  &\quad + 6 {\mathcal{C}} {D_{s}} {(\partial_{y} u^{x})} {\pi_{Q}^{yy}} {\tau} \left(u^{x}\right)^2 - 9 {\mathcal{C}} {D_{s}} {(\partial_{x} u^{y})} {\pi_{Q}{}^{\eta}{}_{\eta}} {\tau} \left(u^{x}\right)^2\\
  &\quad + 12 {\pi_{Q}^{xy}} {T} {\tau} \left(u^{x}\right)^2 + 6 {\mathcal{C}} {D_{s}} {\pi_{Q}^{xy}} {u^{\tau}} \left(u^{x}\right)^2 + 6 {\mathcal{A}} {D_{s}} {(\partial_{\tau} u^{y})} {\nu^{x}} {\tau} {u^{\tau}} \left(u^{x}\right)^2\\
  &\quad + 6 {\mathcal{B}} {D_{s}} {(\partial_{\tau} u^{y})} {\nu^{x}} {\tau} {u^{\tau}} \left(u^{x}\right)^2 + 6 {\mathcal{A}} {D_{s}} {(\partial_{\tau} u^{x})} {\nu^{y}} {\tau} {u^{\tau}} \left(u^{x}\right)^2\\
  &\quad + 6 {\mathcal{B}} {D_{s}} {(\partial_{\tau} u^{x})} {\nu^{y}} {\tau} {u^{\tau}} \left(u^{x}\right)^2 + 6 {\mathcal{C}'(T)} {D_{s}} {(\partial_{\tau} T)} {\pi_{Q}^{xy}} {\tau} {u^{\tau}} \left(u^{x}\right)^2\\
  &\quad + 6 {\mathcal{A}} {D_{s}} {(\partial_{x} u^{y})} {\nu^{x}} {\tau} \left(u^{x}\right)^3 + 6 {\mathcal{B}} {D_{s}} {(\partial_{x} u^{y})} {\nu^{x}} {\tau} \left(u^{x}\right)^3\\
  &\quad - 3 {\mathcal{B}} {D_{s}} {(\partial_{y} u^{x})} {\nu^{x}} {\tau} \left(u^{x}\right)^3 + 6 {\mathcal{A}} {D_{s}} {(\partial_{x} u^{x})} {\nu^{y}} {\tau} \left(u^{x}\right)^3\\
  &\quad + 6 {\mathcal{B}} {D_{s}} {(\partial_{x} u^{x})} {\nu^{y}} {\tau} \left(u^{x}\right)^3 - 6 {\mathcal{B}} {D_{s}} {(\partial_{y} u^{y})} {\nu^{y}} {\tau} \left(u^{x}\right)^3\\
  &\quad + 6 {\mathcal{C}'(T)} {D_{s}} {(\partial_{x} T)} {\pi_{Q}^{xy}} {\tau} \left(u^{x}\right)^3 + 6 {D_{s}} {(\partial_{\tau} u^{y})} {I_{31}} {\tau} {u^{\tau}} \left(u^{x}\right)^3\\
  &\quad + 6 {\mathcal{B}'(T)} {D_{s}} {(\partial_{\tau} T)} {\nu^{y}} {\tau} {u^{\tau}} \left(u^{x}\right)^3 + 6 {\mathcal{C}} {D_{s}} {(\partial_{\tau} u^{y})} {\Pi_{Q}} {\tau} {u^{\tau}} \left(u^{x}\right)^3\\
  &\quad + 3 {\mathcal{C}} {D_{s}} {(\partial_{\tau} u^{x})} {\pi_{Q}^{xy}} {\tau} {u^{\tau}} \left(u^{x}\right)^3 - 12 {\mathcal{C}} {D_{s}} {(\partial_{\tau} u^{y})} {\pi_{Q}^{yy}} {\tau} {u^{\tau}} \left(u^{x}\right)^3\\
  &\quad - 6 {\mathcal{C}} {D_{s}} {(\partial_{\tau} u^{y})} {\pi_{Q}{}^{\eta}{}_{\eta}} {\tau} {u^{\tau}} \left(u^{x}\right)^3 + 9 {D_{s}} {(\partial_{x} u^{y})} {I_{31}} {\tau} \left(u^{x}\right)^4\\
  &\quad + 3 {D_{s}} {(\partial_{y} u^{x})} {I_{31}} {\tau} \left(u^{x}\right)^4 + 3 {\mathcal{B}'(T)} {D_{s}} {(\partial_{y} T)} {\nu^{x}} {\tau} \left(u^{x}\right)^4\\
  &\quad + 9 {\mathcal{B}'(T)} {D_{s}} {(\partial_{x} T)} {\nu^{y}} {\tau} \left(u^{x}\right)^4 + 9 {\mathcal{C}} {D_{s}} {(\partial_{x} u^{y})} {\Pi_{Q}} {\tau} \left(u^{x}\right)^4\\
  &\quad + 3 {\mathcal{C}} {D_{s}} {(\partial_{y} u^{x})} {\Pi_{Q}} {\tau} \left(u^{x}\right)^4 + 3 {\mathcal{C}} {D_{s}} {(\partial_{x} u^{x})} {\pi_{Q}^{xy}} {\tau} \left(u^{x}\right)^4\\
  &\quad + 6 {\mathcal{C}} {D_{s}} {(\partial_{y} u^{y})} {\pi_{Q}^{xy}} {\tau} \left(u^{x}\right)^4 - 15 {\mathcal{C}} {D_{s}} {(\partial_{x} u^{y})} {\pi_{Q}^{yy}} {\tau} \left(u^{x}\right)^4\\
  &\quad + 3 {\mathcal{C}} {D_{s}} {(\partial_{y} u^{x})} {\pi_{Q}^{yy}} {\tau} \left(u^{x}\right)^4 - 9 {\mathcal{C}} {D_{s}} {(\partial_{x} u^{y})} {\pi_{Q}{}^{\eta}{}_{\eta}} {\tau} \left(u^{x}\right)^4\\
  &\quad + 6 {\pi_{Q}^{xy}} {T} {\tau} \left(u^{x}\right)^4 + 3 {\mathcal{C}} {D_{s}} {\pi_{Q}^{xy}} {u^{\tau}} \left(u^{x}\right)^4 + 3 {\mathcal{A}} {D_{s}} {(\partial_{\tau} u^{y})} {\nu^{x}} {\tau} {u^{\tau}} \left(u^{x}\right)^4\\
  &\quad + 3 {\mathcal{B}} {D_{s}} {(\partial_{\tau} u^{y})} {\nu^{x}} {\tau} {u^{\tau}} \left(u^{x}\right)^4 + 3 {\mathcal{A}} {D_{s}} {(\partial_{\tau} u^{x})} {\nu^{y}} {\tau} {u^{\tau}} \left(u^{x}\right)^4\\
  &\quad + 3 {\mathcal{B}} {D_{s}} {(\partial_{\tau} u^{x})} {\nu^{y}} {\tau} {u^{\tau}} \left(u^{x}\right)^4 + 3 {\mathcal{C}'(T)} {D_{s}} {(\partial_{\tau} T)} {\pi_{Q}^{xy}} {\tau} {u^{\tau}} \left(u^{x}\right)^4\\
  &\quad + 3 {\mathcal{A}} {D_{s}} {(\partial_{x} u^{y})} {\nu^{x}} {\tau} \left(u^{x}\right)^5 + 3 {\mathcal{B}} {D_{s}} {(\partial_{x} u^{y})} {\nu^{x}} {\tau} \left(u^{x}\right)^5\\
  &\quad + 3 {\mathcal{A}} {D_{s}} {(\partial_{x} u^{x})} {\nu^{y}} {\tau} \left(u^{x}\right)^5 + 3 {\mathcal{B}} {D_{s}} {(\partial_{x} u^{x})} {\nu^{y}} {\tau} \left(u^{x}\right)^5\\
  &\quad - 3 {\mathcal{B}} {D_{s}} {(\partial_{y} u^{y})} {\nu^{y}} {\tau} \left(u^{x}\right)^5 + 3 {\mathcal{C}'(T)} {D_{s}} {(\partial_{x} T)} {\pi_{Q}^{xy}} {\tau} \left(u^{x}\right)^5\\
  &\quad + 3 {D_{s}} {(\partial_{\tau} u^{y})} {I_{31}} {\tau} {u^{\tau}} \left(u^{x}\right)^5 + 3 {\mathcal{B}'(T)} {D_{s}} {(\partial_{\tau} T)} {\nu^{y}} {\tau} {u^{\tau}} \left(u^{x}\right)^5\\
  &\quad + 3 {\mathcal{C}} {D_{s}} {(\partial_{\tau} u^{y})} {\Pi_{Q}} {\tau} {u^{\tau}} \left(u^{x}\right)^5 - 6 {\mathcal{C}} {D_{s}} {(\partial_{\tau} u^{y})} {\pi_{Q}^{yy}} {\tau} {u^{\tau}} \left(u^{x}\right)^5\\
  &\quad - 3 {\mathcal{C}} {D_{s}} {(\partial_{\tau} u^{y})} {\pi_{Q}{}^{\eta}{}_{\eta}} {\tau} {u^{\tau}} \left(u^{x}\right)^5 + 3 {D_{s}} {(\partial_{x} u^{y})} {I_{31}} {\tau} \left(u^{x}\right)^6\\
  &\quad + 3 {\mathcal{B}'(T)} {D_{s}} {(\partial_{x} T)} {\nu^{y}} {\tau} \left(u^{x}\right)^6 + 3 {\mathcal{C}} {D_{s}} {(\partial_{x} u^{y})} {\Pi_{Q}} {\tau} \left(u^{x}\right)^6\\
  &\quad - 6 {\mathcal{C}} {D_{s}} {(\partial_{x} u^{y})} {\pi_{Q}^{yy}} {\tau} \left(u^{x}\right)^6 - 3 {\mathcal{C}} {D_{s}} {(\partial_{x} u^{y})} {\pi_{Q}{}^{\eta}{}_{\eta}} {\tau} \left(u^{x}\right)^6\\
  &\quad - 3 {\mathcal{B}} {D_{s}} {(\partial_{x} u^{x})} {\nu^{x}} {\tau} {u^{y}} + 3 {\mathcal{A}} {D_{s}} {(\partial_{y} u^{y})} {\nu^{x}} {\tau} {u^{y}} + 3 {\mathcal{B}} {D_{s}} {(\partial_{y} u^{y})} {\nu^{x}} {\tau} {u^{y}}\\
  &\quad - 3 {\mathcal{B}} {D_{s}} {(\partial_{x} u^{y})} {\nu^{y}} {\tau} {u^{y}} + 3 {\mathcal{A}} {D_{s}} {(\partial_{y} u^{x})} {\nu^{y}} {\tau} {u^{y}} + 3 {\mathcal{B}} {D_{s}} {(\partial_{y} u^{x})} {\nu^{y}} {\tau} {u^{y}}\\
  &\quad + 3 {\mathcal{C}'(T)} {D_{s}} {(\partial_{y} T)} {\pi_{Q}^{xy}} {\tau} {u^{y}} + 3 {D_{s}} {(\partial_{\tau} u^{x})} {I_{31}} {\tau} {u^{\tau}} {u^{y}}\\
  &\quad + 3 {\mathcal{B}'(T)} {D_{s}} {(\partial_{\tau} T)} {\nu^{x}} {\tau} {u^{\tau}} {u^{y}} + 3 {\mathcal{C}} {D_{s}} {(\partial_{\tau} u^{x})} {\Pi_{Q}} {\tau} {u^{\tau}} {u^{y}}\\
  &\quad + 3 {\mathcal{C}} {D_{s}} {(\partial_{\tau} u^{y})} {\pi_{Q}^{xy}} {\tau} {u^{\tau}} {u^{y}} + 6 {\mathcal{C}} {D_{s}} {(\partial_{\tau} u^{x})} {\pi_{Q}^{yy}} {\tau} {u^{\tau}} {u^{y}}\\
  &\quad + 3 {\mathcal{C}} {D_{s}} {(\partial_{\tau} u^{x})} {\pi_{Q}{}^{\eta}{}_{\eta}} {\tau} {u^{\tau}} {u^{y}} + {D_{s}} {(\partial_{x} u^{x})} {I_{31}} {\tau} {u^{x}} {u^{y}}\\
  &\quad + {D_{s}} {(\partial_{y} u^{y})} {I_{31}} {\tau} {u^{x}} {u^{y}} + {\mathcal{B}'(T)} {D_{s}} {(\partial_{x} T)} {\nu^{x}} {\tau} {u^{x}} {u^{y}}\\
  &\quad + {\mathcal{B}'(T)} {D_{s}} {(\partial_{y} T)} {\nu^{y}} {\tau} {u^{x}} {u^{y}} + {\mathcal{C}} {D_{s}} {(\partial_{x} u^{x})} {\Pi_{Q}} {\tau} {u^{x}} {u^{y}}\\
  &\quad + {\mathcal{C}} {D_{s}} {(\partial_{y} u^{y})} {\Pi_{Q}} {\tau} {u^{x}} {u^{y}} + {\mathcal{C}} {D_{s}} {(\partial_{x} u^{y})} {\pi_{Q}^{xy}} {\tau} {u^{x}} {u^{y}}\\
  &\quad - 5 {\mathcal{C}} {D_{s}} {(\partial_{y} u^{x})} {\pi_{Q}^{xy}} {\tau} {u^{x}} {u^{y}} + 5 {\mathcal{C}} {D_{s}} {(\partial_{x} u^{x})} {\pi_{Q}^{yy}} {\tau} {u^{x}} {u^{y}}\\
  &\quad - 5 {\mathcal{C}} {D_{s}} {(\partial_{y} u^{y})} {\pi_{Q}^{yy}} {\tau} {u^{x}} {u^{y}} + 5 {\mathcal{C}} {D_{s}} {(\partial_{x} u^{x})} {\pi_{Q}{}^{\eta}{}_{\eta}} {\tau} {u^{x}} {u^{y}}\\
  &\quad - 2 {D_{s}} {I_{31}} {u^{\tau}} {u^{x}} {u^{y}} - 2 {\mathcal{C}} {D_{s}} {\Pi_{Q}} {u^{\tau}} {u^{x}} {u^{y}} - 2 {\mathcal{C}} {D_{s}} {\pi_{Q}{}^{\eta}{}_{\eta}} {u^{\tau}} {u^{x}} {u^{y}}\\
  &\quad - 2 {\mathcal{A}} {D_{s}} {(\partial_{\tau} u^{x})} {\nu^{x}} {\tau} {u^{\tau}} {u^{x}} {u^{y}} - 8 {\mathcal{B}} {D_{s}} {(\partial_{\tau} u^{x})} {\nu^{x}} {\tau} {u^{\tau}} {u^{x}} {u^{y}}\\
  &\quad - 2 {\mathcal{A}} {D_{s}} {(\partial_{\tau} u^{y})} {\nu^{y}} {\tau} {u^{\tau}} {u^{x}} {u^{y}} - 8 {\mathcal{B}} {D_{s}} {(\partial_{\tau} u^{y})} {\nu^{y}} {\tau} {u^{\tau}} {u^{x}} {u^{y}}\\
  &\quad - 2 {\mathcal{A}} {D_{s}} {(\partial_{x} u^{x})} {\nu^{x}} {\tau} \left(u^{x}\right)^2 {u^{y}} - 9 {\mathcal{B}} {D_{s}} {(\partial_{x} u^{x})} {\nu^{x}} {\tau} \left(u^{x}\right)^2 {u^{y}}\\
  &\quad + 6 {\mathcal{A}} {D_{s}} {(\partial_{y} u^{y})} {\nu^{x}} {\tau} \left(u^{x}\right)^2 {u^{y}} + 9 {\mathcal{B}} {D_{s}} {(\partial_{y} u^{y})} {\nu^{x}} {\tau} \left(u^{x}\right)^2 {u^{y}}\\
  &\quad - 2 {\mathcal{A}} {D_{s}} {(\partial_{x} u^{y})} {\nu^{y}} {\tau} \left(u^{x}\right)^2 {u^{y}} - 12 {\mathcal{B}} {D_{s}} {(\partial_{x} u^{y})} {\nu^{y}} {\tau} \left(u^{x}\right)^2 {u^{y}}\\
  &\quad + 6 {\mathcal{A}} {D_{s}} {(\partial_{y} u^{x})} {\nu^{y}} {\tau} \left(u^{x}\right)^2 {u^{y}} + 9 {\mathcal{B}} {D_{s}} {(\partial_{y} u^{x})} {\nu^{y}} {\tau} \left(u^{x}\right)^2 {u^{y}}\\
  &\quad + 6 {\mathcal{C}'(T)} {D_{s}} {(\partial_{y} T)} {\pi_{Q}^{xy}} {\tau} \left(u^{x}\right)^2 {u^{y}} - 2 {\mathcal{B}} {D_{s}} {\nu^{x}} {u^{\tau}} \left(u^{x}\right)^2 {u^{y}}\\
  &\quad + 4 {D_{s}} {(\partial_{\tau} u^{x})} {I_{31}} {\tau} {u^{\tau}} \left(u^{x}\right)^2 {u^{y}} + 4 {\mathcal{B}'(T)} {D_{s}} {(\partial_{\tau} T)} {\nu^{x}} {\tau} {u^{\tau}} \left(u^{x}\right)^2 {u^{y}}\\
  &\quad + 4 {\mathcal{C}} {D_{s}} {(\partial_{\tau} u^{x})} {\Pi_{Q}} {\tau} {u^{\tau}} \left(u^{x}\right)^2 {u^{y}} + 7 {\mathcal{C}} {D_{s}} {(\partial_{\tau} u^{y})} {\pi_{Q}^{xy}} {\tau} {u^{\tau}} \left(u^{x}\right)^2 {u^{y}}\\
  &\quad + 14 {\mathcal{C}} {D_{s}} {(\partial_{\tau} u^{x})} {\pi_{Q}^{yy}} {\tau} {u^{\tau}} \left(u^{x}\right)^2 {u^{y}} + 8 {\mathcal{C}} {D_{s}} {(\partial_{\tau} u^{x})} {\pi_{Q}{}^{\eta}{}_{\eta}} {\tau} {u^{\tau}} \left(u^{x}\right)^2 {u^{y}}\\
  &\quad + 2 {D_{s}} {(\partial_{x} u^{x})} {I_{31}} {\tau} \left(u^{x}\right)^3 {u^{y}} + 2 {D_{s}} {(\partial_{y} u^{y})} {I_{31}} {\tau} \left(u^{x}\right)^3 {u^{y}}\\
  &\quad + 2 {\mathcal{B}'(T)} {D_{s}} {(\partial_{x} T)} {\nu^{x}} {\tau} \left(u^{x}\right)^3 {u^{y}} + 2 {\mathcal{B}'(T)} {D_{s}} {(\partial_{y} T)} {\nu^{y}} {\tau} \left(u^{x}\right)^3 {u^{y}}\\
  &\quad + 2 {\mathcal{C}} {D_{s}} {(\partial_{x} u^{x})} {\Pi_{Q}} {\tau} \left(u^{x}\right)^3 {u^{y}} + 2 {\mathcal{C}} {D_{s}} {(\partial_{y} u^{y})} {\Pi_{Q}} {\tau} \left(u^{x}\right)^3 {u^{y}}\\
  &\quad + 5 {\mathcal{C}} {D_{s}} {(\partial_{x} u^{y})} {\pi_{Q}^{xy}} {\tau} \left(u^{x}\right)^3 {u^{y}} - 5 {\mathcal{C}} {D_{s}} {(\partial_{y} u^{x})} {\pi_{Q}^{xy}} {\tau} \left(u^{x}\right)^3 {u^{y}}\\
  &\quad + 13 {\mathcal{C}} {D_{s}} {(\partial_{x} u^{x})} {\pi_{Q}^{yy}} {\tau} \left(u^{x}\right)^3 {u^{y}} - 10 {\mathcal{C}} {D_{s}} {(\partial_{y} u^{y})} {\pi_{Q}^{yy}} {\tau} \left(u^{x}\right)^3 {u^{y}}\\
  &\quad + 10 {\mathcal{C}} {D_{s}} {(\partial_{x} u^{x})} {\pi_{Q}{}^{\eta}{}_{\eta}} {\tau} \left(u^{x}\right)^3 {u^{y}} - 4 {D_{s}} {I_{31}} {u^{\tau}} \left(u^{x}\right)^3 {u^{y}}\\
  &\quad - 4 {\mathcal{C}} {D_{s}} {\Pi_{Q}} {u^{\tau}} \left(u^{x}\right)^3 {u^{y}} - 4 {\mathcal{C}} {D_{s}} {\pi_{Q}{}^{\eta}{}_{\eta}} {u^{\tau}} \left(u^{x}\right)^3 {u^{y}}\\
  &\quad - 2 {\mathcal{A}} {D_{s}} {(\partial_{\tau} u^{x})} {\nu^{x}} {\tau} {u^{\tau}} \left(u^{x}\right)^3 {u^{y}} - 6 {\mathcal{B}} {D_{s}} {(\partial_{\tau} u^{x})} {\nu^{x}} {\tau} {u^{\tau}} \left(u^{x}\right)^3 {u^{y}}\\
  &\quad - 4 {\mathcal{A}} {D_{s}} {(\partial_{\tau} u^{y})} {\nu^{y}} {\tau} {u^{\tau}} \left(u^{x}\right)^3 {u^{y}} - 14 {\mathcal{B}} {D_{s}} {(\partial_{\tau} u^{y})} {\nu^{y}} {\tau} {u^{\tau}} \left(u^{x}\right)^3 {u^{y}}\\
  &\quad - 2 {\mathcal{A}} {D_{s}} {(\partial_{x} u^{x})} {\nu^{x}} {\tau} \left(u^{x}\right)^4 {u^{y}} - 6 {\mathcal{B}} {D_{s}} {(\partial_{x} u^{x})} {\nu^{x}} {\tau} \left(u^{x}\right)^4 {u^{y}}\\
  &\quad + 3 {\mathcal{A}} {D_{s}} {(\partial_{y} u^{y})} {\nu^{x}} {\tau} \left(u^{x}\right)^4 {u^{y}} + 6 {\mathcal{B}} {D_{s}} {(\partial_{y} u^{y})} {\nu^{x}} {\tau} \left(u^{x}\right)^4 {u^{y}}\\
  &\quad - 4 {\mathcal{A}} {D_{s}} {(\partial_{x} u^{y})} {\nu^{y}} {\tau} \left(u^{x}\right)^4 {u^{y}} - 15 {\mathcal{B}} {D_{s}} {(\partial_{x} u^{y})} {\nu^{y}} {\tau} \left(u^{x}\right)^4 {u^{y}}\\
  &\quad + 3 {\mathcal{A}} {D_{s}} {(\partial_{y} u^{x})} {\nu^{y}} {\tau} \left(u^{x}\right)^4 {u^{y}} + 6 {\mathcal{B}} {D_{s}} {(\partial_{y} u^{x})} {\nu^{y}} {\tau} \left(u^{x}\right)^4 {u^{y}}\\
  &\quad + 3 {\mathcal{C}'(T)} {D_{s}} {(\partial_{y} T)} {\pi_{Q}^{xy}} {\tau} \left(u^{x}\right)^4 {u^{y}} - 2 {\mathcal{B}} {D_{s}} {\nu^{x}} {u^{\tau}} \left(u^{x}\right)^4 {u^{y}}\\
  &\quad + {D_{s}} {(\partial_{\tau} u^{x})} {I_{31}} {\tau} {u^{\tau}} \left(u^{x}\right)^4 {u^{y}} + {\mathcal{B}'(T)} {D_{s}} {(\partial_{\tau} T)} {\nu^{x}} {\tau} {u^{\tau}} \left(u^{x}\right)^4 {u^{y}}\\
  &\quad + {\mathcal{C}} {D_{s}} {(\partial_{\tau} u^{x})} {\Pi_{Q}} {\tau} {u^{\tau}} \left(u^{x}\right)^4 {u^{y}} + 4 {\mathcal{C}} {D_{s}} {(\partial_{\tau} u^{y})} {\pi_{Q}^{xy}} {\tau} {u^{\tau}} \left(u^{x}\right)^4 {u^{y}}\\
  &\quad + 8 {\mathcal{C}} {D_{s}} {(\partial_{\tau} u^{x})} {\pi_{Q}^{yy}} {\tau} {u^{\tau}} \left(u^{x}\right)^4 {u^{y}} + 5 {\mathcal{C}} {D_{s}} {(\partial_{\tau} u^{x})} {\pi_{Q}{}^{\eta}{}_{\eta}} {\tau} {u^{\tau}} \left(u^{x}\right)^4 {u^{y}}\\
  &\quad + {D_{s}} {(\partial_{x} u^{x})} {I_{31}} {\tau} \left(u^{x}\right)^5 {u^{y}} + {D_{s}} {(\partial_{y} u^{y})} {I_{31}} {\tau} \left(u^{x}\right)^5 {u^{y}}\\
  &\quad + {\mathcal{B}'(T)} {D_{s}} {(\partial_{x} T)} {\nu^{x}} {\tau} \left(u^{x}\right)^5 {u^{y}} + {\mathcal{B}'(T)} {D_{s}} {(\partial_{y} T)} {\nu^{y}} {\tau} \left(u^{x}\right)^5 {u^{y}}\\
  &\quad + {\mathcal{C}} {D_{s}} {(\partial_{x} u^{x})} {\Pi_{Q}} {\tau} \left(u^{x}\right)^5 {u^{y}} + {\mathcal{C}} {D_{s}} {(\partial_{y} u^{y})} {\Pi_{Q}} {\tau} \left(u^{x}\right)^5 {u^{y}}\\
  &\quad + 4 {\mathcal{C}} {D_{s}} {(\partial_{x} u^{y})} {\pi_{Q}^{xy}} {\tau} \left(u^{x}\right)^5 {u^{y}} + 8 {\mathcal{C}} {D_{s}} {(\partial_{x} u^{x})} {\pi_{Q}^{yy}} {\tau} \left(u^{x}\right)^5 {u^{y}}\\
  &\quad - 5 {\mathcal{C}} {D_{s}} {(\partial_{y} u^{y})} {\pi_{Q}^{yy}} {\tau} \left(u^{x}\right)^5 {u^{y}} + 5 {\mathcal{C}} {D_{s}} {(\partial_{x} u^{x})} {\pi_{Q}{}^{\eta}{}_{\eta}} {\tau} \left(u^{x}\right)^5 {u^{y}}\\
  &\quad - 2 {D_{s}} {I_{31}} {u^{\tau}} \left(u^{x}\right)^5 {u^{y}} - 2 {\mathcal{C}} {D_{s}} {\Pi_{Q}} {u^{\tau}} \left(u^{x}\right)^5 {u^{y}} - 2 {\mathcal{C}} {D_{s}} {\pi_{Q}{}^{\eta}{}_{\eta}} {u^{\tau}} \left(u^{x}\right)^5 {u^{y}}\\
  &\quad - 2 {\mathcal{A}} {D_{s}} {(\partial_{\tau} u^{y})} {\nu^{y}} {\tau} {u^{\tau}} \left(u^{x}\right)^5 {u^{y}} - 6 {\mathcal{B}} {D_{s}} {(\partial_{\tau} u^{y})} {\nu^{y}} {\tau} {u^{\tau}} \left(u^{x}\right)^5 {u^{y}}\\
  &\quad - 2 {\mathcal{A}} {D_{s}} {(\partial_{x} u^{y})} {\nu^{y}} {\tau} \left(u^{x}\right)^6 {u^{y}} - 6 {\mathcal{B}} {D_{s}} {(\partial_{x} u^{y})} {\nu^{y}} {\tau} \left(u^{x}\right)^6 {u^{y}}\\
  &\quad + 9 {D_{s}} {(\partial_{x} u^{y})} {I_{31}} {\tau} \left(u^{y}\right)^2 + 12 {D_{s}} {(\partial_{y} u^{x})} {I_{31}} {\tau} \left(u^{y}\right)^2\\
  &\quad + 12 {\mathcal{B}'(T)} {D_{s}} {(\partial_{y} T)} {\nu^{x}} {\tau} \left(u^{y}\right)^2 + 9 {\mathcal{B}'(T)} {D_{s}} {(\partial_{x} T)} {\nu^{y}} {\tau} \left(u^{y}\right)^2\\
  &\quad + 9 {\mathcal{C}} {D_{s}} {(\partial_{x} u^{y})} {\Pi_{Q}} {\tau} \left(u^{y}\right)^2 + 12 {\mathcal{C}} {D_{s}} {(\partial_{y} u^{x})} {\Pi_{Q}} {\tau} \left(u^{y}\right)^2\\
  &\quad + 18 {\mathcal{C}} {D_{s}} {(\partial_{x} u^{x})} {\pi_{Q}^{xy}} {\tau} \left(u^{y}\right)^2 + 15 {\mathcal{C}} {D_{s}} {(\partial_{y} u^{y})} {\pi_{Q}^{xy}} {\tau} \left(u^{y}\right)^2\\
  &\quad - 6 {\mathcal{C}} {D_{s}} {(\partial_{x} u^{y})} {\pi_{Q}^{yy}} {\tau} \left(u^{y}\right)^2 + 12 {\mathcal{C}} {D_{s}} {(\partial_{y} u^{x})} {\pi_{Q}^{yy}} {\tau} \left(u^{y}\right)^2\\
  &\quad - 9 {\mathcal{C}} {D_{s}} {(\partial_{x} u^{y})} {\pi_{Q}{}^{\eta}{}_{\eta}} {\tau} \left(u^{y}\right)^2 + 3 {\mathcal{C}} {D_{s}} {(\partial_{y} u^{x})} {\pi_{Q}{}^{\eta}{}_{\eta}} {\tau} \left(u^{y}\right)^2\\
  &\quad + 18 {\pi_{Q}^{xy}} {T} {\tau} \left(u^{y}\right)^2 + 9 {\mathcal{C}} {D_{s}} {\pi_{Q}^{xy}} {u^{\tau}} \left(u^{y}\right)^2 + 9 {\mathcal{A}} {D_{s}} {(\partial_{\tau} u^{y})} {\nu^{x}} {\tau} {u^{\tau}} \left(u^{y}\right)^2\\
  &\quad + 9 {\mathcal{B}} {D_{s}} {(\partial_{\tau} u^{y})} {\nu^{x}} {\tau} {u^{\tau}} \left(u^{y}\right)^2 + 9 {\mathcal{A}} {D_{s}} {(\partial_{\tau} u^{x})} {\nu^{y}} {\tau} {u^{\tau}} \left(u^{y}\right)^2\\
  &\quad + 9 {\mathcal{B}} {D_{s}} {(\partial_{\tau} u^{x})} {\nu^{y}} {\tau} {u^{\tau}} \left(u^{y}\right)^2 + 9 {\mathcal{C}'(T)} {D_{s}} {(\partial_{\tau} T)} {\pi_{Q}^{xy}} {\tau} {u^{\tau}} \left(u^{y}\right)^2\\
  &\quad + 9 {\mathcal{A}} {D_{s}} {(\partial_{x} u^{y})} {\nu^{x}} {\tau} {u^{x}} \left(u^{y}\right)^2 + 12 {\mathcal{B}} {D_{s}} {(\partial_{x} u^{y})} {\nu^{x}} {\tau} {u^{x}} \left(u^{y}\right)^2\\
  &\quad - 2 {\mathcal{A}} {D_{s}} {(\partial_{y} u^{x})} {\nu^{x}} {\tau} {u^{x}} \left(u^{y}\right)^2 - 15 {\mathcal{B}} {D_{s}} {(\partial_{y} u^{x})} {\nu^{x}} {\tau} {u^{x}} \left(u^{y}\right)^2\\
  &\quad + 9 {\mathcal{A}} {D_{s}} {(\partial_{x} u^{x})} {\nu^{y}} {\tau} {u^{x}} \left(u^{y}\right)^2 + 12 {\mathcal{B}} {D_{s}} {(\partial_{x} u^{x})} {\nu^{y}} {\tau} {u^{x}} \left(u^{y}\right)^2\\
  &\quad - 2 {\mathcal{A}} {D_{s}} {(\partial_{y} u^{y})} {\nu^{y}} {\tau} {u^{x}} \left(u^{y}\right)^2 - 12 {\mathcal{B}} {D_{s}} {(\partial_{y} u^{y})} {\nu^{y}} {\tau} {u^{x}} \left(u^{y}\right)^2\\
  &\quad + 9 {\mathcal{C}'(T)} {D_{s}} {(\partial_{x} T)} {\pi_{Q}^{xy}} {\tau} {u^{x}} \left(u^{y}\right)^2 - 2 {\mathcal{B}} {D_{s}} {\nu^{y}} {u^{\tau}} {u^{x}} \left(u^{y}\right)^2\\
  &\quad + 7 {D_{s}} {(\partial_{\tau} u^{y})} {I_{31}} {\tau} {u^{\tau}} {u^{x}} \left(u^{y}\right)^2 + 7 {\mathcal{B}'(T)} {D_{s}} {(\partial_{\tau} T)} {\nu^{y}} {\tau} {u^{\tau}} {u^{x}} \left(u^{y}\right)^2\\
  &\quad + 7 {\mathcal{C}} {D_{s}} {(\partial_{\tau} u^{y})} {\Pi_{Q}} {\tau} {u^{\tau}} {u^{x}} \left(u^{y}\right)^2 - 2 {\mathcal{C}} {D_{s}} {(\partial_{\tau} u^{x})} {\pi_{Q}^{xy}} {\tau} {u^{\tau}} {u^{x}} \left(u^{y}\right)^2\\
  &\quad - 14 {\mathcal{C}} {D_{s}} {(\partial_{\tau} u^{y})} {\pi_{Q}^{yy}} {\tau} {u^{\tau}} {u^{x}} \left(u^{y}\right)^2 - 9 {\mathcal{C}} {D_{s}} {(\partial_{\tau} u^{y})} {\pi_{Q}{}^{\eta}{}_{\eta}} {\tau} {u^{\tau}} {u^{x}} \left(u^{y}\right)^2\\
  &\quad + 21 {D_{s}} {(\partial_{x} u^{y})} {I_{31}} {\tau} \left(u^{x}\right)^2 \left(u^{y}\right)^2 + 18 {D_{s}} {(\partial_{y} u^{x})} {I_{31}} {\tau} \left(u^{x}\right)^2 \left(u^{y}\right)^2\\
  &\quad + 18 {\mathcal{B}'(T)} {D_{s}} {(\partial_{y} T)} {\nu^{x}} {\tau} \left(u^{x}\right)^2 \left(u^{y}\right)^2\\
  &\quad + 21 {\mathcal{B}'(T)} {D_{s}} {(\partial_{x} T)} {\nu^{y}} {\tau} \left(u^{x}\right)^2 \left(u^{y}\right)^2\\
  &\quad + 21 {\mathcal{C}} {D_{s}} {(\partial_{x} u^{y})} {\Pi_{Q}} {\tau} \left(u^{x}\right)^2 \left(u^{y}\right)^2 + 18 {\mathcal{C}} {D_{s}} {(\partial_{y} u^{x})} {\Pi_{Q}} {\tau} \left(u^{x}\right)^2 \left(u^{y}\right)^2\\
  &\quad + 10 {\mathcal{C}} {D_{s}} {(\partial_{x} u^{x})} {\pi_{Q}^{xy}} {\tau} \left(u^{x}\right)^2 \left(u^{y}\right)^2 + 23 {\mathcal{C}} {D_{s}} {(\partial_{y} u^{y})} {\pi_{Q}^{xy}} {\tau} \left(u^{x}\right)^2 \left(u^{y}\right)^2\\
  &\quad - 23 {\mathcal{C}} {D_{s}} {(\partial_{x} u^{y})} {\pi_{Q}^{yy}} {\tau} \left(u^{x}\right)^2 \left(u^{y}\right)^2 + 23 {\mathcal{C}} {D_{s}} {(\partial_{y} u^{x})} {\pi_{Q}^{yy}} {\tau} \left(u^{x}\right)^2 \left(u^{y}\right)^2\\
  &\quad - 23 {\mathcal{C}} {D_{s}} {(\partial_{x} u^{y})} {\pi_{Q}{}^{\eta}{}_{\eta}} {\tau} \left(u^{x}\right)^2 \left(u^{y}\right)^2 + 6 {\mathcal{C}} {D_{s}} {(\partial_{y} u^{x})} {\pi_{Q}{}^{\eta}{}_{\eta}} {\tau} \left(u^{x}\right)^2 \left(u^{y}\right)^2\\
  &\quad + 24 {\pi_{Q}^{xy}} {T} {\tau} \left(u^{x}\right)^2 \left(u^{y}\right)^2 + 12 {\mathcal{C}} {D_{s}} {\pi_{Q}^{xy}} {u^{\tau}} \left(u^{x}\right)^2 \left(u^{y}\right)^2\\
  &\quad + 14 {\mathcal{A}} {D_{s}} {(\partial_{\tau} u^{y})} {\nu^{x}} {\tau} {u^{\tau}} \left(u^{x}\right)^2 \left(u^{y}\right)^2\\
  &\quad + 20 {\mathcal{B}} {D_{s}} {(\partial_{\tau} u^{y})} {\nu^{x}} {\tau} {u^{\tau}} \left(u^{x}\right)^2 \left(u^{y}\right)^2\\
  &\quad + 14 {\mathcal{A}} {D_{s}} {(\partial_{\tau} u^{x})} {\nu^{y}} {\tau} {u^{\tau}} \left(u^{x}\right)^2 \left(u^{y}\right)^2\\
  &\quad + 20 {\mathcal{B}} {D_{s}} {(\partial_{\tau} u^{x})} {\nu^{y}} {\tau} {u^{\tau}} \left(u^{x}\right)^2 \left(u^{y}\right)^2\\
  &\quad + 12 {\mathcal{C}'(T)} {D_{s}} {(\partial_{\tau} T)} {\pi_{Q}^{xy}} {\tau} {u^{\tau}} \left(u^{x}\right)^2 \left(u^{y}\right)^2\\
  &\quad + 14 {\mathcal{A}} {D_{s}} {(\partial_{x} u^{y})} {\nu^{x}} {\tau} \left(u^{x}\right)^3 \left(u^{y}\right)^2 + 21 {\mathcal{B}} {D_{s}} {(\partial_{x} u^{y})} {\nu^{x}} {\tau} \left(u^{x}\right)^3 \left(u^{y}\right)^2\\
  &\quad - 2 {\mathcal{A}} {D_{s}} {(\partial_{y} u^{x})} {\nu^{x}} {\tau} \left(u^{x}\right)^3 \left(u^{y}\right)^2 - 12 {\mathcal{B}} {D_{s}} {(\partial_{y} u^{x})} {\nu^{x}} {\tau} \left(u^{x}\right)^3 \left(u^{y}\right)^2\\
  &\quad + 14 {\mathcal{A}} {D_{s}} {(\partial_{x} u^{x})} {\nu^{y}} {\tau} \left(u^{x}\right)^3 \left(u^{y}\right)^2 + 21 {\mathcal{B}} {D_{s}} {(\partial_{x} u^{x})} {\nu^{y}} {\tau} \left(u^{x}\right)^3 \left(u^{y}\right)^2\\
  &\quad - 4 {\mathcal{A}} {D_{s}} {(\partial_{y} u^{y})} {\nu^{y}} {\tau} \left(u^{x}\right)^3 \left(u^{y}\right)^2 - 21 {\mathcal{B}} {D_{s}} {(\partial_{y} u^{y})} {\nu^{y}} {\tau} \left(u^{x}\right)^3 \left(u^{y}\right)^2\\
  &\quad + 12 {\mathcal{C}'(T)} {D_{s}} {(\partial_{x} T)} {\pi_{Q}^{xy}} {\tau} \left(u^{x}\right)^3 \left(u^{y}\right)^2 - 2 {\mathcal{B}} {D_{s}} {\nu^{y}} {u^{\tau}} \left(u^{x}\right)^3 \left(u^{y}\right)^2\\
  &\quad + 10 {D_{s}} {(\partial_{\tau} u^{y})} {I_{31}} {\tau} {u^{\tau}} \left(u^{x}\right)^3 \left(u^{y}\right)^2\\
  &\quad + 10 {\mathcal{B}'(T)} {D_{s}} {(\partial_{\tau} T)} {\nu^{y}} {\tau} {u^{\tau}} \left(u^{x}\right)^3 \left(u^{y}\right)^2\\
  &\quad + 10 {\mathcal{C}} {D_{s}} {(\partial_{\tau} u^{y})} {\Pi_{Q}} {\tau} {u^{\tau}} \left(u^{x}\right)^3 \left(u^{y}\right)^2\\
  &\quad - 5 {\mathcal{C}} {D_{s}} {(\partial_{\tau} u^{x})} {\pi_{Q}^{xy}} {\tau} {u^{\tau}} \left(u^{x}\right)^3 \left(u^{y}\right)^2\\
  &\quad - 22 {\mathcal{C}} {D_{s}} {(\partial_{\tau} u^{y})} {\pi_{Q}^{yy}} {\tau} {u^{\tau}} \left(u^{x}\right)^3 \left(u^{y}\right)^2\\
  &\quad - 14 {\mathcal{C}} {D_{s}} {(\partial_{\tau} u^{y})} {\pi_{Q}{}^{\eta}{}_{\eta}} {\tau} {u^{\tau}} \left(u^{x}\right)^3 \left(u^{y}\right)^2\\
  &\quad + 15 {D_{s}} {(\partial_{x} u^{y})} {I_{31}} {\tau} \left(u^{x}\right)^4 \left(u^{y}\right)^2 + 6 {D_{s}} {(\partial_{y} u^{x})} {I_{31}} {\tau} \left(u^{x}\right)^4 \left(u^{y}\right)^2\\
  &\quad + 6 {\mathcal{B}'(T)} {D_{s}} {(\partial_{y} T)} {\nu^{x}} {\tau} \left(u^{x}\right)^4 \left(u^{y}\right)^2\\
  &\quad + 15 {\mathcal{B}'(T)} {D_{s}} {(\partial_{x} T)} {\nu^{y}} {\tau} \left(u^{x}\right)^4 \left(u^{y}\right)^2\\
  &\quad + 15 {\mathcal{C}} {D_{s}} {(\partial_{x} u^{y})} {\Pi_{Q}} {\tau} \left(u^{x}\right)^4 \left(u^{y}\right)^2 + 6 {\mathcal{C}} {D_{s}} {(\partial_{y} u^{x})} {\Pi_{Q}} {\tau} \left(u^{x}\right)^4 \left(u^{y}\right)^2\\
  &\quad - 5 {\mathcal{C}} {D_{s}} {(\partial_{x} u^{x})} {\pi_{Q}^{xy}} {\tau} \left(u^{x}\right)^4 \left(u^{y}\right)^2 + 8 {\mathcal{C}} {D_{s}} {(\partial_{y} u^{y})} {\pi_{Q}^{xy}} {\tau} \left(u^{x}\right)^4 \left(u^{y}\right)^2\\
  &\quad - 25 {\mathcal{C}} {D_{s}} {(\partial_{x} u^{y})} {\pi_{Q}^{yy}} {\tau} \left(u^{x}\right)^4 \left(u^{y}\right)^2 + 11 {\mathcal{C}} {D_{s}} {(\partial_{y} u^{x})} {\pi_{Q}^{yy}} {\tau} \left(u^{x}\right)^4 \left(u^{y}\right)^2\\
  &\quad - 19 {\mathcal{C}} {D_{s}} {(\partial_{x} u^{y})} {\pi_{Q}{}^{\eta}{}_{\eta}} {\tau} \left(u^{x}\right)^4 \left(u^{y}\right)^2 + 3 {\mathcal{C}} {D_{s}} {(\partial_{y} u^{x})} {\pi_{Q}{}^{\eta}{}_{\eta}} {\tau} \left(u^{x}\right)^4 \left(u^{y}\right)^2\\
  &\quad + 6 {\pi_{Q}^{xy}} {T} {\tau} \left(u^{x}\right)^4 \left(u^{y}\right)^2 + 3 {\mathcal{C}} {D_{s}} {\pi_{Q}^{xy}} {u^{\tau}} \left(u^{x}\right)^4 \left(u^{y}\right)^2\\
  &\quad + 5 {\mathcal{A}} {D_{s}} {(\partial_{\tau} u^{y})} {\nu^{x}} {\tau} {u^{\tau}} \left(u^{x}\right)^4 \left(u^{y}\right)^2 + 9 {\mathcal{B}} {D_{s}} {(\partial_{\tau} u^{y})} {\nu^{x}} {\tau} {u^{\tau}} \left(u^{x}\right)^4 \left(u^{y}\right)^2\\
  &\quad + 5 {\mathcal{A}} {D_{s}} {(\partial_{\tau} u^{x})} {\nu^{y}} {\tau} {u^{\tau}} \left(u^{x}\right)^4 \left(u^{y}\right)^2 + 9 {\mathcal{B}} {D_{s}} {(\partial_{\tau} u^{x})} {\nu^{y}} {\tau} {u^{\tau}} \left(u^{x}\right)^4 \left(u^{y}\right)^2\\
  &\quad + 3 {\mathcal{C}'(T)} {D_{s}} {(\partial_{\tau} T)} {\pi_{Q}^{xy}} {\tau} {u^{\tau}} \left(u^{x}\right)^4 \left(u^{y}\right)^2\\
  &\quad + 5 {\mathcal{A}} {D_{s}} {(\partial_{x} u^{y})} {\nu^{x}} {\tau} \left(u^{x}\right)^5 \left(u^{y}\right)^2 + 9 {\mathcal{B}} {D_{s}} {(\partial_{x} u^{y})} {\nu^{x}} {\tau} \left(u^{x}\right)^5 \left(u^{y}\right)^2\\
  &\quad + 5 {\mathcal{A}} {D_{s}} {(\partial_{x} u^{x})} {\nu^{y}} {\tau} \left(u^{x}\right)^5 \left(u^{y}\right)^2 + 9 {\mathcal{B}} {D_{s}} {(\partial_{x} u^{x})} {\nu^{y}} {\tau} \left(u^{x}\right)^5 \left(u^{y}\right)^2\\
  &\quad - 2 {\mathcal{A}} {D_{s}} {(\partial_{y} u^{y})} {\nu^{y}} {\tau} \left(u^{x}\right)^5 \left(u^{y}\right)^2 - 9 {\mathcal{B}} {D_{s}} {(\partial_{y} u^{y})} {\nu^{y}} {\tau} \left(u^{x}\right)^5 \left(u^{y}\right)^2\\
  &\quad + 3 {\mathcal{C}'(T)} {D_{s}} {(\partial_{x} T)} {\pi_{Q}^{xy}} {\tau} \left(u^{x}\right)^5 \left(u^{y}\right)^2\\
  &\quad + 3 {D_{s}} {(\partial_{\tau} u^{y})} {I_{31}} {\tau} {u^{\tau}} \left(u^{x}\right)^5 \left(u^{y}\right)^2\\
  &\quad + 3 {\mathcal{B}'(T)} {D_{s}} {(\partial_{\tau} T)} {\nu^{y}} {\tau} {u^{\tau}} \left(u^{x}\right)^5 \left(u^{y}\right)^2\\
  &\quad + 3 {\mathcal{C}} {D_{s}} {(\partial_{\tau} u^{y})} {\Pi_{Q}} {\tau} {u^{\tau}} \left(u^{x}\right)^5 \left(u^{y}\right)^2\\
  &\quad - 8 {\mathcal{C}} {D_{s}} {(\partial_{\tau} u^{y})} {\pi_{Q}^{yy}} {\tau} {u^{\tau}} \left(u^{x}\right)^5 \left(u^{y}\right)^2\\
  &\quad - 5 {\mathcal{C}} {D_{s}} {(\partial_{\tau} u^{y})} {\pi_{Q}{}^{\eta}{}_{\eta}} {\tau} {u^{\tau}} \left(u^{x}\right)^5 \left(u^{y}\right)^2 + 3 {D_{s}} {(\partial_{x} u^{y})} {I_{31}} {\tau} \left(u^{x}\right)^6 \left(u^{y}\right)^2\\
  &\quad + 3 {\mathcal{B}'(T)} {D_{s}} {(\partial_{x} T)} {\nu^{y}} {\tau} \left(u^{x}\right)^6 \left(u^{y}\right)^2 + 3 {\mathcal{C}} {D_{s}} {(\partial_{x} u^{y})} {\Pi_{Q}} {\tau} \left(u^{x}\right)^6 \left(u^{y}\right)^2\\
  &\quad - 8 {\mathcal{C}} {D_{s}} {(\partial_{x} u^{y})} {\pi_{Q}^{yy}} {\tau} \left(u^{x}\right)^6 \left(u^{y}\right)^2 - 5 {\mathcal{C}} {D_{s}} {(\partial_{x} u^{y})} {\pi_{Q}{}^{\eta}{}_{\eta}} {\tau} \left(u^{x}\right)^6 \left(u^{y}\right)^2\\
  &\quad - 9 {\mathcal{B}} {D_{s}} {(\partial_{x} u^{x})} {\nu^{x}} {\tau} \left(u^{y}\right)^3 + 9 {\mathcal{A}} {D_{s}} {(\partial_{y} u^{y})} {\nu^{x}} {\tau} \left(u^{y}\right)^3\\
  &\quad + 9 {\mathcal{B}} {D_{s}} {(\partial_{y} u^{y})} {\nu^{x}} {\tau} \left(u^{y}\right)^3 - 6 {\mathcal{B}} {D_{s}} {(\partial_{x} u^{y})} {\nu^{y}} {\tau} \left(u^{y}\right)^3\\
  &\quad + 9 {\mathcal{A}} {D_{s}} {(\partial_{y} u^{x})} {\nu^{y}} {\tau} \left(u^{y}\right)^3 + 9 {\mathcal{B}} {D_{s}} {(\partial_{y} u^{x})} {\nu^{y}} {\tau} \left(u^{y}\right)^3\\
  &\quad + 9 {\mathcal{C}'(T)} {D_{s}} {(\partial_{y} T)} {\pi_{Q}^{xy}} {\tau} \left(u^{y}\right)^3 + 9 {D_{s}} {(\partial_{\tau} u^{x})} {I_{31}} {\tau} {u^{\tau}} \left(u^{y}\right)^3\\
  &\quad + 9 {\mathcal{B}'(T)} {D_{s}} {(\partial_{\tau} T)} {\nu^{x}} {\tau} {u^{\tau}} \left(u^{y}\right)^3 + 9 {\mathcal{C}} {D_{s}} {(\partial_{\tau} u^{x})} {\Pi_{Q}} {\tau} {u^{\tau}} \left(u^{y}\right)^3\\
  &\quad + 6 {\mathcal{C}} {D_{s}} {(\partial_{\tau} u^{y})} {\pi_{Q}^{xy}} {\tau} {u^{\tau}} \left(u^{y}\right)^3 + 12 {\mathcal{C}} {D_{s}} {(\partial_{\tau} u^{x})} {\pi_{Q}^{yy}} {\tau} {u^{\tau}} \left(u^{y}\right)^3\\
  &\quad + 9 {\mathcal{C}} {D_{s}} {(\partial_{\tau} u^{x})} {\pi_{Q}{}^{\eta}{}_{\eta}} {\tau} {u^{\tau}} \left(u^{y}\right)^3 + 3 {D_{s}} {(\partial_{x} u^{x})} {I_{31}} {\tau} {u^{x}} \left(u^{y}\right)^3\\
  &\quad + 3 {D_{s}} {(\partial_{y} u^{y})} {I_{31}} {\tau} {u^{x}} \left(u^{y}\right)^3 + 3 {\mathcal{B}'(T)} {D_{s}} {(\partial_{x} T)} {\nu^{x}} {\tau} {u^{x}} \left(u^{y}\right)^3\\
  &\quad + 3 {\mathcal{B}'(T)} {D_{s}} {(\partial_{y} T)} {\nu^{y}} {\tau} {u^{x}} \left(u^{y}\right)^3 + 3 {\mathcal{C}} {D_{s}} {(\partial_{x} u^{x})} {\Pi_{Q}} {\tau} {u^{x}} \left(u^{y}\right)^3\\
  &\quad + 3 {\mathcal{C}} {D_{s}} {(\partial_{y} u^{y})} {\Pi_{Q}} {\tau} {u^{x}} \left(u^{y}\right)^3 + 2 {\mathcal{C}} {D_{s}} {(\partial_{x} u^{y})} {\pi_{Q}^{xy}} {\tau} {u^{x}} \left(u^{y}\right)^3\\
  &\quad - 18 {\mathcal{C}} {D_{s}} {(\partial_{y} u^{x})} {\pi_{Q}^{xy}} {\tau} {u^{x}} \left(u^{y}\right)^3 + 10 {\mathcal{C}} {D_{s}} {(\partial_{x} u^{x})} {\pi_{Q}^{yy}} {\tau} {u^{x}} \left(u^{y}\right)^3\\
  &\quad - 13 {\mathcal{C}} {D_{s}} {(\partial_{y} u^{y})} {\pi_{Q}^{yy}} {\tau} {u^{x}} \left(u^{y}\right)^3 + 15 {\mathcal{C}} {D_{s}} {(\partial_{x} u^{x})} {\pi_{Q}{}^{\eta}{}_{\eta}} {\tau} {u^{x}} \left(u^{y}\right)^3\\
  &\quad - 3 {\mathcal{C}} {D_{s}} {(\partial_{y} u^{y})} {\pi_{Q}{}^{\eta}{}_{\eta}} {\tau} {u^{x}} \left(u^{y}\right)^3 - 6 {D_{s}} {I_{31}} {u^{\tau}} {u^{x}} \left(u^{y}\right)^3\\
  &\quad - 6 {\mathcal{C}} {D_{s}} {\Pi_{Q}} {u^{\tau}} {u^{x}} \left(u^{y}\right)^3 - 6 {\mathcal{C}} {D_{s}} {\pi_{Q}{}^{\eta}{}_{\eta}} {u^{\tau}} {u^{x}} \left(u^{y}\right)^3\\
  &\quad - 6 {\mathcal{A}} {D_{s}} {(\partial_{\tau} u^{x})} {\nu^{x}} {\tau} {u^{\tau}} {u^{x}} \left(u^{y}\right)^3 - 22 {\mathcal{B}} {D_{s}} {(\partial_{\tau} u^{x})} {\nu^{x}} {\tau} {u^{\tau}} {u^{x}} \left(u^{y}\right)^3\\
  &\quad - 4 {\mathcal{A}} {D_{s}} {(\partial_{\tau} u^{y})} {\nu^{y}} {\tau} {u^{\tau}} {u^{x}} \left(u^{y}\right)^3 - 14 {\mathcal{B}} {D_{s}} {(\partial_{\tau} u^{y})} {\nu^{y}} {\tau} {u^{\tau}} {u^{x}} \left(u^{y}\right)^3\\
  &\quad - 6 {\mathcal{A}} {D_{s}} {(\partial_{x} u^{x})} {\nu^{x}} {\tau} \left(u^{x}\right)^2 \left(u^{y}\right)^3 - 24 {\mathcal{B}} {D_{s}} {(\partial_{x} u^{x})} {\nu^{x}} {\tau} \left(u^{x}\right)^2 \left(u^{y}\right)^3\\
  &\quad + 14 {\mathcal{A}} {D_{s}} {(\partial_{y} u^{y})} {\nu^{x}} {\tau} \left(u^{x}\right)^2 \left(u^{y}\right)^3 + 24 {\mathcal{B}} {D_{s}} {(\partial_{y} u^{y})} {\nu^{x}} {\tau} \left(u^{x}\right)^2 \left(u^{y}\right)^3\\
  &\quad - 4 {\mathcal{A}} {D_{s}} {(\partial_{x} u^{y})} {\nu^{y}} {\tau} \left(u^{x}\right)^2 \left(u^{y}\right)^3 - 21 {\mathcal{B}} {D_{s}} {(\partial_{x} u^{y})} {\nu^{y}} {\tau} \left(u^{x}\right)^2 \left(u^{y}\right)^3\\
  &\quad + 14 {\mathcal{A}} {D_{s}} {(\partial_{y} u^{x})} {\nu^{y}} {\tau} \left(u^{x}\right)^2 \left(u^{y}\right)^3 + 24 {\mathcal{B}} {D_{s}} {(\partial_{y} u^{x})} {\nu^{y}} {\tau} \left(u^{x}\right)^2 \left(u^{y}\right)^3\\
  &\quad + 12 {\mathcal{C}'(T)} {D_{s}} {(\partial_{y} T)} {\pi_{Q}^{xy}} {\tau} \left(u^{x}\right)^2 \left(u^{y}\right)^3 - 4 {\mathcal{B}} {D_{s}} {\nu^{x}} {u^{\tau}} \left(u^{x}\right)^2 \left(u^{y}\right)^3\\
  &\quad + 8 {D_{s}} {(\partial_{\tau} u^{x})} {I_{31}} {\tau} {u^{\tau}} \left(u^{x}\right)^2 \left(u^{y}\right)^3\\
  &\quad + 8 {\mathcal{B}'(T)} {D_{s}} {(\partial_{\tau} T)} {\nu^{x}} {\tau} {u^{\tau}} \left(u^{x}\right)^2 \left(u^{y}\right)^3\\
  &\quad + 8 {\mathcal{C}} {D_{s}} {(\partial_{\tau} u^{x})} {\Pi_{Q}} {\tau} {u^{\tau}} \left(u^{x}\right)^2 \left(u^{y}\right)^3\\
  &\quad + 15 {\mathcal{C}} {D_{s}} {(\partial_{\tau} u^{y})} {\pi_{Q}^{xy}} {\tau} {u^{\tau}} \left(u^{x}\right)^2 \left(u^{y}\right)^3\\
  &\quad + 22 {\mathcal{C}} {D_{s}} {(\partial_{\tau} u^{x})} {\pi_{Q}^{yy}} {\tau} {u^{\tau}} \left(u^{x}\right)^2 \left(u^{y}\right)^3\\
  &\quad + 16 {\mathcal{C}} {D_{s}} {(\partial_{\tau} u^{x})} {\pi_{Q}{}^{\eta}{}_{\eta}} {\tau} {u^{\tau}} \left(u^{x}\right)^2 \left(u^{y}\right)^3 + 4 {D_{s}} {(\partial_{x} u^{x})} {I_{31}} {\tau} \left(u^{x}\right)^3 \left(u^{y}\right)^3\\
  &\quad + 4 {D_{s}} {(\partial_{y} u^{y})} {I_{31}} {\tau} \left(u^{x}\right)^3 \left(u^{y}\right)^3 + 4 {\mathcal{B}'(T)} {D_{s}} {(\partial_{x} T)} {\nu^{x}} {\tau} \left(u^{x}\right)^3 \left(u^{y}\right)^3\\
  &\quad + 4 {\mathcal{B}'(T)} {D_{s}} {(\partial_{y} T)} {\nu^{y}} {\tau} \left(u^{x}\right)^3 \left(u^{y}\right)^3 + 4 {\mathcal{C}} {D_{s}} {(\partial_{x} u^{x})} {\Pi_{Q}} {\tau} \left(u^{x}\right)^3 \left(u^{y}\right)^3\\
  &\quad + 4 {\mathcal{C}} {D_{s}} {(\partial_{y} u^{y})} {\Pi_{Q}} {\tau} \left(u^{x}\right)^3 \left(u^{y}\right)^3 + 13 {\mathcal{C}} {D_{s}} {(\partial_{x} u^{y})} {\pi_{Q}^{xy}} {\tau} \left(u^{x}\right)^3 \left(u^{y}\right)^3\\
  &\quad - 13 {\mathcal{C}} {D_{s}} {(\partial_{y} u^{x})} {\pi_{Q}^{xy}} {\tau} \left(u^{x}\right)^3 \left(u^{y}\right)^3 + 21 {\mathcal{C}} {D_{s}} {(\partial_{x} u^{x})} {\pi_{Q}^{yy}} {\tau} \left(u^{x}\right)^3 \left(u^{y}\right)^3\\
  &\quad - 21 {\mathcal{C}} {D_{s}} {(\partial_{y} u^{y})} {\pi_{Q}^{yy}} {\tau} \left(u^{x}\right)^3 \left(u^{y}\right)^3\\
  &\quad + 20 {\mathcal{C}} {D_{s}} {(\partial_{x} u^{x})} {\pi_{Q}{}^{\eta}{}_{\eta}} {\tau} \left(u^{x}\right)^3 \left(u^{y}\right)^3 - 6 {\mathcal{C}} {D_{s}} {(\partial_{y} u^{y})} {\pi_{Q}{}^{\eta}{}_{\eta}} {\tau} \left(u^{x}\right)^3 \left(u^{y}\right)^3\\
  &\quad - 8 {D_{s}} {I_{31}} {u^{\tau}} \left(u^{x}\right)^3 \left(u^{y}\right)^3 - 8 {\mathcal{C}} {D_{s}} {\Pi_{Q}} {u^{\tau}} \left(u^{x}\right)^3 \left(u^{y}\right)^3\\
  &\quad - 8 {\mathcal{C}} {D_{s}} {\pi_{Q}{}^{\eta}{}_{\eta}} {u^{\tau}} \left(u^{x}\right)^3 \left(u^{y}\right)^3 - 4 {\mathcal{A}} {D_{s}} {(\partial_{\tau} u^{x})} {\nu^{x}} {\tau} {u^{\tau}} \left(u^{x}\right)^3 \left(u^{y}\right)^3\\
  &\quad - 12 {\mathcal{B}} {D_{s}} {(\partial_{\tau} u^{x})} {\nu^{x}} {\tau} {u^{\tau}} \left(u^{x}\right)^3 \left(u^{y}\right)^3\\
  &\quad - 6 {\mathcal{A}} {D_{s}} {(\partial_{\tau} u^{y})} {\nu^{y}} {\tau} {u^{\tau}} \left(u^{x}\right)^3 \left(u^{y}\right)^3\\
  &\quad - 20 {\mathcal{B}} {D_{s}} {(\partial_{\tau} u^{y})} {\nu^{y}} {\tau} {u^{\tau}} \left(u^{x}\right)^3 \left(u^{y}\right)^3 - 4 {\mathcal{A}} {D_{s}} {(\partial_{x} u^{x})} {\nu^{x}} {\tau} \left(u^{x}\right)^4 \left(u^{y}\right)^3\\
  &\quad - 12 {\mathcal{B}} {D_{s}} {(\partial_{x} u^{x})} {\nu^{x}} {\tau} \left(u^{x}\right)^4 \left(u^{y}\right)^3 + 5 {\mathcal{A}} {D_{s}} {(\partial_{y} u^{y})} {\nu^{x}} {\tau} \left(u^{x}\right)^4 \left(u^{y}\right)^3\\
  &\quad + 12 {\mathcal{B}} {D_{s}} {(\partial_{y} u^{y})} {\nu^{x}} {\tau} \left(u^{x}\right)^4 \left(u^{y}\right)^3 - 6 {\mathcal{A}} {D_{s}} {(\partial_{x} u^{y})} {\nu^{y}} {\tau} \left(u^{x}\right)^4 \left(u^{y}\right)^3\\
  &\quad - 21 {\mathcal{B}} {D_{s}} {(\partial_{x} u^{y})} {\nu^{y}} {\tau} \left(u^{x}\right)^4 \left(u^{y}\right)^3 + 5 {\mathcal{A}} {D_{s}} {(\partial_{y} u^{x})} {\nu^{y}} {\tau} \left(u^{x}\right)^4 \left(u^{y}\right)^3\\
  &\quad + 12 {\mathcal{B}} {D_{s}} {(\partial_{y} u^{x})} {\nu^{y}} {\tau} \left(u^{x}\right)^4 \left(u^{y}\right)^3\\
  &\quad + 3 {\mathcal{C}'(T)} {D_{s}} {(\partial_{y} T)} {\pi_{Q}^{xy}} {\tau} \left(u^{x}\right)^4 \left(u^{y}\right)^3 - 2 {\mathcal{B}} {D_{s}} {\nu^{x}} {u^{\tau}} \left(u^{x}\right)^4 \left(u^{y}\right)^3\\
  &\quad + {D_{s}} {(\partial_{\tau} u^{x})} {I_{31}} {\tau} {u^{\tau}} \left(u^{x}\right)^4 \left(u^{y}\right)^3\\
  &\quad + {\mathcal{B}'(T)} {D_{s}} {(\partial_{\tau} T)} {\nu^{x}} {\tau} {u^{\tau}} \left(u^{x}\right)^4 \left(u^{y}\right)^3\\
  &\quad + {\mathcal{C}} {D_{s}} {(\partial_{\tau} u^{x})} {\Pi_{Q}} {\tau} {u^{\tau}} \left(u^{x}\right)^4 \left(u^{y}\right)^3 + 8 {\mathcal{C}} {D_{s}} {(\partial_{\tau} u^{y})} {\pi_{Q}^{xy}} {\tau} {u^{\tau}} \left(u^{x}\right)^4 \left(u^{y}\right)^3\\
  &\quad + 8 {\mathcal{C}} {D_{s}} {(\partial_{\tau} u^{x})} {\pi_{Q}^{yy}} {\tau} {u^{\tau}} \left(u^{x}\right)^4 \left(u^{y}\right)^3\\
  &\quad + 5 {\mathcal{C}} {D_{s}} {(\partial_{\tau} u^{x})} {\pi_{Q}{}^{\eta}{}_{\eta}} {\tau} {u^{\tau}} \left(u^{x}\right)^4 \left(u^{y}\right)^3 + {D_{s}} {(\partial_{x} u^{x})} {I_{31}} {\tau} \left(u^{x}\right)^5 \left(u^{y}\right)^3\\
  &\quad + {D_{s}} {(\partial_{y} u^{y})} {I_{31}} {\tau} \left(u^{x}\right)^5 \left(u^{y}\right)^3 + {\mathcal{B}'(T)} {D_{s}} {(\partial_{x} T)} {\nu^{x}} {\tau} \left(u^{x}\right)^5 \left(u^{y}\right)^3\\
  &\quad + {\mathcal{B}'(T)} {D_{s}} {(\partial_{y} T)} {\nu^{y}} {\tau} \left(u^{x}\right)^5 \left(u^{y}\right)^3 + {\mathcal{C}} {D_{s}} {(\partial_{x} u^{x})} {\Pi_{Q}} {\tau} \left(u^{x}\right)^5 \left(u^{y}\right)^3\\
  &\quad + {\mathcal{C}} {D_{s}} {(\partial_{y} u^{y})} {\Pi_{Q}} {\tau} \left(u^{x}\right)^5 \left(u^{y}\right)^3 + 8 {\mathcal{C}} {D_{s}} {(\partial_{x} u^{y})} {\pi_{Q}^{xy}} {\tau} \left(u^{x}\right)^5 \left(u^{y}\right)^3\\
  &\quad + 8 {\mathcal{C}} {D_{s}} {(\partial_{x} u^{x})} {\pi_{Q}^{yy}} {\tau} \left(u^{x}\right)^5 \left(u^{y}\right)^3 - 8 {\mathcal{C}} {D_{s}} {(\partial_{y} u^{y})} {\pi_{Q}^{yy}} {\tau} \left(u^{x}\right)^5 \left(u^{y}\right)^3\\
  &\quad + 5 {\mathcal{C}} {D_{s}} {(\partial_{x} u^{x})} {\pi_{Q}{}^{\eta}{}_{\eta}} {\tau} \left(u^{x}\right)^5 \left(u^{y}\right)^3 - 3 {\mathcal{C}} {D_{s}} {(\partial_{y} u^{y})} {\pi_{Q}{}^{\eta}{}_{\eta}} {\tau} \left(u^{x}\right)^5 \left(u^{y}\right)^3\\
  &\quad - 2 {D_{s}} {I_{31}} {u^{\tau}} \left(u^{x}\right)^5 \left(u^{y}\right)^3 - 2 {\mathcal{C}} {D_{s}} {\Pi_{Q}} {u^{\tau}} \left(u^{x}\right)^5 \left(u^{y}\right)^3\\
  &\quad - 2 {\mathcal{C}} {D_{s}} {\pi_{Q}{}^{\eta}{}_{\eta}} {u^{\tau}} \left(u^{x}\right)^5 \left(u^{y}\right)^3 - 2 {\mathcal{A}} {D_{s}} {(\partial_{\tau} u^{y})} {\nu^{y}} {\tau} {u^{\tau}} \left(u^{x}\right)^5 \left(u^{y}\right)^3\\
  &\quad - 6 {\mathcal{B}} {D_{s}} {(\partial_{\tau} u^{y})} {\nu^{y}} {\tau} {u^{\tau}} \left(u^{x}\right)^5 \left(u^{y}\right)^3 - 2 {\mathcal{A}} {D_{s}} {(\partial_{x} u^{y})} {\nu^{y}} {\tau} \left(u^{x}\right)^6 \left(u^{y}\right)^3\\
  &\quad - 6 {\mathcal{B}} {D_{s}} {(\partial_{x} u^{y})} {\nu^{y}} {\tau} \left(u^{x}\right)^6 \left(u^{y}\right)^3 + 9 {D_{s}} {(\partial_{x} u^{y})} {I_{31}} {\tau} \left(u^{y}\right)^4\\
  &\quad + 18 {D_{s}} {(\partial_{y} u^{x})} {I_{31}} {\tau} \left(u^{y}\right)^4 + 18 {\mathcal{B}'(T)} {D_{s}} {(\partial_{y} T)} {\nu^{x}} {\tau} \left(u^{y}\right)^4\\
  &\quad + 9 {\mathcal{B}'(T)} {D_{s}} {(\partial_{x} T)} {\nu^{y}} {\tau} \left(u^{y}\right)^4 + 9 {\mathcal{C}} {D_{s}} {(\partial_{x} u^{y})} {\Pi_{Q}} {\tau} \left(u^{y}\right)^4\\
  &\quad + 18 {\mathcal{C}} {D_{s}} {(\partial_{y} u^{x})} {\Pi_{Q}} {\tau} \left(u^{y}\right)^4 + 18 {\mathcal{C}} {D_{s}} {(\partial_{x} u^{x})} {\pi_{Q}^{xy}} {\tau} \left(u^{y}\right)^4\\
  &\quad + 12 {\mathcal{C}} {D_{s}} {(\partial_{y} u^{y})} {\pi_{Q}^{xy}} {\tau} \left(u^{y}\right)^4 - 3 {\mathcal{C}} {D_{s}} {(\partial_{x} u^{y})} {\pi_{Q}^{yy}} {\tau} \left(u^{y}\right)^4\\
  &\quad + 15 {\mathcal{C}} {D_{s}} {(\partial_{y} u^{x})} {\pi_{Q}^{yy}} {\tau} \left(u^{y}\right)^4 - 9 {\mathcal{C}} {D_{s}} {(\partial_{x} u^{y})} {\pi_{Q}{}^{\eta}{}_{\eta}} {\tau} \left(u^{y}\right)^4\\
  &\quad + 9 {\mathcal{C}} {D_{s}} {(\partial_{y} u^{x})} {\pi_{Q}{}^{\eta}{}_{\eta}} {\tau} \left(u^{y}\right)^4 + 18 {\pi_{Q}^{xy}} {T} {\tau} \left(u^{y}\right)^4 + 9 {\mathcal{C}} {D_{s}} {\pi_{Q}^{xy}} {u^{\tau}} \left(u^{y}\right)^4\\
  &\quad + 9 {\mathcal{A}} {D_{s}} {(\partial_{\tau} u^{y})} {\nu^{x}} {\tau} {u^{\tau}} \left(u^{y}\right)^4 + 9 {\mathcal{B}} {D_{s}} {(\partial_{\tau} u^{y})} {\nu^{x}} {\tau} {u^{\tau}} \left(u^{y}\right)^4\\
  &\quad + 9 {\mathcal{A}} {D_{s}} {(\partial_{\tau} u^{x})} {\nu^{y}} {\tau} {u^{\tau}} \left(u^{y}\right)^4 + 9 {\mathcal{B}} {D_{s}} {(\partial_{\tau} u^{x})} {\nu^{y}} {\tau} {u^{\tau}} \left(u^{y}\right)^4\\
  &\quad + 9 {\mathcal{C}'(T)} {D_{s}} {(\partial_{\tau} T)} {\pi_{Q}^{xy}} {\tau} {u^{\tau}} \left(u^{y}\right)^4 + 9 {\mathcal{A}} {D_{s}} {(\partial_{x} u^{y})} {\nu^{x}} {\tau} {u^{x}} \left(u^{y}\right)^4\\
  &\quad + 15 {\mathcal{B}} {D_{s}} {(\partial_{x} u^{y})} {\nu^{x}} {\tau} {u^{x}} \left(u^{y}\right)^4 - 6 {\mathcal{A}} {D_{s}} {(\partial_{y} u^{x})} {\nu^{x}} {\tau} {u^{x}} \left(u^{y}\right)^4\\
  &\quad - 27 {\mathcal{B}} {D_{s}} {(\partial_{y} u^{x})} {\nu^{x}} {\tau} {u^{x}} \left(u^{y}\right)^4 + 9 {\mathcal{A}} {D_{s}} {(\partial_{x} u^{x})} {\nu^{y}} {\tau} {u^{x}} \left(u^{y}\right)^4\\
  &\quad + 15 {\mathcal{B}} {D_{s}} {(\partial_{x} u^{x})} {\nu^{y}} {\tau} {u^{x}} \left(u^{y}\right)^4 - 4 {\mathcal{A}} {D_{s}} {(\partial_{y} u^{y})} {\nu^{y}} {\tau} {u^{x}} \left(u^{y}\right)^4\\
  &\quad - 15 {\mathcal{B}} {D_{s}} {(\partial_{y} u^{y})} {\nu^{y}} {\tau} {u^{x}} \left(u^{y}\right)^4 + 9 {\mathcal{C}'(T)} {D_{s}} {(\partial_{x} T)} {\pi_{Q}^{xy}} {\tau} {u^{x}} \left(u^{y}\right)^4\\
  &\quad - 4 {\mathcal{B}} {D_{s}} {\nu^{y}} {u^{\tau}} {u^{x}} \left(u^{y}\right)^4 + 5 {D_{s}} {(\partial_{\tau} u^{y})} {I_{31}} {\tau} {u^{\tau}} {u^{x}} \left(u^{y}\right)^4\\
  &\quad + 5 {\mathcal{B}'(T)} {D_{s}} {(\partial_{\tau} T)} {\nu^{y}} {\tau} {u^{\tau}} {u^{x}} \left(u^{y}\right)^4 + 5 {\mathcal{C}} {D_{s}} {(\partial_{\tau} u^{y})} {\Pi_{Q}} {\tau} {u^{\tau}} {u^{x}} \left(u^{y}\right)^4\\
  &\quad - 13 {\mathcal{C}} {D_{s}} {(\partial_{\tau} u^{x})} {\pi_{Q}^{xy}} {\tau} {u^{\tau}} {u^{x}} \left(u^{y}\right)^4 - 8 {\mathcal{C}} {D_{s}} {(\partial_{\tau} u^{y})} {\pi_{Q}^{yy}} {\tau} {u^{\tau}} {u^{x}} \left(u^{y}\right)^4\\
  &\quad - 9 {\mathcal{C}} {D_{s}} {(\partial_{\tau} u^{y})} {\pi_{Q}{}^{\eta}{}_{\eta}} {\tau} {u^{\tau}} {u^{x}} \left(u^{y}\right)^4 + 15 {D_{s}} {(\partial_{x} u^{y})} {I_{31}} {\tau} \left(u^{x}\right)^2 \left(u^{y}\right)^4\\
  &\quad + 18 {D_{s}} {(\partial_{y} u^{x})} {I_{31}} {\tau} \left(u^{x}\right)^2 \left(u^{y}\right)^4\\
  &\quad + 18 {\mathcal{B}'(T)} {D_{s}} {(\partial_{y} T)} {\nu^{x}} {\tau} \left(u^{x}\right)^2 \left(u^{y}\right)^4\\
  &\quad + 15 {\mathcal{B}'(T)} {D_{s}} {(\partial_{x} T)} {\nu^{y}} {\tau} \left(u^{x}\right)^2 \left(u^{y}\right)^4\\
  &\quad + 15 {\mathcal{C}} {D_{s}} {(\partial_{x} u^{y})} {\Pi_{Q}} {\tau} \left(u^{x}\right)^2 \left(u^{y}\right)^4 + 18 {\mathcal{C}} {D_{s}} {(\partial_{y} u^{x})} {\Pi_{Q}} {\tau} \left(u^{x}\right)^2 \left(u^{y}\right)^4\\
  &\quad - 7 {\mathcal{C}} {D_{s}} {(\partial_{x} u^{x})} {\pi_{Q}^{xy}} {\tau} \left(u^{x}\right)^2 \left(u^{y}\right)^4 + 19 {\mathcal{C}} {D_{s}} {(\partial_{y} u^{y})} {\pi_{Q}^{xy}} {\tau} \left(u^{x}\right)^2 \left(u^{y}\right)^4\\
  &\quad - 11 {\mathcal{C}} {D_{s}} {(\partial_{x} u^{y})} {\pi_{Q}^{yy}} {\tau} \left(u^{x}\right)^2 \left(u^{y}\right)^4 + 25 {\mathcal{C}} {D_{s}} {(\partial_{y} u^{x})} {\pi_{Q}^{yy}} {\tau} \left(u^{x}\right)^2 \left(u^{y}\right)^4\\
  &\quad - 19 {\mathcal{C}} {D_{s}} {(\partial_{x} u^{y})} {\pi_{Q}{}^{\eta}{}_{\eta}} {\tau} \left(u^{x}\right)^2 \left(u^{y}\right)^4\\
  &\quad + 12 {\mathcal{C}} {D_{s}} {(\partial_{y} u^{x})} {\pi_{Q}{}^{\eta}{}_{\eta}} {\tau} \left(u^{x}\right)^2 \left(u^{y}\right)^4 + 12 {\pi_{Q}^{xy}} {T} {\tau} \left(u^{x}\right)^2 \left(u^{y}\right)^4\\
  &\quad + 6 {\mathcal{C}} {D_{s}} {\pi_{Q}^{xy}} {u^{\tau}} \left(u^{x}\right)^2 \left(u^{y}\right)^4 + 10 {\mathcal{A}} {D_{s}} {(\partial_{\tau} u^{y})} {\nu^{x}} {\tau} {u^{\tau}} \left(u^{x}\right)^2 \left(u^{y}\right)^4\\
  &\quad + 20 {\mathcal{B}} {D_{s}} {(\partial_{\tau} u^{y})} {\nu^{x}} {\tau} {u^{\tau}} \left(u^{x}\right)^2 \left(u^{y}\right)^4\\
  &\quad + 10 {\mathcal{A}} {D_{s}} {(\partial_{\tau} u^{x})} {\nu^{y}} {\tau} {u^{\tau}} \left(u^{x}\right)^2 \left(u^{y}\right)^4\\
  &\quad + 20 {\mathcal{B}} {D_{s}} {(\partial_{\tau} u^{x})} {\nu^{y}} {\tau} {u^{\tau}} \left(u^{x}\right)^2 \left(u^{y}\right)^4\\
  &\quad + 6 {\mathcal{C}'(T)} {D_{s}} {(\partial_{\tau} T)} {\pi_{Q}^{xy}} {\tau} {u^{\tau}} \left(u^{x}\right)^2 \left(u^{y}\right)^4\\
  &\quad + 10 {\mathcal{A}} {D_{s}} {(\partial_{x} u^{y})} {\nu^{x}} {\tau} \left(u^{x}\right)^3 \left(u^{y}\right)^4 + 21 {\mathcal{B}} {D_{s}} {(\partial_{x} u^{y})} {\nu^{x}} {\tau} \left(u^{x}\right)^3 \left(u^{y}\right)^4\\
  &\quad - 4 {\mathcal{A}} {D_{s}} {(\partial_{y} u^{x})} {\nu^{x}} {\tau} \left(u^{x}\right)^3 \left(u^{y}\right)^4 - 15 {\mathcal{B}} {D_{s}} {(\partial_{y} u^{x})} {\nu^{x}} {\tau} \left(u^{x}\right)^3 \left(u^{y}\right)^4\\
  &\quad + 10 {\mathcal{A}} {D_{s}} {(\partial_{x} u^{x})} {\nu^{y}} {\tau} \left(u^{x}\right)^3 \left(u^{y}\right)^4 + 21 {\mathcal{B}} {D_{s}} {(\partial_{x} u^{x})} {\nu^{y}} {\tau} \left(u^{x}\right)^3 \left(u^{y}\right)^4\\
  &\quad - 6 {\mathcal{A}} {D_{s}} {(\partial_{y} u^{y})} {\nu^{y}} {\tau} \left(u^{x}\right)^3 \left(u^{y}\right)^4 - 21 {\mathcal{B}} {D_{s}} {(\partial_{y} u^{y})} {\nu^{y}} {\tau} \left(u^{x}\right)^3 \left(u^{y}\right)^4\\
  &\quad + 6 {\mathcal{C}'(T)} {D_{s}} {(\partial_{x} T)} {\pi_{Q}^{xy}} {\tau} \left(u^{x}\right)^3 \left(u^{y}\right)^4 - 2 {\mathcal{B}} {D_{s}} {\nu^{y}} {u^{\tau}} \left(u^{x}\right)^3 \left(u^{y}\right)^4\\
  &\quad + 4 {D_{s}} {(\partial_{\tau} u^{y})} {I_{31}} {\tau} {u^{\tau}} \left(u^{x}\right)^3 \left(u^{y}\right)^4\\
  &\quad + 4 {\mathcal{B}'(T)} {D_{s}} {(\partial_{\tau} T)} {\nu^{y}} {\tau} {u^{\tau}} \left(u^{x}\right)^3 \left(u^{y}\right)^4\\
  &\quad + 4 {\mathcal{C}} {D_{s}} {(\partial_{\tau} u^{y})} {\Pi_{Q}} {\tau} {u^{\tau}} \left(u^{x}\right)^3 \left(u^{y}\right)^4\\
  &\quad - 8 {\mathcal{C}} {D_{s}} {(\partial_{\tau} u^{x})} {\pi_{Q}^{xy}} {\tau} {u^{\tau}} \left(u^{x}\right)^3 \left(u^{y}\right)^4\\
  &\quad - 8 {\mathcal{C}} {D_{s}} {(\partial_{\tau} u^{y})} {\pi_{Q}^{yy}} {\tau} {u^{\tau}} \left(u^{x}\right)^3 \left(u^{y}\right)^4\\
  &\quad - 8 {\mathcal{C}} {D_{s}} {(\partial_{\tau} u^{y})} {\pi_{Q}{}^{\eta}{}_{\eta}} {\tau} {u^{\tau}} \left(u^{x}\right)^3 \left(u^{y}\right)^4 + 6 {D_{s}} {(\partial_{x} u^{y})} {I_{31}} {\tau} \left(u^{x}\right)^4 \left(u^{y}\right)^4\\
  &\quad + 3 {D_{s}} {(\partial_{y} u^{x})} {I_{31}} {\tau} \left(u^{x}\right)^4 \left(u^{y}\right)^4 + 3 {\mathcal{B}'(T)} {D_{s}} {(\partial_{y} T)} {\nu^{x}} {\tau} \left(u^{x}\right)^4 \left(u^{y}\right)^4\\
  &\quad + 6 {\mathcal{B}'(T)} {D_{s}} {(\partial_{x} T)} {\nu^{y}} {\tau} \left(u^{x}\right)^4 \left(u^{y}\right)^4 + 6 {\mathcal{C}} {D_{s}} {(\partial_{x} u^{y})} {\Pi_{Q}} {\tau} \left(u^{x}\right)^4 \left(u^{y}\right)^4\\
  &\quad + 3 {\mathcal{C}} {D_{s}} {(\partial_{y} u^{x})} {\Pi_{Q}} {\tau} \left(u^{x}\right)^4 \left(u^{y}\right)^4 - 8 {\mathcal{C}} {D_{s}} {(\partial_{x} u^{x})} {\pi_{Q}^{xy}} {\tau} \left(u^{x}\right)^4 \left(u^{y}\right)^4\\
  &\quad + 8 {\mathcal{C}} {D_{s}} {(\partial_{y} u^{y})} {\pi_{Q}^{xy}} {\tau} \left(u^{x}\right)^4 \left(u^{y}\right)^4 - 8 {\mathcal{C}} {D_{s}} {(\partial_{x} u^{y})} {\pi_{Q}^{yy}} {\tau} \left(u^{x}\right)^4 \left(u^{y}\right)^4\\
  &\quad + 8 {\mathcal{C}} {D_{s}} {(\partial_{y} u^{x})} {\pi_{Q}^{yy}} {\tau} \left(u^{x}\right)^4 \left(u^{y}\right)^4 - 10 {\mathcal{C}} {D_{s}} {(\partial_{x} u^{y})} {\pi_{Q}{}^{\eta}{}_{\eta}} {\tau} \left(u^{x}\right)^4 \left(u^{y}\right)^4\\
  &\quad + 3 {\mathcal{C}} {D_{s}} {(\partial_{y} u^{x})} {\pi_{Q}{}^{\eta}{}_{\eta}} {\tau} \left(u^{x}\right)^4 \left(u^{y}\right)^4 + 2 {\mathcal{A}} {D_{s}} {(\partial_{\tau} u^{y})} {\nu^{x}} {\tau} {u^{\tau}} \left(u^{x}\right)^4 \left(u^{y}\right)^4\\
  &\quad + 6 {\mathcal{B}} {D_{s}} {(\partial_{\tau} u^{y})} {\nu^{x}} {\tau} {u^{\tau}} \left(u^{x}\right)^4 \left(u^{y}\right)^4 + 2 {\mathcal{A}} {D_{s}} {(\partial_{\tau} u^{x})} {\nu^{y}} {\tau} {u^{\tau}} \left(u^{x}\right)^4 \left(u^{y}\right)^4\\
  &\quad + 6 {\mathcal{B}} {D_{s}} {(\partial_{\tau} u^{x})} {\nu^{y}} {\tau} {u^{\tau}} \left(u^{x}\right)^4 \left(u^{y}\right)^4 + 2 {\mathcal{A}} {D_{s}} {(\partial_{x} u^{y})} {\nu^{x}} {\tau} \left(u^{x}\right)^5 \left(u^{y}\right)^4\\
  &\quad + 6 {\mathcal{B}} {D_{s}} {(\partial_{x} u^{y})} {\nu^{x}} {\tau} \left(u^{x}\right)^5 \left(u^{y}\right)^4 + 2 {\mathcal{A}} {D_{s}} {(\partial_{x} u^{x})} {\nu^{y}} {\tau} \left(u^{x}\right)^5 \left(u^{y}\right)^4\\
  &\quad + 6 {\mathcal{B}} {D_{s}} {(\partial_{x} u^{x})} {\nu^{y}} {\tau} \left(u^{x}\right)^5 \left(u^{y}\right)^4 - 2 {\mathcal{A}} {D_{s}} {(\partial_{y} u^{y})} {\nu^{y}} {\tau} \left(u^{x}\right)^5 \left(u^{y}\right)^4\\
  &\quad - 6 {\mathcal{B}} {D_{s}} {(\partial_{y} u^{y})} {\nu^{y}} {\tau} \left(u^{x}\right)^5 \left(u^{y}\right)^4 - 9 {\mathcal{B}} {D_{s}} {(\partial_{x} u^{x})} {\nu^{x}} {\tau} \left(u^{y}\right)^5\\
  &\quad + 9 {\mathcal{A}} {D_{s}} {(\partial_{y} u^{y})} {\nu^{x}} {\tau} \left(u^{y}\right)^5 + 9 {\mathcal{B}} {D_{s}} {(\partial_{y} u^{y})} {\nu^{x}} {\tau} \left(u^{y}\right)^5\\
  &\quad - 3 {\mathcal{B}} {D_{s}} {(\partial_{x} u^{y})} {\nu^{y}} {\tau} \left(u^{y}\right)^5 + 9 {\mathcal{A}} {D_{s}} {(\partial_{y} u^{x})} {\nu^{y}} {\tau} \left(u^{y}\right)^5\\
  &\quad + 9 {\mathcal{B}} {D_{s}} {(\partial_{y} u^{x})} {\nu^{y}} {\tau} \left(u^{y}\right)^5 + 9 {\mathcal{C}'(T)} {D_{s}} {(\partial_{y} T)} {\pi_{Q}^{xy}} {\tau} \left(u^{y}\right)^5\\
  &\quad + 9 {D_{s}} {(\partial_{\tau} u^{x})} {I_{31}} {\tau} {u^{\tau}} \left(u^{y}\right)^5 + 9 {\mathcal{B}'(T)} {D_{s}} {(\partial_{\tau} T)} {\nu^{x}} {\tau} {u^{\tau}} \left(u^{y}\right)^5\\
  &\quad + 9 {\mathcal{C}} {D_{s}} {(\partial_{\tau} u^{x})} {\Pi_{Q}} {\tau} {u^{\tau}} \left(u^{y}\right)^5 + 3 {\mathcal{C}} {D_{s}} {(\partial_{\tau} u^{y})} {\pi_{Q}^{xy}} {\tau} {u^{\tau}} \left(u^{y}\right)^5\\
  &\quad + 6 {\mathcal{C}} {D_{s}} {(\partial_{\tau} u^{x})} {\pi_{Q}^{yy}} {\tau} {u^{\tau}} \left(u^{y}\right)^5 + 9 {\mathcal{C}} {D_{s}} {(\partial_{\tau} u^{x})} {\pi_{Q}{}^{\eta}{}_{\eta}} {\tau} {u^{\tau}} \left(u^{y}\right)^5\\
  &\quad + 3 {D_{s}} {(\partial_{x} u^{x})} {I_{31}} {\tau} {u^{x}} \left(u^{y}\right)^5 + 3 {D_{s}} {(\partial_{y} u^{y})} {I_{31}} {\tau} {u^{x}} \left(u^{y}\right)^5\\
  &\quad + 3 {\mathcal{B}'(T)} {D_{s}} {(\partial_{x} T)} {\nu^{x}} {\tau} {u^{x}} \left(u^{y}\right)^5 + 3 {\mathcal{B}'(T)} {D_{s}} {(\partial_{y} T)} {\nu^{y}} {\tau} {u^{x}} \left(u^{y}\right)^5\\
  &\quad + 3 {\mathcal{C}} {D_{s}} {(\partial_{x} u^{x})} {\Pi_{Q}} {\tau} {u^{x}} \left(u^{y}\right)^5 + 3 {\mathcal{C}} {D_{s}} {(\partial_{y} u^{y})} {\Pi_{Q}} {\tau} {u^{x}} \left(u^{y}\right)^5\\
  &\quad + {\mathcal{C}} {D_{s}} {(\partial_{x} u^{y})} {\pi_{Q}^{xy}} {\tau} {u^{x}} \left(u^{y}\right)^5 - 21 {\mathcal{C}} {D_{s}} {(\partial_{y} u^{x})} {\pi_{Q}^{xy}} {\tau} {u^{x}} \left(u^{y}\right)^5\\
  &\quad + 5 {\mathcal{C}} {D_{s}} {(\partial_{x} u^{x})} {\pi_{Q}^{yy}} {\tau} {u^{x}} \left(u^{y}\right)^5 - 8 {\mathcal{C}} {D_{s}} {(\partial_{y} u^{y})} {\pi_{Q}^{yy}} {\tau} {u^{x}} \left(u^{y}\right)^5\\
  &\quad + 15 {\mathcal{C}} {D_{s}} {(\partial_{x} u^{x})} {\pi_{Q}{}^{\eta}{}_{\eta}} {\tau} {u^{x}} \left(u^{y}\right)^5 - 6 {\mathcal{C}} {D_{s}} {(\partial_{y} u^{y})} {\pi_{Q}{}^{\eta}{}_{\eta}} {\tau} {u^{x}} \left(u^{y}\right)^5\\
  &\quad - 6 {D_{s}} {I_{31}} {u^{\tau}} {u^{x}} \left(u^{y}\right)^5 - 6 {\mathcal{C}} {D_{s}} {\Pi_{Q}} {u^{\tau}} {u^{x}} \left(u^{y}\right)^5 - 6 {\mathcal{C}} {D_{s}} {\pi_{Q}{}^{\eta}{}_{\eta}} {u^{\tau}} {u^{x}} \left(u^{y}\right)^5\\
  &\quad - 6 {\mathcal{A}} {D_{s}} {(\partial_{\tau} u^{x})} {\nu^{x}} {\tau} {u^{\tau}} {u^{x}} \left(u^{y}\right)^5 - 20 {\mathcal{B}} {D_{s}} {(\partial_{\tau} u^{x})} {\nu^{x}} {\tau} {u^{\tau}} {u^{x}} \left(u^{y}\right)^5\\
  &\quad - 2 {\mathcal{A}} {D_{s}} {(\partial_{\tau} u^{y})} {\nu^{y}} {\tau} {u^{\tau}} {u^{x}} \left(u^{y}\right)^5 - 6 {\mathcal{B}} {D_{s}} {(\partial_{\tau} u^{y})} {\nu^{y}} {\tau} {u^{\tau}} {u^{x}} \left(u^{y}\right)^5\\
  &\quad - 6 {\mathcal{A}} {D_{s}} {(\partial_{x} u^{x})} {\nu^{x}} {\tau} \left(u^{x}\right)^2 \left(u^{y}\right)^5 - 21 {\mathcal{B}} {D_{s}} {(\partial_{x} u^{x})} {\nu^{x}} {\tau} \left(u^{x}\right)^2 \left(u^{y}\right)^5\\
  &\quad + 10 {\mathcal{A}} {D_{s}} {(\partial_{y} u^{y})} {\nu^{x}} {\tau} \left(u^{x}\right)^2 \left(u^{y}\right)^5 + 21 {\mathcal{B}} {D_{s}} {(\partial_{y} u^{y})} {\nu^{x}} {\tau} \left(u^{x}\right)^2 \left(u^{y}\right)^5\\
  &\quad - 2 {\mathcal{A}} {D_{s}} {(\partial_{x} u^{y})} {\nu^{y}} {\tau} \left(u^{x}\right)^2 \left(u^{y}\right)^5 - 9 {\mathcal{B}} {D_{s}} {(\partial_{x} u^{y})} {\nu^{y}} {\tau} \left(u^{x}\right)^2 \left(u^{y}\right)^5\\
  &\quad + 10 {\mathcal{A}} {D_{s}} {(\partial_{y} u^{x})} {\nu^{y}} {\tau} \left(u^{x}\right)^2 \left(u^{y}\right)^5 + 21 {\mathcal{B}} {D_{s}} {(\partial_{y} u^{x})} {\nu^{y}} {\tau} \left(u^{x}\right)^2 \left(u^{y}\right)^5\\
  &\quad + 6 {\mathcal{C}'(T)} {D_{s}} {(\partial_{y} T)} {\pi_{Q}^{xy}} {\tau} \left(u^{x}\right)^2 \left(u^{y}\right)^5 - 2 {\mathcal{B}} {D_{s}} {\nu^{x}} {u^{\tau}} \left(u^{x}\right)^2 \left(u^{y}\right)^5\\
  &\quad + 4 {D_{s}} {(\partial_{\tau} u^{x})} {I_{31}} {\tau} {u^{\tau}} \left(u^{x}\right)^2 \left(u^{y}\right)^5\\
  &\quad + 4 {\mathcal{B}'(T)} {D_{s}} {(\partial_{\tau} T)} {\nu^{x}} {\tau} {u^{\tau}} \left(u^{x}\right)^2 \left(u^{y}\right)^5\\
  &\quad + 4 {\mathcal{C}} {D_{s}} {(\partial_{\tau} u^{x})} {\Pi_{Q}} {\tau} {u^{\tau}} \left(u^{x}\right)^2 \left(u^{y}\right)^5\\
  &\quad + 8 {\mathcal{C}} {D_{s}} {(\partial_{\tau} u^{y})} {\pi_{Q}^{xy}} {\tau} {u^{\tau}} \left(u^{x}\right)^2 \left(u^{y}\right)^5\\
  &\quad + 8 {\mathcal{C}} {D_{s}} {(\partial_{\tau} u^{x})} {\pi_{Q}^{yy}} {\tau} {u^{\tau}} \left(u^{x}\right)^2 \left(u^{y}\right)^5\\
  &\quad + 8 {\mathcal{C}} {D_{s}} {(\partial_{\tau} u^{x})} {\pi_{Q}{}^{\eta}{}_{\eta}} {\tau} {u^{\tau}} \left(u^{x}\right)^2 \left(u^{y}\right)^5 + 2 {D_{s}} {(\partial_{x} u^{x})} {I_{31}} {\tau} \left(u^{x}\right)^3 \left(u^{y}\right)^5\\
  &\quad + 2 {D_{s}} {(\partial_{y} u^{y})} {I_{31}} {\tau} \left(u^{x}\right)^3 \left(u^{y}\right)^5 + 2 {\mathcal{B}'(T)} {D_{s}} {(\partial_{x} T)} {\nu^{x}} {\tau} \left(u^{x}\right)^3 \left(u^{y}\right)^5\\
  &\quad + 2 {\mathcal{B}'(T)} {D_{s}} {(\partial_{y} T)} {\nu^{y}} {\tau} \left(u^{x}\right)^3 \left(u^{y}\right)^5 + 2 {\mathcal{C}} {D_{s}} {(\partial_{x} u^{x})} {\Pi_{Q}} {\tau} \left(u^{x}\right)^3 \left(u^{y}\right)^5\\
  &\quad + 2 {\mathcal{C}} {D_{s}} {(\partial_{y} u^{y})} {\Pi_{Q}} {\tau} \left(u^{x}\right)^3 \left(u^{y}\right)^5 + 8 {\mathcal{C}} {D_{s}} {(\partial_{x} u^{y})} {\pi_{Q}^{xy}} {\tau} \left(u^{x}\right)^3 \left(u^{y}\right)^5\\
  &\quad - 8 {\mathcal{C}} {D_{s}} {(\partial_{y} u^{x})} {\pi_{Q}^{xy}} {\tau} \left(u^{x}\right)^3 \left(u^{y}\right)^5 + 8 {\mathcal{C}} {D_{s}} {(\partial_{x} u^{x})} {\pi_{Q}^{yy}} {\tau} \left(u^{x}\right)^3 \left(u^{y}\right)^5\\
  &\quad - 8 {\mathcal{C}} {D_{s}} {(\partial_{y} u^{y})} {\pi_{Q}^{yy}} {\tau} \left(u^{x}\right)^3 \left(u^{y}\right)^5 + 10 {\mathcal{C}} {D_{s}} {(\partial_{x} u^{x})} {\pi_{Q}{}^{\eta}{}_{\eta}} {\tau} \left(u^{x}\right)^3 \left(u^{y}\right)^5\\
  &\quad - 6 {\mathcal{C}} {D_{s}} {(\partial_{y} u^{y})} {\pi_{Q}{}^{\eta}{}_{\eta}} {\tau} \left(u^{x}\right)^3 \left(u^{y}\right)^5 - 4 {D_{s}} {I_{31}} {u^{\tau}} \left(u^{x}\right)^3 \left(u^{y}\right)^5\\
  &\quad - 4 {\mathcal{C}} {D_{s}} {\Pi_{Q}} {u^{\tau}} \left(u^{x}\right)^3 \left(u^{y}\right)^5 - 4 {\mathcal{C}} {D_{s}} {\pi_{Q}{}^{\eta}{}_{\eta}} {u^{\tau}} \left(u^{x}\right)^3 \left(u^{y}\right)^5\\
  &\quad - 2 {\mathcal{A}} {D_{s}} {(\partial_{\tau} u^{x})} {\nu^{x}} {\tau} {u^{\tau}} \left(u^{x}\right)^3 \left(u^{y}\right)^5 - 6 {\mathcal{B}} {D_{s}} {(\partial_{\tau} u^{x})} {\nu^{x}} {\tau} {u^{\tau}} \left(u^{x}\right)^3 \left(u^{y}\right)^5\\
  &\quad - 2 {\mathcal{A}} {D_{s}} {(\partial_{\tau} u^{y})} {\nu^{y}} {\tau} {u^{\tau}} \left(u^{x}\right)^3 \left(u^{y}\right)^5 - 6 {\mathcal{B}} {D_{s}} {(\partial_{\tau} u^{y})} {\nu^{y}} {\tau} {u^{\tau}} \left(u^{x}\right)^3 \left(u^{y}\right)^5\\
  &\quad - 2 {\mathcal{A}} {D_{s}} {(\partial_{x} u^{x})} {\nu^{x}} {\tau} \left(u^{x}\right)^4 \left(u^{y}\right)^5 - 6 {\mathcal{B}} {D_{s}} {(\partial_{x} u^{x})} {\nu^{x}} {\tau} \left(u^{x}\right)^4 \left(u^{y}\right)^5\\
  &\quad + 2 {\mathcal{A}} {D_{s}} {(\partial_{y} u^{y})} {\nu^{x}} {\tau} \left(u^{x}\right)^4 \left(u^{y}\right)^5 + 6 {\mathcal{B}} {D_{s}} {(\partial_{y} u^{y})} {\nu^{x}} {\tau} \left(u^{x}\right)^4 \left(u^{y}\right)^5\\
  &\quad - 2 {\mathcal{A}} {D_{s}} {(\partial_{x} u^{y})} {\nu^{y}} {\tau} \left(u^{x}\right)^4 \left(u^{y}\right)^5 - 6 {\mathcal{B}} {D_{s}} {(\partial_{x} u^{y})} {\nu^{y}} {\tau} \left(u^{x}\right)^4 \left(u^{y}\right)^5\\
  &\quad + 2 {\mathcal{A}} {D_{s}} {(\partial_{y} u^{x})} {\nu^{y}} {\tau} \left(u^{x}\right)^4 \left(u^{y}\right)^5 + 6 {\mathcal{B}} {D_{s}} {(\partial_{y} u^{x})} {\nu^{y}} {\tau} \left(u^{x}\right)^4 \left(u^{y}\right)^5\\
  &\quad + 3 {D_{s}} {(\partial_{x} u^{y})} {I_{31}} {\tau} \left(u^{y}\right)^6 + 12 {D_{s}} {(\partial_{y} u^{x})} {I_{31}} {\tau} \left(u^{y}\right)^6\\
  &\quad + 12 {\mathcal{B}'(T)} {D_{s}} {(\partial_{y} T)} {\nu^{x}} {\tau} \left(u^{y}\right)^6 + 3 {\mathcal{B}'(T)} {D_{s}} {(\partial_{x} T)} {\nu^{y}} {\tau} \left(u^{y}\right)^6\\
  &\quad + 3 {\mathcal{C}} {D_{s}} {(\partial_{x} u^{y})} {\Pi_{Q}} {\tau} \left(u^{y}\right)^6 + 12 {\mathcal{C}} {D_{s}} {(\partial_{y} u^{x})} {\Pi_{Q}} {\tau} \left(u^{y}\right)^6\\
  &\quad + 6 {\mathcal{C}} {D_{s}} {(\partial_{x} u^{x})} {\pi_{Q}^{xy}} {\tau} \left(u^{y}\right)^6 + 3 {\mathcal{C}} {D_{s}} {(\partial_{y} u^{y})} {\pi_{Q}^{xy}} {\tau} \left(u^{y}\right)^6\\
  &\quad + 6 {\mathcal{C}} {D_{s}} {(\partial_{y} u^{x})} {\pi_{Q}^{yy}} {\tau} \left(u^{y}\right)^6 - 3 {\mathcal{C}} {D_{s}} {(\partial_{x} u^{y})} {\pi_{Q}{}^{\eta}{}_{\eta}} {\tau} \left(u^{y}\right)^6\\
  &\quad + 9 {\mathcal{C}} {D_{s}} {(\partial_{y} u^{x})} {\pi_{Q}{}^{\eta}{}_{\eta}} {\tau} \left(u^{y}\right)^6 + 6 {\pi_{Q}^{xy}} {T} {\tau} \left(u^{y}\right)^6 + 3 {\mathcal{C}} {D_{s}} {\pi_{Q}^{xy}} {u^{\tau}} \left(u^{y}\right)^6\\
  &\quad + 3 {\mathcal{A}} {D_{s}} {(\partial_{\tau} u^{y})} {\nu^{x}} {\tau} {u^{\tau}} \left(u^{y}\right)^6 + 3 {\mathcal{B}} {D_{s}} {(\partial_{\tau} u^{y})} {\nu^{x}} {\tau} {u^{\tau}} \left(u^{y}\right)^6\\
  &\quad + 3 {\mathcal{A}} {D_{s}} {(\partial_{\tau} u^{x})} {\nu^{y}} {\tau} {u^{\tau}} \left(u^{y}\right)^6 + 3 {\mathcal{B}} {D_{s}} {(\partial_{\tau} u^{x})} {\nu^{y}} {\tau} {u^{\tau}} \left(u^{y}\right)^6\\
  &\quad + 3 {\mathcal{C}'(T)} {D_{s}} {(\partial_{\tau} T)} {\pi_{Q}^{xy}} {\tau} {u^{\tau}} \left(u^{y}\right)^6 + 3 {\mathcal{A}} {D_{s}} {(\partial_{x} u^{y})} {\nu^{x}} {\tau} {u^{x}} \left(u^{y}\right)^6\\
  &\quad + 6 {\mathcal{B}} {D_{s}} {(\partial_{x} u^{y})} {\nu^{x}} {\tau} {u^{x}} \left(u^{y}\right)^6 - 6 {\mathcal{A}} {D_{s}} {(\partial_{y} u^{x})} {\nu^{x}} {\tau} {u^{x}} \left(u^{y}\right)^6\\
  &\quad - 21 {\mathcal{B}} {D_{s}} {(\partial_{y} u^{x})} {\nu^{x}} {\tau} {u^{x}} \left(u^{y}\right)^6 + 3 {\mathcal{A}} {D_{s}} {(\partial_{x} u^{x})} {\nu^{y}} {\tau} {u^{x}} \left(u^{y}\right)^6\\
  &\quad + 6 {\mathcal{B}} {D_{s}} {(\partial_{x} u^{x})} {\nu^{y}} {\tau} {u^{x}} \left(u^{y}\right)^6 - 2 {\mathcal{A}} {D_{s}} {(\partial_{y} u^{y})} {\nu^{y}} {\tau} {u^{x}} \left(u^{y}\right)^6\\
  &\quad - 6 {\mathcal{B}} {D_{s}} {(\partial_{y} u^{y})} {\nu^{y}} {\tau} {u^{x}} \left(u^{y}\right)^6 + 3 {\mathcal{C}'(T)} {D_{s}} {(\partial_{x} T)} {\pi_{Q}^{xy}} {\tau} {u^{x}} \left(u^{y}\right)^6\\
  &\quad - 2 {\mathcal{B}} {D_{s}} {\nu^{y}} {u^{\tau}} {u^{x}} \left(u^{y}\right)^6 + {D_{s}} {(\partial_{\tau} u^{y})} {I_{31}} {\tau} {u^{\tau}} {u^{x}} \left(u^{y}\right)^6\\
  &\quad + {\mathcal{B}'(T)} {D_{s}} {(\partial_{\tau} T)} {\nu^{y}} {\tau} {u^{\tau}} {u^{x}} \left(u^{y}\right)^6 + {\mathcal{C}} {D_{s}} {(\partial_{\tau} u^{y})} {\Pi_{Q}} {\tau} {u^{\tau}} {u^{x}} \left(u^{y}\right)^6\\
  &\quad - 8 {\mathcal{C}} {D_{s}} {(\partial_{\tau} u^{x})} {\pi_{Q}^{xy}} {\tau} {u^{\tau}} {u^{x}} \left(u^{y}\right)^6 - 3 {\mathcal{C}} {D_{s}} {(\partial_{\tau} u^{y})} {\pi_{Q}{}^{\eta}{}_{\eta}} {\tau} {u^{\tau}} {u^{x}} \left(u^{y}\right)^6\\
  &\quad + 3 {D_{s}} {(\partial_{x} u^{y})} {I_{31}} {\tau} \left(u^{x}\right)^2 \left(u^{y}\right)^6 + 6 {D_{s}} {(\partial_{y} u^{x})} {I_{31}} {\tau} \left(u^{x}\right)^2 \left(u^{y}\right)^6\\
  &\quad + 6 {\mathcal{B}'(T)} {D_{s}} {(\partial_{y} T)} {\nu^{x}} {\tau} \left(u^{x}\right)^2 \left(u^{y}\right)^6\\
  &\quad + 3 {\mathcal{B}'(T)} {D_{s}} {(\partial_{x} T)} {\nu^{y}} {\tau} \left(u^{x}\right)^2 \left(u^{y}\right)^6 + 3 {\mathcal{C}} {D_{s}} {(\partial_{x} u^{y})} {\Pi_{Q}} {\tau} \left(u^{x}\right)^2 \left(u^{y}\right)^6\\
  &\quad + 6 {\mathcal{C}} {D_{s}} {(\partial_{y} u^{x})} {\Pi_{Q}} {\tau} \left(u^{x}\right)^2 \left(u^{y}\right)^6 - 8 {\mathcal{C}} {D_{s}} {(\partial_{x} u^{x})} {\pi_{Q}^{xy}} {\tau} \left(u^{x}\right)^2 \left(u^{y}\right)^6\\
  &\quad + 8 {\mathcal{C}} {D_{s}} {(\partial_{y} u^{y})} {\pi_{Q}^{xy}} {\tau} \left(u^{x}\right)^2 \left(u^{y}\right)^6 + 8 {\mathcal{C}} {D_{s}} {(\partial_{y} u^{x})} {\pi_{Q}^{yy}} {\tau} \left(u^{x}\right)^2 \left(u^{y}\right)^6\\
  &\quad - 5 {\mathcal{C}} {D_{s}} {(\partial_{x} u^{y})} {\pi_{Q}{}^{\eta}{}_{\eta}} {\tau} \left(u^{x}\right)^2 \left(u^{y}\right)^6 + 6 {\mathcal{C}} {D_{s}} {(\partial_{y} u^{x})} {\pi_{Q}{}^{\eta}{}_{\eta}} {\tau} \left(u^{x}\right)^2 \left(u^{y}\right)^6\\
  &\quad + 2 {\mathcal{A}} {D_{s}} {(\partial_{\tau} u^{y})} {\nu^{x}} {\tau} {u^{\tau}} \left(u^{x}\right)^2 \left(u^{y}\right)^6 + 6 {\mathcal{B}} {D_{s}} {(\partial_{\tau} u^{y})} {\nu^{x}} {\tau} {u^{\tau}} \left(u^{x}\right)^2 \left(u^{y}\right)^6\\
  &\quad + 2 {\mathcal{A}} {D_{s}} {(\partial_{\tau} u^{x})} {\nu^{y}} {\tau} {u^{\tau}} \left(u^{x}\right)^2 \left(u^{y}\right)^6 + 6 {\mathcal{B}} {D_{s}} {(\partial_{\tau} u^{x})} {\nu^{y}} {\tau} {u^{\tau}} \left(u^{x}\right)^2 \left(u^{y}\right)^6\\
  &\quad + 2 {\mathcal{A}} {D_{s}} {(\partial_{x} u^{y})} {\nu^{x}} {\tau} \left(u^{x}\right)^3 \left(u^{y}\right)^6 + 6 {\mathcal{B}} {D_{s}} {(\partial_{x} u^{y})} {\nu^{x}} {\tau} \left(u^{x}\right)^3 \left(u^{y}\right)^6\\
  &\quad - 2 {\mathcal{A}} {D_{s}} {(\partial_{y} u^{x})} {\nu^{x}} {\tau} \left(u^{x}\right)^3 \left(u^{y}\right)^6 - 6 {\mathcal{B}} {D_{s}} {(\partial_{y} u^{x})} {\nu^{x}} {\tau} \left(u^{x}\right)^3 \left(u^{y}\right)^6\\
  &\quad + 2 {\mathcal{A}} {D_{s}} {(\partial_{x} u^{x})} {\nu^{y}} {\tau} \left(u^{x}\right)^3 \left(u^{y}\right)^6 + 6 {\mathcal{B}} {D_{s}} {(\partial_{x} u^{x})} {\nu^{y}} {\tau} \left(u^{x}\right)^3 \left(u^{y}\right)^6\\
  &\quad - 2 {\mathcal{A}} {D_{s}} {(\partial_{y} u^{y})} {\nu^{y}} {\tau} \left(u^{x}\right)^3 \left(u^{y}\right)^6 - 6 {\mathcal{B}} {D_{s}} {(\partial_{y} u^{y})} {\nu^{y}} {\tau} \left(u^{x}\right)^3 \left(u^{y}\right)^6\\
  &\quad - 3 {\mathcal{B}} {D_{s}} {(\partial_{x} u^{x})} {\nu^{x}} {\tau} \left(u^{y}\right)^7 + 3 {\mathcal{A}} {D_{s}} {(\partial_{y} u^{y})} {\nu^{x}} {\tau} \left(u^{y}\right)^7\\
  &\quad + 3 {\mathcal{B}} {D_{s}} {(\partial_{y} u^{y})} {\nu^{x}} {\tau} \left(u^{y}\right)^7 + 3 {\mathcal{A}} {D_{s}} {(\partial_{y} u^{x})} {\nu^{y}} {\tau} \left(u^{y}\right)^7\\
  &\quad + 3 {\mathcal{B}} {D_{s}} {(\partial_{y} u^{x})} {\nu^{y}} {\tau} \left(u^{y}\right)^7 + 3 {\mathcal{C}'(T)} {D_{s}} {(\partial_{y} T)} {\pi_{Q}^{xy}} {\tau} \left(u^{y}\right)^7\\
  &\quad + 3 {D_{s}} {(\partial_{\tau} u^{x})} {I_{31}} {\tau} {u^{\tau}} \left(u^{y}\right)^7 + 3 {\mathcal{B}'(T)} {D_{s}} {(\partial_{\tau} T)} {\nu^{x}} {\tau} {u^{\tau}} \left(u^{y}\right)^7\\
  &\quad + 3 {\mathcal{C}} {D_{s}} {(\partial_{\tau} u^{x})} {\Pi_{Q}} {\tau} {u^{\tau}} \left(u^{y}\right)^7 + 3 {\mathcal{C}} {D_{s}} {(\partial_{\tau} u^{x})} {\pi_{Q}{}^{\eta}{}_{\eta}} {\tau} {u^{\tau}} \left(u^{y}\right)^7\\
  &\quad + {D_{s}} {(\partial_{x} u^{x})} {I_{31}} {\tau} {u^{x}} \left(u^{y}\right)^7 + {D_{s}} {(\partial_{y} u^{y})} {I_{31}} {\tau} {u^{x}} \left(u^{y}\right)^7\\
  &\quad + {\mathcal{B}'(T)} {D_{s}} {(\partial_{x} T)} {\nu^{x}} {\tau} {u^{x}} \left(u^{y}\right)^7 + {\mathcal{B}'(T)} {D_{s}} {(\partial_{y} T)} {\nu^{y}} {\tau} {u^{x}} \left(u^{y}\right)^7\\
  &\quad + {\mathcal{C}} {D_{s}} {(\partial_{x} u^{x})} {\Pi_{Q}} {\tau} {u^{x}} \left(u^{y}\right)^7 + {\mathcal{C}} {D_{s}} {(\partial_{y} u^{y})} {\Pi_{Q}} {\tau} {u^{x}} \left(u^{y}\right)^7\\
  &\quad - 8 {\mathcal{C}} {D_{s}} {(\partial_{y} u^{x})} {\pi_{Q}^{xy}} {\tau} {u^{x}} \left(u^{y}\right)^7 + 5 {\mathcal{C}} {D_{s}} {(\partial_{x} u^{x})} {\pi_{Q}{}^{\eta}{}_{\eta}} {\tau} {u^{x}} \left(u^{y}\right)^7\\
  &\quad - 3 {\mathcal{C}} {D_{s}} {(\partial_{y} u^{y})} {\pi_{Q}{}^{\eta}{}_{\eta}} {\tau} {u^{x}} \left(u^{y}\right)^7 - 2 {D_{s}} {I_{31}} {u^{\tau}} {u^{x}} \left(u^{y}\right)^7\\
  &\quad - 2 {\mathcal{C}} {D_{s}} {\Pi_{Q}} {u^{\tau}} {u^{x}} \left(u^{y}\right)^7 - 2 {\mathcal{C}} {D_{s}} {\pi_{Q}{}^{\eta}{}_{\eta}} {u^{\tau}} {u^{x}} \left(u^{y}\right)^7\\
  &\quad - 2 {\mathcal{A}} {D_{s}} {(\partial_{\tau} u^{x})} {\nu^{x}} {\tau} {u^{\tau}} {u^{x}} \left(u^{y}\right)^7 - 6 {\mathcal{B}} {D_{s}} {(\partial_{\tau} u^{x})} {\nu^{x}} {\tau} {u^{\tau}} {u^{x}} \left(u^{y}\right)^7\\
  &\quad - 2 {\mathcal{A}} {D_{s}} {(\partial_{x} u^{x})} {\nu^{x}} {\tau} \left(u^{x}\right)^2 \left(u^{y}\right)^7 - 6 {\mathcal{B}} {D_{s}} {(\partial_{x} u^{x})} {\nu^{x}} {\tau} \left(u^{x}\right)^2 \left(u^{y}\right)^7\\
  &\quad + 2 {\mathcal{A}} {D_{s}} {(\partial_{y} u^{y})} {\nu^{x}} {\tau} \left(u^{x}\right)^2 \left(u^{y}\right)^7 + 6 {\mathcal{B}} {D_{s}} {(\partial_{y} u^{y})} {\nu^{x}} {\tau} \left(u^{x}\right)^2 \left(u^{y}\right)^7\\
  &\quad + 2 {\mathcal{A}} {D_{s}} {(\partial_{y} u^{x})} {\nu^{y}} {\tau} \left(u^{x}\right)^2 \left(u^{y}\right)^7 + 6 {\mathcal{B}} {D_{s}} {(\partial_{y} u^{x})} {\nu^{y}} {\tau} \left(u^{x}\right)^2 \left(u^{y}\right)^7\\
  &\quad + 3 {D_{s}} {(\partial_{y} u^{x})} {I_{31}} {\tau} \left(u^{y}\right)^8 + 3 {\mathcal{B}'(T)} {D_{s}} {(\partial_{y} T)} {\nu^{x}} {\tau} \left(u^{y}\right)^8\\
  &\quad + 3 {\mathcal{C}} {D_{s}} {(\partial_{y} u^{x})} {\Pi_{Q}} {\tau} \left(u^{y}\right)^8 + 3 {\mathcal{C}} {D_{s}} {(\partial_{y} u^{x})} {\pi_{Q}{}^{\eta}{}_{\eta}} {\tau} \left(u^{y}\right)^8\\
  &\quad - 2 {\mathcal{A}} {D_{s}} {(\partial_{y} u^{x})} {\nu^{x}} {\tau} {u^{x}} \left(u^{y}\right)^8 - 6 {\mathcal{B}} {D_{s}} {(\partial_{y} u^{x})} {\nu^{x}} {\tau} {u^{x}} \left(u^{y}\right)^8\Bigr) \\
S_{7} &= \frac{1}{6 {T} {\tau} \left(u^{\tau}\right)^4 \left(\left(u^{y}\right)^2+1\right)}\Bigl(5 {D_{s}} {(\partial_{x} u^{x})} {I_{31}} {\tau} + 5 {D_{s}} {(\partial_{y} u^{y})} {I_{31}} {\tau} + 5 {\mathcal{B}'(T)} {D_{s}} {(\partial_{x} T)} {\nu^{x}} {\tau}\\
  &\quad + 5 {\mathcal{B}'(T)} {D_{s}} {(\partial_{y} T)} {\nu^{y}} {\tau} + 5 {\mathcal{C}} {D_{s}} {(\partial_{x} u^{x})} {\Pi_{Q}} {\tau} + 5 {\mathcal{C}} {D_{s}} {(\partial_{y} u^{y})} {\Pi_{Q}} {\tau}\\
  &\quad + 2 {\mathcal{C}} {D_{s}} {(\partial_{x} u^{y})} {\pi_{Q}^{xy}} {\tau} + 2 {\mathcal{C}} {D_{s}} {(\partial_{y} u^{x})} {\pi_{Q}^{xy}} {\tau} - 2 {\mathcal{C}} {D_{s}} {(\partial_{x} u^{x})} {\pi_{Q}^{yy}} {\tau}\\
  &\quad + 2 {\mathcal{C}} {D_{s}} {(\partial_{y} u^{y})} {\pi_{Q}^{yy}} {\tau} - 2 {\mathcal{C}} {D_{s}} {(\partial_{x} u^{x})} {\pi_{Q}{}^{\eta}{}_{\eta}} {\tau} + 6 {\Pi_{Q}} {T} {\tau} + 5 {D_{s}} {I_{31}} {u^{\tau}}\\
  &\quad + 5 {\mathcal{C}} {D_{s}} {\Pi_{Q}} {u^{\tau}} + 2 {\mathcal{C}} {D_{s}} {\pi_{Q}{}^{\eta}{}_{\eta}} {u^{\tau}} + 3 {(\partial_{T} I_{31})} {D_{s}} {(\partial_{\tau} T)} {\tau} {u^{\tau}}\\
  &\quad + 2 {\mathcal{A}} {D_{s}} {(\partial_{\tau} u^{x})} {\nu^{x}} {\tau} {u^{\tau}} + 5 {\mathcal{B}} {D_{s}} {(\partial_{\tau} u^{x})} {\nu^{x}} {\tau} {u^{\tau}} + 2 {\mathcal{A}} {D_{s}} {(\partial_{\tau} u^{y})} {\nu^{y}} {\tau} {u^{\tau}}\\
  &\quad + 5 {\mathcal{B}} {D_{s}} {(\partial_{\tau} u^{y})} {\nu^{y}} {\tau} {u^{\tau}} + 3 {\mathcal{C}'(T)} {D_{s}} {(\partial_{\tau} T)} {\Pi_{Q}} {\tau} {u^{\tau}}\\
  &\quad + 3 {(\partial_{T} I_{31})} {D_{s}} {(\partial_{x} T)} {\tau} {u^{x}} + 2 {\mathcal{A}} {D_{s}} {(\partial_{x} u^{x})} {\nu^{x}} {\tau} {u^{x}}\\
  &\quad + 2 {\mathcal{A}} {D_{s}} {(\partial_{x} u^{y})} {\nu^{y}} {\tau} {u^{x}} + 3 {\mathcal{C}'(T)} {D_{s}} {(\partial_{x} T)} {\Pi_{Q}} {\tau} {u^{x}} + 5 {\mathcal{B}} {D_{s}} {\nu^{x}} {u^{\tau}} {u^{x}}\\
  &\quad + 5 {D_{s}} {(\partial_{\tau} u^{x})} {I_{31}} {\tau} {u^{\tau}} {u^{x}} + 5 {\mathcal{B}'(T)} {D_{s}} {(\partial_{\tau} T)} {\nu^{x}} {\tau} {u^{\tau}} {u^{x}}\\
  &\quad + 5 {\mathcal{C}} {D_{s}} {(\partial_{\tau} u^{x})} {\Pi_{Q}} {\tau} {u^{\tau}} {u^{x}} + 2 {\mathcal{C}} {D_{s}} {(\partial_{\tau} u^{y})} {\pi_{Q}^{xy}} {\tau} {u^{\tau}} {u^{x}}\\
  &\quad - 2 {\mathcal{C}} {D_{s}} {(\partial_{\tau} u^{x})} {\pi_{Q}^{yy}} {\tau} {u^{\tau}} {u^{x}} - 2 {\mathcal{C}} {D_{s}} {(\partial_{\tau} u^{x})} {\pi_{Q}{}^{\eta}{}_{\eta}} {\tau} {u^{\tau}} {u^{x}}\\
  &\quad + 10 {D_{s}} {(\partial_{x} u^{x})} {I_{31}} {\tau} \left(u^{x}\right)^2 + 10 {D_{s}} {(\partial_{y} u^{y})} {I_{31}} {\tau} \left(u^{x}\right)^2\\
  &\quad + 10 {\mathcal{B}'(T)} {D_{s}} {(\partial_{x} T)} {\nu^{x}} {\tau} \left(u^{x}\right)^2 + 10 {\mathcal{B}'(T)} {D_{s}} {(\partial_{y} T)} {\nu^{y}} {\tau} \left(u^{x}\right)^2\\
  &\quad + 10 {\mathcal{C}} {D_{s}} {(\partial_{x} u^{x})} {\Pi_{Q}} {\tau} \left(u^{x}\right)^2 + 10 {\mathcal{C}} {D_{s}} {(\partial_{y} u^{y})} {\Pi_{Q}} {\tau} \left(u^{x}\right)^2\\
  &\quad + 4 {\mathcal{C}} {D_{s}} {(\partial_{x} u^{y})} {\pi_{Q}^{xy}} {\tau} \left(u^{x}\right)^2 + 2 {\mathcal{C}} {D_{s}} {(\partial_{y} u^{x})} {\pi_{Q}^{xy}} {\tau} \left(u^{x}\right)^2\\
  &\quad - 4 {\mathcal{C}} {D_{s}} {(\partial_{x} u^{x})} {\pi_{Q}^{yy}} {\tau} \left(u^{x}\right)^2 + 4 {\mathcal{C}} {D_{s}} {(\partial_{y} u^{y})} {\pi_{Q}^{yy}} {\tau} \left(u^{x}\right)^2\\
  &\quad - 4 {\mathcal{C}} {D_{s}} {(\partial_{x} u^{x})} {\pi_{Q}{}^{\eta}{}_{\eta}} {\tau} \left(u^{x}\right)^2 + 12 {\Pi_{Q}} {T} {\tau} \left(u^{x}\right)^2 + 10 {D_{s}} {I_{31}} {u^{\tau}} \left(u^{x}\right)^2\\
  &\quad + 10 {\mathcal{C}} {D_{s}} {\Pi_{Q}} {u^{\tau}} \left(u^{x}\right)^2 + 4 {\mathcal{C}} {D_{s}} {\pi_{Q}{}^{\eta}{}_{\eta}} {u^{\tau}} \left(u^{x}\right)^2\\
  &\quad + 6 {(\partial_{T} I_{31})} {D_{s}} {(\partial_{\tau} T)} {\tau} {u^{\tau}} \left(u^{x}\right)^2 + 2 {\mathcal{A}} {D_{s}} {(\partial_{\tau} u^{x})} {\nu^{x}} {\tau} {u^{\tau}} \left(u^{x}\right)^2\\
  &\quad + 4 {\mathcal{A}} {D_{s}} {(\partial_{\tau} u^{y})} {\nu^{y}} {\tau} {u^{\tau}} \left(u^{x}\right)^2 + 5 {\mathcal{B}} {D_{s}} {(\partial_{\tau} u^{y})} {\nu^{y}} {\tau} {u^{\tau}} \left(u^{x}\right)^2\\
  &\quad + 6 {\mathcal{C}'(T)} {D_{s}} {(\partial_{\tau} T)} {\Pi_{Q}} {\tau} {u^{\tau}} \left(u^{x}\right)^2 + 6 {(\partial_{T} I_{31})} {D_{s}} {(\partial_{x} T)} {\tau} \left(u^{x}\right)^3\\
  &\quad + 2 {\mathcal{A}} {D_{s}} {(\partial_{x} u^{x})} {\nu^{x}} {\tau} \left(u^{x}\right)^3 + 4 {\mathcal{A}} {D_{s}} {(\partial_{x} u^{y})} {\nu^{y}} {\tau} \left(u^{x}\right)^3\\
  &\quad + 6 {\mathcal{C}'(T)} {D_{s}} {(\partial_{x} T)} {\Pi_{Q}} {\tau} \left(u^{x}\right)^3 + 5 {\mathcal{B}} {D_{s}} {\nu^{x}} {u^{\tau}} \left(u^{x}\right)^3\\
  &\quad + 5 {D_{s}} {(\partial_{\tau} u^{x})} {I_{31}} {\tau} {u^{\tau}} \left(u^{x}\right)^3 + 5 {\mathcal{B}'(T)} {D_{s}} {(\partial_{\tau} T)} {\nu^{x}} {\tau} {u^{\tau}} \left(u^{x}\right)^3\\
  &\quad + 5 {\mathcal{C}} {D_{s}} {(\partial_{\tau} u^{x})} {\Pi_{Q}} {\tau} {u^{\tau}} \left(u^{x}\right)^3 + 2 {\mathcal{C}} {D_{s}} {(\partial_{\tau} u^{y})} {\pi_{Q}^{xy}} {\tau} {u^{\tau}} \left(u^{x}\right)^3\\
  &\quad - 2 {\mathcal{C}} {D_{s}} {(\partial_{\tau} u^{x})} {\pi_{Q}^{yy}} {\tau} {u^{\tau}} \left(u^{x}\right)^3 - 2 {\mathcal{C}} {D_{s}} {(\partial_{\tau} u^{x})} {\pi_{Q}{}^{\eta}{}_{\eta}} {\tau} {u^{\tau}} \left(u^{x}\right)^3\\
  &\quad + 5 {D_{s}} {(\partial_{x} u^{x})} {I_{31}} {\tau} \left(u^{x}\right)^4 + 5 {D_{s}} {(\partial_{y} u^{y})} {I_{31}} {\tau} \left(u^{x}\right)^4\\
  &\quad + 5 {\mathcal{B}'(T)} {D_{s}} {(\partial_{x} T)} {\nu^{x}} {\tau} \left(u^{x}\right)^4 + 5 {\mathcal{B}'(T)} {D_{s}} {(\partial_{y} T)} {\nu^{y}} {\tau} \left(u^{x}\right)^4\\
  &\quad + 5 {\mathcal{C}} {D_{s}} {(\partial_{x} u^{x})} {\Pi_{Q}} {\tau} \left(u^{x}\right)^4 + 5 {\mathcal{C}} {D_{s}} {(\partial_{y} u^{y})} {\Pi_{Q}} {\tau} \left(u^{x}\right)^4\\
  &\quad + 2 {\mathcal{C}} {D_{s}} {(\partial_{x} u^{y})} {\pi_{Q}^{xy}} {\tau} \left(u^{x}\right)^4 - 2 {\mathcal{C}} {D_{s}} {(\partial_{x} u^{x})} {\pi_{Q}^{yy}} {\tau} \left(u^{x}\right)^4\\
  &\quad + 2 {\mathcal{C}} {D_{s}} {(\partial_{y} u^{y})} {\pi_{Q}^{yy}} {\tau} \left(u^{x}\right)^4 - 2 {\mathcal{C}} {D_{s}} {(\partial_{x} u^{x})} {\pi_{Q}{}^{\eta}{}_{\eta}} {\tau} \left(u^{x}\right)^4\\
  &\quad + 6 {\Pi_{Q}} {T} {\tau} \left(u^{x}\right)^4 + 5 {D_{s}} {I_{31}} {u^{\tau}} \left(u^{x}\right)^4 + 5 {\mathcal{C}} {D_{s}} {\Pi_{Q}} {u^{\tau}} \left(u^{x}\right)^4\\
  &\quad + 2 {\mathcal{C}} {D_{s}} {\pi_{Q}{}^{\eta}{}_{\eta}} {u^{\tau}} \left(u^{x}\right)^4 + 3 {(\partial_{T} I_{31})} {D_{s}} {(\partial_{\tau} T)} {\tau} {u^{\tau}} \left(u^{x}\right)^4\\
  &\quad + 2 {\mathcal{A}} {D_{s}} {(\partial_{\tau} u^{y})} {\nu^{y}} {\tau} {u^{\tau}} \left(u^{x}\right)^4 + 3 {\mathcal{C}'(T)} {D_{s}} {(\partial_{\tau} T)} {\Pi_{Q}} {\tau} {u^{\tau}} \left(u^{x}\right)^4\\
  &\quad + 3 {(\partial_{T} I_{31})} {D_{s}} {(\partial_{x} T)} {\tau} \left(u^{x}\right)^5 + 2 {\mathcal{A}} {D_{s}} {(\partial_{x} u^{y})} {\nu^{y}} {\tau} \left(u^{x}\right)^5\\
  &\quad + 3 {\mathcal{C}'(T)} {D_{s}} {(\partial_{x} T)} {\Pi_{Q}} {\tau} \left(u^{x}\right)^5 + 3 {(\partial_{T} I_{31})} {D_{s}} {(\partial_{y} T)} {\tau} {u^{y}}\\
  &\quad + 2 {\mathcal{A}} {D_{s}} {(\partial_{y} u^{x})} {\nu^{x}} {\tau} {u^{y}} + 2 {\mathcal{A}} {D_{s}} {(\partial_{y} u^{y})} {\nu^{y}} {\tau} {u^{y}}\\
  &\quad + 3 {\mathcal{C}'(T)} {D_{s}} {(\partial_{y} T)} {\Pi_{Q}} {\tau} {u^{y}} + 5 {\mathcal{B}} {D_{s}} {\nu^{y}} {u^{\tau}} {u^{y}} + 5 {D_{s}} {(\partial_{\tau} u^{y})} {I_{31}} {\tau} {u^{\tau}} {u^{y}}\\
  &\quad + 5 {\mathcal{B}'(T)} {D_{s}} {(\partial_{\tau} T)} {\nu^{y}} {\tau} {u^{\tau}} {u^{y}} + 5 {\mathcal{C}} {D_{s}} {(\partial_{\tau} u^{y})} {\Pi_{Q}} {\tau} {u^{\tau}} {u^{y}}\\
  &\quad + 2 {\mathcal{C}} {D_{s}} {(\partial_{\tau} u^{x})} {\pi_{Q}^{xy}} {\tau} {u^{\tau}} {u^{y}} + 2 {\mathcal{C}} {D_{s}} {(\partial_{\tau} u^{y})} {\pi_{Q}^{yy}} {\tau} {u^{\tau}} {u^{y}}\\
  &\quad + 2 {\mathcal{C}} {D_{s}} {(\partial_{x} u^{x})} {\pi_{Q}^{xy}} {\tau} {u^{x}} {u^{y}} - 2 {\mathcal{C}} {D_{s}} {(\partial_{y} u^{y})} {\pi_{Q}^{xy}} {\tau} {u^{x}} {u^{y}}\\
  &\quad + 2 {\mathcal{C}} {D_{s}} {(\partial_{x} u^{y})} {\pi_{Q}^{yy}} {\tau} {u^{x}} {u^{y}} - 2 {\mathcal{C}} {D_{s}} {(\partial_{y} u^{x})} {\pi_{Q}^{yy}} {\tau} {u^{x}} {u^{y}}\\
  &\quad + 2 {\mathcal{C}} {D_{s}} {(\partial_{x} u^{y})} {\pi_{Q}{}^{\eta}{}_{\eta}} {\tau} {u^{x}} {u^{y}} - 2 {\mathcal{A}} {D_{s}} {(\partial_{\tau} u^{y})} {\nu^{x}} {\tau} {u^{\tau}} {u^{x}} {u^{y}}\\
  &\quad - 5 {\mathcal{B}} {D_{s}} {(\partial_{\tau} u^{y})} {\nu^{x}} {\tau} {u^{\tau}} {u^{x}} {u^{y}} - 2 {\mathcal{A}} {D_{s}} {(\partial_{\tau} u^{x})} {\nu^{y}} {\tau} {u^{\tau}} {u^{x}} {u^{y}}\\
  &\quad - 5 {\mathcal{B}} {D_{s}} {(\partial_{\tau} u^{x})} {\nu^{y}} {\tau} {u^{\tau}} {u^{x}} {u^{y}} + 6 {(\partial_{T} I_{31})} {D_{s}} {(\partial_{y} T)} {\tau} \left(u^{x}\right)^2 {u^{y}}\\
  &\quad - 2 {\mathcal{A}} {D_{s}} {(\partial_{x} u^{y})} {\nu^{x}} {\tau} \left(u^{x}\right)^2 {u^{y}} + 2 {\mathcal{A}} {D_{s}} {(\partial_{y} u^{x})} {\nu^{x}} {\tau} \left(u^{x}\right)^2 {u^{y}}\\
  &\quad - 2 {\mathcal{A}} {D_{s}} {(\partial_{x} u^{x})} {\nu^{y}} {\tau} \left(u^{x}\right)^2 {u^{y}} + 4 {\mathcal{A}} {D_{s}} {(\partial_{y} u^{y})} {\nu^{y}} {\tau} \left(u^{x}\right)^2 {u^{y}}\\
  &\quad + 6 {\mathcal{C}'(T)} {D_{s}} {(\partial_{y} T)} {\Pi_{Q}} {\tau} \left(u^{x}\right)^2 {u^{y}} + 5 {\mathcal{B}} {D_{s}} {\nu^{y}} {u^{\tau}} \left(u^{x}\right)^2 {u^{y}}\\
  &\quad + 5 {D_{s}} {(\partial_{\tau} u^{y})} {I_{31}} {\tau} {u^{\tau}} \left(u^{x}\right)^2 {u^{y}} + 5 {\mathcal{B}'(T)} {D_{s}} {(\partial_{\tau} T)} {\nu^{y}} {\tau} {u^{\tau}} \left(u^{x}\right)^2 {u^{y}}\\
  &\quad + 5 {\mathcal{C}} {D_{s}} {(\partial_{\tau} u^{y})} {\Pi_{Q}} {\tau} {u^{\tau}} \left(u^{x}\right)^2 {u^{y}} + 2 {\mathcal{C}} {D_{s}} {(\partial_{\tau} u^{x})} {\pi_{Q}^{xy}} {\tau} {u^{\tau}} \left(u^{x}\right)^2 {u^{y}}\\
  &\quad + 4 {\mathcal{C}} {D_{s}} {(\partial_{\tau} u^{y})} {\pi_{Q}^{yy}} {\tau} {u^{\tau}} \left(u^{x}\right)^2 {u^{y}} + 2 {\mathcal{C}} {D_{s}} {(\partial_{\tau} u^{y})} {\pi_{Q}{}^{\eta}{}_{\eta}} {\tau} {u^{\tau}} \left(u^{x}\right)^2 {u^{y}}\\
  &\quad + 2 {\mathcal{C}} {D_{s}} {(\partial_{x} u^{x})} {\pi_{Q}^{xy}} {\tau} \left(u^{x}\right)^3 {u^{y}} - 2 {\mathcal{C}} {D_{s}} {(\partial_{y} u^{y})} {\pi_{Q}^{xy}} {\tau} \left(u^{x}\right)^3 {u^{y}}\\
  &\quad + 4 {\mathcal{C}} {D_{s}} {(\partial_{x} u^{y})} {\pi_{Q}^{yy}} {\tau} \left(u^{x}\right)^3 {u^{y}} - 2 {\mathcal{C}} {D_{s}} {(\partial_{y} u^{x})} {\pi_{Q}^{yy}} {\tau} \left(u^{x}\right)^3 {u^{y}}\\
  &\quad + 4 {\mathcal{C}} {D_{s}} {(\partial_{x} u^{y})} {\pi_{Q}{}^{\eta}{}_{\eta}} {\tau} \left(u^{x}\right)^3 {u^{y}} - 2 {\mathcal{A}} {D_{s}} {(\partial_{\tau} u^{y})} {\nu^{x}} {\tau} {u^{\tau}} \left(u^{x}\right)^3 {u^{y}}\\
  &\quad - 2 {\mathcal{A}} {D_{s}} {(\partial_{\tau} u^{x})} {\nu^{y}} {\tau} {u^{\tau}} \left(u^{x}\right)^3 {u^{y}} + 3 {(\partial_{T} I_{31})} {D_{s}} {(\partial_{y} T)} {\tau} \left(u^{x}\right)^4 {u^{y}}\\
  &\quad - 2 {\mathcal{A}} {D_{s}} {(\partial_{x} u^{y})} {\nu^{x}} {\tau} \left(u^{x}\right)^4 {u^{y}} - 2 {\mathcal{A}} {D_{s}} {(\partial_{x} u^{x})} {\nu^{y}} {\tau} \left(u^{x}\right)^4 {u^{y}}\\
  &\quad + 2 {\mathcal{A}} {D_{s}} {(\partial_{y} u^{y})} {\nu^{y}} {\tau} \left(u^{x}\right)^4 {u^{y}} + 3 {\mathcal{C}'(T)} {D_{s}} {(\partial_{y} T)} {\Pi_{Q}} {\tau} \left(u^{x}\right)^4 {u^{y}}\\
  &\quad + 2 {\mathcal{C}} {D_{s}} {(\partial_{\tau} u^{y})} {\pi_{Q}^{yy}} {\tau} {u^{\tau}} \left(u^{x}\right)^4 {u^{y}} + 2 {\mathcal{C}} {D_{s}} {(\partial_{\tau} u^{y})} {\pi_{Q}{}^{\eta}{}_{\eta}} {\tau} {u^{\tau}} \left(u^{x}\right)^4 {u^{y}}\\
  &\quad + 2 {\mathcal{C}} {D_{s}} {(\partial_{x} u^{y})} {\pi_{Q}^{yy}} {\tau} \left(u^{x}\right)^5 {u^{y}} + 2 {\mathcal{C}} {D_{s}} {(\partial_{x} u^{y})} {\pi_{Q}{}^{\eta}{}_{\eta}} {\tau} \left(u^{x}\right)^5 {u^{y}}\\
  &\quad + 15 {D_{s}} {(\partial_{x} u^{x})} {I_{31}} {\tau} \left(u^{y}\right)^2 + 15 {D_{s}} {(\partial_{y} u^{y})} {I_{31}} {\tau} \left(u^{y}\right)^2\\
  &\quad + 15 {\mathcal{B}'(T)} {D_{s}} {(\partial_{x} T)} {\nu^{x}} {\tau} \left(u^{y}\right)^2 + 15 {\mathcal{B}'(T)} {D_{s}} {(\partial_{y} T)} {\nu^{y}} {\tau} \left(u^{y}\right)^2\\
  &\quad + 15 {\mathcal{C}} {D_{s}} {(\partial_{x} u^{x})} {\Pi_{Q}} {\tau} \left(u^{y}\right)^2 + 15 {\mathcal{C}} {D_{s}} {(\partial_{y} u^{y})} {\Pi_{Q}} {\tau} \left(u^{y}\right)^2\\
  &\quad + 4 {\mathcal{C}} {D_{s}} {(\partial_{x} u^{y})} {\pi_{Q}^{xy}} {\tau} \left(u^{y}\right)^2 + 6 {\mathcal{C}} {D_{s}} {(\partial_{y} u^{x})} {\pi_{Q}^{xy}} {\tau} \left(u^{y}\right)^2\\
  &\quad - 4 {\mathcal{C}} {D_{s}} {(\partial_{x} u^{x})} {\pi_{Q}^{yy}} {\tau} \left(u^{y}\right)^2 + 4 {\mathcal{C}} {D_{s}} {(\partial_{y} u^{y})} {\pi_{Q}^{yy}} {\tau} \left(u^{y}\right)^2\\
  &\quad - 6 {\mathcal{C}} {D_{s}} {(\partial_{x} u^{x})} {\pi_{Q}{}^{\eta}{}_{\eta}} {\tau} \left(u^{y}\right)^2 + 18 {\Pi_{Q}} {T} {\tau} \left(u^{y}\right)^2 + 15 {D_{s}} {I_{31}} {u^{\tau}} \left(u^{y}\right)^2\\
  &\quad + 15 {\mathcal{C}} {D_{s}} {\Pi_{Q}} {u^{\tau}} \left(u^{y}\right)^2 + 6 {\mathcal{C}} {D_{s}} {\pi_{Q}{}^{\eta}{}_{\eta}} {u^{\tau}} \left(u^{y}\right)^2\\
  &\quad + 9 {(\partial_{T} I_{31})} {D_{s}} {(\partial_{\tau} T)} {\tau} {u^{\tau}} \left(u^{y}\right)^2 + 6 {\mathcal{A}} {D_{s}} {(\partial_{\tau} u^{x})} {\nu^{x}} {\tau} {u^{\tau}} \left(u^{y}\right)^2\\
  &\quad + 10 {\mathcal{B}} {D_{s}} {(\partial_{\tau} u^{x})} {\nu^{x}} {\tau} {u^{\tau}} \left(u^{y}\right)^2 + 4 {\mathcal{A}} {D_{s}} {(\partial_{\tau} u^{y})} {\nu^{y}} {\tau} {u^{\tau}} \left(u^{y}\right)^2\\
  &\quad + 5 {\mathcal{B}} {D_{s}} {(\partial_{\tau} u^{y})} {\nu^{y}} {\tau} {u^{\tau}} \left(u^{y}\right)^2 + 9 {\mathcal{C}'(T)} {D_{s}} {(\partial_{\tau} T)} {\Pi_{Q}} {\tau} {u^{\tau}} \left(u^{y}\right)^2\\
  &\quad + 9 {(\partial_{T} I_{31})} {D_{s}} {(\partial_{x} T)} {\tau} {u^{x}} \left(u^{y}\right)^2 + 6 {\mathcal{A}} {D_{s}} {(\partial_{x} u^{x})} {\nu^{x}} {\tau} {u^{x}} \left(u^{y}\right)^2\\
  &\quad - 2 {\mathcal{A}} {D_{s}} {(\partial_{y} u^{y})} {\nu^{x}} {\tau} {u^{x}} \left(u^{y}\right)^2 + 4 {\mathcal{A}} {D_{s}} {(\partial_{x} u^{y})} {\nu^{y}} {\tau} {u^{x}} \left(u^{y}\right)^2\\
  &\quad - 2 {\mathcal{A}} {D_{s}} {(\partial_{y} u^{x})} {\nu^{y}} {\tau} {u^{x}} \left(u^{y}\right)^2 + 9 {\mathcal{C}'(T)} {D_{s}} {(\partial_{x} T)} {\Pi_{Q}} {\tau} {u^{x}} \left(u^{y}\right)^2\\
  &\quad + 10 {\mathcal{B}} {D_{s}} {\nu^{x}} {u^{\tau}} {u^{x}} \left(u^{y}\right)^2 + 10 {D_{s}} {(\partial_{\tau} u^{x})} {I_{31}} {\tau} {u^{\tau}} {u^{x}} \left(u^{y}\right)^2\\
  &\quad + 10 {\mathcal{B}'(T)} {D_{s}} {(\partial_{\tau} T)} {\nu^{x}} {\tau} {u^{\tau}} {u^{x}} \left(u^{y}\right)^2 + 10 {\mathcal{C}} {D_{s}} {(\partial_{\tau} u^{x})} {\Pi_{Q}} {\tau} {u^{\tau}} {u^{x}} \left(u^{y}\right)^2\\
  &\quad - 4 {\mathcal{C}} {D_{s}} {(\partial_{\tau} u^{x})} {\pi_{Q}^{yy}} {\tau} {u^{\tau}} {u^{x}} \left(u^{y}\right)^2 - 4 {\mathcal{C}} {D_{s}} {(\partial_{\tau} u^{x})} {\pi_{Q}{}^{\eta}{}_{\eta}} {\tau} {u^{\tau}} {u^{x}} \left(u^{y}\right)^2\\
  &\quad + 20 {D_{s}} {(\partial_{x} u^{x})} {I_{31}} {\tau} \left(u^{x}\right)^2 \left(u^{y}\right)^2 + 20 {D_{s}} {(\partial_{y} u^{y})} {I_{31}} {\tau} \left(u^{x}\right)^2 \left(u^{y}\right)^2\\
  &\quad + 20 {\mathcal{B}'(T)} {D_{s}} {(\partial_{x} T)} {\nu^{x}} {\tau} \left(u^{x}\right)^2 \left(u^{y}\right)^2\\
  &\quad + 20 {\mathcal{B}'(T)} {D_{s}} {(\partial_{y} T)} {\nu^{y}} {\tau} \left(u^{x}\right)^2 \left(u^{y}\right)^2\\
  &\quad + 20 {\mathcal{C}} {D_{s}} {(\partial_{x} u^{x})} {\Pi_{Q}} {\tau} \left(u^{x}\right)^2 \left(u^{y}\right)^2 + 20 {\mathcal{C}} {D_{s}} {(\partial_{y} u^{y})} {\Pi_{Q}} {\tau} \left(u^{x}\right)^2 \left(u^{y}\right)^2\\
  &\quad + 2 {\mathcal{C}} {D_{s}} {(\partial_{x} u^{y})} {\pi_{Q}^{xy}} {\tau} \left(u^{x}\right)^2 \left(u^{y}\right)^2 + 4 {\mathcal{C}} {D_{s}} {(\partial_{y} u^{x})} {\pi_{Q}^{xy}} {\tau} \left(u^{x}\right)^2 \left(u^{y}\right)^2\\
  &\quad - 6 {\mathcal{C}} {D_{s}} {(\partial_{x} u^{x})} {\pi_{Q}^{yy}} {\tau} \left(u^{x}\right)^2 \left(u^{y}\right)^2 + 6 {\mathcal{C}} {D_{s}} {(\partial_{y} u^{y})} {\pi_{Q}^{yy}} {\tau} \left(u^{x}\right)^2 \left(u^{y}\right)^2\\
  &\quad - 8 {\mathcal{C}} {D_{s}} {(\partial_{x} u^{x})} {\pi_{Q}{}^{\eta}{}_{\eta}} {\tau} \left(u^{x}\right)^2 \left(u^{y}\right)^2 + 24 {\Pi_{Q}} {T} {\tau} \left(u^{x}\right)^2 \left(u^{y}\right)^2\\
  &\quad + 20 {D_{s}} {I_{31}} {u^{\tau}} \left(u^{x}\right)^2 \left(u^{y}\right)^2 + 20 {\mathcal{C}} {D_{s}} {\Pi_{Q}} {u^{\tau}} \left(u^{x}\right)^2 \left(u^{y}\right)^2\\
  &\quad + 8 {\mathcal{C}} {D_{s}} {\pi_{Q}{}^{\eta}{}_{\eta}} {u^{\tau}} \left(u^{x}\right)^2 \left(u^{y}\right)^2 + 12 {(\partial_{T} I_{31})} {D_{s}} {(\partial_{\tau} T)} {\tau} {u^{\tau}} \left(u^{x}\right)^2 \left(u^{y}\right)^2\\
  &\quad + 4 {\mathcal{A}} {D_{s}} {(\partial_{\tau} u^{x})} {\nu^{x}} {\tau} {u^{\tau}} \left(u^{x}\right)^2 \left(u^{y}\right)^2 + 6 {\mathcal{A}} {D_{s}} {(\partial_{\tau} u^{y})} {\nu^{y}} {\tau} {u^{\tau}} \left(u^{x}\right)^2 \left(u^{y}\right)^2\\
  &\quad + 5 {\mathcal{B}} {D_{s}} {(\partial_{\tau} u^{y})} {\nu^{y}} {\tau} {u^{\tau}} \left(u^{x}\right)^2 \left(u^{y}\right)^2\\
  &\quad + 12 {\mathcal{C}'(T)} {D_{s}} {(\partial_{\tau} T)} {\Pi_{Q}} {\tau} {u^{\tau}} \left(u^{x}\right)^2 \left(u^{y}\right)^2\\
  &\quad + 12 {(\partial_{T} I_{31})} {D_{s}} {(\partial_{x} T)} {\tau} \left(u^{x}\right)^3 \left(u^{y}\right)^2 + 4 {\mathcal{A}} {D_{s}} {(\partial_{x} u^{x})} {\nu^{x}} {\tau} \left(u^{x}\right)^3 \left(u^{y}\right)^2\\
  &\quad - 2 {\mathcal{A}} {D_{s}} {(\partial_{y} u^{y})} {\nu^{x}} {\tau} \left(u^{x}\right)^3 \left(u^{y}\right)^2 + 6 {\mathcal{A}} {D_{s}} {(\partial_{x} u^{y})} {\nu^{y}} {\tau} \left(u^{x}\right)^3 \left(u^{y}\right)^2\\
  &\quad - 2 {\mathcal{A}} {D_{s}} {(\partial_{y} u^{x})} {\nu^{y}} {\tau} \left(u^{x}\right)^3 \left(u^{y}\right)^2\\
  &\quad + 12 {\mathcal{C}'(T)} {D_{s}} {(\partial_{x} T)} {\Pi_{Q}} {\tau} \left(u^{x}\right)^3 \left(u^{y}\right)^2 + 5 {\mathcal{B}} {D_{s}} {\nu^{x}} {u^{\tau}} \left(u^{x}\right)^3 \left(u^{y}\right)^2\\
  &\quad + 5 {D_{s}} {(\partial_{\tau} u^{x})} {I_{31}} {\tau} {u^{\tau}} \left(u^{x}\right)^3 \left(u^{y}\right)^2\\
  &\quad + 5 {\mathcal{B}'(T)} {D_{s}} {(\partial_{\tau} T)} {\nu^{x}} {\tau} {u^{\tau}} \left(u^{x}\right)^3 \left(u^{y}\right)^2\\
  &\quad + 5 {\mathcal{C}} {D_{s}} {(\partial_{\tau} u^{x})} {\Pi_{Q}} {\tau} {u^{\tau}} \left(u^{x}\right)^3 \left(u^{y}\right)^2\\
  &\quad - 2 {\mathcal{C}} {D_{s}} {(\partial_{\tau} u^{y})} {\pi_{Q}^{xy}} {\tau} {u^{\tau}} \left(u^{x}\right)^3 \left(u^{y}\right)^2\\
  &\quad - 2 {\mathcal{C}} {D_{s}} {(\partial_{\tau} u^{x})} {\pi_{Q}^{yy}} {\tau} {u^{\tau}} \left(u^{x}\right)^3 \left(u^{y}\right)^2\\
  &\quad - 2 {\mathcal{C}} {D_{s}} {(\partial_{\tau} u^{x})} {\pi_{Q}{}^{\eta}{}_{\eta}} {\tau} {u^{\tau}} \left(u^{x}\right)^3 \left(u^{y}\right)^2 + 5 {D_{s}} {(\partial_{x} u^{x})} {I_{31}} {\tau} \left(u^{x}\right)^4 \left(u^{y}\right)^2\\
  &\quad + 5 {D_{s}} {(\partial_{y} u^{y})} {I_{31}} {\tau} \left(u^{x}\right)^4 \left(u^{y}\right)^2 + 5 {\mathcal{B}'(T)} {D_{s}} {(\partial_{x} T)} {\nu^{x}} {\tau} \left(u^{x}\right)^4 \left(u^{y}\right)^2\\
  &\quad + 5 {\mathcal{B}'(T)} {D_{s}} {(\partial_{y} T)} {\nu^{y}} {\tau} \left(u^{x}\right)^4 \left(u^{y}\right)^2 + 5 {\mathcal{C}} {D_{s}} {(\partial_{x} u^{x})} {\Pi_{Q}} {\tau} \left(u^{x}\right)^4 \left(u^{y}\right)^2\\
  &\quad + 5 {\mathcal{C}} {D_{s}} {(\partial_{y} u^{y})} {\Pi_{Q}} {\tau} \left(u^{x}\right)^4 \left(u^{y}\right)^2 - 2 {\mathcal{C}} {D_{s}} {(\partial_{x} u^{y})} {\pi_{Q}^{xy}} {\tau} \left(u^{x}\right)^4 \left(u^{y}\right)^2\\
  &\quad - 2 {\mathcal{C}} {D_{s}} {(\partial_{x} u^{x})} {\pi_{Q}^{yy}} {\tau} \left(u^{x}\right)^4 \left(u^{y}\right)^2 + 2 {\mathcal{C}} {D_{s}} {(\partial_{y} u^{y})} {\pi_{Q}^{yy}} {\tau} \left(u^{x}\right)^4 \left(u^{y}\right)^2\\
  &\quad - 2 {\mathcal{C}} {D_{s}} {(\partial_{x} u^{x})} {\pi_{Q}{}^{\eta}{}_{\eta}} {\tau} \left(u^{x}\right)^4 \left(u^{y}\right)^2 + 6 {\Pi_{Q}} {T} {\tau} \left(u^{x}\right)^4 \left(u^{y}\right)^2\\
  &\quad + 5 {D_{s}} {I_{31}} {u^{\tau}} \left(u^{x}\right)^4 \left(u^{y}\right)^2 + 5 {\mathcal{C}} {D_{s}} {\Pi_{Q}} {u^{\tau}} \left(u^{x}\right)^4 \left(u^{y}\right)^2\\
  &\quad + 2 {\mathcal{C}} {D_{s}} {\pi_{Q}{}^{\eta}{}_{\eta}} {u^{\tau}} \left(u^{x}\right)^4 \left(u^{y}\right)^2 + 3 {(\partial_{T} I_{31})} {D_{s}} {(\partial_{\tau} T)} {\tau} {u^{\tau}} \left(u^{x}\right)^4 \left(u^{y}\right)^2\\
  &\quad + 2 {\mathcal{A}} {D_{s}} {(\partial_{\tau} u^{y})} {\nu^{y}} {\tau} {u^{\tau}} \left(u^{x}\right)^4 \left(u^{y}\right)^2\\
  &\quad + 3 {\mathcal{C}'(T)} {D_{s}} {(\partial_{\tau} T)} {\Pi_{Q}} {\tau} {u^{\tau}} \left(u^{x}\right)^4 \left(u^{y}\right)^2\\
  &\quad + 3 {(\partial_{T} I_{31})} {D_{s}} {(\partial_{x} T)} {\tau} \left(u^{x}\right)^5 \left(u^{y}\right)^2 + 2 {\mathcal{A}} {D_{s}} {(\partial_{x} u^{y})} {\nu^{y}} {\tau} \left(u^{x}\right)^5 \left(u^{y}\right)^2\\
  &\quad + 3 {\mathcal{C}'(T)} {D_{s}} {(\partial_{x} T)} {\Pi_{Q}} {\tau} \left(u^{x}\right)^5 \left(u^{y}\right)^2 + 9 {(\partial_{T} I_{31})} {D_{s}} {(\partial_{y} T)} {\tau} \left(u^{y}\right)^3\\
  &\quad + 6 {\mathcal{A}} {D_{s}} {(\partial_{y} u^{x})} {\nu^{x}} {\tau} \left(u^{y}\right)^3 + 4 {\mathcal{A}} {D_{s}} {(\partial_{y} u^{y})} {\nu^{y}} {\tau} \left(u^{y}\right)^3\\
  &\quad + 9 {\mathcal{C}'(T)} {D_{s}} {(\partial_{y} T)} {\Pi_{Q}} {\tau} \left(u^{y}\right)^3 + 10 {\mathcal{B}} {D_{s}} {\nu^{y}} {u^{\tau}} \left(u^{y}\right)^3\\
  &\quad + 10 {D_{s}} {(\partial_{\tau} u^{y})} {I_{31}} {\tau} {u^{\tau}} \left(u^{y}\right)^3 + 10 {\mathcal{B}'(T)} {D_{s}} {(\partial_{\tau} T)} {\nu^{y}} {\tau} {u^{\tau}} \left(u^{y}\right)^3\\
  &\quad + 10 {\mathcal{C}} {D_{s}} {(\partial_{\tau} u^{y})} {\Pi_{Q}} {\tau} {u^{\tau}} \left(u^{y}\right)^3 + 4 {\mathcal{C}} {D_{s}} {(\partial_{\tau} u^{x})} {\pi_{Q}^{xy}} {\tau} {u^{\tau}} \left(u^{y}\right)^3\\
  &\quad + 2 {\mathcal{C}} {D_{s}} {(\partial_{\tau} u^{y})} {\pi_{Q}^{yy}} {\tau} {u^{\tau}} \left(u^{y}\right)^3 + 4 {\mathcal{C}} {D_{s}} {(\partial_{x} u^{x})} {\pi_{Q}^{xy}} {\tau} {u^{x}} \left(u^{y}\right)^3\\
  &\quad - 4 {\mathcal{C}} {D_{s}} {(\partial_{y} u^{y})} {\pi_{Q}^{xy}} {\tau} {u^{x}} \left(u^{y}\right)^3 + 2 {\mathcal{C}} {D_{s}} {(\partial_{x} u^{y})} {\pi_{Q}^{yy}} {\tau} {u^{x}} \left(u^{y}\right)^3\\
  &\quad - 4 {\mathcal{C}} {D_{s}} {(\partial_{y} u^{x})} {\pi_{Q}^{yy}} {\tau} {u^{x}} \left(u^{y}\right)^3 + 4 {\mathcal{C}} {D_{s}} {(\partial_{x} u^{y})} {\pi_{Q}{}^{\eta}{}_{\eta}} {\tau} {u^{x}} \left(u^{y}\right)^3\\
  &\quad - 4 {\mathcal{A}} {D_{s}} {(\partial_{\tau} u^{y})} {\nu^{x}} {\tau} {u^{\tau}} {u^{x}} \left(u^{y}\right)^3 - 5 {\mathcal{B}} {D_{s}} {(\partial_{\tau} u^{y})} {\nu^{x}} {\tau} {u^{\tau}} {u^{x}} \left(u^{y}\right)^3\\
  &\quad - 4 {\mathcal{A}} {D_{s}} {(\partial_{\tau} u^{x})} {\nu^{y}} {\tau} {u^{\tau}} {u^{x}} \left(u^{y}\right)^3 - 5 {\mathcal{B}} {D_{s}} {(\partial_{\tau} u^{x})} {\nu^{y}} {\tau} {u^{\tau}} {u^{x}} \left(u^{y}\right)^3\\
  &\quad + 12 {(\partial_{T} I_{31})} {D_{s}} {(\partial_{y} T)} {\tau} \left(u^{x}\right)^2 \left(u^{y}\right)^3 - 4 {\mathcal{A}} {D_{s}} {(\partial_{x} u^{y})} {\nu^{x}} {\tau} \left(u^{x}\right)^2 \left(u^{y}\right)^3\\
  &\quad + 4 {\mathcal{A}} {D_{s}} {(\partial_{y} u^{x})} {\nu^{x}} {\tau} \left(u^{x}\right)^2 \left(u^{y}\right)^3 - 4 {\mathcal{A}} {D_{s}} {(\partial_{x} u^{x})} {\nu^{y}} {\tau} \left(u^{x}\right)^2 \left(u^{y}\right)^3\\
  &\quad + 6 {\mathcal{A}} {D_{s}} {(\partial_{y} u^{y})} {\nu^{y}} {\tau} \left(u^{x}\right)^2 \left(u^{y}\right)^3\\
  &\quad + 12 {\mathcal{C}'(T)} {D_{s}} {(\partial_{y} T)} {\Pi_{Q}} {\tau} \left(u^{x}\right)^2 \left(u^{y}\right)^3 + 5 {\mathcal{B}} {D_{s}} {\nu^{y}} {u^{\tau}} \left(u^{x}\right)^2 \left(u^{y}\right)^3\\
  &\quad + 5 {D_{s}} {(\partial_{\tau} u^{y})} {I_{31}} {\tau} {u^{\tau}} \left(u^{x}\right)^2 \left(u^{y}\right)^3\\
  &\quad + 5 {\mathcal{B}'(T)} {D_{s}} {(\partial_{\tau} T)} {\nu^{y}} {\tau} {u^{\tau}} \left(u^{x}\right)^2 \left(u^{y}\right)^3\\
  &\quad + 5 {\mathcal{C}} {D_{s}} {(\partial_{\tau} u^{y})} {\Pi_{Q}} {\tau} {u^{\tau}} \left(u^{x}\right)^2 \left(u^{y}\right)^3\\
  &\quad + 2 {\mathcal{C}} {D_{s}} {(\partial_{\tau} u^{x})} {\pi_{Q}^{xy}} {\tau} {u^{\tau}} \left(u^{x}\right)^2 \left(u^{y}\right)^3\\
  &\quad + 2 {\mathcal{C}} {D_{s}} {(\partial_{\tau} u^{y})} {\pi_{Q}^{yy}} {\tau} {u^{\tau}} \left(u^{x}\right)^2 \left(u^{y}\right)^3\\
  &\quad + 2 {\mathcal{C}} {D_{s}} {(\partial_{\tau} u^{y})} {\pi_{Q}{}^{\eta}{}_{\eta}} {\tau} {u^{\tau}} \left(u^{x}\right)^2 \left(u^{y}\right)^3 + 2 {\mathcal{C}} {D_{s}} {(\partial_{x} u^{x})} {\pi_{Q}^{xy}} {\tau} \left(u^{x}\right)^3 \left(u^{y}\right)^3\\
  &\quad - 2 {\mathcal{C}} {D_{s}} {(\partial_{y} u^{y})} {\pi_{Q}^{xy}} {\tau} \left(u^{x}\right)^3 \left(u^{y}\right)^3 + 2 {\mathcal{C}} {D_{s}} {(\partial_{x} u^{y})} {\pi_{Q}^{yy}} {\tau} \left(u^{x}\right)^3 \left(u^{y}\right)^3\\
  &\quad - 2 {\mathcal{C}} {D_{s}} {(\partial_{y} u^{x})} {\pi_{Q}^{yy}} {\tau} \left(u^{x}\right)^3 \left(u^{y}\right)^3 + 4 {\mathcal{C}} {D_{s}} {(\partial_{x} u^{y})} {\pi_{Q}{}^{\eta}{}_{\eta}} {\tau} \left(u^{x}\right)^3 \left(u^{y}\right)^3\\
  &\quad - 2 {\mathcal{A}} {D_{s}} {(\partial_{\tau} u^{y})} {\nu^{x}} {\tau} {u^{\tau}} \left(u^{x}\right)^3 \left(u^{y}\right)^3 - 2 {\mathcal{A}} {D_{s}} {(\partial_{\tau} u^{x})} {\nu^{y}} {\tau} {u^{\tau}} \left(u^{x}\right)^3 \left(u^{y}\right)^3\\
  &\quad + 3 {(\partial_{T} I_{31})} {D_{s}} {(\partial_{y} T)} {\tau} \left(u^{x}\right)^4 \left(u^{y}\right)^3 - 2 {\mathcal{A}} {D_{s}} {(\partial_{x} u^{y})} {\nu^{x}} {\tau} \left(u^{x}\right)^4 \left(u^{y}\right)^3\\
  &\quad - 2 {\mathcal{A}} {D_{s}} {(\partial_{x} u^{x})} {\nu^{y}} {\tau} \left(u^{x}\right)^4 \left(u^{y}\right)^3 + 2 {\mathcal{A}} {D_{s}} {(\partial_{y} u^{y})} {\nu^{y}} {\tau} \left(u^{x}\right)^4 \left(u^{y}\right)^3\\
  &\quad + 3 {\mathcal{C}'(T)} {D_{s}} {(\partial_{y} T)} {\Pi_{Q}} {\tau} \left(u^{x}\right)^4 \left(u^{y}\right)^3 + 15 {D_{s}} {(\partial_{x} u^{x})} {I_{31}} {\tau} \left(u^{y}\right)^4\\
  &\quad + 15 {D_{s}} {(\partial_{y} u^{y})} {I_{31}} {\tau} \left(u^{y}\right)^4 + 15 {\mathcal{B}'(T)} {D_{s}} {(\partial_{x} T)} {\nu^{x}} {\tau} \left(u^{y}\right)^4\\
  &\quad + 15 {\mathcal{B}'(T)} {D_{s}} {(\partial_{y} T)} {\nu^{y}} {\tau} \left(u^{y}\right)^4 + 15 {\mathcal{C}} {D_{s}} {(\partial_{x} u^{x})} {\Pi_{Q}} {\tau} \left(u^{y}\right)^4\\
  &\quad + 15 {\mathcal{C}} {D_{s}} {(\partial_{y} u^{y})} {\Pi_{Q}} {\tau} \left(u^{y}\right)^4 + 2 {\mathcal{C}} {D_{s}} {(\partial_{x} u^{y})} {\pi_{Q}^{xy}} {\tau} \left(u^{y}\right)^4\\
  &\quad + 6 {\mathcal{C}} {D_{s}} {(\partial_{y} u^{x})} {\pi_{Q}^{xy}} {\tau} \left(u^{y}\right)^4 - 2 {\mathcal{C}} {D_{s}} {(\partial_{x} u^{x})} {\pi_{Q}^{yy}} {\tau} \left(u^{y}\right)^4\\
  &\quad + 2 {\mathcal{C}} {D_{s}} {(\partial_{y} u^{y})} {\pi_{Q}^{yy}} {\tau} \left(u^{y}\right)^4 - 6 {\mathcal{C}} {D_{s}} {(\partial_{x} u^{x})} {\pi_{Q}{}^{\eta}{}_{\eta}} {\tau} \left(u^{y}\right)^4\\
  &\quad + 18 {\Pi_{Q}} {T} {\tau} \left(u^{y}\right)^4 + 15 {D_{s}} {I_{31}} {u^{\tau}} \left(u^{y}\right)^4 + 15 {\mathcal{C}} {D_{s}} {\Pi_{Q}} {u^{\tau}} \left(u^{y}\right)^4\\
  &\quad + 6 {\mathcal{C}} {D_{s}} {\pi_{Q}{}^{\eta}{}_{\eta}} {u^{\tau}} \left(u^{y}\right)^4 + 9 {(\partial_{T} I_{31})} {D_{s}} {(\partial_{\tau} T)} {\tau} {u^{\tau}} \left(u^{y}\right)^4\\
  &\quad + 6 {\mathcal{A}} {D_{s}} {(\partial_{\tau} u^{x})} {\nu^{x}} {\tau} {u^{\tau}} \left(u^{y}\right)^4 + 5 {\mathcal{B}} {D_{s}} {(\partial_{\tau} u^{x})} {\nu^{x}} {\tau} {u^{\tau}} \left(u^{y}\right)^4\\
  &\quad + 2 {\mathcal{A}} {D_{s}} {(\partial_{\tau} u^{y})} {\nu^{y}} {\tau} {u^{\tau}} \left(u^{y}\right)^4 + 9 {\mathcal{C}'(T)} {D_{s}} {(\partial_{\tau} T)} {\Pi_{Q}} {\tau} {u^{\tau}} \left(u^{y}\right)^4\\
  &\quad + 9 {(\partial_{T} I_{31})} {D_{s}} {(\partial_{x} T)} {\tau} {u^{x}} \left(u^{y}\right)^4 + 6 {\mathcal{A}} {D_{s}} {(\partial_{x} u^{x})} {\nu^{x}} {\tau} {u^{x}} \left(u^{y}\right)^4\\
  &\quad - 4 {\mathcal{A}} {D_{s}} {(\partial_{y} u^{y})} {\nu^{x}} {\tau} {u^{x}} \left(u^{y}\right)^4 + 2 {\mathcal{A}} {D_{s}} {(\partial_{x} u^{y})} {\nu^{y}} {\tau} {u^{x}} \left(u^{y}\right)^4\\
  &\quad - 4 {\mathcal{A}} {D_{s}} {(\partial_{y} u^{x})} {\nu^{y}} {\tau} {u^{x}} \left(u^{y}\right)^4 + 9 {\mathcal{C}'(T)} {D_{s}} {(\partial_{x} T)} {\Pi_{Q}} {\tau} {u^{x}} \left(u^{y}\right)^4\\
  &\quad + 5 {\mathcal{B}} {D_{s}} {\nu^{x}} {u^{\tau}} {u^{x}} \left(u^{y}\right)^4 + 5 {D_{s}} {(\partial_{\tau} u^{x})} {I_{31}} {\tau} {u^{\tau}} {u^{x}} \left(u^{y}\right)^4\\
  &\quad + 5 {\mathcal{B}'(T)} {D_{s}} {(\partial_{\tau} T)} {\nu^{x}} {\tau} {u^{\tau}} {u^{x}} \left(u^{y}\right)^4 + 5 {\mathcal{C}} {D_{s}} {(\partial_{\tau} u^{x})} {\Pi_{Q}} {\tau} {u^{\tau}} {u^{x}} \left(u^{y}\right)^4\\
  &\quad - 2 {\mathcal{C}} {D_{s}} {(\partial_{\tau} u^{y})} {\pi_{Q}^{xy}} {\tau} {u^{\tau}} {u^{x}} \left(u^{y}\right)^4 - 2 {\mathcal{C}} {D_{s}} {(\partial_{\tau} u^{x})} {\pi_{Q}^{yy}} {\tau} {u^{\tau}} {u^{x}} \left(u^{y}\right)^4\\
  &\quad - 2 {\mathcal{C}} {D_{s}} {(\partial_{\tau} u^{x})} {\pi_{Q}{}^{\eta}{}_{\eta}} {\tau} {u^{\tau}} {u^{x}} \left(u^{y}\right)^4 + 10 {D_{s}} {(\partial_{x} u^{x})} {I_{31}} {\tau} \left(u^{x}\right)^2 \left(u^{y}\right)^4\\
  &\quad + 10 {D_{s}} {(\partial_{y} u^{y})} {I_{31}} {\tau} \left(u^{x}\right)^2 \left(u^{y}\right)^4\\
  &\quad + 10 {\mathcal{B}'(T)} {D_{s}} {(\partial_{x} T)} {\nu^{x}} {\tau} \left(u^{x}\right)^2 \left(u^{y}\right)^4\\
  &\quad + 10 {\mathcal{B}'(T)} {D_{s}} {(\partial_{y} T)} {\nu^{y}} {\tau} \left(u^{x}\right)^2 \left(u^{y}\right)^4\\
  &\quad + 10 {\mathcal{C}} {D_{s}} {(\partial_{x} u^{x})} {\Pi_{Q}} {\tau} \left(u^{x}\right)^2 \left(u^{y}\right)^4 + 10 {\mathcal{C}} {D_{s}} {(\partial_{y} u^{y})} {\Pi_{Q}} {\tau} \left(u^{x}\right)^2 \left(u^{y}\right)^4\\
  &\quad - 2 {\mathcal{C}} {D_{s}} {(\partial_{x} u^{y})} {\pi_{Q}^{xy}} {\tau} \left(u^{x}\right)^2 \left(u^{y}\right)^4 + 2 {\mathcal{C}} {D_{s}} {(\partial_{y} u^{x})} {\pi_{Q}^{xy}} {\tau} \left(u^{x}\right)^2 \left(u^{y}\right)^4\\
  &\quad - 2 {\mathcal{C}} {D_{s}} {(\partial_{x} u^{x})} {\pi_{Q}^{yy}} {\tau} \left(u^{x}\right)^2 \left(u^{y}\right)^4 + 2 {\mathcal{C}} {D_{s}} {(\partial_{y} u^{y})} {\pi_{Q}^{yy}} {\tau} \left(u^{x}\right)^2 \left(u^{y}\right)^4\\
  &\quad - 4 {\mathcal{C}} {D_{s}} {(\partial_{x} u^{x})} {\pi_{Q}{}^{\eta}{}_{\eta}} {\tau} \left(u^{x}\right)^2 \left(u^{y}\right)^4 + 12 {\Pi_{Q}} {T} {\tau} \left(u^{x}\right)^2 \left(u^{y}\right)^4\\
  &\quad + 10 {D_{s}} {I_{31}} {u^{\tau}} \left(u^{x}\right)^2 \left(u^{y}\right)^4 + 10 {\mathcal{C}} {D_{s}} {\Pi_{Q}} {u^{\tau}} \left(u^{x}\right)^2 \left(u^{y}\right)^4\\
  &\quad + 4 {\mathcal{C}} {D_{s}} {\pi_{Q}{}^{\eta}{}_{\eta}} {u^{\tau}} \left(u^{x}\right)^2 \left(u^{y}\right)^4 + 6 {(\partial_{T} I_{31})} {D_{s}} {(\partial_{\tau} T)} {\tau} {u^{\tau}} \left(u^{x}\right)^2 \left(u^{y}\right)^4\\
  &\quad + 2 {\mathcal{A}} {D_{s}} {(\partial_{\tau} u^{x})} {\nu^{x}} {\tau} {u^{\tau}} \left(u^{x}\right)^2 \left(u^{y}\right)^4 + 2 {\mathcal{A}} {D_{s}} {(\partial_{\tau} u^{y})} {\nu^{y}} {\tau} {u^{\tau}} \left(u^{x}\right)^2 \left(u^{y}\right)^4\\
  &\quad + 6 {\mathcal{C}'(T)} {D_{s}} {(\partial_{\tau} T)} {\Pi_{Q}} {\tau} {u^{\tau}} \left(u^{x}\right)^2 \left(u^{y}\right)^4\\
  &\quad + 6 {(\partial_{T} I_{31})} {D_{s}} {(\partial_{x} T)} {\tau} \left(u^{x}\right)^3 \left(u^{y}\right)^4 + 2 {\mathcal{A}} {D_{s}} {(\partial_{x} u^{x})} {\nu^{x}} {\tau} \left(u^{x}\right)^3 \left(u^{y}\right)^4\\
  &\quad - 2 {\mathcal{A}} {D_{s}} {(\partial_{y} u^{y})} {\nu^{x}} {\tau} \left(u^{x}\right)^3 \left(u^{y}\right)^4 + 2 {\mathcal{A}} {D_{s}} {(\partial_{x} u^{y})} {\nu^{y}} {\tau} \left(u^{x}\right)^3 \left(u^{y}\right)^4\\
  &\quad - 2 {\mathcal{A}} {D_{s}} {(\partial_{y} u^{x})} {\nu^{y}} {\tau} \left(u^{x}\right)^3 \left(u^{y}\right)^4 + 6 {\mathcal{C}'(T)} {D_{s}} {(\partial_{x} T)} {\Pi_{Q}} {\tau} \left(u^{x}\right)^3 \left(u^{y}\right)^4\\
  &\quad + 9 {(\partial_{T} I_{31})} {D_{s}} {(\partial_{y} T)} {\tau} \left(u^{y}\right)^5 + 6 {\mathcal{A}} {D_{s}} {(\partial_{y} u^{x})} {\nu^{x}} {\tau} \left(u^{y}\right)^5\\
  &\quad + 2 {\mathcal{A}} {D_{s}} {(\partial_{y} u^{y})} {\nu^{y}} {\tau} \left(u^{y}\right)^5 + 9 {\mathcal{C}'(T)} {D_{s}} {(\partial_{y} T)} {\Pi_{Q}} {\tau} \left(u^{y}\right)^5\\
  &\quad + 5 {\mathcal{B}} {D_{s}} {\nu^{y}} {u^{\tau}} \left(u^{y}\right)^5 + 5 {D_{s}} {(\partial_{\tau} u^{y})} {I_{31}} {\tau} {u^{\tau}} \left(u^{y}\right)^5\\
  &\quad + 5 {\mathcal{B}'(T)} {D_{s}} {(\partial_{\tau} T)} {\nu^{y}} {\tau} {u^{\tau}} \left(u^{y}\right)^5 + 5 {\mathcal{C}} {D_{s}} {(\partial_{\tau} u^{y})} {\Pi_{Q}} {\tau} {u^{\tau}} \left(u^{y}\right)^5\\
  &\quad + 2 {\mathcal{C}} {D_{s}} {(\partial_{\tau} u^{x})} {\pi_{Q}^{xy}} {\tau} {u^{\tau}} \left(u^{y}\right)^5 + 2 {\mathcal{C}} {D_{s}} {(\partial_{x} u^{x})} {\pi_{Q}^{xy}} {\tau} {u^{x}} \left(u^{y}\right)^5\\
  &\quad - 2 {\mathcal{C}} {D_{s}} {(\partial_{y} u^{y})} {\pi_{Q}^{xy}} {\tau} {u^{x}} \left(u^{y}\right)^5 - 2 {\mathcal{C}} {D_{s}} {(\partial_{y} u^{x})} {\pi_{Q}^{yy}} {\tau} {u^{x}} \left(u^{y}\right)^5\\
  &\quad + 2 {\mathcal{C}} {D_{s}} {(\partial_{x} u^{y})} {\pi_{Q}{}^{\eta}{}_{\eta}} {\tau} {u^{x}} \left(u^{y}\right)^5 - 2 {\mathcal{A}} {D_{s}} {(\partial_{\tau} u^{y})} {\nu^{x}} {\tau} {u^{\tau}} {u^{x}} \left(u^{y}\right)^5\\
  &\quad - 2 {\mathcal{A}} {D_{s}} {(\partial_{\tau} u^{x})} {\nu^{y}} {\tau} {u^{\tau}} {u^{x}} \left(u^{y}\right)^5 + 6 {(\partial_{T} I_{31})} {D_{s}} {(\partial_{y} T)} {\tau} \left(u^{x}\right)^2 \left(u^{y}\right)^5\\
  &\quad - 2 {\mathcal{A}} {D_{s}} {(\partial_{x} u^{y})} {\nu^{x}} {\tau} \left(u^{x}\right)^2 \left(u^{y}\right)^5 + 2 {\mathcal{A}} {D_{s}} {(\partial_{y} u^{x})} {\nu^{x}} {\tau} \left(u^{x}\right)^2 \left(u^{y}\right)^5\\
  &\quad - 2 {\mathcal{A}} {D_{s}} {(\partial_{x} u^{x})} {\nu^{y}} {\tau} \left(u^{x}\right)^2 \left(u^{y}\right)^5 + 2 {\mathcal{A}} {D_{s}} {(\partial_{y} u^{y})} {\nu^{y}} {\tau} \left(u^{x}\right)^2 \left(u^{y}\right)^5\\
  &\quad + 6 {\mathcal{C}'(T)} {D_{s}} {(\partial_{y} T)} {\Pi_{Q}} {\tau} \left(u^{x}\right)^2 \left(u^{y}\right)^5 + 5 {D_{s}} {(\partial_{x} u^{x})} {I_{31}} {\tau} \left(u^{y}\right)^6\\
  &\quad + 5 {D_{s}} {(\partial_{y} u^{y})} {I_{31}} {\tau} \left(u^{y}\right)^6 + 5 {\mathcal{B}'(T)} {D_{s}} {(\partial_{x} T)} {\nu^{x}} {\tau} \left(u^{y}\right)^6\\
  &\quad + 5 {\mathcal{B}'(T)} {D_{s}} {(\partial_{y} T)} {\nu^{y}} {\tau} \left(u^{y}\right)^6 + 5 {\mathcal{C}} {D_{s}} {(\partial_{x} u^{x})} {\Pi_{Q}} {\tau} \left(u^{y}\right)^6\\
  &\quad + 5 {\mathcal{C}} {D_{s}} {(\partial_{y} u^{y})} {\Pi_{Q}} {\tau} \left(u^{y}\right)^6 + 2 {\mathcal{C}} {D_{s}} {(\partial_{y} u^{x})} {\pi_{Q}^{xy}} {\tau} \left(u^{y}\right)^6\\
  &\quad - 2 {\mathcal{C}} {D_{s}} {(\partial_{x} u^{x})} {\pi_{Q}{}^{\eta}{}_{\eta}} {\tau} \left(u^{y}\right)^6 + 6 {\Pi_{Q}} {T} {\tau} \left(u^{y}\right)^6 + 5 {D_{s}} {I_{31}} {u^{\tau}} \left(u^{y}\right)^6\\
  &\quad + 5 {\mathcal{C}} {D_{s}} {\Pi_{Q}} {u^{\tau}} \left(u^{y}\right)^6 + 2 {\mathcal{C}} {D_{s}} {\pi_{Q}{}^{\eta}{}_{\eta}} {u^{\tau}} \left(u^{y}\right)^6\\
  &\quad + 3 {(\partial_{T} I_{31})} {D_{s}} {(\partial_{\tau} T)} {\tau} {u^{\tau}} \left(u^{y}\right)^6 + 2 {\mathcal{A}} {D_{s}} {(\partial_{\tau} u^{x})} {\nu^{x}} {\tau} {u^{\tau}} \left(u^{y}\right)^6\\
  &\quad + 3 {\mathcal{C}'(T)} {D_{s}} {(\partial_{\tau} T)} {\Pi_{Q}} {\tau} {u^{\tau}} \left(u^{y}\right)^6 + 3 {(\partial_{T} I_{31})} {D_{s}} {(\partial_{x} T)} {\tau} {u^{x}} \left(u^{y}\right)^6\\
  &\quad + 2 {\mathcal{A}} {D_{s}} {(\partial_{x} u^{x})} {\nu^{x}} {\tau} {u^{x}} \left(u^{y}\right)^6 - 2 {\mathcal{A}} {D_{s}} {(\partial_{y} u^{y})} {\nu^{x}} {\tau} {u^{x}} \left(u^{y}\right)^6\\
  &\quad - 2 {\mathcal{A}} {D_{s}} {(\partial_{y} u^{x})} {\nu^{y}} {\tau} {u^{x}} \left(u^{y}\right)^6 + 3 {\mathcal{C}'(T)} {D_{s}} {(\partial_{x} T)} {\Pi_{Q}} {\tau} {u^{x}} \left(u^{y}\right)^6\\
  &\quad + 3 {(\partial_{T} I_{31})} {D_{s}} {(\partial_{y} T)} {\tau} \left(u^{y}\right)^7 + 2 {\mathcal{A}} {D_{s}} {(\partial_{y} u^{x})} {\nu^{x}} {\tau} \left(u^{y}\right)^7\\
  &\quad + 3 {\mathcal{C}'(T)} {D_{s}} {(\partial_{y} T)} {\Pi_{Q}} {\tau} \left(u^{y}\right)^7\Bigr)
\end{align*}
\subsubsection{S5, two-way}

\texttt{two\_\allowbreak{}d\_\allowbreak{}two\_\allowbreak{}way\_\allowbreak{}matrix} in \texttt{src/Matrix/2d\_\allowbreak{}two\_\allowbreak{}way.jl}, fields:

\[
\phi = ({T}, {u^{x}}, {u^{y}}, {\pi^{yy}}, {\pi^{\eta}{}_{\eta}}, {\pi^{xy}}, {\Pi}, {\alpha}, {\nu^{x}}, {\nu^{y}}, {\pi_{Q}^{yy}}, {\pi_{Q}{}^{\eta}{}_{\eta}}, {\pi_{Q}^{xy}}, {\Pi_{Q}})
\]
\begin{align*}
(A_t)_{1,1} &= - {n} - {(\partial_{T} n)} {T} + {P''(T)} {T} + {(\partial_{T} n)} {h} \left(u^{\tau}\right)^2 + {h'(T)} {n} \left(u^{\tau}\right)^2 + 2 {h'(T)} {\nu^{x}} {u^{x}}\\
  &\quad + {P'(T)} \left(u^{x}\right)^2 + {P''(T)} {T} \left(u^{x}\right)^2 + 2 {h'(T)} {\nu^{y}} {u^{y}} + {P'(T)} \left(u^{y}\right)^2 + {P''(T)} {T} \left(u^{y}\right)^2 \\
(A_t)_{1,2} &= \frac{2}{\left(u^{y}\right)^2+1}\Bigl(- {\pi^{\eta}{}_{\eta}} {u^{x}} - {\pi^{yy}} {u^{x}} + {P'(T)} {T} {u^{x}} + {\pi^{xy}} {u^{y}} + {P'(T)} {T} {u^{x}} \left(u^{y}\right)^2\\
  &\quad + {\Pi} {u^{x}} \left(\left(u^{y}\right)^2+1\right) + {h} \left(\left(u^{y}\right)^2+1\right) ({n} {u^{x}}+{\nu^{x}})\Bigr) \\
(A_t)_{1,3} &= \frac{2}{\left(\left(u^{y}\right)^2+1\right)^2}\Bigl({h} \left(\left(u^{y}\right)^2+1\right)^2 ({n} {u^{y}}+{\nu^{y}}) + {\pi^{xy}} \left({u^{x}}-{u^{x}} \left(u^{y}\right)^2\right)\\
  &\quad + {u^{y}}\Bigl({\pi^{yy}} + {P'(T)} {T} + {\pi^{\eta}{}_{\eta}} \left(u^{x}\right)^2 + {\pi^{yy}} \left(u^{x}\right)^2 + 2 {P'(T)} {T} \left(u^{y}\right)^2\\
  &\quad + {P'(T)} {T} \left(u^{y}\right)^4 + {\Pi} \left(\left(u^{y}\right)^2+1\right)^2\Bigr)\Bigr) \\
(A_t)_{1,4} &= \frac{\left(u^{y}\right)^2-\left(u^{x}\right)^2}{\left(u^{y}\right)^2+1} \\
(A_t)_{1,5} &= -\frac{\left(u^{x}\right)^2}{\left(u^{y}\right)^2+1} \\
(A_t)_{1,6} &= \frac{2 {u^{x}} {u^{y}}}{\left(u^{y}\right)^2+1} \\
(A_t)_{1,7} &= \left(u^{x}\right)^2+\left(u^{y}\right)^2 \\
(A_t)_{1,8} &= {(\partial_{\alpha} n)} \left({h} \left(u^{\tau}\right)^2-{T}\right) \\
(A_t)_{1,9} &= 2 {h} {u^{x}} \\
(A_t)_{1,10} &= 2 {h} {u^{y}} \\
(A_t)_{2,1} &= \frac{1}{{u^{\tau}}}\Bigl(\left(u^{\tau}\right)^2 {u^{x}} ({(\partial_{T} n)} {h}+{P''(T)} {T}+{P'(T)})\\
  &\quad + {h'(T)} \left({n} \left(u^{\tau}\right)^2 {u^{x}}+{\nu^{x}} \left(2 \left(u^{x}\right)^2+\left(u^{y}\right)^2+1\right)+{\nu^{y}} {u^{x}} {u^{y}}\right)\Bigr) \\
(A_t)_{2,2} &= \frac{1}{\left(u^{\tau}\right)^3 \left(\left(u^{y}\right)^2+1\right)}\Bigl(- {\pi^{\eta}{}_{\eta}} - {\pi^{yy}} + {P'(T)} {T} - 3 {\pi^{\eta}{}_{\eta}} \left(u^{x}\right)^2 - 3 {\pi^{yy}} \left(u^{x}\right)^2\\
  &\quad + 3 {P'(T)} {T} \left(u^{x}\right)^2 - 2 {\pi^{\eta}{}_{\eta}} \left(u^{x}\right)^4 - 2 {\pi^{yy}} \left(u^{x}\right)^4 + 2 {P'(T)} {T} \left(u^{x}\right)^4\\
  &\quad + 3 {\pi^{xy}} {u^{x}} {u^{y}} + 2 {\pi^{xy}} \left(u^{x}\right)^3 {u^{y}} - 2 {\pi^{\eta}{}_{\eta}} \left(u^{y}\right)^2 - {\pi^{yy}} \left(u^{y}\right)^2\\
  &\quad + 3 {P'(T)} {T} \left(u^{y}\right)^2 - 3 {\pi^{\eta}{}_{\eta}} \left(u^{x}\right)^2 \left(u^{y}\right)^2 - 3 {\pi^{yy}} \left(u^{x}\right)^2 \left(u^{y}\right)^2\\
  &\quad + 6 {P'(T)} {T} \left(u^{x}\right)^2 \left(u^{y}\right)^2 + 2 {P'(T)} {T} \left(u^{x}\right)^4 \left(u^{y}\right)^2 + 3 {\pi^{xy}} {u^{x}} \left(u^{y}\right)^3\\
  &\quad - {\pi^{\eta}{}_{\eta}} \left(u^{y}\right)^4 + 3 {P'(T)} {T} \left(u^{y}\right)^4 + 3 {P'(T)} {T} \left(u^{x}\right)^2 \left(u^{y}\right)^4\\
  &\quad + {P'(T)} {T} \left(u^{y}\right)^6\\
  &\quad + {\Pi} \left(\left(u^{y}\right)^2+1\right) \left(2 \left(u^{x}\right)^4+3 \left(u^{x}\right)^2 \left(\left(u^{y}\right)^2+1\right)+\left(\left(u^{y}\right)^2+1\right)^2\right)\\
  &\quad + {h} \left(\left(u^{y}\right)^2+1\right)\Bigl({\nu^{y}} {u^{y}} \left(\left(u^{y}\right)^2+1\right)\\
  &\quad + {\nu^{x}} {u^{x}} \left(2 \left(u^{x}\right)^2+3 \left(u^{y}\right)^2+3\right)\\
  &\quad + {n}\Bigl(2 \left(u^{x}\right)^4+3 \left(u^{x}\right)^2 \left(\left(u^{y}\right)^2+1\right)+\left(\left(u^{y}\right)^2+1\right)^2\Bigr)\Bigr)\Bigr) \\
(A_t)_{2,3} &= \frac{1}{\left(u^{\tau}\right)^3 \left(\left(u^{y}\right)^2+1\right)^2}\Bigl({h} \left(\left(u^{y}\right)^2+1\right)^2\Bigl({n} \left(u^{\tau}\right)^2 {u^{x}} {u^{y}}+{\nu^{x}} {u^{y}} \left(\left(u^{y}\right)^2+1\right)+{\nu^{y}} \left(\left(u^{x}\right)^3+{u^{x}}\right)\Bigr)\\
  &\quad + {\pi^{xy}}\Bigl(- 2 \left(u^{x}\right)^4 \left(\left(u^{y}\right)^2-1\right) + \left(\left(u^{y}\right)^2+1\right)^2\\
  &\quad - 3 \left(u^{x}\right)^2 \left(\left(u^{y}\right)^4-1\right)\Bigr)\\
  &\quad + {u^{x}} {u^{y}}\Bigl(3 {\pi^{yy}} + {P'(T)} {T} + 5 {\pi^{yy}} \left(u^{x}\right)^2 + {P'(T)} {T} \left(u^{x}\right)^2\\
  &\quad + 2 {\pi^{yy}} \left(u^{x}\right)^4 + 3 {\pi^{yy}} \left(u^{y}\right)^2 + 3 {P'(T)} {T} \left(u^{y}\right)^2\\
  &\quad + 3 {\pi^{yy}} \left(u^{x}\right)^2 \left(u^{y}\right)^2 + 2 {P'(T)} {T} \left(u^{x}\right)^2 \left(u^{y}\right)^2\\
  &\quad + 3 {P'(T)} {T} \left(u^{y}\right)^4 + {P'(T)} {T} \left(u^{x}\right)^2 \left(u^{y}\right)^4\\
  &\quad + {P'(T)} {T} \left(u^{y}\right)^6 + {\Pi} \left(u^{\tau}\right)^2 \left(\left(u^{y}\right)^2+1\right)^2\\
  &\quad + {\pi^{\eta}{}_{\eta}} \left(2 \left(u^{x}\right)^4+3 \left(u^{x}\right)^2 \left(\left(u^{y}\right)^2+1\right)+\left(\left(u^{y}\right)^2+1\right)^2\right)\Bigr)\Bigr) \\
(A_t)_{2,4} &= -\frac{{u^{x}} \left(\left(u^{x}\right)^2+1\right)}{{u^{\tau}} \left(\left(u^{y}\right)^2+1\right)} \\
(A_t)_{2,5} &= -\frac{{u^{\tau}} {u^{x}}}{\left(u^{y}\right)^2+1} \\
(A_t)_{2,6} &= \frac{{u^{y}} \left(2 \left(u^{x}\right)^2+\left(u^{y}\right)^2+1\right)}{{u^{\tau}} \left(\left(u^{y}\right)^2+1\right)} \\
(A_t)_{2,7} &= {u^{\tau}} {u^{x}} \\
(A_t)_{2,8} &= {(\partial_{\alpha} n)} {h} {u^{\tau}} {u^{x}} \\
(A_t)_{2,9} &= \frac{{h} \left(2 \left(u^{x}\right)^2+\left(u^{y}\right)^2+1\right)}{{u^{\tau}}} \\
(A_t)_{2,10} &= \frac{{h} {u^{x}} {u^{y}}}{{u^{\tau}}} \\
(A_t)_{3,1} &= \frac{1}{{u^{\tau}}}\Bigl(\left(u^{\tau}\right)^2 {u^{y}} ({(\partial_{T} n)} {h}+{P''(T)} {T}+{P'(T)})\\
  &\quad + {h'(T)} \left({n} \left(u^{\tau}\right)^2 {u^{y}}+{\nu^{x}} {u^{x}} {u^{y}}+{\nu^{y}} \left(\left(u^{x}\right)^2+2 \left(u^{y}\right)^2+1\right)\right)\Bigr) \\
(A_t)_{3,2} &= \frac{1}{\left(u^{\tau}\right)^3}\Bigl({u^{x}} {u^{y}} \left({P'(T)} {T} \left(u^{\tau}\right)^2+{\Pi} \left(u^{\tau}\right)^2-{\pi^{yy}}\right) + {\pi^{xy}} \left(\left(u^{y}\right)^2+1\right)\\
  &\quad + {h} \left({n} \left(u^{\tau}\right)^2 {u^{x}} {u^{y}}+{\nu^{x}} {u^{y}} \left(\left(u^{y}\right)^2+1\right)+{\nu^{y}} \left(\left(u^{x}\right)^3+{u^{x}}\right)\right)\Bigr) \\
(A_t)_{3,3} &= \frac{1}{\left(u^{\tau}\right)^3}\Bigl({\pi^{yy}} + {P'(T)} {T} + {\pi^{yy}} \left(u^{x}\right)^2 + 2 {P'(T)} {T} \left(u^{x}\right)^2 + {P'(T)} {T} \left(u^{x}\right)^4 - {\pi^{xy}} {u^{x}} {u^{y}}\\
  &\quad + 3 {P'(T)} {T} \left(u^{y}\right)^2 + 3 {P'(T)} {T} \left(u^{x}\right)^2 \left(u^{y}\right)^2 + 2 {P'(T)} {T} \left(u^{y}\right)^4\\
  &\quad + {\Pi} \left(\left(u^{x}\right)^4+\left(u^{x}\right)^2 \left(3 \left(u^{y}\right)^2+2\right)+2 \left(u^{y}\right)^4+3 \left(u^{y}\right)^2+1\right)\\
  &\quad + {h}\Bigl({\nu^{x}} \left(\left(u^{x}\right)^3+{u^{x}}\right) + {\nu^{y}} {u^{y}} \left(3 \left(u^{x}\right)^2+2 \left(u^{y}\right)^2+3\right)\\
  &\quad + {n} \left(\left(u^{x}\right)^4+\left(u^{x}\right)^2 \left(3 \left(u^{y}\right)^2+2\right)+2 \left(u^{y}\right)^4+3 \left(u^{y}\right)^2+1\right)\Bigr)\Bigr) \\
(A_t)_{3,4} &= \frac{{u^{y}}}{{u^{\tau}}} \\
(A_t)_{3,6} &= \frac{{u^{x}}}{{u^{\tau}}} \\
(A_t)_{3,7} &= {u^{\tau}} {u^{y}} \\
(A_t)_{3,8} &= {(\partial_{\alpha} n)} {h} {u^{\tau}} {u^{y}} \\
(A_t)_{3,9} &= \frac{{h} {u^{x}} {u^{y}}}{{u^{\tau}}} \\
(A_t)_{3,10} &= \frac{{h} \left(\left(u^{x}\right)^2+2 \left(u^{y}\right)^2+1\right)}{{u^{\tau}}} \\
(A_t)_{4,2} &= \frac{3 {\delta_{\pi\pi}} {\pi^{yy}} {u^{x}}-6 {\pi^{xy}} {\tau_{\pi}} {u^{y}} \left(\left(u^{y}\right)^2+1\right)-2 {u^{x}} \left(-3 {\pi^{yy}} {\tau_{\pi}} \left(u^{y}\right)^2+\left(u^{y}\right)^2 {\eta}+{\eta}\right)}{3 {u^{\tau}}} \\
(A_t)_{4,3} &= \frac{{u^{y}} \left(3 {\delta_{\pi\pi}} {\pi^{yy}}+6 {\pi^{xy}} {\tau_{\pi}} {u^{x}} {u^{y}}-6 {\pi^{yy}} {\tau_{\pi}} \left(\left(u^{x}\right)^2+1\right)+6 \left(u^{x}\right)^2 {\eta}+4 \left(u^{y}\right)^2 {\eta}+4 {\eta}\right)}{3 {u^{\tau}}} \\
(A_t)_{4,4} &= {\tau_{\pi}} {u^{\tau}} \\
(A_t)_{5,2} &= \frac{{u^{x}} (3 {\delta_{\pi\pi}} {\pi^{\eta}{}_{\eta}}-2 {\eta})}{3 {\tau}^2 {u^{\tau}}} \\
(A_t)_{5,3} &= \frac{{u^{y}} (3 {\delta_{\pi\pi}} {\pi^{\eta}{}_{\eta}}-2 {\eta})}{3 {\tau}^2 {u^{\tau}}} \\
(A_t)_{5,5} &= \frac{{\tau_{\pi}} {u^{\tau}}}{{\tau}^2} \\
(A_t)_{6,2} &= \frac{1}{3 {u^{\tau}}}\Bigl(3 {\delta_{\pi\pi}} {\pi^{xy}} {u^{x}} - 3 {\pi^{xy}} {\tau_{\pi}} {u^{x}} \left(2 \left(u^{y}\right)^2+1\right)\\
  &\quad + {u^{y}} \left(3 {\pi^{\eta}{}_{\eta}} {\tau_{\pi}} \left(u^{\tau}\right)^2+3 {\pi^{yy}} \left(2 {\tau_{\pi}} \left(u^{x}\right)^2+{\tau_{\pi}}\right)+{\eta} \left(\left(u^{x}\right)^2+3 \left(u^{y}\right)^2+3\right)\right)\Bigr) \\
(A_t)_{6,3} &= \frac{1}{3 {u^{\tau}} \left(\left(u^{y}\right)^2+1\right)}\Bigl(- 3 {\pi^{xy}} {\tau_{\pi}} {u^{y}} - 3 {\pi^{\eta}{}_{\eta}} {\tau_{\pi}} {u^{x}} \left(u^{y}\right)^2 - 3 {\pi^{\eta}{}_{\eta}} {\tau_{\pi}} \left(u^{x}\right)^3 \left(u^{y}\right)^2 - 3 {\pi^{xy}} {\tau_{\pi}} \left(u^{y}\right)^3\\
  &\quad + 6 {\pi^{xy}} {\tau_{\pi}} \left(u^{x}\right)^2 \left(u^{y}\right)^3 - 3 {\pi^{\eta}{}_{\eta}} {\tau_{\pi}} {u^{x}} \left(u^{y}\right)^4 + 3 {\delta_{\pi\pi}} {\pi^{xy}} {u^{y}} \left(\left(u^{y}\right)^2+1\right)\\
  &\quad - 3 {\pi^{yy}} {\tau_{\pi}} {u^{x}} \left(\left(u^{x}\right)^2+1\right) \left(2 \left(u^{y}\right)^2+1\right) + 3 {u^{x}} {\eta} + 3 \left(u^{x}\right)^3 {\eta}\\
  &\quad + 4 {u^{x}} \left(u^{y}\right)^2 {\eta} + 3 \left(u^{x}\right)^3 \left(u^{y}\right)^2 {\eta} + {u^{x}} \left(u^{y}\right)^4 {\eta}\Bigr) \\
(A_t)_{6,6} &= {\tau_{\pi}} {u^{\tau}} \\
(A_t)_{7,2} &= \frac{{u^{x}} {\zeta}}{{u^{\tau}}} \\
(A_t)_{7,3} &= \frac{{u^{y}} {\zeta}}{{u^{\tau}}} \\
(A_t)_{7,7} &= {\tau_{\Pi}} {u^{\tau}} \\
(A_t)_{8,1} &= {(\partial_{T} n)} {u^{\tau}} \\
(A_t)_{8,2} &= \frac{{n} \left(u^{\tau}\right)^2 {u^{x}}+{\nu^{x}} \left(\left(u^{y}\right)^2+1\right)-{\nu^{y}} {u^{x}} {u^{y}}}{\left(u^{\tau}\right)^3} \\
(A_t)_{8,3} &= \frac{{n} \left(u^{\tau}\right)^2 {u^{y}}-{\nu^{x}} {u^{x}} {u^{y}}+{\nu^{y}} \left(\left(u^{x}\right)^2+1\right)}{\left(u^{\tau}\right)^3} \\
(A_t)_{8,8} &= {(\partial_{\alpha} n)} {u^{\tau}} \\
(A_t)_{8,9} &= \frac{{u^{x}}}{{u^{\tau}}} \\
(A_t)_{8,10} &= \frac{{u^{y}}}{{u^{\tau}}} \\
(A_t)_{9,1} &= \frac{{D_{s}} {u^{\tau}} ({h'(T)} {\nu^{x}}+{u^{x}} ({(\partial_{T} n)} {T}+{n}))}{{T}} \\
(A_t)_{9,2} &= \frac{{D_{s}} {h} \left({n} \left(u^{\tau}\right)^2+{\nu^{x}} \left({u^{x}}-{u^{x}} \left(u^{y}\right)^2\right)+{\nu^{y}} \left(\left(u^{x}\right)^2+1\right) {u^{y}}\right)}{{T} {u^{\tau}}} \\
(A_t)_{9,3} &= -\frac{{D_{s}} {h} \left(\left(u^{x}\right)^2+1\right) ({\nu^{y}} {u^{x}}-{\nu^{x}} {u^{y}})}{{T} {u^{\tau}}} \\
(A_t)_{9,8} &= {(\partial_{\alpha} n)} {D_{s}} {u^{\tau}} {u^{x}} \\
(A_t)_{9,9} &= \frac{{D_{s}} {h} {u^{\tau}}}{{T}} \\
(A_t)_{10,1} &= \frac{{D_{s}} {u^{\tau}} ({h'(T)} {\nu^{y}}+{u^{y}} ({(\partial_{T} n)} {T}+{n}))}{{T}} \\
(A_t)_{10,2} &= \frac{{D_{s}} {h} \left(\left(u^{y}\right)^2+1\right) ({\nu^{y}} {u^{x}}-{\nu^{x}} {u^{y}})}{{T} {u^{\tau}}} \\
(A_t)_{10,3} &= \frac{{D_{s}} {h} \left({n} \left(u^{\tau}\right)^2+{\nu^{x}} {u^{x}} \left(\left(u^{y}\right)^2+1\right)-{\nu^{y}} \left(\left(u^{x}\right)^2-1\right) {u^{y}}\right)}{{T} {u^{\tau}}} \\
(A_t)_{10,8} &= {(\partial_{\alpha} n)} {D_{s}} {u^{\tau}} {u^{y}} \\
(A_t)_{10,10} &= \frac{{D_{s}} {h} {u^{\tau}}}{{T}} \\
(A_t)_{11,1} &= \frac{1}{6 {T} {u^{\tau}}}{D_{s}}\Bigl(3 {\mathcal{C}'(T)} {\pi_{Q}^{yy}} \left(u^{\tau}\right)^2 - 2 {\mathcal{B}'(T)} {\nu^{x}} {u^{x}} \left(\left(u^{y}\right)^2+1\right)\\
  &\quad + 2 {\mathcal{B}'(T)} {\nu^{y}} {u^{y}} \left(3 \left(u^{x}\right)^2+2 \left(u^{y}\right)^2+2\right)\Bigr) \\
(A_t)_{11,2} &= -\frac{1}{6 {T} \left(u^{\tau}\right)^3}{D_{s}}\Bigl(2 {\mathcal{A}} \left(u^{\tau}\right)^2 \left(\left(u^{y}\right)^2+1\right) \left({\nu^{x}} \left(u^{y}\right)^2+{\nu^{x}}-{\nu^{y}} {u^{x}} {u^{y}}\right)\\
  &\quad + 2 {\mathcal{B}} \left(\left(3 \left(u^{x}\right)^2+4\right) \left(u^{y}\right)^2+3 \left(u^{y}\right)^4+1\right) \left({\nu^{x}} \left(u^{y}\right)^2+{\nu^{x}}-{\nu^{y}} {u^{x}} {u^{y}}\right)\\
  &\quad + \left(u^{\tau}\right)^2\Bigl(2 {I_{31}} {u^{x}} \left(\left(u^{y}\right)^2+1\right)\\
  &\quad + {\mathcal{C}}\Bigl(2 {\Pi_{Q}} {u^{x}} \left(\left(u^{y}\right)^2+1\right)\\
  &\quad - 2 \left(\left(u^{y}\right)^2+1\right) ({\pi_{Q}{}^{\eta}{}_{\eta}} {u^{x}}-4 {\pi_{Q}^{xy}} {u^{y}})\\
  &\quad - {\pi_{Q}^{yy}} {u^{x}} \left(8 \left(u^{y}\right)^2+5\right)\Bigr)\Bigr)\Bigr) \\
(A_t)_{11,3} &= \frac{1}{6 {T} \left(u^{\tau}\right)^3}{D_{s}}\Bigl(2 {\mathcal{B}} {\nu^{x}} {u^{x}} {u^{y}} \left(\left(3 \left(u^{x}\right)^2+4\right) \left(u^{y}\right)^2+3 \left(u^{y}\right)^4+1\right)\\
  &\quad - 2{\mathcal{A}} \left(u^{\tau}\right)^2\Bigl({\nu^{y}} \left(u^{x}\right)^2 \left(\left(u^{y}\right)^2-2\right) - 2 {\nu^{y}} \left(\left(u^{y}\right)^2+1\right)\\
  &\quad - {\nu^{x}} {u^{x}} {u^{y}} \left(\left(u^{y}\right)^2+1\right)\Bigr)\\
  &\quad - 2{\mathcal{B}} {\nu^{y}}\Bigl(3 \left(u^{x}\right)^4 \left(\left(u^{y}\right)^2-1\right) - 2 \left(\left(u^{y}\right)^2+1\right)\\
  &\quad + \left(u^{x}\right)^2 \left(3 \left(u^{y}\right)^4+\left(u^{y}\right)^2-5\right)\Bigr)\\
  &\quad + \left(u^{\tau}\right)^2\Bigl(2 {I_{31}} {u^{y}} \left(3 \left(u^{x}\right)^2+2 \left(u^{y}\right)^2+2\right) + 4 {\mathcal{C}} {\pi_{Q}^{xy}} \left(2 {u^{x}} \left(u^{y}\right)^2+{u^{x}}\right)\\
  &\quad + {\mathcal{C}} {u^{y}}\Bigl({\pi_{Q}^{yy}} - 8 {\pi_{Q}^{yy}} \left(u^{x}\right)^2 - 2 {\pi_{Q}{}^{\eta}{}_{\eta}} \left(u^{x}\right)^2\\
  &\quad + {\Pi_{Q}} \left(6 \left(u^{x}\right)^2+4 \left(u^{y}\right)^2+4\right)\Bigr)\Bigr)\Bigr) \\
(A_t)_{11,9} &= -\frac{{\mathcal{B}} {D_{s}} {u^{x}} \left(\left(u^{y}\right)^2+1\right)}{3 {T} {u^{\tau}}} \\
(A_t)_{11,10} &= \frac{{\mathcal{B}} {D_{s}} {u^{y}} \left(3 \left(u^{x}\right)^2+2 \left(u^{y}\right)^2+2\right)}{3 {T} {u^{\tau}}} \\
(A_t)_{11,11} &= \frac{{\mathcal{C}} {D_{s}} {u^{\tau}}}{2 {T}} \\
(A_t)_{12,1} &= \frac{{D_{s}} \left(3 {\mathcal{C}'(T)} {\pi_{Q}{}^{\eta}{}_{\eta}} \left(u^{\tau}\right)^2-2 {\mathcal{B}'(T)} ({\nu^{x}} {u^{x}}+{\nu^{y}} {u^{y}})\right)}{6 {T} {u^{\tau}}} \\
(A_t)_{12,2} &= -\frac{1}{6 {T} \left(u^{\tau}\right)^3}{D_{s}}\Bigl(2 {\mathcal{B}} \left({\nu^{x}} \left(u^{y}\right)^2+{\nu^{x}}-{\nu^{y}} {u^{x}} {u^{y}}\right) + 2 {\mathcal{A}} \left(u^{\tau}\right)^2 \left({\nu^{x}} \left(u^{y}\right)^2+{\nu^{x}}-{\nu^{y}} {u^{x}} {u^{y}}\right)\\
  &\quad + \left(u^{\tau}\right)^2 ({\mathcal{C}} (2 {\Pi_{Q}} {u^{x}}+2 {\pi_{Q}^{xy}} {u^{y}}-2 {\pi_{Q}^{yy}} {u^{x}}-5 {\pi_{Q}{}^{\eta}{}_{\eta}} {u^{x}})+2 {I_{31}} {u^{x}})\Bigr) \\
(A_t)_{12,3} &= -\frac{1}{6 {T} \left(u^{\tau}\right)^3 \left(\left(u^{y}\right)^2+1\right)}{D_{s}}\Bigl(2 {\mathcal{B}} \left(\left(u^{y}\right)^2+1\right) \left(-{\nu^{x}} {u^{x}} {u^{y}}+{\nu^{y}} \left(u^{x}\right)^2+{\nu^{y}}\right)\\
  &\quad + 2 {\mathcal{A}} \left(u^{\tau}\right)^2 \left(\left(u^{y}\right)^2+1\right) \left(-{\nu^{x}} {u^{x}} {u^{y}}+{\nu^{y}} \left(u^{x}\right)^2+{\nu^{y}}\right)\\
  &\quad + \left(u^{\tau}\right)^2\Bigl(- 2 {\mathcal{C}} {\pi_{Q}^{xy}} {u^{x}} \left(\left(u^{y}\right)^2-1\right) + 2 {I_{31}} {u^{y}} \left(\left(u^{y}\right)^2+1\right)\\
  &\quad + {\mathcal{C}} {u^{y}}\Bigl(2 {\pi_{Q}^{yy}} \left(\left(u^{x}\right)^2+1\right)\\
  &\quad + {\pi_{Q}{}^{\eta}{}_{\eta}} \left(2 \left(u^{x}\right)^2-3 \left(u^{y}\right)^2-3\right) + 2 {\Pi_{Q}} \left(\left(u^{y}\right)^2+1\right)\Bigr)\Bigr)\Bigr) \\
(A_t)_{12,9} &= -\frac{{\mathcal{B}} {D_{s}} {u^{x}}}{3 {T} {u^{\tau}}} \\
(A_t)_{12,10} &= -\frac{{\mathcal{B}} {D_{s}} {u^{y}}}{3 {T} {u^{\tau}}} \\
(A_t)_{12,12} &= \frac{{\mathcal{C}} {D_{s}} {u^{\tau}}}{2 {T}} \\
(A_t)_{13,1} &= \frac{1}{6 {T} {u^{\tau}}}{D_{s}}\Bigl(3 {\mathcal{C}'(T)} {\pi_{Q}^{xy}} \left(u^{\tau}\right)^2 + {\mathcal{B}'(T)} {\nu^{y}} {u^{x}} \left(3 \left(u^{x}\right)^2+\left(u^{y}\right)^2+3\right)\\
  &\quad + {\mathcal{B}'(T)} {\nu^{x}} {u^{y}} \left(\left(u^{x}\right)^2+3 \left(u^{y}\right)^2+3\right)\Bigr) \\
(A_t)_{13,2} &= \frac{1}{6 {T} \left(u^{\tau}\right)^3}{D_{s}}\Bigl({\mathcal{A}} \left(u^{\tau}\right)^2\Bigl(3 {\nu^{y}} \left(\left(u^{y}\right)^2+1\right) - 2 {\nu^{x}} {u^{x}} {u^{y}} \left(\left(u^{y}\right)^2+1\right)\\
  &\quad + {\nu^{y}} \left(u^{x}\right)^2 \left(2 \left(u^{y}\right)^2+3\right)\Bigr)\\
  &\quad + \left(u^{\tau}\right)^2\Bigl({\mathcal{C}} {\pi_{Q}^{xy}} {u^{x}} \left(3-8 \left(u^{y}\right)^2\right) + {I_{31}} {u^{y}} \left(\left(u^{x}\right)^2+3 \left(u^{y}\right)^2+3\right)\\
  &\quad + {\mathcal{C}} {u^{y}}\Bigl({\pi_{Q}^{yy}} \left(8 \left(u^{x}\right)^2+6\right) + {\Pi_{Q}} \left(\left(u^{x}\right)^2+3 \left(u^{y}\right)^2+3\right)\\
  &\quad + {\pi_{Q}{}^{\eta}{}_{\eta}} \left(5 \left(u^{x}\right)^2+3 \left(u^{y}\right)^2+3\right)\Bigr)\Bigr)\\
  &\quad + {\mathcal{B}}\Bigl(- 2 {\nu^{x}} {u^{x}} {u^{y}} \left(\left(u^{y}\right)^2+1\right) \left(3 \left(u^{x}\right)^2+3 \left(u^{y}\right)^2+4\right)\\
  &\quad + {\nu^{y}}\Bigl(3 \left(\left(u^{y}\right)^2+1\right)^2 + \left(u^{x}\right)^4 \left(6 \left(u^{y}\right)^2+3\right)\\
  &\quad + 2 \left(u^{x}\right)^2 \left(3 \left(u^{y}\right)^4+7 \left(u^{y}\right)^2+3\right)\Bigr)\Bigr)\Bigr) \\
(A_t)_{13,3} &= \frac{1}{6 {T} \left(u^{\tau}\right)^3 \left(\left(u^{y}\right)^2+1\right)}{D_{s}}\Bigl({\mathcal{A}} \left(u^{\tau}\right)^2 \left(\left(u^{y}\right)^2+1\right)\Bigl(- 2 {\nu^{y}} {u^{x}} \left(\left(u^{x}\right)^2+1\right) {u^{y}} + 3 {\nu^{x}} \left(\left(u^{y}\right)^2+1\right)\\
  &\quad + {\nu^{x}} \left(u^{x}\right)^2 \left(2 \left(u^{y}\right)^2+3\right)\Bigr)\\
  &\quad + \left(u^{\tau}\right)^2\Bigl({I_{31}} {u^{x}} \left(\left(u^{y}\right)^2+1\right) \left(3 \left(u^{x}\right)^2+\left(u^{y}\right)^2+3\right)\\
  &\quad + {\mathcal{C}}\Bigl(- 3 {\pi_{Q}{}^{\eta}{}_{\eta}} {u^{x}} - 3 {\pi_{Q}{}^{\eta}{}_{\eta}} \left(u^{x}\right)^3 + 3 {\pi_{Q}^{xy}} {u^{y}} + 4 {\pi_{Q}^{xy}} \left(u^{x}\right)^2 {u^{y}}\\
  &\quad - 6 {\pi_{Q}{}^{\eta}{}_{\eta}} {u^{x}} \left(u^{y}\right)^2 - 5 {\pi_{Q}{}^{\eta}{}_{\eta}} \left(u^{x}\right)^3 \left(u^{y}\right)^2 + 3 {\pi_{Q}^{xy}} \left(u^{y}\right)^3\\
  &\quad + 8 {\pi_{Q}^{xy}} \left(u^{x}\right)^2 \left(u^{y}\right)^3 - 3 {\pi_{Q}{}^{\eta}{}_{\eta}} {u^{x}} \left(u^{y}\right)^4\\
  &\quad + {\Pi_{Q}} {u^{x}} \left(\left(u^{y}\right)^2+1\right) \left(3 \left(u^{x}\right)^2+\left(u^{y}\right)^2+3\right)\\
  &\quad - 2 {\pi_{Q}^{yy}} {u^{x}} \left(\left(u^{x}\right)^2+1\right) \left(4 \left(u^{y}\right)^2+3\right)\Bigr)\Bigr)\\
  &\quad + {\mathcal{B}} \left(\left(u^{y}\right)^2+1\right)\Bigl(- 2 {\nu^{y}} {u^{x}} \left(\left(u^{x}\right)^2+1\right) {u^{y}} \left(3 \left(u^{x}\right)^2+3 \left(u^{y}\right)^2+4\right)\\
  &\quad + {\nu^{x}}\Bigl(3 \left(\left(u^{y}\right)^2+1\right)^2 + \left(u^{x}\right)^4 \left(6 \left(u^{y}\right)^2+3\right)\\
  &\quad + 2 \left(u^{x}\right)^2 \left(3 \left(u^{y}\right)^4+7 \left(u^{y}\right)^2+3\right)\Bigr)\Bigr)\Bigr) \\
(A_t)_{13,9} &= \frac{{\mathcal{B}} {D_{s}} {u^{y}} \left(\left(u^{x}\right)^2+3 \left(u^{y}\right)^2+3\right)}{6 {T} {u^{\tau}}} \\
(A_t)_{13,10} &= \frac{{\mathcal{B}} {D_{s}} {u^{x}} \left(3 \left(u^{x}\right)^2+\left(u^{y}\right)^2+3\right)}{6 {T} {u^{\tau}}} \\
(A_t)_{13,13} &= \frac{{\mathcal{C}} {D_{s}} {u^{\tau}}}{2 {T}} \\
(A_t)_{14,1} &= \frac{{D_{s}} \left(5 {\mathcal{B}'(T)} ({\nu^{x}} {u^{x}}+{\nu^{y}} {u^{y}})+3 {\mathcal{C}'(T)} {\Pi_{Q}} \left(u^{\tau}\right)^2+3 {(\partial_{T} I_{31})} \left(u^{\tau}\right)^2\right)}{6 {T} {u^{\tau}}} \\
(A_t)_{14,2} &= \frac{1}{6 {T} \left(u^{\tau}\right)^3}{D_{s}}\Bigl(5 {\mathcal{B}} \left({\nu^{x}} \left(u^{y}\right)^2+{\nu^{x}}-{\nu^{y}} {u^{x}} {u^{y}}\right) + 2 {\mathcal{A}} \left(u^{\tau}\right)^2 \left({\nu^{x}} \left(u^{y}\right)^2+{\nu^{x}}-{\nu^{y}} {u^{x}} {u^{y}}\right)\\
  &\quad + \left(u^{\tau}\right)^2 ({\mathcal{C}} (5 {\Pi_{Q}} {u^{x}}+2 {\pi_{Q}^{xy}} {u^{y}}-2 {\pi_{Q}^{yy}} {u^{x}}-2 {\pi_{Q}{}^{\eta}{}_{\eta}} {u^{x}})+5 {I_{31}} {u^{x}})\Bigr) \\
(A_t)_{14,3} &= \frac{1}{6 {T} \left(u^{\tau}\right)^3 \left(\left(u^{y}\right)^2+1\right)}{D_{s}}\Bigl(5 {\mathcal{B}} \left(\left(u^{y}\right)^2+1\right) \left(-{\nu^{x}} {u^{x}} {u^{y}}+{\nu^{y}} \left(u^{x}\right)^2+{\nu^{y}}\right)\\
  &\quad + 2 {\mathcal{A}} \left(u^{\tau}\right)^2 \left(\left(u^{y}\right)^2+1\right) \left(-{\nu^{x}} {u^{x}} {u^{y}}+{\nu^{y}} \left(u^{x}\right)^2+{\nu^{y}}\right)\\
  &\quad + \left(u^{\tau}\right)^2\Bigl(5 {I_{31}} {u^{y}} \left(\left(u^{y}\right)^2+1\right)\\
  &\quad + {\mathcal{C}}\Bigl(2 {u^{y}} \left({\pi_{Q}^{yy}} \left(u^{x}\right)^2+{\pi_{Q}^{yy}}+{\pi_{Q}{}^{\eta}{}_{\eta}} \left(u^{x}\right)^2\right)\\
  &\quad - 2 {\pi_{Q}^{xy}} {u^{x}} \left(\left(u^{y}\right)^2-1\right) + 5 {\Pi_{Q}} {u^{y}} \left(\left(u^{y}\right)^2+1\right)\Bigr)\Bigr)\Bigr) \\
(A_t)_{14,8} &= \frac{{(\partial_{\alpha} I_{31})} {D_{s}} {u^{\tau}}}{2 {T}} \\
(A_t)_{14,9} &= \frac{5 {\mathcal{B}} {D_{s}} {u^{x}}}{6 {T} {u^{\tau}}} \\
(A_t)_{14,10} &= \frac{5 {\mathcal{B}} {D_{s}} {u^{y}}}{6 {T} {u^{\tau}}} \\
(A_t)_{14,14} &= \frac{{\mathcal{C}} {D_{s}} {u^{\tau}}}{2 {T}} \\
(A_x)_{1,1} &= \frac{1}{{u^{\tau}}}\Bigl(\left(u^{\tau}\right)^2 {u^{x}} ({(\partial_{T} n)} {h}+{P''(T)} {T}+{P'(T)})\\
  &\quad + {h'(T)} \left({n} \left(u^{\tau}\right)^2 {u^{x}}+{\nu^{x}} \left(2 \left(u^{x}\right)^2+\left(u^{y}\right)^2+1\right)+{\nu^{y}} {u^{x}} {u^{y}}\right)\Bigr) \\
(A_x)_{1,2} &= \frac{1}{\left(u^{\tau}\right)^3 \left(\left(u^{y}\right)^2+1\right)}\Bigl(- {\pi^{\eta}{}_{\eta}} - {\pi^{yy}} + {P'(T)} {T} - 3 {\pi^{\eta}{}_{\eta}} \left(u^{x}\right)^2 - 3 {\pi^{yy}} \left(u^{x}\right)^2\\
  &\quad + 3 {P'(T)} {T} \left(u^{x}\right)^2 - 2 {\pi^{\eta}{}_{\eta}} \left(u^{x}\right)^4 - 2 {\pi^{yy}} \left(u^{x}\right)^4 + 2 {P'(T)} {T} \left(u^{x}\right)^4\\
  &\quad + 3 {\pi^{xy}} {u^{x}} {u^{y}} + 2 {\pi^{xy}} \left(u^{x}\right)^3 {u^{y}} - 2 {\pi^{\eta}{}_{\eta}} \left(u^{y}\right)^2 - {\pi^{yy}} \left(u^{y}\right)^2\\
  &\quad + 3 {P'(T)} {T} \left(u^{y}\right)^2 - 3 {\pi^{\eta}{}_{\eta}} \left(u^{x}\right)^2 \left(u^{y}\right)^2 - 3 {\pi^{yy}} \left(u^{x}\right)^2 \left(u^{y}\right)^2\\
  &\quad + 6 {P'(T)} {T} \left(u^{x}\right)^2 \left(u^{y}\right)^2 + 2 {P'(T)} {T} \left(u^{x}\right)^4 \left(u^{y}\right)^2 + 3 {\pi^{xy}} {u^{x}} \left(u^{y}\right)^3\\
  &\quad - {\pi^{\eta}{}_{\eta}} \left(u^{y}\right)^4 + 3 {P'(T)} {T} \left(u^{y}\right)^4 + 3 {P'(T)} {T} \left(u^{x}\right)^2 \left(u^{y}\right)^4\\
  &\quad + {P'(T)} {T} \left(u^{y}\right)^6\\
  &\quad + {\Pi} \left(\left(u^{y}\right)^2+1\right) \left(2 \left(u^{x}\right)^4+3 \left(u^{x}\right)^2 \left(\left(u^{y}\right)^2+1\right)+\left(\left(u^{y}\right)^2+1\right)^2\right)\\
  &\quad + {h} \left(\left(u^{y}\right)^2+1\right)\Bigl({\nu^{y}} {u^{y}} \left(\left(u^{y}\right)^2+1\right)\\
  &\quad + {\nu^{x}} {u^{x}} \left(2 \left(u^{x}\right)^2+3 \left(u^{y}\right)^2+3\right)\\
  &\quad + {n}\Bigl(2 \left(u^{x}\right)^4+3 \left(u^{x}\right)^2 \left(\left(u^{y}\right)^2+1\right)+\left(\left(u^{y}\right)^2+1\right)^2\Bigr)\Bigr)\Bigr) \\
(A_x)_{1,3} &= \frac{1}{\left(u^{\tau}\right)^3 \left(\left(u^{y}\right)^2+1\right)^2}\Bigl({h} \left(\left(u^{y}\right)^2+1\right)^2\Bigl({n} \left(u^{\tau}\right)^2 {u^{x}} {u^{y}}+{\nu^{x}} {u^{y}} \left(\left(u^{y}\right)^2+1\right)+{\nu^{y}} \left(\left(u^{x}\right)^3+{u^{x}}\right)\Bigr)\\
  &\quad + {\pi^{xy}}\Bigl(- 2 \left(u^{x}\right)^4 \left(\left(u^{y}\right)^2-1\right) + \left(\left(u^{y}\right)^2+1\right)^2\\
  &\quad - 3 \left(u^{x}\right)^2 \left(\left(u^{y}\right)^4-1\right)\Bigr)\\
  &\quad + {u^{x}} {u^{y}}\Bigl(3 {\pi^{yy}} + {P'(T)} {T} + 5 {\pi^{yy}} \left(u^{x}\right)^2 + {P'(T)} {T} \left(u^{x}\right)^2\\
  &\quad + 2 {\pi^{yy}} \left(u^{x}\right)^4 + 3 {\pi^{yy}} \left(u^{y}\right)^2 + 3 {P'(T)} {T} \left(u^{y}\right)^2\\
  &\quad + 3 {\pi^{yy}} \left(u^{x}\right)^2 \left(u^{y}\right)^2 + 2 {P'(T)} {T} \left(u^{x}\right)^2 \left(u^{y}\right)^2\\
  &\quad + 3 {P'(T)} {T} \left(u^{y}\right)^4 + {P'(T)} {T} \left(u^{x}\right)^2 \left(u^{y}\right)^4\\
  &\quad + {P'(T)} {T} \left(u^{y}\right)^6 + {\Pi} \left(u^{\tau}\right)^2 \left(\left(u^{y}\right)^2+1\right)^2\\
  &\quad + {\pi^{\eta}{}_{\eta}} \left(2 \left(u^{x}\right)^4+3 \left(u^{x}\right)^2 \left(\left(u^{y}\right)^2+1\right)+\left(\left(u^{y}\right)^2+1\right)^2\right)\Bigr)\Bigr) \\
(A_x)_{1,4} &= -\frac{{u^{x}} \left(\left(u^{x}\right)^2+1\right)}{{u^{\tau}} \left(\left(u^{y}\right)^2+1\right)} \\
(A_x)_{1,5} &= -\frac{{u^{\tau}} {u^{x}}}{\left(u^{y}\right)^2+1} \\
(A_x)_{1,6} &= \frac{{u^{y}} \left(2 \left(u^{x}\right)^2+\left(u^{y}\right)^2+1\right)}{{u^{\tau}} \left(\left(u^{y}\right)^2+1\right)} \\
(A_x)_{1,7} &= {u^{\tau}} {u^{x}} \\
(A_x)_{1,8} &= {(\partial_{\alpha} n)} {h} {u^{\tau}} {u^{x}} \\
(A_x)_{1,9} &= \frac{{h} \left(2 \left(u^{x}\right)^2+\left(u^{y}\right)^2+1\right)}{{u^{\tau}}} \\
(A_x)_{1,10} &= \frac{{h} {u^{x}} {u^{y}}}{{u^{\tau}}} \\
(A_x)_{2,1} &= {P'(T)} + {n} + {(\partial_{T} n)} {T} + 2 {h'(T)} {\nu^{x}} {u^{x}} + {P'(T)} \left(u^{x}\right)^2 + {(\partial_{T} n)} {h} \left(u^{x}\right)^2 + {h'(T)} {n} \left(u^{x}\right)^2\\
  &\quad + {P''(T)} {T} \left(u^{x}\right)^2 \\
(A_x)_{2,2} &= 2 \left({P'(T)} {T} {u^{x}}+{h} {n} {u^{x}}+{h} {\nu^{x}}+{\Pi} {u^{x}}-\frac{{\pi^{\eta}{}_{\eta}} {u^{x}}-{\pi^{xy}} {u^{y}}+{\pi^{yy}} {u^{x}}}{\left(u^{y}\right)^2+1}\right) \\
(A_x)_{2,3} &= \frac{2 \left({u^{y}} \left({\pi^{\eta}{}_{\eta}} \left(u^{x}\right)^2+{\pi^{yy}} \left(u^{x}\right)^2+{\pi^{yy}}\right)+{\pi^{xy}} \left({u^{x}}-{u^{x}} \left(u^{y}\right)^2\right)\right)}{\left(\left(u^{y}\right)^2+1\right)^2} \\
(A_x)_{2,4} &= -\frac{\left(u^{x}\right)^2+1}{\left(u^{y}\right)^2+1} \\
(A_x)_{2,5} &= -\frac{\left(u^{\tau}\right)^2}{\left(u^{y}\right)^2+1} \\
(A_x)_{2,6} &= \frac{2 {u^{x}} {u^{y}}}{\left(u^{y}\right)^2+1} \\
(A_x)_{2,7} &= \left(u^{x}\right)^2+1 \\
(A_x)_{2,8} &= {(\partial_{\alpha} n)} \left({h} \left(u^{x}\right)^2+{T}\right) \\
(A_x)_{2,9} &= 2 {h} {u^{x}} \\
(A_x)_{3,1} &= {h'(T)} {u^{y}} ({n} {u^{x}}+{\nu^{x}})+{h'(T)} {\nu^{y}} {u^{x}}+{u^{x}} {u^{y}} ({(\partial_{T} n)} {h}+{P''(T)} {T}+{P'(T)}) \\
(A_x)_{3,2} &= {u^{y}} ({P'(T)} {T}+{\Pi})+{h} ({n} {u^{y}}+{\nu^{y}}) \\
(A_x)_{3,3} &= {u^{x}} ({P'(T)} {T}+{\Pi})+{h} ({n} {u^{x}}+{\nu^{x}}) \\
(A_x)_{3,6} &= 1 \\
(A_x)_{3,7} &= {u^{x}} {u^{y}} \\
(A_x)_{3,8} &= {(\partial_{\alpha} n)} {h} {u^{x}} {u^{y}} \\
(A_x)_{3,9} &= {h} {u^{y}} \\
(A_x)_{3,10} &= {h} {u^{x}} \\
(A_x)_{4,2} &= {\delta_{\pi\pi}} {\pi^{yy}}-\frac{2 {\pi^{xy}} {\tau_{\pi}} {u^{x}} \left(\left(u^{y}\right)^2+1\right) {u^{y}}}{\left(u^{\tau}\right)^2}+\frac{2 {\pi^{yy}} {\tau_{\pi}} \left(u^{x}\right)^2 \left(u^{y}\right)^2}{\left(u^{\tau}\right)^2}-\frac{2}{3} \left(\left(u^{y}\right)^2+1\right) {\eta} \\
(A_x)_{4,3} &= \frac{2 {u^{x}} {u^{y}} \left({\pi^{xy}} {\tau_{\pi}} {u^{x}} {u^{y}}-{\pi^{yy}} {\tau_{\pi}} \left(\left(u^{x}\right)^2+1\right)+\left(u^{\tau}\right)^2 {\eta}\right)}{\left(u^{\tau}\right)^2} \\
(A_x)_{4,4} &= {\tau_{\pi}} {u^{x}} \\
(A_x)_{5,2} &= \frac{3 {\delta_{\pi\pi}} {\pi^{\eta}{}_{\eta}}-2 {\eta}}{3 {\tau}^2} \\
(A_x)_{5,5} &= \frac{{\tau_{\pi}} {u^{x}}}{{\tau}^2} \\
(A_x)_{6,2} &= {\delta_{\pi\pi}} {\pi^{xy}}\\
  &\quad + \frac{{u^{x}} \left({u^{y}} \left(3 {\pi^{\eta}{}_{\eta}} {\tau_{\pi}} \left(u^{\tau}\right)^2+3 {\pi^{yy}} \left(2 {\tau_{\pi}} \left(u^{x}\right)^2+{\tau_{\pi}}\right)+\left(u^{\tau}\right)^2 {\eta}\right)-3 {\pi^{xy}} {\tau_{\pi}} {u^{x}} \left(2 \left(u^{y}\right)^2+1\right)\right)}{3 \left(u^{\tau}\right)^2} \\
(A_x)_{6,3} &= \frac{1}{\left(u^{\tau}\right)^2 \left(\left(u^{y}\right)^2+1\right)}\Bigl(- {\pi^{yy}} {\tau_{\pi}} \left(u^{x}\right)^2 \left(\left(u^{x}\right)^2+1\right) \left(2 \left(u^{y}\right)^2+1\right)\\
  &\quad + {\pi^{xy}} {\tau_{\pi}} {u^{x}} {u^{y}} \left(\left(2 \left(u^{x}\right)^2-1\right) \left(u^{y}\right)^2-1\right)\\
  &\quad - \left(u^{\tau}\right)^2 \left({\pi^{\eta}{}_{\eta}} {\tau_{\pi}} \left(u^{x}\right)^2 \left(u^{y}\right)^2-\left(\left(u^{x}\right)^2+1\right) \left(\left(u^{y}\right)^2+1\right) {\eta}\right)\Bigr) \\
(A_x)_{6,6} &= {\tau_{\pi}} {u^{x}} \\
(A_x)_{7,2} &= {\zeta} \\
(A_x)_{7,7} &= {\tau_{\Pi}} {u^{x}} \\
(A_x)_{8,1} &= {(\partial_{T} n)} {u^{x}} \\
(A_x)_{8,2} &= {n} \\
(A_x)_{8,8} &= {(\partial_{\alpha} n)} {u^{x}} \\
(A_x)_{8,9} &= 1 \\
(A_x)_{9,1} &= \frac{{D_{s}} \left({h'(T)} {\nu^{x}} {u^{x}}+{(\partial_{T} n)} {T} \left(\left(u^{x}\right)^2+1\right)+{n} \left(\left(u^{x}\right)^2+1\right)\right)}{{T}} \\
(A_x)_{9,2} &= \frac{{D_{s}} {h} \left({n} \left(u^{\tau}\right)^2 {u^{x}}+{\nu^{x}} \left(\left(u^{x}\right)^2 \left(-\left(u^{y}\right)^2\right)+\left(u^{x}\right)^2+2 \left(u^{y}\right)^2+2\right)+{\nu^{y}} \left(u^{x}\right)^3 {u^{y}}\right)}{{T} \left(u^{\tau}\right)^2} \\
(A_x)_{9,3} &= -\frac{{D_{s}} {h} \left(u^{x}\right)^2 \left(-{\nu^{x}} {u^{x}} {u^{y}}+{\nu^{y}} \left(u^{x}\right)^2+{\nu^{y}}\right)}{{T} \left(u^{\tau}\right)^2} \\
(A_x)_{9,8} &= {(\partial_{\alpha} n)} {D_{s}} \left(\left(u^{x}\right)^2+1\right) \\
(A_x)_{9,9} &= \frac{{D_{s}} {h} {u^{x}}}{{T}} \\
(A_x)_{10,1} &= \frac{{D_{s}} {u^{x}} ({h'(T)} {\nu^{y}}+{u^{y}} ({(\partial_{T} n)} {T}+{n}))}{{T}} \\
(A_x)_{10,2} &= \frac{{D_{s}} {h} \left(\left(u^{y}\right)^2+1\right) \left(-{\nu^{x}} {u^{x}} {u^{y}}+{\nu^{y}} \left(u^{x}\right)^2+{\nu^{y}}\right)}{{T} \left(u^{\tau}\right)^2} \\
(A_x)_{10,3} &= \frac{{D_{s}} {h} \left({n} \left(u^{\tau}\right)^2 {u^{x}}+{\nu^{x}} \left(\left(u^{x}\right)^2+1\right) \left(\left(u^{y}\right)^2+1\right)-{\nu^{y}} {u^{x}} \left(\left(u^{x}\right)^2+1\right) {u^{y}}\right)}{{T} \left(u^{\tau}\right)^2} \\
(A_x)_{10,8} &= {(\partial_{\alpha} n)} {D_{s}} {u^{x}} {u^{y}} \\
(A_x)_{10,10} &= \frac{{D_{s}} {h} {u^{x}}}{{T}} \\
(A_x)_{11,1} &= \frac{{D_{s}} \left(3 {\mathcal{C}'(T)} {\pi_{Q}^{yy}} {u^{x}}-2 {\mathcal{B}'(T)} \left({\nu^{x}} \left(u^{y}\right)^2+{\nu^{x}}-3 {\nu^{y}} {u^{x}} {u^{y}}\right)\right)}{6 {T}} \\
(A_x)_{11,2} &= -\frac{1}{6 {T} \left(u^{\tau}\right)^2}{D_{s}}\Bigl(2 {I_{31}} \left(u^{\tau}\right)^2 \left(\left(u^{y}\right)^2+1\right)\\
  &\quad + 2 {u^{x}} \left({\mathcal{A}} \left(u^{y}\right)^2+{\mathcal{A}}+3 {\mathcal{B}} \left(u^{y}\right)^2\right) \left({\nu^{x}} \left(u^{y}\right)^2+{\nu^{x}}-{\nu^{y}} {u^{x}} {u^{y}}\right)\\
  &\quad + {\mathcal{C}}\Bigl(2 {\Pi_{Q}} \left(u^{\tau}\right)^2 \left(\left(u^{y}\right)^2+1\right)\\
  &\quad - 2 \left(\left(u^{y}\right)^2+1\right) \left({\pi_{Q}{}^{\eta}{}_{\eta}} \left(u^{\tau}\right)^2-4 {\pi_{Q}^{xy}} {u^{x}} {u^{y}}\right)\\
  &\quad - {\pi_{Q}^{yy}} \left(\left(u^{x}\right)^2 \left(8 \left(u^{y}\right)^2+5\right)+5 \left(\left(u^{y}\right)^2+1\right)\right)\Bigr)\Bigr) \\
(A_x)_{11,3} &= \frac{1}{3 {T} \left(u^{\tau}\right)^2}{D_{s}}\Bigl({u^{x}}\Bigl(3 {I_{31}} \left(u^{\tau}\right)^2 {u^{y}} - {\mathcal{A}} {\nu^{y}} \left(u^{x}\right)^2 \left(\left(u^{y}\right)^2-2\right) + 2 {\mathcal{A}} {\nu^{y}} \left(\left(u^{y}\right)^2+1\right)\\
  &\quad + {\mathcal{A}} {\nu^{x}} {u^{x}} {u^{y}} \left(\left(u^{y}\right)^2+1\right)\\
  &\quad + 3 {\mathcal{B}} \left({\nu^{x}} {u^{x}} \left(u^{y}\right)^3-{\nu^{y}} \left(u^{x}\right)^2 \left(\left(u^{y}\right)^2-1\right)+{\nu^{y}}\right)\Bigr)\\
  &\quad + {\mathcal{C}}\Bigl(- {u^{x}} {u^{y}} \left(-3 {\Pi_{Q}} \left(u^{\tau}\right)^2+4 {\pi_{Q}^{yy}} \left(\left(u^{x}\right)^2+1\right)+{\pi_{Q}{}^{\eta}{}_{\eta}} \left(u^{\tau}\right)^2\right)\\
  &\quad + 2 {\pi_{Q}^{xy}} \left(\left(u^{x}\right)^2 \left(2 \left(u^{y}\right)^2+1\right)+\left(u^{y}\right)^2+1\right)\Bigr)\Bigr) \\
(A_x)_{11,9} &= -\frac{{\mathcal{B}} {D_{s}} \left(\left(u^{y}\right)^2+1\right)}{3 {T}} \\
(A_x)_{11,10} &= \frac{{\mathcal{B}} {D_{s}} {u^{x}} {u^{y}}}{{T}} \\
(A_x)_{11,11} &= \frac{{\mathcal{C}} {D_{s}} {u^{x}}}{2 {T}} \\
(A_x)_{12,1} &= -\frac{{D_{s}} (2 {\mathcal{B}'(T)} {\nu^{x}}-3 {\mathcal{C}'(T)} {\pi_{Q}{}^{\eta}{}_{\eta}} {u^{x}})}{6 {T}} \\
(A_x)_{12,2} &= -\frac{1}{6 {T} \left(u^{\tau}\right)^2}{D_{s}}\Bigl(2 {I_{31}} \left(u^{\tau}\right)^2 + 2 {\mathcal{A}} {u^{x}} \left({\nu^{x}} \left(u^{y}\right)^2+{\nu^{x}}-{\nu^{y}} {u^{x}} {u^{y}}\right)\\
  &\quad + {\mathcal{C}}\Bigl(- 5 {\pi_{Q}{}^{\eta}{}_{\eta}} + 2 {\Pi_{Q}} \left(u^{\tau}\right)^2 - 5 {\pi_{Q}{}^{\eta}{}_{\eta}} \left(u^{x}\right)^2 - 2 {\pi_{Q}^{yy}} \left(\left(u^{x}\right)^2+1\right)\\
  &\quad + 2 {\pi_{Q}^{xy}} {u^{x}} {u^{y}} - 5 {\pi_{Q}{}^{\eta}{}_{\eta}} \left(u^{y}\right)^2\Bigr)\Bigr) \\
(A_x)_{12,3} &= \frac{1}{3 {T} \left(u^{\tau}\right)^2 \left(\left(u^{y}\right)^2+1\right)}{D_{s}}\Bigl(- {\mathcal{C}} {u^{x}} {u^{y}} \left({\pi_{Q}^{yy}} \left(\left(u^{x}\right)^2+1\right)+{\pi_{Q}{}^{\eta}{}_{\eta}} \left(u^{\tau}\right)^2\right)\\
  &\quad - {\mathcal{A}} {u^{x}} \left(\left(u^{y}\right)^2+1\right) \left(-{\nu^{x}} {u^{x}} {u^{y}}+{\nu^{y}} \left(u^{x}\right)^2+{\nu^{y}}\right)\\
  &\quad + {\mathcal{C}} {\pi_{Q}^{xy}} \left(\left(u^{x}\right)^2 \left(\left(u^{y}\right)^2-1\right)-\left(u^{y}\right)^2-1\right)\Bigr) \\
(A_x)_{12,9} &= -\frac{{\mathcal{B}} {D_{s}}}{3 {T}} \\
(A_x)_{12,12} &= \frac{{\mathcal{C}} {D_{s}} {u^{x}}}{2 {T}} \\
(A_x)_{13,1} &= \frac{{D_{s}} \left({\mathcal{B}'(T)} {\nu^{x}} {u^{x}} {u^{y}}+3 {\mathcal{B}'(T)} {\nu^{y}} \left(\left(u^{x}\right)^2+1\right)+3 {\mathcal{C}'(T)} {\pi_{Q}^{xy}} {u^{x}}\right)}{6 {T}} \\
(A_x)_{13,2} &= \frac{1}{6 {T} \left(u^{\tau}\right)^2}{D_{s}}\Bigl(3 {\mathcal{B}} {\nu^{y}} {u^{x}} + 3 {\mathcal{B}} {\nu^{y}} \left(u^{x}\right)^3 - 3 {\mathcal{B}} {\nu^{x}} {u^{y}} + {I_{31}} {u^{x}} {u^{y}} - 6 {\mathcal{B}} {\nu^{x}} \left(u^{x}\right)^2 {u^{y}} + {I_{31}} \left(u^{x}\right)^3 {u^{y}}\\
  &\quad + 6 {\mathcal{B}} {\nu^{y}} {u^{x}} \left(u^{y}\right)^2 + 6 {\mathcal{B}} {\nu^{y}} \left(u^{x}\right)^3 \left(u^{y}\right)^2 - 3 {\mathcal{B}} {\nu^{x}} \left(u^{y}\right)^3 + {I_{31}} {u^{x}} \left(u^{y}\right)^3\\
  &\quad - 6 {\mathcal{B}} {\nu^{x}} \left(u^{x}\right)^2 \left(u^{y}\right)^3\\
  &\quad + {\mathcal{A}} {u^{x}}\Bigl(3 {\nu^{y}} \left(\left(u^{y}\right)^2+1\right) - 2 {\nu^{x}} {u^{x}} {u^{y}} \left(\left(u^{y}\right)^2+1\right)\\
  &\quad + {\nu^{y}} \left(u^{x}\right)^2 \left(2 \left(u^{y}\right)^2+3\right)\Bigr)\\
  &\quad + {\mathcal{C}}\Bigl({u^{x}} {u^{y}} \left({\Pi_{Q}} \left(u^{\tau}\right)^2+{\pi_{Q}^{yy}} \left(8 \left(u^{x}\right)^2+5\right)+5 {\pi_{Q}{}^{\eta}{}_{\eta}} \left(u^{\tau}\right)^2\right)\\
  &\quad + {\pi_{Q}^{xy}} \left(\left(u^{x}\right)^2 \left(3-8 \left(u^{y}\right)^2\right)+6 \left(\left(u^{y}\right)^2+1\right)\right)\Bigr)\Bigr) \\
(A_x)_{13,3} &= \frac{1}{6 {T} \left(u^{\tau}\right)^2 \left(\left(u^{y}\right)^2+1\right)}{D_{s}}\Bigl(3 {I_{31}} \left(u^{\tau}\right)^2 \left(\left(u^{x}\right)^2+1\right) \left(\left(u^{y}\right)^2+1\right)\\
  &\quad + {\mathcal{C}}\Bigl(- 3 {\pi_{Q}{}^{\eta}{}_{\eta}} - 6 {\pi_{Q}{}^{\eta}{}_{\eta}} \left(u^{x}\right)^2 - 3 {\pi_{Q}{}^{\eta}{}_{\eta}} \left(u^{x}\right)^4 + {\pi_{Q}^{xy}} {u^{x}} {u^{y}}\\
  &\quad + 4 {\pi_{Q}^{xy}} \left(u^{x}\right)^3 {u^{y}} - 6 {\pi_{Q}{}^{\eta}{}_{\eta}} \left(u^{y}\right)^2 - 11 {\pi_{Q}{}^{\eta}{}_{\eta}} \left(u^{x}\right)^2 \left(u^{y}\right)^2\\
  &\quad - 5 {\pi_{Q}{}^{\eta}{}_{\eta}} \left(u^{x}\right)^4 \left(u^{y}\right)^2 + {\pi_{Q}^{xy}} {u^{x}} \left(u^{y}\right)^3 + 8 {\pi_{Q}^{xy}} \left(u^{x}\right)^3 \left(u^{y}\right)^3\\
  &\quad - 3 {\pi_{Q}{}^{\eta}{}_{\eta}} \left(u^{y}\right)^4 - 5 {\pi_{Q}{}^{\eta}{}_{\eta}} \left(u^{x}\right)^2 \left(u^{y}\right)^4\\
  &\quad + 3 {\Pi_{Q}} \left(u^{\tau}\right)^2 \left(\left(u^{x}\right)^2+1\right) \left(\left(u^{y}\right)^2+1\right)\\
  &\quad - {\pi_{Q}^{yy}} \left(\left(u^{x}\right)^2+1\right)\Bigl(\left(u^{x}\right)^2 \left(8 \left(u^{y}\right)^2+6\right)+3 \left(\left(u^{y}\right)^2+1\right)\Bigr)\Bigr)\\
  &\quad + \left(u^{y}\right)^2+1\Bigl({\mathcal{A}} {u^{x}}\Bigl(- 2 {\nu^{y}} {u^{x}} \left(\left(u^{x}\right)^2+1\right) {u^{y}} + 3 {\nu^{x}} \left(\left(u^{y}\right)^2+1\right)\\
  &\quad + {\nu^{x}} \left(u^{x}\right)^2 \left(2 \left(u^{y}\right)^2+3\right)\Bigr)\\
  &\quad + 3{\mathcal{B}} \left(\left(u^{x}\right)^2+1\right)\Bigl({\nu^{x}} \left(2 {u^{x}} \left(u^{y}\right)^2+{u^{x}}\right)-{\nu^{y}} \left(2 \left(u^{x}\right)^2+1\right) {u^{y}}\Bigr)\Bigr)\Bigr) \\
(A_x)_{13,9} &= \frac{{\mathcal{B}} {D_{s}} {u^{x}} {u^{y}}}{6 {T}} \\
(A_x)_{13,10} &= \frac{{\mathcal{B}} {D_{s}} \left(\left(u^{x}\right)^2+1\right)}{2 {T}} \\
(A_x)_{13,13} &= \frac{{\mathcal{C}} {D_{s}} {u^{x}}}{2 {T}} \\
(A_x)_{14,1} &= \frac{{D_{s}} (5 {\mathcal{B}'(T)} {\nu^{x}}+3 {u^{x}} ({\mathcal{C}'(T)} {\Pi_{Q}}+{(\partial_{T} I_{31})}))}{6 {T}} \\
(A_x)_{14,2} &= \frac{1}{6 {T} \left(u^{\tau}\right)^2}{D_{s}}\Bigl(5 {I_{31}} \left(u^{\tau}\right)^2 + 5 {\mathcal{C}} {\Pi_{Q}} \left(u^{\tau}\right)^2 + 2 {\mathcal{A}} {u^{x}} \left({\nu^{x}} \left(u^{y}\right)^2+{\nu^{x}}-{\nu^{y}} {u^{x}} {u^{y}}\right)\\
  &\quad - 2 {\mathcal{C}} \left(-{\pi_{Q}^{xy}} {u^{x}} {u^{y}}+{\pi_{Q}^{yy}} \left(u^{x}\right)^2+{\pi_{Q}^{yy}}+{\pi_{Q}{}^{\eta}{}_{\eta}} \left(u^{x}\right)^2+{\pi_{Q}{}^{\eta}{}_{\eta}} \left(u^{y}\right)^2+{\pi_{Q}{}^{\eta}{}_{\eta}}\right)\Bigr) \\
(A_x)_{14,3} &= \frac{1}{3 {T} \left(u^{\tau}\right)^2 \left(\left(u^{y}\right)^2+1\right)}{D_{s}}\Bigl({\mathcal{C}} {u^{x}} {u^{y}} \left({\pi_{Q}^{yy}} \left(\left(u^{x}\right)^2+1\right)+{\pi_{Q}{}^{\eta}{}_{\eta}} \left(u^{\tau}\right)^2\right)\\
  &\quad + {\mathcal{A}} {u^{x}} \left(\left(u^{y}\right)^2+1\right) \left(-{\nu^{x}} {u^{x}} {u^{y}}+{\nu^{y}} \left(u^{x}\right)^2+{\nu^{y}}\right)\\
  &\quad + {\mathcal{C}} {\pi_{Q}^{xy}} \left(-\left(\left(u^{x}\right)^2 \left(\left(u^{y}\right)^2-1\right)\right)+\left(u^{y}\right)^2+1\right)\Bigr) \\
(A_x)_{14,8} &= \frac{{(\partial_{\alpha} I_{31})} {D_{s}} {u^{x}}}{2 {T}} \\
(A_x)_{14,9} &= \frac{5 {\mathcal{B}} {D_{s}}}{6 {T}} \\
(A_x)_{14,14} &= \frac{{\mathcal{C}} {D_{s}} {u^{x}}}{2 {T}} \\
(A_y)_{1,1} &= \frac{1}{{u^{\tau}}}\Bigl(\left(u^{\tau}\right)^2 {u^{y}} ({(\partial_{T} n)} {h}+{P''(T)} {T}+{P'(T)})\\
  &\quad + {h'(T)} \left({n} \left(u^{\tau}\right)^2 {u^{y}}+{\nu^{x}} {u^{x}} {u^{y}}+{\nu^{y}} \left(\left(u^{x}\right)^2+2 \left(u^{y}\right)^2+1\right)\right)\Bigr) \\
(A_y)_{1,2} &= \frac{1}{\left(u^{\tau}\right)^3}\Bigl({u^{x}} {u^{y}} \left({P'(T)} {T} \left(u^{\tau}\right)^2+{\Pi} \left(u^{\tau}\right)^2-{\pi^{yy}}\right) + {\pi^{xy}} \left(\left(u^{y}\right)^2+1\right)\\
  &\quad + {h} \left({n} \left(u^{\tau}\right)^2 {u^{x}} {u^{y}}+{\nu^{x}} {u^{y}} \left(\left(u^{y}\right)^2+1\right)+{\nu^{y}} \left(\left(u^{x}\right)^3+{u^{x}}\right)\right)\Bigr) \\
(A_y)_{1,3} &= \frac{1}{\left(u^{\tau}\right)^3}\Bigl({\pi^{yy}} + {P'(T)} {T} + {\pi^{yy}} \left(u^{x}\right)^2 + 2 {P'(T)} {T} \left(u^{x}\right)^2 + {P'(T)} {T} \left(u^{x}\right)^4 - {\pi^{xy}} {u^{x}} {u^{y}}\\
  &\quad + 3 {P'(T)} {T} \left(u^{y}\right)^2 + 3 {P'(T)} {T} \left(u^{x}\right)^2 \left(u^{y}\right)^2 + 2 {P'(T)} {T} \left(u^{y}\right)^4\\
  &\quad + {\Pi} \left(\left(u^{x}\right)^4+\left(u^{x}\right)^2 \left(3 \left(u^{y}\right)^2+2\right)+2 \left(u^{y}\right)^4+3 \left(u^{y}\right)^2+1\right)\\
  &\quad + {h}\Bigl({\nu^{x}} \left(\left(u^{x}\right)^3+{u^{x}}\right) + {\nu^{y}} {u^{y}} \left(3 \left(u^{x}\right)^2+2 \left(u^{y}\right)^2+3\right)\\
  &\quad + {n} \left(\left(u^{x}\right)^4+\left(u^{x}\right)^2 \left(3 \left(u^{y}\right)^2+2\right)+2 \left(u^{y}\right)^4+3 \left(u^{y}\right)^2+1\right)\Bigr)\Bigr) \\
(A_y)_{1,4} &= \frac{{u^{y}}}{{u^{\tau}}} \\
(A_y)_{1,6} &= \frac{{u^{x}}}{{u^{\tau}}} \\
(A_y)_{1,7} &= {u^{\tau}} {u^{y}} \\
(A_y)_{1,8} &= {(\partial_{\alpha} n)} {h} {u^{\tau}} {u^{y}} \\
(A_y)_{1,9} &= \frac{{h} {u^{x}} {u^{y}}}{{u^{\tau}}} \\
(A_y)_{1,10} &= \frac{{h} \left(\left(u^{x}\right)^2+2 \left(u^{y}\right)^2+1\right)}{{u^{\tau}}} \\
(A_y)_{2,1} &= {h'(T)} {u^{y}} ({n} {u^{x}}+{\nu^{x}})+{h'(T)} {\nu^{y}} {u^{x}}+{u^{x}} {u^{y}} ({(\partial_{T} n)} {h}+{P''(T)} {T}+{P'(T)}) \\
(A_y)_{2,2} &= {u^{y}} ({P'(T)} {T}+{\Pi})+{h} ({n} {u^{y}}+{\nu^{y}}) \\
(A_y)_{2,3} &= {u^{x}} ({P'(T)} {T}+{\Pi})+{h} ({n} {u^{x}}+{\nu^{x}}) \\
(A_y)_{2,6} &= 1 \\
(A_y)_{2,7} &= {u^{x}} {u^{y}} \\
(A_y)_{2,8} &= {(\partial_{\alpha} n)} {h} {u^{x}} {u^{y}} \\
(A_y)_{2,9} &= {h} {u^{y}} \\
(A_y)_{2,10} &= {h} {u^{x}} \\
(A_y)_{3,1} &= {P'(T)} + {n} + {(\partial_{T} n)} {T} + 2 {h'(T)} {\nu^{y}} {u^{y}} + {P'(T)} \left(u^{y}\right)^2 + {(\partial_{T} n)} {h} \left(u^{y}\right)^2 + {h'(T)} {n} \left(u^{y}\right)^2\\
  &\quad + {P''(T)} {T} \left(u^{y}\right)^2 \\
(A_y)_{3,3} &= 2 ({u^{y}} ({P'(T)} {T}+{\Pi})+{h} ({n} {u^{y}}+{\nu^{y}})) \\
(A_y)_{3,4} &= 1 \\
(A_y)_{3,7} &= \left(u^{y}\right)^2+1 \\
(A_y)_{3,8} &= {(\partial_{\alpha} n)} \left({h} \left(u^{y}\right)^2+{T}\right) \\
(A_y)_{3,10} &= 2 {h} {u^{y}} \\
(A_y)_{4,2} &= -\frac{2 {\tau_{\pi}} \left(u^{y}\right)^2 \left({\pi^{xy}} \left(u^{y}\right)^2+{\pi^{xy}}-{\pi^{yy}} {u^{x}} {u^{y}}\right)}{\left(u^{\tau}\right)^2} \\
(A_y)_{4,3} &= \frac{1}{3 \left(u^{\tau}\right)^2}\Bigl(3 {\delta_{\pi\pi}} {\pi^{yy}} \left(u^{\tau}\right)^2 - 6 {\pi^{yy}} {\tau_{\pi}} \left(\left(u^{x}\right)^2+1\right) \left(u^{y}\right)^2 + 6 {\pi^{xy}} {\tau_{\pi}} {u^{x}} \left(u^{y}\right)^3 + 4 {\eta}\\
  &\quad + 4 \left(u^{x}\right)^2 {\eta} + 8 \left(u^{y}\right)^2 {\eta} + 4 \left(u^{x}\right)^2 \left(u^{y}\right)^2 {\eta} + 4 \left(u^{y}\right)^4 {\eta}\Bigr) \\
(A_y)_{4,4} &= {\tau_{\pi}} {u^{y}} \\
(A_y)_{5,3} &= \frac{3 {\delta_{\pi\pi}} {\pi^{\eta}{}_{\eta}}-2 {\eta}}{3 {\tau}^2} \\
(A_y)_{5,5} &= \frac{{\tau_{\pi}} {u^{y}}}{{\tau}^2} \\
(A_y)_{6,2} &= \frac{1}{\left(u^{\tau}\right)^2}\Bigl({\pi^{yy}} {\tau_{\pi}} \left(u^{y}\right)^2 + {\pi^{\eta}{}_{\eta}} {\tau_{\pi}} \left(u^{\tau}\right)^2 \left(u^{y}\right)^2 + 2 {\pi^{yy}} {\tau_{\pi}} \left(u^{x}\right)^2 \left(u^{y}\right)^2\\
  &\quad - {\pi^{xy}} {\tau_{\pi}} {u^{x}} {u^{y}} \left(2 \left(u^{y}\right)^2+1\right) + {\eta} + \left(u^{x}\right)^2 {\eta} + 2 \left(u^{y}\right)^2 {\eta} + \left(u^{x}\right)^2 \left(u^{y}\right)^2 {\eta} + \left(u^{y}\right)^4 {\eta}\Bigr) \\
(A_y)_{6,3} &= {\delta_{\pi\pi}} {\pi^{xy}}\\
  &\quad - \frac{1}{3 \left(u^{\tau}\right)^2 \left(\left(u^{y}\right)^2+1\right)}{u^{y}}\Bigl(3 {\pi^{yy}} {\tau_{\pi}} {u^{x}} \left(\left(u^{x}\right)^2+1\right) \left(2 \left(u^{y}\right)^2+1\right)\\
  &\quad + 3 {\pi^{xy}} {\tau_{\pi}} {u^{y}} \left(\left(1-2 \left(u^{x}\right)^2\right) \left(u^{y}\right)^2+1\right)\\
  &\quad + \left(u^{\tau}\right)^2 {u^{x}} \left(3 {\pi^{\eta}{}_{\eta}} {\tau_{\pi}} \left(u^{y}\right)^2-\left(\left(u^{y}\right)^2+1\right) {\eta}\right)\Bigr) \\
(A_y)_{6,6} &= {\tau_{\pi}} {u^{y}} \\
(A_y)_{7,3} &= {\zeta} \\
(A_y)_{7,7} &= {\tau_{\Pi}} {u^{y}} \\
(A_y)_{8,1} &= {(\partial_{T} n)} {u^{y}} \\
(A_y)_{8,3} &= {n} \\
(A_y)_{8,8} &= {(\partial_{\alpha} n)} {u^{y}} \\
(A_y)_{8,10} &= 1 \\
(A_y)_{9,1} &= \frac{{D_{s}} {u^{y}} ({h'(T)} {\nu^{x}}+{u^{x}} ({(\partial_{T} n)} {T}+{n}))}{{T}} \\
(A_y)_{9,2} &= \frac{{D_{s}} {h} \left({n} \left(u^{\tau}\right)^2 {u^{y}}-{\nu^{x}} {u^{x}} {u^{y}} \left(\left(u^{y}\right)^2+1\right)+{\nu^{y}} \left(\left(u^{x}\right)^2+1\right) \left(\left(u^{y}\right)^2+1\right)\right)}{{T} \left(u^{\tau}\right)^2} \\
(A_y)_{9,3} &= \frac{{D_{s}} {h} \left(\left(u^{x}\right)^2+1\right) \left({\nu^{x}} \left(u^{y}\right)^2+{\nu^{x}}-{\nu^{y}} {u^{x}} {u^{y}}\right)}{{T} \left(u^{\tau}\right)^2} \\
(A_y)_{9,8} &= {(\partial_{\alpha} n)} {D_{s}} {u^{x}} {u^{y}} \\
(A_y)_{9,9} &= \frac{{D_{s}} {h} {u^{y}}}{{T}} \\
(A_y)_{10,1} &= \frac{{D_{s}} \left({h'(T)} {\nu^{y}} {u^{y}}+{(\partial_{T} n)} {T} \left(\left(u^{y}\right)^2+1\right)+{n} \left(\left(u^{y}\right)^2+1\right)\right)}{{T}} \\
(A_y)_{10,2} &= -\frac{{D_{s}} {h} \left(u^{y}\right)^2 \left({\nu^{x}} \left(u^{y}\right)^2+{\nu^{x}}-{\nu^{y}} {u^{x}} {u^{y}}\right)}{{T} \left(u^{\tau}\right)^2} \\
(A_y)_{10,3} &= \frac{{D_{s}} {h} \left({n} \left(u^{\tau}\right)^2 {u^{y}}+{\nu^{x}} {u^{x}} \left(u^{y}\right)^3+{\nu^{y}} \left(-\left(\left(u^{x}\right)^2 \left(\left(u^{y}\right)^2-2\right)\right)+\left(u^{y}\right)^2+2\right)\right)}{{T} \left(u^{\tau}\right)^2} \\
(A_y)_{10,8} &= {(\partial_{\alpha} n)} {D_{s}} \left(\left(u^{y}\right)^2+1\right) \\
(A_y)_{10,10} &= \frac{{D_{s}} {h} {u^{y}}}{{T}} \\
(A_y)_{11,1} &= \frac{{D_{s}} \left(4 {\mathcal{B}'(T)} {\nu^{y}} \left(\left(u^{y}\right)^2+1\right)+3 {\mathcal{C}'(T)} {\pi_{Q}^{yy}} {u^{y}}\right)}{6 {T}} \\
(A_y)_{11,2} &= -\frac{1}{3 {T} \left(u^{\tau}\right)^2}{D_{s}}\Bigl({u^{y}} \left(\left(u^{y}\right)^2+1\right) ({\mathcal{A}}+3 {\mathcal{B}}) \left({\nu^{x}} \left(u^{y}\right)^2+{\nu^{x}}-{\nu^{y}} {u^{x}} {u^{y}}\right)\\
  &\quad + {\mathcal{C}} \left(4 \left(u^{y}\right)^2+1\right) \left({\pi_{Q}^{xy}} \left(u^{y}\right)^2+{\pi_{Q}^{xy}}-{\pi_{Q}^{yy}} {u^{x}} {u^{y}}\right)\Bigr) \\
(A_y)_{11,3} &= \frac{1}{6 {T} \left(u^{\tau}\right)^2}{D_{s}}\Bigl(4 {I_{31}} \left(u^{\tau}\right)^2 \left(\left(u^{y}\right)^2+1\right)\\
  &\quad + {\mathcal{C}}\Bigl(4 {\Pi_{Q}} \left(u^{\tau}\right)^2 \left(\left(u^{y}\right)^2+1\right) + 2 {\pi_{Q}^{xy}} {u^{x}} {u^{y}} \left(4 \left(u^{y}\right)^2+1\right)\\
  &\quad + {\pi_{Q}^{yy}} \left(\left(u^{x}\right)^2 \left(7-8 \left(u^{y}\right)^2\right)+\left(u^{y}\right)^2+7\right)\Bigr)\\
  &\quad + 2{u^{y}}\Bigl(3 {\mathcal{B}} {u^{x}} {u^{y}} \left({\nu^{x}} \left(u^{y}\right)^2+{\nu^{x}}-{\nu^{y}} {u^{x}} {u^{y}}\right)\\
  &\quad + {\mathcal{A}}\Bigl(- {\nu^{y}} \left(u^{x}\right)^2 \left(\left(u^{y}\right)^2-2\right) + 2 {\nu^{y}} \left(\left(u^{y}\right)^2+1\right)\\
  &\quad + {\nu^{x}} {u^{x}} {u^{y}} \left(\left(u^{y}\right)^2+1\right)\Bigr)\Bigr)\Bigr) \\
(A_y)_{11,10} &= \frac{2 {\mathcal{B}} {D_{s}} \left(\left(u^{y}\right)^2+1\right)}{3 {T}} \\
(A_y)_{11,11} &= \frac{{\mathcal{C}} {D_{s}} {u^{y}}}{2 {T}} \\
(A_y)_{12,1} &= -\frac{{D_{s}} (2 {\mathcal{B}'(T)} {\nu^{y}}-3 {\mathcal{C}'(T)} {\pi_{Q}{}^{\eta}{}_{\eta}} {u^{y}})}{6 {T}} \\
(A_y)_{12,2} &= -\frac{{D_{s}} \left({\mathcal{A}} {u^{y}} \left({\nu^{x}} \left(u^{y}\right)^2+{\nu^{x}}-{\nu^{y}} {u^{x}} {u^{y}}\right)+{\mathcal{C}} \left({\pi_{Q}^{xy}} \left(u^{y}\right)^2+{\pi_{Q}^{xy}}-{\pi_{Q}^{yy}} {u^{x}} {u^{y}}\right)\right)}{3 {T} \left(u^{\tau}\right)^2} \\
(A_y)_{12,3} &= -\frac{1}{6 {T} \left(u^{\tau}\right)^2}{D_{s}}\Bigl(2 {I_{31}} \left(u^{\tau}\right)^2 + 2 {\mathcal{A}} {u^{y}} \left(-{\nu^{x}} {u^{x}} {u^{y}}+{\nu^{y}} \left(u^{x}\right)^2+{\nu^{y}}\right)\\
  &\quad + {\mathcal{C}}\Bigl(- 3 {\pi_{Q}{}^{\eta}{}_{\eta}} + 2 {\Pi_{Q}} \left(u^{\tau}\right)^2 - 3 {\pi_{Q}{}^{\eta}{}_{\eta}} \left(u^{x}\right)^2 + 2 {\pi_{Q}^{yy}} \left(\left(u^{x}\right)^2+1\right)\\
  &\quad - 2 {\pi_{Q}^{xy}} {u^{x}} {u^{y}} - 3 {\pi_{Q}{}^{\eta}{}_{\eta}} \left(u^{y}\right)^2\Bigr)\Bigr) \\
(A_y)_{12,10} &= -\frac{{\mathcal{B}} {D_{s}}}{3 {T}} \\
(A_y)_{12,12} &= \frac{{\mathcal{C}} {D_{s}} {u^{y}}}{2 {T}} \\
(A_y)_{13,1} &= \frac{{D_{s}} \left(3 {\mathcal{B}'(T)} {\nu^{x}} \left(\left(u^{y}\right)^2+1\right)+{\mathcal{B}'(T)} {\nu^{y}} {u^{x}} {u^{y}}+3 {\mathcal{C}'(T)} {\pi_{Q}^{xy}} {u^{y}}\right)}{6 {T}} \\
(A_y)_{13,2} &= \frac{1}{6 {T} \left(u^{\tau}\right)^2}{D_{s}}\Bigl(- 3 {\mathcal{B}} {\nu^{x}} {u^{x}} + 3 {\mathcal{A}} {\nu^{y}} {u^{y}} + 3 {\mathcal{B}} {\nu^{y}} {u^{y}} + 3 {\mathcal{A}} {\nu^{y}} \left(u^{x}\right)^2 {u^{y}} + 6 {\mathcal{B}} {\nu^{y}} \left(u^{x}\right)^2 {u^{y}}\\
  &\quad - 2 {\mathcal{A}} {\nu^{x}} {u^{x}} \left(u^{y}\right)^2 - 9 {\mathcal{B}} {\nu^{x}} {u^{x}} \left(u^{y}\right)^2 + 3 {\mathcal{A}} {\nu^{y}} \left(u^{y}\right)^3 + 3 {\mathcal{B}} {\nu^{y}} \left(u^{y}\right)^3\\
  &\quad + 2 {\mathcal{A}} {\nu^{y}} \left(u^{x}\right)^2 \left(u^{y}\right)^3 + 6 {\mathcal{B}} {\nu^{y}} \left(u^{x}\right)^2 \left(u^{y}\right)^3 - 2 {\mathcal{A}} {\nu^{x}} {u^{x}} \left(u^{y}\right)^4 - 6 {\mathcal{B}} {\nu^{x}} {u^{x}} \left(u^{y}\right)^4\\
  &\quad + 3 {I_{31}} \left(u^{\tau}\right)^2 \left(\left(u^{y}\right)^2+1\right)\\
  &\quad + {\mathcal{C}}\Bigl(3 {\Pi_{Q}} \left(u^{\tau}\right)^2 \left(\left(u^{y}\right)^2+1\right)\\
  &\quad + {\pi_{Q}^{yy}} \left(\left(u^{x}\right)^2 \left(8 \left(u^{y}\right)^2+3\right)+6 \left(u^{y}\right)^2+3\right)\\
  &\quad + {u^{y}} \left(3 {\pi_{Q}{}^{\eta}{}_{\eta}} \left(u^{\tau}\right)^2 {u^{y}}-{\pi_{Q}^{xy}} {u^{x}} \left(8 \left(u^{y}\right)^2+5\right)\right)\Bigr)\Bigr) \\
(A_y)_{13,3} &= \frac{1}{6 {T} \left(u^{\tau}\right)^2 \left(\left(u^{y}\right)^2+1\right)}{D_{s}}\Bigl(\left(u^{y}\right)^2+1\Bigl({I_{31}} \left(u^{\tau}\right)^2 {u^{x}} {u^{y}}\\
  &\quad + 3 {\mathcal{B}} {\nu^{x}} \left(2 \left(u^{x}\right)^2+1\right) {u^{y}} \left(\left(u^{y}\right)^2+1\right)\\
  &\quad - 3 {\mathcal{B}} {\nu^{y}} {u^{x}} \left(\left(u^{x}\right)^2+1\right) \left(2 \left(u^{y}\right)^2+1\right)\\
  &\quad + {\mathcal{A}} {u^{y}}\Bigl(- 2 {\nu^{y}} {u^{x}} \left(\left(u^{x}\right)^2+1\right) {u^{y}} + 3 {\nu^{x}} \left(\left(u^{y}\right)^2+1\right)\\
  &\quad + {\nu^{x}} \left(u^{x}\right)^2 \left(2 \left(u^{y}\right)^2+3\right)\Bigr)\Bigr)\\
  &\quad + {\mathcal{C}}\Bigl({u^{x}} {u^{y}}\Bigl(- 3 {\pi_{Q}{}^{\eta}{}_{\eta}} \left(u^{\tau}\right)^2 \left(u^{y}\right)^2 + {\Pi_{Q}} \left(u^{\tau}\right)^2 \left(\left(u^{y}\right)^2+1\right)\\
  &\quad - {\pi_{Q}^{yy}} \left(\left(u^{x}\right)^2+1\right) \left(8 \left(u^{y}\right)^2+5\right)\Bigr)\\
  &\quad + {\pi_{Q}^{xy}}\Bigl(3 \left(\left(u^{y}\right)^4+3 \left(u^{y}\right)^2+2\right)\\
  &\quad + \left(u^{x}\right)^2 \left(8 \left(u^{y}\right)^4+8 \left(u^{y}\right)^2+6\right)\Bigr)\Bigr)\Bigr) \\
(A_y)_{13,9} &= \frac{{\mathcal{B}} {D_{s}} \left(\left(u^{y}\right)^2+1\right)}{2 {T}} \\
(A_y)_{13,10} &= \frac{{\mathcal{B}} {D_{s}} {u^{x}} {u^{y}}}{6 {T}} \\
(A_y)_{13,13} &= \frac{{\mathcal{C}} {D_{s}} {u^{y}}}{2 {T}} \\
(A_y)_{14,1} &= \frac{{D_{s}} (5 {\mathcal{B}'(T)} {\nu^{y}}+3 {u^{y}} ({\mathcal{C}'(T)} {\Pi_{Q}}+{(\partial_{T} I_{31})}))}{6 {T}} \\
(A_y)_{14,2} &= \frac{{D_{s}} \left({\mathcal{A}} {u^{y}} \left({\nu^{x}} \left(u^{y}\right)^2+{\nu^{x}}-{\nu^{y}} {u^{x}} {u^{y}}\right)+{\mathcal{C}} \left({\pi_{Q}^{xy}} \left(u^{y}\right)^2+{\pi_{Q}^{xy}}-{\pi_{Q}^{yy}} {u^{x}} {u^{y}}\right)\right)}{3 {T} \left(u^{\tau}\right)^2} \\
(A_y)_{14,3} &= \frac{1}{6 {T} \left(u^{\tau}\right)^2}{D_{s}}\Bigl(5 {I_{31}} \left(u^{\tau}\right)^2 + 5 {\mathcal{C}} {\Pi_{Q}} \left(u^{\tau}\right)^2 + 2 {\mathcal{A}} {u^{y}} \left(-{\nu^{x}} {u^{x}} {u^{y}}+{\nu^{y}} \left(u^{x}\right)^2+{\nu^{y}}\right)\\
  &\quad + 2 {\mathcal{C}} \left(-{\pi_{Q}^{xy}} {u^{x}} {u^{y}}+{\pi_{Q}^{yy}} \left(u^{x}\right)^2+{\pi_{Q}^{yy}}\right)\Bigr) \\
(A_y)_{14,8} &= \frac{{(\partial_{\alpha} I_{31})} {D_{s}} {u^{y}}}{2 {T}} \\
(A_y)_{14,10} &= \frac{5 {\mathcal{B}} {D_{s}}}{6 {T}} \\
(A_y)_{14,14} &= \frac{{\mathcal{C}} {D_{s}} {u^{y}}}{2 {T}} \\
S_{1} &= \frac{1}{{\tau} \left(\left(u^{y}\right)^2+1\right)}\Bigl({\pi^{\eta}{}_{\eta}} + {P'(T)} {T} - {\pi^{\eta}{}_{\eta}} \left(u^{x}\right)^2 - {\pi^{yy}} \left(u^{x}\right)^2 + {P'(T)} {T} \left(u^{x}\right)^2 + 2 {\pi^{xy}} {u^{x}} {u^{y}}\\
  &\quad + {\pi^{\eta}{}_{\eta}} \left(u^{y}\right)^2 + {\pi^{yy}} \left(u^{y}\right)^2 + 2 {P'(T)} {T} \left(u^{y}\right)^2 + {P'(T)} {T} \left(u^{x}\right)^2 \left(u^{y}\right)^2\\
  &\quad + {P'(T)} {T} \left(u^{y}\right)^4 + {\Pi} \left(u^{\tau}\right)^2 \left(\left(u^{y}\right)^2+1\right)\\
  &\quad + {h} \left(\left(u^{y}\right)^2+1\right) \left({n} \left(u^{\tau}\right)^2+2 ({\nu^{x}} {u^{x}}+{\nu^{y}} {u^{y}})\right)\Bigr) \\
S_{2} &= \frac{1}{{\tau} {u^{\tau}} \left(\left(u^{y}\right)^2+1\right)}\Bigl(- {\pi^{\eta}{}_{\eta}} {u^{x}} - {\pi^{yy}} {u^{x}} + {P'(T)} {T} {u^{x}} - {\pi^{\eta}{}_{\eta}} \left(u^{x}\right)^3 - {\pi^{yy}} \left(u^{x}\right)^3\\
  &\quad + {P'(T)} {T} \left(u^{x}\right)^3 + {\pi^{xy}} {u^{y}} + 2 {\pi^{xy}} \left(u^{x}\right)^2 {u^{y}} - {\pi^{\eta}{}_{\eta}} {u^{x}} \left(u^{y}\right)^2\\
  &\quad + 2 {P'(T)} {T} {u^{x}} \left(u^{y}\right)^2 + {P'(T)} {T} \left(u^{x}\right)^3 \left(u^{y}\right)^2 + {\pi^{xy}} \left(u^{y}\right)^3\\
  &\quad + {P'(T)} {T} {u^{x}} \left(u^{y}\right)^4 + {\Pi} \left(u^{\tau}\right)^2 {u^{x}} \left(\left(u^{y}\right)^2+1\right)\\
  &\quad + {h} \left(\left(u^{y}\right)^2+1\right) \left({n} \left(u^{\tau}\right)^2 {u^{x}}+{\nu^{x}} \left(2 \left(u^{x}\right)^2+\left(u^{y}\right)^2+1\right)+{\nu^{y}} {u^{x}} {u^{y}}\right)\Bigr) \\
S_{3} &= \frac{1}{{\tau} {u^{\tau}}}\Bigl({\pi^{xy}} {u^{x}} + {u^{y}} \left({P'(T)} {T} \left(u^{\tau}\right)^2+{\Pi} \left(u^{\tau}\right)^2+{\pi^{yy}}\right)\\
  &\quad + {h} \left({n} \left(u^{\tau}\right)^2 {u^{y}}+{\nu^{x}} {u^{x}} {u^{y}}+{\nu^{y}} \left(\left(u^{x}\right)^2+2 \left(u^{y}\right)^2+1\right)\right)\Bigr) \\
S_{4} &= \frac{{\delta_{\pi\pi}} {\pi^{yy}} {u^{\tau}}}{{\tau}}+{\pi^{yy}}-\frac{2 {u^{\tau}} \left(\left(u^{y}\right)^2+1\right) {\eta}}{3 {\tau}} \\
S_{5} &= \frac{3 {\pi^{\eta}{}_{\eta}} ({\delta_{\pi\pi}} {u^{\tau}}+{\tau})+4 {u^{\tau}} {\eta}}{3 {\tau}^3} \\
S_{6} &= \frac{{\delta_{\pi\pi}} {\pi^{xy}} {u^{\tau}}}{{\tau}}+{\pi^{xy}}-\frac{2 {u^{\tau}} {u^{x}} {u^{y}} {\eta}}{3 {\tau}} \\
S_{7} &= {\Pi}+\frac{{u^{\tau}} {\zeta}}{{\tau}} \\
S_{8} &= \frac{{n} \left(u^{\tau}\right)^2+{\nu^{x}} {u^{x}}+{\nu^{y}} {u^{y}}}{{\tau} {u^{\tau}}} \\
S_{9} &= \frac{{D_{s}} {h} {\nu^{x}} {u^{\tau}}}{{T} {\tau}}+{\nu^{x}} \\
S_{10} &= \frac{{D_{s}} {h} {\nu^{y}} {u^{\tau}}}{{T} {\tau}}+{\nu^{y}} \\
S_{11} &= {\pi_{Q}^{yy}}\\
  &\quad - \frac{1}{6 {T} {\tau} {u^{\tau}}}{D_{s}}\Bigl(2 {I_{31}} \left(u^{\tau}\right)^2 \left(\left(u^{y}\right)^2+1\right) + 2 {\mathcal{B}} \left(\left(u^{y}\right)^2+1\right) ({\nu^{x}} {u^{x}}+{\nu^{y}} {u^{y}})\\
  &\quad + {\mathcal{C}} \left(u^{\tau}\right)^2 \left(2 {\Pi_{Q}} \left(\left(u^{y}\right)^2+1\right)-3 {\pi_{Q}^{yy}}+2 {\pi_{Q}{}^{\eta}{}_{\eta}} \left(\left(u^{y}\right)^2+1\right)\right)\Bigr) \\
S_{12} &= \frac{{D_{s}} \left(4 {\mathcal{B}} ({\nu^{x}} {u^{x}}+{\nu^{y}} {u^{y}})+{\mathcal{C}} \left(u^{\tau}\right)^2 (4 {\Pi_{Q}}+7 {\pi_{Q}{}^{\eta}{}_{\eta}})+4 {I_{31}} \left(u^{\tau}\right)^2\right)}{6 {T} {\tau} {u^{\tau}}}+{\pi_{Q}{}^{\eta}{}_{\eta}} \\
S_{13} &= -\frac{{D_{s}} {u^{x}} {u^{y}} \left({\mathcal{B}} ({\nu^{x}} {u^{x}}+{\nu^{y}} {u^{y}})+{I_{31}} \left(u^{\tau}\right)^2\right)}{3 {T} {\tau} {u^{\tau}}}+\frac{{\mathcal{C}} {D_{s}} {u^{\tau}} (3 {\pi_{Q}^{xy}}-2 {u^{x}} {u^{y}} ({\Pi_{Q}}+{\pi_{Q}{}^{\eta}{}_{\eta}}))}{6 {T} {\tau}}+{\pi_{Q}^{xy}} \\
S_{14} &= \frac{{D_{s}} \left(5 {\mathcal{B}} ({\nu^{x}} {u^{x}}+{\nu^{y}} {u^{y}})+{\mathcal{C}} \left(u^{\tau}\right)^2 (5 {\Pi_{Q}}+2 {\pi_{Q}{}^{\eta}{}_{\eta}})+5 {I_{31}} \left(u^{\tau}\right)^2\right)}{6 {T} {\tau} {u^{\tau}}}+{\Pi_{Q}}
\end{align*}
\subsubsection{S5, two-way, conformal}

\texttt{two\_\allowbreak{}d\_\allowbreak{}two\_\allowbreak{}way\_\allowbreak{}matrix\_\allowbreak{}conformal} in \texttt{src/Matrix/2d\_\allowbreak{}two\_\allowbreak{}way.jl}, fields:

\[
\phi = ({T}, {u^{x}}, {u^{y}}, {\pi^{yy}}, {\pi^{\eta}{}_{\eta}}, {\pi^{xy}}, {\alpha}, {\nu^{x}}, {\nu^{y}}, {\pi_{Q}^{yy}}, {\pi_{Q}{}^{\eta}{}_{\eta}}, {\pi_{Q}^{xy}}, {\Pi_{Q}})
\]
\begin{align*}
(A_t)_{1,1} &= - {n} - {(\partial_{T} n)} {T} + {(\partial_{T} n)} {h} \left(u^{\tau}\right)^2 + {h'(T)} {n} \left(u^{\tau}\right)^2 + 2 {h'(T)} ({\nu^{x}} {u^{x}}+{\nu^{y}} {u^{y}})\\
  &\quad + 4 {a_{\mathrm{SB}}} {T}^3 \left(4 \left(u^{x}\right)^2+4 \left(u^{y}\right)^2+3\right) \\
(A_t)_{1,2} &= \frac{2 \left(4 {a_{\mathrm{SB}}} {T}^4 {u^{x}} \left(u^{y}\right)^2+4 {a_{\mathrm{SB}}} {T}^4 {u^{x}}+{h} \left(\left(u^{y}\right)^2+1\right) ({n} {u^{x}}+{\nu^{x}})-{\pi^{\eta}{}_{\eta}} {u^{x}}+{\pi^{xy}} {u^{y}}-{\pi^{yy}} {u^{x}}\right)}{\left(u^{y}\right)^2+1} \\
(A_t)_{1,3} &= \frac{2}{\left(\left(u^{y}\right)^2+1\right)^2}\Bigl({h} \left(\left(u^{y}\right)^2+1\right)^2 ({n} {u^{y}}+{\nu^{y}}) + {\pi^{xy}} \left({u^{x}}-{u^{x}} \left(u^{y}\right)^2\right)\\
  &\quad + {u^{y}} \left(4 {a_{\mathrm{SB}}} {T}^4 \left(\left(u^{y}\right)^2+1\right)^2+{\pi^{\eta}{}_{\eta}} \left(u^{x}\right)^2+{\pi^{yy}} \left(\left(u^{x}\right)^2+1\right)\right)\Bigr) \\
(A_t)_{1,4} &= \frac{\left(u^{y}\right)^2-\left(u^{x}\right)^2}{\left(u^{y}\right)^2+1} \\
(A_t)_{1,5} &= -\frac{\left(u^{x}\right)^2}{\left(u^{y}\right)^2+1} \\
(A_t)_{1,6} &= \frac{2 {u^{x}} {u^{y}}}{\left(u^{y}\right)^2+1} \\
(A_t)_{1,7} &= {(\partial_{\alpha} n)} \left({h} \left(u^{\tau}\right)^2-{T}\right) \\
(A_t)_{1,8} &= 2 {h} {u^{x}} \\
(A_t)_{1,9} &= 2 {h} {u^{y}} \\
(A_t)_{2,1} &= \frac{\left(u^{\tau}\right)^2 {u^{x}} \left(16 {a_{\mathrm{SB}}} {T}^3+{(\partial_{T} n)} {h}\right)+{h'(T)} \left({n} \left(u^{\tau}\right)^2 {u^{x}}+{\nu^{x}} \left(2 \left(u^{x}\right)^2+\left(u^{y}\right)^2+1\right)+{\nu^{y}} {u^{x}} {u^{y}}\right)}{{u^{\tau}}} \\
(A_t)_{2,2} &= \frac{1}{\left(u^{\tau}\right)^3 \left(\left(u^{y}\right)^2+1\right)}\Bigl(- {\pi^{yy}} + 4 {a_{\mathrm{SB}}} {T}^4 - 3 {\pi^{yy}} \left(u^{x}\right)^2 + 12 {a_{\mathrm{SB}}} {T}^4 \left(u^{x}\right)^2 - 2 {\pi^{yy}} \left(u^{x}\right)^4\\
  &\quad + 8 {a_{\mathrm{SB}}} {T}^4 \left(u^{x}\right)^4 + 3 {\pi^{xy}} {u^{x}} {u^{y}} + 2 {\pi^{xy}} \left(u^{x}\right)^3 {u^{y}} - {\pi^{yy}} \left(u^{y}\right)^2\\
  &\quad + 12 {a_{\mathrm{SB}}} {T}^4 \left(u^{y}\right)^2 - 3 {\pi^{yy}} \left(u^{x}\right)^2 \left(u^{y}\right)^2 + 24 {a_{\mathrm{SB}}} {T}^4 \left(u^{x}\right)^2 \left(u^{y}\right)^2\\
  &\quad + 8 {a_{\mathrm{SB}}} {T}^4 \left(u^{x}\right)^4 \left(u^{y}\right)^2 + 3 {\pi^{xy}} {u^{x}} \left(u^{y}\right)^3 + 12 {a_{\mathrm{SB}}} {T}^4 \left(u^{y}\right)^4\\
  &\quad + 12 {a_{\mathrm{SB}}} {T}^4 \left(u^{x}\right)^2 \left(u^{y}\right)^4 + 4 {a_{\mathrm{SB}}} {T}^4 \left(u^{y}\right)^6\\
  &\quad - {\pi^{\eta}{}_{\eta}} \left(2 \left(u^{x}\right)^4+3 \left(u^{x}\right)^2 \left(\left(u^{y}\right)^2+1\right)+\left(\left(u^{y}\right)^2+1\right)^2\right)\\
  &\quad + {h} \left(\left(u^{y}\right)^2+1\right)\Bigl({\nu^{y}} {u^{y}} \left(\left(u^{y}\right)^2+1\right)\\
  &\quad + {\nu^{x}} {u^{x}} \left(2 \left(u^{x}\right)^2+3 \left(u^{y}\right)^2+3\right)\\
  &\quad + {n}\Bigl(2 \left(u^{x}\right)^4+3 \left(u^{x}\right)^2 \left(\left(u^{y}\right)^2+1\right)+\left(\left(u^{y}\right)^2+1\right)^2\Bigr)\Bigr)\Bigr) \\
(A_t)_{2,3} &= \frac{1}{\left(u^{\tau}\right)^3 \left(\left(u^{y}\right)^2+1\right)^2}\Bigl({h} \left(\left(u^{y}\right)^2+1\right)^2\Bigl({n} \left(u^{\tau}\right)^2 {u^{x}} {u^{y}}+{\nu^{x}} {u^{y}} \left(\left(u^{y}\right)^2+1\right)+{\nu^{y}} \left(\left(u^{x}\right)^3+{u^{x}}\right)\Bigr)\\
  &\quad + {\pi^{xy}}\Bigl(- 2 \left(u^{x}\right)^4 \left(\left(u^{y}\right)^2-1\right) + \left(\left(u^{y}\right)^2+1\right)^2\\
  &\quad - 3 \left(u^{x}\right)^2 \left(\left(u^{y}\right)^4-1\right)\Bigr)\\
  &\quad + {u^{x}} {u^{y}}\Bigl(4 {a_{\mathrm{SB}}} {T}^4 \left(u^{\tau}\right)^2 \left(\left(u^{y}\right)^2+1\right)^2\\
  &\quad + {\pi^{yy}} \left(\left(u^{x}\right)^2+1\right) \left(2 \left(u^{x}\right)^2+3 \left(u^{y}\right)^2+3\right)\\
  &\quad + {\pi^{\eta}{}_{\eta}} \left(2 \left(u^{x}\right)^4+3 \left(u^{x}\right)^2 \left(\left(u^{y}\right)^2+1\right)+\left(\left(u^{y}\right)^2+1\right)^2\right)\Bigr)\Bigr) \\
(A_t)_{2,4} &= -\frac{{u^{x}} \left(\left(u^{x}\right)^2+1\right)}{{u^{\tau}} \left(\left(u^{y}\right)^2+1\right)} \\
(A_t)_{2,5} &= -\frac{{u^{\tau}} {u^{x}}}{\left(u^{y}\right)^2+1} \\
(A_t)_{2,6} &= \frac{{u^{y}} \left(2 \left(u^{x}\right)^2+\left(u^{y}\right)^2+1\right)}{{u^{\tau}} \left(\left(u^{y}\right)^2+1\right)} \\
(A_t)_{2,7} &= {(\partial_{\alpha} n)} {h} {u^{\tau}} {u^{x}} \\
(A_t)_{2,8} &= \frac{{h} \left(2 \left(u^{x}\right)^2+\left(u^{y}\right)^2+1\right)}{{u^{\tau}}} \\
(A_t)_{2,9} &= \frac{{h} {u^{x}} {u^{y}}}{{u^{\tau}}} \\
(A_t)_{3,1} &= \frac{\left(u^{\tau}\right)^2 {u^{y}} \left(16 {a_{\mathrm{SB}}} {T}^3+{(\partial_{T} n)} {h}\right)+{h'(T)} \left({n} \left(u^{\tau}\right)^2 {u^{y}}+{\nu^{x}} {u^{x}} {u^{y}}+{\nu^{y}} \left(\left(u^{x}\right)^2+2 \left(u^{y}\right)^2+1\right)\right)}{{u^{\tau}}} \\
(A_t)_{3,2} &= \frac{1}{\left(u^{\tau}\right)^3}\Bigl({u^{x}} {u^{y}} \left(4 {a_{\mathrm{SB}}} {T}^4 \left(u^{\tau}\right)^2-{\pi^{yy}}\right) + {\pi^{xy}} \left(\left(u^{y}\right)^2+1\right)\\
  &\quad + {h} \left({n} \left(u^{\tau}\right)^2 {u^{x}} {u^{y}}+{\nu^{x}} {u^{y}} \left(\left(u^{y}\right)^2+1\right)+{\nu^{y}} \left(\left(u^{x}\right)^3+{u^{x}}\right)\right)\Bigr) \\
(A_t)_{3,3} &= \frac{1}{\left(u^{\tau}\right)^3}\Bigl({\pi^{yy}} + 4 {a_{\mathrm{SB}}} {T}^4 + {\pi^{yy}} \left(u^{x}\right)^2 + 8 {a_{\mathrm{SB}}} {T}^4 \left(u^{x}\right)^2 + 4 {a_{\mathrm{SB}}} {T}^4 \left(u^{x}\right)^4 - {\pi^{xy}} {u^{x}} {u^{y}}\\
  &\quad + 12 {a_{\mathrm{SB}}} {T}^4 \left(u^{y}\right)^2 + 12 {a_{\mathrm{SB}}} {T}^4 \left(u^{x}\right)^2 \left(u^{y}\right)^2 + 8 {a_{\mathrm{SB}}} {T}^4 \left(u^{y}\right)^4\\
  &\quad + {h}\Bigl({\nu^{x}} \left(\left(u^{x}\right)^3+{u^{x}}\right) + {\nu^{y}} {u^{y}} \left(3 \left(u^{x}\right)^2+2 \left(u^{y}\right)^2+3\right)\\
  &\quad + {n} \left(\left(u^{x}\right)^4+\left(u^{x}\right)^2 \left(3 \left(u^{y}\right)^2+2\right)+2 \left(u^{y}\right)^4+3 \left(u^{y}\right)^2+1\right)\Bigr)\Bigr) \\
(A_t)_{3,4} &= \frac{{u^{y}}}{{u^{\tau}}} \\
(A_t)_{3,6} &= \frac{{u^{x}}}{{u^{\tau}}} \\
(A_t)_{3,7} &= {(\partial_{\alpha} n)} {h} {u^{\tau}} {u^{y}} \\
(A_t)_{3,8} &= \frac{{h} {u^{x}} {u^{y}}}{{u^{\tau}}} \\
(A_t)_{3,9} &= \frac{{h} \left(\left(u^{x}\right)^2+2 \left(u^{y}\right)^2+1\right)}{{u^{\tau}}} \\
(A_t)_{4,2} &= \frac{3 {\delta_{\pi\pi}} {\pi^{yy}} {u^{x}}-6 {\pi^{xy}} {\tau_{\pi}} {u^{y}} \left(\left(u^{y}\right)^2+1\right)-2 {u^{x}} \left(-3 {\pi^{yy}} {\tau_{\pi}} \left(u^{y}\right)^2+\left(u^{y}\right)^2 {\eta}+{\eta}\right)}{3 {u^{\tau}}} \\
(A_t)_{4,3} &= \frac{{u^{y}} \left(3 {\delta_{\pi\pi}} {\pi^{yy}}+6 {\pi^{xy}} {\tau_{\pi}} {u^{x}} {u^{y}}-6 {\pi^{yy}} {\tau_{\pi}} \left(\left(u^{x}\right)^2+1\right)+6 \left(u^{x}\right)^2 {\eta}+4 \left(u^{y}\right)^2 {\eta}+4 {\eta}\right)}{3 {u^{\tau}}} \\
(A_t)_{4,4} &= {\tau_{\pi}} {u^{\tau}} \\
(A_t)_{5,2} &= \frac{{u^{x}} (3 {\delta_{\pi\pi}} {\pi^{\eta}{}_{\eta}}-2 {\eta})}{3 {\tau}^2 {u^{\tau}}} \\
(A_t)_{5,3} &= \frac{{u^{y}} (3 {\delta_{\pi\pi}} {\pi^{\eta}{}_{\eta}}-2 {\eta})}{3 {\tau}^2 {u^{\tau}}} \\
(A_t)_{5,5} &= \frac{{\tau_{\pi}} {u^{\tau}}}{{\tau}^2} \\
(A_t)_{6,2} &= \frac{1}{3 {u^{\tau}}}\Bigl(3 {\delta_{\pi\pi}} {\pi^{xy}} {u^{x}} - 3 {\pi^{xy}} {\tau_{\pi}} {u^{x}} \left(2 \left(u^{y}\right)^2+1\right)\\
  &\quad + {u^{y}} \left(3 {\pi^{\eta}{}_{\eta}} {\tau_{\pi}} \left(u^{\tau}\right)^2+3 {\pi^{yy}} \left(2 {\tau_{\pi}} \left(u^{x}\right)^2+{\tau_{\pi}}\right)+{\eta} \left(\left(u^{x}\right)^2+3 \left(u^{y}\right)^2+3\right)\right)\Bigr) \\
(A_t)_{6,3} &= \frac{1}{3 {u^{\tau}} \left(\left(u^{y}\right)^2+1\right)}\Bigl(- 3 {\pi^{xy}} {\tau_{\pi}} {u^{y}} - 3 {\pi^{\eta}{}_{\eta}} {\tau_{\pi}} {u^{x}} \left(u^{y}\right)^2 - 3 {\pi^{\eta}{}_{\eta}} {\tau_{\pi}} \left(u^{x}\right)^3 \left(u^{y}\right)^2 - 3 {\pi^{xy}} {\tau_{\pi}} \left(u^{y}\right)^3\\
  &\quad + 6 {\pi^{xy}} {\tau_{\pi}} \left(u^{x}\right)^2 \left(u^{y}\right)^3 - 3 {\pi^{\eta}{}_{\eta}} {\tau_{\pi}} {u^{x}} \left(u^{y}\right)^4 + 3 {\delta_{\pi\pi}} {\pi^{xy}} {u^{y}} \left(\left(u^{y}\right)^2+1\right)\\
  &\quad - 3 {\pi^{yy}} {\tau_{\pi}} {u^{x}} \left(\left(u^{x}\right)^2+1\right) \left(2 \left(u^{y}\right)^2+1\right) + 3 {u^{x}} {\eta} + 3 \left(u^{x}\right)^3 {\eta}\\
  &\quad + 4 {u^{x}} \left(u^{y}\right)^2 {\eta} + 3 \left(u^{x}\right)^3 \left(u^{y}\right)^2 {\eta} + {u^{x}} \left(u^{y}\right)^4 {\eta}\Bigr) \\
(A_t)_{6,6} &= {\tau_{\pi}} {u^{\tau}} \\
(A_t)_{7,1} &= {(\partial_{T} n)} {u^{\tau}} \\
(A_t)_{7,2} &= \frac{{n} \left(u^{\tau}\right)^2 {u^{x}}+{\nu^{x}} \left(\left(u^{y}\right)^2+1\right)-{\nu^{y}} {u^{x}} {u^{y}}}{\left(u^{\tau}\right)^3} \\
(A_t)_{7,3} &= \frac{{n} \left(u^{\tau}\right)^2 {u^{y}}-{\nu^{x}} {u^{x}} {u^{y}}+{\nu^{y}} \left(\left(u^{x}\right)^2+1\right)}{\left(u^{\tau}\right)^3} \\
(A_t)_{7,7} &= {(\partial_{\alpha} n)} {u^{\tau}} \\
(A_t)_{7,8} &= \frac{{u^{x}}}{{u^{\tau}}} \\
(A_t)_{7,9} &= \frac{{u^{y}}}{{u^{\tau}}} \\
(A_t)_{8,1} &= \frac{{D_{s}} {u^{\tau}} ({h'(T)} {\nu^{x}}+{u^{x}} ({(\partial_{T} n)} {T}+{n}))}{{T}} \\
(A_t)_{8,2} &= \frac{{D_{s}} {h} \left({n} \left(u^{\tau}\right)^2+{\nu^{x}} \left({u^{x}}-{u^{x}} \left(u^{y}\right)^2\right)+{\nu^{y}} \left(\left(u^{x}\right)^2+1\right) {u^{y}}\right)}{{T} {u^{\tau}}} \\
(A_t)_{8,3} &= -\frac{{D_{s}} {h} \left(\left(u^{x}\right)^2+1\right) ({\nu^{y}} {u^{x}}-{\nu^{x}} {u^{y}})}{{T} {u^{\tau}}} \\
(A_t)_{8,7} &= {(\partial_{\alpha} n)} {D_{s}} {u^{\tau}} {u^{x}} \\
(A_t)_{8,8} &= \frac{{D_{s}} {h} {u^{\tau}}}{{T}} \\
(A_t)_{9,1} &= \frac{{D_{s}} {u^{\tau}} ({h'(T)} {\nu^{y}}+{u^{y}} ({(\partial_{T} n)} {T}+{n}))}{{T}} \\
(A_t)_{9,2} &= \frac{{D_{s}} {h} \left(\left(u^{y}\right)^2+1\right) ({\nu^{y}} {u^{x}}-{\nu^{x}} {u^{y}})}{{T} {u^{\tau}}} \\
(A_t)_{9,3} &= \frac{{D_{s}} {h} \left({n} \left(u^{\tau}\right)^2+{\nu^{x}} {u^{x}} \left(\left(u^{y}\right)^2+1\right)-{\nu^{y}} \left(\left(u^{x}\right)^2-1\right) {u^{y}}\right)}{{T} {u^{\tau}}} \\
(A_t)_{9,7} &= {(\partial_{\alpha} n)} {D_{s}} {u^{\tau}} {u^{y}} \\
(A_t)_{9,9} &= \frac{{D_{s}} {h} {u^{\tau}}}{{T}} \\
(A_t)_{10,1} &= \frac{1}{6 {T} {u^{\tau}}}{D_{s}}\Bigl(3 {\mathcal{C}'(T)} {\pi_{Q}^{yy}} \left(u^{\tau}\right)^2 - 2 {\mathcal{B}'(T)} {\nu^{x}} {u^{x}} \left(\left(u^{y}\right)^2+1\right)\\
  &\quad + 2 {\mathcal{B}'(T)} {\nu^{y}} {u^{y}} \left(3 \left(u^{x}\right)^2+2 \left(u^{y}\right)^2+2\right)\Bigr) \\
(A_t)_{10,2} &= -\frac{1}{6 {T} \left(u^{\tau}\right)^3}{D_{s}}\Bigl(2 {\mathcal{A}} \left(u^{\tau}\right)^2 \left(\left(u^{y}\right)^2+1\right) \left({\nu^{x}} \left(u^{y}\right)^2+{\nu^{x}}-{\nu^{y}} {u^{x}} {u^{y}}\right)\\
  &\quad + 2 {\mathcal{B}} \left(\left(3 \left(u^{x}\right)^2+4\right) \left(u^{y}\right)^2+3 \left(u^{y}\right)^4+1\right) \left({\nu^{x}} \left(u^{y}\right)^2+{\nu^{x}}-{\nu^{y}} {u^{x}} {u^{y}}\right)\\
  &\quad + \left(u^{\tau}\right)^2\Bigl(2 {I_{31}} {u^{x}} \left(\left(u^{y}\right)^2+1\right)\\
  &\quad + {\mathcal{C}}\Bigl(2 {\Pi_{Q}} {u^{x}} \left(\left(u^{y}\right)^2+1\right)\\
  &\quad - 2 \left(\left(u^{y}\right)^2+1\right) ({\pi_{Q}{}^{\eta}{}_{\eta}} {u^{x}}-4 {\pi_{Q}^{xy}} {u^{y}})\\
  &\quad - {\pi_{Q}^{yy}} {u^{x}} \left(8 \left(u^{y}\right)^2+5\right)\Bigr)\Bigr)\Bigr) \\
(A_t)_{10,3} &= \frac{1}{6 {T} \left(u^{\tau}\right)^3}{D_{s}}\Bigl(2 {\mathcal{B}} {\nu^{x}} {u^{x}} {u^{y}} \left(\left(3 \left(u^{x}\right)^2+4\right) \left(u^{y}\right)^2+3 \left(u^{y}\right)^4+1\right)\\
  &\quad - 2{\mathcal{A}} \left(u^{\tau}\right)^2\Bigl({\nu^{y}} \left(u^{x}\right)^2 \left(\left(u^{y}\right)^2-2\right) - 2 {\nu^{y}} \left(\left(u^{y}\right)^2+1\right)\\
  &\quad - {\nu^{x}} {u^{x}} {u^{y}} \left(\left(u^{y}\right)^2+1\right)\Bigr)\\
  &\quad - 2{\mathcal{B}} {\nu^{y}}\Bigl(3 \left(u^{x}\right)^4 \left(\left(u^{y}\right)^2-1\right) - 2 \left(\left(u^{y}\right)^2+1\right)\\
  &\quad + \left(u^{x}\right)^2 \left(3 \left(u^{y}\right)^4+\left(u^{y}\right)^2-5\right)\Bigr)\\
  &\quad + \left(u^{\tau}\right)^2\Bigl(2 {I_{31}} {u^{y}} \left(3 \left(u^{x}\right)^2+2 \left(u^{y}\right)^2+2\right) + 4 {\mathcal{C}} {\pi_{Q}^{xy}} \left(2 {u^{x}} \left(u^{y}\right)^2+{u^{x}}\right)\\
  &\quad + {\mathcal{C}} {u^{y}}\Bigl({\pi_{Q}^{yy}} - 8 {\pi_{Q}^{yy}} \left(u^{x}\right)^2 - 2 {\pi_{Q}{}^{\eta}{}_{\eta}} \left(u^{x}\right)^2\\
  &\quad + {\Pi_{Q}} \left(6 \left(u^{x}\right)^2+4 \left(u^{y}\right)^2+4\right)\Bigr)\Bigr)\Bigr) \\
(A_t)_{10,8} &= -\frac{{\mathcal{B}} {D_{s}} {u^{x}} \left(\left(u^{y}\right)^2+1\right)}{3 {T} {u^{\tau}}} \\
(A_t)_{10,9} &= \frac{{\mathcal{B}} {D_{s}} {u^{y}} \left(3 \left(u^{x}\right)^2+2 \left(u^{y}\right)^2+2\right)}{3 {T} {u^{\tau}}} \\
(A_t)_{10,10} &= \frac{{\mathcal{C}} {D_{s}} {u^{\tau}}}{2 {T}} \\
(A_t)_{11,1} &= \frac{{D_{s}} \left(3 {\mathcal{C}'(T)} {\pi_{Q}{}^{\eta}{}_{\eta}} \left(u^{\tau}\right)^2-2 {\mathcal{B}'(T)} ({\nu^{x}} {u^{x}}+{\nu^{y}} {u^{y}})\right)}{6 {T} {u^{\tau}}} \\
(A_t)_{11,2} &= -\frac{1}{6 {T} \left(u^{\tau}\right)^3}{D_{s}}\Bigl(2 {\mathcal{B}} \left({\nu^{x}} \left(u^{y}\right)^2+{\nu^{x}}-{\nu^{y}} {u^{x}} {u^{y}}\right) + 2 {\mathcal{A}} \left(u^{\tau}\right)^2 \left({\nu^{x}} \left(u^{y}\right)^2+{\nu^{x}}-{\nu^{y}} {u^{x}} {u^{y}}\right)\\
  &\quad + \left(u^{\tau}\right)^2 ({\mathcal{C}} (2 {\Pi_{Q}} {u^{x}}+2 {\pi_{Q}^{xy}} {u^{y}}-2 {\pi_{Q}^{yy}} {u^{x}}-5 {\pi_{Q}{}^{\eta}{}_{\eta}} {u^{x}})+2 {I_{31}} {u^{x}})\Bigr) \\
(A_t)_{11,3} &= -\frac{1}{6 {T} \left(u^{\tau}\right)^3 \left(\left(u^{y}\right)^2+1\right)}{D_{s}}\Bigl(2 {\mathcal{B}} \left(\left(u^{y}\right)^2+1\right) \left(-{\nu^{x}} {u^{x}} {u^{y}}+{\nu^{y}} \left(u^{x}\right)^2+{\nu^{y}}\right)\\
  &\quad + 2 {\mathcal{A}} \left(u^{\tau}\right)^2 \left(\left(u^{y}\right)^2+1\right) \left(-{\nu^{x}} {u^{x}} {u^{y}}+{\nu^{y}} \left(u^{x}\right)^2+{\nu^{y}}\right)\\
  &\quad + \left(u^{\tau}\right)^2\Bigl(- 2 {\mathcal{C}} {\pi_{Q}^{xy}} {u^{x}} \left(\left(u^{y}\right)^2-1\right) + 2 {I_{31}} {u^{y}} \left(\left(u^{y}\right)^2+1\right)\\
  &\quad + {\mathcal{C}} {u^{y}}\Bigl(2 {\pi_{Q}^{yy}} \left(\left(u^{x}\right)^2+1\right)\\
  &\quad + {\pi_{Q}{}^{\eta}{}_{\eta}} \left(2 \left(u^{x}\right)^2-3 \left(u^{y}\right)^2-3\right) + 2 {\Pi_{Q}} \left(\left(u^{y}\right)^2+1\right)\Bigr)\Bigr)\Bigr) \\
(A_t)_{11,8} &= -\frac{{\mathcal{B}} {D_{s}} {u^{x}}}{3 {T} {u^{\tau}}} \\
(A_t)_{11,9} &= -\frac{{\mathcal{B}} {D_{s}} {u^{y}}}{3 {T} {u^{\tau}}} \\
(A_t)_{11,11} &= \frac{{\mathcal{C}} {D_{s}} {u^{\tau}}}{2 {T}} \\
(A_t)_{12,1} &= \frac{1}{6 {T} {u^{\tau}}}{D_{s}}\Bigl(3 {\mathcal{C}'(T)} {\pi_{Q}^{xy}} \left(u^{\tau}\right)^2 + {\mathcal{B}'(T)} {\nu^{y}} {u^{x}} \left(3 \left(u^{x}\right)^2+\left(u^{y}\right)^2+3\right)\\
  &\quad + {\mathcal{B}'(T)} {\nu^{x}} {u^{y}} \left(\left(u^{x}\right)^2+3 \left(u^{y}\right)^2+3\right)\Bigr) \\
(A_t)_{12,2} &= \frac{1}{6 {T} \left(u^{\tau}\right)^3}{D_{s}}\Bigl({\mathcal{A}} \left(u^{\tau}\right)^2\Bigl(3 {\nu^{y}} \left(\left(u^{y}\right)^2+1\right) - 2 {\nu^{x}} {u^{x}} {u^{y}} \left(\left(u^{y}\right)^2+1\right)\\
  &\quad + {\nu^{y}} \left(u^{x}\right)^2 \left(2 \left(u^{y}\right)^2+3\right)\Bigr)\\
  &\quad + \left(u^{\tau}\right)^2\Bigl({\mathcal{C}} {\pi_{Q}^{xy}} {u^{x}} \left(3-8 \left(u^{y}\right)^2\right) + {I_{31}} {u^{y}} \left(\left(u^{x}\right)^2+3 \left(u^{y}\right)^2+3\right)\\
  &\quad + {\mathcal{C}} {u^{y}}\Bigl({\pi_{Q}^{yy}} \left(8 \left(u^{x}\right)^2+6\right) + {\Pi_{Q}} \left(\left(u^{x}\right)^2+3 \left(u^{y}\right)^2+3\right)\\
  &\quad + {\pi_{Q}{}^{\eta}{}_{\eta}} \left(5 \left(u^{x}\right)^2+3 \left(u^{y}\right)^2+3\right)\Bigr)\Bigr)\\
  &\quad + {\mathcal{B}}\Bigl(- 2 {\nu^{x}} {u^{x}} {u^{y}} \left(\left(u^{y}\right)^2+1\right) \left(3 \left(u^{x}\right)^2+3 \left(u^{y}\right)^2+4\right)\\
  &\quad + {\nu^{y}}\Bigl(3 \left(\left(u^{y}\right)^2+1\right)^2 + \left(u^{x}\right)^4 \left(6 \left(u^{y}\right)^2+3\right)\\
  &\quad + 2 \left(u^{x}\right)^2 \left(3 \left(u^{y}\right)^4+7 \left(u^{y}\right)^2+3\right)\Bigr)\Bigr)\Bigr) \\
(A_t)_{12,3} &= \frac{1}{6 {T} \left(u^{\tau}\right)^3 \left(\left(u^{y}\right)^2+1\right)}{D_{s}}\Bigl({\mathcal{A}} \left(u^{\tau}\right)^2 \left(\left(u^{y}\right)^2+1\right)\Bigl(- 2 {\nu^{y}} {u^{x}} \left(\left(u^{x}\right)^2+1\right) {u^{y}} + 3 {\nu^{x}} \left(\left(u^{y}\right)^2+1\right)\\
  &\quad + {\nu^{x}} \left(u^{x}\right)^2 \left(2 \left(u^{y}\right)^2+3\right)\Bigr)\\
  &\quad + \left(u^{\tau}\right)^2\Bigl({I_{31}} {u^{x}} \left(\left(u^{y}\right)^2+1\right) \left(3 \left(u^{x}\right)^2+\left(u^{y}\right)^2+3\right)\\
  &\quad + {\mathcal{C}}\Bigl(- 3 {\pi_{Q}{}^{\eta}{}_{\eta}} {u^{x}} - 3 {\pi_{Q}{}^{\eta}{}_{\eta}} \left(u^{x}\right)^3 + 3 {\pi_{Q}^{xy}} {u^{y}} + 4 {\pi_{Q}^{xy}} \left(u^{x}\right)^2 {u^{y}}\\
  &\quad - 6 {\pi_{Q}{}^{\eta}{}_{\eta}} {u^{x}} \left(u^{y}\right)^2 - 5 {\pi_{Q}{}^{\eta}{}_{\eta}} \left(u^{x}\right)^3 \left(u^{y}\right)^2 + 3 {\pi_{Q}^{xy}} \left(u^{y}\right)^3\\
  &\quad + 8 {\pi_{Q}^{xy}} \left(u^{x}\right)^2 \left(u^{y}\right)^3 - 3 {\pi_{Q}{}^{\eta}{}_{\eta}} {u^{x}} \left(u^{y}\right)^4\\
  &\quad + {\Pi_{Q}} {u^{x}} \left(\left(u^{y}\right)^2+1\right) \left(3 \left(u^{x}\right)^2+\left(u^{y}\right)^2+3\right)\\
  &\quad - 2 {\pi_{Q}^{yy}} {u^{x}} \left(\left(u^{x}\right)^2+1\right) \left(4 \left(u^{y}\right)^2+3\right)\Bigr)\Bigr)\\
  &\quad + {\mathcal{B}} \left(\left(u^{y}\right)^2+1\right)\Bigl(- 2 {\nu^{y}} {u^{x}} \left(\left(u^{x}\right)^2+1\right) {u^{y}} \left(3 \left(u^{x}\right)^2+3 \left(u^{y}\right)^2+4\right)\\
  &\quad + {\nu^{x}}\Bigl(3 \left(\left(u^{y}\right)^2+1\right)^2 + \left(u^{x}\right)^4 \left(6 \left(u^{y}\right)^2+3\right)\\
  &\quad + 2 \left(u^{x}\right)^2 \left(3 \left(u^{y}\right)^4+7 \left(u^{y}\right)^2+3\right)\Bigr)\Bigr)\Bigr) \\
(A_t)_{12,8} &= \frac{{\mathcal{B}} {D_{s}} {u^{y}} \left(\left(u^{x}\right)^2+3 \left(u^{y}\right)^2+3\right)}{6 {T} {u^{\tau}}} \\
(A_t)_{12,9} &= \frac{{\mathcal{B}} {D_{s}} {u^{x}} \left(3 \left(u^{x}\right)^2+\left(u^{y}\right)^2+3\right)}{6 {T} {u^{\tau}}} \\
(A_t)_{12,12} &= \frac{{\mathcal{C}} {D_{s}} {u^{\tau}}}{2 {T}} \\
(A_t)_{13,1} &= \frac{{D_{s}} \left(5 {\mathcal{B}'(T)} ({\nu^{x}} {u^{x}}+{\nu^{y}} {u^{y}})+3 {\mathcal{C}'(T)} {\Pi_{Q}} \left(u^{\tau}\right)^2+3 {(\partial_{T} I_{31})} \left(u^{\tau}\right)^2\right)}{6 {T} {u^{\tau}}} \\
(A_t)_{13,2} &= \frac{1}{6 {T} \left(u^{\tau}\right)^3}{D_{s}}\Bigl(5 {\mathcal{B}} \left({\nu^{x}} \left(u^{y}\right)^2+{\nu^{x}}-{\nu^{y}} {u^{x}} {u^{y}}\right) + 2 {\mathcal{A}} \left(u^{\tau}\right)^2 \left({\nu^{x}} \left(u^{y}\right)^2+{\nu^{x}}-{\nu^{y}} {u^{x}} {u^{y}}\right)\\
  &\quad + \left(u^{\tau}\right)^2 ({\mathcal{C}} (5 {\Pi_{Q}} {u^{x}}+2 {\pi_{Q}^{xy}} {u^{y}}-2 {\pi_{Q}^{yy}} {u^{x}}-2 {\pi_{Q}{}^{\eta}{}_{\eta}} {u^{x}})+5 {I_{31}} {u^{x}})\Bigr) \\
(A_t)_{13,3} &= \frac{1}{6 {T} \left(u^{\tau}\right)^3 \left(\left(u^{y}\right)^2+1\right)}{D_{s}}\Bigl(5 {\mathcal{B}} \left(\left(u^{y}\right)^2+1\right) \left(-{\nu^{x}} {u^{x}} {u^{y}}+{\nu^{y}} \left(u^{x}\right)^2+{\nu^{y}}\right)\\
  &\quad + 2 {\mathcal{A}} \left(u^{\tau}\right)^2 \left(\left(u^{y}\right)^2+1\right) \left(-{\nu^{x}} {u^{x}} {u^{y}}+{\nu^{y}} \left(u^{x}\right)^2+{\nu^{y}}\right)\\
  &\quad + \left(u^{\tau}\right)^2\Bigl(5 {I_{31}} {u^{y}} \left(\left(u^{y}\right)^2+1\right)\\
  &\quad + {\mathcal{C}}\Bigl(2 {u^{y}} \left({\pi_{Q}^{yy}} \left(u^{x}\right)^2+{\pi_{Q}^{yy}}+{\pi_{Q}{}^{\eta}{}_{\eta}} \left(u^{x}\right)^2\right)\\
  &\quad - 2 {\pi_{Q}^{xy}} {u^{x}} \left(\left(u^{y}\right)^2-1\right) + 5 {\Pi_{Q}} {u^{y}} \left(\left(u^{y}\right)^2+1\right)\Bigr)\Bigr)\Bigr) \\
(A_t)_{13,7} &= \frac{{(\partial_{\alpha} I_{31})} {D_{s}} {u^{\tau}}}{2 {T}} \\
(A_t)_{13,8} &= \frac{5 {\mathcal{B}} {D_{s}} {u^{x}}}{6 {T} {u^{\tau}}} \\
(A_t)_{13,9} &= \frac{5 {\mathcal{B}} {D_{s}} {u^{y}}}{6 {T} {u^{\tau}}} \\
(A_t)_{13,13} &= \frac{{\mathcal{C}} {D_{s}} {u^{\tau}}}{2 {T}} \\
(A_x)_{1,1} &= \frac{\left(u^{\tau}\right)^2 {u^{x}} \left(16 {a_{\mathrm{SB}}} {T}^3+{(\partial_{T} n)} {h}\right)+{h'(T)} \left({n} \left(u^{\tau}\right)^2 {u^{x}}+{\nu^{x}} \left(2 \left(u^{x}\right)^2+\left(u^{y}\right)^2+1\right)+{\nu^{y}} {u^{x}} {u^{y}}\right)}{{u^{\tau}}} \\
(A_x)_{1,2} &= \frac{1}{\left(u^{\tau}\right)^3 \left(\left(u^{y}\right)^2+1\right)}\Bigl(- {\pi^{yy}} + 4 {a_{\mathrm{SB}}} {T}^4 - 3 {\pi^{yy}} \left(u^{x}\right)^2 + 12 {a_{\mathrm{SB}}} {T}^4 \left(u^{x}\right)^2 - 2 {\pi^{yy}} \left(u^{x}\right)^4\\
  &\quad + 8 {a_{\mathrm{SB}}} {T}^4 \left(u^{x}\right)^4 + 3 {\pi^{xy}} {u^{x}} {u^{y}} + 2 {\pi^{xy}} \left(u^{x}\right)^3 {u^{y}} - {\pi^{yy}} \left(u^{y}\right)^2\\
  &\quad + 12 {a_{\mathrm{SB}}} {T}^4 \left(u^{y}\right)^2 - 3 {\pi^{yy}} \left(u^{x}\right)^2 \left(u^{y}\right)^2 + 24 {a_{\mathrm{SB}}} {T}^4 \left(u^{x}\right)^2 \left(u^{y}\right)^2\\
  &\quad + 8 {a_{\mathrm{SB}}} {T}^4 \left(u^{x}\right)^4 \left(u^{y}\right)^2 + 3 {\pi^{xy}} {u^{x}} \left(u^{y}\right)^3 + 12 {a_{\mathrm{SB}}} {T}^4 \left(u^{y}\right)^4\\
  &\quad + 12 {a_{\mathrm{SB}}} {T}^4 \left(u^{x}\right)^2 \left(u^{y}\right)^4 + 4 {a_{\mathrm{SB}}} {T}^4 \left(u^{y}\right)^6\\
  &\quad - {\pi^{\eta}{}_{\eta}} \left(2 \left(u^{x}\right)^4+3 \left(u^{x}\right)^2 \left(\left(u^{y}\right)^2+1\right)+\left(\left(u^{y}\right)^2+1\right)^2\right)\\
  &\quad + {h} \left(\left(u^{y}\right)^2+1\right)\Bigl({\nu^{y}} {u^{y}} \left(\left(u^{y}\right)^2+1\right)\\
  &\quad + {\nu^{x}} {u^{x}} \left(2 \left(u^{x}\right)^2+3 \left(u^{y}\right)^2+3\right)\\
  &\quad + {n}\Bigl(2 \left(u^{x}\right)^4+3 \left(u^{x}\right)^2 \left(\left(u^{y}\right)^2+1\right)+\left(\left(u^{y}\right)^2+1\right)^2\Bigr)\Bigr)\Bigr) \\
(A_x)_{1,3} &= \frac{1}{\left(u^{\tau}\right)^3 \left(\left(u^{y}\right)^2+1\right)^2}\Bigl({h} \left(\left(u^{y}\right)^2+1\right)^2\Bigl({n} \left(u^{\tau}\right)^2 {u^{x}} {u^{y}}+{\nu^{x}} {u^{y}} \left(\left(u^{y}\right)^2+1\right)+{\nu^{y}} \left(\left(u^{x}\right)^3+{u^{x}}\right)\Bigr)\\
  &\quad + {\pi^{xy}}\Bigl(- 2 \left(u^{x}\right)^4 \left(\left(u^{y}\right)^2-1\right) + \left(\left(u^{y}\right)^2+1\right)^2\\
  &\quad - 3 \left(u^{x}\right)^2 \left(\left(u^{y}\right)^4-1\right)\Bigr)\\
  &\quad + {u^{x}} {u^{y}}\Bigl(4 {a_{\mathrm{SB}}} {T}^4 \left(u^{\tau}\right)^2 \left(\left(u^{y}\right)^2+1\right)^2\\
  &\quad + {\pi^{yy}} \left(\left(u^{x}\right)^2+1\right) \left(2 \left(u^{x}\right)^2+3 \left(u^{y}\right)^2+3\right)\\
  &\quad + {\pi^{\eta}{}_{\eta}} \left(2 \left(u^{x}\right)^4+3 \left(u^{x}\right)^2 \left(\left(u^{y}\right)^2+1\right)+\left(\left(u^{y}\right)^2+1\right)^2\right)\Bigr)\Bigr) \\
(A_x)_{1,4} &= -\frac{{u^{x}} \left(\left(u^{x}\right)^2+1\right)}{{u^{\tau}} \left(\left(u^{y}\right)^2+1\right)} \\
(A_x)_{1,5} &= -\frac{{u^{\tau}} {u^{x}}}{\left(u^{y}\right)^2+1} \\
(A_x)_{1,6} &= \frac{{u^{y}} \left(2 \left(u^{x}\right)^2+\left(u^{y}\right)^2+1\right)}{{u^{\tau}} \left(\left(u^{y}\right)^2+1\right)} \\
(A_x)_{1,7} &= {(\partial_{\alpha} n)} {h} {u^{\tau}} {u^{x}} \\
(A_x)_{1,8} &= \frac{{h} \left(2 \left(u^{x}\right)^2+\left(u^{y}\right)^2+1\right)}{{u^{\tau}}} \\
(A_x)_{1,9} &= \frac{{h} {u^{x}} {u^{y}}}{{u^{\tau}}} \\
(A_x)_{2,1} &= 16 {a_{\mathrm{SB}}} {T}^3 \left(u^{x}\right)^2+4 {a_{\mathrm{SB}}} {T}^3+{h'(T)} {n} \left(u^{x}\right)^2+2 {h'(T)} {\nu^{x}} {u^{x}}+{(\partial_{T} n)} {h} \left(u^{x}\right)^2+{(\partial_{T} n)} {T}+{n} \\
(A_x)_{2,2} &= 2 \left(4 {a_{\mathrm{SB}}} {T}^4 {u^{x}}+{h} {n} {u^{x}}+{h} {\nu^{x}}-\frac{{\pi^{\eta}{}_{\eta}} {u^{x}}-{\pi^{xy}} {u^{y}}+{\pi^{yy}} {u^{x}}}{\left(u^{y}\right)^2+1}\right) \\
(A_x)_{2,3} &= \frac{2 \left({u^{y}} \left({\pi^{\eta}{}_{\eta}} \left(u^{x}\right)^2+{\pi^{yy}} \left(u^{x}\right)^2+{\pi^{yy}}\right)+{\pi^{xy}} \left({u^{x}}-{u^{x}} \left(u^{y}\right)^2\right)\right)}{\left(\left(u^{y}\right)^2+1\right)^2} \\
(A_x)_{2,4} &= -\frac{\left(u^{x}\right)^2+1}{\left(u^{y}\right)^2+1} \\
(A_x)_{2,5} &= -\frac{\left(u^{\tau}\right)^2}{\left(u^{y}\right)^2+1} \\
(A_x)_{2,6} &= \frac{2 {u^{x}} {u^{y}}}{\left(u^{y}\right)^2+1} \\
(A_x)_{2,7} &= {(\partial_{\alpha} n)} \left({h} \left(u^{x}\right)^2+{T}\right) \\
(A_x)_{2,8} &= 2 {h} {u^{x}} \\
(A_x)_{3,1} &= {u^{x}} {u^{y}} \left(16 {a_{\mathrm{SB}}} {T}^3+{(\partial_{T} n)} {h}\right)+{h'(T)} {u^{y}} ({n} {u^{x}}+{\nu^{x}})+{h'(T)} {\nu^{y}} {u^{x}} \\
(A_x)_{3,2} &= 4 {a_{\mathrm{SB}}} {T}^4 {u^{y}}+{h} ({n} {u^{y}}+{\nu^{y}}) \\
(A_x)_{3,3} &= 4 {a_{\mathrm{SB}}} {T}^4 {u^{x}}+{h} ({n} {u^{x}}+{\nu^{x}}) \\
(A_x)_{3,6} &= 1 \\
(A_x)_{3,7} &= {(\partial_{\alpha} n)} {h} {u^{x}} {u^{y}} \\
(A_x)_{3,8} &= {h} {u^{y}} \\
(A_x)_{3,9} &= {h} {u^{x}} \\
(A_x)_{4,2} &= {\delta_{\pi\pi}} {\pi^{yy}}-\frac{2 {\pi^{xy}} {\tau_{\pi}} {u^{x}} \left(\left(u^{y}\right)^2+1\right) {u^{y}}}{\left(u^{\tau}\right)^2}+\frac{2 {\pi^{yy}} {\tau_{\pi}} \left(u^{x}\right)^2 \left(u^{y}\right)^2}{\left(u^{\tau}\right)^2}-\frac{2}{3} \left(\left(u^{y}\right)^2+1\right) {\eta} \\
(A_x)_{4,3} &= \frac{2 {u^{x}} {u^{y}} \left({\pi^{xy}} {\tau_{\pi}} {u^{x}} {u^{y}}-{\pi^{yy}} {\tau_{\pi}} \left(\left(u^{x}\right)^2+1\right)+\left(u^{\tau}\right)^2 {\eta}\right)}{\left(u^{\tau}\right)^2} \\
(A_x)_{4,4} &= {\tau_{\pi}} {u^{x}} \\
(A_x)_{5,2} &= \frac{3 {\delta_{\pi\pi}} {\pi^{\eta}{}_{\eta}}-2 {\eta}}{3 {\tau}^2} \\
(A_x)_{5,5} &= \frac{{\tau_{\pi}} {u^{x}}}{{\tau}^2} \\
(A_x)_{6,2} &= {\delta_{\pi\pi}} {\pi^{xy}}\\
  &\quad + \frac{{u^{x}} \left({u^{y}} \left(3 {\pi^{\eta}{}_{\eta}} {\tau_{\pi}} \left(u^{\tau}\right)^2+3 {\pi^{yy}} \left(2 {\tau_{\pi}} \left(u^{x}\right)^2+{\tau_{\pi}}\right)+\left(u^{\tau}\right)^2 {\eta}\right)-3 {\pi^{xy}} {\tau_{\pi}} {u^{x}} \left(2 \left(u^{y}\right)^2+1\right)\right)}{3 \left(u^{\tau}\right)^2} \\
(A_x)_{6,3} &= \frac{1}{\left(u^{\tau}\right)^2 \left(\left(u^{y}\right)^2+1\right)}\Bigl(- {\pi^{yy}} {\tau_{\pi}} \left(u^{x}\right)^2 \left(\left(u^{x}\right)^2+1\right) \left(2 \left(u^{y}\right)^2+1\right)\\
  &\quad + {\pi^{xy}} {\tau_{\pi}} {u^{x}} {u^{y}} \left(\left(2 \left(u^{x}\right)^2-1\right) \left(u^{y}\right)^2-1\right)\\
  &\quad - \left(u^{\tau}\right)^2 \left({\pi^{\eta}{}_{\eta}} {\tau_{\pi}} \left(u^{x}\right)^2 \left(u^{y}\right)^2-\left(\left(u^{x}\right)^2+1\right) \left(\left(u^{y}\right)^2+1\right) {\eta}\right)\Bigr) \\
(A_x)_{6,6} &= {\tau_{\pi}} {u^{x}} \\
(A_x)_{7,1} &= {(\partial_{T} n)} {u^{x}} \\
(A_x)_{7,2} &= {n} \\
(A_x)_{7,7} &= {(\partial_{\alpha} n)} {u^{x}} \\
(A_x)_{7,8} &= 1 \\
(A_x)_{8,1} &= \frac{{D_{s}} \left({h'(T)} {\nu^{x}} {u^{x}}+{(\partial_{T} n)} {T} \left(\left(u^{x}\right)^2+1\right)+{n} \left(\left(u^{x}\right)^2+1\right)\right)}{{T}} \\
(A_x)_{8,2} &= \frac{{D_{s}} {h} \left({n} \left(u^{\tau}\right)^2 {u^{x}}+{\nu^{x}} \left(\left(u^{x}\right)^2 \left(-\left(u^{y}\right)^2\right)+\left(u^{x}\right)^2+2 \left(u^{y}\right)^2+2\right)+{\nu^{y}} \left(u^{x}\right)^3 {u^{y}}\right)}{{T} \left(u^{\tau}\right)^2} \\
(A_x)_{8,3} &= -\frac{{D_{s}} {h} \left(u^{x}\right)^2 \left(-{\nu^{x}} {u^{x}} {u^{y}}+{\nu^{y}} \left(u^{x}\right)^2+{\nu^{y}}\right)}{{T} \left(u^{\tau}\right)^2} \\
(A_x)_{8,7} &= {(\partial_{\alpha} n)} {D_{s}} \left(\left(u^{x}\right)^2+1\right) \\
(A_x)_{8,8} &= \frac{{D_{s}} {h} {u^{x}}}{{T}} \\
(A_x)_{9,1} &= \frac{{D_{s}} {u^{x}} ({h'(T)} {\nu^{y}}+{u^{y}} ({(\partial_{T} n)} {T}+{n}))}{{T}} \\
(A_x)_{9,2} &= \frac{{D_{s}} {h} \left(\left(u^{y}\right)^2+1\right) \left(-{\nu^{x}} {u^{x}} {u^{y}}+{\nu^{y}} \left(u^{x}\right)^2+{\nu^{y}}\right)}{{T} \left(u^{\tau}\right)^2} \\
(A_x)_{9,3} &= \frac{{D_{s}} {h} \left({n} \left(u^{\tau}\right)^2 {u^{x}}+{\nu^{x}} \left(\left(u^{x}\right)^2+1\right) \left(\left(u^{y}\right)^2+1\right)-{\nu^{y}} {u^{x}} \left(\left(u^{x}\right)^2+1\right) {u^{y}}\right)}{{T} \left(u^{\tau}\right)^2} \\
(A_x)_{9,7} &= {(\partial_{\alpha} n)} {D_{s}} {u^{x}} {u^{y}} \\
(A_x)_{9,9} &= \frac{{D_{s}} {h} {u^{x}}}{{T}} \\
(A_x)_{10,1} &= \frac{{D_{s}} \left(3 {\mathcal{C}'(T)} {\pi_{Q}^{yy}} {u^{x}}-2 {\mathcal{B}'(T)} \left({\nu^{x}} \left(u^{y}\right)^2+{\nu^{x}}-3 {\nu^{y}} {u^{x}} {u^{y}}\right)\right)}{6 {T}} \\
(A_x)_{10,2} &= -\frac{1}{6 {T} \left(u^{\tau}\right)^2}{D_{s}}\Bigl(2 {I_{31}} \left(u^{\tau}\right)^2 \left(\left(u^{y}\right)^2+1\right)\\
  &\quad + 2 {u^{x}} \left({\mathcal{A}} \left(u^{y}\right)^2+{\mathcal{A}}+3 {\mathcal{B}} \left(u^{y}\right)^2\right) \left({\nu^{x}} \left(u^{y}\right)^2+{\nu^{x}}-{\nu^{y}} {u^{x}} {u^{y}}\right)\\
  &\quad + {\mathcal{C}}\Bigl(2 {\Pi_{Q}} \left(u^{\tau}\right)^2 \left(\left(u^{y}\right)^2+1\right)\\
  &\quad - 2 \left(\left(u^{y}\right)^2+1\right) \left({\pi_{Q}{}^{\eta}{}_{\eta}} \left(u^{\tau}\right)^2-4 {\pi_{Q}^{xy}} {u^{x}} {u^{y}}\right)\\
  &\quad - {\pi_{Q}^{yy}} \left(\left(u^{x}\right)^2 \left(8 \left(u^{y}\right)^2+5\right)+5 \left(\left(u^{y}\right)^2+1\right)\right)\Bigr)\Bigr) \\
(A_x)_{10,3} &= \frac{1}{3 {T} \left(u^{\tau}\right)^2}{D_{s}}\Bigl({u^{x}}\Bigl(3 {I_{31}} \left(u^{\tau}\right)^2 {u^{y}} - {\mathcal{A}} {\nu^{y}} \left(u^{x}\right)^2 \left(\left(u^{y}\right)^2-2\right) + 2 {\mathcal{A}} {\nu^{y}} \left(\left(u^{y}\right)^2+1\right)\\
  &\quad + {\mathcal{A}} {\nu^{x}} {u^{x}} {u^{y}} \left(\left(u^{y}\right)^2+1\right)\\
  &\quad + 3 {\mathcal{B}} \left({\nu^{x}} {u^{x}} \left(u^{y}\right)^3-{\nu^{y}} \left(u^{x}\right)^2 \left(\left(u^{y}\right)^2-1\right)+{\nu^{y}}\right)\Bigr)\\
  &\quad + {\mathcal{C}}\Bigl(- {u^{x}} {u^{y}} \left(-3 {\Pi_{Q}} \left(u^{\tau}\right)^2+4 {\pi_{Q}^{yy}} \left(\left(u^{x}\right)^2+1\right)+{\pi_{Q}{}^{\eta}{}_{\eta}} \left(u^{\tau}\right)^2\right)\\
  &\quad + 2 {\pi_{Q}^{xy}} \left(\left(u^{x}\right)^2 \left(2 \left(u^{y}\right)^2+1\right)+\left(u^{y}\right)^2+1\right)\Bigr)\Bigr) \\
(A_x)_{10,8} &= -\frac{{\mathcal{B}} {D_{s}} \left(\left(u^{y}\right)^2+1\right)}{3 {T}} \\
(A_x)_{10,9} &= \frac{{\mathcal{B}} {D_{s}} {u^{x}} {u^{y}}}{{T}} \\
(A_x)_{10,10} &= \frac{{\mathcal{C}} {D_{s}} {u^{x}}}{2 {T}} \\
(A_x)_{11,1} &= -\frac{{D_{s}} (2 {\mathcal{B}'(T)} {\nu^{x}}-3 {\mathcal{C}'(T)} {\pi_{Q}{}^{\eta}{}_{\eta}} {u^{x}})}{6 {T}} \\
(A_x)_{11,2} &= -\frac{1}{6 {T} \left(u^{\tau}\right)^2}{D_{s}}\Bigl(2 {I_{31}} \left(u^{\tau}\right)^2 + 2 {\mathcal{A}} {u^{x}} \left({\nu^{x}} \left(u^{y}\right)^2+{\nu^{x}}-{\nu^{y}} {u^{x}} {u^{y}}\right)\\
  &\quad + {\mathcal{C}}\Bigl(- 5 {\pi_{Q}{}^{\eta}{}_{\eta}} + 2 {\Pi_{Q}} \left(u^{\tau}\right)^2 - 5 {\pi_{Q}{}^{\eta}{}_{\eta}} \left(u^{x}\right)^2 - 2 {\pi_{Q}^{yy}} \left(\left(u^{x}\right)^2+1\right)\\
  &\quad + 2 {\pi_{Q}^{xy}} {u^{x}} {u^{y}} - 5 {\pi_{Q}{}^{\eta}{}_{\eta}} \left(u^{y}\right)^2\Bigr)\Bigr) \\
(A_x)_{11,3} &= \frac{1}{3 {T} \left(u^{\tau}\right)^2 \left(\left(u^{y}\right)^2+1\right)}{D_{s}}\Bigl(- {\mathcal{C}} {u^{x}} {u^{y}} \left({\pi_{Q}^{yy}} \left(\left(u^{x}\right)^2+1\right)+{\pi_{Q}{}^{\eta}{}_{\eta}} \left(u^{\tau}\right)^2\right)\\
  &\quad - {\mathcal{A}} {u^{x}} \left(\left(u^{y}\right)^2+1\right) \left(-{\nu^{x}} {u^{x}} {u^{y}}+{\nu^{y}} \left(u^{x}\right)^2+{\nu^{y}}\right)\\
  &\quad + {\mathcal{C}} {\pi_{Q}^{xy}} \left(\left(u^{x}\right)^2 \left(\left(u^{y}\right)^2-1\right)-\left(u^{y}\right)^2-1\right)\Bigr) \\
(A_x)_{11,8} &= -\frac{{\mathcal{B}} {D_{s}}}{3 {T}} \\
(A_x)_{11,11} &= \frac{{\mathcal{C}} {D_{s}} {u^{x}}}{2 {T}} \\
(A_x)_{12,1} &= \frac{{D_{s}} \left({\mathcal{B}'(T)} {\nu^{x}} {u^{x}} {u^{y}}+3 {\mathcal{B}'(T)} {\nu^{y}} \left(\left(u^{x}\right)^2+1\right)+3 {\mathcal{C}'(T)} {\pi_{Q}^{xy}} {u^{x}}\right)}{6 {T}} \\
(A_x)_{12,2} &= \frac{1}{6 {T} \left(u^{\tau}\right)^2}{D_{s}}\Bigl(3 {\mathcal{B}} {\nu^{y}} {u^{x}} + 3 {\mathcal{B}} {\nu^{y}} \left(u^{x}\right)^3 - 3 {\mathcal{B}} {\nu^{x}} {u^{y}} + {I_{31}} {u^{x}} {u^{y}} - 6 {\mathcal{B}} {\nu^{x}} \left(u^{x}\right)^2 {u^{y}} + {I_{31}} \left(u^{x}\right)^3 {u^{y}}\\
  &\quad + 6 {\mathcal{B}} {\nu^{y}} {u^{x}} \left(u^{y}\right)^2 + 6 {\mathcal{B}} {\nu^{y}} \left(u^{x}\right)^3 \left(u^{y}\right)^2 - 3 {\mathcal{B}} {\nu^{x}} \left(u^{y}\right)^3 + {I_{31}} {u^{x}} \left(u^{y}\right)^3\\
  &\quad - 6 {\mathcal{B}} {\nu^{x}} \left(u^{x}\right)^2 \left(u^{y}\right)^3\\
  &\quad + {\mathcal{A}} {u^{x}}\Bigl(3 {\nu^{y}} \left(\left(u^{y}\right)^2+1\right) - 2 {\nu^{x}} {u^{x}} {u^{y}} \left(\left(u^{y}\right)^2+1\right)\\
  &\quad + {\nu^{y}} \left(u^{x}\right)^2 \left(2 \left(u^{y}\right)^2+3\right)\Bigr)\\
  &\quad + {\mathcal{C}}\Bigl({u^{x}} {u^{y}} \left({\Pi_{Q}} \left(u^{\tau}\right)^2+{\pi_{Q}^{yy}} \left(8 \left(u^{x}\right)^2+5\right)+5 {\pi_{Q}{}^{\eta}{}_{\eta}} \left(u^{\tau}\right)^2\right)\\
  &\quad + {\pi_{Q}^{xy}} \left(\left(u^{x}\right)^2 \left(3-8 \left(u^{y}\right)^2\right)+6 \left(\left(u^{y}\right)^2+1\right)\right)\Bigr)\Bigr) \\
(A_x)_{12,3} &= \frac{1}{6 {T} \left(u^{\tau}\right)^2 \left(\left(u^{y}\right)^2+1\right)}{D_{s}}\Bigl(3 {I_{31}} \left(u^{\tau}\right)^2 \left(\left(u^{x}\right)^2+1\right) \left(\left(u^{y}\right)^2+1\right)\\
  &\quad + {\mathcal{C}}\Bigl(- 3 {\pi_{Q}{}^{\eta}{}_{\eta}} - 6 {\pi_{Q}{}^{\eta}{}_{\eta}} \left(u^{x}\right)^2 - 3 {\pi_{Q}{}^{\eta}{}_{\eta}} \left(u^{x}\right)^4 + {\pi_{Q}^{xy}} {u^{x}} {u^{y}}\\
  &\quad + 4 {\pi_{Q}^{xy}} \left(u^{x}\right)^3 {u^{y}} - 6 {\pi_{Q}{}^{\eta}{}_{\eta}} \left(u^{y}\right)^2 - 11 {\pi_{Q}{}^{\eta}{}_{\eta}} \left(u^{x}\right)^2 \left(u^{y}\right)^2\\
  &\quad - 5 {\pi_{Q}{}^{\eta}{}_{\eta}} \left(u^{x}\right)^4 \left(u^{y}\right)^2 + {\pi_{Q}^{xy}} {u^{x}} \left(u^{y}\right)^3 + 8 {\pi_{Q}^{xy}} \left(u^{x}\right)^3 \left(u^{y}\right)^3\\
  &\quad - 3 {\pi_{Q}{}^{\eta}{}_{\eta}} \left(u^{y}\right)^4 - 5 {\pi_{Q}{}^{\eta}{}_{\eta}} \left(u^{x}\right)^2 \left(u^{y}\right)^4\\
  &\quad + 3 {\Pi_{Q}} \left(u^{\tau}\right)^2 \left(\left(u^{x}\right)^2+1\right) \left(\left(u^{y}\right)^2+1\right)\\
  &\quad - {\pi_{Q}^{yy}} \left(\left(u^{x}\right)^2+1\right)\Bigl(\left(u^{x}\right)^2 \left(8 \left(u^{y}\right)^2+6\right)+3 \left(\left(u^{y}\right)^2+1\right)\Bigr)\Bigr)\\
  &\quad + \left(u^{y}\right)^2+1\Bigl({\mathcal{A}} {u^{x}}\Bigl(- 2 {\nu^{y}} {u^{x}} \left(\left(u^{x}\right)^2+1\right) {u^{y}} + 3 {\nu^{x}} \left(\left(u^{y}\right)^2+1\right)\\
  &\quad + {\nu^{x}} \left(u^{x}\right)^2 \left(2 \left(u^{y}\right)^2+3\right)\Bigr)\\
  &\quad + 3{\mathcal{B}} \left(\left(u^{x}\right)^2+1\right)\Bigl({\nu^{x}} \left(2 {u^{x}} \left(u^{y}\right)^2+{u^{x}}\right)-{\nu^{y}} \left(2 \left(u^{x}\right)^2+1\right) {u^{y}}\Bigr)\Bigr)\Bigr) \\
(A_x)_{12,8} &= \frac{{\mathcal{B}} {D_{s}} {u^{x}} {u^{y}}}{6 {T}} \\
(A_x)_{12,9} &= \frac{{\mathcal{B}} {D_{s}} \left(\left(u^{x}\right)^2+1\right)}{2 {T}} \\
(A_x)_{12,12} &= \frac{{\mathcal{C}} {D_{s}} {u^{x}}}{2 {T}} \\
(A_x)_{13,1} &= \frac{{D_{s}} (5 {\mathcal{B}'(T)} {\nu^{x}}+3 {u^{x}} ({\mathcal{C}'(T)} {\Pi_{Q}}+{(\partial_{T} I_{31})}))}{6 {T}} \\
(A_x)_{13,2} &= \frac{1}{6 {T} \left(u^{\tau}\right)^2}{D_{s}}\Bigl(5 {I_{31}} \left(u^{\tau}\right)^2 + 5 {\mathcal{C}} {\Pi_{Q}} \left(u^{\tau}\right)^2 + 2 {\mathcal{A}} {u^{x}} \left({\nu^{x}} \left(u^{y}\right)^2+{\nu^{x}}-{\nu^{y}} {u^{x}} {u^{y}}\right)\\
  &\quad - 2 {\mathcal{C}} \left(-{\pi_{Q}^{xy}} {u^{x}} {u^{y}}+{\pi_{Q}^{yy}} \left(u^{x}\right)^2+{\pi_{Q}^{yy}}+{\pi_{Q}{}^{\eta}{}_{\eta}} \left(u^{x}\right)^2+{\pi_{Q}{}^{\eta}{}_{\eta}} \left(u^{y}\right)^2+{\pi_{Q}{}^{\eta}{}_{\eta}}\right)\Bigr) \\
(A_x)_{13,3} &= \frac{1}{3 {T} \left(u^{\tau}\right)^2 \left(\left(u^{y}\right)^2+1\right)}{D_{s}}\Bigl({\mathcal{C}} {u^{x}} {u^{y}} \left({\pi_{Q}^{yy}} \left(\left(u^{x}\right)^2+1\right)+{\pi_{Q}{}^{\eta}{}_{\eta}} \left(u^{\tau}\right)^2\right)\\
  &\quad + {\mathcal{A}} {u^{x}} \left(\left(u^{y}\right)^2+1\right) \left(-{\nu^{x}} {u^{x}} {u^{y}}+{\nu^{y}} \left(u^{x}\right)^2+{\nu^{y}}\right)\\
  &\quad + {\mathcal{C}} {\pi_{Q}^{xy}} \left(-\left(\left(u^{x}\right)^2 \left(\left(u^{y}\right)^2-1\right)\right)+\left(u^{y}\right)^2+1\right)\Bigr) \\
(A_x)_{13,7} &= \frac{{(\partial_{\alpha} I_{31})} {D_{s}} {u^{x}}}{2 {T}} \\
(A_x)_{13,8} &= \frac{5 {\mathcal{B}} {D_{s}}}{6 {T}} \\
(A_x)_{13,13} &= \frac{{\mathcal{C}} {D_{s}} {u^{x}}}{2 {T}} \\
(A_y)_{1,1} &= \frac{\left(u^{\tau}\right)^2 {u^{y}} \left(16 {a_{\mathrm{SB}}} {T}^3+{(\partial_{T} n)} {h}\right)+{h'(T)} \left({n} \left(u^{\tau}\right)^2 {u^{y}}+{\nu^{x}} {u^{x}} {u^{y}}+{\nu^{y}} \left(\left(u^{x}\right)^2+2 \left(u^{y}\right)^2+1\right)\right)}{{u^{\tau}}} \\
(A_y)_{1,2} &= \frac{1}{\left(u^{\tau}\right)^3}\Bigl({u^{x}} {u^{y}} \left(4 {a_{\mathrm{SB}}} {T}^4 \left(u^{\tau}\right)^2-{\pi^{yy}}\right) + {\pi^{xy}} \left(\left(u^{y}\right)^2+1\right)\\
  &\quad + {h} \left({n} \left(u^{\tau}\right)^2 {u^{x}} {u^{y}}+{\nu^{x}} {u^{y}} \left(\left(u^{y}\right)^2+1\right)+{\nu^{y}} \left(\left(u^{x}\right)^3+{u^{x}}\right)\right)\Bigr) \\
(A_y)_{1,3} &= \frac{1}{\left(u^{\tau}\right)^3}\Bigl({\pi^{yy}} + 4 {a_{\mathrm{SB}}} {T}^4 + {\pi^{yy}} \left(u^{x}\right)^2 + 8 {a_{\mathrm{SB}}} {T}^4 \left(u^{x}\right)^2 + 4 {a_{\mathrm{SB}}} {T}^4 \left(u^{x}\right)^4 - {\pi^{xy}} {u^{x}} {u^{y}}\\
  &\quad + 12 {a_{\mathrm{SB}}} {T}^4 \left(u^{y}\right)^2 + 12 {a_{\mathrm{SB}}} {T}^4 \left(u^{x}\right)^2 \left(u^{y}\right)^2 + 8 {a_{\mathrm{SB}}} {T}^4 \left(u^{y}\right)^4\\
  &\quad + {h}\Bigl({\nu^{x}} \left(\left(u^{x}\right)^3+{u^{x}}\right) + {\nu^{y}} {u^{y}} \left(3 \left(u^{x}\right)^2+2 \left(u^{y}\right)^2+3\right)\\
  &\quad + {n} \left(\left(u^{x}\right)^4+\left(u^{x}\right)^2 \left(3 \left(u^{y}\right)^2+2\right)+2 \left(u^{y}\right)^4+3 \left(u^{y}\right)^2+1\right)\Bigr)\Bigr) \\
(A_y)_{1,4} &= \frac{{u^{y}}}{{u^{\tau}}} \\
(A_y)_{1,6} &= \frac{{u^{x}}}{{u^{\tau}}} \\
(A_y)_{1,7} &= {(\partial_{\alpha} n)} {h} {u^{\tau}} {u^{y}} \\
(A_y)_{1,8} &= \frac{{h} {u^{x}} {u^{y}}}{{u^{\tau}}} \\
(A_y)_{1,9} &= \frac{{h} \left(\left(u^{x}\right)^2+2 \left(u^{y}\right)^2+1\right)}{{u^{\tau}}} \\
(A_y)_{2,1} &= {u^{x}} {u^{y}} \left(16 {a_{\mathrm{SB}}} {T}^3+{(\partial_{T} n)} {h}\right)+{h'(T)} {u^{y}} ({n} {u^{x}}+{\nu^{x}})+{h'(T)} {\nu^{y}} {u^{x}} \\
(A_y)_{2,2} &= 4 {a_{\mathrm{SB}}} {T}^4 {u^{y}}+{h} ({n} {u^{y}}+{\nu^{y}}) \\
(A_y)_{2,3} &= 4 {a_{\mathrm{SB}}} {T}^4 {u^{x}}+{h} ({n} {u^{x}}+{\nu^{x}}) \\
(A_y)_{2,6} &= 1 \\
(A_y)_{2,7} &= {(\partial_{\alpha} n)} {h} {u^{x}} {u^{y}} \\
(A_y)_{2,8} &= {h} {u^{y}} \\
(A_y)_{2,9} &= {h} {u^{x}} \\
(A_y)_{3,1} &= 16 {a_{\mathrm{SB}}} {T}^3 \left(u^{y}\right)^2+4 {a_{\mathrm{SB}}} {T}^3+{h'(T)} {n} \left(u^{y}\right)^2+2 {h'(T)} {\nu^{y}} {u^{y}}+{(\partial_{T} n)} {h} \left(u^{y}\right)^2+{(\partial_{T} n)} {T}+{n} \\
(A_y)_{3,3} &= 8 {a_{\mathrm{SB}}} {T}^4 {u^{y}}+2 {h} ({n} {u^{y}}+{\nu^{y}}) \\
(A_y)_{3,4} &= 1 \\
(A_y)_{3,7} &= {(\partial_{\alpha} n)} \left({h} \left(u^{y}\right)^2+{T}\right) \\
(A_y)_{3,9} &= 2 {h} {u^{y}} \\
(A_y)_{4,2} &= -\frac{2 {\tau_{\pi}} \left(u^{y}\right)^2 \left({\pi^{xy}} \left(u^{y}\right)^2+{\pi^{xy}}-{\pi^{yy}} {u^{x}} {u^{y}}\right)}{\left(u^{\tau}\right)^2} \\
(A_y)_{4,3} &= \frac{1}{3 \left(u^{\tau}\right)^2}\Bigl(3 {\delta_{\pi\pi}} {\pi^{yy}} \left(u^{\tau}\right)^2 - 6 {\pi^{yy}} {\tau_{\pi}} \left(\left(u^{x}\right)^2+1\right) \left(u^{y}\right)^2 + 6 {\pi^{xy}} {\tau_{\pi}} {u^{x}} \left(u^{y}\right)^3 + 4 {\eta}\\
  &\quad + 4 \left(u^{x}\right)^2 {\eta} + 8 \left(u^{y}\right)^2 {\eta} + 4 \left(u^{x}\right)^2 \left(u^{y}\right)^2 {\eta} + 4 \left(u^{y}\right)^4 {\eta}\Bigr) \\
(A_y)_{4,4} &= {\tau_{\pi}} {u^{y}} \\
(A_y)_{5,3} &= \frac{3 {\delta_{\pi\pi}} {\pi^{\eta}{}_{\eta}}-2 {\eta}}{3 {\tau}^2} \\
(A_y)_{5,5} &= \frac{{\tau_{\pi}} {u^{y}}}{{\tau}^2} \\
(A_y)_{6,2} &= \frac{1}{\left(u^{\tau}\right)^2}\Bigl({\pi^{yy}} {\tau_{\pi}} \left(u^{y}\right)^2 + {\pi^{\eta}{}_{\eta}} {\tau_{\pi}} \left(u^{\tau}\right)^2 \left(u^{y}\right)^2 + 2 {\pi^{yy}} {\tau_{\pi}} \left(u^{x}\right)^2 \left(u^{y}\right)^2\\
  &\quad - {\pi^{xy}} {\tau_{\pi}} {u^{x}} {u^{y}} \left(2 \left(u^{y}\right)^2+1\right) + {\eta} + \left(u^{x}\right)^2 {\eta} + 2 \left(u^{y}\right)^2 {\eta} + \left(u^{x}\right)^2 \left(u^{y}\right)^2 {\eta} + \left(u^{y}\right)^4 {\eta}\Bigr) \\
(A_y)_{6,3} &= {\delta_{\pi\pi}} {\pi^{xy}}\\
  &\quad - \frac{1}{3 \left(u^{\tau}\right)^2 \left(\left(u^{y}\right)^2+1\right)}{u^{y}}\Bigl(3 {\pi^{yy}} {\tau_{\pi}} {u^{x}} \left(\left(u^{x}\right)^2+1\right) \left(2 \left(u^{y}\right)^2+1\right)\\
  &\quad + 3 {\pi^{xy}} {\tau_{\pi}} {u^{y}} \left(\left(1-2 \left(u^{x}\right)^2\right) \left(u^{y}\right)^2+1\right)\\
  &\quad + \left(u^{\tau}\right)^2 {u^{x}} \left(3 {\pi^{\eta}{}_{\eta}} {\tau_{\pi}} \left(u^{y}\right)^2-\left(\left(u^{y}\right)^2+1\right) {\eta}\right)\Bigr) \\
(A_y)_{6,6} &= {\tau_{\pi}} {u^{y}} \\
(A_y)_{7,1} &= {(\partial_{T} n)} {u^{y}} \\
(A_y)_{7,3} &= {n} \\
(A_y)_{7,7} &= {(\partial_{\alpha} n)} {u^{y}} \\
(A_y)_{7,9} &= 1 \\
(A_y)_{8,1} &= \frac{{D_{s}} {u^{y}} ({h'(T)} {\nu^{x}}+{u^{x}} ({(\partial_{T} n)} {T}+{n}))}{{T}} \\
(A_y)_{8,2} &= \frac{{D_{s}} {h} \left({n} \left(u^{\tau}\right)^2 {u^{y}}-{\nu^{x}} {u^{x}} {u^{y}} \left(\left(u^{y}\right)^2+1\right)+{\nu^{y}} \left(\left(u^{x}\right)^2+1\right) \left(\left(u^{y}\right)^2+1\right)\right)}{{T} \left(u^{\tau}\right)^2} \\
(A_y)_{8,3} &= \frac{{D_{s}} {h} \left(\left(u^{x}\right)^2+1\right) \left({\nu^{x}} \left(u^{y}\right)^2+{\nu^{x}}-{\nu^{y}} {u^{x}} {u^{y}}\right)}{{T} \left(u^{\tau}\right)^2} \\
(A_y)_{8,7} &= {(\partial_{\alpha} n)} {D_{s}} {u^{x}} {u^{y}} \\
(A_y)_{8,8} &= \frac{{D_{s}} {h} {u^{y}}}{{T}} \\
(A_y)_{9,1} &= \frac{{D_{s}} \left({h'(T)} {\nu^{y}} {u^{y}}+{(\partial_{T} n)} {T} \left(\left(u^{y}\right)^2+1\right)+{n} \left(\left(u^{y}\right)^2+1\right)\right)}{{T}} \\
(A_y)_{9,2} &= -\frac{{D_{s}} {h} \left(u^{y}\right)^2 \left({\nu^{x}} \left(u^{y}\right)^2+{\nu^{x}}-{\nu^{y}} {u^{x}} {u^{y}}\right)}{{T} \left(u^{\tau}\right)^2} \\
(A_y)_{9,3} &= \frac{{D_{s}} {h} \left({n} \left(u^{\tau}\right)^2 {u^{y}}+{\nu^{x}} {u^{x}} \left(u^{y}\right)^3+{\nu^{y}} \left(-\left(\left(u^{x}\right)^2 \left(\left(u^{y}\right)^2-2\right)\right)+\left(u^{y}\right)^2+2\right)\right)}{{T} \left(u^{\tau}\right)^2} \\
(A_y)_{9,7} &= {(\partial_{\alpha} n)} {D_{s}} \left(\left(u^{y}\right)^2+1\right) \\
(A_y)_{9,9} &= \frac{{D_{s}} {h} {u^{y}}}{{T}} \\
(A_y)_{10,1} &= \frac{{D_{s}} \left(4 {\mathcal{B}'(T)} {\nu^{y}} \left(\left(u^{y}\right)^2+1\right)+3 {\mathcal{C}'(T)} {\pi_{Q}^{yy}} {u^{y}}\right)}{6 {T}} \\
(A_y)_{10,2} &= -\frac{1}{3 {T} \left(u^{\tau}\right)^2}{D_{s}}\Bigl({u^{y}} \left(\left(u^{y}\right)^2+1\right) ({\mathcal{A}}+3 {\mathcal{B}}) \left({\nu^{x}} \left(u^{y}\right)^2+{\nu^{x}}-{\nu^{y}} {u^{x}} {u^{y}}\right)\\
  &\quad + {\mathcal{C}} \left(4 \left(u^{y}\right)^2+1\right) \left({\pi_{Q}^{xy}} \left(u^{y}\right)^2+{\pi_{Q}^{xy}}-{\pi_{Q}^{yy}} {u^{x}} {u^{y}}\right)\Bigr) \\
(A_y)_{10,3} &= \frac{1}{6 {T} \left(u^{\tau}\right)^2}{D_{s}}\Bigl(4 {I_{31}} \left(u^{\tau}\right)^2 \left(\left(u^{y}\right)^2+1\right)\\
  &\quad + {\mathcal{C}}\Bigl(4 {\Pi_{Q}} \left(u^{\tau}\right)^2 \left(\left(u^{y}\right)^2+1\right) + 2 {\pi_{Q}^{xy}} {u^{x}} {u^{y}} \left(4 \left(u^{y}\right)^2+1\right)\\
  &\quad + {\pi_{Q}^{yy}} \left(\left(u^{x}\right)^2 \left(7-8 \left(u^{y}\right)^2\right)+\left(u^{y}\right)^2+7\right)\Bigr)\\
  &\quad + 2{u^{y}}\Bigl(3 {\mathcal{B}} {u^{x}} {u^{y}} \left({\nu^{x}} \left(u^{y}\right)^2+{\nu^{x}}-{\nu^{y}} {u^{x}} {u^{y}}\right)\\
  &\quad + {\mathcal{A}}\Bigl(- {\nu^{y}} \left(u^{x}\right)^2 \left(\left(u^{y}\right)^2-2\right) + 2 {\nu^{y}} \left(\left(u^{y}\right)^2+1\right)\\
  &\quad + {\nu^{x}} {u^{x}} {u^{y}} \left(\left(u^{y}\right)^2+1\right)\Bigr)\Bigr)\Bigr) \\
(A_y)_{10,9} &= \frac{2 {\mathcal{B}} {D_{s}} \left(\left(u^{y}\right)^2+1\right)}{3 {T}} \\
(A_y)_{10,10} &= \frac{{\mathcal{C}} {D_{s}} {u^{y}}}{2 {T}} \\
(A_y)_{11,1} &= -\frac{{D_{s}} (2 {\mathcal{B}'(T)} {\nu^{y}}-3 {\mathcal{C}'(T)} {\pi_{Q}{}^{\eta}{}_{\eta}} {u^{y}})}{6 {T}} \\
(A_y)_{11,2} &= -\frac{{D_{s}} \left({\mathcal{A}} {u^{y}} \left({\nu^{x}} \left(u^{y}\right)^2+{\nu^{x}}-{\nu^{y}} {u^{x}} {u^{y}}\right)+{\mathcal{C}} \left({\pi_{Q}^{xy}} \left(u^{y}\right)^2+{\pi_{Q}^{xy}}-{\pi_{Q}^{yy}} {u^{x}} {u^{y}}\right)\right)}{3 {T} \left(u^{\tau}\right)^2} \\
(A_y)_{11,3} &= -\frac{1}{6 {T} \left(u^{\tau}\right)^2}{D_{s}}\Bigl(2 {I_{31}} \left(u^{\tau}\right)^2 + 2 {\mathcal{A}} {u^{y}} \left(-{\nu^{x}} {u^{x}} {u^{y}}+{\nu^{y}} \left(u^{x}\right)^2+{\nu^{y}}\right)\\
  &\quad + {\mathcal{C}}\Bigl(- 3 {\pi_{Q}{}^{\eta}{}_{\eta}} + 2 {\Pi_{Q}} \left(u^{\tau}\right)^2 - 3 {\pi_{Q}{}^{\eta}{}_{\eta}} \left(u^{x}\right)^2 + 2 {\pi_{Q}^{yy}} \left(\left(u^{x}\right)^2+1\right)\\
  &\quad - 2 {\pi_{Q}^{xy}} {u^{x}} {u^{y}} - 3 {\pi_{Q}{}^{\eta}{}_{\eta}} \left(u^{y}\right)^2\Bigr)\Bigr) \\
(A_y)_{11,9} &= -\frac{{\mathcal{B}} {D_{s}}}{3 {T}} \\
(A_y)_{11,11} &= \frac{{\mathcal{C}} {D_{s}} {u^{y}}}{2 {T}} \\
(A_y)_{12,1} &= \frac{{D_{s}} \left(3 {\mathcal{B}'(T)} {\nu^{x}} \left(\left(u^{y}\right)^2+1\right)+{\mathcal{B}'(T)} {\nu^{y}} {u^{x}} {u^{y}}+3 {\mathcal{C}'(T)} {\pi_{Q}^{xy}} {u^{y}}\right)}{6 {T}} \\
(A_y)_{12,2} &= \frac{1}{6 {T} \left(u^{\tau}\right)^2}{D_{s}}\Bigl(- 3 {\mathcal{B}} {\nu^{x}} {u^{x}} + 3 {\mathcal{A}} {\nu^{y}} {u^{y}} + 3 {\mathcal{B}} {\nu^{y}} {u^{y}} + 3 {\mathcal{A}} {\nu^{y}} \left(u^{x}\right)^2 {u^{y}} + 6 {\mathcal{B}} {\nu^{y}} \left(u^{x}\right)^2 {u^{y}}\\
  &\quad - 2 {\mathcal{A}} {\nu^{x}} {u^{x}} \left(u^{y}\right)^2 - 9 {\mathcal{B}} {\nu^{x}} {u^{x}} \left(u^{y}\right)^2 + 3 {\mathcal{A}} {\nu^{y}} \left(u^{y}\right)^3 + 3 {\mathcal{B}} {\nu^{y}} \left(u^{y}\right)^3\\
  &\quad + 2 {\mathcal{A}} {\nu^{y}} \left(u^{x}\right)^2 \left(u^{y}\right)^3 + 6 {\mathcal{B}} {\nu^{y}} \left(u^{x}\right)^2 \left(u^{y}\right)^3 - 2 {\mathcal{A}} {\nu^{x}} {u^{x}} \left(u^{y}\right)^4 - 6 {\mathcal{B}} {\nu^{x}} {u^{x}} \left(u^{y}\right)^4\\
  &\quad + 3 {I_{31}} \left(u^{\tau}\right)^2 \left(\left(u^{y}\right)^2+1\right)\\
  &\quad + {\mathcal{C}}\Bigl(3 {\Pi_{Q}} \left(u^{\tau}\right)^2 \left(\left(u^{y}\right)^2+1\right)\\
  &\quad + {\pi_{Q}^{yy}} \left(\left(u^{x}\right)^2 \left(8 \left(u^{y}\right)^2+3\right)+6 \left(u^{y}\right)^2+3\right)\\
  &\quad + {u^{y}} \left(3 {\pi_{Q}{}^{\eta}{}_{\eta}} \left(u^{\tau}\right)^2 {u^{y}}-{\pi_{Q}^{xy}} {u^{x}} \left(8 \left(u^{y}\right)^2+5\right)\right)\Bigr)\Bigr) \\
(A_y)_{12,3} &= \frac{1}{6 {T} \left(u^{\tau}\right)^2 \left(\left(u^{y}\right)^2+1\right)}{D_{s}}\Bigl(\left(u^{y}\right)^2+1\Bigl({I_{31}} \left(u^{\tau}\right)^2 {u^{x}} {u^{y}}\\
  &\quad + 3 {\mathcal{B}} {\nu^{x}} \left(2 \left(u^{x}\right)^2+1\right) {u^{y}} \left(\left(u^{y}\right)^2+1\right)\\
  &\quad - 3 {\mathcal{B}} {\nu^{y}} {u^{x}} \left(\left(u^{x}\right)^2+1\right) \left(2 \left(u^{y}\right)^2+1\right)\\
  &\quad + {\mathcal{A}} {u^{y}}\Bigl(- 2 {\nu^{y}} {u^{x}} \left(\left(u^{x}\right)^2+1\right) {u^{y}} + 3 {\nu^{x}} \left(\left(u^{y}\right)^2+1\right)\\
  &\quad + {\nu^{x}} \left(u^{x}\right)^2 \left(2 \left(u^{y}\right)^2+3\right)\Bigr)\Bigr)\\
  &\quad + {\mathcal{C}}\Bigl({u^{x}} {u^{y}}\Bigl(- 3 {\pi_{Q}{}^{\eta}{}_{\eta}} \left(u^{\tau}\right)^2 \left(u^{y}\right)^2 + {\Pi_{Q}} \left(u^{\tau}\right)^2 \left(\left(u^{y}\right)^2+1\right)\\
  &\quad - {\pi_{Q}^{yy}} \left(\left(u^{x}\right)^2+1\right) \left(8 \left(u^{y}\right)^2+5\right)\Bigr)\\
  &\quad + {\pi_{Q}^{xy}}\Bigl(3 \left(\left(u^{y}\right)^4+3 \left(u^{y}\right)^2+2\right)\\
  &\quad + \left(u^{x}\right)^2 \left(8 \left(u^{y}\right)^4+8 \left(u^{y}\right)^2+6\right)\Bigr)\Bigr)\Bigr) \\
(A_y)_{12,8} &= \frac{{\mathcal{B}} {D_{s}} \left(\left(u^{y}\right)^2+1\right)}{2 {T}} \\
(A_y)_{12,9} &= \frac{{\mathcal{B}} {D_{s}} {u^{x}} {u^{y}}}{6 {T}} \\
(A_y)_{12,12} &= \frac{{\mathcal{C}} {D_{s}} {u^{y}}}{2 {T}} \\
(A_y)_{13,1} &= \frac{{D_{s}} (5 {\mathcal{B}'(T)} {\nu^{y}}+3 {u^{y}} ({\mathcal{C}'(T)} {\Pi_{Q}}+{(\partial_{T} I_{31})}))}{6 {T}} \\
(A_y)_{13,2} &= \frac{{D_{s}} \left({\mathcal{A}} {u^{y}} \left({\nu^{x}} \left(u^{y}\right)^2+{\nu^{x}}-{\nu^{y}} {u^{x}} {u^{y}}\right)+{\mathcal{C}} \left({\pi_{Q}^{xy}} \left(u^{y}\right)^2+{\pi_{Q}^{xy}}-{\pi_{Q}^{yy}} {u^{x}} {u^{y}}\right)\right)}{3 {T} \left(u^{\tau}\right)^2} \\
(A_y)_{13,3} &= \frac{1}{6 {T} \left(u^{\tau}\right)^2}{D_{s}}\Bigl(5 {I_{31}} \left(u^{\tau}\right)^2 + 5 {\mathcal{C}} {\Pi_{Q}} \left(u^{\tau}\right)^2 + 2 {\mathcal{A}} {u^{y}} \left(-{\nu^{x}} {u^{x}} {u^{y}}+{\nu^{y}} \left(u^{x}\right)^2+{\nu^{y}}\right)\\
  &\quad + 2 {\mathcal{C}} \left(-{\pi_{Q}^{xy}} {u^{x}} {u^{y}}+{\pi_{Q}^{yy}} \left(u^{x}\right)^2+{\pi_{Q}^{yy}}\right)\Bigr) \\
(A_y)_{13,7} &= \frac{{(\partial_{\alpha} I_{31})} {D_{s}} {u^{y}}}{2 {T}} \\
(A_y)_{13,9} &= \frac{5 {\mathcal{B}} {D_{s}}}{6 {T}} \\
(A_y)_{13,13} &= \frac{{\mathcal{C}} {D_{s}} {u^{y}}}{2 {T}} \\
S_{1} &= \frac{1}{{\tau} \left(\left(u^{y}\right)^2+1\right)}\Bigl(4 {a_{\mathrm{SB}}} {T}^4 - {\pi^{yy}} \left(u^{x}\right)^2 + 4 {a_{\mathrm{SB}}} {T}^4 \left(u^{x}\right)^2 + 2 {\pi^{xy}} {u^{x}} {u^{y}} + {\pi^{yy}} \left(u^{y}\right)^2\\
  &\quad + 8 {a_{\mathrm{SB}}} {T}^4 \left(u^{y}\right)^2 + 4 {a_{\mathrm{SB}}} {T}^4 \left(u^{x}\right)^2 \left(u^{y}\right)^2 + 4 {a_{\mathrm{SB}}} {T}^4 \left(u^{y}\right)^4\\
  &\quad + {\pi^{\eta}{}_{\eta}} \left(-\left(u^{x}\right)^2+\left(u^{y}\right)^2+1\right)\\
  &\quad + {h} \left(\left(u^{y}\right)^2+1\right) \left({n} \left(u^{\tau}\right)^2+2 ({\nu^{x}} {u^{x}}+{\nu^{y}} {u^{y}})\right)\Bigr) \\
S_{2} &= \frac{1}{{\tau} {u^{\tau}} \left(\left(u^{y}\right)^2+1\right)}\Bigl(- {\pi^{yy}} {u^{x}} + 4 {a_{\mathrm{SB}}} {T}^4 {u^{x}} - {\pi^{\eta}{}_{\eta}} \left(u^{\tau}\right)^2 {u^{x}} - {\pi^{yy}} \left(u^{x}\right)^3 + 4 {a_{\mathrm{SB}}} {T}^4 \left(u^{x}\right)^3\\
  &\quad + {\pi^{xy}} {u^{y}} + 2 {\pi^{xy}} \left(u^{x}\right)^2 {u^{y}} + 8 {a_{\mathrm{SB}}} {T}^4 {u^{x}} \left(u^{y}\right)^2 + 4 {a_{\mathrm{SB}}} {T}^4 \left(u^{x}\right)^3 \left(u^{y}\right)^2\\
  &\quad + {\pi^{xy}} \left(u^{y}\right)^3 + 4 {a_{\mathrm{SB}}} {T}^4 {u^{x}} \left(u^{y}\right)^4\\
  &\quad + {h} \left(\left(u^{y}\right)^2+1\right) \left({n} \left(u^{\tau}\right)^2 {u^{x}}+{\nu^{x}} \left(2 \left(u^{x}\right)^2+\left(u^{y}\right)^2+1\right)+{\nu^{y}} {u^{x}} {u^{y}}\right)\Bigr) \\
S_{3} &= \frac{{u^{y}} \left(4 {a_{\mathrm{SB}}} {T}^4 \left(u^{\tau}\right)^2+{\pi^{yy}}\right)+{h} \left({n} \left(u^{\tau}\right)^2 {u^{y}}+{\nu^{x}} {u^{x}} {u^{y}}+{\nu^{y}} \left(\left(u^{x}\right)^2+2 \left(u^{y}\right)^2+1\right)\right)+{\pi^{xy}} {u^{x}}}{{\tau} {u^{\tau}}} \\
S_{4} &= \frac{{\delta_{\pi\pi}} {\pi^{yy}} {u^{\tau}}}{{\tau}}+{\pi^{yy}}-\frac{2 {u^{\tau}} \left(\left(u^{y}\right)^2+1\right) {\eta}}{3 {\tau}} \\
S_{5} &= \frac{3 {\pi^{\eta}{}_{\eta}} ({\delta_{\pi\pi}} {u^{\tau}}+{\tau})+4 {u^{\tau}} {\eta}}{3 {\tau}^3} \\
S_{6} &= \frac{{\delta_{\pi\pi}} {\pi^{xy}} {u^{\tau}}}{{\tau}}+{\pi^{xy}}-\frac{2 {u^{\tau}} {u^{x}} {u^{y}} {\eta}}{3 {\tau}} \\
S_{7} &= \frac{{n} \left(u^{\tau}\right)^2+{\nu^{x}} {u^{x}}+{\nu^{y}} {u^{y}}}{{\tau} {u^{\tau}}} \\
S_{8} &= \frac{{D_{s}} {h} {\nu^{x}} {u^{\tau}}}{{T} {\tau}}+{\nu^{x}} \\
S_{9} &= \frac{{D_{s}} {h} {\nu^{y}} {u^{\tau}}}{{T} {\tau}}+{\nu^{y}} \\
S_{10} &= {\pi_{Q}^{yy}}\\
  &\quad - \frac{1}{6 {T} {\tau} {u^{\tau}}}{D_{s}}\Bigl(2 {I_{31}} \left(u^{\tau}\right)^2 \left(\left(u^{y}\right)^2+1\right) + 2 {\mathcal{B}} \left(\left(u^{y}\right)^2+1\right) ({\nu^{x}} {u^{x}}+{\nu^{y}} {u^{y}})\\
  &\quad + {\mathcal{C}} \left(u^{\tau}\right)^2 \left(2 {\Pi_{Q}} \left(\left(u^{y}\right)^2+1\right)-3 {\pi_{Q}^{yy}}+2 {\pi_{Q}{}^{\eta}{}_{\eta}} \left(\left(u^{y}\right)^2+1\right)\right)\Bigr) \\
S_{11} &= \frac{{D_{s}} \left(4 {\mathcal{B}} ({\nu^{x}} {u^{x}}+{\nu^{y}} {u^{y}})+{\mathcal{C}} \left(u^{\tau}\right)^2 (4 {\Pi_{Q}}+7 {\pi_{Q}{}^{\eta}{}_{\eta}})+4 {I_{31}} \left(u^{\tau}\right)^2\right)}{6 {T} {\tau} {u^{\tau}}}+{\pi_{Q}{}^{\eta}{}_{\eta}} \\
S_{12} &= -\frac{{D_{s}} {u^{x}} {u^{y}} \left({\mathcal{B}} ({\nu^{x}} {u^{x}}+{\nu^{y}} {u^{y}})+{I_{31}} \left(u^{\tau}\right)^2\right)}{3 {T} {\tau} {u^{\tau}}}+\frac{{\mathcal{C}} {D_{s}} {u^{\tau}} (3 {\pi_{Q}^{xy}}-2 {u^{x}} {u^{y}} ({\Pi_{Q}}+{\pi_{Q}{}^{\eta}{}_{\eta}}))}{6 {T} {\tau}}+{\pi_{Q}^{xy}} \\
S_{13} &= \frac{{D_{s}} \left(5 {\mathcal{B}} ({\nu^{x}} {u^{x}}+{\nu^{y}} {u^{y}})+{\mathcal{C}} \left(u^{\tau}\right)^2 (5 {\Pi_{Q}}+2 {\pi_{Q}{}^{\eta}{}_{\eta}})+5 {I_{31}} \left(u^{\tau}\right)^2\right)}{6 {T} {\tau} {u^{\tau}}}+{\Pi_{Q}}
\end{align*}
\subsubsection{S6, fully coupled}

\texttt{two\_\allowbreak{}d\_\allowbreak{}two\_\allowbreak{}way\_\allowbreak{}cM\_\allowbreak{}matrix} in \texttt{src/Matrix/2d\_\allowbreak{}two\_\allowbreak{}way\_\allowbreak{}cM.jl}, fields:

\[
\phi = ({T}, {u^{x}}, {u^{y}}, {\pi^{yy}}, {\pi^{\eta}{}_{\eta}}, {\pi^{xy}}, {\Pi}, {\alpha}, {\nu^{x}}, {\nu^{y}}, {\pi_{Q}^{yy}}, {\pi_{Q}{}^{\eta}{}_{\eta}}, {\pi_{Q}^{xy}}, {\Pi_{Q}})
\]
\begin{align*}
(A_t)_{1,1} &= - {n} - {(\partial_{T} n)} {T} + {P''(T)} {T} + {(\partial_{T} n)} {h} \left(u^{\tau}\right)^2 + {h'(T)} {n} \left(u^{\tau}\right)^2 + 2 {h'(T)} {\nu^{x}} {u^{x}}\\
  &\quad + {P'(T)} \left(u^{x}\right)^2 + {P''(T)} {T} \left(u^{x}\right)^2 + 2 {h'(T)} {\nu^{y}} {u^{y}} + {P'(T)} \left(u^{y}\right)^2 + {P''(T)} {T} \left(u^{y}\right)^2 \\
(A_t)_{1,2} &= \frac{2}{\left(u^{y}\right)^2+1}\Bigl({\Pi_{Q}} {u^{x}} - {\pi^{\eta}{}_{\eta}} {u^{x}} - {\pi^{yy}} {u^{x}} - {\pi_{Q}^{yy}} {u^{x}} - {\pi_{Q}{}^{\eta}{}_{\eta}} {u^{x}} + {P'(T)} {T} {u^{x}} + {\pi^{xy}} {u^{y}} + {\pi_{Q}^{xy}} {u^{y}}\\
  &\quad + {\Pi_{Q}} {u^{x}} \left(u^{y}\right)^2 + {P'(T)} {T} {u^{x}} \left(u^{y}\right)^2 + {\Pi} {u^{x}} \left(\left(u^{y}\right)^2+1\right)\\
  &\quad + {h} \left(\left(u^{y}\right)^2+1\right) ({n} {u^{x}}+{\nu^{x}})\Bigr) \\
(A_t)_{1,3} &= \frac{2}{\left(\left(u^{y}\right)^2+1\right)^2}\Bigl({\pi_{Q}^{xy}} {u^{x}} + {\Pi} {u^{y}} + {\Pi_{Q}} {u^{y}} + {\pi^{yy}} {u^{y}} + {\pi_{Q}^{yy}} {u^{y}} + {P'(T)} {T} {u^{y}} + {\pi^{\eta}{}_{\eta}} \left(u^{x}\right)^2 {u^{y}}\\
  &\quad + {\pi^{yy}} \left(u^{x}\right)^2 {u^{y}} + {\pi_{Q}^{yy}} \left(u^{x}\right)^2 {u^{y}} + {\pi_{Q}{}^{\eta}{}_{\eta}} \left(u^{x}\right)^2 {u^{y}} - {\pi_{Q}^{xy}} {u^{x}} \left(u^{y}\right)^2 + 2 {\Pi} \left(u^{y}\right)^3\\
  &\quad + 2 {\Pi_{Q}} \left(u^{y}\right)^3 + 2 {P'(T)} {T} \left(u^{y}\right)^3 + {\Pi} \left(u^{y}\right)^5 + {\Pi_{Q}} \left(u^{y}\right)^5 + {P'(T)} {T} \left(u^{y}\right)^5\\
  &\quad + {h} \left(\left(u^{y}\right)^2+1\right)^2 ({n} {u^{y}}+{\nu^{y}}) + {\pi^{xy}} \left({u^{x}}-{u^{x}} \left(u^{y}\right)^2\right)\Bigr) \\
(A_t)_{1,4} &= \frac{\left(u^{y}\right)^2-\left(u^{x}\right)^2}{\left(u^{y}\right)^2+1} \\
(A_t)_{1,5} &= -\frac{\left(u^{x}\right)^2}{\left(u^{y}\right)^2+1} \\
(A_t)_{1,6} &= \frac{2 {u^{x}} {u^{y}}}{\left(u^{y}\right)^2+1} \\
(A_t)_{1,7} &= \left(u^{x}\right)^2+\left(u^{y}\right)^2 \\
(A_t)_{1,8} &= {(\partial_{\alpha} n)} \left({h} \left(u^{\tau}\right)^2-{T}\right) \\
(A_t)_{1,9} &= 2 {h} {u^{x}} \\
(A_t)_{1,10} &= 2 {h} {u^{y}} \\
(A_t)_{1,11} &= \frac{\left(u^{y}\right)^2-\left(u^{x}\right)^2}{\left(u^{y}\right)^2+1} \\
(A_t)_{1,12} &= -\frac{\left(u^{x}\right)^2}{\left(u^{y}\right)^2+1} \\
(A_t)_{1,13} &= \frac{2 {u^{x}} {u^{y}}}{\left(u^{y}\right)^2+1} \\
(A_t)_{1,14} &= \left(u^{x}\right)^2+\left(u^{y}\right)^2 \\
(A_t)_{2,1} &= \frac{1}{{u^{\tau}}}\Bigl(\left(u^{\tau}\right)^2 {u^{x}} ({(\partial_{T} n)} {h}+{P''(T)} {T}+{P'(T)})\\
  &\quad + {h'(T)} \left({n} \left(u^{\tau}\right)^2 {u^{x}}+{\nu^{x}} \left(2 \left(u^{x}\right)^2+\left(u^{y}\right)^2+1\right)+{\nu^{y}} {u^{x}} {u^{y}}\right)\Bigr) \\
(A_t)_{2,2} &= \frac{1}{\left(u^{\tau}\right)^3 \left(\left(u^{y}\right)^2+1\right)}\Bigl({\Pi_{Q}} - {\pi^{\eta}{}_{\eta}} - {\pi^{yy}} - {\pi_{Q}^{yy}} - {\pi_{Q}{}^{\eta}{}_{\eta}} + {P'(T)} {T} + 3 {\Pi_{Q}} \left(u^{x}\right)^2\\
  &\quad - 3 {\pi^{\eta}{}_{\eta}} \left(u^{x}\right)^2 - 3 {\pi^{yy}} \left(u^{x}\right)^2 - 3 {\pi_{Q}^{yy}} \left(u^{x}\right)^2 - 3 {\pi_{Q}{}^{\eta}{}_{\eta}} \left(u^{x}\right)^2\\
  &\quad + 3 {P'(T)} {T} \left(u^{x}\right)^2 + 2 {\Pi_{Q}} \left(u^{x}\right)^4 - 2 {\pi^{\eta}{}_{\eta}} \left(u^{x}\right)^4 - 2 {\pi^{yy}} \left(u^{x}\right)^4\\
  &\quad - 2 {\pi_{Q}^{yy}} \left(u^{x}\right)^4 - 2 {\pi_{Q}{}^{\eta}{}_{\eta}} \left(u^{x}\right)^4 + 2 {P'(T)} {T} \left(u^{x}\right)^4 + 3 {\pi^{xy}} {u^{x}} {u^{y}}\\
  &\quad + 3 {\pi_{Q}^{xy}} {u^{x}} {u^{y}} + 2 {\pi^{xy}} \left(u^{x}\right)^3 {u^{y}} + 2 {\pi_{Q}^{xy}} \left(u^{x}\right)^3 {u^{y}} + 3 {\Pi_{Q}} \left(u^{y}\right)^2\\
  &\quad - 2 {\pi^{\eta}{}_{\eta}} \left(u^{y}\right)^2 - {\pi^{yy}} \left(u^{y}\right)^2 - {\pi_{Q}^{yy}} \left(u^{y}\right)^2 - 2 {\pi_{Q}{}^{\eta}{}_{\eta}} \left(u^{y}\right)^2\\
  &\quad + 3 {P'(T)} {T} \left(u^{y}\right)^2 + 6 {\Pi_{Q}} \left(u^{x}\right)^2 \left(u^{y}\right)^2 - 3 {\pi^{\eta}{}_{\eta}} \left(u^{x}\right)^2 \left(u^{y}\right)^2\\
  &\quad - 3 {\pi^{yy}} \left(u^{x}\right)^2 \left(u^{y}\right)^2 - 3 {\pi_{Q}^{yy}} \left(u^{x}\right)^2 \left(u^{y}\right)^2 - 3 {\pi_{Q}{}^{\eta}{}_{\eta}} \left(u^{x}\right)^2 \left(u^{y}\right)^2\\
  &\quad + 6 {P'(T)} {T} \left(u^{x}\right)^2 \left(u^{y}\right)^2 + 2 {\Pi_{Q}} \left(u^{x}\right)^4 \left(u^{y}\right)^2 + 2 {P'(T)} {T} \left(u^{x}\right)^4 \left(u^{y}\right)^2\\
  &\quad + 3 {\pi^{xy}} {u^{x}} \left(u^{y}\right)^3 + 3 {\pi_{Q}^{xy}} {u^{x}} \left(u^{y}\right)^3 + 3 {\Pi_{Q}} \left(u^{y}\right)^4 - {\pi^{\eta}{}_{\eta}} \left(u^{y}\right)^4\\
  &\quad - {\pi_{Q}{}^{\eta}{}_{\eta}} \left(u^{y}\right)^4 + 3 {P'(T)} {T} \left(u^{y}\right)^4 + 3 {\Pi_{Q}} \left(u^{x}\right)^2 \left(u^{y}\right)^4\\
  &\quad + 3 {P'(T)} {T} \left(u^{x}\right)^2 \left(u^{y}\right)^4 + {\Pi_{Q}} \left(u^{y}\right)^6 + {P'(T)} {T} \left(u^{y}\right)^6\\
  &\quad + {\Pi} \left(\left(u^{y}\right)^2+1\right) \left(2 \left(u^{x}\right)^4+3 \left(u^{x}\right)^2 \left(\left(u^{y}\right)^2+1\right)+\left(\left(u^{y}\right)^2+1\right)^2\right)\\
  &\quad + {h} \left(\left(u^{y}\right)^2+1\right)\Bigl({\nu^{y}} {u^{y}} \left(\left(u^{y}\right)^2+1\right)\\
  &\quad + {\nu^{x}} {u^{x}} \left(2 \left(u^{x}\right)^2+3 \left(u^{y}\right)^2+3\right)\\
  &\quad + {n}\Bigl(2 \left(u^{x}\right)^4+3 \left(u^{x}\right)^2 \left(\left(u^{y}\right)^2+1\right)+\left(\left(u^{y}\right)^2+1\right)^2\Bigr)\Bigr)\Bigr) \\
(A_t)_{2,3} &= \frac{1}{\left(u^{\tau}\right)^3 \left(\left(u^{y}\right)^2+1\right)^2}\Bigl({\Pi_{Q}} \left(u^{\tau}\right)^2 {u^{x}} {u^{y}} \left(\left(u^{y}\right)^2+1\right)^2 + {n} \left(u^{\tau}\right)^2 {u^{x}} {u^{y}} \left(\left(u^{y}\right)^2+1\right)^2 ({h}-{T})\\
  &\quad + {n} {T} \left(u^{\tau}\right)^2 {u^{x}} {u^{y}} \left(\left(u^{y}\right)^2+1\right)^2\\
  &\quad + \left(u^{\tau}\right)^2 {u^{x}} {u^{y}} \left(\left(u^{y}\right)^2+1\right)^2 ({P'(T)} {T}+{\Pi})\\
  &\quad + {h} \left(\left(u^{y}\right)^2+1\right)^2 \left({\nu^{x}} {u^{y}} \left(\left(u^{y}\right)^2+1\right)+{\nu^{y}} \left(\left(u^{x}\right)^3+{u^{x}}\right)\right)\\
  &\quad - {u^{y}} \left(\left(u^{y}\right)^2+1\right)\Bigl(- {\pi^{\eta}{}_{\eta}} \left(u^{\tau}\right)^2 {u^{x}} - {\pi^{yy}} \left(\left(u^{x}\right)^3+{u^{x}}\right)\\
  &\quad + {\pi^{xy}} {u^{y}} \left(2 \left(u^{x}\right)^2+\left(u^{y}\right)^2+1\right)\Bigr)\\
  &\quad - {u^{y}} \left(\left(u^{y}\right)^2+1\right)\Bigl(- {\pi_{Q}{}^{\eta}{}_{\eta}} \left(u^{\tau}\right)^2 {u^{x}} - {\pi_{Q}^{yy}} \left(\left(u^{x}\right)^3+{u^{x}}\right)\\
  &\quad + {\pi_{Q}^{xy}} {u^{y}} \left(2 \left(u^{x}\right)^2+\left(u^{y}\right)^2+1\right)\Bigr)\\
  &\quad + \left(u^{\tau}\right)^2\Bigl(2 {u^{x}} {u^{y}} \left({\pi^{\eta}{}_{\eta}} \left(u^{x}\right)^2+{\pi^{yy}} \left(u^{x}\right)^2+{\pi^{yy}}\right)\\
  &\quad + {\pi^{xy}} \left(\left(\left(u^{y}\right)^2+1\right)^2-2 \left(u^{x}\right)^2 \left(\left(u^{y}\right)^2-1\right)\right)\Bigr)\\
  &\quad + \left(u^{\tau}\right)^2\Bigl(2 {u^{x}} {u^{y}} \left({\pi_{Q}^{yy}} \left(u^{x}\right)^2+{\pi_{Q}^{yy}}+{\pi_{Q}{}^{\eta}{}_{\eta}} \left(u^{x}\right)^2\right)\\
  &\quad + {\pi_{Q}^{xy}} \left(\left(\left(u^{y}\right)^2+1\right)^2-2 \left(u^{x}\right)^2 \left(\left(u^{y}\right)^2-1\right)\right)\Bigr)\Bigr) \\
(A_t)_{2,4} &= -\frac{{u^{x}} \left(\left(u^{x}\right)^2+1\right)}{{u^{\tau}} \left(\left(u^{y}\right)^2+1\right)} \\
(A_t)_{2,5} &= -\frac{{u^{\tau}} {u^{x}}}{\left(u^{y}\right)^2+1} \\
(A_t)_{2,6} &= \frac{{u^{y}} \left(2 \left(u^{x}\right)^2+\left(u^{y}\right)^2+1\right)}{{u^{\tau}} \left(\left(u^{y}\right)^2+1\right)} \\
(A_t)_{2,7} &= {u^{\tau}} {u^{x}} \\
(A_t)_{2,8} &= {(\partial_{\alpha} n)} {h} {u^{\tau}} {u^{x}} \\
(A_t)_{2,9} &= \frac{{h} \left(2 \left(u^{x}\right)^2+\left(u^{y}\right)^2+1\right)}{{u^{\tau}}} \\
(A_t)_{2,10} &= \frac{{h} {u^{x}} {u^{y}}}{{u^{\tau}}} \\
(A_t)_{2,11} &= -\frac{{u^{x}} \left(\left(u^{x}\right)^2+1\right)}{{u^{\tau}} \left(\left(u^{y}\right)^2+1\right)} \\
(A_t)_{2,12} &= -\frac{{u^{\tau}} {u^{x}}}{\left(u^{y}\right)^2+1} \\
(A_t)_{2,13} &= \frac{{u^{y}} \left(2 \left(u^{x}\right)^2+\left(u^{y}\right)^2+1\right)}{{u^{\tau}} \left(\left(u^{y}\right)^2+1\right)} \\
(A_t)_{2,14} &= {u^{\tau}} {u^{x}} \\
(A_t)_{3,1} &= \frac{1}{{u^{\tau}}}\Bigl(\left(u^{\tau}\right)^2 {u^{y}} ({(\partial_{T} n)} {h}+{P''(T)} {T}+{P'(T)})\\
  &\quad + {h'(T)} \left({n} \left(u^{\tau}\right)^2 {u^{y}}+{\nu^{x}} {u^{x}} {u^{y}}+{\nu^{y}} \left(\left(u^{x}\right)^2+2 \left(u^{y}\right)^2+1\right)\right)\Bigr) \\
(A_t)_{3,2} &= \frac{1}{\left(u^{\tau}\right)^3}\Bigl({\pi^{xy}} + {\pi_{Q}^{xy}} + {h} {\nu^{y}} {u^{x}} + {h} {\nu^{y}} \left(u^{x}\right)^3 + {h} {\nu^{x}} {u^{y}} + {h} {n} {u^{x}} {u^{y}} + {\Pi} {u^{x}} {u^{y}} + {\Pi_{Q}} {u^{x}} {u^{y}}\\
  &\quad - {\pi^{yy}} {u^{x}} {u^{y}} - {\pi_{Q}^{yy}} {u^{x}} {u^{y}} + {P'(T)} {T} {u^{x}} {u^{y}} + {h} {n} \left(u^{x}\right)^3 {u^{y}} + {\Pi} \left(u^{x}\right)^3 {u^{y}} + {\Pi_{Q}} \left(u^{x}\right)^3 {u^{y}}\\
  &\quad + {P'(T)} {T} \left(u^{x}\right)^3 {u^{y}} + {\pi^{xy}} \left(u^{y}\right)^2 + {\pi_{Q}^{xy}} \left(u^{y}\right)^2 + {h} {\nu^{x}} \left(u^{y}\right)^3 + {h} {n} {u^{x}} \left(u^{y}\right)^3\\
  &\quad + {\Pi} {u^{x}} \left(u^{y}\right)^3 + {\Pi_{Q}} {u^{x}} \left(u^{y}\right)^3 + {P'(T)} {T} {u^{x}} \left(u^{y}\right)^3\Bigr) \\
(A_t)_{3,3} &= \frac{1}{\left(u^{\tau}\right)^3}\Bigl({\Pi_{Q}} + {\pi^{yy}} + {\pi_{Q}^{yy}} + {P'(T)} {T} + 2 {\Pi_{Q}} \left(u^{x}\right)^2 + {\pi^{yy}} \left(u^{x}\right)^2 + {\pi_{Q}^{yy}} \left(u^{x}\right)^2\\
  &\quad + 2 {P'(T)} {T} \left(u^{x}\right)^2 + {\Pi_{Q}} \left(u^{x}\right)^4 + {P'(T)} {T} \left(u^{x}\right)^4 - {\pi^{xy}} {u^{x}} {u^{y}} - {\pi_{Q}^{xy}} {u^{x}} {u^{y}} + 3 {\Pi_{Q}} \left(u^{y}\right)^2\\
  &\quad + 3 {P'(T)} {T} \left(u^{y}\right)^2 + 3 {\Pi_{Q}} \left(u^{x}\right)^2 \left(u^{y}\right)^2 + 3 {P'(T)} {T} \left(u^{x}\right)^2 \left(u^{y}\right)^2 + 2 {\Pi_{Q}} \left(u^{y}\right)^4\\
  &\quad + 2 {P'(T)} {T} \left(u^{y}\right)^4 + {\Pi} \left(\left(u^{x}\right)^4+\left(u^{x}\right)^2 \left(3 \left(u^{y}\right)^2+2\right)+2 \left(u^{y}\right)^4+3 \left(u^{y}\right)^2+1\right)\\
  &\quad + {h}\Bigl({\nu^{x}} \left(\left(u^{x}\right)^3+{u^{x}}\right) + {\nu^{y}} {u^{y}} \left(3 \left(u^{x}\right)^2+2 \left(u^{y}\right)^2+3\right)\\
  &\quad + {n} \left(\left(u^{x}\right)^4+\left(u^{x}\right)^2 \left(3 \left(u^{y}\right)^2+2\right)+2 \left(u^{y}\right)^4+3 \left(u^{y}\right)^2+1\right)\Bigr)\Bigr) \\
(A_t)_{3,4} &= \frac{{u^{y}}}{{u^{\tau}}} \\
(A_t)_{3,6} &= \frac{{u^{x}}}{{u^{\tau}}} \\
(A_t)_{3,7} &= {u^{\tau}} {u^{y}} \\
(A_t)_{3,8} &= {(\partial_{\alpha} n)} {h} {u^{\tau}} {u^{y}} \\
(A_t)_{3,9} &= \frac{{h} {u^{x}} {u^{y}}}{{u^{\tau}}} \\
(A_t)_{3,10} &= \frac{{h} \left(\left(u^{x}\right)^2+2 \left(u^{y}\right)^2+1\right)}{{u^{\tau}}} \\
(A_t)_{3,11} &= \frac{{u^{y}}}{{u^{\tau}}} \\
(A_t)_{3,13} &= \frac{{u^{x}}}{{u^{\tau}}} \\
(A_t)_{3,14} &= {u^{\tau}} {u^{y}} \\
(A_t)_{4,2} &= \frac{3 {\delta_{\pi\pi}} {\pi^{yy}} {u^{x}}-6 {\pi^{xy}} {\tau_{\pi}} {u^{y}} \left(\left(u^{y}\right)^2+1\right)-2 {u^{x}} \left(-3 {\pi^{yy}} {\tau_{\pi}} \left(u^{y}\right)^2+\left(u^{y}\right)^2 {\eta}+{\eta}\right)}{3 {u^{\tau}}} \\
(A_t)_{4,3} &= \frac{{u^{y}} \left(3 {\delta_{\pi\pi}} {\pi^{yy}}+6 {\pi^{xy}} {\tau_{\pi}} {u^{x}} {u^{y}}-6 {\pi^{yy}} {\tau_{\pi}} \left(\left(u^{x}\right)^2+1\right)+6 \left(u^{x}\right)^2 {\eta}+4 \left(u^{y}\right)^2 {\eta}+4 {\eta}\right)}{3 {u^{\tau}}} \\
(A_t)_{4,4} &= {\tau_{\pi}} {u^{\tau}} \\
(A_t)_{5,2} &= \frac{{u^{x}} (3 {\delta_{\pi\pi}} {\pi^{\eta}{}_{\eta}}-2 {\eta})}{3 {\tau}^2 {u^{\tau}}} \\
(A_t)_{5,3} &= \frac{{u^{y}} (3 {\delta_{\pi\pi}} {\pi^{\eta}{}_{\eta}}-2 {\eta})}{3 {\tau}^2 {u^{\tau}}} \\
(A_t)_{5,5} &= \frac{{\tau_{\pi}} {u^{\tau}}}{{\tau}^2} \\
(A_t)_{6,2} &= \frac{1}{3 {u^{\tau}}}\Bigl(3 {\delta_{\pi\pi}} {\pi^{xy}} {u^{x}} - 3 {\pi^{xy}} {\tau_{\pi}} {u^{x}} \left(2 \left(u^{y}\right)^2+1\right)\\
  &\quad + {u^{y}} \left(3 {\pi^{\eta}{}_{\eta}} {\tau_{\pi}} \left(u^{\tau}\right)^2+3 {\pi^{yy}} \left(2 {\tau_{\pi}} \left(u^{x}\right)^2+{\tau_{\pi}}\right)+{\eta} \left(\left(u^{x}\right)^2+3 \left(u^{y}\right)^2+3\right)\right)\Bigr) \\
(A_t)_{6,3} &= \frac{1}{3 {u^{\tau}} \left(\left(u^{y}\right)^2+1\right)}\Bigl(- 3 {\pi^{xy}} {\tau_{\pi}} {u^{y}} - 3 {\pi^{\eta}{}_{\eta}} {\tau_{\pi}} {u^{x}} \left(u^{y}\right)^2 - 3 {\pi^{\eta}{}_{\eta}} {\tau_{\pi}} \left(u^{x}\right)^3 \left(u^{y}\right)^2 - 3 {\pi^{xy}} {\tau_{\pi}} \left(u^{y}\right)^3\\
  &\quad + 6 {\pi^{xy}} {\tau_{\pi}} \left(u^{x}\right)^2 \left(u^{y}\right)^3 - 3 {\pi^{\eta}{}_{\eta}} {\tau_{\pi}} {u^{x}} \left(u^{y}\right)^4 + 3 {\delta_{\pi\pi}} {\pi^{xy}} {u^{y}} \left(\left(u^{y}\right)^2+1\right)\\
  &\quad - 3 {\pi^{yy}} {\tau_{\pi}} {u^{x}} \left(\left(u^{x}\right)^2+1\right) \left(2 \left(u^{y}\right)^2+1\right) + 3 {u^{x}} {\eta} + 3 \left(u^{x}\right)^3 {\eta}\\
  &\quad + 4 {u^{x}} \left(u^{y}\right)^2 {\eta} + 3 \left(u^{x}\right)^3 \left(u^{y}\right)^2 {\eta} + {u^{x}} \left(u^{y}\right)^4 {\eta}\Bigr) \\
(A_t)_{6,6} &= {\tau_{\pi}} {u^{\tau}} \\
(A_t)_{7,2} &= \frac{{u^{x}} {\zeta}}{{u^{\tau}}} \\
(A_t)_{7,3} &= \frac{{u^{y}} {\zeta}}{{u^{\tau}}} \\
(A_t)_{7,7} &= {\tau_{\Pi}} {u^{\tau}} \\
(A_t)_{8,1} &= {(\partial_{T} n)} {u^{\tau}} \\
(A_t)_{8,2} &= \frac{{n} \left(u^{\tau}\right)^2 {u^{x}}+{\nu^{x}} \left(\left(u^{y}\right)^2+1\right)-{\nu^{y}} {u^{x}} {u^{y}}}{\left(u^{\tau}\right)^3} \\
(A_t)_{8,3} &= \frac{{n} \left(u^{\tau}\right)^2 {u^{y}}-{\nu^{x}} {u^{x}} {u^{y}}+{\nu^{y}} \left(\left(u^{x}\right)^2+1\right)}{\left(u^{\tau}\right)^3} \\
(A_t)_{8,8} &= {(\partial_{\alpha} n)} {u^{\tau}} \\
(A_t)_{8,9} &= \frac{{u^{x}}}{{u^{\tau}}} \\
(A_t)_{8,10} &= \frac{{u^{y}}}{{u^{\tau}}} \\
(A_t)_{9,1} &= \frac{{D_{s}} {u^{\tau}} ({h'(T)} {\nu^{x}}+{u^{x}} ({(\partial_{T} n)} {T}+{n}))}{{T}} \\
(A_t)_{9,2} &= \frac{1}{{T} \left(u^{\tau}\right)^3 \left(\left(u^{y}\right)^2+1\right)}{D_{s}}\Bigl(- {\pi_{Q}^{yy}} - {\pi_{Q}{}^{\eta}{}_{\eta}} - 2 {\pi_{Q}^{yy}} \left(u^{x}\right)^2 - 2 {\pi_{Q}{}^{\eta}{}_{\eta}} \left(u^{x}\right)^2 - {\pi_{Q}^{yy}} \left(u^{x}\right)^4\\
  &\quad - {\pi_{Q}{}^{\eta}{}_{\eta}} \left(u^{x}\right)^4 + 2 {\pi_{Q}^{xy}} {u^{x}} {u^{y}} + {\pi_{Q}^{xy}} \left(u^{x}\right)^3 {u^{y}} - {\pi_{Q}^{yy}} \left(u^{y}\right)^2\\
  &\quad - 2 {\pi_{Q}{}^{\eta}{}_{\eta}} \left(u^{y}\right)^2 - {\pi_{Q}^{yy}} \left(u^{x}\right)^2 \left(u^{y}\right)^2 - {\pi_{Q}{}^{\eta}{}_{\eta}} \left(u^{x}\right)^2 \left(u^{y}\right)^2\\
  &\quad + {\pi_{Q}^{yy}} \left(u^{x}\right)^4 \left(u^{y}\right)^2 + {\pi_{Q}{}^{\eta}{}_{\eta}} \left(u^{x}\right)^4 \left(u^{y}\right)^2 + {\pi_{Q}^{xy}} {u^{x}} \left(u^{y}\right)^3\\
  &\quad - {\pi_{Q}^{xy}} \left(u^{x}\right)^3 \left(u^{y}\right)^3 - {\pi_{Q}{}^{\eta}{}_{\eta}} \left(u^{y}\right)^4 + {\pi_{Q}^{yy}} \left(u^{x}\right)^2 \left(u^{y}\right)^4\\
  &\quad + {\pi_{Q}{}^{\eta}{}_{\eta}} \left(u^{x}\right)^2 \left(u^{y}\right)^4 - {\pi_{Q}^{xy}} {u^{x}} \left(u^{y}\right)^5 + {\Pi_{Q}} \left(u^{\tau}\right)^4 \left(\left(u^{y}\right)^2+1\right)\\
  &\quad + {h} \left(u^{\tau}\right)^2 \left(\left(u^{y}\right)^2+1\right)\Bigl({n} \left(u^{\tau}\right)^2+{\nu^{x}} \left({u^{x}}-{u^{x}} \left(u^{y}\right)^2\right)+{\nu^{y}} \left(\left(u^{x}\right)^2+1\right) {u^{y}}\Bigr)\Bigr) \\
(A_t)_{9,3} &= \frac{1}{{T} \left(u^{\tau}\right)^3 \left(\left(u^{y}\right)^2+1\right)^2}{D_{s}}\Bigl(- {h} \left(u^{\tau}\right)^2 \left(\left(u^{x}\right)^2+1\right) \left(\left(u^{y}\right)^2+1\right)^2 ({\nu^{y}} {u^{x}}-{\nu^{x}} {u^{y}})\\
  &\quad + {\pi_{Q}^{xy}}\Bigl(\left(u^{x}\right)^4 \left(\left(u^{y}\right)^2-1\right)^2 + \left(\left(u^{y}\right)^2+1\right)^2\\
  &\quad + \left(u^{x}\right)^2 \left(\left(u^{y}\right)^6-2 \left(u^{y}\right)^4-\left(u^{y}\right)^2+2\right)\Bigr)\\
  &\quad + {u^{x}} {u^{y}}\Bigl(- {\pi_{Q}^{yy}} \left(\left(u^{x}\right)^2+1\right)\Bigl(\left(u^{x}\right)^2 \left(\left(u^{y}\right)^2-1\right)+\left(u^{y}\right)^4-\left(u^{y}\right)^2-2\Bigr)\\
  &\quad + {\pi_{Q}{}^{\eta}{}_{\eta}}\Bigl(- \left(u^{x}\right)^4 \left(\left(u^{y}\right)^2-1\right) + \left(\left(u^{y}\right)^2+1\right)^2\\
  &\quad + \left(u^{x}\right)^2 \left(-\left(u^{y}\right)^4+\left(u^{y}\right)^2+2\right)\Bigr)\Bigr)\Bigr) \\
(A_t)_{9,8} &= {(\partial_{\alpha} n)} {D_{s}} {u^{\tau}} {u^{x}} \\
(A_t)_{9,9} &= \frac{{D_{s}} {h} {u^{\tau}}}{{T}} \\
(A_t)_{9,11} &= -\frac{{D_{s}} {u^{x}} \left(\left(u^{x}\right)^2+1\right)}{{T} {u^{\tau}} \left(\left(u^{y}\right)^2+1\right)} \\
(A_t)_{9,12} &= -\frac{{D_{s}} {u^{\tau}} {u^{x}}}{{T} \left(u^{y}\right)^2+{T}} \\
(A_t)_{9,13} &= \frac{{D_{s}} \left(2 \left(u^{x}\right)^2 {u^{y}}+\left(u^{y}\right)^3+{u^{y}}\right)}{{T} {u^{\tau}} \left(\left(u^{y}\right)^2+1\right)} \\
(A_t)_{9,14} &= \frac{{D_{s}} {u^{\tau}} {u^{x}}}{{T}} \\
(A_t)_{10,1} &= \frac{{D_{s}} {u^{\tau}} ({h'(T)} {\nu^{y}}+{u^{y}} ({(\partial_{T} n)} {T}+{n}))}{{T}} \\
(A_t)_{10,2} &= \frac{1}{{T} \left(u^{\tau}\right)^3}{D_{s}}\Bigl({h} \left(u^{\tau}\right)^2 \left(\left(u^{y}\right)^2+1\right) ({\nu^{y}} {u^{x}}-{\nu^{x}} {u^{y}}) - {\pi_{Q}^{xy}} \left(\left(u^{x}\right)^2 \left(u^{y}\right)^2+\left(u^{y}\right)^4-1\right)\\
  &\quad + {u^{x}} {u^{y}} \left({\pi_{Q}^{yy}} \left(\left(u^{x}\right)^2+\left(u^{y}\right)^2\right)+{\pi_{Q}{}^{\eta}{}_{\eta}} \left(u^{\tau}\right)^2\right)\Bigr) \\
(A_t)_{10,3} &= \frac{1}{{T} \left(u^{\tau}\right)^3 \left(\left(u^{y}\right)^2+1\right)}{D_{s}}\Bigl({\pi_{Q}^{yy}} + {\pi_{Q}^{yy}} \left(u^{x}\right)^2 - 2 {\pi_{Q}^{xy}} {u^{x}} {u^{y}} - {\pi_{Q}^{xy}} \left(u^{x}\right)^3 {u^{y}} - {\pi_{Q}^{yy}} \left(u^{x}\right)^2 \left(u^{y}\right)^2\\
  &\quad - {\pi_{Q}{}^{\eta}{}_{\eta}} \left(u^{x}\right)^2 \left(u^{y}\right)^2 - {\pi_{Q}^{yy}} \left(u^{x}\right)^4 \left(u^{y}\right)^2 - {\pi_{Q}{}^{\eta}{}_{\eta}} \left(u^{x}\right)^4 \left(u^{y}\right)^2\\
  &\quad - {\pi_{Q}^{xy}} {u^{x}} \left(u^{y}\right)^3 + {\pi_{Q}^{xy}} \left(u^{x}\right)^3 \left(u^{y}\right)^3 - {\pi_{Q}^{yy}} \left(u^{y}\right)^4 - {\pi_{Q}^{yy}} \left(u^{x}\right)^2 \left(u^{y}\right)^4\\
  &\quad - {\pi_{Q}{}^{\eta}{}_{\eta}} \left(u^{x}\right)^2 \left(u^{y}\right)^4 + {\pi_{Q}^{xy}} {u^{x}} \left(u^{y}\right)^5 + {\Pi_{Q}} \left(u^{\tau}\right)^4 \left(\left(u^{y}\right)^2+1\right)\\
  &\quad + {h} \left(u^{\tau}\right)^2 \left(\left(u^{y}\right)^2+1\right)\Bigl({n} \left(u^{\tau}\right)^2+{\nu^{x}} {u^{x}} \left(\left(u^{y}\right)^2+1\right)-{\nu^{y}} \left(\left(u^{x}\right)^2-1\right) {u^{y}}\Bigr)\Bigr) \\
(A_t)_{10,8} &= {(\partial_{\alpha} n)} {D_{s}} {u^{\tau}} {u^{y}} \\
(A_t)_{10,10} &= \frac{{D_{s}} {h} {u^{\tau}}}{{T}} \\
(A_t)_{10,11} &= \frac{{D_{s}} {u^{y}}}{{T} {u^{\tau}}} \\
(A_t)_{10,13} &= \frac{{D_{s}} {u^{x}}}{{T} {u^{\tau}}} \\
(A_t)_{10,14} &= \frac{{D_{s}} {u^{\tau}} {u^{y}}}{{T}} \\
(A_t)_{11,1} &= \frac{1}{6 {T} {u^{\tau}}}{D_{s}}\Bigl(3 {\mathcal{C}'(T)} {\pi_{Q}^{yy}} \left(u^{\tau}\right)^2 - 2 {\mathcal{B}'(T)} {\nu^{x}} {u^{x}} \left(\left(u^{y}\right)^2+1\right)\\
  &\quad + 2 {\mathcal{B}'(T)} {\nu^{y}} {u^{y}} \left(3 \left(u^{x}\right)^2+2 \left(u^{y}\right)^2+2\right)\Bigr) \\
(A_t)_{11,2} &= -\frac{1}{6 {T} \left(u^{\tau}\right)^3}{D_{s}}\Bigl(2 {\mathcal{A}} \left(u^{\tau}\right)^2 \left(\left(u^{y}\right)^2+1\right) \left({\nu^{x}} \left(u^{y}\right)^2+{\nu^{x}}-{\nu^{y}} {u^{x}} {u^{y}}\right)\\
  &\quad + 2 {\mathcal{B}} \left(\left(3 \left(u^{x}\right)^2+4\right) \left(u^{y}\right)^2+3 \left(u^{y}\right)^4+1\right) \left({\nu^{x}} \left(u^{y}\right)^2+{\nu^{x}}-{\nu^{y}} {u^{x}} {u^{y}}\right)\\
  &\quad + \left(u^{\tau}\right)^2\Bigl(2 {I_{31}} {u^{x}} \left(\left(u^{y}\right)^2+1\right)\\
  &\quad + {\mathcal{C}}\Bigl(2 {\Pi_{Q}} {u^{x}} \left(\left(u^{y}\right)^2+1\right)\\
  &\quad - 2 \left(\left(u^{y}\right)^2+1\right) ({\pi_{Q}{}^{\eta}{}_{\eta}} {u^{x}}-4 {\pi_{Q}^{xy}} {u^{y}})\\
  &\quad - {\pi_{Q}^{yy}} {u^{x}} \left(8 \left(u^{y}\right)^2+5\right)\Bigr)\Bigr)\Bigr) \\
(A_t)_{11,3} &= \frac{1}{6 {T} \left(u^{\tau}\right)^3}{D_{s}}\Bigl(2 {\mathcal{B}} {\nu^{x}} {u^{x}} {u^{y}} \left(\left(3 \left(u^{x}\right)^2+4\right) \left(u^{y}\right)^2+3 \left(u^{y}\right)^4+1\right)\\
  &\quad - 2{\mathcal{A}} \left(u^{\tau}\right)^2\Bigl({\nu^{y}} \left(u^{x}\right)^2 \left(\left(u^{y}\right)^2-2\right) - 2 {\nu^{y}} \left(\left(u^{y}\right)^2+1\right)\\
  &\quad - {\nu^{x}} {u^{x}} {u^{y}} \left(\left(u^{y}\right)^2+1\right)\Bigr)\\
  &\quad - 2{\mathcal{B}} {\nu^{y}}\Bigl(3 \left(u^{x}\right)^4 \left(\left(u^{y}\right)^2-1\right) - 2 \left(\left(u^{y}\right)^2+1\right)\\
  &\quad + \left(u^{x}\right)^2 \left(3 \left(u^{y}\right)^4+\left(u^{y}\right)^2-5\right)\Bigr)\\
  &\quad + \left(u^{\tau}\right)^2\Bigl(2 {I_{31}} {u^{y}} \left(3 \left(u^{x}\right)^2+2 \left(u^{y}\right)^2+2\right) + 4 {\mathcal{C}} {\pi_{Q}^{xy}} \left(2 {u^{x}} \left(u^{y}\right)^2+{u^{x}}\right)\\
  &\quad + {\mathcal{C}} {u^{y}}\Bigl({\pi_{Q}^{yy}} - 8 {\pi_{Q}^{yy}} \left(u^{x}\right)^2 - 2 {\pi_{Q}{}^{\eta}{}_{\eta}} \left(u^{x}\right)^2\\
  &\quad + {\Pi_{Q}} \left(6 \left(u^{x}\right)^2+4 \left(u^{y}\right)^2+4\right)\Bigr)\Bigr)\Bigr) \\
(A_t)_{11,9} &= -\frac{{\mathcal{B}} {D_{s}} {u^{x}} \left(\left(u^{y}\right)^2+1\right)}{3 {T} {u^{\tau}}} \\
(A_t)_{11,10} &= \frac{{\mathcal{B}} {D_{s}} {u^{y}} \left(3 \left(u^{x}\right)^2+2 \left(u^{y}\right)^2+2\right)}{3 {T} {u^{\tau}}} \\
(A_t)_{11,11} &= \frac{{\mathcal{C}} {D_{s}} {u^{\tau}}}{2 {T}} \\
(A_t)_{12,1} &= \frac{{D_{s}} \left(3 {\mathcal{C}'(T)} {\pi_{Q}{}^{\eta}{}_{\eta}} \left(u^{\tau}\right)^2-2 {\mathcal{B}'(T)} ({\nu^{x}} {u^{x}}+{\nu^{y}} {u^{y}})\right)}{6 {T} {u^{\tau}}} \\
(A_t)_{12,2} &= -\frac{1}{6 {T} \left(u^{\tau}\right)^3}{D_{s}}\Bigl(2 {\mathcal{B}} \left({\nu^{x}} \left(u^{y}\right)^2+{\nu^{x}}-{\nu^{y}} {u^{x}} {u^{y}}\right) + 2 {\mathcal{A}} \left(u^{\tau}\right)^2 \left({\nu^{x}} \left(u^{y}\right)^2+{\nu^{x}}-{\nu^{y}} {u^{x}} {u^{y}}\right)\\
  &\quad + \left(u^{\tau}\right)^2 ({\mathcal{C}} (2 {\Pi_{Q}} {u^{x}}+2 {\pi_{Q}^{xy}} {u^{y}}-2 {\pi_{Q}^{yy}} {u^{x}}-5 {\pi_{Q}{}^{\eta}{}_{\eta}} {u^{x}})+2 {I_{31}} {u^{x}})\Bigr) \\
(A_t)_{12,3} &= -\frac{1}{6 {T} \left(u^{\tau}\right)^3 \left(\left(u^{y}\right)^2+1\right)}{D_{s}}\Bigl(2 {\mathcal{B}} \left(\left(u^{y}\right)^2+1\right) \left(-{\nu^{x}} {u^{x}} {u^{y}}+{\nu^{y}} \left(u^{x}\right)^2+{\nu^{y}}\right)\\
  &\quad + 2 {\mathcal{A}} \left(u^{\tau}\right)^2 \left(\left(u^{y}\right)^2+1\right) \left(-{\nu^{x}} {u^{x}} {u^{y}}+{\nu^{y}} \left(u^{x}\right)^2+{\nu^{y}}\right)\\
  &\quad + \left(u^{\tau}\right)^2\Bigl(- 2 {\mathcal{C}} {\pi_{Q}^{xy}} {u^{x}} \left(\left(u^{y}\right)^2-1\right) + 2 {I_{31}} {u^{y}} \left(\left(u^{y}\right)^2+1\right)\\
  &\quad + {\mathcal{C}} {u^{y}}\Bigl(2 {\pi_{Q}^{yy}} \left(\left(u^{x}\right)^2+1\right)\\
  &\quad + {\pi_{Q}{}^{\eta}{}_{\eta}} \left(2 \left(u^{x}\right)^2-3 \left(u^{y}\right)^2-3\right) + 2 {\Pi_{Q}} \left(\left(u^{y}\right)^2+1\right)\Bigr)\Bigr)\Bigr) \\
(A_t)_{12,9} &= -\frac{{\mathcal{B}} {D_{s}} {u^{x}}}{3 {T} {u^{\tau}}} \\
(A_t)_{12,10} &= -\frac{{\mathcal{B}} {D_{s}} {u^{y}}}{3 {T} {u^{\tau}}} \\
(A_t)_{12,12} &= \frac{{\mathcal{C}} {D_{s}} {u^{\tau}}}{2 {T}} \\
(A_t)_{13,1} &= \frac{1}{6 {T} {u^{\tau}}}{D_{s}}\Bigl(3 {\mathcal{C}'(T)} {\pi_{Q}^{xy}} \left(u^{\tau}\right)^2 + {\mathcal{B}'(T)} {\nu^{y}} {u^{x}} \left(3 \left(u^{x}\right)^2+\left(u^{y}\right)^2+3\right)\\
  &\quad + {\mathcal{B}'(T)} {\nu^{x}} {u^{y}} \left(\left(u^{x}\right)^2+3 \left(u^{y}\right)^2+3\right)\Bigr) \\
(A_t)_{13,2} &= \frac{1}{6 {T} \left(u^{\tau}\right)^3}{D_{s}}\Bigl({\mathcal{A}} \left(u^{\tau}\right)^2\Bigl(3 {\nu^{y}} \left(\left(u^{y}\right)^2+1\right) - 2 {\nu^{x}} {u^{x}} {u^{y}} \left(\left(u^{y}\right)^2+1\right)\\
  &\quad + {\nu^{y}} \left(u^{x}\right)^2 \left(2 \left(u^{y}\right)^2+3\right)\Bigr)\\
  &\quad + \left(u^{\tau}\right)^2\Bigl({\mathcal{C}} {\pi_{Q}^{xy}} {u^{x}} \left(3-8 \left(u^{y}\right)^2\right) + {I_{31}} {u^{y}} \left(\left(u^{x}\right)^2+3 \left(u^{y}\right)^2+3\right)\\
  &\quad + {\mathcal{C}} {u^{y}}\Bigl({\pi_{Q}^{yy}} \left(8 \left(u^{x}\right)^2+6\right) + {\Pi_{Q}} \left(\left(u^{x}\right)^2+3 \left(u^{y}\right)^2+3\right)\\
  &\quad + {\pi_{Q}{}^{\eta}{}_{\eta}} \left(5 \left(u^{x}\right)^2+3 \left(u^{y}\right)^2+3\right)\Bigr)\Bigr)\\
  &\quad + {\mathcal{B}}\Bigl(- 2 {\nu^{x}} {u^{x}} {u^{y}} \left(\left(u^{y}\right)^2+1\right) \left(3 \left(u^{x}\right)^2+3 \left(u^{y}\right)^2+4\right)\\
  &\quad + {\nu^{y}}\Bigl(3 \left(\left(u^{y}\right)^2+1\right)^2 + \left(u^{x}\right)^4 \left(6 \left(u^{y}\right)^2+3\right)\\
  &\quad + 2 \left(u^{x}\right)^2 \left(3 \left(u^{y}\right)^4+7 \left(u^{y}\right)^2+3\right)\Bigr)\Bigr)\Bigr) \\
(A_t)_{13,3} &= \frac{1}{6 {T} \left(u^{\tau}\right)^3 \left(\left(u^{y}\right)^2+1\right)}{D_{s}}\Bigl({\mathcal{A}} \left(u^{\tau}\right)^2 \left(\left(u^{y}\right)^2+1\right)\Bigl(- 2 {\nu^{y}} {u^{x}} \left(\left(u^{x}\right)^2+1\right) {u^{y}} + 3 {\nu^{x}} \left(\left(u^{y}\right)^2+1\right)\\
  &\quad + {\nu^{x}} \left(u^{x}\right)^2 \left(2 \left(u^{y}\right)^2+3\right)\Bigr)\\
  &\quad + \left(u^{\tau}\right)^2\Bigl({I_{31}} {u^{x}} \left(\left(u^{y}\right)^2+1\right) \left(3 \left(u^{x}\right)^2+\left(u^{y}\right)^2+3\right)\\
  &\quad + {\mathcal{C}}\Bigl(- 3 {\pi_{Q}{}^{\eta}{}_{\eta}} {u^{x}} - 3 {\pi_{Q}{}^{\eta}{}_{\eta}} \left(u^{x}\right)^3 + 3 {\pi_{Q}^{xy}} {u^{y}} + 4 {\pi_{Q}^{xy}} \left(u^{x}\right)^2 {u^{y}}\\
  &\quad - 6 {\pi_{Q}{}^{\eta}{}_{\eta}} {u^{x}} \left(u^{y}\right)^2 - 5 {\pi_{Q}{}^{\eta}{}_{\eta}} \left(u^{x}\right)^3 \left(u^{y}\right)^2 + 3 {\pi_{Q}^{xy}} \left(u^{y}\right)^3\\
  &\quad + 8 {\pi_{Q}^{xy}} \left(u^{x}\right)^2 \left(u^{y}\right)^3 - 3 {\pi_{Q}{}^{\eta}{}_{\eta}} {u^{x}} \left(u^{y}\right)^4\\
  &\quad + {\Pi_{Q}} {u^{x}} \left(\left(u^{y}\right)^2+1\right) \left(3 \left(u^{x}\right)^2+\left(u^{y}\right)^2+3\right)\\
  &\quad - 2 {\pi_{Q}^{yy}} {u^{x}} \left(\left(u^{x}\right)^2+1\right) \left(4 \left(u^{y}\right)^2+3\right)\Bigr)\Bigr)\\
  &\quad + {\mathcal{B}} \left(\left(u^{y}\right)^2+1\right)\Bigl(- 2 {\nu^{y}} {u^{x}} \left(\left(u^{x}\right)^2+1\right) {u^{y}} \left(3 \left(u^{x}\right)^2+3 \left(u^{y}\right)^2+4\right)\\
  &\quad + {\nu^{x}}\Bigl(3 \left(\left(u^{y}\right)^2+1\right)^2 + \left(u^{x}\right)^4 \left(6 \left(u^{y}\right)^2+3\right)\\
  &\quad + 2 \left(u^{x}\right)^2 \left(3 \left(u^{y}\right)^4+7 \left(u^{y}\right)^2+3\right)\Bigr)\Bigr)\Bigr) \\
(A_t)_{13,9} &= \frac{{\mathcal{B}} {D_{s}} {u^{y}} \left(\left(u^{x}\right)^2+3 \left(u^{y}\right)^2+3\right)}{6 {T} {u^{\tau}}} \\
(A_t)_{13,10} &= \frac{{\mathcal{B}} {D_{s}} {u^{x}} \left(3 \left(u^{x}\right)^2+\left(u^{y}\right)^2+3\right)}{6 {T} {u^{\tau}}} \\
(A_t)_{13,13} &= \frac{{\mathcal{C}} {D_{s}} {u^{\tau}}}{2 {T}} \\
(A_t)_{14,1} &= \frac{{D_{s}} \left(5 {\mathcal{B}'(T)} ({\nu^{x}} {u^{x}}+{\nu^{y}} {u^{y}})+3 {\mathcal{C}'(T)} {\Pi_{Q}} \left(u^{\tau}\right)^2+3 {(\partial_{T} I_{31})} \left(u^{\tau}\right)^2\right)}{6 {T} {u^{\tau}}} \\
(A_t)_{14,2} &= \frac{1}{6 {T} \left(u^{\tau}\right)^3}{D_{s}}\Bigl(5 {\mathcal{B}} \left({\nu^{x}} \left(u^{y}\right)^2+{\nu^{x}}-{\nu^{y}} {u^{x}} {u^{y}}\right) + 2 {\mathcal{A}} \left(u^{\tau}\right)^2 \left({\nu^{x}} \left(u^{y}\right)^2+{\nu^{x}}-{\nu^{y}} {u^{x}} {u^{y}}\right)\\
  &\quad + \left(u^{\tau}\right)^2 ({\mathcal{C}} (5 {\Pi_{Q}} {u^{x}}+2 {\pi_{Q}^{xy}} {u^{y}}-2 {\pi_{Q}^{yy}} {u^{x}}-2 {\pi_{Q}{}^{\eta}{}_{\eta}} {u^{x}})+5 {I_{31}} {u^{x}})\Bigr) \\
(A_t)_{14,3} &= \frac{1}{6 {T} \left(u^{\tau}\right)^3 \left(\left(u^{y}\right)^2+1\right)}{D_{s}}\Bigl(5 {\mathcal{B}} \left(\left(u^{y}\right)^2+1\right) \left(-{\nu^{x}} {u^{x}} {u^{y}}+{\nu^{y}} \left(u^{x}\right)^2+{\nu^{y}}\right)\\
  &\quad + 2 {\mathcal{A}} \left(u^{\tau}\right)^2 \left(\left(u^{y}\right)^2+1\right) \left(-{\nu^{x}} {u^{x}} {u^{y}}+{\nu^{y}} \left(u^{x}\right)^2+{\nu^{y}}\right)\\
  &\quad + \left(u^{\tau}\right)^2\Bigl(5 {I_{31}} {u^{y}} \left(\left(u^{y}\right)^2+1\right)\\
  &\quad + {\mathcal{C}}\Bigl(2 {u^{y}} \left({\pi_{Q}^{yy}} \left(u^{x}\right)^2+{\pi_{Q}^{yy}}+{\pi_{Q}{}^{\eta}{}_{\eta}} \left(u^{x}\right)^2\right)\\
  &\quad - 2 {\pi_{Q}^{xy}} {u^{x}} \left(\left(u^{y}\right)^2-1\right) + 5 {\Pi_{Q}} {u^{y}} \left(\left(u^{y}\right)^2+1\right)\Bigr)\Bigr)\Bigr) \\
(A_t)_{14,8} &= \frac{{(\partial_{\alpha} I_{31})} {D_{s}} {u^{\tau}}}{2 {T}} \\
(A_t)_{14,9} &= \frac{5 {\mathcal{B}} {D_{s}} {u^{x}}}{6 {T} {u^{\tau}}} \\
(A_t)_{14,10} &= \frac{5 {\mathcal{B}} {D_{s}} {u^{y}}}{6 {T} {u^{\tau}}} \\
(A_t)_{14,14} &= \frac{{\mathcal{C}} {D_{s}} {u^{\tau}}}{2 {T}} \\
(A_x)_{1,1} &= \frac{1}{{u^{\tau}}}\Bigl(\left(u^{\tau}\right)^2 {u^{x}} ({(\partial_{T} n)} {h}+{P''(T)} {T}+{P'(T)})\\
  &\quad + {h'(T)} \left({n} \left(u^{\tau}\right)^2 {u^{x}}+{\nu^{x}} \left(2 \left(u^{x}\right)^2+\left(u^{y}\right)^2+1\right)+{\nu^{y}} {u^{x}} {u^{y}}\right)\Bigr) \\
(A_x)_{1,2} &= \frac{1}{\left(u^{\tau}\right)^3 \left(\left(u^{y}\right)^2+1\right)}\Bigl({\Pi_{Q}} - {\pi^{\eta}{}_{\eta}} - {\pi^{yy}} - {\pi_{Q}^{yy}} - {\pi_{Q}{}^{\eta}{}_{\eta}} + {P'(T)} {T} + 3 {\Pi_{Q}} \left(u^{x}\right)^2\\
  &\quad - 3 {\pi^{\eta}{}_{\eta}} \left(u^{x}\right)^2 - 3 {\pi^{yy}} \left(u^{x}\right)^2 - 3 {\pi_{Q}^{yy}} \left(u^{x}\right)^2 - 3 {\pi_{Q}{}^{\eta}{}_{\eta}} \left(u^{x}\right)^2\\
  &\quad + 3 {P'(T)} {T} \left(u^{x}\right)^2 + 2 {\Pi_{Q}} \left(u^{x}\right)^4 - 2 {\pi^{\eta}{}_{\eta}} \left(u^{x}\right)^4 - 2 {\pi^{yy}} \left(u^{x}\right)^4\\
  &\quad - 2 {\pi_{Q}^{yy}} \left(u^{x}\right)^4 - 2 {\pi_{Q}{}^{\eta}{}_{\eta}} \left(u^{x}\right)^4 + 2 {P'(T)} {T} \left(u^{x}\right)^4 + 3 {\pi^{xy}} {u^{x}} {u^{y}}\\
  &\quad + 3 {\pi_{Q}^{xy}} {u^{x}} {u^{y}} + 2 {\pi^{xy}} \left(u^{x}\right)^3 {u^{y}} + 2 {\pi_{Q}^{xy}} \left(u^{x}\right)^3 {u^{y}} + 3 {\Pi_{Q}} \left(u^{y}\right)^2\\
  &\quad - 2 {\pi^{\eta}{}_{\eta}} \left(u^{y}\right)^2 - {\pi^{yy}} \left(u^{y}\right)^2 - {\pi_{Q}^{yy}} \left(u^{y}\right)^2 - 2 {\pi_{Q}{}^{\eta}{}_{\eta}} \left(u^{y}\right)^2\\
  &\quad + 3 {P'(T)} {T} \left(u^{y}\right)^2 + 6 {\Pi_{Q}} \left(u^{x}\right)^2 \left(u^{y}\right)^2 - 3 {\pi^{\eta}{}_{\eta}} \left(u^{x}\right)^2 \left(u^{y}\right)^2\\
  &\quad - 3 {\pi^{yy}} \left(u^{x}\right)^2 \left(u^{y}\right)^2 - 3 {\pi_{Q}^{yy}} \left(u^{x}\right)^2 \left(u^{y}\right)^2 - 3 {\pi_{Q}{}^{\eta}{}_{\eta}} \left(u^{x}\right)^2 \left(u^{y}\right)^2\\
  &\quad + 6 {P'(T)} {T} \left(u^{x}\right)^2 \left(u^{y}\right)^2 + 2 {\Pi_{Q}} \left(u^{x}\right)^4 \left(u^{y}\right)^2 + 2 {P'(T)} {T} \left(u^{x}\right)^4 \left(u^{y}\right)^2\\
  &\quad + 3 {\pi^{xy}} {u^{x}} \left(u^{y}\right)^3 + 3 {\pi_{Q}^{xy}} {u^{x}} \left(u^{y}\right)^3 + 3 {\Pi_{Q}} \left(u^{y}\right)^4 - {\pi^{\eta}{}_{\eta}} \left(u^{y}\right)^4\\
  &\quad - {\pi_{Q}{}^{\eta}{}_{\eta}} \left(u^{y}\right)^4 + 3 {P'(T)} {T} \left(u^{y}\right)^4 + 3 {\Pi_{Q}} \left(u^{x}\right)^2 \left(u^{y}\right)^4\\
  &\quad + 3 {P'(T)} {T} \left(u^{x}\right)^2 \left(u^{y}\right)^4 + {\Pi_{Q}} \left(u^{y}\right)^6 + {P'(T)} {T} \left(u^{y}\right)^6\\
  &\quad + {\Pi} \left(\left(u^{y}\right)^2+1\right) \left(2 \left(u^{x}\right)^4+3 \left(u^{x}\right)^2 \left(\left(u^{y}\right)^2+1\right)+\left(\left(u^{y}\right)^2+1\right)^2\right)\\
  &\quad + {h} \left(\left(u^{y}\right)^2+1\right)\Bigl({\nu^{y}} {u^{y}} \left(\left(u^{y}\right)^2+1\right)\\
  &\quad + {\nu^{x}} {u^{x}} \left(2 \left(u^{x}\right)^2+3 \left(u^{y}\right)^2+3\right)\\
  &\quad + {n}\Bigl(2 \left(u^{x}\right)^4+3 \left(u^{x}\right)^2 \left(\left(u^{y}\right)^2+1\right)+\left(\left(u^{y}\right)^2+1\right)^2\Bigr)\Bigr)\Bigr) \\
(A_x)_{1,3} &= \frac{1}{\left(u^{\tau}\right)^3 \left(\left(u^{y}\right)^2+1\right)^2}\Bigl({\Pi_{Q}} \left(u^{\tau}\right)^2 {u^{x}} {u^{y}} \left(\left(u^{y}\right)^2+1\right)^2 + {n} \left(u^{\tau}\right)^2 {u^{x}} {u^{y}} \left(\left(u^{y}\right)^2+1\right)^2 ({h}-{T})\\
  &\quad + {n} {T} \left(u^{\tau}\right)^2 {u^{x}} {u^{y}} \left(\left(u^{y}\right)^2+1\right)^2\\
  &\quad + \left(u^{\tau}\right)^2 {u^{x}} {u^{y}} \left(\left(u^{y}\right)^2+1\right)^2 ({P'(T)} {T}+{\Pi})\\
  &\quad + {h} \left(\left(u^{y}\right)^2+1\right)^2 \left({\nu^{x}} {u^{y}} \left(\left(u^{y}\right)^2+1\right)+{\nu^{y}} \left(\left(u^{x}\right)^3+{u^{x}}\right)\right)\\
  &\quad - {u^{y}} \left(\left(u^{y}\right)^2+1\right)\Bigl(- {\pi^{\eta}{}_{\eta}} \left(u^{\tau}\right)^2 {u^{x}} - {\pi^{yy}} \left(\left(u^{x}\right)^3+{u^{x}}\right)\\
  &\quad + {\pi^{xy}} {u^{y}} \left(2 \left(u^{x}\right)^2+\left(u^{y}\right)^2+1\right)\Bigr)\\
  &\quad - {u^{y}} \left(\left(u^{y}\right)^2+1\right)\Bigl(- {\pi_{Q}{}^{\eta}{}_{\eta}} \left(u^{\tau}\right)^2 {u^{x}} - {\pi_{Q}^{yy}} \left(\left(u^{x}\right)^3+{u^{x}}\right)\\
  &\quad + {\pi_{Q}^{xy}} {u^{y}} \left(2 \left(u^{x}\right)^2+\left(u^{y}\right)^2+1\right)\Bigr)\\
  &\quad + \left(u^{\tau}\right)^2\Bigl(2 {u^{x}} {u^{y}} \left({\pi^{\eta}{}_{\eta}} \left(u^{x}\right)^2+{\pi^{yy}} \left(u^{x}\right)^2+{\pi^{yy}}\right)\\
  &\quad + {\pi^{xy}} \left(\left(\left(u^{y}\right)^2+1\right)^2-2 \left(u^{x}\right)^2 \left(\left(u^{y}\right)^2-1\right)\right)\Bigr)\\
  &\quad + \left(u^{\tau}\right)^2\Bigl(2 {u^{x}} {u^{y}} \left({\pi_{Q}^{yy}} \left(u^{x}\right)^2+{\pi_{Q}^{yy}}+{\pi_{Q}{}^{\eta}{}_{\eta}} \left(u^{x}\right)^2\right)\\
  &\quad + {\pi_{Q}^{xy}} \left(\left(\left(u^{y}\right)^2+1\right)^2-2 \left(u^{x}\right)^2 \left(\left(u^{y}\right)^2-1\right)\right)\Bigr)\Bigr) \\
(A_x)_{1,4} &= -\frac{{u^{x}} \left(\left(u^{x}\right)^2+1\right)}{{u^{\tau}} \left(\left(u^{y}\right)^2+1\right)} \\
(A_x)_{1,5} &= -\frac{{u^{\tau}} {u^{x}}}{\left(u^{y}\right)^2+1} \\
(A_x)_{1,6} &= \frac{{u^{y}} \left(2 \left(u^{x}\right)^2+\left(u^{y}\right)^2+1\right)}{{u^{\tau}} \left(\left(u^{y}\right)^2+1\right)} \\
(A_x)_{1,7} &= {u^{\tau}} {u^{x}} \\
(A_x)_{1,8} &= {(\partial_{\alpha} n)} {h} {u^{\tau}} {u^{x}} \\
(A_x)_{1,9} &= \frac{{h} \left(2 \left(u^{x}\right)^2+\left(u^{y}\right)^2+1\right)}{{u^{\tau}}} \\
(A_x)_{1,10} &= \frac{{h} {u^{x}} {u^{y}}}{{u^{\tau}}} \\
(A_x)_{1,11} &= -\frac{{u^{x}} \left(\left(u^{x}\right)^2+1\right)}{{u^{\tau}} \left(\left(u^{y}\right)^2+1\right)} \\
(A_x)_{1,12} &= -\frac{{u^{\tau}} {u^{x}}}{\left(u^{y}\right)^2+1} \\
(A_x)_{1,13} &= \frac{{u^{y}} \left(2 \left(u^{x}\right)^2+\left(u^{y}\right)^2+1\right)}{{u^{\tau}} \left(\left(u^{y}\right)^2+1\right)} \\
(A_x)_{1,14} &= {u^{\tau}} {u^{x}} \\
(A_x)_{2,1} &= {P'(T)} + {n} + {(\partial_{T} n)} {T} + 2 {h'(T)} {\nu^{x}} {u^{x}} + {P'(T)} \left(u^{x}\right)^2 + {(\partial_{T} n)} {h} \left(u^{x}\right)^2 + {h'(T)} {n} \left(u^{x}\right)^2\\
  &\quad + {P''(T)} {T} \left(u^{x}\right)^2 \\
(A_x)_{2,2} &= \frac{2}{\left(u^{y}\right)^2+1}\Bigl({\Pi_{Q}} {u^{x}} - {\pi^{\eta}{}_{\eta}} {u^{x}} - {\pi^{yy}} {u^{x}} - {\pi_{Q}^{yy}} {u^{x}} - {\pi_{Q}{}^{\eta}{}_{\eta}} {u^{x}} + {P'(T)} {T} {u^{x}} + {\pi^{xy}} {u^{y}} + {\pi_{Q}^{xy}} {u^{y}}\\
  &\quad + {\Pi_{Q}} {u^{x}} \left(u^{y}\right)^2 + {P'(T)} {T} {u^{x}} \left(u^{y}\right)^2 + {\Pi} {u^{x}} \left(\left(u^{y}\right)^2+1\right)\\
  &\quad + {h} \left(\left(u^{y}\right)^2+1\right) ({n} {u^{x}}+{\nu^{x}})\Bigr) \\
(A_x)_{2,3} &= \frac{2}{\left(\left(u^{y}\right)^2+1\right)^2}\Bigl({u^{y}} \left(\left(u^{x}\right)^2 ({\pi^{\eta}{}_{\eta}}+{\pi_{Q}{}^{\eta}{}_{\eta}})+{\pi^{yy}} \left(\left(u^{x}\right)^2+1\right)+{\pi_{Q}^{yy}} \left(\left(u^{x}\right)^2+1\right)\right)\\
  &\quad + {\pi^{xy}} \left({u^{x}}-{u^{x}} \left(u^{y}\right)^2\right) + {\pi_{Q}^{xy}} \left({u^{x}}-{u^{x}} \left(u^{y}\right)^2\right)\Bigr) \\
(A_x)_{2,4} &= -\frac{\left(u^{x}\right)^2+1}{\left(u^{y}\right)^2+1} \\
(A_x)_{2,5} &= -\frac{\left(u^{\tau}\right)^2}{\left(u^{y}\right)^2+1} \\
(A_x)_{2,6} &= \frac{2 {u^{x}} {u^{y}}}{\left(u^{y}\right)^2+1} \\
(A_x)_{2,7} &= \left(u^{x}\right)^2+1 \\
(A_x)_{2,8} &= {(\partial_{\alpha} n)} \left({h} \left(u^{x}\right)^2+{T}\right) \\
(A_x)_{2,9} &= 2 {h} {u^{x}} \\
(A_x)_{2,11} &= -\frac{\left(u^{x}\right)^2+1}{\left(u^{y}\right)^2+1} \\
(A_x)_{2,12} &= -\frac{\left(u^{\tau}\right)^2}{\left(u^{y}\right)^2+1} \\
(A_x)_{2,13} &= \frac{2 {u^{x}} {u^{y}}}{\left(u^{y}\right)^2+1} \\
(A_x)_{2,14} &= \left(u^{x}\right)^2+1 \\
(A_x)_{3,1} &= {h'(T)} {u^{y}} ({n} {u^{x}}+{\nu^{x}})+{h'(T)} {\nu^{y}} {u^{x}}+{u^{x}} {u^{y}} ({(\partial_{T} n)} {h}+{P''(T)} {T}+{P'(T)}) \\
(A_x)_{3,2} &= {u^{y}} ({P'(T)} {T}+{\Pi}+{\Pi_{Q}})+{h} ({n} {u^{y}}+{\nu^{y}}) \\
(A_x)_{3,3} &= {u^{x}} ({P'(T)} {T}+{\Pi}+{\Pi_{Q}})+{h} ({n} {u^{x}}+{\nu^{x}}) \\
(A_x)_{3,6} &= 1 \\
(A_x)_{3,7} &= {u^{x}} {u^{y}} \\
(A_x)_{3,8} &= {(\partial_{\alpha} n)} {h} {u^{x}} {u^{y}} \\
(A_x)_{3,9} &= {h} {u^{y}} \\
(A_x)_{3,10} &= {h} {u^{x}} \\
(A_x)_{3,13} &= 1 \\
(A_x)_{3,14} &= {u^{x}} {u^{y}} \\
(A_x)_{4,2} &= {\delta_{\pi\pi}} {\pi^{yy}}-\frac{2 {\pi^{xy}} {\tau_{\pi}} {u^{x}} \left(\left(u^{y}\right)^2+1\right) {u^{y}}}{\left(u^{\tau}\right)^2}+\frac{2 {\pi^{yy}} {\tau_{\pi}} \left(u^{x}\right)^2 \left(u^{y}\right)^2}{\left(u^{\tau}\right)^2}-\frac{2}{3} \left(\left(u^{y}\right)^2+1\right) {\eta} \\
(A_x)_{4,3} &= \frac{2 {u^{x}} {u^{y}} \left({\pi^{xy}} {\tau_{\pi}} {u^{x}} {u^{y}}-{\pi^{yy}} {\tau_{\pi}} \left(\left(u^{x}\right)^2+1\right)+\left(u^{\tau}\right)^2 {\eta}\right)}{\left(u^{\tau}\right)^2} \\
(A_x)_{4,4} &= {\tau_{\pi}} {u^{x}} \\
(A_x)_{5,2} &= \frac{3 {\delta_{\pi\pi}} {\pi^{\eta}{}_{\eta}}-2 {\eta}}{3 {\tau}^2} \\
(A_x)_{5,5} &= \frac{{\tau_{\pi}} {u^{x}}}{{\tau}^2} \\
(A_x)_{6,2} &= {\delta_{\pi\pi}} {\pi^{xy}}\\
  &\quad + \frac{{u^{x}} \left({u^{y}} \left(3 {\pi^{\eta}{}_{\eta}} {\tau_{\pi}} \left(u^{\tau}\right)^2+3 {\pi^{yy}} \left(2 {\tau_{\pi}} \left(u^{x}\right)^2+{\tau_{\pi}}\right)+\left(u^{\tau}\right)^2 {\eta}\right)-3 {\pi^{xy}} {\tau_{\pi}} {u^{x}} \left(2 \left(u^{y}\right)^2+1\right)\right)}{3 \left(u^{\tau}\right)^2} \\
(A_x)_{6,3} &= \frac{1}{\left(u^{\tau}\right)^2 \left(\left(u^{y}\right)^2+1\right)}\Bigl(- {\pi^{yy}} {\tau_{\pi}} \left(u^{x}\right)^2 \left(\left(u^{x}\right)^2+1\right) \left(2 \left(u^{y}\right)^2+1\right)\\
  &\quad + {\pi^{xy}} {\tau_{\pi}} {u^{x}} {u^{y}} \left(\left(2 \left(u^{x}\right)^2-1\right) \left(u^{y}\right)^2-1\right)\\
  &\quad - \left(u^{\tau}\right)^2 \left({\pi^{\eta}{}_{\eta}} {\tau_{\pi}} \left(u^{x}\right)^2 \left(u^{y}\right)^2-\left(\left(u^{x}\right)^2+1\right) \left(\left(u^{y}\right)^2+1\right) {\eta}\right)\Bigr) \\
(A_x)_{6,6} &= {\tau_{\pi}} {u^{x}} \\
(A_x)_{7,2} &= {\zeta} \\
(A_x)_{7,7} &= {\tau_{\Pi}} {u^{x}} \\
(A_x)_{8,1} &= {(\partial_{T} n)} {u^{x}} \\
(A_x)_{8,2} &= {n} \\
(A_x)_{8,8} &= {(\partial_{\alpha} n)} {u^{x}} \\
(A_x)_{8,9} &= 1 \\
(A_x)_{9,1} &= \frac{{D_{s}} \left({h'(T)} {\nu^{x}} {u^{x}}+{(\partial_{T} n)} {T} \left(\left(u^{x}\right)^2+1\right)+{n} \left(\left(u^{x}\right)^2+1\right)\right)}{{T}} \\
(A_x)_{9,2} &= -\frac{1}{{T} \left(u^{\tau}\right)^2 \left(\left(u^{y}\right)^2+1\right)}{D_{s}}\Bigl({\pi_{Q}^{yy}} {u^{x}} + {\pi_{Q}{}^{\eta}{}_{\eta}} {u^{x}} + {\pi_{Q}^{yy}} \left(u^{x}\right)^3 + {\pi_{Q}{}^{\eta}{}_{\eta}} \left(u^{x}\right)^3 - 2 {\pi_{Q}^{xy}} {u^{y}}\\
  &\quad - {\pi_{Q}^{xy}} \left(u^{x}\right)^2 {u^{y}} + {\pi_{Q}^{yy}} {u^{x}} \left(u^{y}\right)^2 - {\pi_{Q}^{yy}} \left(u^{x}\right)^3 \left(u^{y}\right)^2\\
  &\quad - {\pi_{Q}{}^{\eta}{}_{\eta}} \left(u^{x}\right)^3 \left(u^{y}\right)^2 - 2 {\pi_{Q}^{xy}} \left(u^{y}\right)^3 + {\pi_{Q}^{xy}} \left(u^{x}\right)^2 \left(u^{y}\right)^3\\
  &\quad - {\pi_{Q}{}^{\eta}{}_{\eta}} {u^{x}} \left(u^{y}\right)^4 - {\Pi_{Q}} \left(u^{\tau}\right)^2 {u^{x}} \left(\left(u^{y}\right)^2+1\right)\\
  &\quad + {h} \left(\left(u^{y}\right)^2+1\right)\Bigl(- {n} \left(u^{\tau}\right)^2 {u^{x}} - {\nu^{y}} \left(u^{x}\right)^3 {u^{y}} + {\nu^{x}} \left(u^{x}\right)^2 \left(\left(u^{y}\right)^2-1\right)\\
  &\quad - 2 {\nu^{x}} \left(\left(u^{y}\right)^2+1\right)\Bigr)\Bigr) \\
(A_x)_{9,3} &= \frac{1}{{T} \left(u^{\tau}\right)^2 \left(\left(u^{y}\right)^2+1\right)^2}{D_{s}}\Bigl(- {\pi_{Q}{}^{\eta}{}_{\eta}} \left(u^{\tau}\right)^2 \left(u^{x}\right)^2 {u^{y}} \left(\left(u^{y}\right)^2-1\right)\\
  &\quad - {h} \left(u^{x}\right)^2 \left(\left(u^{y}\right)^2+1\right)^2 \left(-{\nu^{x}} {u^{x}} {u^{y}}+{\nu^{y}} \left(u^{x}\right)^2+{\nu^{y}}\right)\\
  &\quad + {\pi_{Q}^{xy}} {u^{x}} \left(\left(u^{x}\right)^2 \left(\left(u^{y}\right)^2-1\right)^2-3 \left(u^{y}\right)^4-2 \left(u^{y}\right)^2+1\right)\\
  &\quad - {\pi_{Q}^{yy}} \left(\left(u^{x}\right)^2+1\right) {u^{y}} \left(\left(u^{x}\right)^2 \left(\left(u^{y}\right)^2-1\right)-2 \left(\left(u^{y}\right)^2+1\right)\right)\Bigr) \\
(A_x)_{9,8} &= {(\partial_{\alpha} n)} {D_{s}} \left(\left(u^{x}\right)^2+1\right) \\
(A_x)_{9,9} &= \frac{{D_{s}} {h} {u^{x}}}{{T}} \\
(A_x)_{9,11} &= -\frac{{D_{s}} \left(u^{x}\right)^2+{D_{s}}}{{T} \left(u^{y}\right)^2+{T}} \\
(A_x)_{9,12} &= -\frac{{D_{s}} \left(u^{\tau}\right)^2}{{T} \left(\left(u^{y}\right)^2+1\right)} \\
(A_x)_{9,13} &= \frac{2 {D_{s}} {u^{x}} {u^{y}}}{{T} \left(u^{y}\right)^2+{T}} \\
(A_x)_{9,14} &= \frac{{D_{s}} \left(\left(u^{x}\right)^2+1\right)}{{T}} \\
(A_x)_{10,1} &= \frac{{D_{s}} {u^{x}} ({h'(T)} {\nu^{y}}+{u^{y}} ({(\partial_{T} n)} {T}+{n}))}{{T}} \\
(A_x)_{10,2} &= \frac{1}{{T} \left(u^{\tau}\right)^2}{D_{s}}\Bigl({h} \left(\left(u^{y}\right)^2+1\right) \left(-{\nu^{x}} {u^{x}} {u^{y}}+{\nu^{y}} \left(u^{x}\right)^2+{\nu^{y}}\right)\\
  &\quad + {u^{y}} \left(-{\pi_{Q}^{xy}} {u^{x}} {u^{y}}+{\pi_{Q}^{yy}} \left(u^{x}\right)^2+{\pi_{Q}^{yy}}+{\pi_{Q}{}^{\eta}{}_{\eta}} \left(u^{x}\right)^2+{\pi_{Q}{}^{\eta}{}_{\eta}} \left(u^{y}\right)^2+{\pi_{Q}{}^{\eta}{}_{\eta}}\right)\Bigr) \\
(A_x)_{10,3} &= \frac{1}{{T} \left(u^{\tau}\right)^2 \left(\left(u^{y}\right)^2+1\right)}{D_{s}}\Bigl({\Pi_{Q}} \left(u^{\tau}\right)^2 {u^{x}} \left(\left(u^{y}\right)^2+1\right)\\
  &\quad + {h} \left(\left(u^{y}\right)^2+1\right)\Bigl({n} \left(u^{\tau}\right)^2 {u^{x}} - {\nu^{y}} {u^{x}} \left(\left(u^{x}\right)^2+1\right) {u^{y}}\\
  &\quad + {\nu^{x}} \left(\left(u^{x}\right)^2+1\right) \left(\left(u^{y}\right)^2+1\right)\Bigr)\\
  &\quad - {u^{y}}\Bigl({u^{x}} {u^{y}} \left({\pi_{Q}^{yy}} \left(\left(u^{x}\right)^2+1\right)+{\pi_{Q}{}^{\eta}{}_{\eta}} \left(u^{\tau}\right)^2\right)\\
  &\quad + {\pi_{Q}^{xy}} \left(-\left(\left(u^{x}\right)^2 \left(\left(u^{y}\right)^2-1\right)\right)+\left(u^{y}\right)^2+1\right)\Bigr)\Bigr) \\
(A_x)_{10,8} &= {(\partial_{\alpha} n)} {D_{s}} {u^{x}} {u^{y}} \\
(A_x)_{10,10} &= \frac{{D_{s}} {h} {u^{x}}}{{T}} \\
(A_x)_{10,13} &= \frac{{D_{s}}}{{T}} \\
(A_x)_{10,14} &= \frac{{D_{s}} {u^{x}} {u^{y}}}{{T}} \\
(A_x)_{11,1} &= \frac{{D_{s}} \left(3 {\mathcal{C}'(T)} {\pi_{Q}^{yy}} {u^{x}}-2 {\mathcal{B}'(T)} \left({\nu^{x}} \left(u^{y}\right)^2+{\nu^{x}}-3 {\nu^{y}} {u^{x}} {u^{y}}\right)\right)}{6 {T}} \\
(A_x)_{11,2} &= -\frac{1}{6 {T} \left(u^{\tau}\right)^2}{D_{s}}\Bigl(2 {I_{31}} \left(u^{\tau}\right)^2 \left(\left(u^{y}\right)^2+1\right)\\
  &\quad + 2 {u^{x}} \left({\mathcal{A}} \left(u^{y}\right)^2+{\mathcal{A}}+3 {\mathcal{B}} \left(u^{y}\right)^2\right) \left({\nu^{x}} \left(u^{y}\right)^2+{\nu^{x}}-{\nu^{y}} {u^{x}} {u^{y}}\right)\\
  &\quad + {\mathcal{C}}\Bigl(2 {\Pi_{Q}} \left(u^{\tau}\right)^2 \left(\left(u^{y}\right)^2+1\right)\\
  &\quad - 2 \left(\left(u^{y}\right)^2+1\right) \left({\pi_{Q}{}^{\eta}{}_{\eta}} \left(u^{\tau}\right)^2-4 {\pi_{Q}^{xy}} {u^{x}} {u^{y}}\right)\\
  &\quad - {\pi_{Q}^{yy}} \left(\left(u^{x}\right)^2 \left(8 \left(u^{y}\right)^2+5\right)+5 \left(\left(u^{y}\right)^2+1\right)\right)\Bigr)\Bigr) \\
(A_x)_{11,3} &= \frac{1}{3 {T} \left(u^{\tau}\right)^2}{D_{s}}\Bigl({u^{x}}\Bigl(3 {I_{31}} \left(u^{\tau}\right)^2 {u^{y}} - {\mathcal{A}} {\nu^{y}} \left(u^{x}\right)^2 \left(\left(u^{y}\right)^2-2\right) + 2 {\mathcal{A}} {\nu^{y}} \left(\left(u^{y}\right)^2+1\right)\\
  &\quad + {\mathcal{A}} {\nu^{x}} {u^{x}} {u^{y}} \left(\left(u^{y}\right)^2+1\right)\\
  &\quad + 3 {\mathcal{B}} \left({\nu^{x}} {u^{x}} \left(u^{y}\right)^3-{\nu^{y}} \left(u^{x}\right)^2 \left(\left(u^{y}\right)^2-1\right)+{\nu^{y}}\right)\Bigr)\\
  &\quad + {\mathcal{C}}\Bigl(- {u^{x}} {u^{y}} \left(-3 {\Pi_{Q}} \left(u^{\tau}\right)^2+4 {\pi_{Q}^{yy}} \left(\left(u^{x}\right)^2+1\right)+{\pi_{Q}{}^{\eta}{}_{\eta}} \left(u^{\tau}\right)^2\right)\\
  &\quad + 2 {\pi_{Q}^{xy}} \left(\left(u^{x}\right)^2 \left(2 \left(u^{y}\right)^2+1\right)+\left(u^{y}\right)^2+1\right)\Bigr)\Bigr) \\
(A_x)_{11,9} &= -\frac{{\mathcal{B}} {D_{s}} \left(\left(u^{y}\right)^2+1\right)}{3 {T}} \\
(A_x)_{11,10} &= \frac{{\mathcal{B}} {D_{s}} {u^{x}} {u^{y}}}{{T}} \\
(A_x)_{11,11} &= \frac{{\mathcal{C}} {D_{s}} {u^{x}}}{2 {T}} \\
(A_x)_{12,1} &= -\frac{{D_{s}} (2 {\mathcal{B}'(T)} {\nu^{x}}-3 {\mathcal{C}'(T)} {\pi_{Q}{}^{\eta}{}_{\eta}} {u^{x}})}{6 {T}} \\
(A_x)_{12,2} &= -\frac{1}{6 {T} \left(u^{\tau}\right)^2}{D_{s}}\Bigl(2 {I_{31}} \left(u^{\tau}\right)^2 + 2 {\mathcal{A}} {u^{x}} \left({\nu^{x}} \left(u^{y}\right)^2+{\nu^{x}}-{\nu^{y}} {u^{x}} {u^{y}}\right)\\
  &\quad + {\mathcal{C}}\Bigl(- 5 {\pi_{Q}{}^{\eta}{}_{\eta}} + 2 {\Pi_{Q}} \left(u^{\tau}\right)^2 - 5 {\pi_{Q}{}^{\eta}{}_{\eta}} \left(u^{x}\right)^2 - 2 {\pi_{Q}^{yy}} \left(\left(u^{x}\right)^2+1\right)\\
  &\quad + 2 {\pi_{Q}^{xy}} {u^{x}} {u^{y}} - 5 {\pi_{Q}{}^{\eta}{}_{\eta}} \left(u^{y}\right)^2\Bigr)\Bigr) \\
(A_x)_{12,3} &= \frac{1}{3 {T} \left(u^{\tau}\right)^2 \left(\left(u^{y}\right)^2+1\right)}{D_{s}}\Bigl(- {\mathcal{C}} {u^{x}} {u^{y}} \left({\pi_{Q}^{yy}} \left(\left(u^{x}\right)^2+1\right)+{\pi_{Q}{}^{\eta}{}_{\eta}} \left(u^{\tau}\right)^2\right)\\
  &\quad - {\mathcal{A}} {u^{x}} \left(\left(u^{y}\right)^2+1\right) \left(-{\nu^{x}} {u^{x}} {u^{y}}+{\nu^{y}} \left(u^{x}\right)^2+{\nu^{y}}\right)\\
  &\quad + {\mathcal{C}} {\pi_{Q}^{xy}} \left(\left(u^{x}\right)^2 \left(\left(u^{y}\right)^2-1\right)-\left(u^{y}\right)^2-1\right)\Bigr) \\
(A_x)_{12,9} &= -\frac{{\mathcal{B}} {D_{s}}}{3 {T}} \\
(A_x)_{12,12} &= \frac{{\mathcal{C}} {D_{s}} {u^{x}}}{2 {T}} \\
(A_x)_{13,1} &= \frac{{D_{s}} \left({\mathcal{B}'(T)} {\nu^{x}} {u^{x}} {u^{y}}+3 {\mathcal{B}'(T)} {\nu^{y}} \left(\left(u^{x}\right)^2+1\right)+3 {\mathcal{C}'(T)} {\pi_{Q}^{xy}} {u^{x}}\right)}{6 {T}} \\
(A_x)_{13,2} &= \frac{1}{6 {T} \left(u^{\tau}\right)^2}{D_{s}}\Bigl(3 {\mathcal{B}} {\nu^{y}} {u^{x}} + 3 {\mathcal{B}} {\nu^{y}} \left(u^{x}\right)^3 - 3 {\mathcal{B}} {\nu^{x}} {u^{y}} + {I_{31}} {u^{x}} {u^{y}} - 6 {\mathcal{B}} {\nu^{x}} \left(u^{x}\right)^2 {u^{y}} + {I_{31}} \left(u^{x}\right)^3 {u^{y}}\\
  &\quad + 6 {\mathcal{B}} {\nu^{y}} {u^{x}} \left(u^{y}\right)^2 + 6 {\mathcal{B}} {\nu^{y}} \left(u^{x}\right)^3 \left(u^{y}\right)^2 - 3 {\mathcal{B}} {\nu^{x}} \left(u^{y}\right)^3 + {I_{31}} {u^{x}} \left(u^{y}\right)^3\\
  &\quad - 6 {\mathcal{B}} {\nu^{x}} \left(u^{x}\right)^2 \left(u^{y}\right)^3\\
  &\quad + {\mathcal{A}} {u^{x}}\Bigl(3 {\nu^{y}} \left(\left(u^{y}\right)^2+1\right) - 2 {\nu^{x}} {u^{x}} {u^{y}} \left(\left(u^{y}\right)^2+1\right)\\
  &\quad + {\nu^{y}} \left(u^{x}\right)^2 \left(2 \left(u^{y}\right)^2+3\right)\Bigr)\\
  &\quad + {\mathcal{C}}\Bigl({u^{x}} {u^{y}} \left({\Pi_{Q}} \left(u^{\tau}\right)^2+{\pi_{Q}^{yy}} \left(8 \left(u^{x}\right)^2+5\right)+5 {\pi_{Q}{}^{\eta}{}_{\eta}} \left(u^{\tau}\right)^2\right)\\
  &\quad + {\pi_{Q}^{xy}} \left(\left(u^{x}\right)^2 \left(3-8 \left(u^{y}\right)^2\right)+6 \left(\left(u^{y}\right)^2+1\right)\right)\Bigr)\Bigr) \\
(A_x)_{13,3} &= \frac{1}{6 {T} \left(u^{\tau}\right)^2 \left(\left(u^{y}\right)^2+1\right)}{D_{s}}\Bigl(3 {I_{31}} \left(u^{\tau}\right)^2 \left(\left(u^{x}\right)^2+1\right) \left(\left(u^{y}\right)^2+1\right)\\
  &\quad + {\mathcal{C}}\Bigl(- 3 {\pi_{Q}{}^{\eta}{}_{\eta}} - 6 {\pi_{Q}{}^{\eta}{}_{\eta}} \left(u^{x}\right)^2 - 3 {\pi_{Q}{}^{\eta}{}_{\eta}} \left(u^{x}\right)^4 + {\pi_{Q}^{xy}} {u^{x}} {u^{y}}\\
  &\quad + 4 {\pi_{Q}^{xy}} \left(u^{x}\right)^3 {u^{y}} - 6 {\pi_{Q}{}^{\eta}{}_{\eta}} \left(u^{y}\right)^2 - 11 {\pi_{Q}{}^{\eta}{}_{\eta}} \left(u^{x}\right)^2 \left(u^{y}\right)^2\\
  &\quad - 5 {\pi_{Q}{}^{\eta}{}_{\eta}} \left(u^{x}\right)^4 \left(u^{y}\right)^2 + {\pi_{Q}^{xy}} {u^{x}} \left(u^{y}\right)^3 + 8 {\pi_{Q}^{xy}} \left(u^{x}\right)^3 \left(u^{y}\right)^3\\
  &\quad - 3 {\pi_{Q}{}^{\eta}{}_{\eta}} \left(u^{y}\right)^4 - 5 {\pi_{Q}{}^{\eta}{}_{\eta}} \left(u^{x}\right)^2 \left(u^{y}\right)^4\\
  &\quad + 3 {\Pi_{Q}} \left(u^{\tau}\right)^2 \left(\left(u^{x}\right)^2+1\right) \left(\left(u^{y}\right)^2+1\right)\\
  &\quad - {\pi_{Q}^{yy}} \left(\left(u^{x}\right)^2+1\right)\Bigl(\left(u^{x}\right)^2 \left(8 \left(u^{y}\right)^2+6\right)+3 \left(\left(u^{y}\right)^2+1\right)\Bigr)\Bigr)\\
  &\quad + \left(u^{y}\right)^2+1\Bigl({\mathcal{A}} {u^{x}}\Bigl(- 2 {\nu^{y}} {u^{x}} \left(\left(u^{x}\right)^2+1\right) {u^{y}} + 3 {\nu^{x}} \left(\left(u^{y}\right)^2+1\right)\\
  &\quad + {\nu^{x}} \left(u^{x}\right)^2 \left(2 \left(u^{y}\right)^2+3\right)\Bigr)\\
  &\quad + 3{\mathcal{B}} \left(\left(u^{x}\right)^2+1\right)\Bigl({\nu^{x}} \left(2 {u^{x}} \left(u^{y}\right)^2+{u^{x}}\right)-{\nu^{y}} \left(2 \left(u^{x}\right)^2+1\right) {u^{y}}\Bigr)\Bigr)\Bigr) \\
(A_x)_{13,9} &= \frac{{\mathcal{B}} {D_{s}} {u^{x}} {u^{y}}}{6 {T}} \\
(A_x)_{13,10} &= \frac{{\mathcal{B}} {D_{s}} \left(\left(u^{x}\right)^2+1\right)}{2 {T}} \\
(A_x)_{13,13} &= \frac{{\mathcal{C}} {D_{s}} {u^{x}}}{2 {T}} \\
(A_x)_{14,1} &= \frac{{D_{s}} (5 {\mathcal{B}'(T)} {\nu^{x}}+3 {u^{x}} ({\mathcal{C}'(T)} {\Pi_{Q}}+{(\partial_{T} I_{31})}))}{6 {T}} \\
(A_x)_{14,2} &= \frac{1}{6 {T} \left(u^{\tau}\right)^2}{D_{s}}\Bigl(5 {I_{31}} \left(u^{\tau}\right)^2 + 5 {\mathcal{C}} {\Pi_{Q}} \left(u^{\tau}\right)^2 + 2 {\mathcal{A}} {u^{x}} \left({\nu^{x}} \left(u^{y}\right)^2+{\nu^{x}}-{\nu^{y}} {u^{x}} {u^{y}}\right)\\
  &\quad - 2 {\mathcal{C}} \left(-{\pi_{Q}^{xy}} {u^{x}} {u^{y}}+{\pi_{Q}^{yy}} \left(u^{x}\right)^2+{\pi_{Q}^{yy}}+{\pi_{Q}{}^{\eta}{}_{\eta}} \left(u^{x}\right)^2+{\pi_{Q}{}^{\eta}{}_{\eta}} \left(u^{y}\right)^2+{\pi_{Q}{}^{\eta}{}_{\eta}}\right)\Bigr) \\
(A_x)_{14,3} &= \frac{1}{3 {T} \left(u^{\tau}\right)^2 \left(\left(u^{y}\right)^2+1\right)}{D_{s}}\Bigl({\mathcal{C}} {u^{x}} {u^{y}} \left({\pi_{Q}^{yy}} \left(\left(u^{x}\right)^2+1\right)+{\pi_{Q}{}^{\eta}{}_{\eta}} \left(u^{\tau}\right)^2\right)\\
  &\quad + {\mathcal{A}} {u^{x}} \left(\left(u^{y}\right)^2+1\right) \left(-{\nu^{x}} {u^{x}} {u^{y}}+{\nu^{y}} \left(u^{x}\right)^2+{\nu^{y}}\right)\\
  &\quad + {\mathcal{C}} {\pi_{Q}^{xy}} \left(-\left(\left(u^{x}\right)^2 \left(\left(u^{y}\right)^2-1\right)\right)+\left(u^{y}\right)^2+1\right)\Bigr) \\
(A_x)_{14,8} &= \frac{{(\partial_{\alpha} I_{31})} {D_{s}} {u^{x}}}{2 {T}} \\
(A_x)_{14,9} &= \frac{5 {\mathcal{B}} {D_{s}}}{6 {T}} \\
(A_x)_{14,14} &= \frac{{\mathcal{C}} {D_{s}} {u^{x}}}{2 {T}} \\
(A_y)_{1,1} &= \frac{1}{{u^{\tau}}}\Bigl(\left(u^{\tau}\right)^2 {u^{y}} ({(\partial_{T} n)} {h}+{P''(T)} {T}+{P'(T)})\\
  &\quad + {h'(T)} \left({n} \left(u^{\tau}\right)^2 {u^{y}}+{\nu^{x}} {u^{x}} {u^{y}}+{\nu^{y}} \left(\left(u^{x}\right)^2+2 \left(u^{y}\right)^2+1\right)\right)\Bigr) \\
(A_y)_{1,2} &= \frac{1}{\left(u^{\tau}\right)^3}\Bigl({\pi^{xy}} + {\pi_{Q}^{xy}} + {h} {\nu^{y}} {u^{x}} + {h} {\nu^{y}} \left(u^{x}\right)^3 + {h} {\nu^{x}} {u^{y}} + {h} {n} {u^{x}} {u^{y}} + {\Pi} {u^{x}} {u^{y}} + {\Pi_{Q}} {u^{x}} {u^{y}}\\
  &\quad - {\pi^{yy}} {u^{x}} {u^{y}} - {\pi_{Q}^{yy}} {u^{x}} {u^{y}} + {P'(T)} {T} {u^{x}} {u^{y}} + {h} {n} \left(u^{x}\right)^3 {u^{y}} + {\Pi} \left(u^{x}\right)^3 {u^{y}} + {\Pi_{Q}} \left(u^{x}\right)^3 {u^{y}}\\
  &\quad + {P'(T)} {T} \left(u^{x}\right)^3 {u^{y}} + {\pi^{xy}} \left(u^{y}\right)^2 + {\pi_{Q}^{xy}} \left(u^{y}\right)^2 + {h} {\nu^{x}} \left(u^{y}\right)^3 + {h} {n} {u^{x}} \left(u^{y}\right)^3\\
  &\quad + {\Pi} {u^{x}} \left(u^{y}\right)^3 + {\Pi_{Q}} {u^{x}} \left(u^{y}\right)^3 + {P'(T)} {T} {u^{x}} \left(u^{y}\right)^3\Bigr) \\
(A_y)_{1,3} &= \frac{1}{\left(u^{\tau}\right)^3}\Bigl({\Pi_{Q}} + {\pi^{yy}} + {\pi_{Q}^{yy}} + {P'(T)} {T} + 2 {\Pi_{Q}} \left(u^{x}\right)^2 + {\pi^{yy}} \left(u^{x}\right)^2 + {\pi_{Q}^{yy}} \left(u^{x}\right)^2\\
  &\quad + 2 {P'(T)} {T} \left(u^{x}\right)^2 + {\Pi_{Q}} \left(u^{x}\right)^4 + {P'(T)} {T} \left(u^{x}\right)^4 - {\pi^{xy}} {u^{x}} {u^{y}} - {\pi_{Q}^{xy}} {u^{x}} {u^{y}} + 3 {\Pi_{Q}} \left(u^{y}\right)^2\\
  &\quad + 3 {P'(T)} {T} \left(u^{y}\right)^2 + 3 {\Pi_{Q}} \left(u^{x}\right)^2 \left(u^{y}\right)^2 + 3 {P'(T)} {T} \left(u^{x}\right)^2 \left(u^{y}\right)^2 + 2 {\Pi_{Q}} \left(u^{y}\right)^4\\
  &\quad + 2 {P'(T)} {T} \left(u^{y}\right)^4 + {\Pi} \left(\left(u^{x}\right)^4+\left(u^{x}\right)^2 \left(3 \left(u^{y}\right)^2+2\right)+2 \left(u^{y}\right)^4+3 \left(u^{y}\right)^2+1\right)\\
  &\quad + {h}\Bigl({\nu^{x}} \left(\left(u^{x}\right)^3+{u^{x}}\right) + {\nu^{y}} {u^{y}} \left(3 \left(u^{x}\right)^2+2 \left(u^{y}\right)^2+3\right)\\
  &\quad + {n} \left(\left(u^{x}\right)^4+\left(u^{x}\right)^2 \left(3 \left(u^{y}\right)^2+2\right)+2 \left(u^{y}\right)^4+3 \left(u^{y}\right)^2+1\right)\Bigr)\Bigr) \\
(A_y)_{1,4} &= \frac{{u^{y}}}{{u^{\tau}}} \\
(A_y)_{1,6} &= \frac{{u^{x}}}{{u^{\tau}}} \\
(A_y)_{1,7} &= {u^{\tau}} {u^{y}} \\
(A_y)_{1,8} &= {(\partial_{\alpha} n)} {h} {u^{\tau}} {u^{y}} \\
(A_y)_{1,9} &= \frac{{h} {u^{x}} {u^{y}}}{{u^{\tau}}} \\
(A_y)_{1,10} &= \frac{{h} \left(\left(u^{x}\right)^2+2 \left(u^{y}\right)^2+1\right)}{{u^{\tau}}} \\
(A_y)_{1,11} &= \frac{{u^{y}}}{{u^{\tau}}} \\
(A_y)_{1,13} &= \frac{{u^{x}}}{{u^{\tau}}} \\
(A_y)_{1,14} &= {u^{\tau}} {u^{y}} \\
(A_y)_{2,1} &= {h'(T)} {u^{y}} ({n} {u^{x}}+{\nu^{x}})+{h'(T)} {\nu^{y}} {u^{x}}+{u^{x}} {u^{y}} ({(\partial_{T} n)} {h}+{P''(T)} {T}+{P'(T)}) \\
(A_y)_{2,2} &= {u^{y}} ({P'(T)} {T}+{\Pi}+{\Pi_{Q}})+{h} ({n} {u^{y}}+{\nu^{y}}) \\
(A_y)_{2,3} &= {u^{x}} ({P'(T)} {T}+{\Pi}+{\Pi_{Q}})+{h} ({n} {u^{x}}+{\nu^{x}}) \\
(A_y)_{2,6} &= 1 \\
(A_y)_{2,7} &= {u^{x}} {u^{y}} \\
(A_y)_{2,8} &= {(\partial_{\alpha} n)} {h} {u^{x}} {u^{y}} \\
(A_y)_{2,9} &= {h} {u^{y}} \\
(A_y)_{2,10} &= {h} {u^{x}} \\
(A_y)_{2,13} &= 1 \\
(A_y)_{2,14} &= {u^{x}} {u^{y}} \\
(A_y)_{3,1} &= {P'(T)} + {n} + {(\partial_{T} n)} {T} + 2 {h'(T)} {\nu^{y}} {u^{y}} + {P'(T)} \left(u^{y}\right)^2 + {(\partial_{T} n)} {h} \left(u^{y}\right)^2 + {h'(T)} {n} \left(u^{y}\right)^2\\
  &\quad + {P''(T)} {T} \left(u^{y}\right)^2 \\
(A_y)_{3,3} &= 2 ({u^{y}} ({P'(T)} {T}+{\Pi}+{\Pi_{Q}})+{h} ({n} {u^{y}}+{\nu^{y}})) \\
(A_y)_{3,4} &= 1 \\
(A_y)_{3,7} &= \left(u^{y}\right)^2+1 \\
(A_y)_{3,8} &= {(\partial_{\alpha} n)} \left({h} \left(u^{y}\right)^2+{T}\right) \\
(A_y)_{3,10} &= 2 {h} {u^{y}} \\
(A_y)_{3,11} &= 1 \\
(A_y)_{3,14} &= \left(u^{y}\right)^2+1 \\
(A_y)_{4,2} &= -\frac{2 {\tau_{\pi}} \left(u^{y}\right)^2 \left({\pi^{xy}} \left(u^{y}\right)^2+{\pi^{xy}}-{\pi^{yy}} {u^{x}} {u^{y}}\right)}{\left(u^{\tau}\right)^2} \\
(A_y)_{4,3} &= \frac{1}{3 \left(u^{\tau}\right)^2}\Bigl(3 {\delta_{\pi\pi}} {\pi^{yy}} \left(u^{\tau}\right)^2 - 6 {\pi^{yy}} {\tau_{\pi}} \left(\left(u^{x}\right)^2+1\right) \left(u^{y}\right)^2 + 6 {\pi^{xy}} {\tau_{\pi}} {u^{x}} \left(u^{y}\right)^3 + 4 {\eta}\\
  &\quad + 4 \left(u^{x}\right)^2 {\eta} + 8 \left(u^{y}\right)^2 {\eta} + 4 \left(u^{x}\right)^2 \left(u^{y}\right)^2 {\eta} + 4 \left(u^{y}\right)^4 {\eta}\Bigr) \\
(A_y)_{4,4} &= {\tau_{\pi}} {u^{y}} \\
(A_y)_{5,3} &= \frac{3 {\delta_{\pi\pi}} {\pi^{\eta}{}_{\eta}}-2 {\eta}}{3 {\tau}^2} \\
(A_y)_{5,5} &= \frac{{\tau_{\pi}} {u^{y}}}{{\tau}^2} \\
(A_y)_{6,2} &= \frac{1}{\left(u^{\tau}\right)^2}\Bigl({\pi^{yy}} {\tau_{\pi}} \left(u^{y}\right)^2 + {\pi^{\eta}{}_{\eta}} {\tau_{\pi}} \left(u^{\tau}\right)^2 \left(u^{y}\right)^2 + 2 {\pi^{yy}} {\tau_{\pi}} \left(u^{x}\right)^2 \left(u^{y}\right)^2\\
  &\quad - {\pi^{xy}} {\tau_{\pi}} {u^{x}} {u^{y}} \left(2 \left(u^{y}\right)^2+1\right) + {\eta} + \left(u^{x}\right)^2 {\eta} + 2 \left(u^{y}\right)^2 {\eta} + \left(u^{x}\right)^2 \left(u^{y}\right)^2 {\eta} + \left(u^{y}\right)^4 {\eta}\Bigr) \\
(A_y)_{6,3} &= {\delta_{\pi\pi}} {\pi^{xy}}\\
  &\quad - \frac{1}{3 \left(u^{\tau}\right)^2 \left(\left(u^{y}\right)^2+1\right)}{u^{y}}\Bigl(3 {\pi^{yy}} {\tau_{\pi}} {u^{x}} \left(\left(u^{x}\right)^2+1\right) \left(2 \left(u^{y}\right)^2+1\right)\\
  &\quad + 3 {\pi^{xy}} {\tau_{\pi}} {u^{y}} \left(\left(1-2 \left(u^{x}\right)^2\right) \left(u^{y}\right)^2+1\right)\\
  &\quad + \left(u^{\tau}\right)^2 {u^{x}} \left(3 {\pi^{\eta}{}_{\eta}} {\tau_{\pi}} \left(u^{y}\right)^2-\left(\left(u^{y}\right)^2+1\right) {\eta}\right)\Bigr) \\
(A_y)_{6,6} &= {\tau_{\pi}} {u^{y}} \\
(A_y)_{7,3} &= {\zeta} \\
(A_y)_{7,7} &= {\tau_{\Pi}} {u^{y}} \\
(A_y)_{8,1} &= {(\partial_{T} n)} {u^{y}} \\
(A_y)_{8,3} &= {n} \\
(A_y)_{8,8} &= {(\partial_{\alpha} n)} {u^{y}} \\
(A_y)_{8,10} &= 1 \\
(A_y)_{9,1} &= \frac{{D_{s}} {u^{y}} ({h'(T)} {\nu^{x}}+{u^{x}} ({(\partial_{T} n)} {T}+{n}))}{{T}} \\
(A_y)_{9,2} &= \frac{1}{{T} \left(u^{\tau}\right)^2}{D_{s}}\Bigl({\Pi_{Q}} \left(u^{\tau}\right)^2 {u^{y}} + {\pi_{Q}^{yy}} \left(u^{x}\right)^2 {u^{y}} - {\pi_{Q}^{xy}} {u^{x}} \left(\left(u^{y}\right)^2+1\right)\\
  &\quad + {h} \left({n} \left(u^{\tau}\right)^2 {u^{y}}-{\nu^{x}} {u^{x}} {u^{y}} \left(\left(u^{y}\right)^2+1\right)+{\nu^{y}} \left(\left(u^{x}\right)^2+1\right) \left(\left(u^{y}\right)^2+1\right)\right)\Bigr) \\
(A_y)_{9,3} &= \frac{{D_{s}} \left({h} \left(\left(u^{x}\right)^2+1\right) \left({\nu^{x}} \left(u^{y}\right)^2+{\nu^{x}}-{\nu^{y}} {u^{x}} {u^{y}}\right)-{u^{x}} \left(-{\pi_{Q}^{xy}} {u^{x}} {u^{y}}+{\pi_{Q}^{yy}} \left(u^{x}\right)^2+{\pi_{Q}^{yy}}\right)\right)}{{T} \left(u^{\tau}\right)^2} \\
(A_y)_{9,8} &= {(\partial_{\alpha} n)} {D_{s}} {u^{x}} {u^{y}} \\
(A_y)_{9,9} &= \frac{{D_{s}} {h} {u^{y}}}{{T}} \\
(A_y)_{9,13} &= \frac{{D_{s}}}{{T}} \\
(A_y)_{9,14} &= \frac{{D_{s}} {u^{x}} {u^{y}}}{{T}} \\
(A_y)_{10,1} &= \frac{{D_{s}} \left({h'(T)} {\nu^{y}} {u^{y}}+{(\partial_{T} n)} {T} \left(\left(u^{y}\right)^2+1\right)+{n} \left(\left(u^{y}\right)^2+1\right)\right)}{{T}} \\
(A_y)_{10,2} &= -\frac{{D_{s}} {u^{y}} \left({h} {u^{y}} \left({\nu^{x}} \left(u^{y}\right)^2+{\nu^{x}}-{\nu^{y}} {u^{x}} {u^{y}}\right)+{\pi_{Q}^{xy}} \left(\left(u^{y}\right)^2+1\right)-{\pi_{Q}^{yy}} {u^{x}} {u^{y}}\right)}{{T} \left(u^{\tau}\right)^2} \\
(A_y)_{10,3} &= \frac{1}{{T} \left(u^{\tau}\right)^2}{D_{s}}\Bigl({u^{y}} \left({\Pi_{Q}} \left(u^{\tau}\right)^2+{\pi_{Q}^{xy}} {u^{x}} {u^{y}}-{\pi_{Q}^{yy}} \left(\left(u^{x}\right)^2+1\right)\right)\\
  &\quad + {h} \left({n} \left(u^{\tau}\right)^2 {u^{y}}+{\nu^{x}} {u^{x}} \left(u^{y}\right)^3+{\nu^{y}} \left(-\left(\left(u^{x}\right)^2 \left(\left(u^{y}\right)^2-2\right)\right)+\left(u^{y}\right)^2+2\right)\right)\Bigr) \\
(A_y)_{10,8} &= {(\partial_{\alpha} n)} {D_{s}} \left(\left(u^{y}\right)^2+1\right) \\
(A_y)_{10,10} &= \frac{{D_{s}} {h} {u^{y}}}{{T}} \\
(A_y)_{10,11} &= \frac{{D_{s}}}{{T}} \\
(A_y)_{10,14} &= \frac{{D_{s}} \left(\left(u^{y}\right)^2+1\right)}{{T}} \\
(A_y)_{11,1} &= \frac{{D_{s}} \left(4 {\mathcal{B}'(T)} {\nu^{y}} \left(\left(u^{y}\right)^2+1\right)+3 {\mathcal{C}'(T)} {\pi_{Q}^{yy}} {u^{y}}\right)}{6 {T}} \\
(A_y)_{11,2} &= -\frac{1}{3 {T} \left(u^{\tau}\right)^2}{D_{s}}\Bigl({u^{y}} \left(\left(u^{y}\right)^2+1\right) ({\mathcal{A}}+3 {\mathcal{B}}) \left({\nu^{x}} \left(u^{y}\right)^2+{\nu^{x}}-{\nu^{y}} {u^{x}} {u^{y}}\right)\\
  &\quad + {\mathcal{C}} \left(4 \left(u^{y}\right)^2+1\right) \left({\pi_{Q}^{xy}} \left(u^{y}\right)^2+{\pi_{Q}^{xy}}-{\pi_{Q}^{yy}} {u^{x}} {u^{y}}\right)\Bigr) \\
(A_y)_{11,3} &= \frac{1}{6 {T} \left(u^{\tau}\right)^2}{D_{s}}\Bigl(4 {I_{31}} \left(u^{\tau}\right)^2 \left(\left(u^{y}\right)^2+1\right)\\
  &\quad + {\mathcal{C}}\Bigl(4 {\Pi_{Q}} \left(u^{\tau}\right)^2 \left(\left(u^{y}\right)^2+1\right) + 2 {\pi_{Q}^{xy}} {u^{x}} {u^{y}} \left(4 \left(u^{y}\right)^2+1\right)\\
  &\quad + {\pi_{Q}^{yy}} \left(\left(u^{x}\right)^2 \left(7-8 \left(u^{y}\right)^2\right)+\left(u^{y}\right)^2+7\right)\Bigr)\\
  &\quad + 2{u^{y}}\Bigl(3 {\mathcal{B}} {u^{x}} {u^{y}} \left({\nu^{x}} \left(u^{y}\right)^2+{\nu^{x}}-{\nu^{y}} {u^{x}} {u^{y}}\right)\\
  &\quad + {\mathcal{A}}\Bigl(- {\nu^{y}} \left(u^{x}\right)^2 \left(\left(u^{y}\right)^2-2\right) + 2 {\nu^{y}} \left(\left(u^{y}\right)^2+1\right)\\
  &\quad + {\nu^{x}} {u^{x}} {u^{y}} \left(\left(u^{y}\right)^2+1\right)\Bigr)\Bigr)\Bigr) \\
(A_y)_{11,10} &= \frac{2 {\mathcal{B}} {D_{s}} \left(\left(u^{y}\right)^2+1\right)}{3 {T}} \\
(A_y)_{11,11} &= \frac{{\mathcal{C}} {D_{s}} {u^{y}}}{2 {T}} \\
(A_y)_{12,1} &= -\frac{{D_{s}} (2 {\mathcal{B}'(T)} {\nu^{y}}-3 {\mathcal{C}'(T)} {\pi_{Q}{}^{\eta}{}_{\eta}} {u^{y}})}{6 {T}} \\
(A_y)_{12,2} &= -\frac{{D_{s}} \left({\mathcal{A}} {u^{y}} \left({\nu^{x}} \left(u^{y}\right)^2+{\nu^{x}}-{\nu^{y}} {u^{x}} {u^{y}}\right)+{\mathcal{C}} \left({\pi_{Q}^{xy}} \left(u^{y}\right)^2+{\pi_{Q}^{xy}}-{\pi_{Q}^{yy}} {u^{x}} {u^{y}}\right)\right)}{3 {T} \left(u^{\tau}\right)^2} \\
(A_y)_{12,3} &= -\frac{1}{6 {T} \left(u^{\tau}\right)^2}{D_{s}}\Bigl(2 {I_{31}} \left(u^{\tau}\right)^2 + 2 {\mathcal{A}} {u^{y}} \left(-{\nu^{x}} {u^{x}} {u^{y}}+{\nu^{y}} \left(u^{x}\right)^2+{\nu^{y}}\right)\\
  &\quad + {\mathcal{C}}\Bigl(- 3 {\pi_{Q}{}^{\eta}{}_{\eta}} + 2 {\Pi_{Q}} \left(u^{\tau}\right)^2 - 3 {\pi_{Q}{}^{\eta}{}_{\eta}} \left(u^{x}\right)^2 + 2 {\pi_{Q}^{yy}} \left(\left(u^{x}\right)^2+1\right)\\
  &\quad - 2 {\pi_{Q}^{xy}} {u^{x}} {u^{y}} - 3 {\pi_{Q}{}^{\eta}{}_{\eta}} \left(u^{y}\right)^2\Bigr)\Bigr) \\
(A_y)_{12,10} &= -\frac{{\mathcal{B}} {D_{s}}}{3 {T}} \\
(A_y)_{12,12} &= \frac{{\mathcal{C}} {D_{s}} {u^{y}}}{2 {T}} \\
(A_y)_{13,1} &= \frac{{D_{s}} \left(3 {\mathcal{B}'(T)} {\nu^{x}} \left(\left(u^{y}\right)^2+1\right)+{\mathcal{B}'(T)} {\nu^{y}} {u^{x}} {u^{y}}+3 {\mathcal{C}'(T)} {\pi_{Q}^{xy}} {u^{y}}\right)}{6 {T}} \\
(A_y)_{13,2} &= \frac{1}{6 {T} \left(u^{\tau}\right)^2}{D_{s}}\Bigl(- 3 {\mathcal{B}} {\nu^{x}} {u^{x}} + 3 {\mathcal{A}} {\nu^{y}} {u^{y}} + 3 {\mathcal{B}} {\nu^{y}} {u^{y}} + 3 {\mathcal{A}} {\nu^{y}} \left(u^{x}\right)^2 {u^{y}} + 6 {\mathcal{B}} {\nu^{y}} \left(u^{x}\right)^2 {u^{y}}\\
  &\quad - 2 {\mathcal{A}} {\nu^{x}} {u^{x}} \left(u^{y}\right)^2 - 9 {\mathcal{B}} {\nu^{x}} {u^{x}} \left(u^{y}\right)^2 + 3 {\mathcal{A}} {\nu^{y}} \left(u^{y}\right)^3 + 3 {\mathcal{B}} {\nu^{y}} \left(u^{y}\right)^3\\
  &\quad + 2 {\mathcal{A}} {\nu^{y}} \left(u^{x}\right)^2 \left(u^{y}\right)^3 + 6 {\mathcal{B}} {\nu^{y}} \left(u^{x}\right)^2 \left(u^{y}\right)^3 - 2 {\mathcal{A}} {\nu^{x}} {u^{x}} \left(u^{y}\right)^4 - 6 {\mathcal{B}} {\nu^{x}} {u^{x}} \left(u^{y}\right)^4\\
  &\quad + 3 {I_{31}} \left(u^{\tau}\right)^2 \left(\left(u^{y}\right)^2+1\right)\\
  &\quad + {\mathcal{C}}\Bigl(3 {\Pi_{Q}} \left(u^{\tau}\right)^2 \left(\left(u^{y}\right)^2+1\right)\\
  &\quad + {\pi_{Q}^{yy}} \left(\left(u^{x}\right)^2 \left(8 \left(u^{y}\right)^2+3\right)+6 \left(u^{y}\right)^2+3\right)\\
  &\quad + {u^{y}} \left(3 {\pi_{Q}{}^{\eta}{}_{\eta}} \left(u^{\tau}\right)^2 {u^{y}}-{\pi_{Q}^{xy}} {u^{x}} \left(8 \left(u^{y}\right)^2+5\right)\right)\Bigr)\Bigr) \\
(A_y)_{13,3} &= \frac{1}{6 {T} \left(u^{\tau}\right)^2 \left(\left(u^{y}\right)^2+1\right)}{D_{s}}\Bigl(\left(u^{y}\right)^2+1\Bigl({I_{31}} \left(u^{\tau}\right)^2 {u^{x}} {u^{y}}\\
  &\quad + 3 {\mathcal{B}} {\nu^{x}} \left(2 \left(u^{x}\right)^2+1\right) {u^{y}} \left(\left(u^{y}\right)^2+1\right)\\
  &\quad - 3 {\mathcal{B}} {\nu^{y}} {u^{x}} \left(\left(u^{x}\right)^2+1\right) \left(2 \left(u^{y}\right)^2+1\right)\\
  &\quad + {\mathcal{A}} {u^{y}}\Bigl(- 2 {\nu^{y}} {u^{x}} \left(\left(u^{x}\right)^2+1\right) {u^{y}} + 3 {\nu^{x}} \left(\left(u^{y}\right)^2+1\right)\\
  &\quad + {\nu^{x}} \left(u^{x}\right)^2 \left(2 \left(u^{y}\right)^2+3\right)\Bigr)\Bigr)\\
  &\quad + {\mathcal{C}}\Bigl({u^{x}} {u^{y}}\Bigl(- 3 {\pi_{Q}{}^{\eta}{}_{\eta}} \left(u^{\tau}\right)^2 \left(u^{y}\right)^2 + {\Pi_{Q}} \left(u^{\tau}\right)^2 \left(\left(u^{y}\right)^2+1\right)\\
  &\quad - {\pi_{Q}^{yy}} \left(\left(u^{x}\right)^2+1\right) \left(8 \left(u^{y}\right)^2+5\right)\Bigr)\\
  &\quad + {\pi_{Q}^{xy}}\Bigl(3 \left(\left(u^{y}\right)^4+3 \left(u^{y}\right)^2+2\right)\\
  &\quad + \left(u^{x}\right)^2 \left(8 \left(u^{y}\right)^4+8 \left(u^{y}\right)^2+6\right)\Bigr)\Bigr)\Bigr) \\
(A_y)_{13,9} &= \frac{{\mathcal{B}} {D_{s}} \left(\left(u^{y}\right)^2+1\right)}{2 {T}} \\
(A_y)_{13,10} &= \frac{{\mathcal{B}} {D_{s}} {u^{x}} {u^{y}}}{6 {T}} \\
(A_y)_{13,13} &= \frac{{\mathcal{C}} {D_{s}} {u^{y}}}{2 {T}} \\
(A_y)_{14,1} &= \frac{{D_{s}} (5 {\mathcal{B}'(T)} {\nu^{y}}+3 {u^{y}} ({\mathcal{C}'(T)} {\Pi_{Q}}+{(\partial_{T} I_{31})}))}{6 {T}} \\
(A_y)_{14,2} &= \frac{{D_{s}} \left({\mathcal{A}} {u^{y}} \left({\nu^{x}} \left(u^{y}\right)^2+{\nu^{x}}-{\nu^{y}} {u^{x}} {u^{y}}\right)+{\mathcal{C}} \left({\pi_{Q}^{xy}} \left(u^{y}\right)^2+{\pi_{Q}^{xy}}-{\pi_{Q}^{yy}} {u^{x}} {u^{y}}\right)\right)}{3 {T} \left(u^{\tau}\right)^2} \\
(A_y)_{14,3} &= \frac{1}{6 {T} \left(u^{\tau}\right)^2}{D_{s}}\Bigl(5 {I_{31}} \left(u^{\tau}\right)^2 + 5 {\mathcal{C}} {\Pi_{Q}} \left(u^{\tau}\right)^2 + 2 {\mathcal{A}} {u^{y}} \left(-{\nu^{x}} {u^{x}} {u^{y}}+{\nu^{y}} \left(u^{x}\right)^2+{\nu^{y}}\right)\\
  &\quad + 2 {\mathcal{C}} \left(-{\pi_{Q}^{xy}} {u^{x}} {u^{y}}+{\pi_{Q}^{yy}} \left(u^{x}\right)^2+{\pi_{Q}^{yy}}\right)\Bigr) \\
(A_y)_{14,8} &= \frac{{(\partial_{\alpha} I_{31})} {D_{s}} {u^{y}}}{2 {T}} \\
(A_y)_{14,10} &= \frac{5 {\mathcal{B}} {D_{s}}}{6 {T}} \\
(A_y)_{14,14} &= \frac{{\mathcal{C}} {D_{s}} {u^{y}}}{2 {T}} \\
S_{1} &= \frac{1}{{\tau} \left(\left(u^{y}\right)^2+1\right)}\Bigl({\Pi_{Q}} + {\pi^{\eta}{}_{\eta}} + {\pi_{Q}{}^{\eta}{}_{\eta}} + {P'(T)} {T} + {\Pi_{Q}} \left(u^{x}\right)^2 - {\pi^{\eta}{}_{\eta}} \left(u^{x}\right)^2 - {\pi^{yy}} \left(u^{x}\right)^2\\
  &\quad - {\pi_{Q}^{yy}} \left(u^{x}\right)^2 - {\pi_{Q}{}^{\eta}{}_{\eta}} \left(u^{x}\right)^2 + {P'(T)} {T} \left(u^{x}\right)^2 + 2 {\pi^{xy}} {u^{x}} {u^{y}} + 2 {\pi_{Q}^{xy}} {u^{x}} {u^{y}}\\
  &\quad + 2 {\Pi_{Q}} \left(u^{y}\right)^2 + {\pi^{\eta}{}_{\eta}} \left(u^{y}\right)^2 + {\pi^{yy}} \left(u^{y}\right)^2 + {\pi_{Q}^{yy}} \left(u^{y}\right)^2 + {\pi_{Q}{}^{\eta}{}_{\eta}} \left(u^{y}\right)^2\\
  &\quad + 2 {P'(T)} {T} \left(u^{y}\right)^2 + {\Pi_{Q}} \left(u^{x}\right)^2 \left(u^{y}\right)^2 + {P'(T)} {T} \left(u^{x}\right)^2 \left(u^{y}\right)^2 + {\Pi_{Q}} \left(u^{y}\right)^4\\
  &\quad + {P'(T)} {T} \left(u^{y}\right)^4 + {\Pi} \left(u^{\tau}\right)^2 \left(\left(u^{y}\right)^2+1\right)\\
  &\quad + {h} \left(\left(u^{y}\right)^2+1\right) \left({n} \left(u^{\tau}\right)^2+2 ({\nu^{x}} {u^{x}}+{\nu^{y}} {u^{y}})\right)\Bigr) \\
S_{2} &= \frac{1}{{\tau} {u^{\tau}} \left(\left(u^{y}\right)^2+1\right)}\Bigl({\Pi_{Q}} {u^{x}} - {\pi^{\eta}{}_{\eta}} {u^{x}} - {\pi^{yy}} {u^{x}} - {\pi_{Q}^{yy}} {u^{x}} - {\pi_{Q}{}^{\eta}{}_{\eta}} {u^{x}} + {P'(T)} {T} {u^{x}} + {\Pi_{Q}} \left(u^{x}\right)^3\\
  &\quad - {\pi^{\eta}{}_{\eta}} \left(u^{x}\right)^3 - {\pi^{yy}} \left(u^{x}\right)^3 - {\pi_{Q}^{yy}} \left(u^{x}\right)^3 - {\pi_{Q}{}^{\eta}{}_{\eta}} \left(u^{x}\right)^3 + {P'(T)} {T} \left(u^{x}\right)^3\\
  &\quad + {\pi^{xy}} {u^{y}} + {\pi_{Q}^{xy}} {u^{y}} + 2 {\pi^{xy}} \left(u^{x}\right)^2 {u^{y}} + 2 {\pi_{Q}^{xy}} \left(u^{x}\right)^2 {u^{y}} + 2 {\Pi_{Q}} {u^{x}} \left(u^{y}\right)^2\\
  &\quad - {\pi^{\eta}{}_{\eta}} {u^{x}} \left(u^{y}\right)^2 - {\pi_{Q}{}^{\eta}{}_{\eta}} {u^{x}} \left(u^{y}\right)^2 + 2 {P'(T)} {T} {u^{x}} \left(u^{y}\right)^2 + {\Pi_{Q}} \left(u^{x}\right)^3 \left(u^{y}\right)^2\\
  &\quad + {P'(T)} {T} \left(u^{x}\right)^3 \left(u^{y}\right)^2 + {\pi^{xy}} \left(u^{y}\right)^3 + {\pi_{Q}^{xy}} \left(u^{y}\right)^3 + {\Pi_{Q}} {u^{x}} \left(u^{y}\right)^4\\
  &\quad + {P'(T)} {T} {u^{x}} \left(u^{y}\right)^4 + {\Pi} \left(u^{\tau}\right)^2 {u^{x}} \left(\left(u^{y}\right)^2+1\right)\\
  &\quad + {h} \left(\left(u^{y}\right)^2+1\right) \left({n} \left(u^{\tau}\right)^2 {u^{x}}+{\nu^{x}} \left(2 \left(u^{x}\right)^2+\left(u^{y}\right)^2+1\right)+{\nu^{y}} {u^{x}} {u^{y}}\right)\Bigr) \\
S_{3} &= \frac{1}{{\tau} {u^{\tau}}}\Bigl({\pi^{xy}} {u^{x}} + {\pi_{Q}^{xy}} {u^{x}} + {\Pi} {u^{y}} + {\Pi_{Q}} {u^{y}} + {\pi^{yy}} {u^{y}} + {\pi_{Q}^{yy}} {u^{y}} + {P'(T)} {T} {u^{y}} + {\Pi} \left(u^{x}\right)^2 {u^{y}}\\
  &\quad + {\Pi_{Q}} \left(u^{x}\right)^2 {u^{y}} + {P'(T)} {T} \left(u^{x}\right)^2 {u^{y}} + {\Pi} \left(u^{y}\right)^3 + {\Pi_{Q}} \left(u^{y}\right)^3 + {P'(T)} {T} \left(u^{y}\right)^3\\
  &\quad + {h} \left({n} \left(u^{\tau}\right)^2 {u^{y}}+{\nu^{x}} {u^{x}} {u^{y}}+{\nu^{y}} \left(\left(u^{x}\right)^2+2 \left(u^{y}\right)^2+1\right)\right)\Bigr) \\
S_{4} &= \frac{{\delta_{\pi\pi}} {\pi^{yy}} {u^{\tau}}}{{\tau}}+{\pi^{yy}}-\frac{2 {u^{\tau}} \left(\left(u^{y}\right)^2+1\right) {\eta}}{3 {\tau}} \\
S_{5} &= \frac{3 {\pi^{\eta}{}_{\eta}} ({\delta_{\pi\pi}} {u^{\tau}}+{\tau})+4 {u^{\tau}} {\eta}}{3 {\tau}^3} \\
S_{6} &= \frac{{\delta_{\pi\pi}} {\pi^{xy}} {u^{\tau}}}{{\tau}}+{\pi^{xy}}-\frac{2 {u^{\tau}} {u^{x}} {u^{y}} {\eta}}{3 {\tau}} \\
S_{7} &= {\Pi}+\frac{{u^{\tau}} {\zeta}}{{\tau}} \\
S_{8} &= \frac{{n} \left(u^{\tau}\right)^2+{\nu^{x}} {u^{x}}+{\nu^{y}} {u^{y}}}{{\tau} {u^{\tau}}} \\
S_{9} &= {\nu^{x}}\\
  &\quad + \frac{1}{{T} {\tau} {u^{\tau}} \left(\left(u^{y}\right)^2+1\right)}{D_{s}}\Bigl(- 2 {\pi_{Q}{}^{\eta}{}_{\eta}} {u^{x}} - 2 {\pi_{Q}{}^{\eta}{}_{\eta}} \left(u^{x}\right)^3 - {\pi_{Q}^{yy}} \left(\left(u^{x}\right)^3+{u^{x}}\right) + {\pi_{Q}^{xy}} {u^{y}}\\
  &\quad + 2 {\pi_{Q}^{xy}} \left(u^{x}\right)^2 {u^{y}} - 3 {\pi_{Q}{}^{\eta}{}_{\eta}} {u^{x}} \left(u^{y}\right)^2 - {\pi_{Q}{}^{\eta}{}_{\eta}} \left(u^{x}\right)^3 \left(u^{y}\right)^2 + {\pi_{Q}^{xy}} \left(u^{y}\right)^3\\
  &\quad - {\pi_{Q}{}^{\eta}{}_{\eta}} {u^{x}} \left(u^{y}\right)^4 + {h} {\nu^{x}} \left(u^{\tau}\right)^2 \left(\left(u^{y}\right)^2+1\right)\Bigr) \\
S_{10} &= \frac{{D_{s}} \left({h} {\nu^{y}} \left(u^{\tau}\right)^2+{\pi_{Q}^{xy}} {u^{x}}+{u^{y}} \left({\pi_{Q}^{yy}}-{\pi_{Q}{}^{\eta}{}_{\eta}} \left(u^{\tau}\right)^2\right)\right)}{{T} {\tau} {u^{\tau}}}+{\nu^{y}} \\
S_{11} &= {\pi_{Q}^{yy}}\\
  &\quad - \frac{1}{6 {T} {\tau} {u^{\tau}}}{D_{s}}\Bigl(2 {I_{31}} \left(u^{\tau}\right)^2 \left(\left(u^{y}\right)^2+1\right) + 2 {\mathcal{B}} \left(\left(u^{y}\right)^2+1\right) ({\nu^{x}} {u^{x}}+{\nu^{y}} {u^{y}})\\
  &\quad + {\mathcal{C}} \left(u^{\tau}\right)^2 \left(2 {\Pi_{Q}} \left(\left(u^{y}\right)^2+1\right)-3 {\pi_{Q}^{yy}}+2 {\pi_{Q}{}^{\eta}{}_{\eta}} \left(\left(u^{y}\right)^2+1\right)\right)\Bigr) \\
S_{12} &= \frac{{D_{s}} \left(4 {\mathcal{B}} ({\nu^{x}} {u^{x}}+{\nu^{y}} {u^{y}})+{\mathcal{C}} \left(u^{\tau}\right)^2 (4 {\Pi_{Q}}+7 {\pi_{Q}{}^{\eta}{}_{\eta}})+4 {I_{31}} \left(u^{\tau}\right)^2\right)}{6 {T} {\tau} {u^{\tau}}}+{\pi_{Q}{}^{\eta}{}_{\eta}} \\
S_{13} &= -\frac{{D_{s}} {u^{x}} {u^{y}} \left({\mathcal{B}} ({\nu^{x}} {u^{x}}+{\nu^{y}} {u^{y}})+{I_{31}} \left(u^{\tau}\right)^2\right)}{3 {T} {\tau} {u^{\tau}}}+\frac{{\mathcal{C}} {D_{s}} {u^{\tau}} (3 {\pi_{Q}^{xy}}-2 {u^{x}} {u^{y}} ({\Pi_{Q}}+{\pi_{Q}{}^{\eta}{}_{\eta}}))}{6 {T} {\tau}}+{\pi_{Q}^{xy}} \\
S_{14} &= \frac{{D_{s}} \left(5 {\mathcal{B}} ({\nu^{x}} {u^{x}}+{\nu^{y}} {u^{y}})+{\mathcal{C}} \left(u^{\tau}\right)^2 (5 {\Pi_{Q}}+2 {\pi_{Q}{}^{\eta}{}_{\eta}})+5 {I_{31}} \left(u^{\tau}\right)^2\right)}{6 {T} {\tau} {u^{\tau}}}+{\Pi_{Q}}
\end{align*}
\subsubsection{S6, fully coupled, conformal}

\texttt{two\_\allowbreak{}d\_\allowbreak{}two\_\allowbreak{}way\_\allowbreak{}cM\_\allowbreak{}matrix\_\allowbreak{}conformal} in \texttt{src/Matrix/2d\_\allowbreak{}two\_\allowbreak{}way\_\allowbreak{}cM.jl}, fields:

\[
\phi = ({T}, {u^{x}}, {u^{y}}, {\pi^{yy}}, {\pi^{\eta}{}_{\eta}}, {\pi^{xy}}, {\alpha}, {\nu^{x}}, {\nu^{y}}, {\pi_{Q}^{yy}}, {\pi_{Q}{}^{\eta}{}_{\eta}}, {\pi_{Q}^{xy}}, {\Pi_{Q}})
\]
\begin{align*}
(A_t)_{1,1} &= - {n} - {(\partial_{T} n)} {T} + {(\partial_{T} n)} {h} \left(u^{\tau}\right)^2 + {h'(T)} {n} \left(u^{\tau}\right)^2 + 2 {h'(T)} ({\nu^{x}} {u^{x}}+{\nu^{y}} {u^{y}})\\
  &\quad + 4 {a_{\mathrm{SB}}} {T}^3 \left(4 \left(u^{x}\right)^2+4 \left(u^{y}\right)^2+3\right) \\
(A_t)_{1,2} &= \frac{2}{\left(u^{y}\right)^2+1}\Bigl(- {\pi^{\eta}{}_{\eta}} {u^{x}} - {\pi^{yy}} {u^{x}} - {\pi_{Q}^{yy}} {u^{x}} - {\pi_{Q}{}^{\eta}{}_{\eta}} {u^{x}} + 4 {a_{\mathrm{SB}}} {T}^4 {u^{x}} + {\pi^{xy}} {u^{y}} + {\pi_{Q}^{xy}} {u^{y}}\\
  &\quad + 4 {a_{\mathrm{SB}}} {T}^4 {u^{x}} \left(u^{y}\right)^2 + {\Pi_{Q}} {u^{x}} \left(\left(u^{y}\right)^2+1\right) + {h} \left(\left(u^{y}\right)^2+1\right) ({n} {u^{x}}+{\nu^{x}})\Bigr) \\
(A_t)_{1,3} &= \frac{2}{\left(\left(u^{y}\right)^2+1\right)^2}\Bigl({\pi_{Q}^{xy}} {u^{x}} + {\Pi_{Q}} {u^{y}} + {\pi^{yy}} {u^{y}} + {\pi_{Q}^{yy}} {u^{y}} + 4 {a_{\mathrm{SB}}} {T}^4 {u^{y}} + {\pi^{\eta}{}_{\eta}} \left(u^{x}\right)^2 {u^{y}}\\
  &\quad + {\pi^{yy}} \left(u^{x}\right)^2 {u^{y}} + {\pi_{Q}^{yy}} \left(u^{x}\right)^2 {u^{y}} + {\pi_{Q}{}^{\eta}{}_{\eta}} \left(u^{x}\right)^2 {u^{y}} - {\pi_{Q}^{xy}} {u^{x}} \left(u^{y}\right)^2 + 2 {\Pi_{Q}} \left(u^{y}\right)^3\\
  &\quad + 8 {a_{\mathrm{SB}}} {T}^4 \left(u^{y}\right)^3 + {\Pi_{Q}} \left(u^{y}\right)^5 + 4 {a_{\mathrm{SB}}} {T}^4 \left(u^{y}\right)^5 + {h} \left(\left(u^{y}\right)^2+1\right)^2 ({n} {u^{y}}+{\nu^{y}})\\
  &\quad + {\pi^{xy}} \left({u^{x}}-{u^{x}} \left(u^{y}\right)^2\right)\Bigr) \\
(A_t)_{1,4} &= \frac{\left(u^{y}\right)^2-\left(u^{x}\right)^2}{\left(u^{y}\right)^2+1} \\
(A_t)_{1,5} &= -\frac{\left(u^{x}\right)^2}{\left(u^{y}\right)^2+1} \\
(A_t)_{1,6} &= \frac{2 {u^{x}} {u^{y}}}{\left(u^{y}\right)^2+1} \\
(A_t)_{1,7} &= {(\partial_{\alpha} n)} \left({h} \left(u^{\tau}\right)^2-{T}\right) \\
(A_t)_{1,8} &= 2 {h} {u^{x}} \\
(A_t)_{1,9} &= 2 {h} {u^{y}} \\
(A_t)_{1,10} &= \frac{\left(u^{y}\right)^2-\left(u^{x}\right)^2}{\left(u^{y}\right)^2+1} \\
(A_t)_{1,11} &= -\frac{\left(u^{x}\right)^2}{\left(u^{y}\right)^2+1} \\
(A_t)_{1,12} &= \frac{2 {u^{x}} {u^{y}}}{\left(u^{y}\right)^2+1} \\
(A_t)_{1,13} &= \left(u^{x}\right)^2+\left(u^{y}\right)^2 \\
(A_t)_{2,1} &= \frac{\left(u^{\tau}\right)^2 {u^{x}} \left(16 {a_{\mathrm{SB}}} {T}^3+{(\partial_{T} n)} {h}\right)+{h'(T)} \left({n} \left(u^{\tau}\right)^2 {u^{x}}+{\nu^{x}} \left(2 \left(u^{x}\right)^2+\left(u^{y}\right)^2+1\right)+{\nu^{y}} {u^{x}} {u^{y}}\right)}{{u^{\tau}}} \\
(A_t)_{2,2} &= \frac{1}{\left(u^{\tau}\right)^3 \left(\left(u^{y}\right)^2+1\right)}\Bigl(- {\pi^{\eta}{}_{\eta}} - {\pi^{yy}} - {\pi_{Q}^{yy}} - {\pi_{Q}{}^{\eta}{}_{\eta}} + 4 {a_{\mathrm{SB}}} {T}^4 - 3 {\pi^{\eta}{}_{\eta}} \left(u^{x}\right)^2 - 3 {\pi^{yy}} \left(u^{x}\right)^2\\
  &\quad - 3 {\pi_{Q}^{yy}} \left(u^{x}\right)^2 - 3 {\pi_{Q}{}^{\eta}{}_{\eta}} \left(u^{x}\right)^2 + 12 {a_{\mathrm{SB}}} {T}^4 \left(u^{x}\right)^2 - 2 {\pi^{\eta}{}_{\eta}} \left(u^{x}\right)^4\\
  &\quad - 2 {\pi^{yy}} \left(u^{x}\right)^4 - 2 {\pi_{Q}^{yy}} \left(u^{x}\right)^4 - 2 {\pi_{Q}{}^{\eta}{}_{\eta}} \left(u^{x}\right)^4 + 8 {a_{\mathrm{SB}}} {T}^4 \left(u^{x}\right)^4\\
  &\quad + 3 {\pi^{xy}} {u^{x}} {u^{y}} + 3 {\pi_{Q}^{xy}} {u^{x}} {u^{y}} + 2 {\pi^{xy}} \left(u^{x}\right)^3 {u^{y}} + 2 {\pi_{Q}^{xy}} \left(u^{x}\right)^3 {u^{y}}\\
  &\quad - 2 {\pi^{\eta}{}_{\eta}} \left(u^{y}\right)^2 - {\pi^{yy}} \left(u^{y}\right)^2 - {\pi_{Q}^{yy}} \left(u^{y}\right)^2 - 2 {\pi_{Q}{}^{\eta}{}_{\eta}} \left(u^{y}\right)^2\\
  &\quad + 12 {a_{\mathrm{SB}}} {T}^4 \left(u^{y}\right)^2 - 3 {\pi^{\eta}{}_{\eta}} \left(u^{x}\right)^2 \left(u^{y}\right)^2 - 3 {\pi^{yy}} \left(u^{x}\right)^2 \left(u^{y}\right)^2\\
  &\quad - 3 {\pi_{Q}^{yy}} \left(u^{x}\right)^2 \left(u^{y}\right)^2 - 3 {\pi_{Q}{}^{\eta}{}_{\eta}} \left(u^{x}\right)^2 \left(u^{y}\right)^2 + 24 {a_{\mathrm{SB}}} {T}^4 \left(u^{x}\right)^2 \left(u^{y}\right)^2\\
  &\quad + 8 {a_{\mathrm{SB}}} {T}^4 \left(u^{x}\right)^4 \left(u^{y}\right)^2 + 3 {\pi^{xy}} {u^{x}} \left(u^{y}\right)^3 + 3 {\pi_{Q}^{xy}} {u^{x}} \left(u^{y}\right)^3 - {\pi^{\eta}{}_{\eta}} \left(u^{y}\right)^4\\
  &\quad - {\pi_{Q}{}^{\eta}{}_{\eta}} \left(u^{y}\right)^4 + 12 {a_{\mathrm{SB}}} {T}^4 \left(u^{y}\right)^4 + 12 {a_{\mathrm{SB}}} {T}^4 \left(u^{x}\right)^2 \left(u^{y}\right)^4 + 4 {a_{\mathrm{SB}}} {T}^4 \left(u^{y}\right)^6\\
  &\quad + {\Pi_{Q}} \left(\left(u^{y}\right)^2+1\right)\Bigl(2 \left(u^{x}\right)^4+3 \left(u^{x}\right)^2 \left(\left(u^{y}\right)^2+1\right)+\left(\left(u^{y}\right)^2+1\right)^2\Bigr)\\
  &\quad + {h} \left(\left(u^{y}\right)^2+1\right)\Bigl({\nu^{y}} {u^{y}} \left(\left(u^{y}\right)^2+1\right)\\
  &\quad + {\nu^{x}} {u^{x}} \left(2 \left(u^{x}\right)^2+3 \left(u^{y}\right)^2+3\right)\\
  &\quad + {n}\Bigl(2 \left(u^{x}\right)^4+3 \left(u^{x}\right)^2 \left(\left(u^{y}\right)^2+1\right)+\left(\left(u^{y}\right)^2+1\right)^2\Bigr)\Bigr)\Bigr) \\
(A_t)_{2,3} &= \frac{1}{\left(u^{\tau}\right)^3 \left(\left(u^{y}\right)^2+1\right)^2}\Bigl({\Pi_{Q}} \left(u^{\tau}\right)^2 {u^{x}} {u^{y}} \left(\left(u^{y}\right)^2+1\right)^2 + {n} \left(u^{\tau}\right)^2 {u^{x}} {u^{y}} \left(\left(u^{y}\right)^2+1\right)^2 ({h}-{T})\\
  &\quad + {n} {T} \left(u^{\tau}\right)^2 {u^{x}} {u^{y}} \left(\left(u^{y}\right)^2+1\right)^2 + 4 {a_{\mathrm{SB}}} {T}^4 \left(u^{\tau}\right)^2 {u^{x}} {u^{y}} \left(\left(u^{y}\right)^2+1\right)^2\\
  &\quad + {h} \left(\left(u^{y}\right)^2+1\right)^2 \left({\nu^{x}} {u^{y}} \left(\left(u^{y}\right)^2+1\right)+{\nu^{y}} \left(\left(u^{x}\right)^3+{u^{x}}\right)\right)\\
  &\quad - {u^{y}} \left(\left(u^{y}\right)^2+1\right)\Bigl(- {\pi^{\eta}{}_{\eta}} \left(u^{\tau}\right)^2 {u^{x}} - {\pi^{yy}} \left(\left(u^{x}\right)^3+{u^{x}}\right)\\
  &\quad + {\pi^{xy}} {u^{y}} \left(2 \left(u^{x}\right)^2+\left(u^{y}\right)^2+1\right)\Bigr)\\
  &\quad - {u^{y}} \left(\left(u^{y}\right)^2+1\right)\Bigl(- {\pi_{Q}{}^{\eta}{}_{\eta}} \left(u^{\tau}\right)^2 {u^{x}} - {\pi_{Q}^{yy}} \left(\left(u^{x}\right)^3+{u^{x}}\right)\\
  &\quad + {\pi_{Q}^{xy}} {u^{y}} \left(2 \left(u^{x}\right)^2+\left(u^{y}\right)^2+1\right)\Bigr)\\
  &\quad + \left(u^{\tau}\right)^2\Bigl(2 {u^{x}} {u^{y}} \left({\pi^{\eta}{}_{\eta}} \left(u^{x}\right)^2+{\pi^{yy}} \left(u^{x}\right)^2+{\pi^{yy}}\right)\\
  &\quad + {\pi^{xy}} \left(\left(\left(u^{y}\right)^2+1\right)^2-2 \left(u^{x}\right)^2 \left(\left(u^{y}\right)^2-1\right)\right)\Bigr)\\
  &\quad + \left(u^{\tau}\right)^2\Bigl(2 {u^{x}} {u^{y}} \left({\pi_{Q}^{yy}} \left(u^{x}\right)^2+{\pi_{Q}^{yy}}+{\pi_{Q}{}^{\eta}{}_{\eta}} \left(u^{x}\right)^2\right)\\
  &\quad + {\pi_{Q}^{xy}} \left(\left(\left(u^{y}\right)^2+1\right)^2-2 \left(u^{x}\right)^2 \left(\left(u^{y}\right)^2-1\right)\right)\Bigr)\Bigr) \\
(A_t)_{2,4} &= -\frac{{u^{x}} \left(\left(u^{x}\right)^2+1\right)}{{u^{\tau}} \left(\left(u^{y}\right)^2+1\right)} \\
(A_t)_{2,5} &= -\frac{{u^{\tau}} {u^{x}}}{\left(u^{y}\right)^2+1} \\
(A_t)_{2,6} &= \frac{{u^{y}} \left(2 \left(u^{x}\right)^2+\left(u^{y}\right)^2+1\right)}{{u^{\tau}} \left(\left(u^{y}\right)^2+1\right)} \\
(A_t)_{2,7} &= {(\partial_{\alpha} n)} {h} {u^{\tau}} {u^{x}} \\
(A_t)_{2,8} &= \frac{{h} \left(2 \left(u^{x}\right)^2+\left(u^{y}\right)^2+1\right)}{{u^{\tau}}} \\
(A_t)_{2,9} &= \frac{{h} {u^{x}} {u^{y}}}{{u^{\tau}}} \\
(A_t)_{2,10} &= -\frac{{u^{x}} \left(\left(u^{x}\right)^2+1\right)}{{u^{\tau}} \left(\left(u^{y}\right)^2+1\right)} \\
(A_t)_{2,11} &= -\frac{{u^{\tau}} {u^{x}}}{\left(u^{y}\right)^2+1} \\
(A_t)_{2,12} &= \frac{{u^{y}} \left(2 \left(u^{x}\right)^2+\left(u^{y}\right)^2+1\right)}{{u^{\tau}} \left(\left(u^{y}\right)^2+1\right)} \\
(A_t)_{2,13} &= {u^{\tau}} {u^{x}} \\
(A_t)_{3,1} &= \frac{\left(u^{\tau}\right)^2 {u^{y}} \left(16 {a_{\mathrm{SB}}} {T}^3+{(\partial_{T} n)} {h}\right)+{h'(T)} \left({n} \left(u^{\tau}\right)^2 {u^{y}}+{\nu^{x}} {u^{x}} {u^{y}}+{\nu^{y}} \left(\left(u^{x}\right)^2+2 \left(u^{y}\right)^2+1\right)\right)}{{u^{\tau}}} \\
(A_t)_{3,2} &= \frac{1}{\left(u^{\tau}\right)^3}\Bigl({\pi^{xy}} + {\pi_{Q}^{xy}} + {h} {\nu^{y}} {u^{x}} + {h} {\nu^{y}} \left(u^{x}\right)^3 + {h} {\nu^{x}} {u^{y}} + {h} {n} {u^{x}} {u^{y}} + {\Pi_{Q}} {u^{x}} {u^{y}} - {\pi^{yy}} {u^{x}} {u^{y}}\\
  &\quad - {\pi_{Q}^{yy}} {u^{x}} {u^{y}} + 4 {a_{\mathrm{SB}}} {T}^4 {u^{x}} {u^{y}} + {h} {n} \left(u^{x}\right)^3 {u^{y}} + {\Pi_{Q}} \left(u^{x}\right)^3 {u^{y}} + 4 {a_{\mathrm{SB}}} {T}^4 \left(u^{x}\right)^3 {u^{y}}\\
  &\quad + {\pi^{xy}} \left(u^{y}\right)^2 + {\pi_{Q}^{xy}} \left(u^{y}\right)^2 + {h} {\nu^{x}} \left(u^{y}\right)^3 + {h} {n} {u^{x}} \left(u^{y}\right)^3 + {\Pi_{Q}} {u^{x}} \left(u^{y}\right)^3\\
  &\quad + 4 {a_{\mathrm{SB}}} {T}^4 {u^{x}} \left(u^{y}\right)^3\Bigr) \\
(A_t)_{3,3} &= \frac{1}{\left(u^{\tau}\right)^3}\Bigl({\pi^{yy}} + {\pi_{Q}^{yy}} + 4 {a_{\mathrm{SB}}} {T}^4 + {\pi^{yy}} \left(u^{x}\right)^2 + {\pi_{Q}^{yy}} \left(u^{x}\right)^2 + 8 {a_{\mathrm{SB}}} {T}^4 \left(u^{x}\right)^2 + 4 {a_{\mathrm{SB}}} {T}^4 \left(u^{x}\right)^4\\
  &\quad - {\pi^{xy}} {u^{x}} {u^{y}} - {\pi_{Q}^{xy}} {u^{x}} {u^{y}} + 12 {a_{\mathrm{SB}}} {T}^4 \left(u^{y}\right)^2 + 12 {a_{\mathrm{SB}}} {T}^4 \left(u^{x}\right)^2 \left(u^{y}\right)^2 + 8 {a_{\mathrm{SB}}} {T}^4 \left(u^{y}\right)^4\\
  &\quad + {\Pi_{Q}} \left(\left(u^{x}\right)^4+\left(u^{x}\right)^2 \left(3 \left(u^{y}\right)^2+2\right)+2 \left(u^{y}\right)^4+3 \left(u^{y}\right)^2+1\right)\\
  &\quad + {h}\Bigl({\nu^{x}} \left(\left(u^{x}\right)^3+{u^{x}}\right) + {\nu^{y}} {u^{y}} \left(3 \left(u^{x}\right)^2+2 \left(u^{y}\right)^2+3\right)\\
  &\quad + {n} \left(\left(u^{x}\right)^4+\left(u^{x}\right)^2 \left(3 \left(u^{y}\right)^2+2\right)+2 \left(u^{y}\right)^4+3 \left(u^{y}\right)^2+1\right)\Bigr)\Bigr) \\
(A_t)_{3,4} &= \frac{{u^{y}}}{{u^{\tau}}} \\
(A_t)_{3,6} &= \frac{{u^{x}}}{{u^{\tau}}} \\
(A_t)_{3,7} &= {(\partial_{\alpha} n)} {h} {u^{\tau}} {u^{y}} \\
(A_t)_{3,8} &= \frac{{h} {u^{x}} {u^{y}}}{{u^{\tau}}} \\
(A_t)_{3,9} &= \frac{{h} \left(\left(u^{x}\right)^2+2 \left(u^{y}\right)^2+1\right)}{{u^{\tau}}} \\
(A_t)_{3,10} &= \frac{{u^{y}}}{{u^{\tau}}} \\
(A_t)_{3,12} &= \frac{{u^{x}}}{{u^{\tau}}} \\
(A_t)_{3,13} &= {u^{\tau}} {u^{y}} \\
(A_t)_{4,2} &= \frac{3 {\delta_{\pi\pi}} {\pi^{yy}} {u^{x}}-6 {\pi^{xy}} {\tau_{\pi}} {u^{y}} \left(\left(u^{y}\right)^2+1\right)-2 {u^{x}} \left(-3 {\pi^{yy}} {\tau_{\pi}} \left(u^{y}\right)^2+\left(u^{y}\right)^2 {\eta}+{\eta}\right)}{3 {u^{\tau}}} \\
(A_t)_{4,3} &= \frac{{u^{y}} \left(3 {\delta_{\pi\pi}} {\pi^{yy}}+6 {\pi^{xy}} {\tau_{\pi}} {u^{x}} {u^{y}}-6 {\pi^{yy}} {\tau_{\pi}} \left(\left(u^{x}\right)^2+1\right)+6 \left(u^{x}\right)^2 {\eta}+4 \left(u^{y}\right)^2 {\eta}+4 {\eta}\right)}{3 {u^{\tau}}} \\
(A_t)_{4,4} &= {\tau_{\pi}} {u^{\tau}} \\
(A_t)_{5,2} &= \frac{{u^{x}} (3 {\delta_{\pi\pi}} {\pi^{\eta}{}_{\eta}}-2 {\eta})}{3 {\tau}^2 {u^{\tau}}} \\
(A_t)_{5,3} &= \frac{{u^{y}} (3 {\delta_{\pi\pi}} {\pi^{\eta}{}_{\eta}}-2 {\eta})}{3 {\tau}^2 {u^{\tau}}} \\
(A_t)_{5,5} &= \frac{{\tau_{\pi}} {u^{\tau}}}{{\tau}^2} \\
(A_t)_{6,2} &= \frac{1}{3 {u^{\tau}}}\Bigl(3 {\delta_{\pi\pi}} {\pi^{xy}} {u^{x}} - 3 {\pi^{xy}} {\tau_{\pi}} {u^{x}} \left(2 \left(u^{y}\right)^2+1\right)\\
  &\quad + {u^{y}} \left(3 {\pi^{\eta}{}_{\eta}} {\tau_{\pi}} \left(u^{\tau}\right)^2+3 {\pi^{yy}} \left(2 {\tau_{\pi}} \left(u^{x}\right)^2+{\tau_{\pi}}\right)+{\eta} \left(\left(u^{x}\right)^2+3 \left(u^{y}\right)^2+3\right)\right)\Bigr) \\
(A_t)_{6,3} &= \frac{1}{3 {u^{\tau}} \left(\left(u^{y}\right)^2+1\right)}\Bigl(- 3 {\pi^{xy}} {\tau_{\pi}} {u^{y}} - 3 {\pi^{\eta}{}_{\eta}} {\tau_{\pi}} {u^{x}} \left(u^{y}\right)^2 - 3 {\pi^{\eta}{}_{\eta}} {\tau_{\pi}} \left(u^{x}\right)^3 \left(u^{y}\right)^2 - 3 {\pi^{xy}} {\tau_{\pi}} \left(u^{y}\right)^3\\
  &\quad + 6 {\pi^{xy}} {\tau_{\pi}} \left(u^{x}\right)^2 \left(u^{y}\right)^3 - 3 {\pi^{\eta}{}_{\eta}} {\tau_{\pi}} {u^{x}} \left(u^{y}\right)^4 + 3 {\delta_{\pi\pi}} {\pi^{xy}} {u^{y}} \left(\left(u^{y}\right)^2+1\right)\\
  &\quad - 3 {\pi^{yy}} {\tau_{\pi}} {u^{x}} \left(\left(u^{x}\right)^2+1\right) \left(2 \left(u^{y}\right)^2+1\right) + 3 {u^{x}} {\eta} + 3 \left(u^{x}\right)^3 {\eta}\\
  &\quad + 4 {u^{x}} \left(u^{y}\right)^2 {\eta} + 3 \left(u^{x}\right)^3 \left(u^{y}\right)^2 {\eta} + {u^{x}} \left(u^{y}\right)^4 {\eta}\Bigr) \\
(A_t)_{6,6} &= {\tau_{\pi}} {u^{\tau}} \\
(A_t)_{7,1} &= {(\partial_{T} n)} {u^{\tau}} \\
(A_t)_{7,2} &= \frac{{n} \left(u^{\tau}\right)^2 {u^{x}}+{\nu^{x}} \left(\left(u^{y}\right)^2+1\right)-{\nu^{y}} {u^{x}} {u^{y}}}{\left(u^{\tau}\right)^3} \\
(A_t)_{7,3} &= \frac{{n} \left(u^{\tau}\right)^2 {u^{y}}-{\nu^{x}} {u^{x}} {u^{y}}+{\nu^{y}} \left(\left(u^{x}\right)^2+1\right)}{\left(u^{\tau}\right)^3} \\
(A_t)_{7,7} &= {(\partial_{\alpha} n)} {u^{\tau}} \\
(A_t)_{7,8} &= \frac{{u^{x}}}{{u^{\tau}}} \\
(A_t)_{7,9} &= \frac{{u^{y}}}{{u^{\tau}}} \\
(A_t)_{8,1} &= \frac{{D_{s}} {u^{\tau}} ({h'(T)} {\nu^{x}}+{u^{x}} ({(\partial_{T} n)} {T}+{n}))}{{T}} \\
(A_t)_{8,2} &= \frac{1}{{T} \left(u^{\tau}\right)^3 \left(\left(u^{y}\right)^2+1\right)}{D_{s}}\Bigl(- {\pi_{Q}^{yy}} - {\pi_{Q}{}^{\eta}{}_{\eta}} - 2 {\pi_{Q}^{yy}} \left(u^{x}\right)^2 - 2 {\pi_{Q}{}^{\eta}{}_{\eta}} \left(u^{x}\right)^2 - {\pi_{Q}^{yy}} \left(u^{x}\right)^4\\
  &\quad - {\pi_{Q}{}^{\eta}{}_{\eta}} \left(u^{x}\right)^4 + 2 {\pi_{Q}^{xy}} {u^{x}} {u^{y}} + {\pi_{Q}^{xy}} \left(u^{x}\right)^3 {u^{y}} - {\pi_{Q}^{yy}} \left(u^{y}\right)^2\\
  &\quad - 2 {\pi_{Q}{}^{\eta}{}_{\eta}} \left(u^{y}\right)^2 - {\pi_{Q}^{yy}} \left(u^{x}\right)^2 \left(u^{y}\right)^2 - {\pi_{Q}{}^{\eta}{}_{\eta}} \left(u^{x}\right)^2 \left(u^{y}\right)^2\\
  &\quad + {\pi_{Q}^{yy}} \left(u^{x}\right)^4 \left(u^{y}\right)^2 + {\pi_{Q}{}^{\eta}{}_{\eta}} \left(u^{x}\right)^4 \left(u^{y}\right)^2 + {\pi_{Q}^{xy}} {u^{x}} \left(u^{y}\right)^3\\
  &\quad - {\pi_{Q}^{xy}} \left(u^{x}\right)^3 \left(u^{y}\right)^3 - {\pi_{Q}{}^{\eta}{}_{\eta}} \left(u^{y}\right)^4 + {\pi_{Q}^{yy}} \left(u^{x}\right)^2 \left(u^{y}\right)^4\\
  &\quad + {\pi_{Q}{}^{\eta}{}_{\eta}} \left(u^{x}\right)^2 \left(u^{y}\right)^4 - {\pi_{Q}^{xy}} {u^{x}} \left(u^{y}\right)^5 + {\Pi_{Q}} \left(u^{\tau}\right)^4 \left(\left(u^{y}\right)^2+1\right)\\
  &\quad + {h} \left(u^{\tau}\right)^2 \left(\left(u^{y}\right)^2+1\right)\Bigl({n} \left(u^{\tau}\right)^2+{\nu^{x}} \left({u^{x}}-{u^{x}} \left(u^{y}\right)^2\right)+{\nu^{y}} \left(\left(u^{x}\right)^2+1\right) {u^{y}}\Bigr)\Bigr) \\
(A_t)_{8,3} &= \frac{1}{{T} \left(u^{\tau}\right)^3 \left(\left(u^{y}\right)^2+1\right)^2}{D_{s}}\Bigl(- {h} \left(u^{\tau}\right)^2 \left(\left(u^{x}\right)^2+1\right) \left(\left(u^{y}\right)^2+1\right)^2 ({\nu^{y}} {u^{x}}-{\nu^{x}} {u^{y}})\\
  &\quad + {\pi_{Q}^{xy}}\Bigl(\left(u^{x}\right)^4 \left(\left(u^{y}\right)^2-1\right)^2 + \left(\left(u^{y}\right)^2+1\right)^2\\
  &\quad + \left(u^{x}\right)^2 \left(\left(u^{y}\right)^6-2 \left(u^{y}\right)^4-\left(u^{y}\right)^2+2\right)\Bigr)\\
  &\quad + {u^{x}} {u^{y}}\Bigl(- {\pi_{Q}^{yy}} \left(\left(u^{x}\right)^2+1\right)\Bigl(\left(u^{x}\right)^2 \left(\left(u^{y}\right)^2-1\right)+\left(u^{y}\right)^4-\left(u^{y}\right)^2-2\Bigr)\\
  &\quad + {\pi_{Q}{}^{\eta}{}_{\eta}}\Bigl(- \left(u^{x}\right)^4 \left(\left(u^{y}\right)^2-1\right) + \left(\left(u^{y}\right)^2+1\right)^2\\
  &\quad + \left(u^{x}\right)^2 \left(-\left(u^{y}\right)^4+\left(u^{y}\right)^2+2\right)\Bigr)\Bigr)\Bigr) \\
(A_t)_{8,7} &= {(\partial_{\alpha} n)} {D_{s}} {u^{\tau}} {u^{x}} \\
(A_t)_{8,8} &= \frac{{D_{s}} {h} {u^{\tau}}}{{T}} \\
(A_t)_{8,10} &= -\frac{{D_{s}} {u^{x}} \left(\left(u^{x}\right)^2+1\right)}{{T} {u^{\tau}} \left(\left(u^{y}\right)^2+1\right)} \\
(A_t)_{8,11} &= -\frac{{D_{s}} {u^{\tau}} {u^{x}}}{{T} \left(u^{y}\right)^2+{T}} \\
(A_t)_{8,12} &= \frac{{D_{s}} \left(2 \left(u^{x}\right)^2 {u^{y}}+\left(u^{y}\right)^3+{u^{y}}\right)}{{T} {u^{\tau}} \left(\left(u^{y}\right)^2+1\right)} \\
(A_t)_{8,13} &= \frac{{D_{s}} {u^{\tau}} {u^{x}}}{{T}} \\
(A_t)_{9,1} &= \frac{{D_{s}} {u^{\tau}} ({h'(T)} {\nu^{y}}+{u^{y}} ({(\partial_{T} n)} {T}+{n}))}{{T}} \\
(A_t)_{9,2} &= \frac{1}{{T} \left(u^{\tau}\right)^3}{D_{s}}\Bigl({h} \left(u^{\tau}\right)^2 \left(\left(u^{y}\right)^2+1\right) ({\nu^{y}} {u^{x}}-{\nu^{x}} {u^{y}}) - {\pi_{Q}^{xy}} \left(\left(u^{x}\right)^2 \left(u^{y}\right)^2+\left(u^{y}\right)^4-1\right)\\
  &\quad + {u^{x}} {u^{y}} \left({\pi_{Q}^{yy}} \left(\left(u^{x}\right)^2+\left(u^{y}\right)^2\right)+{\pi_{Q}{}^{\eta}{}_{\eta}} \left(u^{\tau}\right)^2\right)\Bigr) \\
(A_t)_{9,3} &= \frac{1}{{T} \left(u^{\tau}\right)^3 \left(\left(u^{y}\right)^2+1\right)}{D_{s}}\Bigl({\pi_{Q}^{yy}} + {\pi_{Q}^{yy}} \left(u^{x}\right)^2 - 2 {\pi_{Q}^{xy}} {u^{x}} {u^{y}} - {\pi_{Q}^{xy}} \left(u^{x}\right)^3 {u^{y}} - {\pi_{Q}^{yy}} \left(u^{x}\right)^2 \left(u^{y}\right)^2\\
  &\quad - {\pi_{Q}{}^{\eta}{}_{\eta}} \left(u^{x}\right)^2 \left(u^{y}\right)^2 - {\pi_{Q}^{yy}} \left(u^{x}\right)^4 \left(u^{y}\right)^2 - {\pi_{Q}{}^{\eta}{}_{\eta}} \left(u^{x}\right)^4 \left(u^{y}\right)^2\\
  &\quad - {\pi_{Q}^{xy}} {u^{x}} \left(u^{y}\right)^3 + {\pi_{Q}^{xy}} \left(u^{x}\right)^3 \left(u^{y}\right)^3 - {\pi_{Q}^{yy}} \left(u^{y}\right)^4 - {\pi_{Q}^{yy}} \left(u^{x}\right)^2 \left(u^{y}\right)^4\\
  &\quad - {\pi_{Q}{}^{\eta}{}_{\eta}} \left(u^{x}\right)^2 \left(u^{y}\right)^4 + {\pi_{Q}^{xy}} {u^{x}} \left(u^{y}\right)^5 + {\Pi_{Q}} \left(u^{\tau}\right)^4 \left(\left(u^{y}\right)^2+1\right)\\
  &\quad + {h} \left(u^{\tau}\right)^2 \left(\left(u^{y}\right)^2+1\right)\Bigl({n} \left(u^{\tau}\right)^2+{\nu^{x}} {u^{x}} \left(\left(u^{y}\right)^2+1\right)-{\nu^{y}} \left(\left(u^{x}\right)^2-1\right) {u^{y}}\Bigr)\Bigr) \\
(A_t)_{9,7} &= {(\partial_{\alpha} n)} {D_{s}} {u^{\tau}} {u^{y}} \\
(A_t)_{9,9} &= \frac{{D_{s}} {h} {u^{\tau}}}{{T}} \\
(A_t)_{9,10} &= \frac{{D_{s}} {u^{y}}}{{T} {u^{\tau}}} \\
(A_t)_{9,12} &= \frac{{D_{s}} {u^{x}}}{{T} {u^{\tau}}} \\
(A_t)_{9,13} &= \frac{{D_{s}} {u^{\tau}} {u^{y}}}{{T}} \\
(A_t)_{10,1} &= \frac{1}{6 {T} {u^{\tau}}}{D_{s}}\Bigl(3 {\mathcal{C}'(T)} {\pi_{Q}^{yy}} \left(u^{\tau}\right)^2 - 2 {\mathcal{B}'(T)} {\nu^{x}} {u^{x}} \left(\left(u^{y}\right)^2+1\right)\\
  &\quad + 2 {\mathcal{B}'(T)} {\nu^{y}} {u^{y}} \left(3 \left(u^{x}\right)^2+2 \left(u^{y}\right)^2+2\right)\Bigr) \\
(A_t)_{10,2} &= -\frac{1}{6 {T} \left(u^{\tau}\right)^3}{D_{s}}\Bigl(2 {\mathcal{A}} \left(u^{\tau}\right)^2 \left(\left(u^{y}\right)^2+1\right) \left({\nu^{x}} \left(u^{y}\right)^2+{\nu^{x}}-{\nu^{y}} {u^{x}} {u^{y}}\right)\\
  &\quad + 2 {\mathcal{B}} \left(\left(3 \left(u^{x}\right)^2+4\right) \left(u^{y}\right)^2+3 \left(u^{y}\right)^4+1\right) \left({\nu^{x}} \left(u^{y}\right)^2+{\nu^{x}}-{\nu^{y}} {u^{x}} {u^{y}}\right)\\
  &\quad + \left(u^{\tau}\right)^2\Bigl(2 {I_{31}} {u^{x}} \left(\left(u^{y}\right)^2+1\right)\\
  &\quad + {\mathcal{C}}\Bigl(2 {\Pi_{Q}} {u^{x}} \left(\left(u^{y}\right)^2+1\right)\\
  &\quad - 2 \left(\left(u^{y}\right)^2+1\right) ({\pi_{Q}{}^{\eta}{}_{\eta}} {u^{x}}-4 {\pi_{Q}^{xy}} {u^{y}})\\
  &\quad - {\pi_{Q}^{yy}} {u^{x}} \left(8 \left(u^{y}\right)^2+5\right)\Bigr)\Bigr)\Bigr) \\
(A_t)_{10,3} &= \frac{1}{6 {T} \left(u^{\tau}\right)^3}{D_{s}}\Bigl(2 {\mathcal{B}} {\nu^{x}} {u^{x}} {u^{y}} \left(\left(3 \left(u^{x}\right)^2+4\right) \left(u^{y}\right)^2+3 \left(u^{y}\right)^4+1\right)\\
  &\quad - 2{\mathcal{A}} \left(u^{\tau}\right)^2\Bigl({\nu^{y}} \left(u^{x}\right)^2 \left(\left(u^{y}\right)^2-2\right) - 2 {\nu^{y}} \left(\left(u^{y}\right)^2+1\right)\\
  &\quad - {\nu^{x}} {u^{x}} {u^{y}} \left(\left(u^{y}\right)^2+1\right)\Bigr)\\
  &\quad - 2{\mathcal{B}} {\nu^{y}}\Bigl(3 \left(u^{x}\right)^4 \left(\left(u^{y}\right)^2-1\right) - 2 \left(\left(u^{y}\right)^2+1\right)\\
  &\quad + \left(u^{x}\right)^2 \left(3 \left(u^{y}\right)^4+\left(u^{y}\right)^2-5\right)\Bigr)\\
  &\quad + \left(u^{\tau}\right)^2\Bigl(2 {I_{31}} {u^{y}} \left(3 \left(u^{x}\right)^2+2 \left(u^{y}\right)^2+2\right) + 4 {\mathcal{C}} {\pi_{Q}^{xy}} \left(2 {u^{x}} \left(u^{y}\right)^2+{u^{x}}\right)\\
  &\quad + {\mathcal{C}} {u^{y}}\Bigl({\pi_{Q}^{yy}} - 8 {\pi_{Q}^{yy}} \left(u^{x}\right)^2 - 2 {\pi_{Q}{}^{\eta}{}_{\eta}} \left(u^{x}\right)^2\\
  &\quad + {\Pi_{Q}} \left(6 \left(u^{x}\right)^2+4 \left(u^{y}\right)^2+4\right)\Bigr)\Bigr)\Bigr) \\
(A_t)_{10,8} &= -\frac{{\mathcal{B}} {D_{s}} {u^{x}} \left(\left(u^{y}\right)^2+1\right)}{3 {T} {u^{\tau}}} \\
(A_t)_{10,9} &= \frac{{\mathcal{B}} {D_{s}} {u^{y}} \left(3 \left(u^{x}\right)^2+2 \left(u^{y}\right)^2+2\right)}{3 {T} {u^{\tau}}} \\
(A_t)_{10,10} &= \frac{{\mathcal{C}} {D_{s}} {u^{\tau}}}{2 {T}} \\
(A_t)_{11,1} &= \frac{{D_{s}} \left(3 {\mathcal{C}'(T)} {\pi_{Q}{}^{\eta}{}_{\eta}} \left(u^{\tau}\right)^2-2 {\mathcal{B}'(T)} ({\nu^{x}} {u^{x}}+{\nu^{y}} {u^{y}})\right)}{6 {T} {u^{\tau}}} \\
(A_t)_{11,2} &= -\frac{1}{6 {T} \left(u^{\tau}\right)^3}{D_{s}}\Bigl(2 {\mathcal{B}} \left({\nu^{x}} \left(u^{y}\right)^2+{\nu^{x}}-{\nu^{y}} {u^{x}} {u^{y}}\right) + 2 {\mathcal{A}} \left(u^{\tau}\right)^2 \left({\nu^{x}} \left(u^{y}\right)^2+{\nu^{x}}-{\nu^{y}} {u^{x}} {u^{y}}\right)\\
  &\quad + \left(u^{\tau}\right)^2 ({\mathcal{C}} (2 {\Pi_{Q}} {u^{x}}+2 {\pi_{Q}^{xy}} {u^{y}}-2 {\pi_{Q}^{yy}} {u^{x}}-5 {\pi_{Q}{}^{\eta}{}_{\eta}} {u^{x}})+2 {I_{31}} {u^{x}})\Bigr) \\
(A_t)_{11,3} &= -\frac{1}{6 {T} \left(u^{\tau}\right)^3 \left(\left(u^{y}\right)^2+1\right)}{D_{s}}\Bigl(2 {\mathcal{B}} \left(\left(u^{y}\right)^2+1\right) \left(-{\nu^{x}} {u^{x}} {u^{y}}+{\nu^{y}} \left(u^{x}\right)^2+{\nu^{y}}\right)\\
  &\quad + 2 {\mathcal{A}} \left(u^{\tau}\right)^2 \left(\left(u^{y}\right)^2+1\right) \left(-{\nu^{x}} {u^{x}} {u^{y}}+{\nu^{y}} \left(u^{x}\right)^2+{\nu^{y}}\right)\\
  &\quad + \left(u^{\tau}\right)^2\Bigl(- 2 {\mathcal{C}} {\pi_{Q}^{xy}} {u^{x}} \left(\left(u^{y}\right)^2-1\right) + 2 {I_{31}} {u^{y}} \left(\left(u^{y}\right)^2+1\right)\\
  &\quad + {\mathcal{C}} {u^{y}}\Bigl(2 {\pi_{Q}^{yy}} \left(\left(u^{x}\right)^2+1\right)\\
  &\quad + {\pi_{Q}{}^{\eta}{}_{\eta}} \left(2 \left(u^{x}\right)^2-3 \left(u^{y}\right)^2-3\right) + 2 {\Pi_{Q}} \left(\left(u^{y}\right)^2+1\right)\Bigr)\Bigr)\Bigr) \\
(A_t)_{11,8} &= -\frac{{\mathcal{B}} {D_{s}} {u^{x}}}{3 {T} {u^{\tau}}} \\
(A_t)_{11,9} &= -\frac{{\mathcal{B}} {D_{s}} {u^{y}}}{3 {T} {u^{\tau}}} \\
(A_t)_{11,11} &= \frac{{\mathcal{C}} {D_{s}} {u^{\tau}}}{2 {T}} \\
(A_t)_{12,1} &= \frac{1}{6 {T} {u^{\tau}}}{D_{s}}\Bigl(3 {\mathcal{C}'(T)} {\pi_{Q}^{xy}} \left(u^{\tau}\right)^2 + {\mathcal{B}'(T)} {\nu^{y}} {u^{x}} \left(3 \left(u^{x}\right)^2+\left(u^{y}\right)^2+3\right)\\
  &\quad + {\mathcal{B}'(T)} {\nu^{x}} {u^{y}} \left(\left(u^{x}\right)^2+3 \left(u^{y}\right)^2+3\right)\Bigr) \\
(A_t)_{12,2} &= \frac{1}{6 {T} \left(u^{\tau}\right)^3}{D_{s}}\Bigl({\mathcal{A}} \left(u^{\tau}\right)^2\Bigl(3 {\nu^{y}} \left(\left(u^{y}\right)^2+1\right) - 2 {\nu^{x}} {u^{x}} {u^{y}} \left(\left(u^{y}\right)^2+1\right)\\
  &\quad + {\nu^{y}} \left(u^{x}\right)^2 \left(2 \left(u^{y}\right)^2+3\right)\Bigr)\\
  &\quad + \left(u^{\tau}\right)^2\Bigl({\mathcal{C}} {\pi_{Q}^{xy}} {u^{x}} \left(3-8 \left(u^{y}\right)^2\right) + {I_{31}} {u^{y}} \left(\left(u^{x}\right)^2+3 \left(u^{y}\right)^2+3\right)\\
  &\quad + {\mathcal{C}} {u^{y}}\Bigl({\pi_{Q}^{yy}} \left(8 \left(u^{x}\right)^2+6\right) + {\Pi_{Q}} \left(\left(u^{x}\right)^2+3 \left(u^{y}\right)^2+3\right)\\
  &\quad + {\pi_{Q}{}^{\eta}{}_{\eta}} \left(5 \left(u^{x}\right)^2+3 \left(u^{y}\right)^2+3\right)\Bigr)\Bigr)\\
  &\quad + {\mathcal{B}}\Bigl(- 2 {\nu^{x}} {u^{x}} {u^{y}} \left(\left(u^{y}\right)^2+1\right) \left(3 \left(u^{x}\right)^2+3 \left(u^{y}\right)^2+4\right)\\
  &\quad + {\nu^{y}}\Bigl(3 \left(\left(u^{y}\right)^2+1\right)^2 + \left(u^{x}\right)^4 \left(6 \left(u^{y}\right)^2+3\right)\\
  &\quad + 2 \left(u^{x}\right)^2 \left(3 \left(u^{y}\right)^4+7 \left(u^{y}\right)^2+3\right)\Bigr)\Bigr)\Bigr) \\
(A_t)_{12,3} &= \frac{1}{6 {T} \left(u^{\tau}\right)^3 \left(\left(u^{y}\right)^2+1\right)}{D_{s}}\Bigl({\mathcal{A}} \left(u^{\tau}\right)^2 \left(\left(u^{y}\right)^2+1\right)\Bigl(- 2 {\nu^{y}} {u^{x}} \left(\left(u^{x}\right)^2+1\right) {u^{y}} + 3 {\nu^{x}} \left(\left(u^{y}\right)^2+1\right)\\
  &\quad + {\nu^{x}} \left(u^{x}\right)^2 \left(2 \left(u^{y}\right)^2+3\right)\Bigr)\\
  &\quad + \left(u^{\tau}\right)^2\Bigl({I_{31}} {u^{x}} \left(\left(u^{y}\right)^2+1\right) \left(3 \left(u^{x}\right)^2+\left(u^{y}\right)^2+3\right)\\
  &\quad + {\mathcal{C}}\Bigl(- 3 {\pi_{Q}{}^{\eta}{}_{\eta}} {u^{x}} - 3 {\pi_{Q}{}^{\eta}{}_{\eta}} \left(u^{x}\right)^3 + 3 {\pi_{Q}^{xy}} {u^{y}} + 4 {\pi_{Q}^{xy}} \left(u^{x}\right)^2 {u^{y}}\\
  &\quad - 6 {\pi_{Q}{}^{\eta}{}_{\eta}} {u^{x}} \left(u^{y}\right)^2 - 5 {\pi_{Q}{}^{\eta}{}_{\eta}} \left(u^{x}\right)^3 \left(u^{y}\right)^2 + 3 {\pi_{Q}^{xy}} \left(u^{y}\right)^3\\
  &\quad + 8 {\pi_{Q}^{xy}} \left(u^{x}\right)^2 \left(u^{y}\right)^3 - 3 {\pi_{Q}{}^{\eta}{}_{\eta}} {u^{x}} \left(u^{y}\right)^4\\
  &\quad + {\Pi_{Q}} {u^{x}} \left(\left(u^{y}\right)^2+1\right) \left(3 \left(u^{x}\right)^2+\left(u^{y}\right)^2+3\right)\\
  &\quad - 2 {\pi_{Q}^{yy}} {u^{x}} \left(\left(u^{x}\right)^2+1\right) \left(4 \left(u^{y}\right)^2+3\right)\Bigr)\Bigr)\\
  &\quad + {\mathcal{B}} \left(\left(u^{y}\right)^2+1\right)\Bigl(- 2 {\nu^{y}} {u^{x}} \left(\left(u^{x}\right)^2+1\right) {u^{y}} \left(3 \left(u^{x}\right)^2+3 \left(u^{y}\right)^2+4\right)\\
  &\quad + {\nu^{x}}\Bigl(3 \left(\left(u^{y}\right)^2+1\right)^2 + \left(u^{x}\right)^4 \left(6 \left(u^{y}\right)^2+3\right)\\
  &\quad + 2 \left(u^{x}\right)^2 \left(3 \left(u^{y}\right)^4+7 \left(u^{y}\right)^2+3\right)\Bigr)\Bigr)\Bigr) \\
(A_t)_{12,8} &= \frac{{\mathcal{B}} {D_{s}} {u^{y}} \left(\left(u^{x}\right)^2+3 \left(u^{y}\right)^2+3\right)}{6 {T} {u^{\tau}}} \\
(A_t)_{12,9} &= \frac{{\mathcal{B}} {D_{s}} {u^{x}} \left(3 \left(u^{x}\right)^2+\left(u^{y}\right)^2+3\right)}{6 {T} {u^{\tau}}} \\
(A_t)_{12,12} &= \frac{{\mathcal{C}} {D_{s}} {u^{\tau}}}{2 {T}} \\
(A_t)_{13,1} &= \frac{{D_{s}} \left(5 {\mathcal{B}'(T)} ({\nu^{x}} {u^{x}}+{\nu^{y}} {u^{y}})+3 {\mathcal{C}'(T)} {\Pi_{Q}} \left(u^{\tau}\right)^2+3 {(\partial_{T} I_{31})} \left(u^{\tau}\right)^2\right)}{6 {T} {u^{\tau}}} \\
(A_t)_{13,2} &= \frac{1}{6 {T} \left(u^{\tau}\right)^3}{D_{s}}\Bigl(5 {\mathcal{B}} \left({\nu^{x}} \left(u^{y}\right)^2+{\nu^{x}}-{\nu^{y}} {u^{x}} {u^{y}}\right) + 2 {\mathcal{A}} \left(u^{\tau}\right)^2 \left({\nu^{x}} \left(u^{y}\right)^2+{\nu^{x}}-{\nu^{y}} {u^{x}} {u^{y}}\right)\\
  &\quad + \left(u^{\tau}\right)^2 ({\mathcal{C}} (5 {\Pi_{Q}} {u^{x}}+2 {\pi_{Q}^{xy}} {u^{y}}-2 {\pi_{Q}^{yy}} {u^{x}}-2 {\pi_{Q}{}^{\eta}{}_{\eta}} {u^{x}})+5 {I_{31}} {u^{x}})\Bigr) \\
(A_t)_{13,3} &= \frac{1}{6 {T} \left(u^{\tau}\right)^3 \left(\left(u^{y}\right)^2+1\right)}{D_{s}}\Bigl(5 {\mathcal{B}} \left(\left(u^{y}\right)^2+1\right) \left(-{\nu^{x}} {u^{x}} {u^{y}}+{\nu^{y}} \left(u^{x}\right)^2+{\nu^{y}}\right)\\
  &\quad + 2 {\mathcal{A}} \left(u^{\tau}\right)^2 \left(\left(u^{y}\right)^2+1\right) \left(-{\nu^{x}} {u^{x}} {u^{y}}+{\nu^{y}} \left(u^{x}\right)^2+{\nu^{y}}\right)\\
  &\quad + \left(u^{\tau}\right)^2\Bigl(5 {I_{31}} {u^{y}} \left(\left(u^{y}\right)^2+1\right)\\
  &\quad + {\mathcal{C}}\Bigl(2 {u^{y}} \left({\pi_{Q}^{yy}} \left(u^{x}\right)^2+{\pi_{Q}^{yy}}+{\pi_{Q}{}^{\eta}{}_{\eta}} \left(u^{x}\right)^2\right)\\
  &\quad - 2 {\pi_{Q}^{xy}} {u^{x}} \left(\left(u^{y}\right)^2-1\right) + 5 {\Pi_{Q}} {u^{y}} \left(\left(u^{y}\right)^2+1\right)\Bigr)\Bigr)\Bigr) \\
(A_t)_{13,7} &= \frac{{(\partial_{\alpha} I_{31})} {D_{s}} {u^{\tau}}}{2 {T}} \\
(A_t)_{13,8} &= \frac{5 {\mathcal{B}} {D_{s}} {u^{x}}}{6 {T} {u^{\tau}}} \\
(A_t)_{13,9} &= \frac{5 {\mathcal{B}} {D_{s}} {u^{y}}}{6 {T} {u^{\tau}}} \\
(A_t)_{13,13} &= \frac{{\mathcal{C}} {D_{s}} {u^{\tau}}}{2 {T}} \\
(A_x)_{1,1} &= \frac{\left(u^{\tau}\right)^2 {u^{x}} \left(16 {a_{\mathrm{SB}}} {T}^3+{(\partial_{T} n)} {h}\right)+{h'(T)} \left({n} \left(u^{\tau}\right)^2 {u^{x}}+{\nu^{x}} \left(2 \left(u^{x}\right)^2+\left(u^{y}\right)^2+1\right)+{\nu^{y}} {u^{x}} {u^{y}}\right)}{{u^{\tau}}} \\
(A_x)_{1,2} &= \frac{1}{\left(u^{\tau}\right)^3 \left(\left(u^{y}\right)^2+1\right)}\Bigl(- {\pi^{\eta}{}_{\eta}} - {\pi^{yy}} - {\pi_{Q}^{yy}} - {\pi_{Q}{}^{\eta}{}_{\eta}} + 4 {a_{\mathrm{SB}}} {T}^4 - 3 {\pi^{\eta}{}_{\eta}} \left(u^{x}\right)^2 - 3 {\pi^{yy}} \left(u^{x}\right)^2\\
  &\quad - 3 {\pi_{Q}^{yy}} \left(u^{x}\right)^2 - 3 {\pi_{Q}{}^{\eta}{}_{\eta}} \left(u^{x}\right)^2 + 12 {a_{\mathrm{SB}}} {T}^4 \left(u^{x}\right)^2 - 2 {\pi^{\eta}{}_{\eta}} \left(u^{x}\right)^4\\
  &\quad - 2 {\pi^{yy}} \left(u^{x}\right)^4 - 2 {\pi_{Q}^{yy}} \left(u^{x}\right)^4 - 2 {\pi_{Q}{}^{\eta}{}_{\eta}} \left(u^{x}\right)^4 + 8 {a_{\mathrm{SB}}} {T}^4 \left(u^{x}\right)^4\\
  &\quad + 3 {\pi^{xy}} {u^{x}} {u^{y}} + 3 {\pi_{Q}^{xy}} {u^{x}} {u^{y}} + 2 {\pi^{xy}} \left(u^{x}\right)^3 {u^{y}} + 2 {\pi_{Q}^{xy}} \left(u^{x}\right)^3 {u^{y}}\\
  &\quad - 2 {\pi^{\eta}{}_{\eta}} \left(u^{y}\right)^2 - {\pi^{yy}} \left(u^{y}\right)^2 - {\pi_{Q}^{yy}} \left(u^{y}\right)^2 - 2 {\pi_{Q}{}^{\eta}{}_{\eta}} \left(u^{y}\right)^2\\
  &\quad + 12 {a_{\mathrm{SB}}} {T}^4 \left(u^{y}\right)^2 - 3 {\pi^{\eta}{}_{\eta}} \left(u^{x}\right)^2 \left(u^{y}\right)^2 - 3 {\pi^{yy}} \left(u^{x}\right)^2 \left(u^{y}\right)^2\\
  &\quad - 3 {\pi_{Q}^{yy}} \left(u^{x}\right)^2 \left(u^{y}\right)^2 - 3 {\pi_{Q}{}^{\eta}{}_{\eta}} \left(u^{x}\right)^2 \left(u^{y}\right)^2 + 24 {a_{\mathrm{SB}}} {T}^4 \left(u^{x}\right)^2 \left(u^{y}\right)^2\\
  &\quad + 8 {a_{\mathrm{SB}}} {T}^4 \left(u^{x}\right)^4 \left(u^{y}\right)^2 + 3 {\pi^{xy}} {u^{x}} \left(u^{y}\right)^3 + 3 {\pi_{Q}^{xy}} {u^{x}} \left(u^{y}\right)^3 - {\pi^{\eta}{}_{\eta}} \left(u^{y}\right)^4\\
  &\quad - {\pi_{Q}{}^{\eta}{}_{\eta}} \left(u^{y}\right)^4 + 12 {a_{\mathrm{SB}}} {T}^4 \left(u^{y}\right)^4 + 12 {a_{\mathrm{SB}}} {T}^4 \left(u^{x}\right)^2 \left(u^{y}\right)^4 + 4 {a_{\mathrm{SB}}} {T}^4 \left(u^{y}\right)^6\\
  &\quad + {\Pi_{Q}} \left(\left(u^{y}\right)^2+1\right)\Bigl(2 \left(u^{x}\right)^4+3 \left(u^{x}\right)^2 \left(\left(u^{y}\right)^2+1\right)+\left(\left(u^{y}\right)^2+1\right)^2\Bigr)\\
  &\quad + {h} \left(\left(u^{y}\right)^2+1\right)\Bigl({\nu^{y}} {u^{y}} \left(\left(u^{y}\right)^2+1\right)\\
  &\quad + {\nu^{x}} {u^{x}} \left(2 \left(u^{x}\right)^2+3 \left(u^{y}\right)^2+3\right)\\
  &\quad + {n}\Bigl(2 \left(u^{x}\right)^4+3 \left(u^{x}\right)^2 \left(\left(u^{y}\right)^2+1\right)+\left(\left(u^{y}\right)^2+1\right)^2\Bigr)\Bigr)\Bigr) \\
(A_x)_{1,3} &= \frac{1}{\left(u^{\tau}\right)^3 \left(\left(u^{y}\right)^2+1\right)^2}\Bigl({\Pi_{Q}} \left(u^{\tau}\right)^2 {u^{x}} {u^{y}} \left(\left(u^{y}\right)^2+1\right)^2 + {n} \left(u^{\tau}\right)^2 {u^{x}} {u^{y}} \left(\left(u^{y}\right)^2+1\right)^2 ({h}-{T})\\
  &\quad + {n} {T} \left(u^{\tau}\right)^2 {u^{x}} {u^{y}} \left(\left(u^{y}\right)^2+1\right)^2 + 4 {a_{\mathrm{SB}}} {T}^4 \left(u^{\tau}\right)^2 {u^{x}} {u^{y}} \left(\left(u^{y}\right)^2+1\right)^2\\
  &\quad + {h} \left(\left(u^{y}\right)^2+1\right)^2 \left({\nu^{x}} {u^{y}} \left(\left(u^{y}\right)^2+1\right)+{\nu^{y}} \left(\left(u^{x}\right)^3+{u^{x}}\right)\right)\\
  &\quad - {u^{y}} \left(\left(u^{y}\right)^2+1\right)\Bigl(- {\pi^{\eta}{}_{\eta}} \left(u^{\tau}\right)^2 {u^{x}} - {\pi^{yy}} \left(\left(u^{x}\right)^3+{u^{x}}\right)\\
  &\quad + {\pi^{xy}} {u^{y}} \left(2 \left(u^{x}\right)^2+\left(u^{y}\right)^2+1\right)\Bigr)\\
  &\quad - {u^{y}} \left(\left(u^{y}\right)^2+1\right)\Bigl(- {\pi_{Q}{}^{\eta}{}_{\eta}} \left(u^{\tau}\right)^2 {u^{x}} - {\pi_{Q}^{yy}} \left(\left(u^{x}\right)^3+{u^{x}}\right)\\
  &\quad + {\pi_{Q}^{xy}} {u^{y}} \left(2 \left(u^{x}\right)^2+\left(u^{y}\right)^2+1\right)\Bigr)\\
  &\quad + \left(u^{\tau}\right)^2\Bigl(2 {u^{x}} {u^{y}} \left({\pi^{\eta}{}_{\eta}} \left(u^{x}\right)^2+{\pi^{yy}} \left(u^{x}\right)^2+{\pi^{yy}}\right)\\
  &\quad + {\pi^{xy}} \left(\left(\left(u^{y}\right)^2+1\right)^2-2 \left(u^{x}\right)^2 \left(\left(u^{y}\right)^2-1\right)\right)\Bigr)\\
  &\quad + \left(u^{\tau}\right)^2\Bigl(2 {u^{x}} {u^{y}} \left({\pi_{Q}^{yy}} \left(u^{x}\right)^2+{\pi_{Q}^{yy}}+{\pi_{Q}{}^{\eta}{}_{\eta}} \left(u^{x}\right)^2\right)\\
  &\quad + {\pi_{Q}^{xy}} \left(\left(\left(u^{y}\right)^2+1\right)^2-2 \left(u^{x}\right)^2 \left(\left(u^{y}\right)^2-1\right)\right)\Bigr)\Bigr) \\
(A_x)_{1,4} &= -\frac{{u^{x}} \left(\left(u^{x}\right)^2+1\right)}{{u^{\tau}} \left(\left(u^{y}\right)^2+1\right)} \\
(A_x)_{1,5} &= -\frac{{u^{\tau}} {u^{x}}}{\left(u^{y}\right)^2+1} \\
(A_x)_{1,6} &= \frac{{u^{y}} \left(2 \left(u^{x}\right)^2+\left(u^{y}\right)^2+1\right)}{{u^{\tau}} \left(\left(u^{y}\right)^2+1\right)} \\
(A_x)_{1,7} &= {(\partial_{\alpha} n)} {h} {u^{\tau}} {u^{x}} \\
(A_x)_{1,8} &= \frac{{h} \left(2 \left(u^{x}\right)^2+\left(u^{y}\right)^2+1\right)}{{u^{\tau}}} \\
(A_x)_{1,9} &= \frac{{h} {u^{x}} {u^{y}}}{{u^{\tau}}} \\
(A_x)_{1,10} &= -\frac{{u^{x}} \left(\left(u^{x}\right)^2+1\right)}{{u^{\tau}} \left(\left(u^{y}\right)^2+1\right)} \\
(A_x)_{1,11} &= -\frac{{u^{\tau}} {u^{x}}}{\left(u^{y}\right)^2+1} \\
(A_x)_{1,12} &= \frac{{u^{y}} \left(2 \left(u^{x}\right)^2+\left(u^{y}\right)^2+1\right)}{{u^{\tau}} \left(\left(u^{y}\right)^2+1\right)} \\
(A_x)_{1,13} &= {u^{\tau}} {u^{x}} \\
(A_x)_{2,1} &= 16 {a_{\mathrm{SB}}} {T}^3 \left(u^{x}\right)^2+4 {a_{\mathrm{SB}}} {T}^3+{h'(T)} {n} \left(u^{x}\right)^2+2 {h'(T)} {\nu^{x}} {u^{x}}+{(\partial_{T} n)} {h} \left(u^{x}\right)^2+{(\partial_{T} n)} {T}+{n} \\
(A_x)_{2,2} &= \frac{2}{\left(u^{y}\right)^2+1}\Bigl(- {\pi^{\eta}{}_{\eta}} {u^{x}} - {\pi^{yy}} {u^{x}} - {\pi_{Q}^{yy}} {u^{x}} - {\pi_{Q}{}^{\eta}{}_{\eta}} {u^{x}} + 4 {a_{\mathrm{SB}}} {T}^4 {u^{x}} + {\pi^{xy}} {u^{y}} + {\pi_{Q}^{xy}} {u^{y}}\\
  &\quad + 4 {a_{\mathrm{SB}}} {T}^4 {u^{x}} \left(u^{y}\right)^2 + {\Pi_{Q}} {u^{x}} \left(\left(u^{y}\right)^2+1\right) + {h} \left(\left(u^{y}\right)^2+1\right) ({n} {u^{x}}+{\nu^{x}})\Bigr) \\
(A_x)_{2,3} &= \frac{2}{\left(\left(u^{y}\right)^2+1\right)^2}\Bigl({u^{y}} \left(\left(u^{x}\right)^2 ({\pi^{\eta}{}_{\eta}}+{\pi_{Q}{}^{\eta}{}_{\eta}})+{\pi^{yy}} \left(\left(u^{x}\right)^2+1\right)+{\pi_{Q}^{yy}} \left(\left(u^{x}\right)^2+1\right)\right)\\
  &\quad + {\pi^{xy}} \left({u^{x}}-{u^{x}} \left(u^{y}\right)^2\right) + {\pi_{Q}^{xy}} \left({u^{x}}-{u^{x}} \left(u^{y}\right)^2\right)\Bigr) \\
(A_x)_{2,4} &= -\frac{\left(u^{x}\right)^2+1}{\left(u^{y}\right)^2+1} \\
(A_x)_{2,5} &= -\frac{\left(u^{\tau}\right)^2}{\left(u^{y}\right)^2+1} \\
(A_x)_{2,6} &= \frac{2 {u^{x}} {u^{y}}}{\left(u^{y}\right)^2+1} \\
(A_x)_{2,7} &= {(\partial_{\alpha} n)} \left({h} \left(u^{x}\right)^2+{T}\right) \\
(A_x)_{2,8} &= 2 {h} {u^{x}} \\
(A_x)_{2,10} &= -\frac{\left(u^{x}\right)^2+1}{\left(u^{y}\right)^2+1} \\
(A_x)_{2,11} &= -\frac{\left(u^{\tau}\right)^2}{\left(u^{y}\right)^2+1} \\
(A_x)_{2,12} &= \frac{2 {u^{x}} {u^{y}}}{\left(u^{y}\right)^2+1} \\
(A_x)_{2,13} &= \left(u^{x}\right)^2+1 \\
(A_x)_{3,1} &= {u^{x}} {u^{y}} \left(16 {a_{\mathrm{SB}}} {T}^3+{(\partial_{T} n)} {h}\right)+{h'(T)} {u^{y}} ({n} {u^{x}}+{\nu^{x}})+{h'(T)} {\nu^{y}} {u^{x}} \\
(A_x)_{3,2} &= {u^{y}} \left(4 {a_{\mathrm{SB}}} {T}^4+{\Pi_{Q}}\right)+{h} ({n} {u^{y}}+{\nu^{y}}) \\
(A_x)_{3,3} &= {u^{x}} \left(4 {a_{\mathrm{SB}}} {T}^4+{\Pi_{Q}}\right)+{h} ({n} {u^{x}}+{\nu^{x}}) \\
(A_x)_{3,6} &= 1 \\
(A_x)_{3,7} &= {(\partial_{\alpha} n)} {h} {u^{x}} {u^{y}} \\
(A_x)_{3,8} &= {h} {u^{y}} \\
(A_x)_{3,9} &= {h} {u^{x}} \\
(A_x)_{3,12} &= 1 \\
(A_x)_{3,13} &= {u^{x}} {u^{y}} \\
(A_x)_{4,2} &= {\delta_{\pi\pi}} {\pi^{yy}}-\frac{2 {\pi^{xy}} {\tau_{\pi}} {u^{x}} \left(\left(u^{y}\right)^2+1\right) {u^{y}}}{\left(u^{\tau}\right)^2}+\frac{2 {\pi^{yy}} {\tau_{\pi}} \left(u^{x}\right)^2 \left(u^{y}\right)^2}{\left(u^{\tau}\right)^2}-\frac{2}{3} \left(\left(u^{y}\right)^2+1\right) {\eta} \\
(A_x)_{4,3} &= \frac{2 {u^{x}} {u^{y}} \left({\pi^{xy}} {\tau_{\pi}} {u^{x}} {u^{y}}-{\pi^{yy}} {\tau_{\pi}} \left(\left(u^{x}\right)^2+1\right)+\left(u^{\tau}\right)^2 {\eta}\right)}{\left(u^{\tau}\right)^2} \\
(A_x)_{4,4} &= {\tau_{\pi}} {u^{x}} \\
(A_x)_{5,2} &= \frac{3 {\delta_{\pi\pi}} {\pi^{\eta}{}_{\eta}}-2 {\eta}}{3 {\tau}^2} \\
(A_x)_{5,5} &= \frac{{\tau_{\pi}} {u^{x}}}{{\tau}^2} \\
(A_x)_{6,2} &= {\delta_{\pi\pi}} {\pi^{xy}}\\
  &\quad + \frac{{u^{x}} \left({u^{y}} \left(3 {\pi^{\eta}{}_{\eta}} {\tau_{\pi}} \left(u^{\tau}\right)^2+3 {\pi^{yy}} \left(2 {\tau_{\pi}} \left(u^{x}\right)^2+{\tau_{\pi}}\right)+\left(u^{\tau}\right)^2 {\eta}\right)-3 {\pi^{xy}} {\tau_{\pi}} {u^{x}} \left(2 \left(u^{y}\right)^2+1\right)\right)}{3 \left(u^{\tau}\right)^2} \\
(A_x)_{6,3} &= \frac{1}{\left(u^{\tau}\right)^2 \left(\left(u^{y}\right)^2+1\right)}\Bigl(- {\pi^{yy}} {\tau_{\pi}} \left(u^{x}\right)^2 \left(\left(u^{x}\right)^2+1\right) \left(2 \left(u^{y}\right)^2+1\right)\\
  &\quad + {\pi^{xy}} {\tau_{\pi}} {u^{x}} {u^{y}} \left(\left(2 \left(u^{x}\right)^2-1\right) \left(u^{y}\right)^2-1\right)\\
  &\quad - \left(u^{\tau}\right)^2 \left({\pi^{\eta}{}_{\eta}} {\tau_{\pi}} \left(u^{x}\right)^2 \left(u^{y}\right)^2-\left(\left(u^{x}\right)^2+1\right) \left(\left(u^{y}\right)^2+1\right) {\eta}\right)\Bigr) \\
(A_x)_{6,6} &= {\tau_{\pi}} {u^{x}} \\
(A_x)_{7,1} &= {(\partial_{T} n)} {u^{x}} \\
(A_x)_{7,2} &= {n} \\
(A_x)_{7,7} &= {(\partial_{\alpha} n)} {u^{x}} \\
(A_x)_{7,8} &= 1 \\
(A_x)_{8,1} &= \frac{{D_{s}} \left({h'(T)} {\nu^{x}} {u^{x}}+{(\partial_{T} n)} {T} \left(\left(u^{x}\right)^2+1\right)+{n} \left(\left(u^{x}\right)^2+1\right)\right)}{{T}} \\
(A_x)_{8,2} &= -\frac{1}{{T} \left(u^{\tau}\right)^2 \left(\left(u^{y}\right)^2+1\right)}{D_{s}}\Bigl({\pi_{Q}^{yy}} {u^{x}} + {\pi_{Q}{}^{\eta}{}_{\eta}} {u^{x}} + {\pi_{Q}^{yy}} \left(u^{x}\right)^3 + {\pi_{Q}{}^{\eta}{}_{\eta}} \left(u^{x}\right)^3 - 2 {\pi_{Q}^{xy}} {u^{y}}\\
  &\quad - {\pi_{Q}^{xy}} \left(u^{x}\right)^2 {u^{y}} + {\pi_{Q}^{yy}} {u^{x}} \left(u^{y}\right)^2 - {\pi_{Q}^{yy}} \left(u^{x}\right)^3 \left(u^{y}\right)^2\\
  &\quad - {\pi_{Q}{}^{\eta}{}_{\eta}} \left(u^{x}\right)^3 \left(u^{y}\right)^2 - 2 {\pi_{Q}^{xy}} \left(u^{y}\right)^3 + {\pi_{Q}^{xy}} \left(u^{x}\right)^2 \left(u^{y}\right)^3\\
  &\quad - {\pi_{Q}{}^{\eta}{}_{\eta}} {u^{x}} \left(u^{y}\right)^4 - {\Pi_{Q}} \left(u^{\tau}\right)^2 {u^{x}} \left(\left(u^{y}\right)^2+1\right)\\
  &\quad + {h} \left(\left(u^{y}\right)^2+1\right)\Bigl(- {n} \left(u^{\tau}\right)^2 {u^{x}} - {\nu^{y}} \left(u^{x}\right)^3 {u^{y}} + {\nu^{x}} \left(u^{x}\right)^2 \left(\left(u^{y}\right)^2-1\right)\\
  &\quad - 2 {\nu^{x}} \left(\left(u^{y}\right)^2+1\right)\Bigr)\Bigr) \\
(A_x)_{8,3} &= \frac{1}{{T} \left(u^{\tau}\right)^2 \left(\left(u^{y}\right)^2+1\right)^2}{D_{s}}\Bigl(- {\pi_{Q}{}^{\eta}{}_{\eta}} \left(u^{\tau}\right)^2 \left(u^{x}\right)^2 {u^{y}} \left(\left(u^{y}\right)^2-1\right)\\
  &\quad - {h} \left(u^{x}\right)^2 \left(\left(u^{y}\right)^2+1\right)^2 \left(-{\nu^{x}} {u^{x}} {u^{y}}+{\nu^{y}} \left(u^{x}\right)^2+{\nu^{y}}\right)\\
  &\quad + {\pi_{Q}^{xy}} {u^{x}} \left(\left(u^{x}\right)^2 \left(\left(u^{y}\right)^2-1\right)^2-3 \left(u^{y}\right)^4-2 \left(u^{y}\right)^2+1\right)\\
  &\quad - {\pi_{Q}^{yy}} \left(\left(u^{x}\right)^2+1\right) {u^{y}} \left(\left(u^{x}\right)^2 \left(\left(u^{y}\right)^2-1\right)-2 \left(\left(u^{y}\right)^2+1\right)\right)\Bigr) \\
(A_x)_{8,7} &= {(\partial_{\alpha} n)} {D_{s}} \left(\left(u^{x}\right)^2+1\right) \\
(A_x)_{8,8} &= \frac{{D_{s}} {h} {u^{x}}}{{T}} \\
(A_x)_{8,10} &= -\frac{{D_{s}} \left(u^{x}\right)^2+{D_{s}}}{{T} \left(u^{y}\right)^2+{T}} \\
(A_x)_{8,11} &= -\frac{{D_{s}} \left(u^{\tau}\right)^2}{{T} \left(\left(u^{y}\right)^2+1\right)} \\
(A_x)_{8,12} &= \frac{2 {D_{s}} {u^{x}} {u^{y}}}{{T} \left(u^{y}\right)^2+{T}} \\
(A_x)_{8,13} &= \frac{{D_{s}} \left(\left(u^{x}\right)^2+1\right)}{{T}} \\
(A_x)_{9,1} &= \frac{{D_{s}} {u^{x}} ({h'(T)} {\nu^{y}}+{u^{y}} ({(\partial_{T} n)} {T}+{n}))}{{T}} \\
(A_x)_{9,2} &= \frac{1}{{T} \left(u^{\tau}\right)^2}{D_{s}}\Bigl({h} \left(\left(u^{y}\right)^2+1\right) \left(-{\nu^{x}} {u^{x}} {u^{y}}+{\nu^{y}} \left(u^{x}\right)^2+{\nu^{y}}\right)\\
  &\quad + {u^{y}} \left(-{\pi_{Q}^{xy}} {u^{x}} {u^{y}}+{\pi_{Q}^{yy}} \left(u^{x}\right)^2+{\pi_{Q}^{yy}}+{\pi_{Q}{}^{\eta}{}_{\eta}} \left(u^{x}\right)^2+{\pi_{Q}{}^{\eta}{}_{\eta}} \left(u^{y}\right)^2+{\pi_{Q}{}^{\eta}{}_{\eta}}\right)\Bigr) \\
(A_x)_{9,3} &= \frac{1}{{T} \left(u^{\tau}\right)^2 \left(\left(u^{y}\right)^2+1\right)}{D_{s}}\Bigl({\Pi_{Q}} \left(u^{\tau}\right)^2 {u^{x}} \left(\left(u^{y}\right)^2+1\right)\\
  &\quad + {h} \left(\left(u^{y}\right)^2+1\right)\Bigl({n} \left(u^{\tau}\right)^2 {u^{x}} - {\nu^{y}} {u^{x}} \left(\left(u^{x}\right)^2+1\right) {u^{y}}\\
  &\quad + {\nu^{x}} \left(\left(u^{x}\right)^2+1\right) \left(\left(u^{y}\right)^2+1\right)\Bigr)\\
  &\quad - {u^{y}}\Bigl({u^{x}} {u^{y}} \left({\pi_{Q}^{yy}} \left(\left(u^{x}\right)^2+1\right)+{\pi_{Q}{}^{\eta}{}_{\eta}} \left(u^{\tau}\right)^2\right)\\
  &\quad + {\pi_{Q}^{xy}} \left(-\left(\left(u^{x}\right)^2 \left(\left(u^{y}\right)^2-1\right)\right)+\left(u^{y}\right)^2+1\right)\Bigr)\Bigr) \\
(A_x)_{9,7} &= {(\partial_{\alpha} n)} {D_{s}} {u^{x}} {u^{y}} \\
(A_x)_{9,9} &= \frac{{D_{s}} {h} {u^{x}}}{{T}} \\
(A_x)_{9,12} &= \frac{{D_{s}}}{{T}} \\
(A_x)_{9,13} &= \frac{{D_{s}} {u^{x}} {u^{y}}}{{T}} \\
(A_x)_{10,1} &= \frac{{D_{s}} \left(3 {\mathcal{C}'(T)} {\pi_{Q}^{yy}} {u^{x}}-2 {\mathcal{B}'(T)} \left({\nu^{x}} \left(u^{y}\right)^2+{\nu^{x}}-3 {\nu^{y}} {u^{x}} {u^{y}}\right)\right)}{6 {T}} \\
(A_x)_{10,2} &= -\frac{1}{6 {T} \left(u^{\tau}\right)^2}{D_{s}}\Bigl(2 {I_{31}} \left(u^{\tau}\right)^2 \left(\left(u^{y}\right)^2+1\right)\\
  &\quad + 2 {u^{x}} \left({\mathcal{A}} \left(u^{y}\right)^2+{\mathcal{A}}+3 {\mathcal{B}} \left(u^{y}\right)^2\right) \left({\nu^{x}} \left(u^{y}\right)^2+{\nu^{x}}-{\nu^{y}} {u^{x}} {u^{y}}\right)\\
  &\quad + {\mathcal{C}}\Bigl(2 {\Pi_{Q}} \left(u^{\tau}\right)^2 \left(\left(u^{y}\right)^2+1\right)\\
  &\quad - 2 \left(\left(u^{y}\right)^2+1\right) \left({\pi_{Q}{}^{\eta}{}_{\eta}} \left(u^{\tau}\right)^2-4 {\pi_{Q}^{xy}} {u^{x}} {u^{y}}\right)\\
  &\quad - {\pi_{Q}^{yy}} \left(\left(u^{x}\right)^2 \left(8 \left(u^{y}\right)^2+5\right)+5 \left(\left(u^{y}\right)^2+1\right)\right)\Bigr)\Bigr) \\
(A_x)_{10,3} &= \frac{1}{3 {T} \left(u^{\tau}\right)^2}{D_{s}}\Bigl({u^{x}}\Bigl(3 {I_{31}} \left(u^{\tau}\right)^2 {u^{y}} - {\mathcal{A}} {\nu^{y}} \left(u^{x}\right)^2 \left(\left(u^{y}\right)^2-2\right) + 2 {\mathcal{A}} {\nu^{y}} \left(\left(u^{y}\right)^2+1\right)\\
  &\quad + {\mathcal{A}} {\nu^{x}} {u^{x}} {u^{y}} \left(\left(u^{y}\right)^2+1\right)\\
  &\quad + 3 {\mathcal{B}} \left({\nu^{x}} {u^{x}} \left(u^{y}\right)^3-{\nu^{y}} \left(u^{x}\right)^2 \left(\left(u^{y}\right)^2-1\right)+{\nu^{y}}\right)\Bigr)\\
  &\quad + {\mathcal{C}}\Bigl(- {u^{x}} {u^{y}} \left(-3 {\Pi_{Q}} \left(u^{\tau}\right)^2+4 {\pi_{Q}^{yy}} \left(\left(u^{x}\right)^2+1\right)+{\pi_{Q}{}^{\eta}{}_{\eta}} \left(u^{\tau}\right)^2\right)\\
  &\quad + 2 {\pi_{Q}^{xy}} \left(\left(u^{x}\right)^2 \left(2 \left(u^{y}\right)^2+1\right)+\left(u^{y}\right)^2+1\right)\Bigr)\Bigr) \\
(A_x)_{10,8} &= -\frac{{\mathcal{B}} {D_{s}} \left(\left(u^{y}\right)^2+1\right)}{3 {T}} \\
(A_x)_{10,9} &= \frac{{\mathcal{B}} {D_{s}} {u^{x}} {u^{y}}}{{T}} \\
(A_x)_{10,10} &= \frac{{\mathcal{C}} {D_{s}} {u^{x}}}{2 {T}} \\
(A_x)_{11,1} &= -\frac{{D_{s}} (2 {\mathcal{B}'(T)} {\nu^{x}}-3 {\mathcal{C}'(T)} {\pi_{Q}{}^{\eta}{}_{\eta}} {u^{x}})}{6 {T}} \\
(A_x)_{11,2} &= -\frac{1}{6 {T} \left(u^{\tau}\right)^2}{D_{s}}\Bigl(2 {I_{31}} \left(u^{\tau}\right)^2 + 2 {\mathcal{A}} {u^{x}} \left({\nu^{x}} \left(u^{y}\right)^2+{\nu^{x}}-{\nu^{y}} {u^{x}} {u^{y}}\right)\\
  &\quad + {\mathcal{C}}\Bigl(- 5 {\pi_{Q}{}^{\eta}{}_{\eta}} + 2 {\Pi_{Q}} \left(u^{\tau}\right)^2 - 5 {\pi_{Q}{}^{\eta}{}_{\eta}} \left(u^{x}\right)^2 - 2 {\pi_{Q}^{yy}} \left(\left(u^{x}\right)^2+1\right)\\
  &\quad + 2 {\pi_{Q}^{xy}} {u^{x}} {u^{y}} - 5 {\pi_{Q}{}^{\eta}{}_{\eta}} \left(u^{y}\right)^2\Bigr)\Bigr) \\
(A_x)_{11,3} &= \frac{1}{3 {T} \left(u^{\tau}\right)^2 \left(\left(u^{y}\right)^2+1\right)}{D_{s}}\Bigl(- {\mathcal{C}} {u^{x}} {u^{y}} \left({\pi_{Q}^{yy}} \left(\left(u^{x}\right)^2+1\right)+{\pi_{Q}{}^{\eta}{}_{\eta}} \left(u^{\tau}\right)^2\right)\\
  &\quad - {\mathcal{A}} {u^{x}} \left(\left(u^{y}\right)^2+1\right) \left(-{\nu^{x}} {u^{x}} {u^{y}}+{\nu^{y}} \left(u^{x}\right)^2+{\nu^{y}}\right)\\
  &\quad + {\mathcal{C}} {\pi_{Q}^{xy}} \left(\left(u^{x}\right)^2 \left(\left(u^{y}\right)^2-1\right)-\left(u^{y}\right)^2-1\right)\Bigr) \\
(A_x)_{11,8} &= -\frac{{\mathcal{B}} {D_{s}}}{3 {T}} \\
(A_x)_{11,11} &= \frac{{\mathcal{C}} {D_{s}} {u^{x}}}{2 {T}} \\
(A_x)_{12,1} &= \frac{{D_{s}} \left({\mathcal{B}'(T)} {\nu^{x}} {u^{x}} {u^{y}}+3 {\mathcal{B}'(T)} {\nu^{y}} \left(\left(u^{x}\right)^2+1\right)+3 {\mathcal{C}'(T)} {\pi_{Q}^{xy}} {u^{x}}\right)}{6 {T}} \\
(A_x)_{12,2} &= \frac{1}{6 {T} \left(u^{\tau}\right)^2}{D_{s}}\Bigl(3 {\mathcal{B}} {\nu^{y}} {u^{x}} + 3 {\mathcal{B}} {\nu^{y}} \left(u^{x}\right)^3 - 3 {\mathcal{B}} {\nu^{x}} {u^{y}} + {I_{31}} {u^{x}} {u^{y}} - 6 {\mathcal{B}} {\nu^{x}} \left(u^{x}\right)^2 {u^{y}} + {I_{31}} \left(u^{x}\right)^3 {u^{y}}\\
  &\quad + 6 {\mathcal{B}} {\nu^{y}} {u^{x}} \left(u^{y}\right)^2 + 6 {\mathcal{B}} {\nu^{y}} \left(u^{x}\right)^3 \left(u^{y}\right)^2 - 3 {\mathcal{B}} {\nu^{x}} \left(u^{y}\right)^3 + {I_{31}} {u^{x}} \left(u^{y}\right)^3\\
  &\quad - 6 {\mathcal{B}} {\nu^{x}} \left(u^{x}\right)^2 \left(u^{y}\right)^3\\
  &\quad + {\mathcal{A}} {u^{x}}\Bigl(3 {\nu^{y}} \left(\left(u^{y}\right)^2+1\right) - 2 {\nu^{x}} {u^{x}} {u^{y}} \left(\left(u^{y}\right)^2+1\right)\\
  &\quad + {\nu^{y}} \left(u^{x}\right)^2 \left(2 \left(u^{y}\right)^2+3\right)\Bigr)\\
  &\quad + {\mathcal{C}}\Bigl({u^{x}} {u^{y}} \left({\Pi_{Q}} \left(u^{\tau}\right)^2+{\pi_{Q}^{yy}} \left(8 \left(u^{x}\right)^2+5\right)+5 {\pi_{Q}{}^{\eta}{}_{\eta}} \left(u^{\tau}\right)^2\right)\\
  &\quad + {\pi_{Q}^{xy}} \left(\left(u^{x}\right)^2 \left(3-8 \left(u^{y}\right)^2\right)+6 \left(\left(u^{y}\right)^2+1\right)\right)\Bigr)\Bigr) \\
(A_x)_{12,3} &= \frac{1}{6 {T} \left(u^{\tau}\right)^2 \left(\left(u^{y}\right)^2+1\right)}{D_{s}}\Bigl(3 {I_{31}} \left(u^{\tau}\right)^2 \left(\left(u^{x}\right)^2+1\right) \left(\left(u^{y}\right)^2+1\right)\\
  &\quad + {\mathcal{C}}\Bigl(- 3 {\pi_{Q}{}^{\eta}{}_{\eta}} - 6 {\pi_{Q}{}^{\eta}{}_{\eta}} \left(u^{x}\right)^2 - 3 {\pi_{Q}{}^{\eta}{}_{\eta}} \left(u^{x}\right)^4 + {\pi_{Q}^{xy}} {u^{x}} {u^{y}}\\
  &\quad + 4 {\pi_{Q}^{xy}} \left(u^{x}\right)^3 {u^{y}} - 6 {\pi_{Q}{}^{\eta}{}_{\eta}} \left(u^{y}\right)^2 - 11 {\pi_{Q}{}^{\eta}{}_{\eta}} \left(u^{x}\right)^2 \left(u^{y}\right)^2\\
  &\quad - 5 {\pi_{Q}{}^{\eta}{}_{\eta}} \left(u^{x}\right)^4 \left(u^{y}\right)^2 + {\pi_{Q}^{xy}} {u^{x}} \left(u^{y}\right)^3 + 8 {\pi_{Q}^{xy}} \left(u^{x}\right)^3 \left(u^{y}\right)^3\\
  &\quad - 3 {\pi_{Q}{}^{\eta}{}_{\eta}} \left(u^{y}\right)^4 - 5 {\pi_{Q}{}^{\eta}{}_{\eta}} \left(u^{x}\right)^2 \left(u^{y}\right)^4\\
  &\quad + 3 {\Pi_{Q}} \left(u^{\tau}\right)^2 \left(\left(u^{x}\right)^2+1\right) \left(\left(u^{y}\right)^2+1\right)\\
  &\quad - {\pi_{Q}^{yy}} \left(\left(u^{x}\right)^2+1\right)\Bigl(\left(u^{x}\right)^2 \left(8 \left(u^{y}\right)^2+6\right)+3 \left(\left(u^{y}\right)^2+1\right)\Bigr)\Bigr)\\
  &\quad + \left(u^{y}\right)^2+1\Bigl({\mathcal{A}} {u^{x}}\Bigl(- 2 {\nu^{y}} {u^{x}} \left(\left(u^{x}\right)^2+1\right) {u^{y}} + 3 {\nu^{x}} \left(\left(u^{y}\right)^2+1\right)\\
  &\quad + {\nu^{x}} \left(u^{x}\right)^2 \left(2 \left(u^{y}\right)^2+3\right)\Bigr)\\
  &\quad + 3{\mathcal{B}} \left(\left(u^{x}\right)^2+1\right)\Bigl({\nu^{x}} \left(2 {u^{x}} \left(u^{y}\right)^2+{u^{x}}\right)-{\nu^{y}} \left(2 \left(u^{x}\right)^2+1\right) {u^{y}}\Bigr)\Bigr)\Bigr) \\
(A_x)_{12,8} &= \frac{{\mathcal{B}} {D_{s}} {u^{x}} {u^{y}}}{6 {T}} \\
(A_x)_{12,9} &= \frac{{\mathcal{B}} {D_{s}} \left(\left(u^{x}\right)^2+1\right)}{2 {T}} \\
(A_x)_{12,12} &= \frac{{\mathcal{C}} {D_{s}} {u^{x}}}{2 {T}} \\
(A_x)_{13,1} &= \frac{{D_{s}} (5 {\mathcal{B}'(T)} {\nu^{x}}+3 {u^{x}} ({\mathcal{C}'(T)} {\Pi_{Q}}+{(\partial_{T} I_{31})}))}{6 {T}} \\
(A_x)_{13,2} &= \frac{1}{6 {T} \left(u^{\tau}\right)^2}{D_{s}}\Bigl(5 {I_{31}} \left(u^{\tau}\right)^2 + 5 {\mathcal{C}} {\Pi_{Q}} \left(u^{\tau}\right)^2 + 2 {\mathcal{A}} {u^{x}} \left({\nu^{x}} \left(u^{y}\right)^2+{\nu^{x}}-{\nu^{y}} {u^{x}} {u^{y}}\right)\\
  &\quad - 2 {\mathcal{C}} \left(-{\pi_{Q}^{xy}} {u^{x}} {u^{y}}+{\pi_{Q}^{yy}} \left(u^{x}\right)^2+{\pi_{Q}^{yy}}+{\pi_{Q}{}^{\eta}{}_{\eta}} \left(u^{x}\right)^2+{\pi_{Q}{}^{\eta}{}_{\eta}} \left(u^{y}\right)^2+{\pi_{Q}{}^{\eta}{}_{\eta}}\right)\Bigr) \\
(A_x)_{13,3} &= \frac{1}{3 {T} \left(u^{\tau}\right)^2 \left(\left(u^{y}\right)^2+1\right)}{D_{s}}\Bigl({\mathcal{C}} {u^{x}} {u^{y}} \left({\pi_{Q}^{yy}} \left(\left(u^{x}\right)^2+1\right)+{\pi_{Q}{}^{\eta}{}_{\eta}} \left(u^{\tau}\right)^2\right)\\
  &\quad + {\mathcal{A}} {u^{x}} \left(\left(u^{y}\right)^2+1\right) \left(-{\nu^{x}} {u^{x}} {u^{y}}+{\nu^{y}} \left(u^{x}\right)^2+{\nu^{y}}\right)\\
  &\quad + {\mathcal{C}} {\pi_{Q}^{xy}} \left(-\left(\left(u^{x}\right)^2 \left(\left(u^{y}\right)^2-1\right)\right)+\left(u^{y}\right)^2+1\right)\Bigr) \\
(A_x)_{13,7} &= \frac{{(\partial_{\alpha} I_{31})} {D_{s}} {u^{x}}}{2 {T}} \\
(A_x)_{13,8} &= \frac{5 {\mathcal{B}} {D_{s}}}{6 {T}} \\
(A_x)_{13,13} &= \frac{{\mathcal{C}} {D_{s}} {u^{x}}}{2 {T}} \\
(A_y)_{1,1} &= \frac{\left(u^{\tau}\right)^2 {u^{y}} \left(16 {a_{\mathrm{SB}}} {T}^3+{(\partial_{T} n)} {h}\right)+{h'(T)} \left({n} \left(u^{\tau}\right)^2 {u^{y}}+{\nu^{x}} {u^{x}} {u^{y}}+{\nu^{y}} \left(\left(u^{x}\right)^2+2 \left(u^{y}\right)^2+1\right)\right)}{{u^{\tau}}} \\
(A_y)_{1,2} &= \frac{1}{\left(u^{\tau}\right)^3}\Bigl({\pi^{xy}} + {\pi_{Q}^{xy}} + {h} {\nu^{y}} {u^{x}} + {h} {\nu^{y}} \left(u^{x}\right)^3 + {h} {\nu^{x}} {u^{y}} + {h} {n} {u^{x}} {u^{y}} + {\Pi_{Q}} {u^{x}} {u^{y}} - {\pi^{yy}} {u^{x}} {u^{y}}\\
  &\quad - {\pi_{Q}^{yy}} {u^{x}} {u^{y}} + 4 {a_{\mathrm{SB}}} {T}^4 {u^{x}} {u^{y}} + {h} {n} \left(u^{x}\right)^3 {u^{y}} + {\Pi_{Q}} \left(u^{x}\right)^3 {u^{y}} + 4 {a_{\mathrm{SB}}} {T}^4 \left(u^{x}\right)^3 {u^{y}}\\
  &\quad + {\pi^{xy}} \left(u^{y}\right)^2 + {\pi_{Q}^{xy}} \left(u^{y}\right)^2 + {h} {\nu^{x}} \left(u^{y}\right)^3 + {h} {n} {u^{x}} \left(u^{y}\right)^3 + {\Pi_{Q}} {u^{x}} \left(u^{y}\right)^3\\
  &\quad + 4 {a_{\mathrm{SB}}} {T}^4 {u^{x}} \left(u^{y}\right)^3\Bigr) \\
(A_y)_{1,3} &= \frac{1}{\left(u^{\tau}\right)^3}\Bigl({\pi^{yy}} + {\pi_{Q}^{yy}} + 4 {a_{\mathrm{SB}}} {T}^4 + {\pi^{yy}} \left(u^{x}\right)^2 + {\pi_{Q}^{yy}} \left(u^{x}\right)^2 + 8 {a_{\mathrm{SB}}} {T}^4 \left(u^{x}\right)^2 + 4 {a_{\mathrm{SB}}} {T}^4 \left(u^{x}\right)^4\\
  &\quad - {\pi^{xy}} {u^{x}} {u^{y}} - {\pi_{Q}^{xy}} {u^{x}} {u^{y}} + 12 {a_{\mathrm{SB}}} {T}^4 \left(u^{y}\right)^2 + 12 {a_{\mathrm{SB}}} {T}^4 \left(u^{x}\right)^2 \left(u^{y}\right)^2 + 8 {a_{\mathrm{SB}}} {T}^4 \left(u^{y}\right)^4\\
  &\quad + {\Pi_{Q}} \left(\left(u^{x}\right)^4+\left(u^{x}\right)^2 \left(3 \left(u^{y}\right)^2+2\right)+2 \left(u^{y}\right)^4+3 \left(u^{y}\right)^2+1\right)\\
  &\quad + {h}\Bigl({\nu^{x}} \left(\left(u^{x}\right)^3+{u^{x}}\right) + {\nu^{y}} {u^{y}} \left(3 \left(u^{x}\right)^2+2 \left(u^{y}\right)^2+3\right)\\
  &\quad + {n} \left(\left(u^{x}\right)^4+\left(u^{x}\right)^2 \left(3 \left(u^{y}\right)^2+2\right)+2 \left(u^{y}\right)^4+3 \left(u^{y}\right)^2+1\right)\Bigr)\Bigr) \\
(A_y)_{1,4} &= \frac{{u^{y}}}{{u^{\tau}}} \\
(A_y)_{1,6} &= \frac{{u^{x}}}{{u^{\tau}}} \\
(A_y)_{1,7} &= {(\partial_{\alpha} n)} {h} {u^{\tau}} {u^{y}} \\
(A_y)_{1,8} &= \frac{{h} {u^{x}} {u^{y}}}{{u^{\tau}}} \\
(A_y)_{1,9} &= \frac{{h} \left(\left(u^{x}\right)^2+2 \left(u^{y}\right)^2+1\right)}{{u^{\tau}}} \\
(A_y)_{1,10} &= \frac{{u^{y}}}{{u^{\tau}}} \\
(A_y)_{1,12} &= \frac{{u^{x}}}{{u^{\tau}}} \\
(A_y)_{1,13} &= {u^{\tau}} {u^{y}} \\
(A_y)_{2,1} &= {u^{x}} {u^{y}} \left(16 {a_{\mathrm{SB}}} {T}^3+{(\partial_{T} n)} {h}\right)+{h'(T)} {u^{y}} ({n} {u^{x}}+{\nu^{x}})+{h'(T)} {\nu^{y}} {u^{x}} \\
(A_y)_{2,2} &= {u^{y}} \left(4 {a_{\mathrm{SB}}} {T}^4+{\Pi_{Q}}\right)+{h} ({n} {u^{y}}+{\nu^{y}}) \\
(A_y)_{2,3} &= {u^{x}} \left(4 {a_{\mathrm{SB}}} {T}^4+{\Pi_{Q}}\right)+{h} ({n} {u^{x}}+{\nu^{x}}) \\
(A_y)_{2,6} &= 1 \\
(A_y)_{2,7} &= {(\partial_{\alpha} n)} {h} {u^{x}} {u^{y}} \\
(A_y)_{2,8} &= {h} {u^{y}} \\
(A_y)_{2,9} &= {h} {u^{x}} \\
(A_y)_{2,12} &= 1 \\
(A_y)_{2,13} &= {u^{x}} {u^{y}} \\
(A_y)_{3,1} &= 16 {a_{\mathrm{SB}}} {T}^3 \left(u^{y}\right)^2+4 {a_{\mathrm{SB}}} {T}^3+{h'(T)} {n} \left(u^{y}\right)^2+2 {h'(T)} {\nu^{y}} {u^{y}}+{(\partial_{T} n)} {h} \left(u^{y}\right)^2+{(\partial_{T} n)} {T}+{n} \\
(A_y)_{3,3} &= 2 \left({u^{y}} \left(4 {a_{\mathrm{SB}}} {T}^4+{\Pi_{Q}}\right)+{h} ({n} {u^{y}}+{\nu^{y}})\right) \\
(A_y)_{3,4} &= 1 \\
(A_y)_{3,7} &= {(\partial_{\alpha} n)} \left({h} \left(u^{y}\right)^2+{T}\right) \\
(A_y)_{3,9} &= 2 {h} {u^{y}} \\
(A_y)_{3,10} &= 1 \\
(A_y)_{3,13} &= \left(u^{y}\right)^2+1 \\
(A_y)_{4,2} &= -\frac{2 {\tau_{\pi}} \left(u^{y}\right)^2 \left({\pi^{xy}} \left(u^{y}\right)^2+{\pi^{xy}}-{\pi^{yy}} {u^{x}} {u^{y}}\right)}{\left(u^{\tau}\right)^2} \\
(A_y)_{4,3} &= \frac{1}{3 \left(u^{\tau}\right)^2}\Bigl(3 {\delta_{\pi\pi}} {\pi^{yy}} \left(u^{\tau}\right)^2 - 6 {\pi^{yy}} {\tau_{\pi}} \left(\left(u^{x}\right)^2+1\right) \left(u^{y}\right)^2 + 6 {\pi^{xy}} {\tau_{\pi}} {u^{x}} \left(u^{y}\right)^3 + 4 {\eta}\\
  &\quad + 4 \left(u^{x}\right)^2 {\eta} + 8 \left(u^{y}\right)^2 {\eta} + 4 \left(u^{x}\right)^2 \left(u^{y}\right)^2 {\eta} + 4 \left(u^{y}\right)^4 {\eta}\Bigr) \\
(A_y)_{4,4} &= {\tau_{\pi}} {u^{y}} \\
(A_y)_{5,3} &= \frac{3 {\delta_{\pi\pi}} {\pi^{\eta}{}_{\eta}}-2 {\eta}}{3 {\tau}^2} \\
(A_y)_{5,5} &= \frac{{\tau_{\pi}} {u^{y}}}{{\tau}^2} \\
(A_y)_{6,2} &= \frac{1}{\left(u^{\tau}\right)^2}\Bigl({\pi^{yy}} {\tau_{\pi}} \left(u^{y}\right)^2 + {\pi^{\eta}{}_{\eta}} {\tau_{\pi}} \left(u^{\tau}\right)^2 \left(u^{y}\right)^2 + 2 {\pi^{yy}} {\tau_{\pi}} \left(u^{x}\right)^2 \left(u^{y}\right)^2\\
  &\quad - {\pi^{xy}} {\tau_{\pi}} {u^{x}} {u^{y}} \left(2 \left(u^{y}\right)^2+1\right) + {\eta} + \left(u^{x}\right)^2 {\eta} + 2 \left(u^{y}\right)^2 {\eta} + \left(u^{x}\right)^2 \left(u^{y}\right)^2 {\eta} + \left(u^{y}\right)^4 {\eta}\Bigr) \\
(A_y)_{6,3} &= {\delta_{\pi\pi}} {\pi^{xy}}\\
  &\quad - \frac{1}{3 \left(u^{\tau}\right)^2 \left(\left(u^{y}\right)^2+1\right)}{u^{y}}\Bigl(3 {\pi^{yy}} {\tau_{\pi}} {u^{x}} \left(\left(u^{x}\right)^2+1\right) \left(2 \left(u^{y}\right)^2+1\right)\\
  &\quad + 3 {\pi^{xy}} {\tau_{\pi}} {u^{y}} \left(\left(1-2 \left(u^{x}\right)^2\right) \left(u^{y}\right)^2+1\right)\\
  &\quad + \left(u^{\tau}\right)^2 {u^{x}} \left(3 {\pi^{\eta}{}_{\eta}} {\tau_{\pi}} \left(u^{y}\right)^2-\left(\left(u^{y}\right)^2+1\right) {\eta}\right)\Bigr) \\
(A_y)_{6,6} &= {\tau_{\pi}} {u^{y}} \\
(A_y)_{7,1} &= {(\partial_{T} n)} {u^{y}} \\
(A_y)_{7,3} &= {n} \\
(A_y)_{7,7} &= {(\partial_{\alpha} n)} {u^{y}} \\
(A_y)_{7,9} &= 1 \\
(A_y)_{8,1} &= \frac{{D_{s}} {u^{y}} ({h'(T)} {\nu^{x}}+{u^{x}} ({(\partial_{T} n)} {T}+{n}))}{{T}} \\
(A_y)_{8,2} &= \frac{1}{{T} \left(u^{\tau}\right)^2}{D_{s}}\Bigl({\Pi_{Q}} \left(u^{\tau}\right)^2 {u^{y}} + {\pi_{Q}^{yy}} \left(u^{x}\right)^2 {u^{y}} - {\pi_{Q}^{xy}} {u^{x}} \left(\left(u^{y}\right)^2+1\right)\\
  &\quad + {h} \left({n} \left(u^{\tau}\right)^2 {u^{y}}-{\nu^{x}} {u^{x}} {u^{y}} \left(\left(u^{y}\right)^2+1\right)+{\nu^{y}} \left(\left(u^{x}\right)^2+1\right) \left(\left(u^{y}\right)^2+1\right)\right)\Bigr) \\
(A_y)_{8,3} &= \frac{{D_{s}} \left({h} \left(\left(u^{x}\right)^2+1\right) \left({\nu^{x}} \left(u^{y}\right)^2+{\nu^{x}}-{\nu^{y}} {u^{x}} {u^{y}}\right)-{u^{x}} \left(-{\pi_{Q}^{xy}} {u^{x}} {u^{y}}+{\pi_{Q}^{yy}} \left(u^{x}\right)^2+{\pi_{Q}^{yy}}\right)\right)}{{T} \left(u^{\tau}\right)^2} \\
(A_y)_{8,7} &= {(\partial_{\alpha} n)} {D_{s}} {u^{x}} {u^{y}} \\
(A_y)_{8,8} &= \frac{{D_{s}} {h} {u^{y}}}{{T}} \\
(A_y)_{8,12} &= \frac{{D_{s}}}{{T}} \\
(A_y)_{8,13} &= \frac{{D_{s}} {u^{x}} {u^{y}}}{{T}} \\
(A_y)_{9,1} &= \frac{{D_{s}} \left({h'(T)} {\nu^{y}} {u^{y}}+{(\partial_{T} n)} {T} \left(\left(u^{y}\right)^2+1\right)+{n} \left(\left(u^{y}\right)^2+1\right)\right)}{{T}} \\
(A_y)_{9,2} &= -\frac{{D_{s}} {u^{y}} \left({h} {u^{y}} \left({\nu^{x}} \left(u^{y}\right)^2+{\nu^{x}}-{\nu^{y}} {u^{x}} {u^{y}}\right)+{\pi_{Q}^{xy}} \left(\left(u^{y}\right)^2+1\right)-{\pi_{Q}^{yy}} {u^{x}} {u^{y}}\right)}{{T} \left(u^{\tau}\right)^2} \\
(A_y)_{9,3} &= \frac{1}{{T} \left(u^{\tau}\right)^2}{D_{s}}\Bigl({u^{y}} \left({\Pi_{Q}} \left(u^{\tau}\right)^2+{\pi_{Q}^{xy}} {u^{x}} {u^{y}}-{\pi_{Q}^{yy}} \left(\left(u^{x}\right)^2+1\right)\right)\\
  &\quad + {h} \left({n} \left(u^{\tau}\right)^2 {u^{y}}+{\nu^{x}} {u^{x}} \left(u^{y}\right)^3+{\nu^{y}} \left(-\left(\left(u^{x}\right)^2 \left(\left(u^{y}\right)^2-2\right)\right)+\left(u^{y}\right)^2+2\right)\right)\Bigr) \\
(A_y)_{9,7} &= {(\partial_{\alpha} n)} {D_{s}} \left(\left(u^{y}\right)^2+1\right) \\
(A_y)_{9,9} &= \frac{{D_{s}} {h} {u^{y}}}{{T}} \\
(A_y)_{9,10} &= \frac{{D_{s}}}{{T}} \\
(A_y)_{9,13} &= \frac{{D_{s}} \left(\left(u^{y}\right)^2+1\right)}{{T}} \\
(A_y)_{10,1} &= \frac{{D_{s}} \left(4 {\mathcal{B}'(T)} {\nu^{y}} \left(\left(u^{y}\right)^2+1\right)+3 {\mathcal{C}'(T)} {\pi_{Q}^{yy}} {u^{y}}\right)}{6 {T}} \\
(A_y)_{10,2} &= -\frac{1}{3 {T} \left(u^{\tau}\right)^2}{D_{s}}\Bigl({u^{y}} \left(\left(u^{y}\right)^2+1\right) ({\mathcal{A}}+3 {\mathcal{B}}) \left({\nu^{x}} \left(u^{y}\right)^2+{\nu^{x}}-{\nu^{y}} {u^{x}} {u^{y}}\right)\\
  &\quad + {\mathcal{C}} \left(4 \left(u^{y}\right)^2+1\right) \left({\pi_{Q}^{xy}} \left(u^{y}\right)^2+{\pi_{Q}^{xy}}-{\pi_{Q}^{yy}} {u^{x}} {u^{y}}\right)\Bigr) \\
(A_y)_{10,3} &= \frac{1}{6 {T} \left(u^{\tau}\right)^2}{D_{s}}\Bigl(4 {I_{31}} \left(u^{\tau}\right)^2 \left(\left(u^{y}\right)^2+1\right)\\
  &\quad + {\mathcal{C}}\Bigl(4 {\Pi_{Q}} \left(u^{\tau}\right)^2 \left(\left(u^{y}\right)^2+1\right) + 2 {\pi_{Q}^{xy}} {u^{x}} {u^{y}} \left(4 \left(u^{y}\right)^2+1\right)\\
  &\quad + {\pi_{Q}^{yy}} \left(\left(u^{x}\right)^2 \left(7-8 \left(u^{y}\right)^2\right)+\left(u^{y}\right)^2+7\right)\Bigr)\\
  &\quad + 2{u^{y}}\Bigl(3 {\mathcal{B}} {u^{x}} {u^{y}} \left({\nu^{x}} \left(u^{y}\right)^2+{\nu^{x}}-{\nu^{y}} {u^{x}} {u^{y}}\right)\\
  &\quad + {\mathcal{A}}\Bigl(- {\nu^{y}} \left(u^{x}\right)^2 \left(\left(u^{y}\right)^2-2\right) + 2 {\nu^{y}} \left(\left(u^{y}\right)^2+1\right)\\
  &\quad + {\nu^{x}} {u^{x}} {u^{y}} \left(\left(u^{y}\right)^2+1\right)\Bigr)\Bigr)\Bigr) \\
(A_y)_{10,9} &= \frac{2 {\mathcal{B}} {D_{s}} \left(\left(u^{y}\right)^2+1\right)}{3 {T}} \\
(A_y)_{10,10} &= \frac{{\mathcal{C}} {D_{s}} {u^{y}}}{2 {T}} \\
(A_y)_{11,1} &= -\frac{{D_{s}} (2 {\mathcal{B}'(T)} {\nu^{y}}-3 {\mathcal{C}'(T)} {\pi_{Q}{}^{\eta}{}_{\eta}} {u^{y}})}{6 {T}} \\
(A_y)_{11,2} &= -\frac{{D_{s}} \left({\mathcal{A}} {u^{y}} \left({\nu^{x}} \left(u^{y}\right)^2+{\nu^{x}}-{\nu^{y}} {u^{x}} {u^{y}}\right)+{\mathcal{C}} \left({\pi_{Q}^{xy}} \left(u^{y}\right)^2+{\pi_{Q}^{xy}}-{\pi_{Q}^{yy}} {u^{x}} {u^{y}}\right)\right)}{3 {T} \left(u^{\tau}\right)^2} \\
(A_y)_{11,3} &= -\frac{1}{6 {T} \left(u^{\tau}\right)^2}{D_{s}}\Bigl(2 {I_{31}} \left(u^{\tau}\right)^2 + 2 {\mathcal{A}} {u^{y}} \left(-{\nu^{x}} {u^{x}} {u^{y}}+{\nu^{y}} \left(u^{x}\right)^2+{\nu^{y}}\right)\\
  &\quad + {\mathcal{C}}\Bigl(- 3 {\pi_{Q}{}^{\eta}{}_{\eta}} + 2 {\Pi_{Q}} \left(u^{\tau}\right)^2 - 3 {\pi_{Q}{}^{\eta}{}_{\eta}} \left(u^{x}\right)^2 + 2 {\pi_{Q}^{yy}} \left(\left(u^{x}\right)^2+1\right)\\
  &\quad - 2 {\pi_{Q}^{xy}} {u^{x}} {u^{y}} - 3 {\pi_{Q}{}^{\eta}{}_{\eta}} \left(u^{y}\right)^2\Bigr)\Bigr) \\
(A_y)_{11,9} &= -\frac{{\mathcal{B}} {D_{s}}}{3 {T}} \\
(A_y)_{11,11} &= \frac{{\mathcal{C}} {D_{s}} {u^{y}}}{2 {T}} \\
(A_y)_{12,1} &= \frac{{D_{s}} \left(3 {\mathcal{B}'(T)} {\nu^{x}} \left(\left(u^{y}\right)^2+1\right)+{\mathcal{B}'(T)} {\nu^{y}} {u^{x}} {u^{y}}+3 {\mathcal{C}'(T)} {\pi_{Q}^{xy}} {u^{y}}\right)}{6 {T}} \\
(A_y)_{12,2} &= \frac{1}{6 {T} \left(u^{\tau}\right)^2}{D_{s}}\Bigl(- 3 {\mathcal{B}} {\nu^{x}} {u^{x}} + 3 {\mathcal{A}} {\nu^{y}} {u^{y}} + 3 {\mathcal{B}} {\nu^{y}} {u^{y}} + 3 {\mathcal{A}} {\nu^{y}} \left(u^{x}\right)^2 {u^{y}} + 6 {\mathcal{B}} {\nu^{y}} \left(u^{x}\right)^2 {u^{y}}\\
  &\quad - 2 {\mathcal{A}} {\nu^{x}} {u^{x}} \left(u^{y}\right)^2 - 9 {\mathcal{B}} {\nu^{x}} {u^{x}} \left(u^{y}\right)^2 + 3 {\mathcal{A}} {\nu^{y}} \left(u^{y}\right)^3 + 3 {\mathcal{B}} {\nu^{y}} \left(u^{y}\right)^3\\
  &\quad + 2 {\mathcal{A}} {\nu^{y}} \left(u^{x}\right)^2 \left(u^{y}\right)^3 + 6 {\mathcal{B}} {\nu^{y}} \left(u^{x}\right)^2 \left(u^{y}\right)^3 - 2 {\mathcal{A}} {\nu^{x}} {u^{x}} \left(u^{y}\right)^4 - 6 {\mathcal{B}} {\nu^{x}} {u^{x}} \left(u^{y}\right)^4\\
  &\quad + 3 {I_{31}} \left(u^{\tau}\right)^2 \left(\left(u^{y}\right)^2+1\right)\\
  &\quad + {\mathcal{C}}\Bigl(3 {\Pi_{Q}} \left(u^{\tau}\right)^2 \left(\left(u^{y}\right)^2+1\right)\\
  &\quad + {\pi_{Q}^{yy}} \left(\left(u^{x}\right)^2 \left(8 \left(u^{y}\right)^2+3\right)+6 \left(u^{y}\right)^2+3\right)\\
  &\quad + {u^{y}} \left(3 {\pi_{Q}{}^{\eta}{}_{\eta}} \left(u^{\tau}\right)^2 {u^{y}}-{\pi_{Q}^{xy}} {u^{x}} \left(8 \left(u^{y}\right)^2+5\right)\right)\Bigr)\Bigr) \\
(A_y)_{12,3} &= \frac{1}{6 {T} \left(u^{\tau}\right)^2 \left(\left(u^{y}\right)^2+1\right)}{D_{s}}\Bigl(\left(u^{y}\right)^2+1\Bigl({I_{31}} \left(u^{\tau}\right)^2 {u^{x}} {u^{y}}\\
  &\quad + 3 {\mathcal{B}} {\nu^{x}} \left(2 \left(u^{x}\right)^2+1\right) {u^{y}} \left(\left(u^{y}\right)^2+1\right)\\
  &\quad - 3 {\mathcal{B}} {\nu^{y}} {u^{x}} \left(\left(u^{x}\right)^2+1\right) \left(2 \left(u^{y}\right)^2+1\right)\\
  &\quad + {\mathcal{A}} {u^{y}}\Bigl(- 2 {\nu^{y}} {u^{x}} \left(\left(u^{x}\right)^2+1\right) {u^{y}} + 3 {\nu^{x}} \left(\left(u^{y}\right)^2+1\right)\\
  &\quad + {\nu^{x}} \left(u^{x}\right)^2 \left(2 \left(u^{y}\right)^2+3\right)\Bigr)\Bigr)\\
  &\quad + {\mathcal{C}}\Bigl({u^{x}} {u^{y}}\Bigl(- 3 {\pi_{Q}{}^{\eta}{}_{\eta}} \left(u^{\tau}\right)^2 \left(u^{y}\right)^2 + {\Pi_{Q}} \left(u^{\tau}\right)^2 \left(\left(u^{y}\right)^2+1\right)\\
  &\quad - {\pi_{Q}^{yy}} \left(\left(u^{x}\right)^2+1\right) \left(8 \left(u^{y}\right)^2+5\right)\Bigr)\\
  &\quad + {\pi_{Q}^{xy}}\Bigl(3 \left(\left(u^{y}\right)^4+3 \left(u^{y}\right)^2+2\right)\\
  &\quad + \left(u^{x}\right)^2 \left(8 \left(u^{y}\right)^4+8 \left(u^{y}\right)^2+6\right)\Bigr)\Bigr)\Bigr) \\
(A_y)_{12,8} &= \frac{{\mathcal{B}} {D_{s}} \left(\left(u^{y}\right)^2+1\right)}{2 {T}} \\
(A_y)_{12,9} &= \frac{{\mathcal{B}} {D_{s}} {u^{x}} {u^{y}}}{6 {T}} \\
(A_y)_{12,12} &= \frac{{\mathcal{C}} {D_{s}} {u^{y}}}{2 {T}} \\
(A_y)_{13,1} &= \frac{{D_{s}} (5 {\mathcal{B}'(T)} {\nu^{y}}+3 {u^{y}} ({\mathcal{C}'(T)} {\Pi_{Q}}+{(\partial_{T} I_{31})}))}{6 {T}} \\
(A_y)_{13,2} &= \frac{{D_{s}} \left({\mathcal{A}} {u^{y}} \left({\nu^{x}} \left(u^{y}\right)^2+{\nu^{x}}-{\nu^{y}} {u^{x}} {u^{y}}\right)+{\mathcal{C}} \left({\pi_{Q}^{xy}} \left(u^{y}\right)^2+{\pi_{Q}^{xy}}-{\pi_{Q}^{yy}} {u^{x}} {u^{y}}\right)\right)}{3 {T} \left(u^{\tau}\right)^2} \\
(A_y)_{13,3} &= \frac{1}{6 {T} \left(u^{\tau}\right)^2}{D_{s}}\Bigl(5 {I_{31}} \left(u^{\tau}\right)^2 + 5 {\mathcal{C}} {\Pi_{Q}} \left(u^{\tau}\right)^2 + 2 {\mathcal{A}} {u^{y}} \left(-{\nu^{x}} {u^{x}} {u^{y}}+{\nu^{y}} \left(u^{x}\right)^2+{\nu^{y}}\right)\\
  &\quad + 2 {\mathcal{C}} \left(-{\pi_{Q}^{xy}} {u^{x}} {u^{y}}+{\pi_{Q}^{yy}} \left(u^{x}\right)^2+{\pi_{Q}^{yy}}\right)\Bigr) \\
(A_y)_{13,7} &= \frac{{(\partial_{\alpha} I_{31})} {D_{s}} {u^{y}}}{2 {T}} \\
(A_y)_{13,9} &= \frac{5 {\mathcal{B}} {D_{s}}}{6 {T}} \\
(A_y)_{13,13} &= \frac{{\mathcal{C}} {D_{s}} {u^{y}}}{2 {T}} \\
S_{1} &= \frac{1}{{\tau} \left(\left(u^{y}\right)^2+1\right)}\Bigl({\pi^{\eta}{}_{\eta}} + {\pi_{Q}{}^{\eta}{}_{\eta}} + 4 {a_{\mathrm{SB}}} {T}^4 - {\pi^{\eta}{}_{\eta}} \left(u^{x}\right)^2 - {\pi^{yy}} \left(u^{x}\right)^2 - {\pi_{Q}^{yy}} \left(u^{x}\right)^2\\
  &\quad - {\pi_{Q}{}^{\eta}{}_{\eta}} \left(u^{x}\right)^2 + 4 {a_{\mathrm{SB}}} {T}^4 \left(u^{x}\right)^2 + 2 {\pi^{xy}} {u^{x}} {u^{y}} + 2 {\pi_{Q}^{xy}} {u^{x}} {u^{y}} + {\pi^{\eta}{}_{\eta}} \left(u^{y}\right)^2\\
  &\quad + {\pi^{yy}} \left(u^{y}\right)^2 + {\pi_{Q}^{yy}} \left(u^{y}\right)^2 + {\pi_{Q}{}^{\eta}{}_{\eta}} \left(u^{y}\right)^2 + 8 {a_{\mathrm{SB}}} {T}^4 \left(u^{y}\right)^2\\
  &\quad + 4 {a_{\mathrm{SB}}} {T}^4 \left(u^{x}\right)^2 \left(u^{y}\right)^2 + 4 {a_{\mathrm{SB}}} {T}^4 \left(u^{y}\right)^4 + {\Pi_{Q}} \left(u^{\tau}\right)^2 \left(\left(u^{y}\right)^2+1\right)\\
  &\quad + {h} \left(\left(u^{y}\right)^2+1\right) \left({n} \left(u^{\tau}\right)^2+2 ({\nu^{x}} {u^{x}}+{\nu^{y}} {u^{y}})\right)\Bigr) \\
S_{2} &= \frac{1}{{\tau} {u^{\tau}} \left(\left(u^{y}\right)^2+1\right)}\Bigl(- {\pi^{\eta}{}_{\eta}} {u^{x}} - {\pi^{yy}} {u^{x}} - {\pi_{Q}^{yy}} {u^{x}} - {\pi_{Q}{}^{\eta}{}_{\eta}} {u^{x}} + 4 {a_{\mathrm{SB}}} {T}^4 {u^{x}} - {\pi^{\eta}{}_{\eta}} \left(u^{x}\right)^3\\
  &\quad - {\pi^{yy}} \left(u^{x}\right)^3 - {\pi_{Q}^{yy}} \left(u^{x}\right)^3 - {\pi_{Q}{}^{\eta}{}_{\eta}} \left(u^{x}\right)^3 + 4 {a_{\mathrm{SB}}} {T}^4 \left(u^{x}\right)^3 + {\pi^{xy}} {u^{y}} + {\pi_{Q}^{xy}} {u^{y}}\\
  &\quad + 2 {\pi^{xy}} \left(u^{x}\right)^2 {u^{y}} + 2 {\pi_{Q}^{xy}} \left(u^{x}\right)^2 {u^{y}} - {\pi^{\eta}{}_{\eta}} {u^{x}} \left(u^{y}\right)^2 - {\pi_{Q}{}^{\eta}{}_{\eta}} {u^{x}} \left(u^{y}\right)^2\\
  &\quad + 8 {a_{\mathrm{SB}}} {T}^4 {u^{x}} \left(u^{y}\right)^2 + 4 {a_{\mathrm{SB}}} {T}^4 \left(u^{x}\right)^3 \left(u^{y}\right)^2 + {\pi^{xy}} \left(u^{y}\right)^3 + {\pi_{Q}^{xy}} \left(u^{y}\right)^3\\
  &\quad + 4 {a_{\mathrm{SB}}} {T}^4 {u^{x}} \left(u^{y}\right)^4 + {\Pi_{Q}} \left(u^{\tau}\right)^2 {u^{x}} \left(\left(u^{y}\right)^2+1\right)\\
  &\quad + {h} \left(\left(u^{y}\right)^2+1\right) \left({n} \left(u^{\tau}\right)^2 {u^{x}}+{\nu^{x}} \left(2 \left(u^{x}\right)^2+\left(u^{y}\right)^2+1\right)+{\nu^{y}} {u^{x}} {u^{y}}\right)\Bigr) \\
S_{3} &= \frac{1}{{\tau} {u^{\tau}}}\Bigl({\pi^{xy}} {u^{x}} + {\pi_{Q}^{xy}} {u^{x}} + {\Pi_{Q}} {u^{y}} + {\pi^{yy}} {u^{y}} + {\pi_{Q}^{yy}} {u^{y}} + 4 {a_{\mathrm{SB}}} {T}^4 {u^{y}} + {\Pi_{Q}} \left(u^{x}\right)^2 {u^{y}}\\
  &\quad + 4 {a_{\mathrm{SB}}} {T}^4 \left(u^{x}\right)^2 {u^{y}} + {\Pi_{Q}} \left(u^{y}\right)^3 + 4 {a_{\mathrm{SB}}} {T}^4 \left(u^{y}\right)^3\\
  &\quad + {h} \left({n} \left(u^{\tau}\right)^2 {u^{y}}+{\nu^{x}} {u^{x}} {u^{y}}+{\nu^{y}} \left(\left(u^{x}\right)^2+2 \left(u^{y}\right)^2+1\right)\right)\Bigr) \\
S_{4} &= \frac{{\delta_{\pi\pi}} {\pi^{yy}} {u^{\tau}}}{{\tau}}+{\pi^{yy}}-\frac{2 {u^{\tau}} \left(\left(u^{y}\right)^2+1\right) {\eta}}{3 {\tau}} \\
S_{5} &= \frac{3 {\pi^{\eta}{}_{\eta}} ({\delta_{\pi\pi}} {u^{\tau}}+{\tau})+4 {u^{\tau}} {\eta}}{3 {\tau}^3} \\
S_{6} &= \frac{{\delta_{\pi\pi}} {\pi^{xy}} {u^{\tau}}}{{\tau}}+{\pi^{xy}}-\frac{2 {u^{\tau}} {u^{x}} {u^{y}} {\eta}}{3 {\tau}} \\
S_{7} &= \frac{{n} \left(u^{\tau}\right)^2+{\nu^{x}} {u^{x}}+{\nu^{y}} {u^{y}}}{{\tau} {u^{\tau}}} \\
S_{8} &= {\nu^{x}}\\
  &\quad + \frac{1}{{T} {\tau} {u^{\tau}} \left(\left(u^{y}\right)^2+1\right)}{D_{s}}\Bigl(- 2 {\pi_{Q}{}^{\eta}{}_{\eta}} {u^{x}} - 2 {\pi_{Q}{}^{\eta}{}_{\eta}} \left(u^{x}\right)^3 - {\pi_{Q}^{yy}} \left(\left(u^{x}\right)^3+{u^{x}}\right) + {\pi_{Q}^{xy}} {u^{y}}\\
  &\quad + 2 {\pi_{Q}^{xy}} \left(u^{x}\right)^2 {u^{y}} - 3 {\pi_{Q}{}^{\eta}{}_{\eta}} {u^{x}} \left(u^{y}\right)^2 - {\pi_{Q}{}^{\eta}{}_{\eta}} \left(u^{x}\right)^3 \left(u^{y}\right)^2 + {\pi_{Q}^{xy}} \left(u^{y}\right)^3\\
  &\quad - {\pi_{Q}{}^{\eta}{}_{\eta}} {u^{x}} \left(u^{y}\right)^4 + {h} {\nu^{x}} \left(u^{\tau}\right)^2 \left(\left(u^{y}\right)^2+1\right)\Bigr) \\
S_{9} &= \frac{{D_{s}} \left({h} {\nu^{y}} \left(u^{\tau}\right)^2+{\pi_{Q}^{xy}} {u^{x}}+{u^{y}} \left({\pi_{Q}^{yy}}-{\pi_{Q}{}^{\eta}{}_{\eta}} \left(u^{\tau}\right)^2\right)\right)}{{T} {\tau} {u^{\tau}}}+{\nu^{y}} \\
S_{10} &= {\pi_{Q}^{yy}}\\
  &\quad - \frac{1}{6 {T} {\tau} {u^{\tau}}}{D_{s}}\Bigl(2 {I_{31}} \left(u^{\tau}\right)^2 \left(\left(u^{y}\right)^2+1\right) + 2 {\mathcal{B}} \left(\left(u^{y}\right)^2+1\right) ({\nu^{x}} {u^{x}}+{\nu^{y}} {u^{y}})\\
  &\quad + {\mathcal{C}} \left(u^{\tau}\right)^2 \left(2 {\Pi_{Q}} \left(\left(u^{y}\right)^2+1\right)-3 {\pi_{Q}^{yy}}+2 {\pi_{Q}{}^{\eta}{}_{\eta}} \left(\left(u^{y}\right)^2+1\right)\right)\Bigr) \\
S_{11} &= \frac{{D_{s}} \left(4 {\mathcal{B}} ({\nu^{x}} {u^{x}}+{\nu^{y}} {u^{y}})+{\mathcal{C}} \left(u^{\tau}\right)^2 (4 {\Pi_{Q}}+7 {\pi_{Q}{}^{\eta}{}_{\eta}})+4 {I_{31}} \left(u^{\tau}\right)^2\right)}{6 {T} {\tau} {u^{\tau}}}+{\pi_{Q}{}^{\eta}{}_{\eta}} \\
S_{12} &= -\frac{{D_{s}} {u^{x}} {u^{y}} \left({\mathcal{B}} ({\nu^{x}} {u^{x}}+{\nu^{y}} {u^{y}})+{I_{31}} \left(u^{\tau}\right)^2\right)}{3 {T} {\tau} {u^{\tau}}}+\frac{{\mathcal{C}} {D_{s}} {u^{\tau}} (3 {\pi_{Q}^{xy}}-2 {u^{x}} {u^{y}} ({\Pi_{Q}}+{\pi_{Q}{}^{\eta}{}_{\eta}}))}{6 {T} {\tau}}+{\pi_{Q}^{xy}} \\
S_{13} &= \frac{{D_{s}} \left(5 {\mathcal{B}} ({\nu^{x}} {u^{x}}+{\nu^{y}} {u^{y}})+{\mathcal{C}} \left(u^{\tau}\right)^2 (5 {\Pi_{Q}}+2 {\pi_{Q}{}^{\eta}{}_{\eta}})+5 {I_{31}} \left(u^{\tau}\right)^2\right)}{6 {T} {\tau} {u^{\tau}}}+{\Pi_{Q}}
\end{align*}


\end{document}
